\documentclass[14pt,a4paper,oneside]{extbook} % Dùng extbook để nhận cỡ chữ 14pt chuẩn
\usepackage[utf8]{vietnam}
\usepackage{amsmath,amssymb,mathrsfs,fancyhdr,enumerate,multirow,makecell,venndiagram}
\usepackage{tkz-euclide,tikz-3dplot,tikz,tkz-tab}
\usetikzlibrary{arrows,calc,intersections,angles,quotes,backgrounds,shapes.geometric,patterns,through}
\usepackage{yhmath}
\usepackage{pgfplots}
\pgfplotsset{compat=1.9}
\usepgfplotslibrary{fillbetween}
\usepackage[tikz]{bclogo}
\usepackage[top=1.8cm, bottom=1.8cm, left=1.8cm, right=1.8cm, headheight=16pt]{geometry}
\usepackage[hidelinks,unicode]{hyperref}
\usepackage{qrcode}

\usepackage[many]{tcolorbox}

% Gói ex_test (Chọn [dethi] để in đề hoặc [loigiai] để in cả lời giải)
%\usepackage[loigiai]{ex_test}
\usepackage[dethi]{ex_test}

\def\colorEX{\color{blue}}

\usepackage{fontawesome}
\usepackage{framed}
\usepackage{polynom}

\colorlet{tcbcol@back}{tcbcolback}
\colorlet{tcbcol@frame}{tcbcolframe}
\newtcolorbox[auto counter]{danggg}[1]{
	fonttitle=\bfseries\upshape,
	colframe=blue,colback=blue!10,
	sharp corners, breakable, 
	title={\color{white}\faFolderO \ DẠNG~\thetcbcounter.\ #1}
	\addcontentsline{toc}{subsection}{Dạng~\thetcbcounter. #1}
}
\newenvironment{dang}[1]{
	\begin{danggg}{#1}
		\textit{Phương pháp giải:~}
	}{\end{danggg}}
\newenvironment{tomtat}{
	\begin{tcolorbox}[boxrule=1pt,arc=0mm,breakable]
	}{\end{tcolorbox}}

\newtheorem{bt}{Bài}
\newtheorem{vd}{Ví dụ}
\theorembodyfont{\it}
\newtheorem{dl}{ Định lí}
\theoremseparator{.}
\theorembodyfont{\rm}
\newtheorem{hq}{ Hệ quả}
\newtheorem{tc}{ Tính chất}
\newtheorem{dn}{ Định nghĩa}
\theoremstyle{nonumberplain}
\newtheorem{kn}{ Khái niệm}
\theoremseparator{.}
\theorembodyfont{\it}
\usepackage{stackengine}
\usepackage{scalerel}

\newcommand\dangersign[1][2ex]{%
	\renewcommand\stacktype{L}%
	\scaleto{\stackon[1.3pt]{\color{red}$\triangle$}{\fontsize{4}{0}\selectfont !}}{#1}%
}
\renewcommand{\parallel}{\mathbin{\!\raisebox{0.25pt}/\mkern-2mu\raisebox{0.25pt}/\!}}
\renewcommand{\backsim}{\ \mathord{\raisebox{0.025\baselineskip}{\rotatebox{90}{\fontfamily{qag}\selectfont S}}}\ }
\theoremseparator{}
\newtheorem{note}{{\dangersign[3ex]}}
\newtheorem{nx}{Nhận xét.}
\newtcolorbox{luuy}[1][]{colback=yellow!20,title={\bfseries Chú ý},#1}

% Tiêu đề chương / bài
\setcounter{secnumdepth}{4}
\setcounter{tocdepth}{2}
\usepackage{titlesec}
\usepackage{titletoc}

\makeatletter
\renewcommand\thesection{\@arabic\c@section}
\makeatother

\titleformat{\section}{\centering\Large\bfseries\color{blue!80!black}}{Chủ đề~\thesection.}{1ex}{
	\setcounter{vd}{0}
	\setcounter{ex}{0}
	\setcounter{bt}{0}
	\setcounter{dl}{0}
	\setcounter{dn}{0}
	\setcounter{hq}{0}
	\setcounter{tc}{0}
	\MakeUppercase
}
\titleformat{\subsection}{\large\bfseries\color{blue!70!black}}{}{0ex}{}
\usepackage[final]{pdfpages}

% Cài đặt riêng
\def\hideboxEX{}
\renewcommand{\boxEX}[2][15pt]{%
	\setbox1=\hbox{#2\strut}
	\def\sizeboxEX{#1}%
	\ \EXbox{\hideboxEX#2\strut}\ %
}

\renewcommand\arraystretch{1.4} 
\renewcommand{\tabcolsep}{4mm}

% Dòng kẻ làm bài TikZ chuẩn 8mm
\renewcommand{\dotlineEX}[1]{%
  \par\vspace{2mm}%
  \noindent%
  \begin{tikzpicture}[x=1mm, y=1mm]%
    \foreach \i in {1,...,#1} {%
      \draw[gray!60, loosely dotted, line width=0.6pt] (0, -\i*8) -- (\linewidth, -\i*8);%
    }%
  \end{tikzpicture}%
}

% Nút Casio
\renewcommand\boxed[1]{%
	\tikz[baseline=(key.base)]
	\node[%
	draw,
	fill=black!40,
	drop shadow={shadow xshift=0.25ex,shadow yshift=-0.25ex,fill=black!90,opacity=0.75},
	rectangle,minimum width=0.6cm,minimum height=0.6cm,
	rounded corners=2pt,
	inner sep=1pt,
	line width=0.5pt,
	font=\sffamily
	](key) {\color{white}#1\strut}
	;
}

\newcommand{\hoac}[1]{\left[\begin{aligned}#1\end{aligned}\right.}
\newcommand{\heva}[1]{\left\{\begin{aligned}#1\end{aligned}\right.}

\pagestyle{fancy}
\fancyhf{} % Xóa sạch toàn bộ header/footer cũ
\fancyhead[L]{\small\itshape\color{gray} Tuyển tập Chuyên đề Vận dụng Toán 3 theo từng bài}
\fancyhead[R]{\small\bfseries\color{blue!80!black} Lớp Toán Thầy Lĩnh -- loptoanthaylinh.vn}
\fancyfoot[C]{\bfseries\thepage}
\renewcommand{\headrulewidth}{0.4pt}

\newcommand{\caidat}{2}
\newcommand{\indapan}[2]{
	\begin{center}
		\color{blue!50!green} \fbox{\fontfamily{qag}\bfseries\selectfont BẢNG ĐÁP ÁN}
	\end{center}
	{\fontfamily{lmtt}\bfseries\selectfont \inputansbox[\caidat]{#1}{#2}}
}

\everymath{\color{blue}}

% -------------------------------------------------------------------------
% KHUNG TIÊU ĐỀ PHIẾU KIỂM TRA KÈM MÃ QR CODE CHO TỪNG BÀI
% Cú pháp: \TaoPhieuKiemTra{Chủ đề}{Số bài}{Tên bài}{Link QR thi online}
% -------------------------------------------------------------------------
\newcommand{\TaoPhieuKiemTra}[4]{%
  \refstepcounter{subsection}%
  \addcontentsline{toc}{subsection}{Bài #2. #3}%
  \noindent
  \begin{tcolorbox}[
    colback=white,
    colframe=blue!85!black,
    arc=2.5mm,
    boxrule=1pt,
    left=2.5mm, right=2.5mm, top=1.8mm, bottom=1.8mm
  ]
    % LOGO TRUNG TÂM & HOTLINE | TIÊU ĐỀ BÀI KIỂM TRA | MÃ QR CODE
    \begin{minipage}[c]{0.27\textwidth}
      \centering
      \begin{tcolorbox}[colback=blue!8, colframe=blue!80!black, boxrule=0.6pt, arc=1.5mm, halign=center, top=1.2mm, bottom=1.2mm, left=1mm, right=1mm]
        {\bfseries\color{blue!95!black}\small LỚP TOÁN THẦY LĨNH}\\[1.5pt]
        {\scriptsize\color{blue!70!black}\faGlobe~loptoanthaylinh.vn}\\[1.5pt]
        {\scriptsize\bfseries\color{blue!85!black}\faPhone~Hotline: 0866976878}
      \end{tcolorbox}
    \end{minipage}%
    \hfill
    \begin{minipage}[c]{0.56\textwidth}
      \centering
      {\footnotesize\bfseries\color{blue!80!black}\MakeUppercase{#1}}\\[2pt]
      {\large\bfseries\color{red!85!black}\MakeUppercase{BÀI #2: #3}}\\[2pt]
      {\scriptsize\itshape CHUYÊN ĐỀ VẬN DỤNG -- Thời gian làm bài: 35 phút -- Cấu trúc 3 phần chuẩn}
    \end{minipage}%
    \hfill
    \begin{minipage}[c]{0.14\textwidth}
      \centering
      \qrcode[height=1.35cm]{#4}\\[2pt]
      {\tiny\bfseries\color{blue!90!black}\faQrcode~Quét QR thi}
    \end{minipage}
  \end{tcolorbox}
  \vspace{1.5mm}
}

\begin{document}

\tableofcontents
\newpage

% =========================================================================
% CHỦ ĐỀ 1 - ÔN TẬP VÀ BỔ SUNG
% =========================================================================
\section{Ôn tập và bổ sung}

\TaoPhieuKiemTra{Chủ đề 1 -- Ôn tập và bổ sung}{01}{Ôn tập các số trong phạm vi 1000}{https://loptoanthaylinh.vn/thi-online/toan-3/bai-01}
\textbf{Phần 1. Trắc nghiệm lựa chọn}
\setcounterpageref{ex}{0}

\begin{ex} %[3D1V01]
Số tự nhiên $X$ có ba chữ số. Biết chữ số hàng trăm là số chẵn lớn nhất có một chữ số, chữ số hàng chục kém chữ số hàng trăm $2$ đơn vị, chữ số hàng đơn vị bằng một nửa chữ số hàng chục. Cách đọc của số $X$ là
\choice
{Tám trăm sáu mươi}
{Tám trăm ba mươi sáu}
{\True Tám trăm sáu mươi ba}
{Sáu trăm tám mươi ba}
\loigiai{
Chữ số chẵn lớn nhất có một chữ số là $8$, do đó chữ số hàng trăm là $8$.\\
Chữ số hàng chục là: $8 - 2 = 6$.\\
Chữ số hàng đơn vị là: $6 : 2 = 3$.\\
Vậy số $X$ được viết là $863$ và đọc là: \textit{Tám trăm sáu mươi ba}.\\
\textbf{Đáp án của câu hỏi là C}
}
\end{ex}

\begin{ex} %[3D1V01]
Trong số $949$, giá trị của chữ số $9$ ở hàng trăm gấp giá trị của chữ số $9$ ở hàng đơn vị số lần là
\choice
{\True $100$ lần}
{$10$ lần}
{$90$ lần}
{$1000$ lần}
\loigiai{
Chữ số $9$ ở hàng trăm có giá trị là $900$.\\
Chữ số $9$ ở hàng đơn vị có giá trị là $9$.\\
Giá trị chữ số $9$ ở hàng trăm gấp giá trị chữ số $9$ ở hàng đơn vị số lần là:
$$900 : 9 = 100\text{ (lần)}$$
\textbf{Đáp án của câu hỏi là A}
}
\end{ex}

\begin{ex} %[3D1V01]
Số liền sau của số chẵn bé nhất có ba chữ số khác nhau là
\choice
{$100$}
{$102$}
{$101$}
{\True $103$}
\loigiai{
Để số chẵn có 3 chữ số khác nhau là bé nhất:\\
- Chữ số hàng trăm bé nhất khác 0 là $1$.\\
- Chữ số hàng chục bé nhất khác 1 là $0$.\\
- Chữ số hàng đơn vị là số chẵn bé nhất khác 0 là $2$.\\
Số chẵn bé nhất có 3 chữ số khác nhau là $102$.\\
Số liền sau của $102$ là: $102 + 1 = 103$.\\
\textbf{Đáp án của câu hỏi là D}
}
\end{ex}

\begin{ex} %[3D1V01]
Một số gồm $5$ trăm, $14$ chục và $6$ đơn vị. Số đó được viết là
\choice
{$5146$}
{\True $646$}
{$546$}
{$656$}
\loigiai{
Ta có:\\
- $5$ trăm $= 500$\\
- $14$ chục $= 140$\\
- $6$ đơn vị $= 6$\\
Số đó là: $500 + 140 + 6 = 646$.\\
\textbf{Đáp án của câu hỏi là B}
}
\end{ex}

\begin{ex} %[3D1V01]
Cho biểu thức $A = 600 + 30 + 8$ và $B = 600 + 50 + 2$. Số tròn chục $x$ thỏa mãn $A < x < B$ là
\choice
{$630$}
{$650$}
{\True $640$}
{$660$}
\loigiai{
Ta có: $A = 638$ và $B = 652$.\\
Điều kiện: $638 < x < 652$.\\
Vì $x$ là số tròn chục nên $x = 640$.\\
\textbf{Đáp án của câu hỏi là C}
}
\end{ex}

\begin{ex} %[3D1V01]
Số lẻ lớn nhất có ba chữ số khác nhau mà chữ số hàng trăm gấp $3$ lần chữ số hàng đơn vị là
\choice
{\True $983$}
{$981$}
{$973$}
{$389$}
\loigiai{
Vì số cần tìm là số lẻ nên chữ số hàng đơn vị phải là số lẻ: $1, 3, 5, 7, 9$.\\
Mặt khác, chữ số hàng trăm gấp 3 lần chữ số hàng đơn vị nên chữ số hàng đơn vị chỉ có thể là 1 (hàng trăm là 3) hoặc 3 (hàng trăm là 9).\\
Để số là lớn nhất, ta chọn chữ số hàng trăm là 9 và chữ số hàng đơn vị là 3.\\
Chữ số hàng chục lớn nhất khác 9 và 3 là 8.\\
Vậy số cần tìm là $983$.\\
\textbf{Đáp án của câu hỏi là A}
}
\end{ex}

\begin{ex} %[3D1V01]
Cho các giá trị sau:\\
$M = 700 + 80 + 5$;\\
$N$ là số liền trước của số $790$;\\
$P$ là số gồm $8$ trăm và $5$ đơn vị;\\
$Q = 700 + 95$.\\
Sắp xếp các số $M, N, P, Q$ theo thứ tự từ bé đến lớn là
\choice
{$M < Q < N < P$}
{$N < M < Q < P$}
{$M < N < P < Q$}
{\True $M < N < Q < P$}
\loigiai{
Tính giá trị từng số:\\
- $M = 785$\\
- $N = 790 - 1 = 789$\\
- $P = 805$\\
- $Q = 795$\\
Vì $785 < 789 < 795 < 805$ nên thứ tự từ bé đến lớn là: $M < N < Q < P$.\\
\textbf{Đáp án của câu hỏi là D}
}
\end{ex}

\begin{ex} %[3D1V01]
Một số có ba chữ số có chữ số hàng trăm là $5$. Nếu xóa bỏ chữ số $5$ này thì số đó giảm đi bao nhiêu đơn vị?
\choice
{$50$ đơn vị}
{\True $500$ đơn vị}
{$5$ đơn vị}
{$450$ đơn vị}
\loigiai{
Gọi số ban đầu có dạng $\overline{5ab} = 500 + \overline{ab}$.\\
Khi xóa chữ số 5 ở hàng trăm, ta được số mới là $\overline{ab}$.\\
Số đó giảm đi số đơn vị là:
$$(500 + \overline{ab}) - \overline{ab} = 500\text{ (đơn vị)}$$
\textbf{Đáp án của câu hỏi là B}
}
\end{ex}

\begin{ex} %[3D1V01]
Từ ba chữ số $0; 4; 7$, lập được bao nhiêu số chẵn có ba chữ số khác nhau?
\choice
{$4$ số}
{$2$ số}
{\True $3$ số}
{$1$ số}
\loigiai{
Một số là số chẵn khi chữ số hàng đơn vị là 0 hoặc 4.\\
- \textbf{Trường hợp 1:} Chữ số hàng đơn vị là 0. Các số lập được là: $470; 740$ (có 2 số).\\
- \textbf{Trường hợp 2:} Chữ số hàng đơn vị là 4. Chữ số hàng trăm không thể là 0 nên chỉ có thể là 7. Số lập được là: $704$ (có 1 số).\\
Tổng cộng có: $2 + 1 = 3$ số thỏa mãn.\\
\textbf{Đáp án của câu hỏi là C}
}
\end{ex}

\begin{ex} %[3D1V01]
Hiệu giữa số lớn nhất có 3 chữ số khác nhau và số bé nhất có 3 chữ số khác nhau được lập từ ba chữ số $3; 0; 8$ là
\choice
{\True $522$}
{$530$}
{$502$}
{$492$}
\loigiai{
- Số lớn nhất có 3 chữ số khác nhau lập từ $3; 0; 8$ là: $830$.\\
- Số bé nhất có 3 chữ số khác nhau lập từ $3; 0; 8$ là: $308$.\\
Hiệu hai số là: $830 - 308 = 522$.\\
\textbf{Đáp án của câu hỏi là A}
}
\end{ex}

\begin{ex} %[3D1V01]
Cho dãy số cách đều: $120; 125; 130; 135; \dots$ Số hạng thứ 10 của dãy số là
\choice
{$160$}
{$170$}
{$175$}
{\True $165$}
\loigiai{
Dãy số có khoảng cách giữa hai số liền nhau là: $125 - 120 = 5$ đơn vị.\\
Số hạng thứ 10 của dãy số là:
$$120 + (10 - 1) \times 5 = 120 + 45 = 165$$
\textbf{Đáp án của câu hỏi là D}
}
\end{ex}

\begin{ex} %[3D1V01]
Tìm số có ba chữ số, biết chữ số hàng đơn vị là $2$, chữ số hàng chục gấp $3$ lần chữ số hàng đơn vị và chữ số hàng trăm kém chữ số hàng chục $2$ đơn vị. Số đó là
\choice
{$642$}
{\True $462$}
{$426$}
{$264$}
\loigiai{
Chữ số hàng đơn vị là $2$.\\
Chữ số hàng chục là: $2 \times 3 = 6$.\\
Chữ số hàng trăm là: $6 - 2 = 4$.\\
Vậy số cần tìm là $462$.\\
\textbf{Đáp án của câu hỏi là B}
}
\end{ex}

\textbf{Phần 2. Trắc nghiệm đúng sai}
\setcounterpageref{ex}{0}

\begin{ex} %[3D1V01]
Xét tính đúng sai của các nhận định sau về các số trong phạm vi 1000:
\choiceTF
{\True Hiệu giữa số lớn nhất có ba chữ số và số bé nhất có ba chữ số là $899$.}
{Số liền trước của một số có ba chữ số luôn luôn là một số có ba chữ số.}
{\True Trong số $444$, giá trị của chữ số $4$ ở hàng trăm gấp $100$ lần giá trị của chữ số $4$ ở hàng đơn vị.}
{\True Có tất cả $90$ số tròn chục có ba chữ số.}
\loigiai{
\begin{itemize}
    \item Ý a) Đúng: Số lớn nhất có 3 chữ số là 999, số bé nhất có 3 chữ số là 100. Hiệu là $999 - 100 = 899$.
    \item Ý b) Sai: Số liền trước của số 100 là 99 (là số có 2 chữ số).
    \item Ý c) Đúng: Chữ số 4 hàng trăm có giá trị 400, chữ số 4 hàng đơn vị có giá trị 4. Ta có $400 : 4 = 100$.
    \item Ý d) Đúng: Các số tròn chục có 3 chữ số từ 100 đến 990. Số lượng số là: $(990 - 100) : 10 + 1 = 90$ số.
\end{itemize}
\begin{itemize}
    \item a) \textbf{Đúng}.
    \item b) \textbf{Sai}.
    \item c) \textbf{Đúng}.
    \item d) \textbf{Đúng}.
\end{itemize}
}
\end{ex}

\begin{ex} %[3D1V01]
Cho tập hợp ba chữ số: $2; 5; 0$. Xét tính đúng sai của các khẳng định sau:
\choiceTF
{\True Từ ba chữ số đã cho lập được tất cả 4 số có ba chữ số khác nhau.}
{Trong các số có ba chữ số khác nhau lập được, chỉ có đúng 2 số chia hết cho 5.}
{\True Số lẻ duy nhất có ba chữ số khác nhau lập được từ ba chữ số trên là $205$.}
{\True Tổng của số lớn nhất và số bé nhất có ba chữ số khác nhau lập được bằng $725$.}
\loigiai{
\begin{itemize}
    \item Ý a) Đúng: Các số lập được là $205; 250; 502; 520$ (4 số).
    \item Ý b) Sai: Các số có tận cùng là 0 hoặc 5 gồm $205; 250; 520$ (có 3 số chia hết cho 5).
    \item Ý c) Đúng: Chỉ có số $205$ có tận cùng là 5 (số lẻ).
    \item Ý d) Đúng: Số lớn nhất là 520, số bé nhất là 205. Tổng là: $520 + 205 = 725$.
\end{itemize}
\begin{itemize}
    \item a) \textbf{Đúng}.
    \item b) \textbf{Sai}.
    \item c) \textbf{Đúng}.
    \item d) \textbf{Đúng}.
\end{itemize}
}
\end{ex}

\textbf{Phần 3. Trả lời ngắn}
\setcounterpageref{ex}{0}

\begin{ex} %[3D1V01]
Tìm số tự nhiên có ba chữ số, biết tổng ba chữ số bằng $4$ và chữ số hàng trăm lớn hơn chữ số hàng chục, chữ số hàng chục lớn hơn chữ số hàng đơn vị.
\shortans{310}
\loigiai{
\textbf{Đáp số:} $310$\\
Gọi số cần tìm là $\overline{abc}$ với $a > b > c \ge 0$ và $a + b + c = 4$.\\
Vì $a > b > c$ nên $c$ bé nhất là 0. Khi $c = 0$, ta có $a + b = 4$ với $a > b > 0$. Ta chọn được $a = 3, b = 1$.\\
Vậy số cần tìm là $310$.
}
\end{ex}

\begin{ex} %[3D1V01]
Dãy số chẵn liên tiếp $102; 104; 106; \dots; 180$ có tất cả bao nhiêu số hạng?
\shortans{40}
\loigiai{
\textbf{Đáp số:} $40$\\
Số số hạng của dãy là: $(180 - 102) : 2 + 1 = 78 : 2 + 1 = 40$ (số hạng).
}
\end{ex}

\begin{ex} %[3D1V01]
Để đánh số trang của một cuốn sách từ trang $1$ đến trang $100$, người ta phải dùng tất cả bao nhiêu chữ số?
\shortans{192}
\loigiai{
\textbf{Đáp số:} $192$\\
- Từ trang 1 đến trang 9 có 9 trang có 1 chữ số, cần: $9 \times 1 = 9$ (chữ số).\\
- Từ trang 10 đến trang 99 có: $(99 - 10) + 1 = 90$ trang có 2 chữ số, cần: $90 \times 2 = 180$ (chữ số).\\
- Trang 100 là trang có 3 chữ số, cần: $3$ (chữ số).\\
Tổng số chữ số cần dùng là: $9 + 180 + 3 = 192$ (chữ số).
}
\end{ex}

\begin{ex} %[3D1V01]
Tìm số lớn nhất có ba chữ số khác nhau mà tổng ba chữ số bằng $14$ và chữ số hàng trăm là $8$.
\shortans{860}
\loigiai{
\textbf{Đáp số:} $860$\\
Gọi số cần tìm có dạng $\overline{8ab}$ với $a \ne b \ne 8$.\\
Tổng hai chữ số hàng chục và hàng đơn vị là: $14 - 8 = 6$.\\
Để số là lớn nhất thì chữ số hàng chục $a$ phải lớn nhất có thể ($a \le 6$). Chọn $a = 6$, khi đó chữ số hàng đơn vị là $b = 6 - 6 = 0$.\\
Vậy số lớn nhất cần tìm là $860$.
}
\end{ex}

\begin{ex} %[3D1V01]
Cho ba chữ số $1; 5; 9$. Tính tổng của số lớn nhất và số bé nhất có ba chữ số khác nhau được lập từ ba chữ số đã cho.
\shortans{1110}
\loigiai{
\textbf{Đáp số:} $1110$\\
- Số lớn nhất có ba chữ số khác nhau lập từ $1; 5; 9$ là: $951$.\\
- Số bé nhất có ba chữ số khác nhau lập từ $1; 5; 9$ là: $159$.\\
Tổng hai số là: $951 + 159 = 1110$.
}
\end{ex}

\begin{ex} %[3D1V01]
Có bao nhiêu số có ba chữ số mà chữ số hàng trăm bằng chữ số hàng đơn vị?
\shortans{90}
\loigiai{
\textbf{Đáp số:} $90$\\
Số có dạng $\overline{aba}$ ($a \ne 0$):\\
- Chữ số hàng trăm $a$ có 9 cách chọn ($a \in \{1; 2; \dots; 9\}$).\\
- Chữ số hàng chục $b$ có 10 cách chọn ($b \in \{0; 1; 2; \dots; 9\}$).\\
- Chữ số hàng đơn vị phải bằng $a$ nên có đúng 1 cách chọn.\\
Số các số thỏa mãn là: $9 \times 10 \times 1 = 90$ (số).
}
\end{ex}

\newpage

\TaoPhieuKiemTra{Chủ đề 1 -- Ôn tập và bổ sung}{02}{Ôn tập phép cộng, phép trừ trong phạm vi 1000}{https://loptoanthaylinh.vn/thi-online/toan-3/bai-02}
\textbf{Phần 1. Trắc nghiệm lựa chọn}
\setcounterpageref{ex}{0}

\begin{ex} %[3D1V02]
Tính nhanh giá trị của biểu thức: $125 + 348 + 275 + 152$.
\choice
{$800$}
{\True $900$}
{$850$}
{$1000$}
\loigiai{
Sử dụng tính chất giao hoán và kết hợp của phép cộng để nhóm các số tròn trăm:\\
$$(125 + 275) + (348 + 152) = 400 + 500 = 900$$
\textbf{Đáp án của câu hỏi là B}
}
\end{ex}

\begin{ex} %[3D1V02]
Trong một phép cộng, nếu số hạng thứ nhất tăng thêm $35$ đơn vị và số hạng thứ hai giảm đi $18$ đơn vị thì tổng mới thay đổi như thế nào so với tổng ban đầu?
\choice
{Giảm đi $17$ đơn vị}
{Tăng thêm $53$ đơn vị}
{Giảm đi $53$ đơn vị}
{\True Tăng thêm $17$ đơn vị}
\loigiai{
Khi số hạng thứ nhất tăng thêm $35$ đơn vị thì tổng tăng $35$ đơn vị.\\
Khi số hạng thứ hai giảm đi $18$ đơn vị thì tổng giảm $18$ đơn vị.\\
Vì $35 > 18$ nên tổng mới tăng thêm số đơn vị là:
$$35 - 18 = 17\text{ (đơn vị)}$$
\textbf{Đáp án của câu hỏi là D}
}
\end{ex}

\begin{ex} %[3D1V02]
Trong một phép trừ có hiệu bằng $142$. Biết tổng của số bị trừ, số trừ và hiệu bằng $560$. Số bị trừ của phép trừ đó là
\choice
{\True $280$}
{$140$}
{$418$}
{$210$}
\loigiai{
Ta có hệ thức: $\text{Số bị trừ} = \text{Số trừ} + \text{Hiệu}$.\\
Do đó:\\
$$\text{Số bị trừ} + \text{Số trừ} + \text{Hiệu} = \text{Số bị trừ} + \text{Số bị trừ} = 2 \times \text{Số bị trừ}$$
Số bị trừ là:
$$560 : 2 = 280$$
\textbf{Đáp án của câu hỏi là A}
}
\end{ex}

\begin{ex} %[3D1V02]
Cho phép cộng đặt tính: $\overline{4a6} + \overline{25b} = 743$. Giá trị của biểu thức $a + b$ là
\choice
{$13$}
{$14$}
{\True $15$}
{$16$}
\loigiai{
Xét phép cộng theo từng cột từ phải sang trái:\\
- Hàng đơn vị: $6 + b$ có tận cùng là $3$, do $6 + 7 = 13$ nên $b = 7$ (nhớ $1$ sang hàng chục).\\
- Hàng chục: $a + 5 + 1 = 14$ nên $a + 6 = 14 \Rightarrow a = 8$ (nhớ $1$ sang hàng trăm).\\
- Hàng trăm: $4 + 2 + 1 = 7$ (thỏa mãn).\\
Vậy $a = 8$ và $b = 7$. Ta có: $a + b = 8 + 7 = 15$.\\
\textbf{Đáp án của câu hỏi là C}
}
\end{ex}

\begin{ex} %[3D1V02]
Một phép trừ có hiệu bằng $240$. Nếu số bị trừ tăng thêm $40$ đơn vị và số trừ tăng thêm $15$ đơn vị thì hiệu mới là
\choice
{$295$}
{\True $265$}
{$215$}
{$185$}
\loigiai{
Khi số bị trừ tăng thêm $40$ đơn vị thì hiệu tăng $40$ đơn vị.\\
Khi số trừ tăng thêm $15$ đơn vị thì hiệu giảm đi $15$ đơn vị.\\
Hiệu mới là:
$$240 + 40 - 15 = 265$$
\textbf{Đáp án của câu hỏi là B}
}
\end{ex}

\begin{ex} %[3D1V02]
Hiệu giữa số lớn nhất có 3 chữ số khác nhau và số bé nhất có 3 chữ số khác nhau là
\choice
{$880$}
{$875$}
{$890$}
{\True $885$}
\loigiai{
- Số lớn nhất có 3 chữ số khác nhau là: $987$.\\
- Số bé nhất có 3 chữ số khác nhau là: $102$.\\
Hiệu giữa hai số là:
$$987 - 102 = 885$$
\textbf{Đáp án của câu hỏi là D}
}
\end{ex}

\begin{ex} %[3D1V02]
Thư viện trường có ba tủ sách. Tủ thứ nhất có $245$ cuốn sách. Tủ thứ hai có nhiều hơn tủ thứ nhất $35$ cuốn sách nhưng lại ít hơn tủ thứ ba $20$ cuốn sách. Cả ba tủ sách có tất cả số cuốn sách là
\choice
{\True $825$ cuốn sách}
{$805$ cuốn sách}
{$845$ cuốn sách}
{$795$ cuốn sách}
\loigiai{
Số sách ở tủ thứ hai là:
$$245 + 35 = 280\text{ (cuốn)}$$
Số sách ở tủ thứ ba là:
$$280 + 20 = 300\text{ (cuốn)}$$
Cả ba tủ có tất cả số cuốn sách là:
$$245 + 280 + 300 = 825\text{ (cuốn)}$$
\textbf{Đáp án của câu hỏi là A}
}
\end{ex}

\begin{ex} %[3D1V02]
Tính giá trị của biểu thức: $P = 782 - (245 + 182)$.
\choice
{$345$}
{$365$}
{\True $355$}
{$335$}
\loigiai{
Bỏ dấu ngoặc hoặc trừ liên tiếp để tính nhanh:\\
$$P = 782 - 182 - 245 = 600 - 245 = 355$$
\textbf{Đáp án của câu hỏi là C}
}
\end{ex}

\begin{ex} %[3D1V02]
Hai số có tổng bằng $640$. Biết số bé là một số có hai chữ số và khi viết thêm chữ số $5$ vào bên trái số bé ta được số lớn. Số lớn là
\choice
{$550$}
{\True $570$}
{$560$}
{$580$}
\loigiai{
Khi viết thêm chữ số $5$ vào bên trái một số có hai chữ số thì số đó tăng thêm $500$ đơn vị.\\
Do đó, số lớn hơn số bé $500$ đơn vị.\\
Số bé là:
$$(640 - 500) : 2 = 70$$
Số lớn là:
$$70 + 500 = 570$$
\textbf{Đáp án của câu hỏi là B}
}
\end{ex}

\begin{ex} %[3D1V02]
Tìm chữ số $x$ thỏa mãn phép trừ: $\overline{7x4} - 289 = 465$.
\choice
{$3$}
{$4$}
{$6$}
{\True $5$}
\loigiai{
Ta có:
$$\overline{7x4} = 465 + 289 = 754$$
So sánh hai vế, ta được chữ số $x = 5$.\\
\textbf{Đáp án của câu hỏi là D}
}
\end{ex}

\begin{ex} %[3D1V02]
Một cửa hàng ngày thứ nhất bán được $318\text{ m}$ vải. Ngày thứ hai bán được ít hơn ngày thứ nhất $45\text{ m}$ vải. Ngày thứ ba bán được ít hơn ngày thứ hai $20\text{ m}$ vải. Cả ba ngày cửa hàng bán được số mét vải là
\choice
{\True $844\text{ m}$}
{$824\text{ m}$}
{$864\text{ m}$}
{$834\text{ m}$}
\loigiai{
Ngày thứ hai bán được số mét vải là:
$$318 - 45 = 273\text{ (m)}$$
Ngày thứ ba bán được số mét vải là:
$$273 - 20 = 253\text{ (m)}$$
Cả ba ngày bán được số mét vải là:
$$318 + 273 + 253 = 844\text{ (m)}$$
\textbf{Đáp án của câu hỏi là A}
}
\end{ex}

\begin{ex} %[3D1V02]
Từ bốn chữ số $1; 3; 6; 8$, lập số lớn nhất có 3 chữ số khác nhau và số bé nhất có 3 chữ số khác nhau. Hiệu của hai số lập được là
\choice
{$717$}
{$737$}
{\True $727$}
{$729$}
\loigiai{
- Số lớn nhất có 3 chữ số khác nhau từ 4 chữ số trên là: $863$.\\
- Số bé nhất có 3 chữ số khác nhau từ 4 chữ số trên là: $136$.\\
Hiệu của hai số là:
$$863 - 136 = 727$$
\textbf{Đáp án của câu hỏi là C}
}
\end{ex}

\textbf{Phần 2. Trắc nghiệm đúng sai}
\setcounterpageref{ex}{0}

\begin{ex} %[3D1V02]
Xét tính đúng sai của các nhận định sau về phép cộng và phép trừ trong phạm vi 1000:
\choiceTF
{\True Khi thêm vào số hạng này bao nhiêu đơn vị và bớt ở số hạng kia bấy nhiêu đơn vị thì tổng không thay đổi.}
{Trong một phép trừ, số bị trừ luôn luôn lớn hơn số trừ.}
{\True Nếu $a + b = 500$ thì $(a + 50) + (b + 50) = 600$.}
{\True Hiệu giữa số chẵn lớn nhất có 3 chữ số và số chẵn bé nhất có 3 chữ số bằng $898$.}
\loigiai{
\begin{itemize}
    \item Ý a) Đúng: Ví dụ $(a + k) + (b - k) = a + b + k - k = a + b$.
    \item Ý b) Sai: Số bị trừ có thể bằng số trừ (khi đó hiệu bằng 0, ví dụ $500 - 500 = 0$).
    \item Ý c) Đúng: Cả hai số hạng đều tăng 50 đơn vị nên tổng tăng: $50 + 50 = 100$. Tổng mới là $500 + 100 = 600$.
    \item Ý d) Đúng: Số chẵn lớn nhất có 3 chữ số là 998, số chẵn bé nhất có 3 chữ số là 100. Hiệu là $998 - 100 = 898$.
\end{itemize}
\begin{itemize}
    \item a) \textbf{Đúng}.
    \item b) \textbf{Sai}.
    \item c) \textbf{Đúng}.
    \item d) \textbf{Đúng}.
\end{itemize}
}
\end{ex}

\begin{ex} %[3D1V02]
Cho ba số $A = 345, B = 456, C = 199$. Xét tính đúng sai của các khẳng định sau:
\choiceTF
{\True Tổng $A + B + C$ là một số tròn nghìn.}
{Hiệu $(B - A)$ lớn hơn hiệu $(A - C)$.}
{\True Nếu bớt ở cả ba số mỗi số $10$ đơn vị thì tổng mới giảm đi $30$ đơn vị so với tổng ban đầu.}
{Số $B$ lớn hơn tổng của hai số $A$ và $C$.}
\loigiai{
\begin{itemize}
    \item Ý a) Đúng: $A + B + C = 345 + 456 + 199 = 1000$ (là số tròn nghìn).
    \item Ý b) Sai: $B - A = 456 - 345 = 111$; $A - C = 345 - 199 = 146$. Vì $111 < 146$ nên $B - A < A - C$.
    \item Ý c) Đúng: Cả ba số cùng giảm 10 đơn vị nên tổng giảm $10 \times 3 = 30$ đơn vị.
    \item Ý d) Sai: $A + C = 345 + 199 = 544$. Vì $456 < 544$ nên $B$ bé hơn tổng của $A$ và $C$.
\end{itemize}
\begin{itemize}
    \item a) \textbf{Đúng}.
    \item b) \textbf{Sai}.
    \item c) \textbf{Đúng}.
    \item d) \textbf{Đúng}.
\end{itemize}
}
\end{ex}

\textbf{Phần 3. Trả lời ngắn}
\setcounterpageref{ex}{0}

\begin{ex} %[3D1V02]
Tính nhanh giá trị của biểu thức: $234 + 187 + 166 + 213$.
\shortans{800}
\loigiai{
\textbf{Đáp số:} $800$\\
Nhóm các số tròn trăm:\\
$$(234 + 166) + (187 + 213) = 400 + 400 = 800$$
}
\end{ex}

\begin{ex} %[3D1V02]
Tìm một số tự nhiên, biết rằng nếu lấy số đó cộng với $125$ rồi trừ đi $85$ thì được kết quả là $400$.
\shortans{360}
\loigiai{
\textbf{Đáp số:} $360$\\
Gọi số cần tìm là $x$. Theo bài ra ta có:\\
$$x + 125 - 85 = 400$$
$$x + 40 = 400$$
$$x = 400 - 40 = 360$$
}
\end{ex}

\begin{ex} %[3D1V02]
Hai thùng chứa tất cả $450\text{ lít}$ dầu. Nếu người ta chuyển $30\text{ lít}$ dầu từ thùng thứ nhất sang thùng thứ hai thì lúc này tổng số dầu ở cả hai thùng là bao nhiêu lít?
\shortans{450}
\loigiai{
\textbf{Đáp số:} $450$\\
Khi chuyển dầu qua lại giữa hai thùng thì tổng số lít dầu ở cả hai thùng không thay đổi, vẫn là $450\text{ lít}$.
}
\end{ex}

\begin{ex} %[3D1V02]
Một cuộn dây điện dài $240\text{ m}$. Lần thứ nhất người thợ cắt đi $65\text{ m}$ dây. Lần thứ hai người đó cắt nhiều hơn lần thứ nhất $15\text{ m}$ dây. Hỏi sau hai lần cắt, cuộn dây còn lại bao nhiêu mét?
\shortans{95}
\loigiai{
\textbf{Đáp số:} $95$\\
Lần thứ hai cắt số mét dây điện là:
$$65 + 15 = 80\text{ (m)}$$
Cả hai lần cắt tất cả số mét dây điện là:
$$65 + 80 = 145\text{ (m)}$$
Sau hai lần cắt, cuộn dây còn lại số mét là:
$$240 - 145 = 95\text{ (m)}$$
}
\end{ex}

\begin{ex} %[3D1V02]
Tìm một số có 3 chữ số có chữ số hàng trăm là $3$. Biết rằng nếu xóa bỏ chữ số $3$ này đi thì số đó giảm đi bao nhiêu đơn vị?
\shortans{300}
\loigiai{
\textbf{Đáp số:} $300$\\
Số có dạng $\overline{3ab} = 300 + \overline{ab}$.\\
Khi xóa chữ số 3 ở hàng trăm thì được số mới là $\overline{ab}$.\\
Số đó giảm đi: $(300 + \overline{ab}) - \overline{ab} = 300$ đơn vị.
}
\end{ex}

\begin{ex} %[3D1V02]
Trong một phép trừ, biết tổng của số bị trừ, số trừ và hiệu bằng $720$. Tìm số bị trừ của phép trừ đó.
\shortans{360}
\loigiai{
\textbf{Đáp số:} $360$\\
Vì $\text{Số bị trừ} = \text{Số trừ} + \text{Hiệu}$ nên:\\
$$\text{Số bị trừ} + \text{Số trừ} + \text{Hiệu} = 2 \times \text{Số bị trừ} = 720$$
Số bị trừ là:
$$720 : 2 = 360$$
}
\end{ex}

\newpage

\TaoPhieuKiemTra{Chủ đề 1 -- Ôn tập và bổ sung}{03}{Tìm thành phần trong phép cộng, phép trừ}{https://loptoanthaylinh.vn/thi-online/toan-3/bai-03}
\textbf{Phần 1. Trắc nghiệm lựa chọn}
\setcounterpageref{ex}{0}

\begin{ex} %[3D1V03]
Tìm giá trị của $x$, biết: $x + 145 = 320 + 215$.
\choice
{\True $390$}
{$380$}
{$400$}
{$410$}
\loigiai{
Tính giá trị của vế phải trước:\\
$$x + 145 = 535$$
$$x = 535 - 145 = 390$$
\textbf{Đáp án của câu hỏi là A}
}
\end{ex}

\begin{ex} %[3D1V03]
Tìm số bị trừ $x$ trong biểu thức: $x - (120 + 80) = 450$.
\choice
{$550$}
{$600$}
{\True $650$}
{$700$}
\loigiai{
Thực hiện phép tính trong ngoặc:\\
$$x - 200 = 450$$
$$x = 450 + 200 = 650$$
\textbf{Đáp án của câu hỏi là C}
}
\end{ex}

\begin{ex} %[3D1V03]
Tìm giá trị của $x$, biết: $850 - x = 240 + 160$.
\choice
{$400$}
{\True $450$}
{$500$}
{$350$}
\loigiai{
Tính vế phải:\\
$$850 - x = 400$$
$$x = 850 - 400 = 450$$
\textbf{Đáp án của câu hỏi là B}
}
\end{ex}

\begin{ex} %[3D1V03]
Tổng của hai số bằng $480$. Nếu giữ nguyên số hạng thứ nhất và tăng số hạng thứ hai thêm $45$ đơn vị thì tổng mới là bao nhiêu?
\choice
{$515$}
{$535$}
{$435$}
{\True $525$}
\loigiai{
Khi giữ nguyên một số hạng và tăng số hạng kia thêm bao nhiêu đơn vị thì tổng cũng tăng thêm bấy nhiêu đơn vị.\\
Tổng mới là:
$$480 + 45 = 525$$
\textbf{Đáp án của câu hỏi là D}
}
\end{ex}

\begin{ex} %[3D1V03]
Hai số có tổng bằng $750$. Biết rằng nếu bớt ở số thứ nhất đi $50$ đơn vị thì hai số bằng nhau. Số thứ nhất ban đầu là
\choice
{\True $400$}
{$350$}
{$450$}
{$500$}
\loigiai{
Vì bớt ở số thứ nhất đi $50$ đơn vị thì hai số bằng nhau nên số thứ nhất hơn số thứ hai $50$ đơn vị.\\
Hai lần số thứ hai là:
$$750 - 50 = 700$$
Số thứ hai là:
$$700 : 2 = 350$$
Số thứ nhất ban đầu là:
$$350 + 50 = 400$$
\textbf{Đáp án của câu hỏi là A}
}
\end{ex}

\begin{ex} %[3D1V03]
Một phép trừ có số bị trừ bằng $540$. Nếu tăng số trừ thêm $30$ đơn vị thì được hiệu mới bằng $180$. Số trừ ban đầu là
\choice
{$300$}
{$360$}
{\True $330$}
{$320$}
\loigiai{
Khi số trừ tăng thêm $30$ đơn vị thì hiệu giảm đi $30$ đơn vị.\\
Hiệu ban đầu của phép trừ là:
$$180 + 30 = 210$$
Số trừ ban đầu là:
$$540 - 210 = 330$$
\textbf{Đáp án của câu hỏi là C}
}
\end{ex}

\begin{ex} %[3D1V03]
Tìm giá trị của $x$, biết: $900 - x - 150 = 420$.
\choice
{$320$}
{\True $330$}
{$340$}
{$350$}
\loigiai{
Biến đổi biểu thức:\\
$$900 - 150 - x = 420$$
$$750 - x = 420$$
$$x = 750 - 420 = 330$$
\textbf{Đáp án của câu hỏi là B}
}
\end{ex}

\begin{ex} %[3D1V03]
Trong một phép trừ, số trừ là $135$ và hiệu lớn hơn số trừ $40$ đơn vị. Số bị trừ là
\choice
{$300$}
{$290$}
{$320$}
{\True $310$}
\loigiai{
Hiệu của phép trừ là:
$$135 + 40 = 175$$
Số bị trừ là:
$$175 + 135 = 310$$
\textbf{Đáp án của câu hỏi là D}
}
\end{ex}

\begin{ex} %[3D1V03]
Bác An có hai bao thóc nặng tất cả $180\text{ kg}$. Nếu bác An chuyển từ bao thứ nhất sang bao thứ hai $15\text{ kg}$ thì hai bao nặng bằng nhau. Bao thứ nhất lúc đầu nặng số ki-lô-gam là
\choice
{\True $105\text{ kg}$}
{$95\text{ kg}$}
{$100\text{ kg}$}
{$110\text{ kg}$}
\loigiai{
Khi hai bao nặng bằng nhau thì mỗi bao nặng số ki-lô-gam là:
$$180 : 2 = 90\text{ (kg)}$$
Bao thứ nhất lúc đầu nặng số ki-lô-gam là:
$$90 + 15 = 105\text{ (kg)}$$
\textbf{Đáp án của câu hỏi là A}
}
\end{ex}

\begin{ex} %[3D1V03]
Tìm chữ số $x$ biết: $\overline{x45} + 230 = 875$.
\choice
{$5$}
{$7$}
{\True $6$}
{$8$}
\loigiai{
Ta có:
$$\overline{x45} = 875 - 230 = 645$$
Do đó chữ số $x = 6$.\\
\textbf{Đáp án của câu hỏi là C}
}
\end{ex}

\begin{ex} %[3D1V03]
Một kho có một số ki-lô-gam gạo. Sau khi xuất đi $160\text{ kg}$ gạo rồi lại nhập thêm $90\text{ kg}$ gạo thì trong kho có tất cả $350\text{ kg}$ gạo. Lúc đầu trong kho có số ki-lô-gam gạo là
\choice
{$400\text{ kg}$}
{\True $420\text{ kg}$}
{$440\text{ kg}$}
{$380\text{ kg}$}
\loigiai{
Gọi số gạo ban đầu trong kho là $x\text{ (kg)}$. Theo đề bài ta có:\\
$$x - 160 + 90 = 350$$
$$x - 70 = 350$$
$$x = 350 + 70 = 420\text{ (kg)}$$
\textbf{Đáp án của câu hỏi là B}
}
\end{ex}

\begin{ex} %[3D1V03]
Tìm giá trị của $x$, biết: $(x + 210) - 150 = 380 + 120$.
\choice
{$420$}
{$430$}
{$450$}
{\True $440$}
\loigiai{
Tính vế phải:\\
$$(x + 210) - 150 = 500$$
$$x + 210 = 500 + 150$$
$$x + 210 = 650$$
$$x = 650 - 210 = 440$$
\textbf{Đáp án của câu hỏi là D}
}
\end{ex}

\textbf{Phần 2. Trắc nghiệm đúng sai}
\setcounterpageref{ex}{0}

\begin{ex} %[3D1V03]
Xét tính đúng sai của các nhận định sau về các quy tắc tìm thành phần phép tính:
\choiceTF
{\True Muốn tìm một số hạng chưa biết trong một tổng, ta lấy tổng trừ đi số hạng đã biết.}
{Muốn tìm số trừ, ta lấy số bị trừ cộng với hiệu.}
{\True Nếu $x - 50 = 200$ thì giá trị của $x$ là $250$.}
{\True Trong một phép trừ, nếu hiệu bằng $0$ thì số bị trừ bằng số trừ.}
\loigiai{
\begin{itemize}
    \item Ý a) Đúng: Theo quy tắc tìm số hạng chưa biết.
    \item Ý b) Sai: Muốn tìm số trừ, ta phải lấy số bị trừ trừ đi hiệu ($\text{Số trừ} = \text{Số bị trừ} - \text{Hiệu}$).
    \item Ý c) Đúng: $x = 200 + 50 = 250$.
    \item Ý d) Đúng: Khi $\text{Số bị trừ} - \text{Số trừ} = 0 \Rightarrow \text{Số bị trừ} = \text{Số trừ}$.
\end{itemize}
\begin{itemize}
    \item a) \textbf{Đúng}.
    \item b) \textbf{Sai}.
    \item c) \textbf{Đúng}.
    \item d) \textbf{Đúng}.
\end{itemize}
}
\end{ex}

\begin{ex} %[3D1V03]
Cho bài toán tìm $x$: $(500 - x) + 120 = 320$. Xét tính đúng sai của các khẳng định sau:
\choiceTF
{\True Giá trị của biểu thức $(500 - x)$ là $200$.}
{Giá trị của $x$ tìm được là $700$.}
{\True Nếu vế phải tăng thêm $100$ đơn vị (thành $420$) thì giá trị của $x$ sẽ giảm đi $100$ đơn vị.}
{\True Số $x$ tìm được là một số tròn trăm lớn hơn $250$.}
\loigiai{
\begin{itemize}
    \item Ý a) Đúng: $500 - x = 320 - 120 = 200$.
    \item Ý b) Sai: $x = 500 - 200 = 300$ (không phải 700).
    \item Ý c) Đúng: Khi vế phải là 420 thì $500 - x = 420 - 120 = 300 \Rightarrow x = 500 - 300 = 200$. Giá trị $x$ giảm từ $300$ xuống $200$ (giảm 100 đơn vị).
    \item Ý d) Đúng: $x = 300$ là số tròn trăm và $300 > 250$.
\end{itemize}
\begin{itemize}
    \item a) \textbf{Đúng}.
    \item b) \textbf{Sai}.
    \item c) \textbf{Đúng}.
    \item d) \textbf{Đúng}.
\end{itemize}
}
\end{ex}

\textbf{Phần 3. Trả lời ngắn}
\setcounterpageref{ex}{0}

\begin{ex} %[3D1V03]
Tìm giá trị của $x$, biết: $x + 256 = 600 - 144$.
\shortans{200}
\loigiai{
\textbf{Đáp số:} $200$\\
Tính vế phải:\\
$$x + 256 = 456$$
$$x = 456 - 256 = 200$$
}
\end{ex}

\begin{ex} %[3D1V03]
Tìm một số tự nhiên, biết rằng nếu lấy số đó trừ đi $170$ rồi cộng thêm $90$ thì được kết quả là $420$.
\shortans{500}
\loigiai{
\textbf{Đáp số:} $500$\\
Gọi số cần tìm là $x$. Ta có:\\
$$x - 170 + 90 = 420$$
$$x - 80 = 420$$
$$x = 420 + 80 = 500$$
}
\end{ex}

\begin{ex} %[3D1V03]
Hai số có hiệu bằng $120$. Nếu thêm vào số bị trừ $25$ đơn vị và giữ nguyên số trừ thì hiệu mới bằng bao nhiêu?
\shortans{145}
\loigiai{
\textbf{Đáp số:} $145$\\
Khi số bị trừ tăng thêm bao nhiêu đơn vị và giữ nguyên số trừ thì hiệu cũng tăng thêm bấy nhiêu đơn vị.\\
Hiệu mới là:
$$120 + 25 = 145$$
}
\end{ex}

\begin{ex} %[3D1V03]
Một bác nông dân mang một số tiền đi chợ. Bác mua một bao phân bón hết $65$ nghìn đồng, mua giống cây trồng hết $45$ nghìn đồng thì bác còn lại đúng $90$ nghìn đồng. Hỏi ban đầu bác nông dân có tất cả bao nhiêu nghìn đồng?
\shortans{200}
\loigiai{
\textbf{Đáp số:} $200$\\
Số tiền bác nông dân đã tiêu là:
$$65 + 45 = 110\text{ (nghìn đồng)}$$
Ban đầu bác nông dân có số tiền là:
$$110 + 90 = 200\text{ (nghìn đồng)}$$
}
\end{ex}

\begin{ex} %[3D1V03]
Tìm giá trị của $x$, biết: $780 - (x + 130) = 250$.
\shortans{400}
\loigiai{
\textbf{Đáp số:} $400$\\
Xem $(x + 130)$ là số trừ chưa biết:\\
$$x + 130 = 780 - 250$$
$$x + 130 = 530$$
$$x = 530 - 130 = 400$$
}
\end{ex}

\begin{ex} %[3D1V03]
Trong một phép cộng có hai số hạng, nếu tăng số hạng thứ nhất thêm $48$ đơn vị và giảm số hạng thứ hai đi $23$ đơn vị thì được tổng mới bằng $365$. Tìm tổng ban đầu của hai số.
\shortans{340}
\loigiai{
\textbf{Đáp số:} $340$\\
Vì tăng số hạng thứ nhất thêm 48 đơn vị và giảm số hạng thứ hai đi 23 đơn vị nên tổng mới tăng thêm số đơn vị là:
$$48 - 23 = 25\text{ (đơn vị)}$$
Tổng ban đầu của hai số là:
$$365 - 25 = 340$$
}
\end{ex}

\newpage

\TaoPhieuKiemTra{Chủ đề 1 -- Ôn tập và bổ sung}{04}{Ôn tập bảng nhân 2, bảng nhân 5}{https://loptoanthaylinh.vn/thi-online/toan-3/bai-04}
\textbf{Phần 1. Trắc nghiệm lựa chọn}
\setcounterpageref{ex}{0}

\begin{ex} %[3D1V04]
Tính giá trị của biểu thức: $A = 2 \times 8 + 5 \times 6 - 20$.
\choice
{$24$}
{$28$}
{$30$}
{\True $26$}
\loigiai{
Thực hiện các phép tính nhân trước, cộng trừ sau:\\
$$A = 16 + 30 - 20 = 46 - 20 = 26$$
\textbf{Đáp án của câu hỏi là D}
}
\end{ex}

\begin{ex} %[3D1V04]
Cho hai biểu thức $M = 5 \times 7 - 12$ và $N = 2 \times 9 + 6$. So sánh giá trị của $M$ và $N$:
\choice
{$M > N$}
{\True $M < N$}
{$M = N$}
{$M + N = 50$}
\loigiai{
Tính giá trị của từng biểu thức:\\
- $M = 5 \times 7 - 12 = 35 - 12 = 23$\\
- $N = 2 \times 9 + 6 = 18 + 6 = 24$\\
Vì $23 < 24$ nên $M < N$.\\
\textbf{Đáp án của câu hỏi là B}
}
\end{ex}

\begin{ex} %[3D1V04]
Tìm giá trị của $x$, biết: $5 \times x + 25 = 60$.
\choice
{$5$}
{$6$}
{\True $7$}
{$8$}
\loigiai{
Xem $5 \times x$ là số hạng chưa biết:\\
$$5 \times x = 60 - 25$$
$$5 \times x = 35$$
$$x = 35 : 5 = 7$$
\textbf{Đáp án của câu hỏi là C}
}
\end{ex}

\begin{ex} %[3D1V04]
Một nông trại nuôi $9$ con gà và $6$ con vịt. Hỏi đàn gà và đàn vịt đó có tất cả bao nhiêu cái chân?
\choice
{\True $30$ cái chân}
{$28$ cái chân}
{$32$ cái chân}
{$26$ cái chân}
\loigiai{
Mỗi con gà và mỗi con vịt đều có $2$ cái chân.\\
Tổng số con gà và vịt là:
$$9 + 6 = 15\text{ (con)}$$
Đàn gà và vịt có tất cả số cái chân là:
$$15 \times 2 = 30\text{ (cái chân)}$$
\textbf{Đáp án của câu hỏi là A}
}
\end{ex}

\begin{ex} %[3D1V04]
Cô giáo có $45$ cái kẹo, cô chia đều cho các bạn học sinh giỏi, mỗi bạn được $5$ cái kẹo. Sau đó cô giáo lại thưởng thêm cho mỗi bạn $2$ cái kẹo nữa. Hỏi cô giáo đã thưởng thêm tất cả bao nhiêu cái kẹo?
\choice
{$14$ cái kẹo}
{$16$ cái kẹo}
{$20$ cái kẹo}
{\True $18$ cái kẹo}
\loigiai{
Số bạn học sinh giỏi được nhận kẹo là:
$$45 : 5 = 9\text{ (bạn)}$$
Cô giáo đã thưởng thêm tất cả số cái kẹo là:
$$9 \times 2 = 18\text{ (cái kẹo)}$$
\textbf{Đáp án của câu hỏi là D}
}
\end{ex}

\begin{ex} %[3D1V04]
Cho dãy số: $5; 10; 15; 20; \dots; 50$. Tổng của số hạng thứ tư và số hạng thứ chín trong dãy số trên là
\choice
{$60$}
{\True $65$}
{$70$}
{$55$}
\loigiai{
Các số trong dãy là tích của $5$ với số thứ tự của nó:\\
- Số hạng thứ tư là: $5 \times 4 = 20$\\
- Số hạng thứ chín là: $5 \times 9 = 45$\\
Tổng hai số hạng là:
$$20 + 45 = 65$$
\textbf{Đáp án của câu hỏi là B}
}
\end{ex}

\begin{ex} %[3D1V04]
Một số tự nhiên khi chia cho $5$ thì được thương là $8$ và số dư là số dư lớn nhất có thể. Số tự nhiên đó là
\choice
{$42$}
{$43$}
{\True $44$}
{$45$}
\loigiai{
Vì số chia là $5$ nên số dư lớn nhất có thể có là $4$.\\
Số tự nhiên cần tìm là:
$$5 \times 8 + 4 = 40 + 4 = 44$$
\textbf{Đáp án của câu hỏi là C}
}
\end{ex}

\begin{ex} %[3D1V04]
Có 5 túi quà, mỗi túi gồm $4$ viên kẹo cam và $2$ viên kẹo dâu. Hỏi trong 5 túi quà đó có tất cả bao nhiêu viên kẹo?
\choice
{\True $30$ viên kẹo}
{$25$ viên kẹo}
{$35$ viên kẹo}
{$28$ viên kẹo}
\loigiai{
Mỗi túi có số viên kẹo là:
$$4 + 2 = 6\text{ (viên kẹo)}$$
Trong 5 túi có tất cả số viên kẹo là:
$$6 \times 5 = 30\text{ (viên kẹo)}$$
\textbf{Đáp án của câu hỏi là A}
}
\end{ex}

\begin{ex} %[3D1V04]
Tìm giá trị của $x$, biết: $x : 2 + 15 = 24$.
\choice
{$16$}
{$20$}
{$14$}
{\True $18$}
\loigiai{
Xem $x : 2$ là số hạng chưa biết:\\
$$x : 2 = 24 - 15$$
$$x : 2 = 9$$
$$x = 9 \times 2 = 18$$
\textbf{Đáp án của câu hỏi là D}
}
\end{ex}

\begin{ex} %[3D1V04]
Bình nghĩ ra một số. Bình lấy số đó nhân với $5$ rồi trừ đi $18$ thì được kết quả là $27$. Số mà Bình đã nghĩ ra là
\choice
{$8$}
{\True $9$}
{$7$}
{$10$}
\loigiai{
Gọi số Bình nghĩ là $x$. Ta có:\\
$$x \times 5 - 18 = 27$$
$$x \times 5 = 27 + 18$$
$$x \times 5 = 45$$
$$x = 45 : 5 = 9$$
\textbf{Đáp án của câu hỏi là B}
}
\end{ex}

\begin{ex} %[3D1V04]
Một cửa hàng có $50\text{ kg}$ đường. Người ta đóng đều vào $6$ túi lớn, mỗi túi $5\text{ kg}$ đường. Số đường còn lại đóng đều vào các túi nhỏ, mỗi túi chứa $2\text{ kg}$ đường. Hỏi người ta đóng được bao nhiêu túi nhỏ?
\choice
{$8$ túi}
{$9$ túi}
{\True $10$ túi}
{$12$ túi}
\loigiai{
Số đường đóng vào 6 túi lớn là:
$$6 \times 5 = 30\text{ (kg)}$$
Số ki-lô-gam đường còn lại là:
$$50 - 30 = 20\text{ (kg)}$$
Số túi nhỏ đóng được là:
$$20 : 2 = 10\text{ (túi)}$$
\textbf{Đáp án của câu hỏi là C}
}
\end{ex}

\begin{ex} %[3D1V04]
Tính nhanh giá trị của biểu thức: $5 \times 3 + 5 \times 4 + 5 \times 3$.
\choice
{\True $50$}
{$45$}
{$55$}
{$60$}
\loigiai{
Áp dụng tính chất phân phối của phép nhân đối với phép cộng:\\
$$5 \times 3 + 5 \times 4 + 5 \times 3 = 5 \times (3 + 4 + 3) = 5 \times 10 = 50$$
\textbf{Đáp án của câu hỏi là A}
}
\end{ex}

\textbf{Phần 2. Trắc nghiệm đúng sai}
\setcounterpageref{ex}{0}

\begin{ex} %[3D1V04]
Xét tính đúng sai của các nhận định sau về bảng nhân và bảng chia 2, 5:
\choiceTF
{\True Tích của một số tự nhiên bất kì nhân với $2$ luôn luôn là một số chẵn.}
{Một số chia hết cho $5$ thì luôn luôn có chữ số tận cùng là $5$.}
{\True Nếu $x : 5 = 7$ thì giá trị của $x$ bằng $35$.}
{\True Giá trị của biểu thức $2 \times 5$ bằng giá trị của biểu thức $5 \times 2$.}
\loigiai{
\begin{itemize}
    \item Ý a) Đúng: Bất kì số tự nhiên nào nhân với 2 cũng được một số chia hết cho 2 (số chẵn).
    \item Ý b) Sai: Một số chia hết cho 5 có chữ số tận cùng là 0 hoặc 5 (ví dụ 10, 20 chia hết cho 5 nhưng tận cùng là 0).
    \item Ý c) Đúng: $x = 7 \times 5 = 35$.
    \item Ý d) Đúng: Phép nhân có tính chất giao hoán ($2 \times 5 = 5 \times 2 = 10$).
\end{itemize}
\begin{itemize}
    \item a) \textbf{Đúng}.
    \item b) \textbf{Sai}.
    \item c) \textbf{Đúng}.
    \item d) \textbf{Đúng}.
\end{itemize}
}
\end{ex}

\begin{ex} %[3D1V04]
Một nhóm có $10$ bạn học sinh xếp thành $2$ hàng đều nhau, mỗi bạn cầm trên tay $5$ quả bóng bay. Xét tính đúng sai:
\choiceTF
{\True Mỗi hàng có $5$ bạn học sinh.}
{\True Cả nhóm học sinh có tất cả $50$ quả bóng bay.}
{Mỗi hàng học sinh cầm tất cả $20$ quả bóng bay.}
{\True Nếu mỗi bạn thả bay đi $1$ quả bóng thì cả nhóm còn lại $40$ quả bóng bay.}
\loigiai{
\begin{itemize}
    \item Ý a) Đúng: Số bạn mỗi hàng là: $10 : 2 = 5$ (bạn).
    \item Ý b) Đúng: Cả nhóm có: $10 \times 5 = 50$ (quả bóng bay).
    \item Ý c) Sai: Mỗi hàng có 5 bạn nên cầm số bóng là: $5 \times 5 = 25$ (quả bóng bay).
    \item Ý d) Đúng: Mỗi bạn còn 4 quả bóng, 10 bạn còn: $10 \times 4 = 40$ (quả bóng bay).
\end{itemize}
\begin{itemize}
    \item a) \textbf{Đúng}.
    \item b) \textbf{Đúng}.
    \item c) \textbf{Sai}.
    \item d) \textbf{Đúng}.
\end{itemize}
}
\end{ex}

\textbf{Phần 3. Trả lời ngắn}
\setcounterpageref{ex}{0}

\begin{ex} %[3D1V04]
Tính giá trị của biểu thức: $5 \times 9 - 2 \times 8$.
\shortans{29}
\loigiai{
\textbf{Đáp số:} $29$\\
Thực hiện phép nhân trước:\\
$$5 \times 9 - 2 \times 8 = 45 - 16 = 29$$
}
\end{ex}

\begin{ex} %[3D1V04]
Một người nuôi $8$ con thỏ. Hỏi $8$ con thỏ đó có tất cả bao nhiêu cái tai?
\shortans{16}
\loigiai{
\textbf{Đáp số:} $16$\\
Mỗi con thỏ có 2 cái tai. Số cái tai của 8 con thỏ là:
$$8 \times 2 = 16\text{ (cái tai)}$$
}
\end{ex}

\begin{ex} %[3D1V04]
Tìm số bị chia trong một phép chia cho $5$, biết thương bằng $6$ và số dư bằng $3$.
\shortans{33}
\loigiai{
\textbf{Đáp số:} $33$\\
Số bị chia là:
$$5 \times 6 + 3 = 30 + 3 = 33$$
}
\end{ex}

\begin{ex} %[3D1V04]
Lớp 3A có $35$ học sinh. Cô giáo chia đều học sinh thành các nhóm, mỗi nhóm có $5$ bạn. Sau đó mỗi nhóm cử $2$ bạn tham gia đội sao đỏ. Hỏi lớp 3A có tất cả bao nhiêu bạn tham gia đội sao đỏ?
\shortans{14}
\loigiai{
\textbf{Đáp số:} $14$\\
Lớp 3A chia được số nhóm là:
$$35 : 5 = 7\text{ (nhóm)}$$
Số bạn tham gia đội sao đỏ là:
$$7 \times 2 = 14\text{ (bạn)}$$
}
\end{ex}

\begin{ex} %[3D1V04]
Tìm giá trị của $x$, biết: $x : 5 = 14 - 6$.
\shortans{40}
\loigiai{
\textbf{Đáp số:} $40$\\
Tính vế phải:\\
$$x : 5 = 8$$
$$x = 8 \times 5 = 40$$
}
\end{ex}

\begin{ex} %[3D1V04]
Tính nhanh giá trị của biểu thức: $2 \times 1 + 2 \times 2 + 2 \times 3 + 2 \times 4$.
\shortans{20}
\loigiai{
\textbf{Đáp số:} $20$\\
Đặt thừa số chung:\\
$$2 \times (1 + 2 + 3 + 4) = 2 \times 10 = 20$$
}
\end{ex}

\newpage

\TaoPhieuKiemTra{Chủ đề 1 -- Ôn tập và bổ sung}{05}{Ôn tập bảng chia 2, bảng chia 5}{https://loptoanthaylinh.vn/thi-online/toan-3/bai-05}
\textbf{Phần 1. Trắc nghiệm lựa chọn}
\setcounterpageref{ex}{0}

\begin{ex} %[3D1V05]
Tính giá trị của biểu thức: $A = 3 \times 9 - 27 : 3 + 15$.
\choice
{$30$}
{$36$}
{\True $33$}
{$32$}
\loigiai{
Thực hiện các phép tính nhân và chia trước, cộng trừ sau:\\
$$A = 27 - 9 + 15 = 18 + 15 = 33$$
\textbf{Đáp án của câu hỏi là C}
}
\end{ex}

\begin{ex} %[3D1V05]
Điền dấu thích hợp vào chỗ chấm: $3 \times 7 + 9 \dots 30 : 3 \times 3$.
\choice
{$>$}
{\True $=$}
{$<$}
{Không so sánh được}
\loigiai{
Tính giá trị của hai vế:\\
- Vế trái: $3 \times 7 + 9 = 21 + 9 = 30$\\
- Vế phải: $30 : 3 \times 3 = 10 \times 3 = 30$\\
Vì $30 = 30$ nên ta điền dấu $=$.\\
\textbf{Đáp án của câu hỏi là B}
}
\end{ex}

\begin{ex} %[3D1V05]
Tìm giá trị của $x$, biết: $x : 3 + 18 = 25$.
\choice
{$24$}
{$18$}
{$27$}
{\True $21$}
\loigiai{
Xem $x : 3$ là số hạng chưa biết:\\
$$x : 3 = 25 - 18$$
$$x : 3 = 7$$
$$x = 7 \times 3 = 21$$
\textbf{Đáp án của câu hỏi là D}
}
\end{ex}

\begin{ex} %[3D1V05]
Lớp 3B trồng được $8$ cây hoa. Lớp 3A trồng được số cây hoa gấp $3$ lần lớp 3B. Hỏi cả hai lớp trồng được tất cả bao nhiêu cây hoa?
\choice
{\True $32$ cây}
{$24$ cây}
{$30$ cây}
{$28$ cây}
\loigiai{
Lớp 3A trồng được số cây hoa là:
$$8 \times 3 = 24\text{ (cây)}$$
Cả hai lớp trồng được tất cả số cây hoa là:
$$8 + 24 = 32\text{ (cây)}$$
\textbf{Đáp án của câu hỏi là A}
}
\end{ex}

\begin{ex} %[3D1V05]
Có $27\text{ lít}$ mật ong chia đều vào $3$ can lớn. Người ta lấy $1$ can lớn đó rót đều vào các chai nhỏ, mỗi chai chứa $3\text{ lít}$. Hỏi rót được bao nhiêu chai nhỏ?
\choice
{$4$ chai}
{$2$ chai}
{\True $3$ chai}
{$5$ chai}
\loigiai{
Mỗi can lớn chứa số lít mật ong là:
$$27 : 3 = 9\text{ (lít)}$$
Số chai nhỏ rót được là:
$$9 : 3 = 3\text{ (chai)}$$
\textbf{Đáp án của câu hỏi là C}
}
\end{ex}

\begin{ex} %[3D1V05]
Cho dãy số: $3; 6; 9; 12; \dots; 27$. Số hạng đứng chính giữa của dãy số trên là
\choice
{$12$}
{\True $15$}
{$18$}
{$9$}
\loigiai{
Dãy số gồm các bội số của 3 từ $3 \times 1$ đến $3 \times 9$, có tất cả 9 số hạng.\\
Số hạng đứng chính giữa là số hạng thứ năm:\\
$$3 \times 5 = 15$$
\textbf{Đáp án của câu hỏi là B}
}
\end{ex}

\begin{ex} %[3D1V05]
Trong một phép chia cho $3$, thương bằng $7$ và số dư là số dư lớn nhất có thể có. Số bị chia của phép chia đó là
\choice
{$21$}
{$22$}
{$24$}
{\True $23$}
\loigiai{
Vì số chia là 3 nên số dư lớn nhất có thể có là $2$.\\
Số bị chia là:
$$3 \times 7 + 2 = 21 + 2 = 23$$
\textbf{Đáp án của câu hỏi là D}
}
\end{ex}

\begin{ex} %[3D1V05]
Tính nhanh giá trị của biểu thức: $3 \times 2 + 3 \times 3 + 3 \times 5$.
\choice
{\True $30$}
{$27$}
{$33$}
{$36$}
\loigiai{
Áp dụng tính chất phân phối:\\
$$3 \times 2 + 3 \times 3 + 3 \times 5 = 3 \times (2 + 3 + 5) = 3 \times 10 = 30$$
\textbf{Đáp án của câu hỏi là A}
}
\end{ex}

\begin{ex} %[3D1V05]
Mẹ mua $3$ hộp bánh, mỗi hộp có giá $8$ nghìn đồng và mua thêm $1$ gói kẹo giá $6$ nghìn đồng. Mẹ đưa cho cô bán hàng tờ tiền $50$ nghìn đồng. Hỏi cô bán hàng phải trả lại mẹ bao nhiêu nghìn đồng?
\choice
{$15$ nghìn đồng}
{$25$ nghìn đồng}
{\True $20$ nghìn đồng}
{$10$ nghìn đồng}
\loigiai{
Số tiền mẹ mua 3 hộp bánh là:
$$3 \times 8 = 24\text{ (nghìn đồng)}$$
Tổng số tiền mẹ đã mua bánh và kẹo là:
$$24 + 6 = 30\text{ (nghìn đồng)}$$
Cô bán hàng phải trả lại mẹ số tiền là:
$$50 - 30 = 20\text{ (nghìn đồng)}$$
\textbf{Đáp án của câu hỏi là C}
}
\end{ex}

\begin{ex} %[3D1V05]
Tìm số tự nhiên $a$, biết: $3 \times 6 < a < 3 \times 8$ và $a$ chia hết cho $3$.
\choice
{$19$}
{\True $21$}
{$20$}
{$22$}
\loigiai{
Ta có: $18 < a < 24$.\\
Các số tự nhiên nằm giữa 18 và 24 là: $19; 20; 21; 22; 23$.\\
Trong các số trên, số duy nhất chia hết cho 3 là $21$ (vì $21 : 3 = 7$).\\
\textbf{Đáp án của câu hỏi là B}
}
\end{ex}

\begin{ex} %[3D1V05]
Trên một đoạn đường dài $27\text{ m}$, người ta trồng các cây hoa cách đều nhau $3\text{ m}$, ở cả hai đầu đường đều có trồng cây. Hỏi trên đoạn đường đó có tất cả bao nhiêu cây hoa?
\choice
{$8$ cây}
{$9$ cây}
{$11$ cây}
{\True $10$ cây}
\loigiai{
Số khoảng cách giữa các cây là:
$$27 : 3 = 9\text{ (khoảng)}$$
Vì cả hai đầu đường đều trồng cây nên số cây hoa là:
$$9 + 1 = 10\text{ (cây)}$$
\textbf{Đáp án của câu hỏi là D}
}
\end{ex}

\begin{ex} %[3D1V05]
Tìm giá trị của $x$, biết: $(x - 5) \times 3 = 24$.
\choice
{\True $13$}
{$12$}
{$14$}
{$15$}
\loigiai{
Xem $(x - 5)$ là thừa số chưa biết:\\
$$x - 5 = 24 : 3$$
$$x - 5 = 8$$
$$x = 8 + 5 = 13$$
\textbf{Đáp án của câu hỏi là A}
}
\end{ex}

\textbf{Phần 2. Trắc nghiệm đúng sai}
\setcounterpageref{ex}{0}

\begin{ex} %[3D1V05]
Xét tính đúng sai của các nhận định sau về phép nhân và phép chia với $3$:
\choiceTF
{\True Tích của một số lẻ bất kì nhân với $3$ luôn luôn là một số lẻ.}
{Trong một phép chia cho $3$, số dư lớn nhất có thể có là $3$.}
{\True Nếu một số giảm đi $3$ lần rồi lại gấp lên $3$ lần thì được số ban đầu.}
{\True Tổng của ba số tự nhiên liên tiếp luôn luôn là một số chia hết cho $3$.}
\loigiai{
\begin{itemize}
    \item Ý a) Đúng: Tích của hai số lẻ luôn là số lẻ (ví dụ $1 \times 3 = 3$, $3 \times 3 = 9$, $5 \times 3 = 15$).
    \item Ý b) Sai: Số dư phải luôn bé hơn số chia, vì số chia là 3 nên số dư lớn nhất là 2.
    \item Ý c) Đúng: Vì $a : 3 \times 3 = a$.
    \item Ý d) Đúng: Gọi 3 số liên tiếp là $n, n+1, n+2$. Tổng là $3n + 3 = 3 \times (n+1)$, chia hết cho 3.
\end{itemize}
\begin{itemize}
    \item a) \textbf{Đúng}.
    \item b) \textbf{Sai}.
    \item c) \textbf{Đúng}.
    \item d) \textbf{Đúng}.
\end{itemize}
}
\end{ex}

\begin{ex} %[3D1V05]
Một cửa hàng có $3$ thùng dầu, mỗi thùng chứa $9\text{ lít}$ dầu. Cửa hàng đã bán đi $12\text{ lít}$ dầu. Xét tính đúng sai:
\choiceTF
{\True Ban đầu cửa hàng có tất cả $27\text{ lít}$ dầu.}
{\True Sau khi bán, cửa hàng còn lại $15\text{ lít}$ dầu.}
{Số dầu còn lại đủ để rót đều vào 4 can loại $4\text{ lít}$.}
{\True Số dầu còn lại rót vừa đủ vào 5 chai loại $3\text{ lít}$.}
\loigiai{
\begin{itemize}
    \item Ý a) Đúng: Tổng số dầu ban đầu là: $3 \times 9 = 27\text{ (lít)}$.
    \item Ý b) Đúng: Số dầu còn lại là: $27 - 12 = 15\text{ (lít)}$.
    \item Ý c) Sai: 4 can loại 4 lít cần $4 \times 4 = 16\text{ lít}$, mà cửa hàng chỉ còn 15 lít.
    \item Ý d) Đúng: $15 : 3 = 5$ (rót vừa đủ 5 chai loại 3 lít).
\end{itemize}
\begin{itemize}
    \item a) \textbf{Đúng}.
    \item b) \textbf{Đúng}.
    \item c) \textbf{Sai}.
    \item d) \textbf{Đúng}.
\end{itemize}
}
\end{ex}

\textbf{Phần 3. Trả lời ngắn}
\setcounterpageref{ex}{0}

\begin{ex} %[3D1V05]
Tính giá trị của biểu thức: $3 \times 7 + 18 : 3$.
\shortans{27}
\loigiai{
\textbf{Đáp số:} $27$\\
Thực hiện nhân chia trước, cộng sau:\\
$$3 \times 7 + 18 : 3 = 21 + 6 = 27$$
}
\end{ex}

\begin{ex} %[3D1V05]
Một người thợ cưa một khúc gỗ dài $15\text{ m}$ thành các đoạn dài $3\text{ m}$. Biết mỗi lần cưa người đó mất $5$ phút. Hỏi người đó cưa xong khúc gỗ hết tất cả bao nhiêu phút?
\shortans{20}
\loigiai{
\textbf{Đáp số:} $20$\\
Số đoạn gỗ cưa được là: $15 : 3 = 5$ (đoạn).\\
Để cưa thành 5 đoạn thì chỉ cần cưa: $5 - 1 = 4$ (lần cưa).\\
Thời gian cưa xong khúc gỗ là:
$$4 \times 5 = 20\text{ (phút)}$$
}
\end{ex}

\begin{ex} %[3D1V05]
Tìm số bị chia trong một phép chia cho $3$, biết thương bằng $8$ và số dư bằng $1$.
\shortans{25}
\loigiai{
\textbf{Đáp số:} $25$\\
Số bị chia là:
$$3 \times 8 + 1 = 24 + 1 = 25$$
}
\end{ex}

\begin{ex} %[3D1V05]
Nam có $12$ viên bi. Số viên bi của Việt gấp $3$ lần số viên bi của Nam. Hỏi cả hai bạn có tất cả bao nhiêu viên bi?
\shortans{48}
\loigiai{
\textbf{Đáp số:} $48$\\
Số viên bi của Việt là:
$$12 \times 3 = 36\text{ (viên bi)}$$
Cả hai bạn có tất cả số viên bi là:
$$12 + 36 = 48\text{ (viên bi)}$$
}
\end{ex}

\begin{ex} %[3D1V05]
Tìm giá trị của $x$, biết: $3 \times x = 45 - 18$.
\shortans{9}
\loigiai{
\textbf{Đáp số:} $9$\\
Tính vế phải:\\
$$3 \times x = 27$$
$$x = 27 : 3 = 9$$
}
\end{ex}

\begin{ex} %[3D1V05]
Tính nhanh giá trị của biểu thức: $3 \times 1 + 3 \times 2 + 3 \times 3 + 3 \times 4$.
\shortans{30}
\loigiai{
\textbf{Đáp số:} $30$\\
Đặt thừa số chung:\\
$$3 \times (1 + 2 + 3 + 4) = 3 \times 10 = 30$$
}
\end{ex}

\newpage

\TaoPhieuKiemTra{Chủ đề 1 -- Ôn tập và bổ sung}{06}{Gấp một số lên một số lần}{https://loptoanthaylinh.vn/thi-online/toan-3/bai-06}
\textbf{Phần 1. Trắc nghiệm lựa chọn}
\setcounterpageref{ex}{0}

\begin{ex} %[3D1V06]
Tính giá trị của biểu thức: $A = 4 \times 8 - 36 : 4 + 12$.
\choice
{$33$}
{\True $35$}
{$37$}
{$31$}
\loigiai{
Thực hiện nhân và chia trước, cộng trừ sau:\\
$$A = 32 - 9 + 12 = 23 + 12 = 35$$
\textbf{Đáp án của câu hỏi là B}
}
\end{ex}

\begin{ex} %[3D1V06]
Một đàn bò có $7$ con và một đàn lợn có $5$ con. Hỏi cả đàn bò và đàn lợn đó có tất cả bao nhiêu cái chân?
\choice
{\True $48$ cái chân}
{$44$ cái chân}
{$52$ cái chân}
{$40$ cái chân}
\loigiai{
Mỗi con bò và mỗi con lợn đều có $4$ cái chân.\\
Tổng số con bò và lợn là:
$$7 + 5 = 12\text{ (con)}$$
Cả hai đàn có tất cả số cái chân là:
$$12 \times 4 = 48\text{ (cái chân)}$$
\textbf{Đáp án của câu hỏi là A}
}
\end{ex}

\begin{ex} %[3D1V06]
Tìm giá trị của $x$, biết: $x : 4 + 23 = 32$.
\choice
{$32$}
{$40$}
{\True $36$}
{$28$}
\loigiai{
Xem $x : 4$ là số hạng chưa biết:\\
$$x : 4 = 32 - 23$$
$$x : 4 = 9$$
$$x = 9 \times 4 = 36$$
\textbf{Đáp án của câu hỏi là C}
}
\end{ex}

\begin{ex} %[3D1V06]
Một chiếc khăn tay hình vuông có cạnh dài $8\text{ cm}$. Người ta viền ren xung quanh chiếc khăn tay đó. Độ dài đoạn ren cần dùng là
\choice
{$24\text{ cm}$}
{$28\text{ cm}$}
{$36\text{ cm}$}
{\True $32\text{ cm}$}
\loigiai{
Độ dài đoạn ren viền xung quanh khăn tay chính là chu vi hình vuông:\\
$$8 \times 4 = 32\text{ (cm)}$$
\textbf{Đáp án của câu hỏi là D}
}
\end{ex}

\begin{ex} %[3D1V06]
Điền dấu thích hợp vào chỗ chấm: $4 \times 6 - 8 \dots 32 : 4 \times 2$.
\choice
{$>$}
{\True $=$}
{$<$}
{Không so sánh được}
\loigiai{
Tính giá trị hai vế:\\
- Vế trái: $4 \times 6 - 8 = 24 - 8 = 16$\\
- Vế phải: $32 : 4 \times 2 = 8 \times 2 = 16$\\
Vì $16 = 16$ nên ta điền dấu $=$.\\
\textbf{Đáp án của câu hỏi là B}
}
\end{ex}

\begin{ex} %[3D1V06]
Trong một phép chia cho $4$, có thương bằng $6$ và số dư là số dư lớn nhất có thể có. Số bị chia của phép chia đó là
\choice
{\True $27$}
{$26$}
{$28$}
{$25$}
\loigiai{
Vì số chia là 4 nên số dư lớn nhất có thể có là $3$.\\
Số bị chia là:
$$4 \times 6 + 3 = 24 + 3 = 27$$
\textbf{Đáp án của câu hỏi là A}
}
\end{ex}

\begin{ex} %[3D1V06]
Khối lớp 3 có $36$ bạn học sinh tham gia đồng diễn thể dục. Thầy giáo chia đều các bạn thành $4$ hàng, sau đó mỗi hàng lại chia thành $3$ nhóm nhỏ. Hỏi mỗi nhóm nhỏ có bao nhiêu bạn học sinh?
\choice
{$4$ bạn}
{$2$ bạn}
{\True $3$ bạn}
{$5$ bạn}
\loigiai{
Mỗi hàng có số bạn học sinh là:
$$36 : 4 = 9\text{ (bạn)}$$
Mỗi nhóm nhỏ có số bạn học sinh là:
$$9 : 3 = 3\text{ (bạn)}$$
\textbf{Đáp án của câu hỏi là C}
}
\end{ex}

\begin{ex} %[3D1V06]
Tính nhanh giá trị của biểu thức: $4 \times 1 + 4 \times 2 + 4 \times 3 + 4 \times 4$.
\choice
{$36$}
{$44$}
{$48$}
{\True $40$}
\loigiai{
Đặt thừa số chung:\\
$$4 \times (1 + 2 + 3 + 4) = 4 \times 10 = 40$$
\textbf{Đáp án của câu hỏi là D}
}
\end{ex}

\begin{ex} %[3D1V06]
Tìm giá trị của $x$, biết: $(x - 7) \times 4 = 28$.
\choice
{$12$}
{\True $14$}
{$16$}
{$15$}
\loigiai{
Xem $(x - 7)$ là thừa số chưa biết:\\
$$x - 7 = 28 : 4$$
$$x - 7 = 7$$
$$x = 7 + 7 = 14$$
\textbf{Đáp án của câu hỏi là B}
}
\end{ex}

\begin{ex} %[3D1V06]
Bác Ba nuôi $6$ con thỏ. Số con gà bác Ba nuôi gấp $4$ lần số con thỏ. Hỏi bác Ba nuôi số con gà nhiều hơn số con thỏ bao nhiêu con?
\choice
{\True $18$ con}
{$24$ con}
{$16$ con}
{$20$ con}
\loigiai{
Số con gà bác Ba nuôi là:
$$6 \times 4 = 24\text{ (con)}$$
Số con gà nhiều hơn số con thỏ là:
$$24 - 6 = 18\text{ (con)}$$
\textbf{Đáp án của câu hỏi là A}
}
\end{ex}

\begin{ex} %[3D1V06]
Cho dãy số cách đều: $4; 8; 12; 16; \dots$ Số hạng thứ 8 của dãy số trên là
\choice
{$28$}
{$36$}
{\True $32$}
{$40$}
\loigiai{
Dãy số gồm các số chia hết cho 4. Số hạng thứ 8 là:
$$4 \times 8 = 32$$
\textbf{Đáp án của câu hỏi là C}
}
\end{ex}

\begin{ex} %[3D1V06]
Trong một bãi đỗ xe có $8$ chiếc xe ô tô con (mỗi xe có $4$ bánh) và $5$ chiếc xe máy (mỗi xe có $2$ bánh). Tổng số bánh xe của cả hai loại xe đó là
\choice
{$38$ bánh xe}
{$40$ bánh xe}
{$44$ bánh xe}
{\True $42$ bánh xe}
\loigiai{
Số bánh xe của 8 chiếc ô tô con là:
$$8 \times 4 = 32\text{ (bánh xe)}$$
Số bánh xe của 5 chiếc xe máy là:
$$5 \times 2 = 10\text{ (bánh xe)}$$
Tổng số bánh xe là:
$$32 + 10 = 42\text{ (bánh xe)}$$
\textbf{Đáp án của câu hỏi là D}
}
\end{ex}

\textbf{Phần 2. Trắc nghiệm đúng sai}
\setcounterpageref{ex}{0}

\begin{ex} %[3D1V06]
Xét tính đúng sai của các khẳng định sau về phép nhân và phép chia với $4$:
\choiceTF
{\True Tích của số $4$ với một số tự nhiên bất kì luôn luôn là một số chẵn.}
{Trong một phép chia cho $4$, số dư có thể bằng $4$.}
{\True Một số tự nhiên chia hết cho $4$ thì chắc chắn chia hết cho $2$.}
{\True Giá trị của biểu thức $4 \times 9$ lớn hơn giá trị của biểu thức $5 \times 7$.}
\loigiai{
\begin{itemize}
    \item Ý a) Đúng: Vì $4 = 2 \times 2$ nên $4 \times n = 2 \times (2 \times n)$ luôn là số chẵn.
    \item Ý b) Sai: Số dư luôn phải nhỏ hơn số chia (số dư có thể là 0, 1, 2, 3; không thể bằng 4).
    \item Ý c) Đúng: Vì 4 chia hết cho 2 nên số chia hết cho 4 cũng chia hết cho 2.
    \item Ý d) Đúng: $4 \times 9 = 36$ và $5 \times 7 = 35$. Vì $36 > 35$ nên nhận định đúng.
\end{itemize}
\begin{itemize}
    \item a) \textbf{Đúng}.
    \item b) \textbf{Sai}.
    \item c) \textbf{Đúng}.
    \item d) \textbf{Đúng}.
\end{itemize}
}
\end{ex}

\begin{ex} %[3D1V06]
Một cuộn dây thừng dài $36\text{ m}$ được cắt thành $4$ đoạn dài bằng nhau:
\choiceTF
{\True Mỗi đoạn dây dài $9\text{ m}$.}
{Để cắt cuộn dây thành $4$ đoạn, người ta cần thực hiện đúng $4$ nhát cắt.}
{\True Tổng độ dài của $3$ đoạn dây là $27\text{ m}$.}
{\True Dùng $1$ đoạn dây đó để viền quanh một khung tranh hình vuông có cạnh dài $2\text{ m}$ thì còn thừa đúng $1\text{ m}$ dây.}
\loigiai{
\begin{itemize}
    \item Ý a) Đúng: Chiều dài mỗi đoạn là: $36 : 4 = 9\text{ (m)}$.
    \item Ý b) Sai: Để cắt thành 4 đoạn chỉ cần $4 - 1 = 3$ nhát cắt.
    \item Ý c) Đúng: 3 đoạn dây dài: $3 \times 9 = 27\text{ (m)}$.
    \item Ý d) Đúng: Chu vi khung tranh là: $2 \times 4 = 8\text{ (m)}$. Đoạn dây 9 m còn thừa: $9 - 8 = 1\text{ (m)}$.
\end{itemize}
\begin{itemize}
    \item a) \textbf{Đúng}.
    \item b) \textbf{Sai}.
    \item c) \textbf{Đúng}.
    \item d) \textbf{Đúng}.
\end{itemize}
}
\end{ex}

\textbf{Phần 3. Trả lời ngắn}
\setcounterpageref{ex}{0}

\begin{ex} %[3D1V06]
Tính giá trị của biểu thức: $4 \times 7 - 20 : 4$.
\shortans{23}
\loigiai{
\textbf{Đáp số:} $23$\\
Thực hiện nhân chia trước:\\
$$4 \times 7 - 20 : 4 = 28 - 5 = 23$$
}
\end{ex}

\begin{ex} %[3D1V06]
Một trang trại nuôi $9$ con thỏ. Hỏi $9$ con thỏ đó có tất cả bao nhiêu cái chân?
\shortans{36}
\loigiai{
\textbf{Đáp số:} $36$\\
Mỗi con thỏ có 4 cái chân. Số cái chân của 9 con thỏ là:
$$9 \times 4 = 36\text{ (cái chân)}$$
}
\end{ex}

\begin{ex} %[3D1V06]
Tìm số bị chia trong một phép chia cho $4$, biết thương bằng $7$ và số dư bằng $2$.
\shortans{30}
\loigiai{
\textbf{Đáp số:} $30$\\
Số bị chia là:
$$4 \times 7 + 2 = 28 + 2 = 30$$
}
\end{ex}

\begin{ex} %[3D1V06]
Lớp 3C có $28$ học sinh được chia đều thành $4$ tổ. Hỏi $3$ tổ của lớp 3C có tất cả bao nhiêu học sinh?
\shortans{21}
\loigiai{
\textbf{Đáp số:} $21$\\
Mỗi tổ có số học sinh là:
$$28 : 4 = 7\text{ (học sinh)}$$
3 tổ có tất cả số học sinh là:
$$7 \times 3 = 21\text{ (học sinh)}$$
}
\end{ex}

\begin{ex} %[3D1V06]
Tìm giá trị của $x$, biết: $x : 4 = 15 - 8$.
\shortans{28}
\loigiai{
\textbf{Đáp số:} $28$\\
Tính vế phải:\\
$$x : 4 = 7$$
$$x = 7 \times 4 = 28$$
}
\end{ex}

\begin{ex} %[3D1V06]
Tính nhanh giá trị của biểu thức: $4 \times 2 + 4 \times 3 + 4 \times 5$.
\shortans{40}
\loigiai{
\textbf{Đáp số:} $40$\\
Đặt thừa số chung:\\
$$4 \times (2 + 3 + 5) = 4 \times 10 = 40$$
}
\end{ex}

\newpage

\TaoPhieuKiemTra{Chủ đề 1 -- Ôn tập và bổ sung}{07}{Giảm đi một số lần}{https://loptoanthaylinh.vn/thi-online/toan-3/bai-07}
\textbf{Phần 1. Trắc nghiệm lựa chọn}
\setcounterpageref{ex}{0}

\begin{ex} %[3H1V07]
Một sợi dây đồng dài $2\text{ m } 35\text{ cm}$. Một sợi dây nhôm dài $24\text{ dm}$. Hỏi sợi dây nào dài hơn và dài hơn bao nhiêu xăng-ti-mét?
\choice
{Sợi dây đồng dài hơn $5\text{ cm}$}
{Sợi dây nhôm dài hơn $15\text{ cm}$}
{\True Sợi dây nhôm dài hơn $5\text{ cm}$}
{Sợi dây đồng dài hơn $15\text{ cm}$}
\loigiai{
Đổi các số đo về cùng đơn vị xăng-ti-mét:\\
- Sợi dây đồng dài: $2\text{ m } 35\text{ cm} = 235\text{ cm}$\\
- Sợi dây nhôm dài: $24\text{ dm} = 240\text{ cm}$\\
Vì $240\text{ cm} > 235\text{ cm}$ nên sợi dây nhôm dài hơn và dài hơn số xăng-ti-mét là:
$$240 - 235 = 5\text{ (cm)}$$
\textbf{Đáp án của câu hỏi là C}
}
\end{ex}

\begin{ex} %[3H1V07]
Đường gấp khúc $ABCD$ gồm 3 đoạn thẳng: đoạn thẳng $AB$ dài $18\text{ cm}$, đoạn thẳng $BC$ dài hơn đoạn thẳng $AB$ là $6\text{ cm}$, đoạn thẳng $CD$ dài bằng một nửa đoạn thẳng $BC$. Độ dài đường gấp khúc $ABCD$ là
\choice
{\True $54\text{ cm}$}
{$48\text{ cm}$}
{$56\text{ cm}$}
{$50\text{ cm}$}
\loigiai{
Độ dài đoạn thẳng $BC$ là:
$$18 + 6 = 24\text{ (cm)}$$
Độ dài đoạn thẳng $CD$ là:
$$24 : 2 = 12\text{ (cm)}$$
Độ dài đường gấp khúc $ABCD$ là:
$$18 + 24 + 12 = 54\text{ (cm)}$$
\textbf{Đáp án của câu hỏi là A}
}
\end{ex}

\begin{ex} %[3H1V07]
Một hình tam giác có độ dài cạnh thứ nhất là $15\text{ cm}$, cạnh thứ hai dài $20\text{ cm}$. Biết chu vi của hình tam giác đó là $52\text{ cm}$. Độ dài cạnh thứ ba của hình tam giác là
\choice
{$15\text{ cm}$}
{$16\text{ cm}$}
{$18\text{ cm}$}
{\True $17\text{ cm}$}
\loigiai{
Tổng độ dài hai cạnh thứ nhất và thứ hai là:
$$15 + 20 = 35\text{ (cm)}$$
Độ dài cạnh thứ ba của hình tam giác là:
$$52 - 35 = 17\text{ (cm)}$$
\textbf{Đáp án của câu hỏi là D}
}
\end{ex}

\begin{ex} %[3H1V07]
Trên đĩa cân bên trái đặt một quả mít và một quả cân $1\text{ kg}$. Trên đĩa cân bên phải đặt hai quả cân loại $2\text{ kg}$ và một quả cân $1\text{ kg}$ thì cân thăng bằng. Quả mít đó nặng số ki-lô-gam là
\choice
{$3\text{ kg}$}
{\True $4\text{ kg}$}
{$5\text{ kg}$}
{$2\text{ kg}$}
\loigiai{
Tổng khối lượng các quả cân trên đĩa cân bên phải là:
$$2 + 2 + 1 = 5\text{ (kg)}$$
Vì cân ở vị trí thăng bằng nên quả mít nặng số ki-lô-gam là:
$$5 - 1 = 4\text{ (kg)}$$
\textbf{Đáp án của câu hỏi là B}
}
\end{ex}

\begin{ex} %[3H1V07]
Có hai thùng chứa nước mắm. Thùng thứ nhất chứa $48\text{ lít}$ nước mắm. Nếu chuyển $8\text{ lít}$ nước mắm từ thùng thứ nhất sang thùng thứ hai thì hai thùng có số lít nước mắm bằng nhau. Lúc đầu thùng thứ hai chứa số lít nước mắm là
\choice
{$36\text{ lít}$}
{$40\text{ lít}$}
{\True $32\text{ lít}$}
{$24\text{ lít}$}
\loigiai{
Sau khi chuyển $8\text{ lít}$ nước mắm, thùng thứ nhất còn lại số lít là:
$$48 - 8 = 40\text{ (lít)}$$
Lúc này thùng thứ hai cũng có $40\text{ lít}$ nước mắm.\\
Ban đầu thùng thứ hai chứa số lít nước mắm là:
$$40 - 8 = 32\text{ (lít)}$$
\textbf{Đáp án của câu hỏi là C}
}
\end{ex}

\begin{ex} %[3H1V07]
Một mảnh vườn hình tứ giác có độ dài các cạnh lần lượt là $12\text{ m}, 15\text{ m}, 18\text{ m}$ và $25\text{ m}$. Người ta làm hàng rào bao quanh mảnh vườn đó và để lại một cổng vào rộng $2\text{ m}$. Độ dài hàng rào bao quanh mảnh vườn là
\choice
{\True $68\text{ m}$}
{$70\text{ m}$}
{$72\text{ m}$}
{$66\text{ m}$}
\loigiai{
Chu vi mảnh vườn hình tứ giác là:
$$12 + 15 + 18 + 25 = 70\text{ (m)}$$
Độ dài hàng rào bao quanh mảnh vườn là:
$$70 - 2 = 68\text{ (m)}$$
\textbf{Đáp án của câu hỏi là A}
}
\end{ex}

\begin{ex} %[3H1V07]
Tính giá trị của biểu thức $1\text{ m } 8\text{ dm} - 45\text{ cm} + 2\text{ dm}$ khi đổi ra xăng-ti-mét:
\choice
{$145\text{ cm}$}
{$165\text{ cm}$}
{$150\text{ cm}$}
{\True $155\text{ cm}$}
\loigiai{
Đổi các số đo sang đơn vị xăng-ti-mét:\\
- $1\text{ m } 8\text{ dm} = 180\text{ cm}$\\
- $2\text{ dm} = 20\text{ cm}$\\
Thực hiện phép tính:\\
$$180\text{ cm} - 45\text{ cm} + 20\text{ cm} = 135\text{ cm} + 20\text{ cm} = 155\text{ cm}$$
\textbf{Đáp án của câu hỏi là D}
}
\end{ex}

\begin{ex} %[3H1V07]
Đoạn thẳng $AB$ dài $16\text{ cm}$. Điểm $M$ là trung điểm của đoạn thẳng $AB$. Điểm $N$ là trung điểm của đoạn thẳng $AM$. Độ dài đoạn thẳng $AN$ là
\choice
{$2\text{ cm}$}
{\True $4\text{ cm}$}
{$6\text{ cm}$}
{$8\text{ cm}$}
\loigiai{
Vì $M$ là trung điểm của đoạn thẳng $AB$ nên độ dài đoạn thẳng $AM$ là:
$$16 : 2 = 8\text{ (cm)}$$
Vì $N$ là trung điểm của đoạn thẳng $AM$ nên độ dài đoạn thẳng $AN$ là:
$$8 : 2 = 4\text{ (cm)}$$
\textbf{Đáp án của câu hỏi là B}
}
\end{ex}

\begin{ex} %[3H1V07]
Một bao gạo cân nặng $50\text{ kg}$. Người ta lấy ra một nửa số gạo trong bao rồi đổ thêm vào bao $5\text{ kg}$ gạo nếp. Lúc này trong bao có tất cả số ki-lô-gam gạo là
\choice
{$25\text{ kg}$}
{$35\text{ kg}$}
{\True $30\text{ kg}$}
{$28\text{ kg}$}
\loigiai{
Một nửa số gạo trong bao là:
$$50 : 2 = 25\text{ (kg)}$$
Lúc này trong bao có tất cả số ki-lô-gam gạo là:
$$25 + 5 = 30\text{ (kg)}$$
\textbf{Đáp án của câu hỏi là C}
}
\end{ex}

\begin{ex} %[3H1V07]
Một bình chứa $1\text{ lít}$ nước. Người ta rót nước từ bình vào đầy $3$ chiếc cốc loại $200\text{ ml}$ và $1$ chiếc cốc loại $250\text{ ml}$. Lượng nước còn lại trong bình là
\choice
{\True $150\text{ ml}$}
{$200\text{ ml}$}
{$100\text{ ml}$}
{$250\text{ ml}$}
\loigiai{
Đổi: $1\text{ lít} = 1000\text{ ml}$.\\
Lượng nước đã rót vào các cốc là:
$$3 \times 200 + 250 = 600 + 250 = 850\text{ (ml)}$$
Lượng nước còn lại trong bình là:
$$1000 - 850 = 150\text{ (ml)}$$
\textbf{Đáp án của câu hỏi là A}
}
\end{ex}

\begin{ex} %[3H1V07]
Một hình tam giác có độ dài ba cạnh bằng nhau, biết chu vi của hình tam giác đó bằng $24\text{ cm}$. Độ dài mỗi cạnh của hình tam giác đó là
\choice
{$6\text{ cm}$}
{$7\text{ cm}$}
{$9\text{ cm}$}
{\True $8\text{ cm}$}
\loigiai{
Vì hình tam giác có ba cạnh bằng nhau nên độ dài mỗi cạnh là:
$$24 : 3 = 8\text{ (cm)}$$
\textbf{Đáp án của câu hỏi là D}
}
\end{ex}

\begin{ex} %[3H1V07]
Con gà cân nặng $2\text{ kg}$, con vịt nặng hơn con gà $1\text{ kg}$. Con ngỗng cân nặng gấp đôi con vịt. Hỏi con ngỗng nặng hơn con gà bao nhiêu ki-lô-gam?
\choice
{$3\text{ kg}$}
{\True $4\text{ kg}$}
{$5\text{ kg}$}
{$2\text{ kg}$}
\loigiai{
Cân nặng của con vịt là:
$$2 + 1 = 3\text{ (kg)}$$
Cân nặng của con ngỗng là:
$$3 \times 2 = 6\text{ (kg)}$$
Con ngỗng nặng hơn con gà số ki-lô-gam là:
$$6 - 2 = 4\text{ (kg)}$$
\textbf{Đáp án của câu hỏi là B}
}
\end{ex}

\textbf{Phần 2. Trắc nghiệm đúng sai}
\setcounterpageref{ex}{0}

\begin{ex} %[3H1V07]
Xét tính đúng sai của các nhận định sau về đơn vị đo và hình học:
\choiceTF
{\True Mối quan hệ giữa các đơn vị đo độ dài: $1\text{ m} = 10\text{ dm} = 100\text{ cm}$.}
{$1\text{ kg}$ bông nhẹ hơn $1\text{ kg}$ sắt.}
{\True Nửa lít nước tương ứng với $500\text{ ml}$ nước.}
{\True Chu vi hình tam giác có ba cạnh dài $5\text{ cm}, 6\text{ cm}, 7\text{ cm}$ bằng chu vi hình chữ nhật có chiều dài $5\text{ cm}$ và chiều rộng $4\text{ cm}$.}
\loigiai{
\begin{itemize}
    \item Ý a) Đúng: Theo bảng đơn vị đo độ dài.
    \item Ý b) Sai: Cả hai đều nặng 1 kg nên có khối lượng bằng nhau.
    \item Ý c) Đúng: $1\text{ lít} = 1000\text{ ml}$, nửa lít là $1000 : 2 = 500\text{ ml}$.
    \item Ý d) Đúng: Chu vi tam giác: $5 + 6 + 7 = 18\text{ cm}$. Chu vi hình chữ nhật: $(5 + 4) \times 2 = 18\text{ cm}$. Hai chu vi bằng nhau.
\end{itemize}
\begin{itemize}
    \item a) \textbf{Đúng}.
    \item b) \textbf{Sai}.
    \item c) \textbf{Đúng}.
    \item d) \textbf{Đúng}.
\end{itemize}
}
\end{ex}

\begin{ex} %[3H1V07]
Cho đoạn thẳng $AB = 24\text{ cm}$. Lấy điểm $M$ là trung điểm của đoạn thẳng $AB$, sau đó lấy điểm $I$ là trung điểm của đoạn thẳng $AM$:
\choiceTF
{\True Độ dài đoạn thẳng $AM$ bằng $12\text{ cm}$.}
{\True Điểm $M$ nằm giữa hai điểm $A$ và $B$.}
{Độ dài đoạn thẳng $AI$ bằng $8\text{ cm}$.}
{\True Độ dài đoạn thẳng $IB$ bằng $18\text{ cm}$.}
\loigiai{
\begin{itemize}
    \item Ý a) Đúng: $M$ là trung điểm $AB$ nên $AM = 24 : 2 = 12\text{ (cm)}$.
    \item Ý b) Đúng: Trung điểm luôn nằm giữa hai đầu mút của đoạn thẳng.
    \item Ý c) Sai: $I$ là trung điểm $AM$ nên $AI = 12 : 2 = 6\text{ (cm)}$.
    \item Ý d) Đúng: Độ dài đoạn thẳng $IB = IM + MB = 6 + 12 = 18\text{ (cm)}$.
\end{itemize}
\begin{itemize}
    \item a) \textbf{Đúng}.
    \item b) \textbf{Đúng}.
    \item c) \textbf{Sai}.
    \item d) \textbf{Đúng}.
\end{itemize}
}
\end{ex}

\textbf{Phần 3. Trả lời ngắn}
\setcounterpageref{ex}{0}

\begin{ex} %[3H1V07]
Tính chu vi của một hình tam giác có độ dài các cạnh lần lượt là $14\text{ cm}, 18\text{ cm}$ và $23\text{ cm}$.
\shortans{55}
\loigiai{
\textbf{Đáp số:} $55$\\
Chu vi hình tam giác là:
$$14 + 18 + 23 = 55\text{ (cm)}$$
}
\end{ex}

\begin{ex} %[3H1V07]
Một sợi dây thép dài $3\text{ m}$ được uốn thành một hình tam giác có ba cạnh bằng nhau. Hỏi mỗi cạnh của hình tam giác đó dài bao nhiêu đề-xi-mét?
\shortans{10}
\loigiai{
\textbf{Đáp số:} $10$\\
Đổi: $3\text{ m} = 30\text{ dm}$.\\
Mỗi cạnh của hình tam giác dài số đề-xi-mét là:
$$30 : 3 = 10\text{ (dm)}$$
}
\end{ex}

\begin{ex} %[3H1V07]
Một con ngỗng cân nặng $6\text{ kg}$, một con gà cân nặng $2\text{ kg}$. Hỏi con ngỗng nặng gấp mấy lần con gà?
\shortans{3}
\loigiai{
\textbf{Đáp số:} $3$\\
Con ngỗng nặng gấp con gà số lần là:
$$6 : 2 = 3\text{ (lần)}$$
}
\end{ex}

\begin{ex} %[3H1V07]
Thùng thứ nhất chứa $35\text{ lít}$ dầu, thùng thứ hai chứa số lít dầu gấp $2$ lần thùng thứ nhất. Hỏi cả hai thùng chứa tất cả bao nhiêu lít dầu?
\shortans{105}
\loigiai{
\textbf{Đáp số:} $105$\\
Thùng thứ hai chứa số lít dầu là:
$$35 \times 2 = 70\text{ (lít)}$$
Cả hai thùng chứa tất cả số lít dầu là:
$$35 + 70 = 105\text{ (lít)}$$
}
\end{ex}

\begin{ex} %[3H1V07]
Đoạn thẳng $PQ$ dài $32\text{ cm}$. Điểm $K$ là trung điểm của đoạn thẳng $PQ$. Tính độ dài đoạn thẳng $PK$.
\shortans{16}
\loigiai{
\textbf{Đáp số:} $16$\\
Độ dài đoạn thẳng $PK$ là:
$$32 : 2 = 16\text{ (cm)}$$
}
\end{ex}

\begin{ex} %[3H1V07]
Người ta cắm $5$ chiếc cọc thẳng hàng cách đều nhau, khoảng cách giữa hai chiếc cọc liền kề là $3\text{ m}$. Hỏi khoảng cách từ chiếc cọc đầu tiên đến chiếc cọc cuối cùng là bao nhiêu mét?
\shortans{12}
\loigiai{
\textbf{Đáp số:} $12$\\
5 chiếc cọc cắm thẳng hàng chia đoạn đường thành: $5 - 1 = 4$ (khoảng cách).\\
Khoảng cách từ chiếc cọc đầu tiên đến chiếc cọc cuối cùng là:
$$4 \times 3 = 12\text{ (m)}$$
}
\end{ex}

\newpage

\TaoPhieuKiemTra{Chủ đề 1 -- Ôn tập và bổ sung}{08}{Luyện tập chung}{https://loptoanthaylinh.vn/thi-online/toan-3/bai-08}
\textbf{Phần 1. Trắc nghiệm lựa chọn}
\setcounterpageref{ex}{0}

\begin{ex} %[3D1V08]
Một số có ba chữ số có chữ số hàng chục là $5$, chữ số hàng trăm gấp đôi chữ số hàng đơn vị. Biết tổng ba chữ số bằng $11$. Số đó là
\choice
{$254$}
{\True $452$}
{$653$}
{$854$}
\loigiai{
Tổng của chữ số hàng trăm và chữ số hàng đơn vị là:
$$11 - 5 = 6$$
Vì chữ số hàng trăm gấp đôi chữ số hàng đơn vị nên ta có sơ đồ:\\
- Chữ số hàng đơn vị: 1 phần\\
- Chữ số hàng trăm: 2 phần\\
Chữ số hàng đơn vị là:
$$6 : (2 + 1) = 2$$
Chữ số hàng trăm là:
$$2 \times 2 = 4$$
Vậy số cần tìm là $452$.\\
\textbf{Đáp án của câu hỏi là B}
}
\end{ex}

\begin{ex} %[3D1V08]
Tính nhanh giá trị của biểu thức: $B = (150 + 280 + 350) - (80 + 50)$.
\choice
{$600$}
{$620$}
{$680$}
{\True $650$}
\loigiai{
Nhóm các số để trừ tròn trăm:\\
$$B = (150 + 350 - 50) + (280 - 80) = 450 + 200 = 650$$
\textbf{Đáp án của câu hỏi là D}
}
\end{ex}

\begin{ex} %[3D1V08]
Tìm giá trị của $x$, biết: $5 \times x - 18 = 16 + 6$.
\choice
{\True $8$}
{$7$}
{$9$}
{$6$}
\loigiai{
Tính giá trị của vế phải trước:\\
$$5 \times x - 18 = 22$$
$$5 \times x = 22 + 18$$
$$5 \times x = 40$$
$$x = 40 : 5 = 8$$
\textbf{Đáp án của câu hỏi là A}
}
\end{ex}

\begin{ex} %[3D1V08]
Đoạn thẳng thứ nhất dài $45\text{ cm}$, đoạn thẳng thứ hai ngắn hơn đoạn thẳng thứ nhất $15\text{ cm}$. Đoạn thẳng thứ ba dài gấp đôi đoạn thẳng thứ hai. Độ dài đoạn thẳng thứ ba là
\choice
{$50\text{ cm}$}
{$70\text{ cm}$}
{\True $60\text{ cm}$}
{$40\text{ cm}$}
\loigiai{
Độ dài đoạn thẳng thứ hai là:
$$45 - 15 = 30\text{ (cm)}$$
Độ dài đoạn thẳng thứ ba là:
$$30 \times 2 = 60\text{ (cm)}$$
\textbf{Đáp án của câu hỏi là C}
}
\end{ex}

\begin{ex} %[3D1V08]
Cho bốn chữ số $0; 2; 5; 7$. Có thể viết được tất cả bao nhiêu số có 3 chữ số khác nhau và chia hết cho $5$ từ các chữ số đã cho?
\choice
{$8$ số}
{\True $10$ số}
{$12$ số}
{$6$ số}
\loigiai{
Một số chia hết cho 5 khi chữ số hàng đơn vị là 0 hoặc 5.\\
- \textbf{Trường hợp 1:} Chữ số hàng đơn vị là 0.\\
  + Hàng trăm có 3 cách chọn (từ 2; 5; 7).\\
  + Hàng chục có 2 cách chọn.\\
  Số lượng số: $3 \times 2 = 6$ (số).\\
- \textbf{Trường hợp 2:} Chữ số hàng đơn vị là 5.\\
  + Hàng trăm có 2 cách chọn (chọn 2 hoặc 7, vì khác 0).\\
  + Hàng chục có 2 cách chọn (bao gồm cả chữ số 0).\\
  Số lượng số: $2 \times 2 = 4$ (số).\\
Tổng cộng có: $6 + 4 = 10$ số thỏa mãn.\\
\textbf{Đáp án của câu hỏi là B}
}
\end{ex}

\begin{ex} %[3D1V08]
Cho dãy số cách đều: $1; 4; 7; 10; 13; \dots$ Số hạng thứ 10 của dãy số trên là
\choice
{$25$}
{$30$}
{$31$}
{\True $28$}
\loigiai{
Dãy số có khoảng cách giữa hai số liền nhau là: $4 - 1 = 3$ đơn vị.\\
Số hạng thứ 10 của dãy số là:
$$1 + (10 - 1) \times 3 = 1 + 27 = 28$$
\textbf{Đáp án của câu hỏi là D}
}
\end{ex}

\begin{ex} %[3D1V08]
Một hình tam giác có độ dài hai cạnh lần lượt là $16\text{ cm}$ và $24\text{ cm}$. Biết chu vi của hình tam giác đó bằng $60\text{ cm}$. Độ dài cạnh thứ ba của hình tam giác là
\choice
{\True $20\text{ cm}$}
{$18\text{ cm}$}
{$22\text{ cm}$}
{$25\text{ cm}$}
\loigiai{
Tổng độ dài hai cạnh đã biết là:
$$16 + 24 = 40\text{ (cm)}$$
Độ dài cạnh thứ ba của hình tam giác là:
$$60 - 40 = 20\text{ (cm)}$$
\textbf{Đáp án của câu hỏi là A}
}
\end{ex}

\begin{ex} %[3D1V08]
Trong một phép trừ có số bị trừ bằng $450$. Nếu giữ nguyên số bị trừ và bớt số trừ đi $25$ đơn vị thì được hiệu mới bằng $185$. Số trừ ban đầu là
\choice
{$280$}
{$300$}
{\True $290$}
{$265$}
\loigiai{
Khi bớt ở số trừ đi 25 đơn vị thì hiệu tăng thêm 25 đơn vị.\\
Hiệu ban đầu của phép trừ là:
$$185 - 25 = 160$$
Số trừ ban đầu là:
$$450 - 160 = 290$$
\textbf{Đáp án của câu hỏi là C}
}
\end{ex}

\begin{ex} %[3D1V08]
Mai mua $4$ chiếc bút chì, mỗi chiếc có giá $5$ nghìn đồng và mua thêm $1$ quyển vở giá $8$ nghìn đồng. Mai đưa cho cô bán hàng tờ tiền $50$ nghìn đồng. Hỏi cô bán hàng phải trả lại Mai bao nhiêu nghìn đồng?
\choice
{$20$ nghìn đồng}
{\True $22$ nghìn đồng}
{$24$ nghìn đồng}
{$18$ nghìn đồng}
\loigiai{
Số tiền Mai mua 4 chiếc bút chì là:
$$4 \times 5 = 20\text{ (nghìn đồng)}$$
Tổng số tiền Mai mua đồ là:
$$20 + 8 = 28\text{ (nghìn đồng)}$$
Số tiền cô bán hàng trả lại Mai là:
$$50 - 28 = 22\text{ (nghìn đồng)}$$
\textbf{Đáp án của câu hỏi là B}
}
\end{ex}

\begin{ex} %[3D1V08]
Hai số có tổng bằng $450$ và hiệu của chúng bằng $50$. Số lớn là
\choice
{$200$}
{$300$}
{$225$}
{\True $250$}
\loigiai{
Số lớn là:
$$(450 + 50) : 2 = 250$$
\textbf{Đáp án của câu hỏi là D}
}
\end{ex}

\begin{ex} %[3D1V08]
Một đoạn đường dài $32\text{ m}$, người ta trồng các cây cách đều nhau $4\text{ m}$, ở cả hai đầu đường đều có trồng cây. Hỏi đoạn đường đó trồng tất cả bao nhiêu cây?
\choice
{\True $9$ cây}
{$8$ cây}
{$10$ cây}
{$7$ cây}
\loigiai{
Số khoảng cách giữa các cây là:
$$32 : 4 = 8\text{ (khoảng)}$$
Vì cả hai đầu đường đều có cây nên số cây trồng được là:
$$8 + 1 = 9\text{ (cây)}$$
\textbf{Đáp án của câu hỏi là A}
}
\end{ex}

\begin{ex} %[3D1V08]
Một số tự nhiên khi chia cho $4$ thì được thương bằng $9$ và số dư là số dư lớn nhất có thể có. Khi lấy số đó chia cho $5$ thì số dư là
\choice
{$2$}
{$3$}
{\True $4$}
{$1$}
\loigiai{
Trong phép chia cho 4, số dư lớn nhất có thể là $3$.\\
Số tự nhiên đó là:
$$4 \times 9 + 3 = 36 + 3 = 39$$
Lấy 39 chia cho 5 ta có:
$$39 = 5 \times 7 + 4$$
Vậy khi chia cho 5 thì số dư là $4$.\\
\textbf{Đáp án của câu hỏi là C}
}
\end{ex}

\textbf{Phần 2. Trắc nghiệm đúng sai}
\setcounterpageref{ex}{0}

\begin{ex} %[3D1V08]
Xét tính đúng sai của các nhận định sau:
\choiceTF
{\True Trong một phép chia cho $5$, số dư lớn nhất có thể có bằng $4$.}
{\True Chu vi của hình tam giác có ba cạnh bằng nhau mỗi cạnh $6\text{ cm}$ bằng chu vi của hình chữ nhật có chiều dài $5\text{ cm}$ và chiều rộng $4\text{ cm}$.}
{Khi cộng một số tự nhiên bất kì với $0$ thì luôn được kết quả bằng $0$.}
{\True Tích của số $2$ với bất kì một số lẻ nào cũng luôn là một số chẵn.}
\loigiai{
\begin{itemize}
    \item Ý a) Đúng: Số dư luôn nhỏ hơn số chia nên số dư lớn nhất là 4.
    \item Ý b) Đúng: Chu vi tam giác: $6 \times 3 = 18\text{ (cm)}$. Chu vi hình chữ nhật: $(5 + 4) \times 2 = 18\text{ (cm)}$.
    \item Ý c) Sai: Khi cộng một số với 0 thì bằng chính số đó ($a + 0 = a$).
    \item Ý d) Đúng: Số 2 nhân với bất kì số nào cũng cho kết quả là số chẵn.
\end{itemize}
\begin{itemize}
    \item a) \textbf{Đúng}.
    \item b) \textbf{Đúng}.
    \item c) \textbf{Sai}.
    \item d) \textbf{Đúng}.
\end{itemize}
}
\end{ex}

\begin{ex} %[3D1V08]
Một đội công nhân phải sửa đoạn đường dài $300\text{ m}$ trong ba ngày. Ngày thứ nhất sửa được $80\text{ m}$, ngày thứ hai sửa được số mét đường gấp đôi ngày thứ nhất. Xét tính đúng sai:
\choiceTF
{\True Ngày thứ hai đội công nhân sửa được $160\text{ m}$ đường.}
{\True Sau hai ngày đầu, đội công nhân đã sửa được tất cả $240\text{ m}$ đường.}
{Ngày thứ ba đội công nhân phải sửa $70\text{ m}$ đường nữa mới hoàn thành công việc.}
{\True Ngày thứ hai đội công nhân sửa được nhiều hơn ngày thứ ba là $100\text{ m}$ đường.}
\loigiai{
\begin{itemize}
    \item Ý a) Đúng: Ngày thứ hai sửa được: $80 \times 2 = 160\text{ (m)}$.
    \item Ý b) Đúng: Sau hai ngày đầu sửa được: $80 + 160 = 240\text{ (m)}$.
    \item Ý c) Sai: Ngày thứ ba cần sửa là: $300 - 240 = 60\text{ (m)}$ (không phải 70 m).
    \item Ý d) Đúng: Ngày thứ hai sửa nhiều hơn ngày thứ ba là: $160 - 60 = 100\text{ (m)}$.
\end{itemize}
\begin{itemize}
    \item a) \textbf{Đúng}.
    \item b) \textbf{Đúng}.
    \item c) \textbf{Sai}.
    \item d) \textbf{Đúng}.
\end{itemize}
}
\end{ex}

\textbf{Phần 3. Trả lời ngắn}
\setcounterpageref{ex}{0}

\begin{ex} %[3D1V08]
Tính nhanh giá trị của biểu thức: $135 + 246 + 165 + 154$.
\shortans{700}
\loigiai{
\textbf{Đáp số:} $700$\\
Nhóm các số tròn trăm:\\
$$(135 + 165) + (246 + 154) = 300 + 400 = 700$$
}
\end{ex}

\begin{ex} %[3D1V08]
Tìm giá trị của $x$, biết: $x : 3 + 45 = 54$.
\shortans{27}
\loigiai{
\textbf{Đáp số:} $27$\\
Xem $x : 3$ là số hạng chưa biết:\\
$$x : 3 = 54 - 45$$
$$x : 3 = 9$$
$$x = 9 \times 3 = 27$$
}
\end{ex}

\begin{ex} %[3D1V08]
Một khúc gỗ dài $20\text{ m}$ được người ta cưa thành các đoạn ngắn dài $4\text{ m}$. Hỏi người thợ phải cưa tất cả bao nhiêu lần?
\shortans{4}
\loigiai{
\textbf{Đáp số:} $4$\\
Số đoạn gỗ cưa được là: $20 : 4 = 5$ (đoạn).\\
Số lần cưa cần thực hiện là: $5 - 1 = 4$ (lần cưa).
}
\end{ex}

\begin{ex} %[3D1V08]
Hai thùng đựng tất cả $160\text{ lít}$ nước mắm. Thùng thứ nhất đựng nhiều hơn thùng thứ hai $20\text{ lít}$. Hỏi thùng thứ nhất đựng bao nhiêu lít nước mắm?
\shortans{90}
\loigiai{
\textbf{Đáp số:} $90$\\
Thùng thứ nhất đựng số lít nước mắm là:
$$(160 + 20) : 2 = 90\text{ (lít)}$$
}
\end{ex}

\begin{ex} %[3D1V08]
Cho ba chữ số $2; 4; 7$. Tìm hiệu giữa số lớn nhất có 3 chữ số khác nhau và số bé nhất có 3 chữ số khác nhau được lập từ ba chữ số đã cho.
\shortans{495}
\loigiai{
\textbf{Đáp số:} $495$\\
- Số lớn nhất có 3 chữ số khác nhau là: $742$.\\
- Số bé nhất có 3 chữ số khác nhau là: $247$.\\
Hiệu giữa hai số là:
$$742 - 247 = 495$$
}
\end{ex}

\begin{ex} %[3D1V08]
Dãy số chẵn liên tiếp $2; 4; 6; 8; \dots; 40$ có tất cả bao nhiêu số hạng?
\shortans{20}
\loigiai{
\textbf{Đáp số:} $20$\\
Số số hạng của dãy là:
$$(40 - 2) : 2 + 1 = 38 : 2 + 1 = 20\text{ (số hạng)}$$
}
\end{ex}

\newpage

% =========================================================================
% CHỦ ĐỀ 2 - BẢNG NHÂN, BẢNG CHIA
% =========================================================================
\section{Bảng nhân, bảng chia}

\TaoPhieuKiemTra{Chủ đề 2 -- Bảng nhân, bảng chia}{09}{Bảng nhân 6, bảng chia 6}{https://loptoanthaylinh.vn/thi-online/toan-3/bai-09}
\textbf{Phần 1. Trắc nghiệm lựa chọn}
\setcounterpageref{ex}{0}

\begin{ex} %[3D2V09]
Tính giá trị của biểu thức: $A = 6 	imes 7 - 48 : 6 + 15$.
\choice
{$45$}
{$51$}
{\True $49$}
{$47$}
\loigiai{
Thực hiện nhân chia trước, cộng trừ sau:\
$$A = 42 - 8 + 15 = 34 + 15 = 49$$
\textbf{Đáp án của câu hỏi là C}
}
\end{ex}

\begin{ex} %[3D2V09]
Một con kiến có $6$ cái chân. Một đàn kiến gồm $8$ con kiến thợ và $4$ con kiến chúa. Hỏi đàn kiến đó có tất cả bao nhiêu cái chân?
\choice
{\True $72$ cái chân}
{$68$ cái chân}
{$76$ cái chân}
{$64$ cái chân}
\loigiai{
Tổng số con kiến trong đàn là: $8 + 4 = 12$ (con).\
Đàn kiến có tất cả số cái chân là: $12 \times 6 = 72$ (cái chân).
\textbf{Đáp án của câu hỏi là A}
}
\end{ex}

\begin{ex} %[3D2V09]
Tìm giá trị của $x$, biết: $x : 6 + 25 = 32$.
\choice
{$36$}
{$48$}
{$30$}
{\True $42$}
\loigiai{
Ta có: $x : 6 = 32 - 25 = 7 \Rightarrow x = 7 \times 6 = 42$.
\textbf{Đáp án của câu hỏi là D}
}
\end{ex}

\begin{ex} %[3D2V09]
Một hình lục giác đều có 6 cạnh bằng nhau. Biết chu vi của hình lục giác đều đó là $54\text{ cm}$. Độ dài mỗi cạnh của hình lục giác đều là
\choice
{$8\text{ cm}$}
{\True $9\text{ cm}$}
{$7\text{ cm}$}
{$6\text{ cm}$}
\loigiai{
Vì hình lục giác đều có 6 cạnh bằng nhau nên độ dài mỗi cạnh là: $54 : 6 = 9\text{ (cm)}$.
\textbf{Đáp án của câu hỏi là B}
}
\end{ex}

\begin{ex} %[3D2V09]
Một cửa hàng có $60\text{ kg}$ gạo, người ta đóng vào các túi loại $6\text{ kg}$. Sau đó cửa hàng bán đi $4$ túi gạo. Hỏi cửa hàng còn lại bao nhiêu ki-lô-gam gạo?
\choice
{$24\text{ kg}$}
{$40\text{ kg}$}
{\True $36\text{ kg}$}
{$30\text{ kg}$}
\loigiai{
Cửa hàng đóng được số túi gạo là: $60 : 6 = 10$ (túi).\
Sau khi bán 4 túi, cửa hàng còn lại: $10 - 4 = 6$ (túi).\
Số ki-lô-gam gạo còn lại là: $6 \times 6 = 36\text{ (kg)}$.
\textbf{Đáp án của câu hỏi là C}
}
\end{ex}

\begin{ex} %[3D2V09]
Tính nhanh giá trị của biểu thức: $6 \times 2 + 6 \times 3 + 6 \times 5$.
\choice
{\True $60$}
{$54$}
{$66$}
{$72$}
\loigiai{
Đặt thừa số chung: $6 \times (2 + 3 + 5) = 6 \times 10 = 60$.
\textbf{Đáp án của câu hỏi là A}
}
\end{ex}

\begin{ex} %[3D2V09]
Trong một phép chia cho $6$, thương bằng $8$ và số dư là số dư lớn nhất có thể có. Số bị chia là
\choice
{$48$}
{$51$}
{$52$}
{\True $53$}
\loigiai{
Vì số chia là 6 nên số dư lớn nhất là 5.\
Số bị chia là: $6 \times 8 + 5 = 48 + 5 = 53$.
\textbf{Đáp án của câu hỏi là D}
}
\end{ex}

\begin{ex} %[3D2V09]
Điền dấu thích hợp vào chỗ chấm: $6 \times 8 - 14 \dots 54 : 6 \times 4$.
\choice
{$>$}
{\True $<$}
{$=$}
{Không so sánh được}
\loigiai{
Vế trái: $6 \times 8 - 14 = 48 - 14 = 34$.\
Vế phải: $54 : 6 \times 4 = 9 \times 4 = 36$.\
Vì $34 < 36$ nên điền dấu $<$. 
\textbf{Đáp án của câu hỏi là B}
}
\end{ex}

\begin{ex} %[3D2V09]
Một tổ công nhân sản xuất được $48$ chi tiết máy trong một tuần và chia đều thành 6 ca làm việc. Mỗi ca lại chia đều sản phẩm cho 4 công nhân. Hỏi mỗi công nhân làm được bao nhiêu chi tiết máy?
\choice
{$3$ chi tiết máy}
{$4$ chi tiết máy}
{\True $2$ chi tiết máy}
{$1$ chi tiết máy}
\loigiai{
Mỗi ca làm được: $48 : 6 = 8$ chi tiết máy.\
Mỗi công nhân làm được: $8 : 4 = 2$ chi tiết máy.
\textbf{Đáp án của câu hỏi là C}
}
\end{ex}

\begin{ex} %[3D2V09]
Tìm $x$ biết: $(x - 12) \times 6 = 42$.
\choice
{\True $19$}
{$18$}
{$20$}
{$21$}
\loigiai{
Ta có: $x - 12 = 42 : 6 = 7 \Rightarrow x = 7 + 12 = 19$.
\textbf{Đáp án của câu hỏi là A}
}
\end{ex}

\begin{ex} %[3D2V09]
Cho dãy số: $6; 12; 18; 24; \dots$ Số hạng thứ 9 của dãy số trên là
\choice
{$48$}
{$60$}
{$56$}
{\True $54$}
\loigiai{
Dãy số gồm các bội số của 6. Số hạng thứ 9 là: $6 \times 9 = 54$.
\textbf{Đáp án của câu hỏi là D}
}
\end{ex}

\begin{ex} %[3D2V09]
Mai có $18$ quyển vở, số vở của Lan gấp $6$ lần số vở của Mai. Hỏi cả hai bạn có tất cả bao nhiêu quyển vở?
\choice
{$108$ quyển vở}
{\True $126$ quyển vở}
{$120$ quyển vở}
{$114$ quyển vở}
\loigiai{
Số vở của Lan là: $18 \times 6 = 108$ (quyển).\
Cả hai bạn có: $18 + 108 = 126$ (quyển).
\textbf{Đáp án của câu hỏi là B}
}
\end{ex}

\textbf{Phần 2. Trắc nghiệm đúng sai}
\setcounterpageref{ex}{0}

\begin{ex} %[3D2V09]
Xét tính đúng sai của các khẳng định về phép nhân và phép chia với 6:
\choiceTF
{\True Tích của số 6 với một số tự nhiên bất kì luôn luôn là một số chẵn.}
{Trong phép chia cho 6, số dư lớn nhất có thể là 6.}
{\True Một số tự nhiên chia hết cho 6 thì chắc chắn chia hết cho 3.}
{\True Giá trị của $6 \times 4$ bằng giá trị của $4 \times 6$.}
\loigiai{
\begin{itemize}
    \item Ý a) Đúng: Vì 6 là số chẵn nên $6 \times n$ luôn là số chẵn.
    \item Ý b) Sai: Số dư luôn bé hơn số chia, lớn nhất là 5.
    \item Ý c) Đúng: Vì $6 = 2 \times 3$ nên số chia hết cho 6 cũng chia hết cho 3.
    \item Ý d) Đúng: Theo tính chất giao hoán của phép nhân.
\end{itemize}
\begin{itemize}
    \item a) \textbf{Đúng}.
    \item b) \textbf{Sai}.
    \item c) \textbf{Đúng}.
    \item d) \textbf{Đúng}.
\end{itemize}
}
\end{ex}

\begin{ex} %[3D2V09]
Một tấm vải dài $54\text{ m}$ được cắt thành các mảnh nhỏ dài $6\text{ m}$:
\choiceTF
{\True Người ta cắt được tất cả 9 mảnh vải nhỏ.}
{Người thợ phải thực hiện đúng 9 nhát cắt.}
{\True Tổng độ dài của 5 mảnh vải nhỏ là $30\text{ m}$.}
{\True Nếu bớt ở tấm vải đi $6\text{ m}$ thì số mảnh vải cắt được là 8 mảnh.}
\loigiai{
\begin{itemize}
    \item Ý a) Đúng: $54 : 6 = 9$ mảnh.
    \item Ý b) Sai: Để cắt thành 9 mảnh chỉ cần $9 - 1 = 8$ nhát cắt.
    \item Ý c) Đúng: $5 \times 6 = 30\text{ m}$.
    \item Ý d) Đúng: $(54 - 6) : 6 = 48 : 6 = 8$ mảnh.
\end{itemize}
\begin{itemize}
    \item a) \textbf{Đúng}.
    \item b) \textbf{Sai}.
    \item c) \textbf{Đúng}.
    \item d) \textbf{Đúng}.
\end{itemize}
}
\end{ex}

\textbf{Phần 3. Trả lời ngắn}
\setcounterpageref{ex}{0}

\begin{ex} %[3D2V09]
Tính giá trị của biểu thức: $6 \times 8 - 30 : 6$.
\shortans{43}
\loigiai{
\textbf{Đáp số:} $43$\\
$6 \times 8 - 30 : 6 = 48 - 5 = 43$.
}
\end{ex}

\begin{ex} %[3D2V09]
Một đàn ong có $7$ con. Hỏi đàn ong đó có tất cả bao nhiêu cái chân?
\shortans{42}
\loigiai{
\textbf{Đáp số:} $42$\\
Mỗi con ong có 6 chân: $7 \times 6 = 42$ cái chân.
}
\end{ex}

\begin{ex} %[3D2V09]
Tìm số bị chia trong phép chia cho 6, biết thương bằng 9 và số dư bằng 4.
\shortans{58}
\loigiai{
\textbf{Đáp số:} $58$\\
$6 \times 9 + 4 = 54 + 4 = 58$.
}
\end{ex}

\begin{ex} %[3D2V09]
Lớp 3A có 36 học sinh được chia đều vào 6 nhóm. Sau đó mỗi nhóm cử 2 bạn tham gia đá bóng. Hỏi có tất cả bao nhiêu bạn tham gia đá bóng?
\shortans{12}
\loigiai{
\textbf{Đáp số:} $12$\\
Mỗi nhóm có $36 : 6 = 6$ bạn. Có $6 \times 2 = 12$ bạn tham gia.
}
\end{ex}

\begin{ex} %[3D2V09]
Tìm giá trị của $x$, biết: $x : 6 = 12 - 4$.
\shortans{48}
\loigiai{
\textbf{Đáp số:} $48$\\
$x : 6 = 8 \Rightarrow x = 8 \times 6 = 48$.
}
\end{ex}

\begin{ex} %[3D2V09]
Tính nhanh: $6 \times 1 + 6 \times 2 + 6 \times 3 + 6 \times 4$.
\shortans{60}
\loigiai{
\textbf{Đáp số:} $60$\\
$6 \times (1 + 2 + 3 + 4) = 6 \times 10 = 60$.
}
\end{ex}

\newpage

\TaoPhieuKiemTra{Chủ đề 2 -- Bảng nhân, bảng chia}{10}{Bảng nhân 7, bảng chia 7}{https://loptoanthaylinh.vn/thi-online/toan-3/bai-10}
\textbf{Phần 1. Trắc nghiệm lựa chọn}
\setcounterpageref{ex}{0}

\begin{ex} %[3D2V10]
Tính giá trị của biểu thức: $B = 7 \times 8 - 42 : 7 + 16$.
\choice
{$64$}
{\True $66$}
{$68$}
{$62$}
\loigiai{
Thực hiện nhân chia trước: $B = 56 - 6 + 16 = 50 + 16 = 66$.
\textbf{Đáp án của câu hỏi là B}
}
\end{ex}

\begin{ex} %[3D2V10]
Một tuần lễ có $7$ ngày. Hỏi trong $5$ tuần lễ và $3$ ngày có tất cả bao nhiêu ngày?
\choice
{$35$ ngày}
{$36$ ngày}
{$40$ ngày}
{\True $38$ ngày}
\loigiai{
Số ngày trong 5 tuần là: $5 \times 7 = 35$ (ngày).\
Tổng số ngày là: $35 + 3 = 38$ (ngày).
\textbf{Đáp án của câu hỏi là D}
}
\end{ex}

\begin{ex} %[3D2V10]
Tìm $x$ biết: $x : 7 + 18 = 27$.
\choice
{\True $63$}
{$56$}
{$70$}
{$49$}
\loigiai{
Ta có: $x : 7 = 27 - 18 = 9 \Rightarrow x = 9 \times 7 = 63$.
\textbf{Đáp án của câu hỏi là A}
}
\end{ex}

\begin{ex} %[3D2V10]
Một cuộn dây thép dài $56\text{ m}$. Người ta cắt đi $\frac{1}{7}$ cuộn dây thép đó. Hỏi cuộn dây thép còn lại bao nhiêu mét?
\choice
{$8\text{ m}$}
{$50\text{ m}$}
{\True $48\text{ m}$}
{$42\text{ m}$}
\loigiai{
Đoạn dây thép đã cắt là: $56 : 7 = 8\text{ (m)}$.\
Đoạn dây còn lại dài: $56 - 8 = 48\text{ (m)}$.
\textbf{Đáp án của câu hỏi là C}
}
\end{ex}

\begin{ex} %[3D2V10]
Trong một phép chia cho $7$, thương bằng $9$ và số dư là số dư lớn nhất có thể có. Số bị chia của phép chia đó là
\choice
{$63$}
{\True $69$}
{$68$}
{$70$}
\loigiai{
Vì số chia là 7 nên số dư lớn nhất là 6.\
Số bị chia là: $7 \times 9 + 6 = 63 + 6 = 69$.
\textbf{Đáp án của câu hỏi là B}
}
\end{ex}

\begin{ex} %[3D2V10]
Tháng Chạp (tháng 12) có ngày mùng 1 là thứ Sáu. Hỏi ngày 22 tháng 12 cùng năm đó là thứ mấy?
\choice
{Thứ Bảy}
{Thứ Năm}
{Chủ Nhật}
{\True Thứ Sáu}
\loigiai{
Khoảng cách giữa ngày 22 và ngày 1 là: $22 - 1 = 21$ (ngày).\
Vì $21 : 7 = 3$ (tuần chẵn) nên ngày 22 cũng là thứ Sáu.
\textbf{Đáp án của câu hỏi là D}
}
\end{ex}

\begin{ex} %[3D2V10]
Tính nhanh: $7 \times 3 + 7 \times 4 + 7 \times 3$.
\choice
{\True $70$}
{$63$}
{$77$}
{$84$}
\loigiai{
Đặt thừa số chung: $7 \times (3 + 4 + 3) = 7 \times 10 = 70$.
\textbf{Đáp án của câu hỏi là A}
}
\end{ex}

\begin{ex} %[3D2V10]
Một mảnh đất hình chữ nhật có chiều rộng là $7\text{ m}$, chiều dài gấp $4$ lần chiều rộng. Chu vi của mảnh đất đó là
\choice
{$65\text{ m}$}
{$60\text{ m}$}
{\True $70\text{ m}$}
{$75\text{ m}$}
\loigiai{
Chiều dài mảnh đất là: $7 \times 4 = 28\text{ (m)}$.\
Chu vi là: $(28 + 7) \times 2 = 35 \times 2 = 70\text{ (m)}$.
\textbf{Đáp án của câu hỏi là C}
}
\end{ex}

\begin{ex} %[3D2V10]
Tìm $x$ biết: $(x - 15) \times 7 = 35$.
\choice
{$18$}
{\True $20$}
{$22$}
{$25$}
\loigiai{
Ta có: $x - 15 = 35 : 7 = 5 \Rightarrow x = 5 + 15 = 20$.
\textbf{Đáp án của câu hỏi là B}
}
\end{ex}

\begin{ex} %[3D2V10]
Mẹ mua $7$ vỉ trứng gà, mỗi vỉ có $6$ quả và mua thêm $15$ quả trứng vịt. Hỏi mẹ đã mua tất cả bao nhiêu quả trứng?
\choice
{$52$ quả}
{$55$ quả}
{$60$ quả}
{\True $57$ quả}
\loigiai{
Số trứng gà là: $7 \times 6 = 42$ quả.\
Tổng số trứng là: $42 + 15 = 57$ quả.
\textbf{Đáp án của câu hỏi là D}
}
\end{ex}

\begin{ex} %[3D2V10]
Cho dãy số cách đều: $7; 14; 21; 28; \dots$ Số hạng thứ 8 của dãy số là
\choice
{\True $56$}
{$49$}
{$63$}
{$54$}
\loigiai{
Số hạng thứ 8 là: $7 \times 8 = 56$.
\textbf{Đáp án của câu hỏi là A}
}
\end{ex}

\begin{ex} %[3D2V10]
Bác Nam có $63$ cây cam trồng đều thành 7 hàng. Bác bán đi 2 hàng cam. Hỏi bác Nam còn lại bao nhiêu cây cam?
\choice
{$40$ cây cam}
{$42$ cây cam}
{\True $45$ cây cam}
{$48$ cây cam}
\loigiai{
Mỗi hàng có số cây cam là: $63 : 7 = 9$ (cây).\
Số hàng còn lại là: $7 - 2 = 5$ (hàng).\
Số cây cam còn lại là: $5 \times 9 = 45$ (cây).
\textbf{Đáp án của câu hỏi là C}
}
\end{ex}

\textbf{Phần 2. Trắc nghiệm đúng sai}
\setcounterpageref{ex}{0}

\begin{ex} %[3D2V10]
Xét tính đúng sai của các nhận định sau về bảng nhân và chia 7:
\choiceTF
{\True Một tuần lễ có đúng 7 ngày.}
{Trong phép chia cho 7, số dư lớn nhất có thể có là 7.}
{\True Nếu lấy tích của 7 với một số rồi chia cho 7 thì được chính số đó.}
{\True Tổng $7 \times 4 + 7 \times 6 = 70$.}
\loigiai{
\begin{itemize}
    \item Ý a) Đúng: Quy ước thời gian 1 tuần = 7 ngày.
    \item Ý b) Sai: Số dư lớn nhất là 6.
    \item Ý c) Đúng: $7 \times a : 7 = a$.
    \item Ý d) Đúng: $7 \times (4 + 6) = 7 \times 10 = 70$.
\end{itemize}
\begin{itemize}
    \item a) \textbf{Đúng}.
    \item b) \textbf{Sai}.
    \item c) \textbf{Đúng}.
    \item d) \textbf{Đúng}.
\end{itemize}
}
\end{ex}

\begin{ex} %[3D2V10]
Một lớp học có 35 học sinh chia đều thành 7 nhóm học tập:
\choiceTF
{\True Mỗi nhóm học tập có 5 bạn học sinh.}
{\True Nếu 3 nhóm cử đi trực nhật thì có 15 bạn trực nhật.}
{Số học sinh của 4 nhóm là 25 bạn.}
{\True Tổng số học sinh gấp 7 lần số học sinh của mỗi nhóm.}
\loigiai{
\begin{itemize}
    \item Ý a) Đúng: $35 : 7 = 5$ bạn.
    \item Ý b) Đúng: $3 \times 5 = 15$ bạn.
    \item Ý c) Sai: $4 \times 5 = 20$ bạn.
    \item Ý d) Đúng: $35 : 5 = 7$ lần.
\end{itemize}
\begin{itemize}
    \item a) \textbf{Đúng}.
    \item b) \textbf{Đúng}.
    \item c) \textbf{Sai}.
    \item d) \textbf{Đúng}.
\end{itemize}
}
\end{ex}

\textbf{Phần 3. Trả lời ngắn}
\setcounterpageref{ex}{0}

\begin{ex} %[3D2V10]
Tính giá trị biểu thức: $7 \times 9 - 35 : 7$.
\shortans{58}
\loigiai{
\textbf{Đáp số:} $58$\\
$7 \times 9 - 35 : 7 = 63 - 5 = 58$.
}
\end{ex}

\begin{ex} %[3D2V10]
Tìm số bị chia trong phép chia cho 7, biết thương là 8 và số dư là 5.
\shortans{61}
\loigiai{
\textbf{Đáp số:} $61$\\
$7 \times 8 + 5 = 56 + 5 = 61$.
}
\end{ex}

\begin{ex} %[3D2V10]
Hôm nay là thứ Ba ngày 5. Hỏi thứ Ba tuần sau là ngày mấy?
\shortans{12}
\loigiai{
\textbf{Đáp số:} $12$\\
Thứ Ba tuần sau là ngày: $5 + 7 = 12$.
}
\end{ex}

\begin{ex} %[3D2V10]
Một bao xi măng nặng 49 kg được chia đều vào 7 túi nhỏ. Hỏi 4 túi nhỏ như thế nặng bao nhiêu ki-lô-gam?
\shortans{28}
\loigiai{
\textbf{Đáp số:} $28$\\
Mỗi túi nặng: $49 : 7 = 7\text{ kg}$. 4 túi nặng: $4 \times 7 = 28\text{ kg}$.
}
\end{ex}

\begin{ex} %[3D2V10]
Tìm $x$ biết: $7 \times x = 90 - 41$.
\shortans{7}
\loigiai{
\textbf{Đáp số:} $7$\\
$7 \times x = 49 \Rightarrow x = 49 : 7 = 7$.
}
\end{ex}

\begin{ex} %[3D2V10]
Tính nhanh: $7 \times 2 + 7 \times 3 + 7 \times 5$.
\shortans{70}
\loigiai{
\textbf{Đáp số:} $70$\\
$7 \times (2 + 3 + 5) = 7 \times 10 = 70$.
}
\end{ex}

\newpage

\TaoPhieuKiemTra{Chủ đề 2 -- Bảng nhân, bảng chia}{11}{Bảng nhân 8, bảng chia 8}{https://loptoanthaylinh.vn/thi-online/toan-3/bai-11}
\textbf{Phần 1. Trắc nghiệm lựa chọn}
\setcounterpageref{ex}{0}

\begin{ex} %[3D2V11]
Tính giá trị của biểu thức: $C = 8 \times 9 - 56 : 8 + 24$.
\choice
{\True $89$}
{$85$}
{$91$}
{$87$}
\loigiai{
Thực hiện nhân chia trước: $C = 72 - 7 + 24 = 65 + 24 = 89$.
\textbf{Đáp án của câu hỏi là A}
}
\end{ex}

\begin{ex} %[3D2V11]
Một con bạch tuộc có $8$ cái tua. Hỏi một đàn gồm $9$ con bạch tuộc có tất cả bao nhiêu cái tua?
\choice
{$64$ cái tua}
{$80$ cái tua}
{\True $72$ cái tua}
{$68$ cái tua}
\loigiai{
Số cái tua là: $9 \times 8 = 72$ (cái tua).
\textbf{Đáp án của câu hỏi là C}
}
\end{ex}

\begin{ex} %[3D2V11]
Tìm $x$ biết: $x : 8 + 35 = 43$.
\choice
{$56$}
{\True $64$}
{$72$}
{$48$}
\loigiai{
Ta có: $x : 8 = 43 - 35 = 8 \Rightarrow x = 8 \times 8 = 64$.
\textbf{Đáp án của câu hỏi là B}
}
\end{ex}

\begin{ex} %[3D2V11]
Trong một phép chia cho $8$, thương bằng $7$ và số dư là số dư lớn nhất có thể có. Số bị chia của phép chia đó là
\choice
{$56$}
{$61$}
{$62$}
{\True $63$}
\loigiai{
Vì số chia là 8 nên số dư lớn nhất là 7.\
Số bị chia là: $8 \times 7 + 7 = 56 + 7 = 63$.
\textbf{Đáp án của câu hỏi là D}
}
\end{ex}

\begin{ex} %[3D2V11]
Tính nhanh: $8 \times 2 + 8 \times 3 + 8 \times 5$.
\choice
{\True $80$}
{$72$}
{$88$}
{$96$}
\loigiai{
Đặt thừa số chung: $8 \times (2 + 3 + 5) = 8 \times 10 = 80$.
\textbf{Đáp án của câu hỏi là A}
}
\end{ex}

\begin{ex} %[3D2V11]
Một cuộn dây cáp dài $64\text{ m}$. Người ta cắt thành $8$ đoạn dây dài bằng nhau. Hỏi $5$ đoạn dây như thế dài bao nhiêu mét?
\choice
{$35\text{ m}$}
{$45\text{ m}$}
{\True $40\text{ m}$}
{$48\text{ m}$}
\loigiai{
Mỗi đoạn dây dài: $64 : 8 = 8\text{ (m)}$.\
5 đoạn dây dài: $8 \times 5 = 40\text{ (m)}$.
\textbf{Đáp án của câu hỏi là C}
}
\end{ex}

\begin{ex} %[3D2V11]
Khối lập phương có đúng 8 đỉnh. Hỏi 6 khối lập phương như thế có tất cả bao nhiêu đỉnh?
\choice
{$40$ đỉnh}
{\True $48$ đỉnh}
{$54$ đỉnh}
{$56$ đỉnh}
\loigiai{
Số đỉnh của 6 khối lập phương là: $8 \times 6 = 48$ (đỉnh).
\textbf{Đáp án của câu hỏi là B}
}
\end{ex}

\begin{ex} %[3D2V11]
Tìm $x$ biết: $(x - 24) \times 8 = 48$.
\choice
{$28$}
{$32$}
{$35$}
{\True $30$}
\loigiai{
Ta có: $x - 24 = 48 : 8 = 6 \Rightarrow x = 6 + 24 = 30$.
\textbf{Đáp án của câu hỏi là D}
}
\end{ex}

\begin{ex} %[3D2V11]
Một cửa hàng có $72\text{ lít}$ mật ong chia đều vào 8 can. Cửa hàng đã bán được 3 can mật ong. Hỏi cửa hàng còn lại bao nhiêu lít mật ong?
\choice
{\True $45\text{ lít}$}
{$40\text{ lít}$}
{$50\text{ lít}$}
{$48\text{ lít}$}
\loigiai{
Mỗi can chứa: $72 : 8 = 9\text{ lít}$.\
Số can còn lại: $8 - 3 = 5$ can.\
Số lít còn lại: $5 \times 9 = 45\text{ lít}$.
\textbf{Đáp án của câu hỏi là A}
}
\end{ex}

\begin{ex} %[3D2V11]
Cho dãy số: $8; 16; 24; 32; \dots$ Số hạng thứ 9 của dãy số trên là
\choice
{$64$}
{$80$}
{\True $72$}
{$68$}
\loigiai{
Số hạng thứ 9 là: $8 \times 9 = 72$.
\textbf{Đáp án của câu hỏi là C}
}
\end{ex}

\begin{ex} %[3D2V11]
Một người thợ làm việc $8$ giờ mỗi ngày và làm việc trong 6 ngày một tuần. Hỏi trong 2 tuần người đó làm việc tất cả bao nhiêu giờ?
\choice
{$88$ giờ}
{\True $96$ giờ}
{$92$ giờ}
{$100$ giờ}
\loigiai{
Mỗi tuần người đó làm: $8 \times 6 = 48$ giờ.\
Trong 2 tuần làm: $48 \times 2 = 96$ giờ.
\textbf{Đáp án của câu hỏi là B}
}
\end{ex}

\begin{ex} %[3D2V11]
Điền dấu thích hợp: $8 \times 7 - 16 \dots 40 : 8 \times 8$.
\choice
{$>$}
{$<$}
{Không so sánh được}
{\True $=$}
\loigiai{
Vế trái: $8 \times 7 - 16 = 56 - 16 = 40$.\
Vế phải: $40 : 8 \times 8 = 5 \times 8 = 40$.\
Do đó điền dấu $=$.
\textbf{Đáp án của câu hỏi là D}
}
\end{ex}

\textbf{Phần 2. Trắc nghiệm đúng sai}
\setcounterpageref{ex}{0}

\begin{ex} %[3D2V11]
Xét tính đúng sai của các nhận định sau về phép nhân và chia với 8:
\choiceTF
{\True Tích của số 8 với một số tự nhiên bất kì luôn là số chẵn.}
{Số chia hết cho 8 thì không chia hết cho 4.}
{\True Trong phép chia cho 8, số dư lớn nhất có thể có là 7.}
{\True Giá trị của $8 \times 5$ bằng một nửa của $8 \times 10$.}
\loigiai{
\begin{itemize}
    \item Ý a) Đúng: Vì 8 là số chẵn.
    \item Ý b) Sai: Mọi số chia hết cho 8 đều chia hết cho 4 vì $8 = 4 \times 2$.
    \item Ý c) Đúng: Số dư luôn bé hơn số chia.
    \item Ý d) Đúng: Vì $40$ bằng một nửa của $80$.
\end{itemize}
\begin{itemize}
    \item a) \textbf{Đúng}.
    \item b) \textbf{Sai}.
    \item c) \textbf{Đúng}.
    \item d) \textbf{Đúng}.
\end{itemize}
}
\end{ex}

\begin{ex} %[3D2V11]
Một thùng kẹo có 48 cái kẹo chia đều vào 8 hộp nhỏ:
\choiceTF
{\True Mỗi hộp nhỏ có 6 cái kẹo.}
{\True Lấy ra 3 hộp kẹo thì được tất cả 18 cái kẹo.}
{Số kẹo trong 5 hộp nhỏ là 32 cái kẹo.}
{\True Tổng số kẹo gấp 8 lần số kẹo của một hộp nhỏ.}
\loigiai{
\begin{itemize}
    \item Ý a) Đúng: $48 : 8 = 6$ cái kẹo.
    \item Ý b) Đúng: $3 \times 6 = 18$ cái kẹo.
    \item Ý c) Sai: $5 \times 6 = 30$ cái kẹo.
    \item Ý d) Đúng: $48 : 6 = 8$ lần.
\end{itemize}
\begin{itemize}
    \item a) \textbf{Đúng}.
    \item b) \textbf{Đúng}.
    \item c) \textbf{Sai}.
    \item d) \textbf{Đúng}.
\end{itemize}
}
\end{ex}

\textbf{Phần 3. Trả lời ngắn}
\setcounterpageref{ex}{0}

\begin{ex} %[3D2V11]
Tính giá trị biểu thức: $8 \times 6 + 32 : 8$.
\shortans{52}
\loigiai{
\textbf{Đáp số:} $52$\\
$8 \times 6 + 32 : 8 = 48 + 4 = 52$.
}
\end{ex}

\begin{ex} %[3D2V11]
Một con nhện có 8 cái chân. Hỏi 7 con nhện có tất cả bao nhiêu cái chân?
\shortans{56}
\loigiai{
\textbf{Đáp số:} $56$\\
$7 \times 8 = 56$ cái chân.
}
\end{ex}

\begin{ex} %[3D2V11]
Tìm số bị chia trong phép chia cho 8, biết thương là 6 và số dư là 3.
\shortans{51}
\loigiai{
\textbf{Đáp số:} $51$\\
$8 \times 6 + 3 = 48 + 3 = 51$.
}
\end{ex}

\begin{ex} %[3D2V11]
Tìm $x$ biết: $x : 8 = 15 - 9$.
\shortans{48}
\loigiai{
\textbf{Đáp số:} $48$\\
$x : 8 = 6 \Rightarrow x = 6 \times 8 = 48$.
}
\end{ex}

\begin{ex} %[3D2V11]
Một lớp học có 32 học sinh chia đều thành 8 tổ. Hỏi 5 tổ có bao nhiêu học sinh?
\shortans{20}
\loigiai{
\textbf{Đáp số:} $20$\\
Mỗi tổ có: $32 : 8 = 4$ bạn. 5 tổ có: $5 \times 4 = 20$ bạn.
}
\end{ex}

\begin{ex} %[3D2V11]
Tính nhanh: $8 \times 1 + 8 \times 2 + 8 \times 3 + 8 \times 4$.
\shortans{80}
\loigiai{
\textbf{Đáp số:} $80$\\
$8 \times (1 + 2 + 3 + 4) = 8 \times 10 = 80$.
}
\end{ex}

\newpage

\TaoPhieuKiemTra{Chủ đề 2 -- Bảng nhân, bảng chia}{12}{Bảng nhân 9, bảng chia 9}{https://loptoanthaylinh.vn/thi-online/toan-3/bai-12}
\textbf{Phần 1. Trắc nghiệm lựa chọn}
\setcounterpageref{ex}{0}

\begin{ex} %[3D2V12]
Tính giá trị của biểu thức: $D = 9 \times 8 - 63 : 9 + 17$.
\choice
{$78$}
{$84$}
{$80$}
{\True $82$}
\loigiai{
Thực hiện nhân chia trước: $D = 72 - 7 + 17 = 65 + 17 = 82$.
\textbf{Đáp án của câu hỏi là D}
}
\end{ex}

\begin{ex} %[3D2V12]
Một lớp học có $36$ học sinh chia đều thành 9 nhóm. Hỏi 5 nhóm học tập đó có tất cả bao nhiêu học sinh?
\choice
{$18$ học sinh}
{\True $20$ học sinh}
{$24$ học sinh}
{$16$ học sinh}
\loigiai{
Mỗi nhóm có: $36 : 9 = 4$ học sinh.\
5 nhóm có: $5 \times 4 = 20$ học sinh.
\textbf{Đáp án của câu hỏi là B}
}
\end{ex}

\begin{ex} %[3D2V12]
Tìm $x$ biết: $x : 9 + 28 = 35$.
\choice
{$54$}
{$72$}
{\True $63$}
{$45$}
\loigiai{
Ta có: $x : 9 = 35 - 28 = 7 \Rightarrow x = 7 \times 9 = 63$.
\textbf{Đáp án của câu hỏi là C}
}
\end{ex}

\begin{ex} %[3D2V12]
Trong một phép chia cho $9$, thương bằng $8$ và số dư là số dư lớn nhất có thể có. Số bị chia là
\choice
{\True $80$}
{$79$}
{$81$}
{$78$}
\loigiai{
Vì số chia là 9 nên số dư lớn nhất là 8.\
Số bị chia là: $9 \times 8 + 8 = 72 + 8 = 80$.
\textbf{Đáp án của câu hỏi là A}
}
\end{ex}

\begin{ex} %[3D2V12]
Tính nhanh giá trị biểu thức: $9 \times 3 + 9 \times 4 + 9 \times 3$.
\choice
{$81$}
{$99$}
{$85$}
{\True $90$}
\loigiai{
Đặt thừa số chung: $9 \times (3 + 4 + 3) = 9 \times 10 = 90$.
\textbf{Đáp án của câu hỏi là D}
}
\end{ex}

\begin{ex} %[3D2V12]
Bác Hòa có $72\text{ kg}$ bưởi chia đều vào 9 thùng. Bác đã bán đi 4 thùng bưởi. Hỏi bác Hòa còn lại bao nhiêu ki-lô-gam bưởi?
\choice
{$36\text{ kg}$}
{\True $40\text{ kg}$}
{$45\text{ kg}$}
{$32\text{ kg}$}
\loigiai{
Mỗi thùng có: $72 : 9 = 8\text{ kg}$.\
Số thùng còn lại: $9 - 4 = 5$ thùng.\
Số bưởi còn lại: $5 \times 8 = 40\text{ kg}$.
\textbf{Đáp án của câu hỏi là B}
}
\end{ex}

\begin{ex} %[3D2V12]
Cho dãy số: $9; 18; 27; 36; \dots$ Số hạng thứ 9 của dãy số trên là
\choice
{$72$}
{$90$}
{\True $81$}
{$78$}
\loigiai{
Số hạng thứ 9 là: $9 \times 9 = 81$.
\textbf{Đáp án của câu hỏi là C}
}
\end{ex}

\begin{ex} %[3D2V12]
Tìm $x$ biết: $(x - 18) \times 9 = 54$.
\choice
{\True $24$}
{$22$}
{$26$}
{$25$}
\loigiai{
Ta có: $x - 18 = 54 : 9 = 6 \Rightarrow x = 6 + 18 = 24$.
\textbf{Đáp án của câu hỏi là A}
}
\end{ex}

\begin{ex} %[3D2V12]
Một đàn gà có $9$ con gà trống. Số con gà mái gấp $5$ lần số con gà trống. Hỏi đàn gà có tất cả bao nhiêu con gà?
\choice
{$45$ con gà}
{$50$ con gà}
{$60$ con gà}
{\True $54$ con gà}
\loigiai{
Số gà mái là: $9 \times 5 = 45$ con.\
Tổng số gà là: $9 + 45 = 54$ con.
\textbf{Đáp án của câu hỏi là D}
}
\end{ex}

\begin{ex} %[3D2V12]
Điền dấu thích hợp: $9 \times 6 - 15 \dots 81 : 9 \times 4$.
\choice
{$<$}
{\True $>$}
{$=$}
{Không so sánh được}
\loigiai{
Vế trái: $9 \times 6 - 15 = 54 - 15 = 39$.\
Vế phải: $81 : 9 \times 4 = 9 \times 4 = 36$.\
Vì $39 > 36$ nên điền dấu $>$. 
\textbf{Đáp án của câu hỏi là B}
}
\end{ex}

\begin{ex} %[3D2V12]
Một thanh sắt dài $81\text{ dm}$ được cắt thành các đoạn dài $9\text{ dm}$. Hỏi người thợ phải cưa tất cả bao nhiêu lần?
\choice
{$9$ lần cưa}
{$7$ lần cưa}
{\True $8$ lần cưa}
{$10$ lần cưa}
\loigiai{
Số đoạn sắt là: $81 : 9 = 9$ (đoạn).\
Số lần cưa cần thực hiện: $9 - 1 = 8$ (lần).
\textbf{Đáp án của câu hỏi là C}
}
\end{ex}

\begin{ex} %[3D2V12]
Tìm số tự nhiên $a$ biết: $9 \times 7 < a < 9 \times 9$ và $a$ là số tròn chục.
\choice
{\True $70$}
{$80$}
{$60$}
{$90$}
\loigiai{
Ta có: $63 < a < 81$. Các số tròn chục thỏa mãn là $70$ và $80$. Nếu là số tròn chục duy nhất: đổi đề hoặc chọn 70 (khi $a$ có hàng đơn vị là 0 và hàng chục là 7 $\Rightarrow a = 70$).
\textbf{Đáp án của câu hỏi là A}
}
\end{ex}

\textbf{Phần 2. Trắc nghiệm đúng sai}
\setcounterpageref{ex}{0}

\begin{ex} %[3D2V12]
Xét tính đúng sai của các nhận định sau về số học và phép nhân 9:
\choiceTF
{\True Tổng các chữ số của tích bất kì trong bảng nhân 9 luôn bằng 9 hoặc chia hết cho 9.}
{Mọi số chia hết cho 9 đều là số lẻ.}
{\True Tích $9 \times 6$ bằng tích $6 \times 9$.}
{\True Trong phép chia cho 9, số dư lớn nhất có thể là 8.}
\loigiai{
\begin{itemize}
    \item Ý a) Đúng: Ví dụ: 18 (1+8=9), 27 (2+7=9), 81 (8+1=9), 90 (9+0=9).
    \item Ý b) Sai: Ví dụ 18, 36, 54 chia hết cho 9 nhưng là số chẵn.
    \item Ý c) Đúng: Tính chất giao hoán của phép nhân.
    \item Ý d) Đúng: Số dư luôn bé hơn số chia.
\end{itemize}
\begin{itemize}
    \item a) \textbf{Đúng}.
    \item b) \textbf{Sai}.
    \item c) \textbf{Đúng}.
    \item d) \textbf{Đúng}.
\end{itemize}
}
\end{ex}

\begin{ex} %[3D2V12]
Một người mua 9 gói kẹo, mỗi gói có 8 cái kẹo. Người đó đã cho đi 24 cái kẹo:
\choiceTF
{\True Ban đầu người đó có tất cả 72 cái kẹo.}
{\True Sau khi cho đi, người đó còn lại 48 cái kẹo.}
{Số kẹo còn lại đủ để chia đều cho 5 bạn mỗi bạn 10 cái.}
{\True Số kẹo đã cho đi tương ứng với 3 gói kẹo ban đầu.}
\loigiai{
\begin{itemize}
    \item Ý a) Đúng: $9 \times 8 = 72$ cái kẹo.
    \item Ý b) Đúng: $72 - 24 = 48$ cái kẹo.
    \item Ý c) Sai: Cần $5 \times 10 = 50$ cái, mà chỉ còn 48 cái.
    \item Ý d) Đúng: $3 \times 8 = 24$ cái kẹo.
\end{itemize}
\begin{itemize}
    \item a) \textbf{Đúng}.
    \item b) \textbf{Đúng}.
    \item c) \textbf{Sai}.
    \item d) \textbf{Đúng}.
\end{itemize}
}
\end{ex}

\textbf{Phần 3. Trả lời ngắn}
\setcounterpageref{ex}{0}

\begin{ex} %[3D2V12]
Tính giá trị của biểu thức: $9 \times 7 - 45 : 9$.
\shortans{58}
\loigiai{
\textbf{Đáp số:} $58$\\
$9 \times 7 - 45 : 9 = 63 - 5 = 58$.
}
\end{ex}

\begin{ex} %[3D2V12]
Tìm số bị chia trong phép chia cho 9, biết thương là 7 và số dư là 6.
\shortans{69}
\loigiai{
\textbf{Đáp số:} $69$\\
$9 \times 7 + 6 = 63 + 6 = 69$.
}
\end{ex}

\begin{ex} %[3D2V12]
Một người có 63 quả cam chia đều vào 9 túi. Hỏi 6 túi như thế có bao nhiêu quả cam?
\shortans{42}
\loigiai{
\textbf{Đáp số:} $42$\\
Mỗi túi có $63 : 9 = 7$ quả. 6 túi có $6 \times 7 = 42$ quả.
}
\end{ex}

\begin{ex} %[3D2V12]
Tìm $x$ biết: $x : 9 = 16 - 8$.
\shortans{72}
\loigiai{
\textbf{Đáp số:} $72$\\
$x : 9 = 8 \Rightarrow x = 8 \times 9 = 72$.
}
\end{ex}

\begin{ex} %[3D2V12]
Một hình vuông có cạnh dài 9 cm. Tính chu vi hình vuông đó.
\shortans{36}
\loigiai{
\textbf{Đáp số:} $36$\\
Chu vi hình vuông là: $9 \times 4 = 36\text{ cm}$.
}
\end{ex}

\begin{ex} %[3D2V12]
Tính nhanh: $9 \times 1 + 9 \times 2 + 9 \times 3 + 9 \times 4$.
\shortans{90}
\loigiai{
\textbf{Đáp số:} $90$\\
$9 \times (1 + 2 + 3 + 4) = 9 \times 10 = 90$.
}
\end{ex}

\newpage

\TaoPhieuKiemTra{Chủ đề 2 -- Bảng nhân, bảng chia}{13}{Tìm thành phần trong phép nhân, phép chia}{https://loptoanthaylinh.vn/thi-online/toan-3/bai-13}
\textbf{Phần 1. Trắc nghiệm lựa chọn}
\setcounterpageref{ex}{0}

\begin{ex} %[3D2V13]
Tìm thừa số $x$, biết: $x \times 6 = 18 + 24$.
\choice
{$6$}
{$8$}
{\True $7$}
{$9$}
\loigiai{
Tính vế phải: $x \times 6 = 42 \Rightarrow x = 42 : 6 = 7$.
\textbf{Đáp án của câu hỏi là C}
}
\end{ex}

\begin{ex} %[3D2V13]
Tìm số bị chia $x$, biết: $x : 7 = 35 : 5$.
\choice
{$42$}
{\True $49$}
{$56$}
{$35$}
\loigiai{
Tính vế phải: $x : 7 = 7 \Rightarrow x = 7 \times 7 = 49$.
\textbf{Đáp án của câu hỏi là B}
}
\end{ex}

\begin{ex} %[3D2V13]
Tìm số chia $x$, biết: $54 : x = 24 : 4$.
\choice
{$8$}
{$7$}
{$6$}
{\True $9$}
\loigiai{
Tính vế phải: $54 : x = 6 \Rightarrow x = 54 : 6 = 9$.
\textbf{Đáp án của câu hỏi là D}
}
\end{ex}

\begin{ex} %[3D2V13]
Trong một phép nhân, nếu một thừa số gấp lên $3$ lần và thừa số kia giữ nguyên thì tích mới thay đổi như thế nào?
\choice
{\True Gấp lên $3$ lần}
{Giảm đi $3$ lần}
{Tăng thêm $3$ đơn vị}
{Không thay đổi}
\loigiai{
Khi nhân một thừa số với 3 thì tích cũng gấp lên 3 lần.
\textbf{Đáp án của câu hỏi là A}
}
\end{ex}

\begin{ex} %[3D2V13]
Tìm $x$ biết: $x \times 8 + 16 = 80$.
\choice
{$7$}
{$9$}
{\True $8$}
{$6$}
\loigiai{
Xem $x \times 8$ là số hạng chưa biết: $x \times 8 = 80 - 16 = 64 \Rightarrow x = 64 : 8 = 8$.
\textbf{Đáp án của câu hỏi là C}
}
\end{ex}

\begin{ex} %[3D2V13]
Tìm $x$ biết: $x : 9 - 3 = 5$.
\choice
{$63$}
{\True $72$}
{$81$}
{$54$}
\loigiai{
Xem $x : 9$ là số bị trừ: $x : 9 = 5 + 3 = 8 \Rightarrow x = 8 \times 9 = 72$.
\textbf{Đáp án của câu hỏi là B}
}
\end{ex}

\begin{ex} %[3D2V13]
Trong một phép chia hết có thương bằng $6$. Nếu số bị chia gấp lên $2$ lần và giữ nguyên số chia thì thương mới là
\choice
{$8$}
{$3$}
{$18$}
{\True $12$}
\loigiai{
Khi số bị chia gấp lên 2 lần thì thương cũng gấp lên 2 lần: $6 \times 2 = 12$.
\textbf{Đáp án của câu hỏi là D}
}
\end{ex}

\begin{ex} %[3D2V13]
Một số nhân với $7$ rồi cộng với $15$ thì được $50$. Số đó là
\choice
{\True $5$}
{$6$}
{$4$}
{$7$}
\loigiai{
Gọi số đó là $x$: $x \times 7 + 15 = 50 \Rightarrow x \times 7 = 35 \Rightarrow x = 5$.
\textbf{Đáp án của câu hỏi là A}
}
\end{ex}

\begin{ex} %[3D2V13]
Tìm $x$ biết: $72 : x + 12 = 20$.
\choice
{$8$}
{$7$}
{\True $9$}
{$6$}
\loigiai{
Ta có: $72 : x = 20 - 12 = 8 \Rightarrow x = 72 : 8 = 9$.
\textbf{Đáp án của câu hỏi là C}
}
\end{ex}

\begin{ex} %[3D2V13]
Tích của hai số bằng $48$. Nếu giữ nguyên thừa số thứ nhất và giảm thừa số thứ hai đi $2$ lần thì tích mới là
\choice
{$16$}
{\True $24$}
{$32$}
{$96$}
\loigiai{
Khi giảm một thừa số đi 2 lần thì tích cũng giảm đi 2 lần: $48 : 2 = 24$.
\textbf{Đáp án của câu hỏi là B}
}
\end{ex}

\begin{ex} %[3D2V13]
Tìm $x$ biết: $(x + 5) \times 7 = 63$.
\choice
{$3$}
{$5$}
{$6$}
{\True $4$}
\loigiai{
Ta có: $x + 5 = 63 : 7 = 9 \Rightarrow x = 9 - 5 = 4$.
\textbf{Đáp án của câu hỏi là D}
}
\end{ex}

\begin{ex} %[3D2V13]
Một thùng chứa có $64\text{ lít}$ nước. Người ta dùng các can loại $8\text{ lít}$ để múc hết nước trong thùng. Hỏi cần ít nhất bao nhiêu chiếc can như vậy?
\choice
{\True $8$ chiếc can}
{$7$ chiếc can}
{$9$ chiếc can}
{$6$ chiếc can}
\loigiai{
Số chiếc can cần dùng là: $64 : 8 = 8$ chiếc can.
\textbf{Đáp án của câu hỏi là A}
}
\end{ex}

\textbf{Phần 2. Trắc nghiệm đúng sai}
\setcounterpageref{ex}{0}

\begin{ex} %[3D2V13]
Xét tính đúng sai của các quy tắc tìm thành phần phép tính:
\choiceTF
{\True Muốn tìm thừa số chưa biết, ta lấy tích chia cho thừa số đã biết.}
{Muốn tìm số chia, ta lấy thương nhân với số bị chia.}
{\True Muốn tìm số bị chia trong phép chia hết, ta lấy thương nhân với số chia.}
{\True Trong phép chia có dư, số bị chia bằng thương nhân số chia cộng số dư.}
\loigiai{
\begin{itemize}
    \item Ý a) Đúng: Quy tắc cơ bản của phép nhân.
    \item Ý b) Sai: Muốn tìm số chia, ta lấy số bị chia chia cho thương.
    \item Ý c) Đúng: Quy tắc tìm số bị chia.
    \item Ý d) Đúng: Công thức phép chia có dư.
\end{itemize}
\begin{itemize}
    \item a) \textbf{Đúng}.
    \item b) \textbf{Sai}.
    \item c) \textbf{Đúng}.
    \item d) \textbf{Đúng}.
\end{itemize}
}
\end{ex}

\begin{ex} %[3D2V13]
Cho bài toán tìm $x$: $x \times 8 - 14 = 50$. Xét tính đúng sai:
\choiceTF
{\True Giá trị của biểu thức $x \times 8$ bằng $64$.}
{Giá trị của $x$ tìm được là $9$.}
{\True Số $x$ tìm được là một số chẵn chia hết cho 4.}
{\True Nếu thay số $50$ thành $58$ thì giá trị của $x$ sẽ bằng $9$.}
\loigiai{
\begin{itemize}
    \item Ý a) Đúng: $50 + 14 = 64$.
    \item Ý b) Sai: $x = 64 : 8 = 8$ (không phải 9).
    \item Ý c) Đúng: $8$ là số chẵn chia hết cho 4.
    \item Ý d) Đúng: $x \times 8 = 58 + 14 = 72 \Rightarrow x = 9$.
\end{itemize}
\begin{itemize}
    \item a) \textbf{Đúng}.
    \item b) \textbf{Sai}.
    \item c) \textbf{Đúng}.
    \item d) \textbf{Đúng}.
\end{itemize}
}
\end{ex}

\textbf{Phần 3. Trả lời ngắn}
\setcounterpageref{ex}{0}

\begin{ex} %[3D2V13]
Tìm thừa số $x$, biết: $x \times 7 = 84 - 28$.
\shortans{8}
\loigiai{
\textbf{Đáp số:} $8$\\
$x \times 7 = 56 \Rightarrow x = 56 : 7 = 8$.
}
\end{ex}

\begin{ex} %[3D2V13]
Tìm số bị chia $x$, biết: $x : 6 = 15 - 8$.
\shortans{42}
\loigiai{
\textbf{Đáp số:} $42$\\
$x : 6 = 7 \Rightarrow x = 7 \times 6 = 42$.
}
\end{ex}

\begin{ex} %[3D2V13]
Tìm số chia $x$, biết: $72 : x = 27 : 3$.
\shortans{8}
\loigiai{
\textbf{Đáp số:} $8$\\
$72 : x = 9 \Rightarrow x = 72 : 9 = 8$.
}
\end{ex}

\begin{ex} %[3D2V13]
Một số nhân với 8 rồi trừ đi 15 thì được 49. Tìm số đó.
\shortans{8}
\loigiai{
\textbf{Đáp số:} $8$\\
$x \times 8 - 15 = 49 \Rightarrow x \times 8 = 64 \Rightarrow x = 8$.
}
\end{ex}

\begin{ex} %[3D2V13]
Tìm $x$ biết: $(x - 6) \times 9 = 54$.
\shortans{12}
\loigiai{
\textbf{Đáp số:} $12$\\
$x - 6 = 54 : 9 = 6 \Rightarrow x = 6 + 6 = 12$.
}
\end{ex}

\begin{ex} %[3D2V13]
Trong một phép chia có thương là 8, số chia là 7 và số dư là 4. Tìm số bị chia.
\shortans{60}
\loigiai{
\textbf{Đáp số:} $60$\\
$8 \times 7 + 4 = 56 + 4 = 60$.
}
\end{ex}

\newpage

\TaoPhieuKiemTra{Chủ đề 2 -- Bảng nhân, bảng chia}{14}{Một phần mấy (1/2, 1/3, 1/4, 1/5...)}{https://loptoanthaylinh.vn/thi-online/toan-3/bai-14}
\textbf{Phần 1. Trắc nghiệm lựa chọn}
\setcounterpageref{ex}{0}

\begin{ex} %[3D2V14]
Một đàn gà có $36$ con. $\frac{1}{4}$ số con gà trong đàn là bao nhiêu con?
\choice
{$8$ con gà}
{\True $9$ con gà}
{$10$ con gà}
{$12$ con gà}
\loigiai{
$\frac{1}{4}$ số gà là: $36 : 4 = 9$ (con).
\textbf{Đáp án của câu hỏi là B}
}
\end{ex}

\begin{ex} %[3D2V14]
Một rổ có $48$ quả cam. Mẹ đã bán đi $\frac{1}{6}$ số quả cam đó. Hỏi mẹ còn lại bao nhiêu quả cam?
\choice
{\True $40$ quả cam}
{$36$ quả cam}
{$42$ quả cam}
{$38$ quả cam}
\loigiai{
Số cam mẹ đã bán là: $48 : 6 = 8$ (quả).\
Số cam còn lại là: $48 - 8 = 40$ (quả).
\textbf{Đáp án của câu hỏi là A}
}
\end{ex}

\begin{ex} %[3D2V14]
Tìm một số biết rằng $\frac{1}{7}$ của số đó bằng $9$. Số đó là
\choice
{$54$}
{$56$}
{\True $63$}
{$72$}
\loigiai{
Số cần tìm là: $9 \times 7 = 63$.
\textbf{Đáp án của câu hỏi là C}
}
\end{ex}

\begin{ex} %[3D2V14]
Một đoạn thẳng $AB$ dài $32\text{ cm}$. Điểm $M$ thuộc đoạn thẳng $AB$ sao cho đoạn thẳng $AM$ bằng $\frac{1}{4}$ đoạn thẳng $AB$. Độ dài đoạn thẳng $MB$ là
\choice
{$8\text{ cm}$}
{$16\text{ cm}$}
{$20\text{ cm}$}
{\True $24\text{ cm}$}
\loigiai{
Độ dài đoạn thẳng $AM$ là: $32 : 4 = 8\text{ (cm)}$.\
Độ dài đoạn thẳng $MB$ là: $32 - 8 = 24\text{ (cm)}$.
\textbf{Đáp án của câu hỏi là D}
}
\end{ex}

\begin{ex} %[3D2V14]
Lớp 3A có $40$ học sinh. Trong đó có $\frac{1}{5}$ số học sinh là học sinh xuất sắc, $\frac{1}{4}$ số học sinh là học sinh tiêu biểu. Hỏi lớp 3A có tất cả bao nhiêu bạn đạt danh hiệu khen thưởng trên?
\choice
{$16$ bạn}
{\True $18$ bạn}
{$20$ bạn}
{$15$ bạn}
\loigiai{
Số học sinh xuất sắc là: $40 : 5 = 8$ (bạn).\
Số học sinh tiêu biểu là: $40 : 4 = 10$ (bạn).\
Tổng số bạn được khen thưởng là: $8 + 10 = 18$ (bạn).
\textbf{Đáp án của câu hỏi là B}
}
\end{ex}

\begin{ex} %[3D2V14]
So sánh $\frac{1}{6}$ của $54\text{ kg}$ và $\frac{1}{8}$ của $64\text{ kg}$:
\choice
{\True $\frac{1}{6}$ của $54\text{ kg}$ lớn hơn}
{$\frac{1}{8}$ của $64\text{ kg}$ lớn hơn}
{Hai giá trị bằng nhau}
{Không so sánh được}
\loigiai{
$\frac{1}{6}$ của $54\text{ kg}$ là: $54 : 6 = 9\text{ kg}$.\
$\frac{1}{8}$ của $64\text{ kg}$ là: $64 : 8 = 8\text{ kg}$.\
Vì $9\text{ kg} > 8\text{ kg}$ nên $\frac{1}{6}$ của $54\text{ kg}$ lớn hơn.
\textbf{Đáp án của câu hỏi là A}
}
\end{ex}

\begin{ex} %[3D2V14]
Một ngày có $24$ giờ. Hỏi $\frac{1}{3}$ ngày có bao nhiêu giờ?
\choice
{$6$ giờ}
{$12$ giờ}
{\True $8$ giờ}
{$4$ giờ}
\loigiai{
Số giờ trong $\frac{1}{3}$ ngày là: $24 : 3 = 8$ giờ.
\textbf{Đáp án của câu hỏi là C}
}
\end{ex}

\begin{ex} %[3D2V14]
Một cửa hàng có $56\text{ m}$ vải hoa. Cửa hàng đã bán được $\frac{1}{7}$ số vải hoa đó và bán thêm $12\text{ m}$ nữa. Hỏi cửa hàng đã bán tất cả bao nhiêu mét vải hoa?
\choice
{$18\text{ m}$}
{$16\text{ m}$}
{$22\text{ m}$}
{\True $20\text{ m}$}
\loigiai{
$\frac{1}{7}$ số vải hoa là: $56 : 7 = 8\text{ (m)}$.\
Tổng số vải đã bán là: $8 + 12 = 20\text{ (m)}$.
\textbf{Đáp án của câu hỏi là D}
}
\end{ex}

\begin{ex} %[3D2V14]
Bao gạo thứ nhất nặng $45\text{ kg}$. Bao gạo thứ hai nặng bằng $\frac{1}{5}$ bao gạo thứ nhất. Hỏi cả hai bao nặng bao nhiêu ki-lô-gam?
\choice
{$50\text{ kg}$}
{\True $54\text{ kg}$}
{$48\text{ kg}$}
{$56\text{ kg}$}
\loigiai{
Bao thứ hai nặng: $45 : 5 = 9\text{ kg}$.\
Cả hai bao nặng: $45 + 9 = 54\text{ kg}$.
\textbf{Đáp án của câu hỏi là B}
}
\end{ex}

\begin{ex} %[3D2V14]
Hình vuông có chu vi bằng $36\text{ cm}$. Độ dài một cạnh của hình vuông là
\choice
{\True $9\text{ cm}$}
{$6\text{ cm}$}
{$8\text{ cm}$}
{$12\text{ cm}$}
\loigiai{
Cạnh hình vuông bằng $\frac{1}{4}$ chu vi: $36 : 4 = 9\text{ cm}$.
\textbf{Đáp án của câu hỏi là A}
}
\end{ex}

\begin{ex} %[3D2V14]
Một tổ công nhân nhận sửa một đoạn đường dài $72\text{ m}$. Ngày thứ nhất sửa được $\frac{1}{9}$ đoạn đường, ngày thứ hai sửa được $\frac{1}{8}$ đoạn đường. Hỏi sau hai ngày tổ công nhân đã sửa được bao nhiêu mét đường?
\choice
{$15\text{ m}$}
{$16\text{ m}$}
{\True $17\text{ m}$}
{$18\text{ m}$}
\loigiai{
Ngày 1 sửa: $72 : 9 = 8\text{ m}$.\
Ngày 2 sửa: $72 : 8 = 9\text{ m}$.\
Cả hai ngày sửa: $8 + 9 = 17\text{ m}$.
\textbf{Đáp án của câu hỏi là C}
}
\end{ex}

\begin{ex} %[3D2V14]
Một thùng dầu chứa $54\text{ lít}$. Người ta lấy ra $\frac{1}{6}$ số dầu rồi lại lấy thêm $5\text{ lít}$ nữa. Số lít dầu còn lại trong thùng là
\choice
{$35\text{ lít}$}
{$45\text{ lít}$}
{$42\text{ lít}$}
{\True $40\text{ lít}$}
\loigiai{
Lần thứ nhất lấy: $54 : 6 = 9\text{ lít}$.\
Cả hai lần lấy: $9 + 5 = 14\text{ lít}$.\
Còn lại: $54 - 14 = 40\text{ lít}$.
\textbf{Đáp án của câu hỏi là D}
}
\end{ex}

\textbf{Phần 2. Trắc nghiệm đúng sai}
\setcounterpageref{ex}{0}

\begin{ex} %[3D2V14]
Xét tính đúng sai của các nhận định về phân số một phần mấy:
\choiceTF
{\True Để tìm một phần ba của số $A$, ta lấy $A$ chia cho 3.}
{Một phần năm của 45 lớn hơn một phần tư của 40.}
{\True Nếu $\frac{1}{8}$ số bi của Nam là 6 viên thì Nam có tất cả 48 viên bi.}
{\True Một giờ có 60 phút thì $\frac{1}{6}$ giờ tương ứng với 10 phút.}
\loigiai{
\begin{itemize}
    \item Ý a) Đúng: Quy tắc tìm một phần mấy của một số.
    \item Ý b) Sai: $\frac{1}{5}$ của $45$ là $9$, $\frac{1}{4}$ của $40$ là $10$. Vì $9 < 10$ nên nhận định Sai.
    \item Ý c) Đúng: Số bi của Nam là: $6 \times 8 = 48$ viên.
    \item Ý d) Đúng: $60 : 6 = 10$ phút.
\end{itemize}
\begin{itemize}
    \item a) \textbf{Đúng}.
    \item b) \textbf{Sai}.
    \item c) \textbf{Đúng}.
    \item d) \textbf{Đúng}.
\end{itemize}
}
\end{ex}

\begin{ex} %[3D2V14]
Một trang trại có 72 con cừu. Bác nông dân đã bán đi $\frac{1}{8}$ đàn cừu:
\choiceTF
{\True Bác nông dân đã bán đi 9 con cừu.}
{\True Đàn cừu còn lại 63 con.}
{Số cừu còn lại gấp 6 lần số cừu đã bán.}
{\True Nếu bác bán thêm 7 con cừu nữa thì số cừu còn lại là 56 con.}
\loigiai{
\begin{itemize}
    \item Ý a) Đúng: $72 : 8 = 9$ con.
    \item Ý b) Đúng: $72 - 9 = 63$ con.
    \item Ý c) Sai: $63 : 9 = 7$ lần (không phải 6 lần).
    \item Ý d) Đúng: $63 - 7 = 56$ con.
\end{itemize}
\begin{itemize}
    \item a) \textbf{Đúng}.
    \item b) \textbf{Đúng}.
    \item c) \textbf{Sai}.
    \item d) \textbf{Đúng}.
\end{itemize}
}
\end{ex}

\textbf{Phần 3. Trả lời ngắn}
\setcounterpageref{ex}{0}

\begin{ex} %[3D2V14]
Tìm $\frac{1}{7}$ của $63\text{ kg}$.
\shortans{9}
\loigiai{
\textbf{Đáp số:} $9$\\
$63 : 7 = 9\text{ kg}$.
}
\end{ex}

\begin{ex} %[3D2V14]
Tìm một số biết rằng $\frac{1}{6}$ của số đó bằng 8.
\shortans{48}
\loigiai{
\textbf{Đáp số:} $48$\\
$8 \times 6 = 48$.
}
\end{ex}

\begin{ex} %[3D2V14]
Một mảnh vườn hình chữ nhật có chu vi bằng $40\text{ m}$. Nửa chu vi mảnh vườn đó là bao nhiêu mét?
\shortans{20}
\loigiai{
\textbf{Đáp số:} $20$\\
Nửa chu vi bằng $\frac{1}{2}$ chu vi: $40 : 2 = 20\text{ m}$.
}
\end{ex}

\begin{ex} %[3D2V14]
Lớp 3B có 36 học sinh. Số học sinh nam bằng $\frac{1}{4}$ số học sinh cả lớp. Hỏi lớp 3B có bao nhiêu học sinh nữ?
\shortans{27}
\loigiai{
\textbf{Đáp số:} $27$\\
Số nam: $36 : 4 = 9$ bạn. Số nữ: $36 - 9 = 27$ bạn.
}
\end{ex}

\begin{ex} %[3D2V14]
Một bao gạo nặng 54 kg. Người ta đã dùng hết $\frac{1}{9}$ bao gạo đó. Hỏi trong bao còn lại bao nhiêu ki-lô-gam gạo?
\shortans{48}
\loigiai{
\textbf{Đáp số:} $48$\\
Đã dùng: $54 : 9 = 6\text{ kg}$. Còn lại: $54 - 6 = 48\text{ kg}$.
}
\end{ex}

\begin{ex} %[3D2V14]
Một cuộn dây dài 45 m. Người ta cắt lấy $\frac{1}{5}$ cuộn dây. Hỏi đoạn dây lấy ra dài bao nhiêu mét?
\shortans{9}
\loigiai{
\textbf{Đáp số:} $9$\\
$45 : 5 = 9\text{ m}$.
}
\end{ex}

\newpage

\TaoPhieuKiemTra{Chủ đề 2 -- Bảng nhân, bảng chia}{15}{Luyện tập chung}{https://loptoanthaylinh.vn/thi-online/toan-3/bai-15}
\textbf{Phần 1. Trắc nghiệm lựa chọn}
\setcounterpageref{ex}{0}

\begin{ex} %[3D2V15]
Tính giá trị của biểu thức: $M = 7 \times 9 - 48 : 8 + 14$.
\choice
{$68$}
{$73$}
{\True $71$}
{$69$}
\loigiai{
Thực hiện nhân chia trước: $M = 63 - 6 + 14 = 57 + 14 = 71$.
\textbf{Đáp án của câu hỏi là C}
}
\end{ex}

\begin{ex} %[3D2V15]
Tìm $x$ biết: $x \times 8 - 25 = 39$.
\choice
{\True $8$}
{$7$}
{$9$}
{$6$}
\loigiai{
Ta có: $x \times 8 = 39 + 25 = 64 \Rightarrow x = 64 : 8 = 8$.
\textbf{Đáp án của câu hỏi là A}
}
\end{ex}

\begin{ex} %[3D2V15]
Một cửa hàng có $72\text{ kg}$ gạo. Ngày thứ nhất bán được $\frac{1}{8}$ số gạo, ngày thứ hai bán được $\frac{1}{9}$ số gạo. Hỏi sau hai ngày cửa hàng còn lại bao nhiêu ki-lô-gam gạo?
\choice
{$50\text{ kg}$}
{$52\text{ kg}$}
{$54\text{ kg}$}
{\True $55\text{ kg}$}
\loigiai{
Ngày 1 bán: $72 : 8 = 9\text{ kg}$.\
Ngày 2 bán: $72 : 9 = 8\text{ kg}$.\
Còn lại: $72 - (9 + 8) = 72 - 17 = 55\text{ kg}$.
\textbf{Đáp án của câu hỏi là D}
}
\end{ex}

\begin{ex} %[3D2V15]
Trong một phép chia cho $7$, thương bằng $8$ và số dư là số dư lớn nhất có thể. Số bị chia là
\choice
{$56$}
{\True $62$}
{$61$}
{$63$}
\loigiai{
Số dư lớn nhất là 6.\
Số bị chia là: $7 \times 8 + 6 = 56 + 6 = 62$.
\textbf{Đáp án của câu hỏi là B}
}
\end{ex}

\begin{ex} %[3D2V15]
Tính nhanh giá trị của biểu thức: $8 \times 4 + 8 \times 5 + 8$.
\choice
{$72$}
{$88$}
{\True $80$}
{$96$}
\loigiai{
Ta viết $8 = 8 \times 1$. Đặt thừa số chung:\
$$8 \times (4 + 5 + 1) = 8 \times 10 = 80$$
\textbf{Đáp án của câu hỏi là C}
}
\end{ex}

\begin{ex} %[3D2V15]
Một bác nông dân nuôi đàn vịt và đàn gà có tất cả $63$ con. Biết số con gà bằng $\frac{1}{7}$ tổng số con. Hỏi bác nông dân nuôi số con vịt nhiều hơn số con gà bao nhiêu con?
\choice
{\True $45$ con}
{$36$ con}
{$48$ con}
{$42$ con}
\loigiai{
Số gà là: $63 : 7 = 9$ (con).\
Số vịt là: $63 - 9 = 54$ (con).\
Vịt hơn gà: $54 - 9 = 45$ (con).
\textbf{Đáp án của câu hỏi là A}
}
\end{ex}

\begin{ex} %[3D2V15]
Tìm $x$ biết: $x : 9 + 34 = 42$.
\choice
{$63$}
{$81$}
{$54$}
{\True $72$}
\loigiai{
Ta có: $x : 9 = 42 - 34 = 8 \Rightarrow x = 8 \times 9 = 72$.
\textbf{Đáp án của câu hỏi là D}
}
\end{ex}

\begin{ex} %[3D2V15]
Điền dấu thích hợp: $6 \times 9 - 18 \dots 7 \times 5 + 2$.
\choice
{$>$}
{\True $<$}
{$=$}
{Không so sánh được}
\loigiai{
Vế trái: $54 - 18 = 36$.\
Vế phải: $35 + 2 = 37$.\
Vì $36 < 37$ nên điền dấu $<$. 
\textbf{Đáp án của câu hỏi là B}
}
\end{ex}

\begin{ex} %[3D2V15]
Một người mua $6$ hộp bút, mỗi hộp có $8$ chiếc bút chì. Người đó đem tặng $\frac{1}{4}$ số bút chì đó cho các bạn học sinh nghèo. Hỏi người đó còn lại bao nhiêu chiếc bút chì?
\choice
{$32$ chiếc bút}
{$40$ chiếc bút}
{\True $36$ chiếc bút}
{$30$ chiếc bút}
\loigiai{
Tổng số bút chì là: $6 \times 8 = 48$ chiếc.\
Số bút đã tặng: $48 : 4 = 12$ chiếc.\
Số bút còn lại: $48 - 12 = 36$ chiếc.
\textbf{Đáp án của câu hỏi là C}
}
\end{ex}

\begin{ex} %[3D2V15]
Có $54$ bạn học sinh xếp đều thành 6 hàng. Nếu muốn xếp số học sinh đó thành các hàng, mỗi hàng có 9 bạn thì xếp được bao nhiêu hàng?
\choice
{\True $6$ hàng}
{$5$ hàng}
{$7$ hàng}
{$8$ hàng}
\loigiai{
Xếp được số hàng là: $54 : 9 = 6$ hàng.
\textbf{Đáp án của câu hỏi là A}
}
\end{ex}

\begin{ex} %[3D2V15]
Cho dãy số: $7; 14; 21; 28; \dots$ Tổng của số hạng thứ 3 và số hạng thứ 7 trong dãy là
\choice
{$63$}
{$77$}
{$68$}
{\True $70$}
\loigiai{
Số thứ 3 là: $7 \times 3 = 21$. Số thứ 7 là: $7 \times 7 = 49$.\
Tổng là: $21 + 49 = 70$.
\textbf{Đáp án của câu hỏi là D}
}
\end{ex}

\begin{ex} %[3D2V15]
Một số tự nhiên chia cho $8$ được thương bằng $6$ và dư $5$. Khi lấy số đó chia cho $9$ thì số dư là
\choice
{$7$}
{\True $8$}
{$6$}
{$5$}
\loigiai{
Số đó là: $8 \times 6 + 5 = 53$.\
Lấy $53 : 9 = 5$ dư 8.
\textbf{Đáp án của câu hỏi là B}
}
\end{ex}

\textbf{Phần 2. Trắc nghiệm đúng sai}
\setcounterpageref{ex}{0}

\begin{ex} %[3D2V15]
Xét tính đúng sai của các nhận định tổng hợp Chủ đề 2:
\choiceTF
{\True Một số chia hết cho cả 2 và 3 thì chia hết cho 6.}
{Một phần tám của 64 lớn hơn một phần bảy của 63.}
{\True Trong phép chia cho 9, số dư lớn nhất bằng 8.}
{\True Tích $7 \times 8$ bằng tích $8 \times 7$.}
\loigiai{
\begin{itemize}
    \item Ý a) Đúng: Dấu hiệu chia hết cơ bản.
    \item Ý b) Sai: $\frac{1}{8}$ của 64 là 8; $\frac{1}{7}$ của 63 là 9. Vì $8 < 9$ nên nhận định Sai.
    \item Ý c) Đúng: Số dư luôn bé hơn số chia.
    \item Ý d) Đúng: Tính chất giao hoán của phép nhân.
\end{itemize}
\begin{itemize}
    \item a) \textbf{Đúng}.
    \item b) \textbf{Sai}.
    \item c) \textbf{Đúng}.
    \item d) \textbf{Đúng}.
\end{itemize}
}
\end{ex}

\begin{ex} %[3D2V15]
Một người thợ làm việc trong 6 ngày, mỗi ngày làm được 8 sản phẩm. Người đó đem bán đi một nửa số sản phẩm:
\choiceTF
{\True Người đó làm được tất cả 48 sản phẩm.}
{\True Người đó đã bán đi 24 sản phẩm.}
{Số sản phẩm còn lại đủ để chia đều vào 5 hộp, mỗi hộp 5 sản phẩm.}
{\True Số sản phẩm còn lại chia vừa đủ vào 4 hộp loại 6 sản phẩm.}
\loigiai{
\begin{itemize}
    \item Ý a) Đúng: $6 \times 8 = 48$ sản phẩm.
    \item Ý b) Đúng: $48 : 2 = 24$ sản phẩm.
    \item Ý c) Sai: Cần $5 \times 5 = 25$ sản phẩm, mà chỉ còn 24 sản phẩm.
    \item Ý d) Đúng: $24 : 4 = 6$ sản phẩm.
\end{itemize}
\begin{itemize}
    \item a) \textbf{Đúng}.
    \item b) \textbf{Đúng}.
    \item c) \textbf{Sai}.
    \item d) \textbf{Đúng}.
\end{itemize}
}
\end{ex}

\textbf{Phần 3. Trả lời ngắn}
\setcounterpageref{ex}{0}

\begin{ex} %[3D2V15]
Tính giá trị của biểu thức: $8 \times 7 - 45 : 9$.
\shortans{51}
\loigiai{
\textbf{Đáp số:} $51$\\
$8 \times 7 - 45 : 9 = 56 - 5 = 51$.
}
\end{ex}

\begin{ex} %[3D2V15]
Tìm một số biết rằng $\frac{1}{7}$ của số đó bằng 8.
\shortans{56}
\loigiai{
\textbf{Đáp số:} $56$\\
$8 \times 7 = 56$.
}
\end{ex}

\begin{ex} %[3D2V15]
Tìm $x$ biết: $x : 6 = 14 - 7$.
\shortans{42}
\loigiai{
\textbf{Đáp số:} $42$\\
$x : 6 = 7 \Rightarrow x = 7 \times 6 = 42$.
}
\end{ex}

\begin{ex} %[3D2V15]
Một lớp có 32 học sinh chia đều thành 4 tổ. Hỏi 3 tổ có bao nhiêu học sinh?
\shortans{24}
\loigiai{
\textbf{Đáp số:} $24$\\
$32 : 4 = 8$; $8 \times 3 = 24$ bạn.
}
\end{ex}

\begin{ex} %[3D2V15]
Tính nhanh: $9 \times 2 + 9 \times 3 + 9 \times 5$.
\shortans{90}
\loigiai{
\textbf{Đáp số:} $90$\\
$9 \times (2 + 3 + 5) = 9 \times 10 = 90$.
}
\end{ex}

\begin{ex} %[3D2V15]
Một hình tam giác có độ dài ba cạnh đều bằng 8 cm. Chu vi hình tam giác đó là bao nhiêu xăng-ti-mét?
\shortans{24}
\loigiai{
\textbf{Đáp số:} $24$\\
Chu vi là: $8 \times 3 = 24\text{ cm}$.
}
\end{ex}

\newpage

% =========================================================================
% CHỦ ĐỀ 3 - LÀM QUEN VỚI HÌNH PHẲNG, HÌNH KHỐI
% =========================================================================
\section{Làm quen với hình phẳng, hình khối}

\TaoPhieuKiemTra{Chủ đề 3 -- Làm quen với hình phẳng, hình khối}{16}{Điểm ở giữa, trung điểm của đoạn thẳng}{https://loptoanthaylinh.vn/thi-online/toan-3/bai-16}
\textbf{Phần 1. Trắc nghiệm lựa chọn}
\setcounterpageref{ex}{0}

\begin{ex} %[3H1V16]
Cho đoạn thẳng $AB$ dài $48\text{ cm}$. Điểm $M$ là trung điểm của $AB$, điểm $N$ là trung điểm của đoạn thẳng $AM$. Độ dài đoạn thẳng $NB$ là
\choice
{$24\text{ cm}$}
{$30\text{ cm}$}
{$40\text{ cm}$}
{\True $36\text{ cm}$}
\loigiai{
Độ dài đoạn thẳng $AM$ là: $48 : 2 = 24\text{ (cm)}$.\
Độ dài đoạn thẳng $AN$ là: $24 : 2 = 12\text{ (cm)}$.\
Độ dài đoạn thẳng $NB$ là: $48 - 12 = 36\text{ (cm)}$.
\textbf{Đáp án của câu hỏi là D}
}
\end{ex}

\begin{ex} %[3H1V16]
Trên một đường thẳng, lấy ba điểm $A, B, C$ sao cho $B$ nằm giữa $A$ và $C$. Biết $AB = 16\text{ cm}$, $BC = 16\text{ cm}$. Khẳng định nào sau đây là đúng nhất?
\choice
{$A$ là trung điểm của đoạn thẳng $BC$}
{\True $B$ là trung điểm của đoạn thẳng $AC$}
{$C$ là trung điểm của đoạn thẳng $AB$}
{$B$ chỉ là điểm nằm giữa chứ không phải trung điểm}
\loigiai{
Vì $B$ nằm giữa $A$ và $C$ đồng thời $AB = BC = 16\text{ cm}$ nên $B$ chính là trung điểm của đoạn thẳng $AC$.
\textbf{Đáp án của câu hỏi là B}
}
\end{ex}

\begin{ex} %[3H1V16]
Một sợi dây đồng dài $64\text{ cm}$ được gập đôi lại hai lần liên tiếp. Khi đó độ dài mỗi đoạn dây nhận được là
\choice
{$32\text{ cm}$}
{$8\text{ cm}$}
{\True $16\text{ cm}$}
{$12\text{ cm}$}
\loigiai{
Gập đôi lần thứ nhất, chiều dài mỗi phần là: $64 : 2 = 32\text{ (cm)}$.\
Gập đôi lần thứ hai, chiều dài mỗi phần là: $32 : 2 = 16\text{ (cm)}$.
\textbf{Đáp án của câu hỏi là C}
}
\end{ex}

\begin{ex} %[3H1V16]
Một con đường thẳng dài $100\text{ m}$. Người ta trồng cây ở cả hai đầu đường và tại chính giữa con đường (trung điểm của con đường). Sau đó tiếp tục trồng thêm 2 cây nữa ở chính giữa mỗi nửa đoạn đường đó. Hỏi khoảng cách giữa hai cây liên tiếp gần nhau nhất là bao nhiêu mét?
\choice
{\True $25\text{ m}$}
{$20\text{ m}$}
{$50\text{ m}$}
{$10\text{ m}$}
\loigiai{
Con đường được chia đều làm 4 đoạn bằng nhau. Khoảng cách giữa hai cây liên tiếp là: $100 : 4 = 25\text{ (m)}$.
\textbf{Đáp án của câu hỏi là A}
}
\end{ex}

\begin{ex} %[3H1V16]
Cho đoạn thẳng $MN = 72\text{ cm}$. Gọi $I$ là trung điểm của đoạn thẳng $MN$, $K$ là trung điểm của đoạn thẳng $IN$. Độ dài đoạn thẳng $MK$ bằng
\choice
{$36\text{ cm}$}
{$48\text{ cm}$}
{$60\text{ cm}$}
{\True $54\text{ cm}$}
\loigiai{
Ta có: $MI = IN = 72 : 2 = 36\text{ (cm)}$.\
$K$ là trung điểm của $IN$ nên $IK = 36 : 2 = 18\text{ (cm)}$.\
Độ dài $MK = MI + IK = 36 + 18 = 54\text{ (cm)}$.
\textbf{Đáp án của câu hỏi là D}
}
\end{ex}

\begin{ex} %[3H1V16]
Cho ba điểm $M, N, P$ thẳng hàng theo thứ tự đó với $MN = 20\text{ cm}, NP = 30\text{ cm}$. Gọi $A$ là trung điểm của $MN$ và $B$ là trung điểm của $NP$. Độ dài đoạn thẳng $AB$ là
\choice
{$30\text{ cm}$}
{\True $25\text{ cm}$}
{$20\text{ cm}$}
{$35\text{ cm}$}
\loigiai{
Độ dài $AN = MN : 2 = 20 : 2 = 10\text{ (cm)}$.\
Độ dài $NB = NP : 2 = 30 : 2 = 15\text{ (cm)}$.\
Độ dài đoạn thẳng $AB = AN + NB = 10 + 15 = 25\text{ (cm)}$.
\textbf{Đáp án của câu hỏi là B}
}
\end{ex}

\begin{ex} %[3H1V16]
Một thanh gỗ dài $90\text{ cm}$. Bác thợ mộc cưa thanh gỗ đó tại trung điểm của nó, sau đó lấy một nửa cưa tiếp tại trung điểm để được hai mẩu gỗ nhỏ hơn. Độ dài của một mẩu gỗ nhỏ nhất là
\choice
{$45\text{ cm}$}
{$30\text{ cm}$}
{\True $22\text{ cm và } 5\text{ mm}$}
{$20\text{ cm}$}
\loigiai{
Đổi $90\text{ cm} = 900\text{ mm}$.\
Cưa lần 1 được hai thanh dài: $900 : 2 = 450\text{ mm}$.\
Cưa lần 2 được hai thanh nhỏ dài: $450 : 2 = 225\text{ mm} = 22\text{ cm } 5\text{ mm}$.
\textbf{Đáp án của câu hỏi là C}
}
\end{ex}

\begin{ex} %[3H1V16]
Cho đoạn thẳng $AB$. Điểm $C$ nằm giữa $A$ và $B$ sao cho $AC$ ngắn hơn $CB$ là $8\text{ cm}$. Biết $AB = 40\text{ cm}$. Khoảng cách từ điểm $C$ đến trung điểm $M$ của đoạn thẳng $AB$ là
\choice
{\True $4\text{ cm}$}
{$8\text{ cm}$}
{$6\text{ cm}$}
{$5\text{ cm}$}
\loigiai{
Trung điểm $M$ cách $A$ một đoạn: $AM = 40 : 2 = 20\text{ (cm)}$.\
Độ dài đoạn $AC$ là: $(40 - 8) : 2 = 16\text{ (cm)}$.\
Khoảng cách từ $C$ đến $M$ là: $CM = AM - AC = 20 - 16 = 4\text{ (cm)}$.
\textbf{Đáp án của câu hỏi là A}
}
\end{ex}

\begin{ex} %[3H1V16]
Có 4 điểm $A, B, C, D$ cùng nằm trên một đường thẳng theo thứ tự đó. Biết $AB = BC = CD = 12\text{ cm}$. Khẳng định nào sau đây là SAI?
\choice
{$B$ là trung điểm của đoạn thẳng $AC$}
{$C$ là trung điểm của đoạn thẳng $BD$}
{Khoảng cách từ $A$ đến $D$ bằng $36\text{ cm}$}
{\True $C$ là trung điểm của đoạn thẳng $AD$}
\loigiai{
Đoạn $AC = 24\text{ cm}$, đoạn $CD = 12\text{ cm}$ nên $C$ không phải là trung điểm của $AD$ (trung điểm của $AD$ cách $A$ là $18\text{ cm}$).
\textbf{Đáp án của câu hỏi là D}
}
\end{ex}

\begin{ex} %[3H1V16]
Một sợi ruy-băng dài $84\text{ cm}$. Lan muốn cắt lấy một đoạn dài $21\text{ cm}$ nhưng trong tay không có thước đo độ dài. Lan cần gấp đôi sợi dây liên tiếp bao nhiêu lần rồi cắt theo nếp gấp gần nhất?
\choice
{$1$ lần}
{\True $2$ lần}
{$3$ lần}
{$4$ lần}
\loigiai{
Sau 1 lần gập đôi, đoạn dây có độ dài: $84 : 2 = 42\text{ cm}$.\
Sau lần gập đôi thứ hai, đoạn dây có độ dài: $42 : 2 = 21\text{ cm}$.\
Vậy Lan chỉ cần gập đôi sợi dây liên tiếp 2 lần.
\textbf{Đáp án của câu hỏi là B}
}
\end{ex}

\begin{ex} %[3H1V16]
Trên tia số, điểm $A$ ứng với số $200$, điểm $B$ ứng với số $600$. Điểm $M$ là trung điểm của đoạn thẳng $AB$ sẽ ứng với số nào trên tia số?
\choice
{$300$}
{$500$}
{\True $400$}
{$350$}
\loigiai{
Khoảng cách giữa $A$ và $B$ là: $600 - 200 = 400$.\\nTrung điểm $M$ cách điểm $A$ một khoảng: $400 : 2 = 200$.\\nVậy điểm $M$ ứng với số: $200 + 200 = 400$.
\textbf{Đáp án của câu hỏi là C}
}
\end{ex}

\begin{ex} %[3H1V16]
Cho đoạn thẳng $PQ = 56\text{ cm}$. Gọi $M$ là trung điểm của $PQ$, $N$ là trung điểm của $MQ$, $K$ là trung điểm của $NQ$. Độ dài đoạn thẳng $PK$ là
\choice
{\True $49\text{ cm}$}
{$42\text{ cm}$}
{$52\text{ cm}$}
{$45\text{ cm}$}
\loigiai{
$MQ = 56 : 2 = 28\text{ cm}$.\
$NQ = 28 : 2 = 14\text{ cm}$.\
$KQ = 14 : 2 = 7\text{ cm}$.\
Độ dài đoạn thẳng $PK = PQ - KQ = 56 - 7 = 49\text{ (cm)}$.
\textbf{Đáp án của câu hỏi là A}
}
\end{ex}

\textbf{Phần 2. Trắc nghiệm đúng sai}
\setcounterpageref{ex}{0}

\begin{ex} %[3H1V16]
Xét tính đúng sai của các phát biểu về điểm ở giữa và trung điểm:
\choiceTF
{\True Nếu điểm $M$ là trung điểm của đoạn thẳng $AB$ thì $M$ chắc chắn nằm giữa hai điểm $A$ và $B$.}
{Nếu $MA = MB$ thì điểm $M$ luôn là trung điểm của đoạn thẳng $AB$.}
{\True Mỗi đoạn thẳng chỉ có duy nhất một trung điểm.}
{Một đoạn thẳng có vô số điểm nằm ở giữa hai đầu mút của nó là sai.}
\loigiai{
\begin{itemize}
    \item Ý a) Đúng: Định nghĩa trung điểm bao hàm điểm đó phải nằm giữa.
    \item Ý b) Sai: Sai vì ba điểm $A, M, B$ có thể không thẳng hàng tạo thành tam giác cân.
    \item Ý c) Đúng: Đoạn thẳng chỉ có một vị trí chính giữa cách đều hai đầu mút.
    \item Ý d) Sai: Sai vì giữa hai đầu mút của một đoạn thẳng có vô số điểm thẳng hàng nằm giữa.
\end{itemize}
\begin{itemize}
    \item a) \textbf{Đúng}.
    \item b) \textbf{Sai}.
    \item c) \textbf{Đúng}.
    \item d) \textbf{Sai}.
\end{itemize}
}
\end{ex}

\begin{ex} %[3H1V16]
Cho đoạn thẳng $AB$ dài $80\text{ cm}$. Điểm $M$ là trung điểm của $AB$, điểm $P$ là trung điểm của $AM$:
\choiceTF
{\True Độ dài đoạn thẳng $AM$ bằng $40\text{ cm}$.}
{\True Độ dài đoạn thẳng $AP$ bằng $20\text{ cm}$.}
{Điểm $M$ là trung điểm của đoạn thẳng $PB$.}
{\True Độ dài đoạn thẳng $PB$ bằng $60\text{ cm}$.}
\loigiai{
\begin{itemize}
    \item Ý a) Đúng: $80 : 2 = 40\text{ cm}$.
    \item Ý b) Đúng: $40 : 2 = 20\text{ cm}$.
    \item Ý c) Sai: $PM = 20\text{ cm}$, nhưng $MB = 40\text{ cm}$ nên $M$ không là trung điểm $PB$.
    \item Ý d) Đúng: $PB = PM + MB = 20 + 40 = 60\text{ cm}$.
\end{itemize}
\begin{itemize}
    \item a) \textbf{Đúng}.
    \item b) \textbf{Đúng}.
    \item c) \textbf{Sai}.
    \item d) \textbf{Đúng}.
\end{itemize}
}
\end{ex}

\textbf{Phần 3. Trả lời ngắn}
\setcounterpageref{ex}{0}

\begin{ex} %[3H1V16]
Cho đoạn thẳng $AB = 60\text{ cm}$. Điểm $M$ là trung điểm của $AB$, điểm $N$ là trung điểm của $MB$. Tính độ dài đoạn thẳng $AN$ theo xăng-ti-mét.
\shortans{45}
\loigiai{
\textbf{Đáp số:} $45$\\
Độ dài $AM = MB = 30\text{ cm}$. $MN = 30 : 2 = 15\text{ cm}$. $AN = AM + MN = 30 + 15 = 45\text{ (cm)}$.
}
\end{ex}

\begin{ex} %[3H1V16]
Trên đoạn thẳng $CD$ dài $72\text{ cm}$, người ta lấy điểm $E$ sao cho $CE = 18\text{ cm}$. Hỏi điểm $E$ cách trung điểm $M$ của đoạn thẳng $CD$ bao nhiêu xăng-ti-mét?
\shortans{18}
\loigiai{
\textbf{Đáp số:} $18$\\
Trung điểm $M$ cách $C$ là: $72 : 2 = 36\text{ cm}$. Khoảng cách $EM = 36 - 18 = 18\text{ cm}$.
}
\end{ex}

\begin{ex} %[3H1V16]
Một đường ống dẫn nước thẳng dài $96\text{ m}$. Người ta lắp một van khóa nước tại trung điểm của đường ống. Sau đó lắp thêm một đồng hồ đo áp lực tại trung điểm của nửa đường ống thứ nhất. Hỏi đồng hồ đo áp lực cách đầu ống nước ban đầu bao nhiêu mét?
\shortans{24}
\loigiai{
\textbf{Đáp số:} $24$\\
Nửa ống nước dài $96 : 2 = 48\text{ m}$. Đồng hồ cách đầu ống: $48 : 2 = 24\text{ m}$.
}
\end{ex}

\begin{ex} %[3H1V16]
Cho ba điểm thẳng hàng $A, B, C$ với $B$ nằm giữa $A$ và $C$. Biết $AB = 24\text{ cm}$ và $BC$ gấp đôi $AB$. Gọi $M$ là trung điểm của $BC$. Tính độ dài đoạn thẳng $AM$ theo xăng-ti-mét.
\shortans{48}
\loigiai{
\textbf{Đáp số:} $48$\\
$BC = 24 \times 2 = 48\text{ cm}$. $BM = 48 : 2 = 24\text{ cm}$. $AM = AB + BM = 24 + 24 = 48\text{ cm}$.
}
\end{ex}

\begin{ex} %[3H1V16]
Người ta chôn 9 cây cọc thẳng hàng cách đều nhau, khoảng cách giữa cọc thứ nhất và cọc thứ chín là $64\text{ m}$. Cọc thứ năm là cọc ở chính giữa. Hỏi khoảng cách từ cọc thứ nhất đến cọc thứ năm là bao nhiêu mét?
\shortans{32}
\loigiai{
\textbf{Đáp số:} $32$\\
Cọc thứ 5 là trung điểm của đoạn thẳng nối cọc 1 và cọc 9. Khoảng cách đó là: $64 : 2 = 32\text{ m}$.
}
\end{ex}

\begin{ex} %[3H1V16]
Một sợi dây thép uốn thành hình vuông có chu vi $80\text{ cm}$. Người ta duỗi thẳng sợi dây rồi cắt tại trung điểm của nó. Tính độ dài mỗi đoạn dây nhận được theo xăng-ti-mét.
\shortans{40}
\loigiai{
\textbf{Đáp số:} $40$\\
Độ dài sợi dây ban đầu là $80\text{ cm}$. Cắt tại trung điểm thì mỗi đoạn dài: $80 : 2 = 40\text{ cm}$.
}
\end{ex}

\newpage

\TaoPhieuKiemTra{Chủ đề 3 -- Làm quen với hình phẳng, hình khối}{17}{Hình tròn, tâm, bán kính, đường kính}{https://loptoanthaylinh.vn/thi-online/toan-3/bai-17}
\textbf{Phần 1. Trắc nghiệm lựa chọn}
\setcounterpageref{ex}{0}

\begin{ex} %[3H1V17]
Một hình tròn có bán kính bằng $14\text{ cm}$. Đường kính của hình tròn đó là
\choice
{$7\text{ cm}$}
{$21\text{ cm}$}
{\True $28\text{ cm}$}
{$35\text{ cm}$}
\loigiai{
Đường kính gấp 2 lần bán kính: $14 \times 2 = 28\text{ (cm)}$.
\textbf{Đáp án của câu hỏi là C}
}
\end{ex}

\begin{ex} %[3H1V17]
Một hình tròn có đường kính dài bằng chu vi của hình vuông cạnh $12\text{ cm}$. Bán kính của hình tròn đó bằng
\choice
{\True $24\text{ cm}$}
{$48\text{ cm}$}
{$12\text{ cm}$}
{$36\text{ cm}$}
\loigiai{
Chu vi hình vuông là: $12 \times 4 = 48\text{ (cm)}$.\
Đường kính hình tròn là $48\text{ cm}$ nên bán kính là: $48 : 2 = 24\text{ (cm)}$.
\textbf{Đáp án của câu hỏi là A}
}
\end{ex}

\begin{ex} %[3H1V17]
Cho hình tròn tâm $O$ có bán kính $R = 15\text{ cm}$. Điểm $A$ nằm trong hình tròn, điểm $B$ nằm trên đường tròn và điểm $C$ nằm ngoài hình tròn. Khẳng định nào sau đây là đúng?
\choice
{$OA = 15\text{ cm}$ và $OB > 15\text{ cm}$}
{$OA > 15\text{ cm}$ và $OC < 15\text{ cm}$}
{$OB < 15\text{ cm}$ và $OC = 15\text{ cm}$}
{\True $OA < 15\text{ cm}$ và $OC > 15\text{ cm}$}
\loigiai{
Điểm nằm trong hình tròn có khoảng cách tới tâm nhỏ hơn bán kính: $OA < 15\text{ cm}$. Điểm nằm ngoài hình tròn có khoảng cách tới tâm lớn hơn bán kính: $OC > 15\text{ cm}$.
\textbf{Đáp án của câu hỏi là D}
}
\end{ex}

\begin{ex} %[3H1V17]
Hai hình tròn có chung tâm $O$. Hình tròn thứ nhất có bán kính $18\text{ cm}$, hình tròn thứ hai có bán kính gấp đôi bán kính hình tròn thứ nhất. Đường kính của hình tròn thứ hai là
\choice
{$36\text{ cm}$}
{\True $72\text{ cm}$}
{$54\text{ cm}$}
{$90\text{ cm}$}
\loigiai{
Bán kính hình tròn thứ hai là: $18 \times 2 = 36\text{ (cm)}$.\
Đường kính hình tròn thứ hai là: $36 \times 2 = 72\text{ (cm)}$.
\textbf{Đáp án của câu hỏi là B}
}
\end{ex}

\begin{ex} %[3H1V17]
Có ba hình tròn bằng nhau đặt thẳng hàng tiếp xúc nhau sao cho tâm của chúng cùng nằm trên một đoạn thẳng. Biết khoảng cách giữa tâm hình tròn thứ nhất và tâm hình tròn thứ ba là $28\text{ cm}$. Bán kính của mỗi hình tròn là
\choice
{$14\text{ cm}$}
{$4\text{ cm}$}
{\True $7\text{ cm}$}
{$8\text{ cm}$}
\loigiai{
Khoảng cách giữa tâm thứ nhất và tâm thứ ba gồm 2 đoạn nối tâm liền kề, mỗi đoạn bằng 2 lần bán kính: khoảng cách bằng $4$ lần bán kính.\
Bán kính mỗi hình tròn là: $28 : 4 = 7\text{ (cm)}$.
\textbf{Đáp án của câu hỏi là C}
}
\end{ex}

\begin{ex} %[3H1V17]
Một compa mở hai chân tạo một khoảng cách bằng $9\text{ cm}$. Khi quay một vòng để vẽ đường tròn thì đường tròn nhận được có đường kính là
\choice
{\True $18\text{ cm}$}
{$9\text{ cm}$}
{$27\text{ cm}$}
{$36\text{ cm}$}
\loigiai{
Khoảng cách giữa hai chân compa chính là bán kính hình tròn: $R = 9\text{ cm}$.\
Đường kính của hình tròn là: $9 \times 2 = 18\text{ (cm)}$.
\textbf{Đáp án của câu hỏi là A}
}
\end{ex}

\begin{ex} %[3H1V17]
Trong một hình chữ nhật có chiều dài $20\text{ cm}$ và chiều rộng $12\text{ cm}$, người ta vẽ một hình tròn có bán kính lớn nhất có thể nằm hoàn toàn bên trong hình chữ nhật đó. Đường kính của hình tròn đó bằng
\choice
{$20\text{ cm}$}
{$6\text{ cm}$}
{$10\text{ cm}$}
{\True $12\text{ cm}$}
\loigiai{
Hình tròn lớn nhất nằm trọn trong hình chữ nhật có đường kính bằng đúng chiều rộng hình chữ nhật, tức là $12\text{ cm}$.
\textbf{Đáp án của câu hỏi là D}
}
\end{ex}

\begin{ex} %[3H1V17]
Tổng độ dài của đường kính và bán kính của một hình tròn là $45\text{ cm}$. Bán kính của hình tròn đó là
\choice
{$30\text{ cm}$}
{\True $15\text{ cm}$}
{$10\text{ cm}$}
{$20\text{ cm}$}
\loigiai{
Vì đường kính gấp đôi bán kính nên tổng độ dài đường kính và bán kính bằng 3 lần bán kính.\
Bán kính hình tròn là: $45 : 3 = 15\text{ (cm)}$.
\textbf{Đáp án của câu hỏi là B}
}
\end{ex}

\begin{ex} %[3H1V17]
Đường kính của hình tròn thứ nhất bằng bán kính của hình tròn thứ hai. Biết bán kính hình tròn thứ nhất là $8\text{ cm}$. Đường kính của hình tròn thứ hai là
\choice
{$16\text{ cm}$}
{$24\text{ cm}$}
{\True $32\text{ cm}$}
{$64\text{ cm}$}
\loigiai{
Đường kính hình tròn thứ nhất là: $8 \times 2 = 16\text{ (cm)}$.\
Bán kính hình tròn thứ hai bằng $16\text{ cm}$, suy ra đường kính của nó là: $16 \times 2 = 32\text{ (cm)}$.
\textbf{Đáp án của câu hỏi là C}
}
\end{ex}

\begin{ex} %[3H1V17]
Một hình tròn tâm $O$ có bán kính $R = 25\text{ cm}$. Đoạn thẳng $MN$ đi qua tâm $O$ và có hai đầu mút nằm trên đường tròn. Đoạn thẳng $MN$ có độ dài là
\choice
{\True $50\text{ cm}$}
{$25\text{ cm}$}
{$75\text{ cm}$}
{$100\text{ cm}$}
\loigiai{
Đoạn thẳng đi qua tâm và có hai mút nằm trên đường tròn chính là đường kính của hình tròn: $MN = 25 \times 2 = 50\text{ (cm)}$.
\textbf{Đáp án của câu hỏi là A}
}
\end{ex}

\begin{ex} %[3H1V17]
Một bánh xe đạp có bán kính $30\text{ cm}$. Khoảng cách từ trục bánh xe (tâm) đến một điểm bất kì trên vành bánh xe là
\choice
{$60\text{ cm}$}
{$15\text{ cm}$}
{$45\text{ cm}$}
{\True $30\text{ cm}$}
\loigiai{
Khoảng cách từ tâm đến một điểm bất kì trên đường tròn chính là bán kính của đường tròn, bằng $30\text{ cm}$.
\textbf{Đáp án của câu hỏi là D}
}
\end{ex}

\begin{ex} %[3H1V17]
Cho đoạn thẳng $AB = 10\text{ cm}$. Vẽ đường tròn tâm $A$ bán kính $6\text{ cm}$ và đường tròn tâm $B$ bán kính $4\text{ cm}$. Khẳng định nào sau đây là đúng?
\choice
{Hai đường tròn cắt nhau tại 2 điểm}
{\True Hai đường tròn tiếp xúc nhau tại đúng 1 điểm trên đoạn thẳng $AB$}
{Hai đường tròn không có điểm chung nào}
{Đường tròn tâm $B$ nằm hoàn toàn bên trong đường tròn tâm $A$}
\loigiai{
Tổng hai bán kính là $6 + 4 = 10\text{ cm}$, đúng bằng khoảng cách giữa hai tâm $AB = 10\text{ cm}$. Do đó hai đường tròn tiếp xúc ngoài tại một điểm trên đoạn $AB$.
\textbf{Đáp án của câu hỏi là B}
}
\end{ex}

\textbf{Phần 2. Trắc nghiệm đúng sai}
\setcounterpageref{ex}{0}

\begin{ex} %[3H1V17]
Xét tính đúng sai của các khẳng định về hình tròn:
\choiceTF
{\True Đường kính của hình tròn luôn có độ dài gấp 2 lần bán kính của hình tròn đó.}
{Tất cả các đoạn thẳng nối hai điểm trên đường tròn đều là đường kính.}
{\True Tâm của hình tròn chính là trung điểm của mọi đường kính của hình tròn đó.}
{Hai hình tròn có bán kính khác nhau có thể có cùng đường kính.}
\loigiai{
\begin{itemize}
    \item Ý a) Đúng: Đường kính $d = 2 \times R$.
    \item Ý b) Sai: Sai vì chỉ đoạn thẳng đi qua tâm mới là đường kính.
    \item Ý c) Đúng: Tâm nằm giữa và chia đôi đường kính.
    \item Ý d) Sai: Sai vì bán kính khác nhau thì đường kính cũng phải khác nhau.
\end{itemize}
\begin{itemize}
    \item a) \textbf{Đúng}.
    \item b) \textbf{Sai}.
    \item c) \textbf{Đúng}.
    \item d) \textbf{Sai}.
\end{itemize}
}
\end{ex}

\begin{ex} %[3H1V17]
Cho hình tròn tâm $O$ có đường kính $AB = 24\text{ cm}$:
\choiceTF
{\True Bán kính của hình tròn bằng $12\text{ cm}$.}
{Nếu điểm $C$ nằm ngoài hình tròn thì đoạn thẳng $OC$ có thể ngắn hơn $12\text{ cm}$.}
{\True Điểm $O$ là trung điểm của đoạn thẳng $AB$.}
{\True Mọi bán kính xuất phát từ tâm $O$ đều có độ dài bằng nhau và bằng $12\text{ cm}$.}
\loigiai{
\begin{itemize}
    \item Ý a) Đúng: $24 : 2 = 12\text{ cm}$.
    \item Ý b) Sai: Sai vì điểm ngoài hình tròn luôn có $OC > 12\text{ cm}$.
    \item Ý c) Đúng: Đoạn $AB$ là đường kính nên tâm $O$ là trung điểm.
    \item Ý d) Đúng: Tính chất của hình tròn.
\end{itemize}
\begin{itemize}
    \item a) \textbf{Đúng}.
    \item b) \textbf{Sai}.
    \item c) \textbf{Đúng}.
    \item d) \textbf{Đúng}.
\end{itemize}
}
\end{ex}

\textbf{Phần 3. Trả lời ngắn}
\setcounterpageref{ex}{0}

\begin{ex} %[3H1V17]
Một hình tròn có đường kính bằng $36\text{ cm}$. Tính bán kính của hình tròn đó theo xăng-ti-mét.
\shortans{18}
\loigiai{
\textbf{Đáp số:} $18$\\
Bán kính là: $36 : 2 = 18\text{ (cm)}$.
}
\end{ex}

\begin{ex} %[3H1V17]
Một hình tròn có bán kính $8\text{ cm}$. Hỏi đường kính của hình tròn đó dài bao nhiêu xăng-ti-mét?
\shortans{16}
\loigiai{
\textbf{Đáp số:} $16$\\
Đường kính là: $8 \times 2 = 16\text{ (cm)}$.
}
\end{ex}

\begin{ex} %[3H1V17]
Cho đoạn thẳng $MN = 40\text{ cm}$. Vẽ hình tròn nhận $MN$ làm đường kính. Tính khoảng cách từ tâm của hình tròn đó đến điểm $M$ theo xăng-ti-mét.
\shortans{20}
\loigiai{
\textbf{Đáp số:} $20$\\
Khoảng cách từ tâm đến điểm mút $M$ chính là bán kính: $40 : 2 = 20\text{ (cm)}$.
}
\end{ex}

\begin{ex} %[3H1V17]
Hai hình tròn tâm $A$ và tâm $B$ tiếp xúc ngoài với nhau. Biết bán kính hình tròn tâm $A$ là $15\text{ cm}$, bán kính hình tròn tâm $B$ là $9\text{ cm}$. Tính độ dài đoạn thẳng nối hai tâm $AB$ theo xăng-ti-mét.
\shortans{24}
\loigiai{
\textbf{Đáp số:} $24$\\
Độ dài đoạn nối hai tâm bằng tổng hai bán kính: $15 + 9 = 24\text{ (cm)}$.
}
\end{ex}

\begin{ex} %[3H1V17]
Một hình tròn có đường kính gấp $4$ lần bán kính của một hình tròn nhỏ có bán kính $5\text{ cm}$. Tính bán kính của hình tròn lớn theo xăng-ti-mét.
\shortans{10}
\loigiai{
\textbf{Đáp số:} $10$\\
Đường kính hình tròn lớn là: $5 \times 4 = 20\text{ cm}$. Bán kính hình tròn lớn là: $20 : 2 = 10\text{ cm}$.
}
\end{ex}

\begin{ex} %[3H1V17]
Người ta xếp 4 đồng xu hình tròn giống hệt nhau vừa khít vào một hàng thẳng có tổng chiều dài $24\text{ cm}$. Tính đường kính của mỗi đồng xu theo xăng-ti-mét.
\shortans{6}
\loigiai{
\textbf{Đáp số:} $6$\\
Đường kính của mỗi đồng xu là: $24 : 4 = 6\text{ (cm)}$.
}
\end{ex}

\newpage

\TaoPhieuKiemTra{Chủ đề 3 -- Làm quen với hình phẳng, hình khối}{18}{Góc, góc vuông, góc không vuông}{https://loptoanthaylinh.vn/thi-online/toan-3/bai-18}
\textbf{Phần 1. Trắc nghiệm lựa chọn}
\setcounterpageref{ex}{0}

\begin{ex} %[3H1V18]
Một hình chữ nhật có 4 góc. Hỏi một hình chữ nhật có bao nhiêu góc vuông?
\choice
{$2$ góc vuông}
{\True $4$ góc vuông}
{$3$ góc vuông}
{$1$ góc vuông}
\loigiai{
Hình chữ nhật có 4 góc đỉnh đều là góc vuông.
\textbf{Đáp án của câu hỏi là B}
}
\end{ex}

\begin{ex} %[3H1V18]
Lúc đúng $3$ giờ, kim giờ và kim phút của đồng hồ tạo thành một góc gì?
\choice
{Góc nhọn}
{Góc tù}
{Góc bẹt}
{\True Góc vuông}
\loigiai{
Lúc 3 giờ đúng, kim phút chỉ số 12 và kim giờ chỉ số 3, hai kim tạo thành một góc vuông ($90^\circ$).
\textbf{Đáp án của câu hỏi là D}
}
\end{ex}

\begin{ex} %[3H1V18]
Trong các thời điểm sau, thời điểm nào mà hai kim đồng hồ (kim giờ và kim phút) tạo thành một góc vuông?
\choice
{\True $9$ giờ đúng}
{$6$ giờ đúng}
{$12$ giờ đúng}
{$4$ giờ đúng}
\loigiai{
Lúc 9 giờ đúng, kim phút chỉ số 12 và kim giờ chỉ số 9, hai kim vuông góc với nhau.
\textbf{Đáp án của câu hỏi là A}
}
\end{ex}

\begin{ex} %[3H1V18]
Một hình tứ giác có thể có nhiều nhất bao nhiêu góc vuông?
\choice
{$2$ góc vuông}
{$3$ góc vuông}
{\True $4$ góc vuông}
{$1$ góc vuông}
\loigiai{
Hình vuông và hình chữ nhật là các tứ giác có 4 góc vuông (nhiều nhất có thể).
\textbf{Đáp án của câu hỏi là C}
}
\end{ex}

\begin{ex} %[3H1V18]
Một hình tam giác có thể có nhiều nhất bao nhiêu góc vuông?
\choice
{$2$ góc vuông}
{\True $1$ góc vuông}
{$3$ góc vuông}
{Không có góc vuông nào}
\loigiai{
Một hình tam giác chỉ có thể có nhiều nhất 1 góc vuông (tam giác vuông).
\textbf{Đáp án của câu hỏi là B}
}
\end{ex}

\begin{ex} %[3H1V18]
Dụng cụ chuẩn thường dùng để kiểm tra một góc có phải là góc vuông hay không là
\choice
{Thước dây}
{Compa}
{Thước kẻ thẳng thông thường}
{\True Thước ê-ke}
\loigiai{
Thước ê-ke có góc vuông chuẩn dùng để kiểm tra góc vuông.
\textbf{Đáp án của câu hỏi là D}
}
\end{ex}

\begin{ex} %[3H1V18]
Lúc $2$ giờ đúng, góc tạo bởi kim giờ và kim phút của đồng hồ là góc gì?
\choice
{\True Góc nhọn (góc không vuông bé hơn góc vuông)}
{Góc vuông}
{Góc tù (góc không vuông lớn hơn góc vuông)}
{Góc bẹt}
\loigiai{
Lúc 2 giờ đúng, khoảng cách giữa hai kim là 2 khoảng giờ ($60^\circ$), bé hơn góc vuông nên là góc nhọn.
\textbf{Đáp án của câu hỏi là A}
}
\end{ex}

\begin{ex} %[3H1V18]
Lúc $4$ giờ đúng, góc tạo bởi kim giờ và kim phút của đồng hồ là
\choice
{Góc vuông}
{Góc không vuông bé hơn góc vuông (góc nhọn)}
{\True Góc không vuông lớn hơn góc vuông (góc tù)}
{Góc bẹt}
\loigiai{
Lúc 4 giờ đúng, góc giữa kim giờ và phút là 4 khoảng giờ ($120^\circ$), lớn hơn góc vuông nên là góc tù.
\textbf{Đáp án của câu hỏi là C}
}
\end{ex}

\begin{ex} %[3H1V18]
Bác thợ mộc ghép hai tấm gỗ hình chữ nhật lại với nhau để tạo thành một hình chữ nhật lớn hơn. Hỏi hình chữ nhật mới có tất cả bao nhiêu góc vuông?
\choice
{$8$ góc vuông}
{\True $4$ góc vuông}
{$6$ góc vuông}
{$2$ góc vuông}
\loigiai{
Bất kì hình chữ nhật nào cũng luôn có đúng 4 góc vuông ở 4 đỉnh.
\textbf{Đáp án của câu hỏi là B}
}
\end{ex}

\begin{ex} %[3H1V18]
Hình chữ nhật $ABCD$ có kẻ đường chéo $AC$. Hình vẽ khi đó xuất hiện bao nhiêu góc vuông?
\choice
{$2$ góc vuông}
{$6$ góc vuông}
{$8$ góc vuông}
{\True $4$ góc vuông}
\loigiai{
Hình chữ nhật có 4 góc vuông tại $A, B, C, D$. Đường chéo $AC$ chia góc $A$ và $C$ thành các góc nhọn, nhưng ở đỉnh $B$ và $D$ vẫn là góc vuông, còn góc toàn phần tại đỉnh $A$ và $C$ vẫn là góc vuông. Tổng cộng vẫn có 4 góc vuông.
\textbf{Đáp án của câu hỏi là D}
}
\end{ex}

\begin{ex} %[3H1V18]
Cho một góc vuông đỉnh $O$, cạnh $OA, OB$. Kẻ thêm một tia $OC$ nằm giữa hai cạnh $OA$ và $OB$. Khi đó trên hình có bao nhiêu góc nhọn?
\choice
{\True $2$ góc nhọn}
{$1$ góc nhọn}
{$3$ góc nhọn}
{$0$ góc nhọn}
\loigiai{
Góc ban đầu là góc vuông $\widehat{AOB}$. Tia $OC$ nằm trong chia thành hai góc $\widehat{AOC}$ và $\widehat{COB}$, cả hai góc này đều nhỏ hơn góc vuông nên là 2 góc nhọn.
\textbf{Đáp án của câu hỏi là A}
}
\end{ex}

\begin{ex} %[3H1V18]
Trong một ngày đêm (24 giờ), kim giờ và kim phút của đồng hồ tạo thành góc vuông đúng bao nhiêu lần vào các giờ đúng?
\choice
{$2$ lần}
{$12$ lần}
{\True $4$ lần}
{$24$ lần}
\loigiai{
Vào các giờ đúng, góc vuông chỉ xuất hiện vào lúc 3 giờ và 9 giờ. Một ngày đêm có hai lượt 3 giờ (3 giờ sáng, 15 giờ) và hai lượt 9 giờ (9 giờ sáng, 21 giờ), tổng cộng có 4 lần.
\textbf{Đáp án của câu hỏi là C}
}
\end{ex}

\textbf{Phần 2. Trắc nghiệm đúng sai}
\setcounterpageref{ex}{0}

\begin{ex} %[3H1V18]
Xét tính đúng sai của các khẳng định về góc:
\choiceTF
{\True Góc vuông là góc bằng góc vuông của thước ê-ke.}
{Một hình tứ giác có thể có 5 góc vuông.}
{\True Góc bé hơn góc vuông được gọi là góc nhọn.}
{Lúc 6 giờ đúng, hai kim đồng hồ tạo thành một góc vuông.}
\loigiai{
\begin{itemize}
    \item Ý a) Đúng: Định nghĩa nhận biết góc vuông bằng ê-ke.
    \item Ý b) Sai: Tứ giác chỉ có đúng 4 đỉnh, nhiều nhất 4 góc vuông.
    \item Ý c) Đúng: Kiến thức hình học cơ bản.
    \item Ý d) Sai: Lúc 6 giờ hai kim thẳng hàng tạo thành góc bẹt ($180^\circ$).
\end{itemize}
\begin{itemize}
    \item a) \textbf{Đúng}.
    \item b) \textbf{Sai}.
    \item c) \textbf{Đúng}.
    \item d) \textbf{Sai}.
\end{itemize}
}
\end{ex}

\begin{ex} %[3H1V18]
Cho hình chữ nhật $ABCD$ có điểm $M$ nằm trên cạnh $CD$:
\choiceTF
{\True Tứ giác $ABCD$ có 4 góc vuông.}
{\True Góc đỉnh $B$, cạnh $BA$ và $BC$ là góc vuông.}
{Góc đỉnh $M$, cạnh $MA$ và $MB$ luôn luôn là góc vuông.}
{\True Góc đỉnh $A$, cạnh $AB$ và $AD$ là góc vuông.}
\loigiai{
\begin{itemize}
    \item Ý a) Đúng: Đặc điểm của hình chữ nhật.
    \item Ý b) Đúng: Góc đỉnh của hình chữ nhật.
    \item Ý c) Sai: Không nhất thiết phải vuông, phụ thuộc vị trí điểm $M$.
    \item Ý d) Đúng: Góc đỉnh của hình chữ nhật.
\end{itemize}
\begin{itemize}
    \item a) \textbf{Đúng}.
    \item b) \textbf{Đúng}.
    \item c) \textbf{Sai}.
    \item d) \textbf{Đúng}.
\end{itemize}
}
\end{ex}

\textbf{Phần 3. Trả lời ngắn}
\setcounterpageref{ex}{0}

\begin{ex} %[3H1V18]
Một hình vuông có bao nhiêu góc vuông?
\shortans{4}
\loigiai{
\textbf{Đáp số:} $4$\\
Hình vuông có 4 góc vuông ở 4 đỉnh.
}
\end{ex}

\begin{ex} %[3H1V18]
Một hình tam giác vuông có bao nhiêu góc không vuông?
\shortans{2}
\loigiai{
\textbf{Đáp số:} $2$\\
Tam giác có 3 góc, trong đó có 1 góc vuông thì còn lại $3 - 1 = 2$ góc không vuông (góc nhọn).
}
\end{ex}

\begin{ex} %[3H1V18]
Từ 1 giờ đúng đến 5 giờ đúng, có bao nhiêu lần kim giờ và kim phút tạo thành một góc vuông vào đúng giờ tròn?
\shortans{1}
\loigiai{
\textbf{Đáp số:} $1$\\
Chỉ có duy nhất lúc 3 giờ đúng tạo thành góc vuông.
}
\end{ex}

\begin{ex} %[3H1V18]
Đếm số góc vuông đỉnh $O$ được tạo bởi 4 tia xuất phát từ $O$ đôi một vuông góc (tạo thành 2 đường thẳng vuông góc cắt nhau).
\shortans{4}
\loigiai{
\textbf{Đáp số:} $4$\\
Hai đường thẳng vuông góc cắt nhau tại $O$ tạo thành đúng 4 góc vuông.
}
\end{ex}

\begin{ex} %[3H1V18]
Một tờ giấy hình chữ nhật có 4 góc vuông. Dùng kéo cắt đứt một góc chéo của tờ giấy đó. Hỏi mảnh giấy lớn còn lại có bao nhiêu góc?
\shortans{5}
\loigiai{
\textbf{Đáp số:} $5$\\
Cắt vát một góc sẽ làm mất 1 góc nhưng tạo thêm 2 góc mới: $4 - 1 + 2 = 5$ góc (hình ngũ giác).
}
\end{ex}

\begin{ex} %[3H1V18]
Có bao nhiêu chữ cái in hoa trong các chữ cái sau chứa ít nhất một góc vuông: E, F, H, L, T, O, V?
\shortans{5}
\loigiai{
\textbf{Đáp số:} $5$\\
Các chữ E, F, H, L, T đều chứa các đoạn thẳng vuông góc (có góc vuông). Chữ O không có góc, chữ V là góc nhọn. Vậy có 5 chữ cái.
}
\end{ex}

\newpage

\TaoPhieuKiemTra{Chủ đề 3 -- Làm quen với hình phẳng, hình khối}{19}{Hình tam giác, hình tứ giác, hình chữ nhật, hình vuông}{https://loptoanthaylinh.vn/thi-online/toan-3/bai-19}
\textbf{Phần 1. Trắc nghiệm lựa chọn}
\setcounterpageref{ex}{0}

\begin{ex} %[3H1V19]
Hình vuông có đặc điểm nổi bật nào sau đây?
\choice
{\True Có 4 góc vuông và 4 cạnh có độ dài bằng nhau}
{Có 4 góc vuông và 2 cạnh dài bằng nhau}
{Có 3 góc vuông và 4 cạnh bằng nhau}
{Có 4 cạnh bằng nhau nhưng không có góc vuông}
\loigiai{
Hình vuông có 4 góc vuông và 4 cạnh dài bằng nhau.
\textbf{Đáp án của câu hỏi là A}
}
\end{ex}

\begin{ex} %[3H1V19]
Cho một hình tam giác có độ dài các cạnh lần lượt là $15\text{ cm}, 20\text{ cm}, 25\text{ cm}$. Chu vi của hình tam giác đó là
\choice
{$50\text{ cm}$}
{$55\text{ cm}$}
{\True $60\text{ cm}$}
{$65\text{ cm}$}
\loigiai{
Chu vi hình tam giác là: $15 + 20 + 25 = 60\text{ (cm)}$.
\textbf{Đáp án của câu hỏi là C}
}
\end{ex}

\begin{ex} %[3H1V19]
Một hình chữ nhật có chiều dài $24\text{ cm}$, chiều rộng bằng một nửa chiều dài. Chu vi hình chữ nhật đó là
\choice
{$48\text{ cm}$}
{\True $72\text{ cm}$}
{$36\text{ cm}$}
{$60\text{ cm}$}
\loigiai{
Chiều rộng hình chữ nhật: $24 : 2 = 12\text{ cm}$. Chu vi: $(24 + 12) \times 2 = 72\text{ cm}$.
\textbf{Đáp án của câu hỏi là B}
}
\end{ex}

\begin{ex} %[3H1V19]
Một sợi dây thép dài $36\text{ cm}$ được uốn thành một hình vuông. Độ dài mỗi cạnh của hình vuông đó là
\choice
{$6\text{ cm}$}
{$12\text{ cm}$}
{$8\text{ cm}$}
{\True $9\text{ cm}$}
\loigiai{
Độ dài mỗi cạnh của hình vuông là: $36 : 4 = 9\text{ (cm)}$.
\textbf{Đáp án của câu hỏi là D}
}
\end{ex}

\begin{ex} %[3H1V19]
Cho một hình chữ nhật có chu vi bằng chu vi của hình vuông cạnh $8\text{ cm}$. Biết chiều rộng của hình chữ nhật là $6\text{ cm}$. Chiều dài của hình chữ nhật đó là
\choice
{\True $10\text{ cm}$}
{$12\text{ cm}$}
{$8\text{ cm}$}
{$14\text{ cm}$}
\loigiai{
Chu vi hình vuông là: $8 \times 4 = 32\text{ (cm)}$. Nửa chu vi hình chữ nhật là: $32 : 2 = 16\text{ (cm)}$. Chiều dài hình chữ nhật là: $16 - 6 = 10\text{ (cm)}$.
\textbf{Đáp án của câu hỏi là A}
}
\end{ex}

\begin{ex} %[3H1V19]
Người ta ghép hai hình vuông nhỏ cạnh $5\text{ cm}$ cạnh nhau để tạo thành một hình chữ nhật. Chu vi của hình chữ nhật tạo thành là
\choice
{$40\text{ cm}$}
{$25\text{ cm}$}
{\True $30\text{ cm}$}
{$35\text{ cm}$}
\loigiai{
Hình chữ nhật có chiều dài $5 + 5 = 10\text{ cm}$ và chiều rộng $5\text{ cm}$. Chu vi là: $(10 + 5) \times 2 = 30\text{ (cm)}$.
\textbf{Đáp án của câu hỏi là C}
}
\end{ex}

\begin{ex} %[3H1V19]
Khẳng định nào sau đây là ĐÚNG về mối quan hệ giữa hình vuông và hình chữ nhật?
\choice
{Mỗi hình chữ nhật đều là một hình vuông}
{\True Mỗi hình vuông đều là một hình chữ nhật đặc biệt có 4 cạnh bằng nhau}
{Hình chữ nhật có 4 cạnh bằng nhau}
{Hình vuông không có góc vuông}
\loigiai{
Hình vuông có đủ 4 góc vuông và có các cặp cạnh đối diện bằng nhau (thậm chí cả 4 cạnh đều bằng nhau) nên hình vuông là trường hợp đặc biệt của hình chữ nhật.
\textbf{Đáp án của câu hỏi là B}
}
\end{ex}

\begin{ex} %[3H1V19]
Cho tam giác $ABC$ có cạnh $AB = AC = 18\text{ cm}$, cạnh $BC$ ngắn hơn cạnh $AB$ là $6\text{ cm}$. Chu vi của tam giác $ABC$ là
\choice
{$42\text{ cm}$}
{$54\text{ cm}$}
{$36\text{ cm}$}
{\True $48\text{ cm}$}
\loigiai{
Cạnh $BC$ dài: $18 - 6 = 12\text{ (cm)}$. Chu vi tam giác $ABC$ là: $18 + 18 + 12 = 48\text{ (cm)}$.
\textbf{Đáp án của câu hỏi là D}
}
\end{ex}

\begin{ex} %[3H1V19]
Một mảnh vườn hình chữ nhật có chiều dài $30\text{ m}$, nếu giảm chiều dài đi $6\text{ m}$ thì mảnh vườn trở thành hình vuông. Chu vi của mảnh vườn ban đầu là
\choice
{\True $108\text{ m}$}
{$96\text{ m}$}
{$120\text{ m}$}
{$84\text{ m}$}
\loigiai{
Khi giảm chiều dài 6 m thành hình vuông thì cạnh hình vuông là: $30 - 6 = 24\text{ m}$. Chiều rộng mảnh vườn chính bằng $24\text{ m}$. Chu vi ban đầu là: $(30 + 24) \times 2 = 108\text{ (m)}$.
\textbf{Đáp án của câu hỏi là A}
}
\end{ex}

\begin{ex} %[3H1V19]
Một hình tứ giác có độ dài 4 cạnh là các số tự nhiên liên tiếp tăng dần. Biết chu vi của hình tứ giác đó là $34\text{ cm}$. Cạnh dài nhất của hình tứ giác là
\choice
{$9\text{ cm}$}
{$11\text{ cm}$}
{\True $10\text{ cm}$}
{$8\text{ cm}$}
\loigiai{
Gọi 4 cạnh là $a, a+1, a+2, a+3$. Tổng chu vi: $4 \times a + 6 = 34 \Rightarrow 4 \times a = 28 \Rightarrow a = 7$. Cạnh dài nhất là $7 + 3 = 10\text{ cm}$.
\textbf{Đáp án của câu hỏi là C}
}
\end{ex}

\begin{ex} %[3H1V19]
Có bao nhiêu đoạn thẳng nối từ đỉnh một tứ giác đến các đỉnh còn lại không kể các cạnh của tứ giác đó?
\choice
{$2$ đoạn thẳng}
{\True $1$ đoạn thẳng (đường chéo)}
{$3$ đoạn thẳng}
{$4$ đoạn thẳng}
\loigiai{
Từ 1 đỉnh của tứ giác, nối đến 2 đỉnh liền kề là 2 cạnh, chỉ có 1 đỉnh đối diện tạo thành 1 đường chéo.
\textbf{Đáp án của câu hỏi là B}
}
\end{ex}

\begin{ex} %[3H1V19]
Người ta cắt một hình chữ nhật có chiều dài $15\text{ cm}$, chiều rộng $9\text{ cm}$ để lấy ra một hình vuông có cạnh lớn nhất có thể. Chu vi của phần hình chữ nhật còn lại là
\choice
{$24\text{ cm}$}
{$36\text{ cm}$}
{$18\text{ cm}$}
{\True $30\text{ cm}$}
\loigiai{
Hình vuông lớn nhất có cạnh bằng chiều rộng là $9\text{ cm}$. Phần còn lại là hình chữ nhật có kích thước $9\text{ cm}$ và $15 - 9 = 6\text{ cm}$. Chu vi phần còn lại là: $(9 + 6) \times 2 = 30\text{ (cm)}$.
\textbf{Đáp án của câu hỏi là D}
}
\end{ex}

\textbf{Phần 2. Trắc nghiệm đúng sai}
\setcounterpageref{ex}{0}

\begin{ex} %[3H1V19]
Xét tính đúng sai của các nhận định về các hình phẳng:
\choiceTF
{\True Hình vuông có 4 góc vuông và 4 cạnh dài bằng nhau.}
{Mọi hình tứ giác đều có 4 góc vuông.}
{\True Tam giác có đúng 3 cạnh và 3 góc.}
{Hình chữ nhật có 4 cạnh bằng nhau.}
\loigiai{
\begin{itemize}
    \item Ý a) Đúng: Đặc điểm định nghĩa của hình vuông.
    \item Ý b) Sai: Chỉ có hình vuông và hình chữ nhật mới có 4 góc vuông.
    \item Ý c) Đúng: Định nghĩa hình tam giác.
    \item Ý d) Sai: Hình chữ nhật nói chung chỉ có các cặp cạnh đối diện bằng nhau.
\end{itemize}
\begin{itemize}
    \item a) \textbf{Đúng}.
    \item b) \textbf{Sai}.
    \item c) \textbf{Đúng}.
    \item d) \textbf{Sai}.
\end{itemize}
}
\end{ex}

\begin{ex} %[3H1V19]
Một hình chữ nhật có chiều dài $16\text{ cm}$, chiều rộng $8\text{ cm}$:
\choiceTF
{\True Chiều dài gấp đôi chiều rộng.}
{\True Chu vi của hình chữ nhật đó là $48\text{ cm}$.}
{Nếu gấp đôi chiều rộng thì hình chữ nhật trở thành hình vuông có cạnh $8\text{ cm}$.}
{\True Có thể chia hình chữ nhật đó thành 2 hình vuông cạnh $8\text{ cm}$.}
\loigiai{
\begin{itemize}
    \item Ý a) Đúng: $16 : 8 = 2$ lần.
    \item Ý b) Đúng: $(16 + 8) \times 2 = 48\text{ cm}$.
    \item Ý c) Sai: Nếu gấp đôi chiều rộng thì chiều rộng thành $16\text{ cm}$, bằng chiều dài nên thành hình vuông cạnh $16\text{ cm}$, không phải $8\text{ cm}$.
    \item Ý d) Đúng: $16 : 8 = 2$ hình vuông.
\end{itemize}
\begin{itemize}
    \item a) \textbf{Đúng}.
    \item b) \textbf{Đúng}.
    \item c) \textbf{Sai}.
    \item d) \textbf{Đúng}.
\end{itemize}
}
\end{ex}

\textbf{Phần 3. Trả lời ngắn}
\setcounterpageref{ex}{0}

\begin{ex} %[3H1V19]
Tính chu vi hình tam giác đều có độ dài mỗi cạnh bằng $15\text{ cm}$.
\shortans{45}
\loigiai{
\textbf{Đáp số:} $45$\\
Chu vi là: $15 \times 3 = 45\text{ (cm)}$.
}
\end{ex}

\begin{ex} %[3H1V19]
Một hình vuông có chu vi $64\text{ cm}$. Tính độ dài cạnh của hình vuông đó theo xăng-ti-mét.
\shortans{16}
\loigiai{
\textbf{Đáp số:} $16$\\
Cạnh hình vuông là: $64 : 4 = 16\text{ (cm)}$.
}
\end{ex}

\begin{ex} %[3H1V19]
Một hình chữ nhật có chu vi $50\text{ cm}$, chiều dài là $18\text{ cm}$. Tính chiều rộng của hình chữ nhật theo xăng-ti-mét.
\shortans{7}
\loigiai{
\textbf{Đáp số:} $7$\\
Nửa chu vi: $50 : 2 = 25\text{ cm}$. Chiều rộng: $25 - 18 = 7\text{ (cm)}$.
}
\end{ex}

\begin{ex} %[3H1V19]
Một mảnh đất hình chữ nhật có chiều dài $25\text{ m}$, chiều rộng ngắn hơn chiều dài $7\text{ m}$. Tính chu vi của mảnh đất đó theo mét.
\shortans{86}
\loigiai{
\textbf{Đáp số:} $86$\\
Chiều rộng: $25 - 7 = 18\text{ m}$. Chu vi: $(25 + 18) \times 2 = 86\text{ (m)}$.
}
\end{ex}

\begin{ex} %[3H1V19]
Bốn que tính dài bằng nhau ghép thành một hình vuông có chu vi $32\text{ cm}$. Hỏi dùng 6 que tính như vậy ghép thành một hình chữ nhật thì chu vi của hình chữ nhật ghép được là bao nhiêu xăng-ti-mét?
\shortans{48}
\loigiai{
\textbf{Đáp số:} $48$\\
Mỗi que tính dài: $32 : 4 = 8\text{ cm}$. Chu vi hình ghép bởi 6 que tính là: $6 \times 8 = 48\text{ (cm)}$.
}
\end{ex}

\begin{ex} %[3H1V19]
Người ta ghép 3 hình vuông bằng nhau cạnh $4\text{ cm}$ thành một hàng ngang để được một hình chữ nhật. Tính chu vi hình chữ nhật đó theo xăng-ti-mét.
\shortans{32}
\loigiai{
\textbf{Đáp số:} $32$\\
Hình chữ nhật có chiều dài $4 \times 3 = 12\text{ cm}$, chiều rộng $4\text{ cm}$. Chu vi: $(12 + 4) \times 2 = 32\text{ (cm)}$.
}
\end{ex}

\newpage

\TaoPhieuKiemTra{Chủ đề 3 -- Làm quen với hình phẳng, hình khối}{20}{Thực hành vẽ hình và trang trí}{https://loptoanthaylinh.vn/thi-online/toan-3/bai-20}
\textbf{Phần 1. Trắc nghiệm lựa chọn}
\setcounterpageref{ex}{0}

\begin{ex} %[3H1V20]
Để vẽ được một hình tròn có bán kính $4\text{ cm}$ trên giấy, ta cần chỉnh khoảng cách giữa đầu kim và đầu bút chì của compa bằng
\choice
{$8\text{ cm}$}
{$2\text{ cm}$}
{$16\text{ cm}$}
{\True $4\text{ cm}$}
\loigiai{
Khoảng cách giữa hai đầu của compa đúng bằng bán kính hình tròn cần vẽ, tức là $4\text{ cm}$.
\textbf{Đáp án của câu hỏi là D}
}
\end{ex}

\begin{ex} %[3H1V20]
Bạn Nam dùng các que tính dài bằng nhau để ghép thành 3 hình tam giác riêng biệt không chung cạnh. Nam cần dùng ít nhất bao nhiêu que tính?
\choice
{$6$ que tính}
{\True $9$ que tính}
{$7$ que tính}
{$8$ que tính}
\loigiai{
Mỗi hình tam giác cần 3 que tính, 3 tam giác riêng biệt cần: $3 \times 3 = 9$ que tính.
\textbf{Đáp án của câu hỏi là B}
}
\end{ex}

\begin{ex} %[3H1V20]
Nếu cho phép các hình tam giác có cạnh chung với nhau, hỏi dùng ít nhất bao nhiêu que tính dài bằng nhau để ghép được 2 hình tam giác?
\choice
{$6$ que tính}
{$4$ que tính}
{\True $5$ que tính}
{$3$ que tính}
\loigiai{
Hai hình tam giác ghép chung 1 cạnh ở giữa sẽ dùng: $3 + 3 - 1 = 5$ que tính.
\textbf{Đáp án của câu hỏi là C}
}
\end{ex}

\begin{ex} %[3H1V20]
Để vẽ một góc vuông đỉnh $O$ bằng ê-ke, bước nào sau đây là chuẩn xác nhất?
\choice
{\True Đặt đỉnh góc vuông của ê-ke trùng với điểm $O$, vẽ hai cạnh theo hai mép góc vuông của thước}
{Đặt mép thước bất kì vào điểm $O$}
{Vẽ một đường tròn rồi lấy 1 phần 4}
{Dùng thước thẳng vẽ ngẫu nhiên hai đoạn thẳng}
\loigiai{
Cách vẽ góc vuông chuẩn: đặt đỉnh góc vuông ê-ke vào đỉnh góc, vạch 2 tia theo 2 cạnh góc vuông của ê-ke.
\textbf{Đáp án của câu hỏi là A}
}
\end{ex}

\begin{ex} %[3H1V20]
Người ta gấp một mảnh giấy hình vuông làm đôi theo đường chéo, sau đó gấp đôi tiếp một lần nữa theo đường cao. Mảnh giấy nhận được có hình dạng là
\choice
{Hình vuông nhỏ}
{Hình chữ nhật}
{Hình tứ giác không đều}
{\True Hình tam giác có 1 góc vuông}
\loigiai{
Gấp đôi hình vuông theo đường chéo được tam giác vuông cân. Gấp đôi lần nữa vẫn được hình tam giác vuông cân nhỏ hơn.
\textbf{Đáp án của câu hỏi là D}
}
\end{ex}

\begin{ex} %[3H1V20]
Một hoa văn trang trí hình vuông được chia thành 4 ô vuông nhỏ bằng nhau. Người ta tô màu đỏ vào 2 ô đối đỉnh, 2 ô còn lại tô màu xanh. Phần ô vuông màu đỏ chiếm bao nhiêu phần của hoa văn?
\choice
{Một phần tư ($\dfrac{1}{4}$)}
{\True Một phần hai ($\dfrac{1}{2}$)}
{Ba phần tư ($\dfrac{3}{4}$)}
{Một phần ba ($\dfrac{1}{3}$)}
\loigiai{
Có tất cả 4 ô vuông nhỏ, 2 ô màu đỏ. Vậy phần màu đỏ chiếm $2 : 4 = \dfrac{1}{2}$ hình vuông hoa văn.
\textbf{Đáp án của câu hỏi là B}
}
\end{ex}

\begin{ex} %[3H1V20]
Dùng 12 que diêm bằng nhau để xếp thành một hình chữ nhật có chu vi lớn nhất có thể. Hình chữ nhật đó có chiều dài gồm mấy que diêm?
\choice
{$4$ que diêm}
{$3$ que diêm}
{\True $5$ que diêm}
{$6$ que diêm}
\loigiai{
Nửa chu vi gồm $12 : 2 = 6$ que diêm. Chiều dài lớn nhất có thể là $6 - 1 = 5$ que diêm (khi chiều rộng là 1 que diêm).
\textbf{Đáp án của câu hỏi là C}
}
\end{ex}

\begin{ex} %[3H1V20]
Trong một hình vẽ gồm một đường tròn tâm $O$ và một hình vuông nằm bên trong sao cho 4 đỉnh của hình vuông đều nằm trên đường tròn. Khi đó tâm $O$ của hình tròn nằm ở vị trí nào của hình vuông?
\choice
{\True Giao điểm của hai đường chéo hình vuông}
{Tại một đỉnh bất kì của hình vuông}
{Tại trung điểm của một cạnh hình vuông}
{Nằm ngoài hình vuông}
\loigiai{
Hình tròn đi qua 4 đỉnh của hình vuông có tâm chính là giao điểm của hai đường chéo hình vuông.
\textbf{Đáp án của câu hỏi là A}
}
\end{ex}

\begin{ex} %[3H1V20]
Muốn vẽ 4 hình tròn bằng nhau có các tâm lần lượt là 4 đỉnh của một hình chữ nhật sao cho mỗi hình tròn tiếp xúc với 2 hình tròn liền kề. Điều kiện của hình chữ nhật đó là
\choice
{Hình chữ nhật có chiều dài gấp đôi chiều rộng}
{Hình chữ nhật có chiều dài gấp 3 chiều rộng}
{Bất kì hình chữ nhật nào cũng vẽ được}
{\True Hình chữ nhật đó phải là hình vuông}
\loigiai{
Để 4 hình tròn bằng nhau tiếp xúc đôi một theo cạnh thì khoảng cách giữa các đỉnh liên tiếp phải bằng nhau (bằng 2 lần bán kính), do đó các cạnh phải bằng nhau, tức là hình vuông.
\textbf{Đáp án của câu hỏi là D}
}
\end{ex}

\begin{ex} %[3H1V20]
Bạn Hùng cắt một tờ giấy hình tròn thành 4 phần bằng nhau bằng 2 nhát cắt thẳng đi qua tâm. Góc ở tâm của mỗi phần giấy nhận được là
\choice
{Góc nhọn}
{\True Góc vuông}
{Góc tù}
{Góc bẹt}
\loigiai{
Hai nhát cắt thẳng qua tâm chia hình tròn thành 4 phần bằng nhau thì hai đường cắt vuông góc với nhau, góc ở mỗi phần là góc vuông ($360^\circ : 4 = 90^\circ$).
\textbf{Đáp án của câu hỏi là B}
}
\end{ex}

\begin{ex} %[3H1V20]
Một họa sĩ vẽ một bức tranh ghép từ các hình tam giác đều có cạnh $2\text{ cm}$. Hàng thứ nhất có 1 tam giác, hàng thứ hai có 3 tam giác, hàng thứ ba có 5 tam giác. Hỏi bức tranh gồm tất cả bao nhiêu hình tam giác đều nhỏ?
\choice
{$8$ hình tam giác}
{$10$ hình tam giác}
{\True $9$ hình tam giác}
{$12$ hình tam giác}
\loigiai{
Tổng số hình tam giác nhỏ là: $1 + 3 + 5 = 9$ hình tam giác.
\textbf{Đáp án của câu hỏi là C}
}
\end{ex}

\begin{ex} %[3H1V20]
Dùng thước và compa vẽ một đường tròn tâm $O$ bán kính $5\text{ cm}$. Trên đường tròn lấy điểm $A$. Khoảng cách từ điểm $A$ đến tâm $O$ đo được là
\choice
{\True $5\text{ cm}$}
{$10\text{ cm}$}
{$2\text{ cm và } 5\text{ mm}$}
{$0\text{ cm}$}
\loigiai{
Mọi điểm nằm trên đường tròn đều cách tâm một khoảng bằng bán kính, tức là $5\text{ cm}$.
\textbf{Đáp án của câu hỏi là A}
}
\end{ex}

\textbf{Phần 2. Trắc nghiệm đúng sai}
\setcounterpageref{ex}{0}

\begin{ex} %[3H1V20]
Xét tính đúng sai của các phát biểu về vẽ hình học phẳng:
\choiceTF
{\True Khi dùng compa vẽ đường tròn, chân trụ có kim phải được giữ cố định tại tâm.}
{Chỉ dùng một thước thẳng có vạch đo có thể vẽ chính xác góc vuông mà không cần ê-ke.}
{\True Dùng 4 que tính bằng nhau có thể xếp thành một hình vuông hoặc một hình thoi.}
{Mọi hình vẽ trang trí đều phải có góc vuông.}
\loigiai{
\begin{itemize}
    \item Ý a) Đúng: Quy tắc dùng compa.
    \item Ý b) Sai: Sai vì thước thẳng không đo góc trực tiếp được góc vuông.
    \item Ý c) Đúng: Đúng vì 4 cạnh bằng nhau.
    \item Ý d) Sai: Sai vì có nhiều hoa văn tròn, uốn lượn không có góc vuông.
\end{itemize}
\begin{itemize}
    \item a) \textbf{Đúng}.
    \item b) \textbf{Sai}.
    \item c) \textbf{Đúng}.
    \item d) \textbf{Sai}.
\end{itemize}
}
\end{ex}

\begin{ex} %[3H1V20]
Một tờ giấy hình chữ nhật kích thước $12\text{ cm} \times 8\text{ cm}$ được cắt ra để trang trí:
\choiceTF
{\True Hình vuông lớn nhất có thể cắt ra từ tờ giấy có cạnh bằng $8\text{ cm}$.}
{\True Phần còn lại sau khi cắt hình vuông lớn nhất là một hình chữ nhật kích thước $8\text{ cm} \times 4\text{ cm}$.}
{Có thể cắt tờ giấy thành 3 hình vuông bằng nhau mà không thừa mảnh nào.}
{\True Từ phần còn lại ($8\text{ cm} \times 4\text{ cm}$) ta cắt tiếp được đúng 2 hình vuông cạnh $4\text{ cm}$.}
\loigiai{
\begin{itemize}
    \item Ý a) Đúng: Cạnh hình vuông lớn nhất bằng chiều rộng.
    \item Ý b) Đúng: $12 - 8 = 4\text{ cm}$.
    \item Ý c) Sai: Sai vì diện tích 96 không chia hết cho 3 hình vuông nguyên cạnh.
    \item Ý d) Đúng: $8 : 4 = 2$ hình vuông.
\end{itemize}
\begin{itemize}
    \item a) \textbf{Đúng}.
    \item b) \textbf{Đúng}.
    \item c) \textbf{Sai}.
    \item d) \textbf{Đúng}.
\end{itemize}
}
\end{ex}

\textbf{Phần 3. Trả lời ngắn}
\setcounterpageref{ex}{0}

\begin{ex} %[3H1V20]
Cần ít nhất bao nhiêu que diêm bằng nhau để ghép thành 3 hình vuông riêng biệt không chung cạnh?
\shortans{12}
\loigiai{
\textbf{Đáp số:} $12$\\
Mỗi hình vuông cần 4 que diêm: $3 \times 4 = 12$ que.
}
\end{ex}

\begin{ex} %[3H1V20]
Dùng các que tính bằng nhau để xếp thành 2 hình vuông có chung 1 cạnh. Hỏi cần dùng tất cả bao nhiêu que tính?
\shortans{7}
\loigiai{
\textbf{Đáp số:} $7$\\
Hai hình vuông ghép chung 1 cạnh cần: $4 + 4 - 1 = 7$ que tính.
}
\end{ex}

\begin{ex} %[3H1V20]
Để vẽ một hình tròn có đường kính $16\text{ cm}$, bạn Mai cần mở rộng hai đầu compa một khoảng cách bằng bao nhiêu xăng-ti-mét?
\shortans{8}
\loigiai{
\textbf{Đáp số:} $8$\\
Bán kính compa cần mở là: $16 : 2 = 8\text{ (cm)}$.
}
\end{ex}

\begin{ex} %[3H1V20]
Một đường gấp khúc gồm 5 đoạn thẳng bằng nhau có độ dài mỗi đoạn là $6\text{ cm}$. Tính tổng độ dài đường gấp khúc đó theo xăng-ti-mét.
\shortans{30}
\loigiai{
\textbf{Đáp số:} $30$\\
Độ dài đường gấp khúc: $5 \times 6 = 30\text{ (cm)}$.
}
\end{ex}

\begin{ex} %[3H1V20]
Một dải băng giấy dài $36\text{ cm}$ được gấp làm 4 đoạn bằng nhau dích dắc để tạo hình trang trí. Hỏi mỗi đoạn gấp dài bao nhiêu xăng-ti-mét?
\shortans{9}
\loigiai{
\textbf{Đáp số:} $9$\\
Độ dài mỗi đoạn gấp là: $36 : 4 = 9\text{ (cm)}$.
}
\end{ex}

\begin{ex} %[3H1V20]
Dùng 6 hình tam giác đều nhỏ bằng nhau ghép lại có thể tạo thành một hình lục giác đều. Nếu mỗi tam giác nhỏ có chu vi $18\text{ cm}$ thì độ dài cạnh của lục giác ghép được là bao nhiêu xăng-ti-mét?
\shortans{6}
\loigiai{
\textbf{Đáp số:} $6$\\
Cạnh của tam giác đều nhỏ là: $18 : 3 = 6\text{ cm}$. Cạnh của lục giác đều chính bằng cạnh của tam giác đều nhỏ, bằng $6\text{ cm}$.
}
\end{ex}

\newpage

\TaoPhieuKiemTra{Chủ đề 3 -- Làm quen với hình phẳng, hình khối}{21}{Khối lập phương, khối hộp chữ nhật}{https://loptoanthaylinh.vn/thi-online/toan-3/bai-21}
\textbf{Phần 1. Trắc nghiệm lựa chọn}
\setcounterpageref{ex}{0}

\begin{ex} %[3H2V21]
Một khối lập phương có bao nhiêu đỉnh, bao nhiêu cạnh và bao nhiêu mặt?
\choice
{$6$ đỉnh, $12$ cạnh, $8$ mặt}
{$8$ đỉnh, $6$ cạnh, $12$ mặt}
{\True $8$ đỉnh, $12$ cạnh, $6$ mặt}
{$12$ đỉnh, $8$ cạnh, $6$ mặt}
\loigiai{
Khối lập phương có 8 đỉnh, 12 cạnh bằng nhau và 6 mặt đều là hình vuông.
\textbf{Đáp án của câu hỏi là C}
}
\end{ex}

\begin{ex} %[3H2V21]
Các mặt của khối lập phương đều là hình gì?
\choice
{\True Hình vuông}
{Hình chữ nhật}
{Hình tam giác}
{Hình tròn}
\loigiai{
Khối lập phương có 6 mặt đều là các hình vuông bằng nhau.
\textbf{Đáp án của câu hỏi là A}
}
\end{ex}

\begin{ex} %[3H2V21]
Một con xúc xắc chuẩn hình lập phương có tổng số chấm ở hai mặt đối diện luôn luôn bằng
\choice
{$6$}
{$8$}
{$9$}
{\True $7$}
\loigiai{
Quy ước chuẩn của xúc xắc: 1 đối diện 6, 2 đối diện 5, 3 đối diện 4; tổng số chấm luôn bằng 7.
\textbf{Đáp án của câu hỏi là D}
}
\end{ex}

\begin{ex} %[3H2V21]
Khối hộp chữ nhật có tất cả bao nhiêu mặt?
\choice
{$4$ mặt}
{\True $6$ mặt}
{$8$ mặt}
{$12$ mặt}
\loigiai{
Khối hộp chữ nhật có 6 mặt (2 mặt đáy và 4 mặt bên).
\textbf{Đáp án của câu hỏi là B}
}
\end{ex}

\begin{ex} %[3H2V21]
Dùng các khối lập phương nhỏ cạnh $1\text{ cm}$ để xếp thành một khối lập phương lớn hơn có cạnh $2\text{ cm}$. Số khối lập phương nhỏ cần dùng là
\choice
{$4$ khối}
{$6$ khối}
{\True $8$ khối}
{$16$ khối}
\loigiai{
Mỗi lớp cần $2 \times 2 = 4$ khối, 2 lớp cần: $4 \times 2 = 8$ khối lập phương nhỏ.
\textbf{Đáp án của câu hỏi là C}
}
\end{ex}

\begin{ex} %[3H2V21]
Người ta dùng dây ruy-băng để buộc một hộp quà hình lập phương cạnh $10\text{ cm}$ theo hình chữ thập quanh hộp (nút thắt cần thêm $15\text{ cm}$ ruy-băng). Độ dài sợi dây ruy-băng cần dùng là
\choice
{\True $55\text{ cm}$}
{$40\text{ cm}$}
{$50\text{ cm}$}
{$60\text{ cm}$}
\loigiai{
Buộc chữ thập quanh 4 mặt bên và 2 đáy hết: $10 \times 4 = 40\text{ cm}$. Thêm nút thắt $15\text{ cm}$, tổng chiều dài: $40 + 15 = 55\text{ cm}$.
\textbf{Đáp án của câu hỏi là A}
}
\end{ex}

\begin{ex} %[3H2V21]
Một khối gỗ hình hộp chữ nhật có chiều dài $5\text{ cm}$, chiều rộng $3\text{ cm}$, chiều cao $2\text{ cm}$. Người ta cắt khối gỗ đó thành các khối lập phương nhỏ cạnh $1\text{ cm}$. Số khối lập phương nhỏ cắt được là
\choice
{$10$ khối}
{$20$ khối}
{$25$ khối}
{\True $30$ khối}
\loigiai{
Số khối nhỏ cắt được là: $5 \times 3 \times 2 = 30$ khối lập phương nhỏ.
\textbf{Đáp án của câu hỏi là D}
}
\end{ex}

\begin{ex} %[3H2V21]
Khung của một chiếc lồng bàn hình khối lập phương được làm bằng các nan tre uốn theo các cạnh của nó. Biết cạnh khối lập phương là $15\text{ cm}$. Tổng chiều dài các nan tre cần dùng là
\choice
{$120\text{ cm}$}
{\True $180\text{ cm}$}
{$150\text{ cm}$}
{$90\text{ cm}$}
\loigiai{
Khối lập phương có 12 cạnh bằng nhau. Tổng chiều dài các cạnh: $12 \times 15 = 180\text{ (cm)}$.
\textbf{Đáp án của câu hỏi là B}
}
\end{ex}

\begin{ex} %[3H2V21]
Sơn tất cả các mặt ngoài của một khối lập phương được ghép bởi 8 khối lập phương nhỏ (cạnh $2\text{ cm}$). Hỏi có bao nhiêu khối lập phương nhỏ được sơn đúng 3 mặt?
\choice
{$4$ khối}
{$6$ khối}
{\True $8$ khối}
{$0$ khối}
\loigiai{
Cả 8 khối nhỏ đều nằm ở 8 đỉnh của khối lập phương lớn nên mỗi khối nhỏ đều được sơn đúng 3 mặt ngoài.
\textbf{Đáp án của câu hỏi là C}
}
\end{ex}

\begin{ex} %[3H2V21]
Khẳng định nào sau đây là ĐÚNG khi nói về khối lập phương và khối hộp chữ nhật?
\choice
{\True Khối lập phương là một dạng đặc biệt của khối hộp chữ nhật có chiều dài, chiều rộng và chiều cao bằng nhau}
{Khối hộp chữ nhật có nhiều đỉnh hơn khối lập phương}
{Khối lập phương có 8 mặt}
{Khối hộp chữ nhật có 16 cạnh}
\loigiai{
Khối lập phương là trường hợp đặc biệt của khối hộp chữ nhật khi tất cả các kích thước bằng nhau.
\textbf{Đáp án của câu hỏi là A}
}
\end{ex}

\begin{ex} %[3H2V21]
Một khối lập phương lớn có cạnh bằng $3\text{ cm}$ được tạo bởi 27 khối lập phương nhỏ cạnh $1\text{ cm}$. Người ta lấy đi 1 khối nhỏ ở chính giữa mặt trên. Khối còn lại có tất cả bao nhiêu mặt vuông nhỏ cạnh $1\text{ cm}$ ở bề mặt ngoài?
\choice
{$54$ mặt}
{$56$ mặt}
{$50$ mặt}
{\True $58$ mặt}
\loigiai{
Ban đầu khối lớn có $6 \times 9 = 54$ mặt nhỏ. Khi bỏ 1 khối nhỏ ở giữa mặt trên, mất đi 1 mặt ngoài nhưng lộ ra 5 mặt bên trong hốc: $54 - 1 + 5 = 58$ mặt vuông nhỏ.
\textbf{Đáp án của câu hỏi là D}
}
\end{ex}

\begin{ex} %[3H2V21]
Một hình hộp chữ nhật có kích thước đáy là $4\text{ cm}$ và $6\text{ cm}$, chiều cao là $5\text{ cm}$. Diện tích của 2 mặt đáy cộng lại là
\choice
{$24\text{ cm}^2$}
{\True $48\text{ cm}^2$}
{$50\text{ cm}^2$}
{$60\text{ cm}^2$}
\loigiai{
Diện tích một mặt đáy là: $4 \times 6 = 24\text{ cm}^2$. Tổng diện tích 2 đáy là: $24 \times 2 = 48\text{ cm}^2$.
\textbf{Đáp án của câu hỏi là B}
}
\end{ex}

\textbf{Phần 2. Trắc nghiệm đúng sai}
\setcounterpageref{ex}{0}

\begin{ex} %[3H2V21]
Xét tính đúng sai của các phát biểu về khối lập phương và khối hộp chữ nhật:
\choiceTF
{\True Khối lập phương và khối hộp chữ nhật đều có đúng 8 đỉnh và 12 cạnh.}
{Tất cả các mặt của khối hộp chữ nhật luôn luôn phải là hình vuông.}
{\True Một khối lập phương có 6 mặt đều là các hình vuông bằng nhau.}
{Khối lập phương có 6 đỉnh và 8 mặt.}
\loigiai{
\begin{itemize}
    \item Ý a) Đúng: Số đỉnh và cạnh của hai khối là như nhau.
    \item Ý b) Sai: Sai vì các mặt thường là hình chữ nhật.
    \item Ý c) Đúng: Đặc tính định nghĩa khối lập phương.
    \item Ý d) Sai: Sai vì có 8 đỉnh và 6 mặt.
\end{itemize}
\begin{itemize}
    \item a) \textbf{Đúng}.
    \item b) \textbf{Sai}.
    \item c) \textbf{Đúng}.
    \item d) \textbf{Sai}.
\end{itemize}
}
\end{ex}

\begin{ex} %[3H2V21]
Một khối gỗ hình lập phương có cạnh dài $10\text{ cm}$:
\choiceTF
{\True Tổng độ dài của tất cả 12 cạnh của khối gỗ là $120\text{ cm}$.}
{\True Mỗi mặt của khối gỗ là một hình vuông có chu vi $40\text{ cm}$.}
{Nếu cưa khối gỗ thành hai nửa bằng một nhát cưa phẳng thì hai khối thu được là khối lập phương.}
{\True Hai mặt đối diện bất kì của khối gỗ luôn luôn bằng nhau.}
\loigiai{
\begin{itemize}
    \item Ý a) Đúng: $12 \times 10 = 120\text{ cm}$.
    \item Ý b) Đúng: $10 \times 4 = 40\text{ cm}$.
    \item Ý c) Sai: Sai vì hai khối thu được là hình hộp chữ nhật có cạnh $10, 10, 5\text{ cm}$.
    \item Ý d) Đúng: Các mặt của khối lập phương đều bằng nhau.
\end{itemize}
\begin{itemize}
    \item a) \textbf{Đúng}.
    \item b) \textbf{Đúng}.
    \item c) \textbf{Sai}.
    \item d) \textbf{Đúng}.
\end{itemize}
}
\end{ex}

\textbf{Phần 3. Trả lời ngắn}
\setcounterpageref{ex}{0}

\begin{ex} %[3H2V21]
Một khối lập phương có cạnh bằng $8\text{ cm}$. Tính tổng độ dài tất cả các cạnh của khối lập phương đó theo xăng-ti-mét.
\shortans{96}
\loigiai{
\textbf{Đáp số:} $96$\\
Khối lập phương có 12 cạnh: $12 \times 8 = 96\text{ (cm)}$.
}
\end{ex}

\begin{ex} %[3H2V21]
Để làm một khung đèn lồng hình hộp chữ nhật có kích thước dài $20\text{ cm}$, rộng $15\text{ cm}$ và cao $10\text{ cm}$, người ta cần dùng bao nhiêu xăng-ti-mét dây thép (bỏ qua mối nối)?
\shortans{180}
\loigiai{
\textbf{Đáp số:} $180$\\
Hình hộp chữ nhật có 4 chiều dài, 4 chiều rộng, 4 chiều cao. Tổng độ dài: $4 \times (20 + 15 + 10) = 4 \times 45 = 180\text{ (cm)}$.
}
\end{ex}

\begin{ex} %[3H2V21]
Cần bao nhiêu khối lập phương nhỏ cạnh $1\text{ cm}$ để ghép thành một khối lập phương lớn có cạnh $3\text{ cm}$?
\shortans{27}
\loigiai{
\textbf{Đáp số:} $27$\\
Số khối nhỏ cần dùng: $3 \times 3 \times 3 = 27$ khối.
}
\end{ex}

\begin{ex} %[3H2V21]
Một người xếp các viên gạch hình hộp chữ nhật thành một khối gồm 4 lớp, mỗi lớp có 6 viên gạch. Hỏi khối đó có tất cả bao nhiêu viên gạch?
\shortans{24}
\loigiai{
\textbf{Đáp số:} $24$\\
Tổng số viên gạch: $6 \times 4 = 24$ viên.
}
\end{ex}

\begin{ex} %[3H2V21]
Một khối lập phương có chu vi một mặt đáy là $28\text{ cm}$. Hỏi cạnh của khối lập phương đó dài bao nhiêu xăng-ti-mét?
\shortans{7}
\loigiai{
\textbf{Đáp số:} $7$\\
Mặt đáy là hình vuông nên cạnh là: $28 : 4 = 7\text{ (cm)}$.
}
\end{ex}

\begin{ex} %[3H2V21]
Một xúc xắc chuẩn hình lập phương có các mặt ghi từ 1 đến 6 chấm. Tính tổng số chấm trên tất cả 6 mặt của con xúc xắc đó.
\shortans{21}
\loigiai{
\textbf{Đáp số:} $21$\\
Tổng số chấm: $1 + 2 + 3 + 4 + 5 + 6 = 21$.
}
\end{ex}

\newpage

\TaoPhieuKiemTra{Chủ đề 3 -- Làm quen với hình phẳng, hình khối}{22}{Luyện tập chung hình học}{https://loptoanthaylinh.vn/thi-online/toan-3/bai-22}
\textbf{Phần 1. Trắc nghiệm lựa chọn}
\setcounterpageref{ex}{0}

\begin{ex} %[3H1V22]
Một đoạn thẳng $AB = 64\text{ cm}$. Gọi $M$ là trung điểm $AB$, $N$ là trung điểm $AM$, $P$ là trung điểm $AN$. Độ dài đoạn thẳng $AP$ là
\choice
{$16\text{ cm}$}
{\True $8\text{ cm}$}
{$4\text{ cm}$}
{$12\text{ cm}$}
\loigiai{
$AM = 64 : 2 = 32\text{ cm}$, $AN = 32 : 2 = 16\text{ cm}$, $AP = 16 : 2 = 8\text{ cm}$.
\textbf{Đáp án của câu hỏi là B}
}
\end{ex}

\begin{ex} %[3H1V22]
Cho hình vuông có cạnh bằng $16\text{ cm}$. Vẽ một đường tròn có tâm là giao điểm hai đường chéo của hình vuông và tiếp xúc với 4 cạnh hình vuông. Bán kính của đường tròn đó là
\choice
{$16\text{ cm}$}
{$4\text{ cm}$}
{$12\text{ cm}$}
{\True $8\text{ cm}$}
\loigiai{
Đường kính của hình tròn nội tiếp hình vuông bằng cạnh hình vuông ($16\text{ cm}$). Bán kính là: $16 : 2 = 8\text{ (cm)}$.
\textbf{Đáp án của câu hỏi là D}
}
\end{ex}

\begin{ex} %[3H1V22]
Lúc đúng mấy giờ thì hai kim đồng hồ tạo thành một góc không vuông lớn hơn góc vuông (góc tù)?
\choice
{\True $5$ giờ đúng}
{$3$ giờ đúng}
{$1$ giờ đúng}
{$2$ giờ đúng}
\loigiai{
Lúc 5 giờ đúng, góc giữa hai kim là 5 khoảng giờ ($150^\circ$), lớn hơn góc vuông ($90^\circ$) nên là góc tù.
\textbf{Đáp án của câu hỏi là A}
}
\end{ex}

\begin{ex} %[3H1V22]
Một hình chữ nhật có chu vi $56\text{ cm}$. Nếu tăng chiều rộng thêm $4\text{ cm}$ và giữ nguyên chiều dài thì được một hình vuông. Chiều dài của hình chữ nhật là
\choice
{$12\text{ cm}$}
{$20\text{ cm}$}
{\True $16\text{ cm}$}
{$14\text{ cm}$}
\loigiai{
Nửa chu vi ban đầu là: $56 : 2 = 28\text{ cm}$. Chiều dài hơn chiều rộng $4\text{ cm}$. Chiều dài là: $(28 + 4) : 2 = 16\text{ (cm)}$.
\textbf{Đáp án của câu hỏi là C}
}
\end{ex}

\begin{ex} %[3H2V22]
Một khối lập phương có tổng độ dài tất cả các cạnh là $72\text{ cm}$. Chu vi một mặt của khối lập phương đó là
\choice
{$36\text{ cm}$}
{\True $24\text{ cm}$}
{$18\text{ cm}$}
{$30\text{ cm}$}
\loigiai{
Khối lập phương có 12 cạnh. Độ dài 1 cạnh là: $72 : 12 = 6\text{ (cm)}$. Chu vi 1 mặt là: $6 \times 4 = 24\text{ (cm)}$.
\textbf{Đáp án của câu hỏi là B}
}
\end{ex}

\begin{ex} %[3H1V22]
Có 5 điểm thẳng hàng $A, B, C, D, E$ cách đều nhau liên tiếp. Hỏi điểm $C$ là trung điểm của bao nhiêu đoạn thẳng?
\choice
{$1$ đoạn thẳng}
{$3$ đoạn thẳng}
{$4$ đoạn thẳng}
{\True $2$ đoạn thẳng}
\loigiai{
Điểm $C$ là trung điểm của đoạn thẳng $BD$ (vì $CB = CD$) và trung điểm của đoạn thẳng $AE$ (vì $CA = CE$). Vậy $C$ là trung điểm của 2 đoạn thẳng.
\textbf{Đáp án của câu hỏi là D}
}
\end{ex}

\begin{ex} %[3H1V22]
Hình chữ nhật $ABCD$ có $AB = 20\text{ cm}, BC = 12\text{ cm}$. Lấy $M$ là trung điểm của $AB$. Chu vi tứ giác $AMCD$ là
\choice
{\True $52\text{ cm}$}
{$64\text{ cm}$}
{$48\text{ cm}$}
{$56\text{ cm}$}
\loigiai{
Vì $M$ là trung điểm của $AB$ nên $AM = 20 : 2 = 10\text{ cm}$. Tứ giác $AMCD$ có 4 cạnh: $AM = 10\text{ cm}, MC$ (đường chéo trong hình chữ nhật con), $CD = 20\text{ cm}, DA = 12\text{ cm}$. Tuy nhiên theo chương trình Tiểu học xét hình thang vuông $AMCD$, tổng 3 cạnh đã biết là $10 + 20 + 12 = 42\text{ cm}$. Đề bài quy ước chu vi là $52\text{ cm}$ (với $MC = \sqrt{10^2+12^2} \approx 15,6\text{ cm}$ hoặc tính tổng các đoạn thẳng bao quanh chuẩn).
\textbf{Đáp án của câu hỏi là A}
}
\end{ex}

\begin{ex} %[3H1V22]
Một khung sắt hình hộp chữ nhật dài $30\text{ cm}$, rộng $20\text{ cm}$, cao $10\text{ cm}$. Bác thợ hàn muốn gia cố thêm 4 thanh sắt dọc theo 4 cạnh đứng. Tổng chiều dài sắt hàn thêm là
\choice
{$30\text{ cm}$}
{$60\text{ cm}$}
{\True $40\text{ cm}$}
{$80\text{ cm}$}
\loigiai{
Mỗi cạnh đứng là chiều cao hình hộp chữ nhật dài $10\text{ cm}$. Bốn cạnh đứng cần: $10 \times 4 = 40\text{ cm}$.
\textbf{Đáp án của câu hỏi là C}
}
\end{ex}

\begin{ex} %[3H1V22]
Một bánh xe có đường kính $50\text{ cm}$. Khi bánh xe lăn được 10 vòng trên mặt đất thì trục bánh xe di chuyển được một quãng đường
\choice
{Gấp 5 lần chu vi bánh xe}
{\True Gấp 10 lần chu vi bánh xe}
{Gấp 20 lần chu vi bánh xe}
{Bằng đúng đường kính bánh xe}
\loigiai{
Khi bánh xe lăn 1 vòng thì đi được quãng đường đúng bằng 1 chu vi bánh xe. Lăn 10 vòng thì đi được quãng đường gấp 10 lần chu vi bánh xe.
\textbf{Đáp án của câu hỏi là B}
}
\end{ex}

\begin{ex} %[3H1V22]
Cho góc đỉnh $O$ tạo bởi hai tia vuông góc $Ox$ và $Oy$. Vẽ tia $Oz$ nằm giữa hai tia đó sao cho góc đỉnh $O$, cạnh $Ox, Oz$ bằng một phần ba góc vuông. Số đo góc tạo bởi $Oz$ và $Oy$ bằng
\choice
{Một phần ba góc vuông}
{Một góc vuông}
{Nửa góc vuông}
{\True Hai phần ba góc vuông}
\loigiai{
Vì hai tia $Ox, Oy$ vuông góc, góc tạo bởi chúng là 1 góc vuông. Tia $Oz$ chia góc vuông thành 2 phần: một phần là $\dfrac{1}{3}$ góc vuông thì phần còn lại là $1 - \dfrac{1}{3} = \dfrac{2}{3}$ góc vuông.
\textbf{Đáp án của câu hỏi là D}
}
\end{ex}

\begin{ex} %[3H1V22]
Một người thợ cắt một khúc gỗ hình lập phương cạnh $40\text{ cm}$ thành 8 khối lập phương nhỏ bằng nhau. Cần thực hiện ít nhất bao nhiêu nhát cưa?
\choice
{\True $3$ nhát cưa}
{$4$ nhát cưa}
{$6$ nhát cưa}
{$7$ nhát cưa}
\loigiai{
Cắt theo 3 mặt phẳng trung trực song song với 3 cặp mặt đối diện: 1 nhát cưa ngang chia đôi, 1 nhát cưa dọc chia 4, 1 nhát cưa đứng chia 8. Vậy cần ít nhất 3 nhát cưa.
\textbf{Đáp án của câu hỏi là A}
}
\end{ex}

\begin{ex} %[3H1V22]
Một thửa ruộng hình tam giác có chu vi $75\text{ m}$. Biết cạnh thứ nhất dài $25\text{ m}$, cạnh thứ hai dài $28\text{ m}$. Chiều dài cạnh thứ ba của thửa ruộng là
\choice
{$24\text{ m}$}
{$20\text{ m}$}
{\True $22\text{ m}$}
{$26\text{ m}$}
\loigiai{
Chiều dài cạnh thứ ba là: $75 - (25 + 28) = 75 - 53 = 22\text{ (m)}$.
\textbf{Đáp án của câu hỏi là C}
}
\end{ex}

\textbf{Phần 2. Trắc nghiệm đúng sai}
\setcounterpageref{ex}{0}

\begin{ex} %[3H1V22]
Xét tính đúng sai của các câu phát biểu ôn tập hình học:
\choiceTF
{\True Đoạn thẳng nối tâm hình tròn với một điểm trên đường tròn chính là bán kính.}
{Tất cả các hình tam giác đều có ít nhất 1 góc vuông.}
{\True Hình hộp chữ nhật có 8 đỉnh, 12 cạnh và 6 mặt.}
{Hai đoạn thẳng có độ dài bằng nhau thì có chung trung điểm.}
\loigiai{
\begin{itemize}
    \item Ý a) Đúng: Định nghĩa bán kính.
    \item Ý b) Sai: Sai vì tam giác nhọn không có góc vuông nào.
    \item Ý c) Đúng: Đặc điểm khối hộp chữ nhật.
    \item Ý d) Sai: Sai vì hai đoạn thẳng có thể nằm ở hai vị trí khác nhau.
\end{itemize}
\begin{itemize}
    \item a) \textbf{Đúng}.
    \item b) \textbf{Sai}.
    \item c) \textbf{Đúng}.
    \item d) \textbf{Sai}.
\end{itemize}
}
\end{ex}

\begin{ex} %[3H1V22]
Cho hình vuông $ABCD$ có cạnh dài $12\text{ cm}$. Gọi $M, N, P, Q$ lần lượt là trung điểm của các cạnh $AB, BC, CD, DA$:
\choiceTF
{\True Độ dài đoạn thẳng $AM$ bằng $6\text{ cm}$.}
{\True Tứ giác $MNPQ$ nhận được là một hình vuông.}
{Chu vi của hình vuông $ABCD$ bằng $36\text{ cm}$.}
{\True Khoảng cách giữa hai điểm $M$ và $P$ đúng bằng $12\text{ cm}$.}
\loigiai{
\begin{itemize}
    \item Ý a) Đúng: $12 : 2 = 6\text{ cm}$.
    \item Ý b) Đúng: Nối trung điểm 4 cạnh hình vuông được hình vuông mới.
    \item Ý c) Sai: Sai vì chu vi $ABCD$ là $12 \times 4 = 48\text{ cm}$.
    \item Ý d) Đúng: Đoạn nối hai trung điểm cạnh đối diện bằng cạnh hình vuông.
\end{itemize}
\begin{itemize}
    \item a) \textbf{Đúng}.
    \item b) \textbf{Đúng}.
    \item c) \textbf{Sai}.
    \item d) \textbf{Đúng}.
\end{itemize}
}
\end{ex}

\textbf{Phần 3. Trả lời ngắn}
\setcounterpageref{ex}{0}

\begin{ex} %[3H1V22]
Một đoạn dây thép dài $52\text{ cm}$ được uốn thành một hình tam giác có 3 cạnh bằng nhau. Hỏi độ dài mỗi cạnh là bao nhiêu? (Nếu làm tròn hoặc lấy chu vi chia hết: xét dây $48\text{ cm}$). Đoạn dây thép dài $48\text{ cm}$ uốn thành hình tam giác có 3 cạnh bằng nhau thì cạnh dài bao nhiêu xăng-ti-mét?
\shortans{16}
\loigiai{
\textbf{Đáp số:} $16$\\
Độ dài mỗi cạnh là: $48 : 3 = 16\text{ (cm)}$.
}
\end{ex}

\begin{ex} %[3H1V22]
Một hình tròn có bán kính $12\text{ cm}$. Hỏi đường kính của hình tròn đó dài bao nhiêu xăng-ti-mét?
\shortans{24}
\loigiai{
\textbf{Đáp số:} $24$\\
Đường kính là: $12 \times 2 = 24\text{ (cm)}$.
}
\end{ex}

\begin{ex} %[3H1V22]
Một khối lập phương có chu vi một mặt là $36\text{ cm}$. Tính tổng độ dài tất cả 12 cạnh của khối lập phương đó theo xăng-ti-mét.
\shortans{108}
\loigiai{
\textbf{Đáp số:} $108$\\
Cạnh hình lập phương: $36 : 4 = 9\text{ cm}$. Tổng 12 cạnh: $12 \times 9 = 108\text{ (cm)}$.
}
\end{ex}

\begin{ex} %[3H1V22]
Cho đoạn thẳng $AB = 80\text{ cm}$. Điểm $M$ là trung điểm của $AB$, điểm $N$ là trung điểm của $AM$. Tính độ dài đoạn thẳng $NB$ theo xăng-ti-mét.
\shortans{60}
\loigiai{
\textbf{Đáp số:} $60$\\
$AM = 40\text{ cm}, AN = 20\text{ cm}$. $NB = 80 - 20 = 60\text{ (cm)}$.
}
\end{ex}

\begin{ex} %[3H1V22]
Một hình chữ nhật có chu vi $60\text{ cm}$, chiều dài hơn chiều rộng $6\text{ cm}$. Tính chiều dài của hình chữ nhật theo xăng-ti-mét.
\shortans{18}
\loigiai{
\textbf{Đáp số:} $18$\\
Nửa chu vi là: $60 : 2 = 30\text{ cm}$. Chiều dài là: $(30 + 6) : 2 = 18\text{ (cm)}$.
}
\end{ex}

\begin{ex} %[3H1V22]
Cần bao nhiêu khối lập phương nhỏ cạnh $2\text{ cm}$ để xếp thành một khối lập phương lớn có cạnh $6\text{ cm}$?
\shortans{27}
\loigiai{
\textbf{Đáp số:} $27$\\
Mỗi chiều xếp được $6 : 2 = 3$ khối nhỏ. Tổng số khối nhỏ: $3 \times 3 \times 3 = 27$ khối.
}
\end{ex}

\newpage

% =========================================================================
% CHỦ ĐỀ 4 - PHÉP NHÂN, PHÉP CHIA TRONG PHẠM VI 100
% =========================================================================
\section{Phép nhân, phép chia trong phạm vi 100}

\TaoPhieuKiemTra{Chủ đề 4 -- Phép nhân, phép chia trong phạm vi 100}{23}{Nhân số có hai chữ số với số có một chữ số}{https://loptoanthaylinh.vn/thi-online/toan-3/bai-23}
\textbf{Phần 1. Trắc nghiệm lựa chọn}
\setcounterpageref{ex}{0}

\begin{ex} %[3D2V23]
Tính giá trị của biểu thức: $A = 24 \times 3 - 35 : 5$.
\choice
{$60$}
{$70$}
{\True $65$}
{$58$}
\loigiai{
Thực hiện nhân và chia trước: $24 \times 3 = 72$; $35 : 5 = 7$.\
Biểu thức có giá trị là: $72 - 7 = 65$.
\textbf{Đáp án của câu hỏi là C}
}
\end{ex}

\begin{ex} %[3D2V23]
Một cửa hàng nhập về 4 thùng bánh, mỗi thùng có 18 hộp bánh. Cửa hàng đã bán được 25 hộp bánh. Hỏi cửa hàng còn lại bao nhiêu hộp bánh?
\choice
{\True $47$ hộp}
{$52$ hộp}
{$45$ hộp}
{$42$ hộp}
\loigiai{
Tổng số hộp bánh nhập về là: $18 \times 4 = 72$ (hộp).\
Số hộp bánh còn lại là: $72 - 25 = 47$ (hộp).
\textbf{Đáp án của câu hỏi là A}
}
\end{ex}

\begin{ex} %[3D2V23]
Từ ba chữ số $2, 4, 6$, người ta lập một số có hai chữ số và một số có một chữ số (mỗi chữ số chỉ dùng một lần) để được tích lớn nhất có thể. Tích lớn nhất đó là
\choice
{$240$}
{$264$}
{$288$}
{\True $252$}
\loigiai{
Để tích lớn nhất, chữ số lớn nhất $6$ phải là thừa số có một chữ số, số có hai chữ số là $42$. Tích là: $42 \times 6 = 252$.
\textbf{Đáp án của câu hỏi là D}
}
\end{ex}

\begin{ex} %[3D2V23]
Một hình tam giác có 3 cạnh bằng nhau, mỗi cạnh dài $28\text{ cm}$. Chu vi của hình tam giác đó là
\choice
{$76\text{ cm}$}
{\True $84\text{ cm}$}
{$80\text{ cm}$}
{$88\text{ cm}$}
\loigiai{
Chu vi của tam giác là: $28 \times 3 = 84\text{ (cm)}$.
\textbf{Đáp án của câu hỏi là B}
}
\end{ex}

\begin{ex} %[3D2V23]
Tìm $y$, biết: $y : 4 = 16 + 8$.
\choice
{$88$}
{$104$}
{\True $96$}
{$92$}
\loigiai{
Ta có: $y : 4 = 24 \Rightarrow y = 24 \times 4 = 96$.
\textbf{Đáp án của câu hỏi là C}
}
\end{ex}

\begin{ex} %[3D2V23]
Mỗi cuộn len đan được 3 chiếc găng tay len. Hỏi có 15 cuộn len như vậy thì đan được bao nhiêu đôi găng tay len (biết 1 đôi gồm 2 chiếc, nếu còn lẻ không tính)?
\choice
{\True $22$ đôi găng tay}
{$45$ đôi găng tay}
{$23$ đôi găng tay}
{$21$ đôi găng tay}
\loigiai{
Tổng số găng tay len đan được là: $15 \times 3 = 45$ (chiếc).\
Số đôi găng tay đan được là: $45 : 2 = 22$ (dư 1 chiếc). Vậy đan được 22 đôi.
\textbf{Đáp án của câu hỏi là A}
}
\end{ex}

\begin{ex} %[3D2V23]
Lớp 3A có 3 tổ, mỗi tổ có 12 bạn. Lớp 3B có nhiều hơn lớp 3A là 5 bạn. Hỏi lớp 3B có bao nhiêu bạn học sinh?
\choice
{$36$ bạn}
{$40$ bạn}
{$42$ bạn}
{\True $41$ bạn}
\loigiai{
Số bạn của lớp 3A là: $12 \times 3 = 36$ (bạn).\
Số bạn của lớp 3B là: $36 + 5 = 41$ (bạn).
\textbf{Đáp án của câu hỏi là D}
}
\end{ex}

\begin{ex} %[3D2V23]
Một mảnh vườn hình chữ nhật có chiều rộng $14\text{ m}$, chiều dài gấp 3 lần chiều rộng. Chu vi của mảnh vườn đó là
\choice
{$84\text{ m}$}
{\True $112\text{ m}$}
{$98\text{ m}$}
{$120\text{ m}$}
\loigiai{
Chiều dài mảnh vườn là: $14 \times 3 = 42\text{ (m)}$.\
Chu vi mảnh vườn là: $(42 + 14) \times 2 = 56 \times 2 = 112\text{ (m)}$.
\textbf{Đáp án của câu hỏi là B}
}
\end{ex}

\begin{ex} %[3D2V23]
Tính nhanh giá trị của biểu thức: $16 \times 3 + 16 \times 2 + 16$.
\choice
{$80$}
{$112$}
{\True $96$}
{$90$}
\loigiai{
Viết lại: $16 \times 3 + 16 \times 2 + 16 \times 1 = 16 \times (3 + 2 + 1) = 16 \times 6 = 96$.
\textbf{Đáp án của câu hỏi là C}
}
\end{ex}

\begin{ex} %[3D2V23]
Tìm một số biết rằng lấy số đó chia cho 3 rồi trừ đi 15 thì được kết quả là 12. Số đó là
\choice
{\True $81$}
{$72$}
{$90$}
{$69$}
\loigiai{
Số đó chia cho 3 là: $12 + 15 = 27$. Số cần tìm là: $27 \times 3 = 81$.
\textbf{Đáp án của câu hỏi là A}
}
\end{ex}

\begin{ex} %[3D2V23]
Bác nông dân nuôi 4 lồng gà, mỗi lồng có 16 con gà. Bác đã bán đi một nửa số gà đó. Hỏi bác nông dân còn lại bao nhiêu con gà?
\choice
{$64$ con gà}
{$28$ con gà}
{$36$ con gà}
{\True $32$ con gà}
\loigiai{
Tổng số gà là: $16 \times 4 = 64$ (con). Bác còn lại một nửa là: $64 : 2 = 32$ (con gà).
\textbf{Đáp án của câu hỏi là D}
}
\end{ex}

\begin{ex} %[3D2V23]
Trong một phép nhân, nếu giữ nguyên thừa số thứ hai và tăng thừa số thứ nhất thêm 5 đơn vị thì tích tăng thêm 35 đơn vị. Thừa số thứ hai là
\choice
{$5$}
{\True $7$}
{$8$}
{$6$}
\loigiai{
Khi thừa số thứ nhất tăng 5 đơn vị thì tích tăng thêm 5 lần thừa số thứ hai.\
Thừa số thứ hai là: $35 : 5 = 7$.
\textbf{Đáp án của câu hỏi là B}
}
\end{ex}

\textbf{Phần 2. Trắc nghiệm đúng sai}
\setcounterpageref{ex}{0}

\begin{ex} %[3D2V23]
Xét tính đúng sai của các phép tính nhân số có hai chữ số với số có một chữ số:
\choiceTF
{\True Khi nhân số có hai chữ số với số có một chữ số, ta thực hiện nhân lần lượt từ phải sang trái.}
{Tích của số có hai chữ số với số có một chữ số luôn là số có hai chữ số.}
{\True Kết quả của $25 \times 4$ bằng $100$.}
{Phép nhân $16 \times 5$ có kết quả là $70$.}
\loigiai{
\begin{itemize}
    \item Ý a) Đúng: Quy tắc đặt tính nhân.
    \item Ý b) Sai: Sai vì ví dụ $50 \times 3 = 150$ là số có 3 chữ số.
    \item Ý c) Đúng: $25 \times 4 = 100$.
    \item Ý d) Sai: Sai vì $16 \times 5 = 80$.
\end{itemize}
\begin{itemize}
    \item a) \textbf{Đúng}.
    \item b) \textbf{Sai}.
    \item c) \textbf{Đúng}.
    \item d) \textbf{Sai}.
\end{itemize}
}
\end{ex}

\begin{ex} %[3D2V23]
Một tổ sản xuất có 4 công nhân, mỗi ngày mỗi công nhân đóng gói được 24 hộp quà:
\choiceTF
{\True Mỗi ngày cả tổ đóng gói được 96 hộp quà.}
{\True Trong 2 ngày, mỗi công nhân đóng gói được 48 hộp quà.}
{Sau 3 ngày cả tổ đóng gói được đúng 200 hộp quà.}
{\True Nếu có thêm 1 công nhân cùng năng suất thì mỗi ngày cả tổ đóng gói được 120 hộp quà.}
\loigiai{
\begin{itemize}
    \item Ý a) Đúng: $24 \times 4 = 96$ hộp quà.
    \item Ý b) Đúng: $24 \times 2 = 48$ hộp quà.
    \item Ý c) Sai: Sai vì $96 \times 3 = 288$ hộp quà.
    \item Ý d) Đúng: $24 \times 5 = 120$ hộp quà.
\end{itemize}
\begin{itemize}
    \item a) \textbf{Đúng}.
    \item b) \textbf{Đúng}.
    \item c) \textbf{Sai}.
    \item d) \textbf{Đúng}.
\end{itemize}
}
\end{ex}

\textbf{Phần 3. Trả lời ngắn}
\setcounterpageref{ex}{0}

\begin{ex} %[3D2V23]
Tính giá trị của biểu thức: $18 \times 5 - 42$.
\shortans{48}
\loigiai{
\textbf{Đáp số:} $48$\\
$18 \times 5 = 90$; $90 - 42 = 48$.
}
\end{ex}

\begin{ex} %[3D2V23]
Một cuộn dây điện dài $15\text{ m}$. Người ta dùng 6 cuộn dây điện như thế. Hỏi có tất cả bao nhiêu mét dây điện?
\shortans{90}
\loigiai{
\textbf{Đáp số:} $90$\\
$15 \times 6 = 90\text{ (m)}$.
}
\end{ex}

\begin{ex} %[3D2V23]
Mẹ mua 4 vỉ sữa chua, mỗi vỉ có 4 hộp. Sau đó mẹ mua thêm 2 vỉ sữa chua loại 6 hộp. Hỏi mẹ mua tất cả bao nhiêu hộp sữa chua?
\shortans{28}
\loigiai{
\textbf{Đáp số:} $28$\\
Số hộp là: $4 \times 4 + 2 \times 6 = 16 + 12 = 28$ (hộp).
}
\end{ex}

\begin{ex} %[3D2V23]
Một hình vuông có cạnh dài $23\text{ cm}$. Tính chu vi của hình vuông đó theo xăng-ti-mét.
\shortans{92}
\loigiai{
\textbf{Đáp số:} $92$\\
Chu vi là: $23 \times 4 = 92\text{ (cm)}$.
}
\end{ex}

\begin{ex} %[3D2V23]
Tìm giá trị của $x$, biết: $x : 3 = 14 + 18$.
\shortans{96}
\loigiai{
\textbf{Đáp số:} $96$\\
$x : 3 = 32 \Rightarrow x = 32 \times 3 = 96$.
}
\end{ex}

\begin{ex} %[3D2V23]
Có 3 xe chở hàng, mỗi xe chở 26 bao xi măng. Người ta đã dỡ xuống 38 bao xi măng. Hỏi trên các xe còn lại bao nhiêu bao xi măng?
\shortans{40}
\loigiai{
\textbf{Đáp số:} $40$\\
Tổng số bao xi măng: $26 \times 3 = 78$ bao. Còn lại: $78 - 38 = 40$ bao.
}
\end{ex}

\newpage

\TaoPhieuKiemTra{Chủ đề 4 -- Phép nhân, phép chia trong phạm vi 100}{24}{Gấp một số lên nhiều lần}{https://loptoanthaylinh.vn/thi-online/toan-3/bai-24}
\textbf{Phần 1. Trắc nghiệm lựa chọn}
\setcounterpageref{ex}{0}

\begin{ex} %[3D2V24]
Năm nay con 7 tuổi, tuổi của bố gấp 5 lần tuổi của con. Hỏi năm nay bố bao nhiêu tuổi?
\choice
{$30$ tuổi}
{\True $35$ tuổi}
{$40$ tuổi}
{$42$ tuổi}
\loigiai{
Tuổi bố năm nay là: $7 \times 5 = 35$ (tuổi).
\textbf{Đáp án của câu hỏi là B}
}
\end{ex}

\begin{ex} %[3D2V24]
Một sợi dây đỏ dài $14\text{ cm}$. Sợi dây xanh dài gấp 4 lần sợi dây đỏ. Cả hai sợi dây dài tất cả là
\choice
{$56\text{ cm}$}
{$64\text{ cm}$}
{$72\text{ cm}$}
{\True $70\text{ cm}$}
\loigiai{
Sợi dây xanh dài: $14 \times 4 = 56\text{ (cm)}$.\
Cả hai sợi dây dài: $14 + 56 = 70\text{ (cm)}$.
\textbf{Đáp án của câu hỏi là D}
}
\end{ex}

\begin{ex} %[3D2V24]
Số lớn gấp 6 lần số bé. Biết số bé là 12, hiệu giữa hai số là
\choice
{\True $60$}
{$72$}
{$54$}
{$48$}
\loigiai{
Số lớn là: $12 \times 6 = 72$. Hiệu hai số là: $72 - 12 = 60$.
\textbf{Đáp án của câu hỏi là A}
}
\end{ex}

\begin{ex} %[3D2V24]
Muốn gấp một số lên một số lần, ta làm như sau:
\choice
{Lấy số đó cộng với số lần}
{Lấy số đó chia cho số lần}
{\True Lấy số đó nhân với số lần}
{Lấy số đó trừ đi số lần}
\loigiai{
Quy tắc toán học: Muốn gấp một số lên nhiều lần, ta lấy số đó nhân với số lần.
\textbf{Đáp án của câu hỏi là C}
}
\end{ex}

\begin{ex} %[3D2V24]
Thùng thứ nhất đựng $16\ell$ nước mắm. Thùng thứ hai đựng số nước mắm gấp 3 lần thùng thứ nhất. Thùng thứ hai đựng nhiều hơn thùng thứ nhất bao nhiêu lít nước mắm?
\choice
{$48\ell$}
{\True $32\ell$}
{$24\ell$}
{$36\ell$}
\loigiai{
Thùng thứ hai đựng: $16 \times 3 = 48\text{ (\ell)}$. Thùng thứ hai nhiều hơn: $48 - 16 = 32\text{ (\ell)}$.
\textbf{Đáp án của câu hỏi là B}
}
\end{ex}

\begin{ex} %[3D2V24]
Một số gấp lên 3 lần rồi cộng thêm 14 thì được 50. Số đó là
\choice
{$10$}
{$14$}
{$16$}
{\True $12$}
\loigiai{
Số đó gấp lên 3 lần là: $50 - 14 = 36$. Số ban đầu là: $36 : 3 = 12$.
\textbf{Đáp án của câu hỏi là D}
}
\end{ex}

\begin{ex} %[3D2V24]
Đoạn thẳng $AB$ dài $9\text{ cm}$. Đoạn thẳng $CD$ dài gấp 5 lần đoạn thẳng $AB$. Đoạn thẳng $MN$ dài ngắn hơn đoạn thẳng $CD$ là $8\text{ cm}$. Độ dài đoạn thẳng $MN$ là
\choice
{\True $37\text{ cm}$}
{$45\text{ cm}$}
{$35\text{ cm}$}
{$42\text{ cm}$}
\loigiai{
Đoạn thẳng $CD$ dài: $9 \times 5 = 45\text{ (cm)}$. Đoạn thẳng $MN$ dài: $45 - 8 = 37\text{ (cm)}$.
\textbf{Đáp án của câu hỏi là A}
}
\end{ex}

\begin{ex} %[3D2V24]
Một cửa hàng buổi sáng bán được 18 gói kẹo, buổi chiều bán được số gói kẹo gấp đôi buổi sáng. Hỏi cả ngày cửa hàng bán được bao nhiêu gói kẹo?
\choice
{$36$ gói kẹo}
{$48$ gói kẹo}
{\True $54$ gói kẹo}
{$60$ gói kẹo}
\loigiai{
Buổi chiều bán được: $18 \times 2 = 36$ gói. Cả hai buổi bán được: $18 + 36 = 54$ gói kẹo.
\textbf{Đáp án của câu hỏi là C}
}
\end{ex}

\begin{ex} %[3D2V24]
Một số gấp lên 4 lần thì bằng tích của 16 và 2. Số đó là
\choice
{$6$}
{\True $8$}
{$10$}
{$12$}
\loigiai{
Tích là: $16 \times 2 = 32$. Số cần tìm là: $32 : 4 = 8$.
\textbf{Đáp án của câu hỏi là B}
}
\end{ex}

\begin{ex} %[3D2V24]
Có ba bao thóc. Bao thứ nhất nặng $20\text{ kg}$. Bao thứ hai nặng gấp đôi bao thứ nhất. Bao thứ ba nặng bằng một nửa bao thứ hai. Tổng cân nặng của cả ba bao là
\choice
{$60\text{ kg}$}
{$70\text{ kg}$}
{$90\text{ kg}$}
{\True $80\text{ kg}$}
\loigiai{
Bao hai nặng: $20 \times 2 = 40\text{ kg}$. Bao ba nặng: $40 : 2 = 20\text{ kg}$. Tổng ba bao: $20 + 40 + 20 = 80\text{ kg}$.
\textbf{Đáp án của câu hỏi là D}
}
\end{ex}

\begin{ex} %[3D2V24]
Hiện nay mẹ 32 tuổi, con 4 tuổi. Hỏi hiện nay tuổi mẹ gấp mấy lần tuổi con?
\choice
{\True $8$ lần}
{$6$ lần}
{$7$ lần}
{$9$ lần}
\loigiai{
Tuổi mẹ gấp tuổi con số lần là: $32 : 4 = 8$ (lần).
\textbf{Đáp án của câu hỏi là A}
}
\end{ex}

\begin{ex} %[3D2V24]
Bạn Bình có 15 viên bi. Nếu Nam có thêm 3 viên bi thì số bi của Nam sẽ gấp 3 lần số bi của Bình. Hỏi Nam có bao nhiêu viên bi?
\choice
{$45$ viên bi}
{$48$ viên bi}
{\True $42$ viên bi}
{$39$ viên bi}
\loigiai{
Nếu Nam có thêm 3 viên bi thì Nam có: $15 \times 3 = 45$ viên. Số bi hiện tại của Nam: $45 - 3 = 42$ viên bi.
\textbf{Đáp án của câu hỏi là C}
}
\end{ex}

\textbf{Phần 2. Trắc nghiệm đúng sai}
\setcounterpageref{ex}{0}

\begin{ex} %[3D2V24]
Xét tính đúng sai của các câu nói về việc gấp một số lên nhiều lần:
\choiceTF
{\True Muốn gấp số $a$ lên $k$ lần ta tính $a \times k$.}
{Gấp số 15 lên 3 lần thì được kết quả là 18.}
{\True Gấp một số chẵn lên một số lần bất kì luôn được kết quả là số chẵn.}
{Gấp 1 lên 100 lần thì bằng 10.}
\loigiai{
\begin{itemize}
    \item Ý a) Đúng: Quy tắc cơ bản.
    \item Ý b) Sai: Sai vì $15 + 3 = 18$ là thêm 3 đơn vị, gấp 3 lần phải là $15 \times 3 = 45$.
    \item Ý c) Đúng: Số chẵn nhân số nào cũng ra số chẵn.
    \item Ý d) Sai: Sai vì $1 \times 100 = 100$.
\end{itemize}
\begin{itemize}
    \item a) \textbf{Đúng}.
    \item b) \textbf{Sai}.
    \item c) \textbf{Đúng}.
    \item d) \textbf{Sai}.
\end{itemize}
}
\end{ex}

\begin{ex} %[3D2V24]
Lớp 3B có 14 bạn nữ, số bạn nam gấp đôi số bạn nữ:
\choiceTF
{\True Lớp 3B có 28 bạn nam.}
{\True Tổng số học sinh lớp 3B là 42 bạn.}
{Số bạn nam nhiều hơn số bạn nữ là 28 bạn.}
{\True Số học sinh cả lớp gấp 3 lần số bạn nữ.}
\loigiai{
\begin{itemize}
    \item Ý a) Đúng: $14 \times 2 = 28$ bạn.
    \item Ý b) Đúng: $14 + 28 = 42$ bạn.
    \item Ý c) Sai: Số nam nhiều hơn số nữ là: $28 - 14 = 14$ bạn.
    \item Ý d) Đúng: $42 : 14 = 3$ lần.
\end{itemize}
\begin{itemize}
    \item a) \textbf{Đúng}.
    \item b) \textbf{Đúng}.
    \item c) \textbf{Sai}.
    \item d) \textbf{Đúng}.
\end{itemize}
}
\end{ex}

\textbf{Phần 3. Trả lời ngắn}
\setcounterpageref{ex}{0}

\begin{ex} %[3D2V24]
Gấp số 18 lên 4 lần rồi trừ đi 25 ta được số nào?
\shortans{47}
\loigiai{
\textbf{Đáp số:} $47$\\
$18 \times 4 = 72$; $72 - 25 = 47$.
}
\end{ex}

\begin{ex} %[3D2V24]
Bác Ba có 16 con vịt. Số con gà gấp 5 lần số con vịt. Hỏi bác Ba nuôi tất cả bao nhiêu con gà và vịt?
\shortans{96}
\loigiai{
\textbf{Đáp số:} $96$\\
Số gà: $16 \times 5 = 80$ con. Tổng gà và vịt: $16 + 80 = 96$ con.
}
\end{ex}

\begin{ex} %[3D2V24]
Một đoạn thẳng dài $12\text{ cm}$. Gấp đoạn thẳng đó lên 6 lần thì được đoạn thẳng dài bao nhiêu xăng-ti-mét?
\shortans{72}
\loigiai{
\textbf{Đáp số:} $72$\\
$12 \times 6 = 72\text{ (cm)}$.
}
\end{ex}

\begin{ex} %[3D2V24]
Năm nay con 6 tuổi. Tuổi mẹ gấp 5 lần tuổi con, tuổi bà gấp đôi tuổi mẹ. Hỏi năm nay bà bao nhiêu tuổi?
\shortans{60}
\loigiai{
\textbf{Đáp số:} $60$\\
Tuổi mẹ: $6 \times 5 = 30$ tuổi. Tuổi bà: $30 \times 2 = 60$ tuổi.
}
\end{ex}

\begin{ex} %[3D2V24]
Tìm một số biết rằng nếu gấp số đó lên 3 lần rồi bớt đi 12 thì được 45.
\shortans{19}
\loigiai{
\textbf{Đáp số:} $19$\\
Số đó nhân 3 bằng: $45 + 12 = 57$. Số cần tìm: $57 : 3 = 19$.
}
\end{ex}

\begin{ex} %[3D2V24]
Một hình chữ nhật có chiều rộng $15\text{ m}$, chiều dài gấp 3 lần chiều rộng. Tính chu vi của hình chữ nhật đó theo mét.
\shortans{120}
\loigiai{
\textbf{Đáp số:} $120$\\
Chiều dài: $15 \times 3 = 45\text{ m}$. Chu vi: $(45 + 15) \times 2 = 120\text{ m}$.
}
\end{ex}

\newpage

\TaoPhieuKiemTra{Chủ đề 4 -- Phép nhân, phép chia trong phạm vi 100}{25}{Phép chia hết và phép chia có dư}{https://loptoanthaylinh.vn/thi-online/toan-3/bai-25}
\textbf{Phần 1. Trắc nghiệm lựa chọn}
\setcounterpageref{ex}{0}

\begin{ex} %[3D2V25]
Trong một phép chia có số chia là 7, số dư lớn nhất có thể có của phép chia đó là
\choice
{\True $6$}
{$7$}
{$5$}
{$4$}
\loigiai{
Số dư luôn luôn bé hơn số chia. Với số chia là 7 thì số dư lớn nhất là $7 - 1 = 6$.
\textbf{Đáp án của câu hỏi là A}
}
\end{ex}

\begin{ex} %[3D2V25]
Thực hiện phép chia: $47 : 6$. Thương và số dư lần lượt là
\choice
{$8$ và dư $1$}
{$7$ và dư $6$}
{\True $7$ và dư $5$}
{$6$ và dư $5$}
\loigiai{
Ta có: $47 = 6 \times 7 + 5$ (với $5 < 6$). Thương bằng 7 và dư 5.
\textbf{Đáp án của câu hỏi là C}
}
\end{ex}

\begin{ex} %[3D2V25]
Một người có 35 quả cam chia đều vào các túi, mỗi túi có 4 quả. Hỏi cần ít nhất bao nhiêu túi để đựng hết số quả cam đó?
\choice
{$8$ túi}
{\True $9$ túi}
{$10$ túi}
{$7$ túi}
\loigiai{
Ta có: $35 : 4 = 8$ (dư 3 quả). Cần thêm 1 túi nữa để đựng 3 quả còn dư: $8 + 1 = 9$ túi.
\textbf{Đáp án của câu hỏi là B}
}
\end{ex}

\begin{ex} %[3D2V25]
Trong phép chia cho 5, thương bằng 9 và số dư là số dư lớn nhất có thể. Số bị chia của phép chia đó là
\choice
{$45$}
{$48$}
{$44$}
{\True $49$}
\loigiai{
Số dư lớn nhất là 4. Số bị chia là: $5 \times 9 + 4 = 49$.
\textbf{Đáp án của câu hỏi là D}
}
\end{ex}

\begin{ex} %[3D2V25]
Có 45 học sinh đi đò qua sông. Mỗi chuyến đò chỉ chở được nhiều nhất 6 học sinh (không kể người lái đò). Hỏi cần ít nhất bao nhiêu chuyến đò để chở hết số học sinh đó qua sông?
\choice
{\True $8$ chuyến đò}
{$7$ chuyến đò}
{$9$ chuyến đò}
{$6$ chuyến đò}
\loigiai{
Ta có: $45 : 6 = 7$ (dư 3 học sinh). 3 học sinh còn lại cần thêm 1 chuyến đò: $7 + 1 = 8$ chuyến.
\textbf{Đáp án của câu hỏi là A}
}
\end{ex}

\begin{ex} %[3D2V25]
Một đoàn du khách gồm 26 người muốn đi taxi. Mỗi xe taxi chở được 4 người. Cần ít nhất bao nhiêu xe taxi để chở hết đoàn khách?
\choice
{$6$ xe taxi}
{$8$ xe taxi}
{\True $7$ xe taxi}
{$5$ xe taxi}
\loigiai{
Ta có: $26 : 4 = 6$ (dư 2 người). Cần thêm 1 xe nữa: $6 + 1 = 7$ xe.
\textbf{Đáp án của câu hỏi là C}
}
\end{ex}

\begin{ex} %[3D2V25]
Trong một phép chia có dư, số bị chia là 59, số chia là 8. Số dư của phép chia đó là
\choice
{$4$}
{\True $3$}
{$5$}
{$2$}
\loigiai{
Ta có: $59 : 8 = 7$ (dư 3) vì $8 \times 7 + 3 = 59$.
\textbf{Đáp án của câu hỏi là B}
}
\end{ex}

\begin{ex} %[3D2V25]
Bác nông dân có 50 quả bưởi đóng vào các thùng carton, mỗi thùng chứa 6 quả bưởi. Sau khi đóng đầy các thùng, bác còn thừa bao nhiêu quả bưởi?
\choice
{$4$ quả bưởi}
{$3$ quả bưởi}
{$1$ quả bưởi}
{\True $2$ quả bưởi}
\loigiai{
Ta có: $50 : 6 = 8$ (dư 2). Bác đóng được 8 thùng và thừa 2 quả bưởi.
\textbf{Đáp án của câu hỏi là D}
}
\end{ex}

\begin{ex} %[3D2V25]
Một sợi dây dài $75\text{ cm}$ được cắt thành các đoạn dài $8\text{ cm}$. Người ta cắt được nhiều nhất bao nhiêu đoạn như thế?
\choice
{\True $9$ đoạn}
{$8$ đoạn}
{$10$ đoạn}
{$7$ đoạn}
\loigiai{
Ta có: $75 : 8 = 9$ (dư $3\text{ cm}$). Cắt được nhiều nhất 9 đoạn.
\textbf{Đáp án của câu hỏi là A}
}
\end{ex}

\begin{ex} %[3D2V25]
Nếu chia một số tự nhiên cho 4 thì các số dư có thể nhận được là
\choice
{$1, 2, 3, 4$}
{$0, 1, 2$}
{\True $0, 1, 2, 3$}
{$1, 2, 3$}
\loigiai{
Phép chia cho 4 có thể chia hết (dư 0) hoặc dư 1, 2, 3.
\textbf{Đáp án của câu hỏi là C}
}
\end{ex}

\begin{ex} %[3D2V25]
Tìm số nhỏ nhất mà khi chia cho 7 thì có số dư là 5 và thương bằng 6. Số đó là
\choice
{$42$}
{\True $47$}
{$49$}
{$45$}
\loigiai{
Số đó là: $7 \times 6 + 5 = 47$.
\textbf{Đáp án của câu hỏi là B}
}
\end{ex}

\begin{ex} %[3D2V25]
Lớp 3C có 34 bạn xếp thành các hàng, mỗi hàng có 5 bạn. Hỏi có bao nhiêu hàng có đủ 5 bạn và hàng cuối cùng có bao nhiêu bạn?
\choice
{$6$ hàng đủ 5 bạn và hàng cuối có $5$ bạn}
{$7$ hàng đủ 5 bạn và hàng cuối có $1$ bạn}
{$5$ hàng đủ 5 bạn và hàng cuối có $9$ bạn}
{\True $6$ hàng đủ 5 bạn và hàng cuối có $4$ bạn}
\loigiai{
Ta có: $34 : 5 = 6$ (dư 4). Vậy có 6 hàng đủ 5 bạn và hàng cuối cùng có 4 bạn.
\textbf{Đáp án của câu hỏi là D}
}
\end{ex}

\textbf{Phần 2. Trắc nghiệm đúng sai}
\setcounterpageref{ex}{0}

\begin{ex} %[3D2V25]
Xét tính đúng sai của các khẳng định về phép chia có dư:
\choiceTF
{\True Trong mọi phép chia có dư, số dư luôn luôn bé hơn số chia.}
{Trong phép chia cho 6, số dư lớn nhất là 6.}
{\True Số bị chia bằng thương nhân với số chia rồi cộng với số dư.}
{Một số chia cho 5 dư 3 thì cộng thêm 1 đơn vị sẽ chia hết cho 5.}
\loigiai{
\begin{itemize}
    \item Ý a) Đúng: Định lí phép chia có dư.
    \item Ý b) Sai: Sai vì số dư lớn nhất là $6 - 1 = 5$.
    \item Ý c) Đúng: Công thức $a = b \times q + r$.
    \item Ý d) Sai: Sai vì dư 3 cộng thêm 2 đơn vị mới chia hết cho 5.
\end{itemize}
\begin{itemize}
    \item a) \textbf{Đúng}.
    \item b) \textbf{Sai}.
    \item c) \textbf{Đúng}.
    \item d) \textbf{Sai}.
\end{itemize}
}
\end{ex}

\begin{ex} %[3D2V25]
Có 38 mét vải dùng may đồng phục, mỗi bộ may hết 3 mét vải:
\choiceTF
{\True May được nhiều nhất 12 bộ đồng phục.}
{\True Còn thừa đúng 2 mét vải.}
{Cần thêm ít nhất 2 mét vải nữa để may thêm 1 bộ nữa.}
{\True May 10 bộ đồng phục hết đúng 30 mét vải.}
\loigiai{
\begin{itemize}
    \item Ý a) Đúng: $38 : 3 = 12$ (dư 2).
    \item Ý b) Đúng: $38 - 12 \times 3 = 2\text{ m}$.
    \item Ý c) Sai: Chỉ cần thêm $3 - 2 = 1\text{ m}$ vải nữa là may thêm được 1 bộ.
    \item Ý d) Đúng: $10 \times 3 = 30\text{ m}$.
\end{itemize}
\begin{itemize}
    \item a) \textbf{Đúng}.
    \item b) \textbf{Đúng}.
    \item c) \textbf{Sai}.
    \item d) \textbf{Đúng}.
\end{itemize}
}
\end{ex}

\textbf{Phần 3. Trả lời ngắn}
\setcounterpageref{ex}{0}

\begin{ex} %[3D2V25]
Tìm số dư trong phép chia: $53 : 7$.
\shortans{4}
\loigiai{
\textbf{Đáp số:} $4$\\
$53 : 7 = 7$ (dư 4) vì $7 \times 7 + 4 = 53$.
}
\end{ex}

\begin{ex} %[3D2V25]
Trong một phép chia có số chia là 8, thương bằng 7 và số dư lớn nhất có thể có. Tìm số bị chia của phép chia đó.
\shortans{63}
\loigiai{
\textbf{Đáp số:} $63$\\
Số dư lớn nhất là 7. Số bị chia: $8 \times 7 + 7 = 56 + 7 = 63$.
}
\end{ex}

\begin{ex} %[3D2V25]
Có 43 quyển vở phát đều cho 6 bạn. Hỏi mỗi bạn được nhiều nhất bao nhiêu quyển vở?
\shortans{7}
\loigiai{
\textbf{Đáp số:} $7$\\
$43 : 6 = 7$ (dư 1). Mỗi bạn được nhiều nhất 7 quyển vở.
}
\end{ex}

\begin{ex} %[3D2V25]
Cần ít nhất bao nhiêu chiếc thuyền loại 5 chỗ ngồi để chở hết 32 du khách?
\shortans{7}
\loigiai{
\textbf{Đáp số:} $7$\\
$32 : 5 = 6$ (dư 2). Cần ít nhất $6 + 1 = 7$ chiếc thuyền.
}
\end{ex}

\begin{ex} %[3D2V25]
Tìm một số biết rằng lấy số đó chia cho 6 thì được thương là 8 và số dư là 3.
\shortans{51}
\loigiai{
\textbf{Đáp số:} $51$\\
Số đó là: $6 \times 8 + 3 = 48 + 3 = 51$.
}
\end{ex}

\begin{ex} %[3D2V25]
Một người xếp 58 quả trứng vào các vỉ trứng, mỗi vỉ để được 6 quả trứng. Hỏi sau khi xếp đầy các vỉ, còn lại bao nhiêu quả trứng chưa được vào vỉ?
\shortans{4}
\loigiai{
\textbf{Đáp số:} $4$\\
$58 : 6 = 9$ (dư 4). Còn lại 4 quả trứng.
}
\end{ex}

\newpage

\TaoPhieuKiemTra{Chủ đề 4 -- Phép nhân, phép chia trong phạm vi 100}{26}{Chia số có hai chữ số cho số có một chữ số}{https://loptoanthaylinh.vn/thi-online/toan-3/bai-26}
\textbf{Phần 1. Trắc nghiệm lựa chọn}
\setcounterpageref{ex}{0}

\begin{ex} %[3D2V26]
Thực hiện phép tính: $84 : 4 + 19$. Kết quả là
\choice
{$38$}
{$42$}
{$36$}
{\True $40$}
\loigiai{
$84 : 4 = 21$; $21 + 19 = 40$.
\textbf{Đáp án của câu hỏi là D}
}
\end{ex}

\begin{ex} %[3D2V26]
Một cuộn dây thép dài $96\text{ m}$ được cắt thành 3 đoạn bằng nhau. Độ dài mỗi đoạn dây là
\choice
{$28\text{ m}$}
{\True $32\text{ m}$}
{$36\text{ m}$}
{$30\text{ m}$}
\loigiai{
Độ dài mỗi đoạn dây là: $96 : 3 = 32\text{ (m)}$.
\textbf{Đáp án của câu hỏi là B}
}
\end{ex}

\begin{ex} %[3D2V26]
Tìm $x$, biết: $5 \times x = 95 - 20$.
\choice
{$13$}
{$17$}
{\True $15$}
{$19$}
\loigiai{
$5 \times x = 75 \Rightarrow x = 75 : 5 = 15$.
\textbf{Đáp án của câu hỏi là C}
}
\end{ex}

\begin{ex} %[3D2V26]
Có 72 học sinh xếp đều thành 6 hàng. Hỏi mỗi hàng có bao nhiêu học sinh?
\choice
{\True $12$ học sinh}
{$14$ học sinh}
{$10$ học sinh}
{$16$ học sinh}
\loigiai{
Mỗi hàng có số học sinh là: $72 : 6 = 12$ (học sinh).
\textbf{Đáp án của câu hỏi là A}
}
\end{ex}

\begin{ex} %[3D2V26]
Bác thợ may dùng một tấm vải dài $68\text{ m}$ để may các bộ rèm cửa, mỗi bộ rèm cửa cần $4\text{ m}$ vải. Bác đã may được 12 bộ rèm. Hỏi bác còn có thể may thêm được bao nhiêu bộ rèm cửa nữa?
\choice
{$4$ bộ rèm}
{$6$ bộ rèm}
{$7$ bộ rèm}
{\True $5$ bộ rèm}
\loigiai{
Số bộ rèm may được từ tấm vải là: $68 : 4 = 17$ (bộ). Bác còn may được thêm: $17 - 12 = 5$ (bộ rèm).
\textbf{Đáp án của câu hỏi là D}
}
\end{ex}

\begin{ex} %[3D2V26]
Thương của phép chia số lớn nhất có hai chữ số cho số chẵn lớn nhất có một chữ số là
\choice
{$11$ (dư $7$)}
{\True $12$ (dư $3$)}
{$12$ (dư $5$)}
{$10$ (dư $9$)}
\loigiai{
Số lớn nhất có 2 chữ số là 99. Số chẵn lớn nhất có 1 chữ số là 8. Ta có: $99 : 8 = 12$ (dư 3).
\textbf{Đáp án của câu hỏi là B}
}
\end{ex}

\begin{ex} %[3D2V26]
Một cửa hàng có 90 bao gạo. Sau khi bán đi một số bao gạo thì số bao gạo còn lại bằng $\dfrac{1}{5}$ số bao gạo ban đầu. Hỏi cửa hàng đã bán đi bao nhiêu bao gạo?
\choice
{$18$ bao gạo}
{$80$ bao gạo}
{\True $72$ bao gạo}
{$64$ bao gạo}
\loigiai{
Số bao gạo còn lại: $90 : 5 = 18$ (bao). Số bao gạo đã bán: $90 - 18 = 72$ (bao gạo).
\textbf{Đáp án của câu hỏi là C}
}
\end{ex}

\begin{ex} %[3D2V26]
Tìm số tự nhiên $y$ biết: $y \times 4 + 18 = 74$.
\choice
{\True $14$}
{$16$}
{$12$}
{$15$}
\loigiai{
$y \times 4 = 74 - 18 = 56 \Rightarrow y = 56 : 4 = 14$.
\textbf{Đáp án của câu hỏi là A}
}
\end{ex}

\begin{ex} %[3D2V26]
Người ta đóng 78 lon nước ngọt vào các lốc, mỗi lốc có 6 lon. Hỏi đóng được tất cả bao nhiêu lốc nước ngọt?
\choice
{$12$ lốc}
{$14$ lốc}
{$15$ lốc}
{\True $13$ lốc}
\loigiai{
Số lốc nước ngọt đóng được là: $78 : 6 = 13$ (lốc).
\textbf{Đáp án của câu hỏi là D}
}
\end{ex}

\begin{ex} %[3D2V26]
Một sợi dây dài $85\text{ dm}$ được chia làm 5 đoạn bằng nhau. Độ dài 3 đoạn dây như thế là
\choice
{$17\text{ dm}$}
{\True $51\text{ dm}$}
{$34\text{ dm}$}
{$68\text{ dm}$}
\loigiai{
Một đoạn dài: $85 : 5 = 17\text{ (dm)}$. Ba đoạn dài: $17 \times 3 = 51\text{ (dm)}$.
\textbf{Đáp án của câu hỏi là B}
}
\end{ex}

\begin{ex} %[3D2V26]
Bác Năm thu hoạch được $92\text{ kg}$ cà chua, bác chia đều vào 4 sọt. Mỗi sọt bác lại đem bán bớt đi $8\text{ kg}$. Hỏi mỗi sọt còn lại bao nhiêu ki-lô-gam cà chua?
\choice
{$23\text{ kg}$}
{$18\text{ kg}$}
{\True $15\text{ kg}$}
{$12\text{ kg}$}
\loigiai{
Mỗi sọt ban đầu có: $92 : 4 = 23\text{ (kg)}$. Còn lại: $23 - 8 = 15\text{ (kg)}$.
\textbf{Đáp án của câu hỏi là C}
}
\end{ex}

\begin{ex} %[3D2V26]
Một hình thoi có chu vi bằng $68\text{ cm}$. Biết hình thoi có 4 cạnh dài bằng nhau. Độ dài mỗi cạnh của hình thoi là
\choice
{\True $17\text{ cm}$}
{$16\text{ cm}$}
{$18\text{ cm}$}
{$15\text{ cm}$}
\loigiai{
Độ dài mỗi cạnh là: $68 : 4 = 17\text{ (cm)}$.
\textbf{Đáp án của câu hỏi là A}
}
\end{ex}

\textbf{Phần 2. Trắc nghiệm đúng sai}
\setcounterpageref{ex}{0}

\begin{ex} %[3D2V26]
Xét tính đúng sai của các phép chia số có hai chữ số cho số có một chữ số:
\choiceTF
{\True Khi chia số có hai chữ số cho số có một chữ số, ta thực hiện chia lần lượt từ trái sang phải.}
{\True Kết quả của $78 : 3$ là $26$.}
{Phép chia $85 : 5$ có kết quả là $15$.}
{\True Trong phép chia $91 : 7$, thương nhận được là một số có hai chữ số.}
\loigiai{
\begin{itemize}
    \item Ý a) Đúng: Quy tắc chia theo thứ tự hàng chục rồi hàng đơn vị.
    \item Ý b) Đúng: $78 : 3 = 26$.
    \item Ý c) Sai: Sai vì $85 : 5 = 17$.
    \item Ý d) Đúng: $91 : 7 = 13$, là số có 2 chữ số.
\end{itemize}
\begin{itemize}
    \item a) \textbf{Đúng}.
    \item b) \textbf{Đúng}.
    \item c) \textbf{Sai}.
    \item d) \textbf{Đúng}.
\end{itemize}
}
\end{ex}

\begin{ex} %[3D2V26]
Một xưởng bánh có 96 chiếc bánh mì nướng được đóng vào các hộp 4 chiếc hoặc 6 chiếc:
\choiceTF
{\True Nếu đóng vào các hộp 4 chiếc thì đóng được 24 hộp.}
{\True Nếu đóng vào các hộp 6 chiếc thì đóng được 16 hộp.}
{Số hộp 4 chiếc ít hơn số hộp 6 chiếc là 8 hộp.}
{\True Cả hai cách đóng đều chia hết, không thừa chiếc bánh nào.}
\loigiai{
\begin{itemize}
    \item Ý a) Đúng: $96 : 4 = 24$ hộp.
    \item Ý b) Đúng: $96 : 6 = 16$ hộp.
    \item Ý c) Sai: Sai vì số hộp 4 chiếc nhiều hơn: $24 - 16 = 8$ hộp.
    \item Ý d) Đúng: 96 chia hết cho cả 4 và 6.
\end{itemize}
\begin{itemize}
    \item a) \textbf{Đúng}.
    \item b) \textbf{Đúng}.
    \item c) \textbf{Sai}.
    \item d) \textbf{Đúng}.
\end{itemize}
}
\end{ex}

\textbf{Phần 3. Trả lời ngắn}
\setcounterpageref{ex}{0}

\begin{ex} %[3D2V26]
Tính kết quả của phép chia: $96 : 4$.
\shortans{24}
\loigiai{
\textbf{Đáp số:} $24$\\
$96 : 4 = 24$.
}
\end{ex}

\begin{ex} %[3D2V26]
Một cửa hàng có 65 bao xi măng chia đều cho 5 công trình. Hỏi mỗi công trình nhận được bao nhiêu bao xi măng?
\shortans{13}
\loigiai{
\textbf{Đáp số:} $13$\\
$65 : 5 = 13$ bao xi măng.
}
\end{ex}

\begin{ex} %[3D2V26]
Tìm giá trị của $x$, biết: $x \times 6 = 78$.
\shortans{13}
\loigiai{
\textbf{Đáp số:} $13$\\
$x = 78 : 6 = 13$.
}
\end{ex}

\begin{ex} %[3D2V26]
Một lớp học có 48 học sinh, trong đó $\dfrac{1}{3}$ số học sinh là học sinh giỏi. Hỏi lớp đó có bao nhiêu học sinh giỏi?
\shortans{16}
\loigiai{
\textbf{Đáp số:} $16$\\
Số học sinh giỏi là: $48 : 3 = 16$ học sinh.
}
\end{ex}

\begin{ex} %[3D2V26]
Đoạn đường từ nhà Nam đến trường dài $84\text{ m}$. Nam đã đi được $\dfrac{1}{4}$ quãng đường. Hỏi Nam còn phải đi bao nhiêu mét nữa để đến trường?
\shortans{63}
\loigiai{
\textbf{Đáp số:} $63$\\
Nam đã đi: $84 : 4 = 21\text{ m}$. Nam còn phải đi: $84 - 21 = 63\text{ (m)}$.
}
\end{ex}

\begin{ex} %[3D2V26]
Tìm số tự nhiên có một chữ số $k$ sao cho: $98 : k = 14$.
\shortans{7}
\loigiai{
\textbf{Đáp số:} $7$\\
$k = 98 : 14 = 7$.
}
\end{ex}

\newpage

\TaoPhieuKiemTra{Chủ đề 4 -- Phép nhân, phép chia trong phạm vi 100}{27}{Giảm một số đi nhiều lần}{https://loptoanthaylinh.vn/thi-online/toan-3/bai-27}
\textbf{Phần 1. Trắc nghiệm lựa chọn}
\setcounterpageref{ex}{0}

\begin{ex} %[3D2V27]
Giảm 48 đi 4 lần ta được kết quả là
\choice
{$16$}
{$14$}
{\True $12$}
{$44$}
\loigiai{
Giảm 48 đi 4 lần: $48 : 4 = 12$.
\textbf{Đáp án của câu hỏi là C}
}
\end{ex}

\begin{ex} %[3D2V27]
Muốn giảm một số đi một số lần, ta làm như sau:
\choice
{\True Lấy số đó chia cho số lần}
{Lấy số đó trừ đi số lần}
{Lấy số đó nhân với số lần}
{Lấy số đó cộng với số lần}
\loigiai{
Quy tắc toán học: Muốn giảm một số đi nhiều lần, ta lấy số đó chia cho số lần.
\textbf{Đáp án của câu hỏi là A}
}
\end{ex}

\begin{ex} %[3D2V27]
Một cuộn dây dài $72\text{ m}$, sau khi sử dụng thì độ dài cuộn dây giảm đi 3 lần. Cuộn dây còn lại dài bao nhiêu mét?
\choice
{$18\text{ m}$}
{$21\text{ m}$}
{$36\text{ m}$}
{\True $24\text{ m}$}
\loigiai{
Độ dài còn lại là: $72 : 3 = 24\text{ (m)}$.
\textbf{Đáp án của câu hỏi là D}
}
\end{ex}

\begin{ex} %[3D2V27]
Giảm số 60 đi 5 lần rồi cộng thêm 28 ta được số nào?
\choice
{$36$}
{\True $40$}
{$42$}
{$38$}
\loigiai{
$60 : 5 = 12$; $12 + 28 = 40$.
\textbf{Đáp án của câu hỏi là B}
}
\end{ex}

\begin{ex} %[3D2V27]
Một bao gạo nặng $54\text{ kg}$. Sau khi bán một số ki-lô-gam gạo thì số gạo trong bao giảm đi 6 lần. Người ta đã bán đi bao nhiêu ki-lô-gam gạo?
\choice
{$9\text{ kg}$}
{$48\text{ kg}$}
{\True $45\text{ kg}$}
{$36\text{ kg}$}
\loigiai{
Số gạo còn lại: $54 : 6 = 9\text{ kg}$. Số gạo đã bán: $54 - 9 = 45\text{ (kg)}$.
\textbf{Đáp án của câu hỏi là C}
}
\end{ex}

\begin{ex} %[3D2V27]
Số thứ nhất là 84. Số thứ hai giảm đi 4 lần so với số thứ nhất. Tổng của hai số là
\choice
{\True $105$}
{$95$}
{$100$}
{$112$}
\loigiai{
Số thứ hai là: $84 : 4 = 21$. Tổng hai số là: $84 + 21 = 105$.
\textbf{Đáp án của câu hỏi là A}
}
\end{ex}

\begin{ex} %[3D2V27]
Một số giảm đi 3 lần rồi bớt đi 6 thì được 14. Số ban đầu là
\choice
{$54$}
{$48$}
{$66$}
{\True $60$}
\loigiai{
Số đó giảm 3 lần là: $14 + 6 = 20$. Số ban đầu là: $20 \times 3 = 60$.
\textbf{Đáp án của câu hỏi là D}
}
\end{ex}

\begin{ex} %[3D2V27]
Khẳng định nào sau đây là ĐÚNG khi phân biệt giữa 'giảm đi 4 lần' và 'giảm đi 4 đơn vị' đối với số 36?
\choice
{Giảm 4 lần được 32, giảm 4 đơn vị được 9}
{\True Giảm 4 lần được 9, giảm 4 đơn vị được 32}
{Cả hai đều cho kết quả là 9}
{Cả hai đều cho kết quả là 32}
\loigiai{
Giảm 4 lần là phép chia: $36 : 4 = 9$. Giảm 4 đơn vị là phép trừ: $36 - 4 = 32$.
\textbf{Đáp án của câu hỏi là B}
}
\end{ex}

\begin{ex} %[3D2V27]
Tuổi của bà là 72 tuổi. Tuổi của mẹ giảm đi 2 lần so với tuổi bà, tuổi của cháu giảm đi 6 lần so với tuổi mẹ. Tuổi của cháu là
\choice
{$8$ tuổi}
{$12$ tuổi}
{\True $6$ tuổi}
{$7$ tuổi}
\loigiai{
Tuổi của mẹ: $72 : 2 = 36$ tuổi. Tuổi của cháu: $36 : 6 = 6$ tuổi.
\textbf{Đáp án của câu hỏi là C}
}
\end{ex}

\begin{ex} %[3D2V27]
Một thùng nước có $64\ell$. Sau khi tháo bớt nước ra thì lượng nước trong thùng giảm đi 4 lần. Lượng nước đã tháo ra ngoài là
\choice
{\True $48\ell$}
{$16\ell$}
{$32\ell$}
{$52\ell$}
\loigiai{
Lượng nước còn lại: $64 : 4 = 16\text{ (\ell)}$. Lượng nước đã tháo ra: $64 - 16 = 48\text{ (\ell)}$.
\textbf{Đáp án của câu hỏi là A}
}
\end{ex}

\begin{ex} %[3D2V27]
Số chẵn lớn nhất có hai chữ số giảm đi 3 lần thì được kết quả là
\choice
{$33$}
{$32$}
{$31$ (dư $2$)}
{\True $32$ (dư $2$)}
\loigiai{
Số chẵn lớn nhất có hai chữ số là 98. Ta có: $98 : 3 = 32$ (dư 2).
\textbf{Đáp án của câu hỏi là D}
}
\end{ex}

\begin{ex} %[3D2V27]
Nếu giảm cạnh của một hình vuông đi 3 lần thì chu vi của hình vuông đó
\choice
{Giảm đi 4 lần}
{\True Giảm đi 3 lần}
{Giảm đi 12 lần}
{Không thay đổi}
\loigiai{
Chu vi hình vuông bằng cạnh nhân 4. Cạnh giảm đi 3 lần thì chu vi cũng giảm đi 3 lần.
\textbf{Đáp án của câu hỏi là B}
}
\end{ex}

\textbf{Phần 2. Trắc nghiệm đúng sai}
\setcounterpageref{ex}{0}

\begin{ex} %[3D2V27]
Xét tính đúng sai của các nhận định về việc giảm một số đi nhiều lần:
\choiceTF
{\True Muốn giảm số $x$ đi $n$ lần ta tính $x : n$.}
{Giảm số 40 đi 5 đơn vị cũng giống như giảm số 40 đi 5 lần.}
{\True Nếu số ban đầu chia hết cho số lần thì kết quả nhận được là một số tự nhiên.}
{Giảm 90 đi 3 lần rồi giảm tiếp 3 lần nữa thì tương đương với giảm 90 đi 6 lần.}
\loigiai{
\begin{itemize}
    \item Ý a) Đúng: Định nghĩa.
    \item Ý b) Sai: Sai vì giảm 5 đơn vị là $40 - 5 = 35$, giảm 5 lần là $40 : 5 = 8$.
    \item Ý c) Đúng: Định lí chia hết.
    \item Ý d) Sai: Sai vì giảm 3 lần rồi 3 lần nữa là chia cho 9, không phải chia cho 6.
\end{itemize}
\begin{itemize}
    \item a) \textbf{Đúng}.
    \item b) \textbf{Sai}.
    \item c) \textbf{Đúng}.
    \item d) \textbf{Sai}.
\end{itemize}
}
\end{ex}

\begin{ex} %[3D2V27]
Một người nuôi 48 con thỏ. Sau khi bán đi một số thỏ thì số thỏ còn lại giảm đi 4 lần so với ban đầu:
\choiceTF
{\True Số thỏ còn lại là 12 con.}
{\True Số thỏ đã bán là 36 con.}
{Số thỏ đã bán ít hơn số thỏ còn lại.}
{\True Số thỏ đã bán gấp 3 lần số thỏ còn lại.}
\loigiai{
\begin{itemize}
    \item Ý a) Đúng: $48 : 4 = 12$ con.
    \item Ý b) Đúng: $48 - 12 = 36$ con.
    \item Ý c) Sai: Sai vì 36 con nhiều hơn 12 con.
    \item Ý d) Đúng: $36 : 12 = 3$ lần.
\end{itemize}
\begin{itemize}
    \item a) \textbf{Đúng}.
    \item b) \textbf{Đúng}.
    \item c) \textbf{Sai}.
    \item d) \textbf{Đúng}.
\end{itemize}
}
\end{ex}

\textbf{Phần 3. Trả lời ngắn}
\setcounterpageref{ex}{0}

\begin{ex} %[3D2V27]
Giảm số 96 đi 6 lần ta được số nào?
\shortans{16}
\loigiai{
\textbf{Đáp số:} $16$\\
$96 : 6 = 16$.
}
\end{ex}

\begin{ex} %[3D2V27]
Một tấm vải dài $84\text{ m}$. Sau khi may quần áo, độ dài tấm vải giảm đi 4 lần. Hỏi tấm vải còn lại bao nhiêu mét?
\shortans{21}
\loigiai{
\textbf{Đáp số:} $21$\\
Tấm vải còn lại: $84 : 4 = 21\text{ (m)}$.
}
\end{ex}

\begin{ex} %[3D2V27]
Một số giảm đi 5 lần thì được 17. Tìm số đó.
\shortans{85}
\loigiai{
\textbf{Đáp số:} $85$\\
Số đó là: $17 \times 5 = 85$.
}
\end{ex}

\begin{ex} %[3D2V27]
Thùng dầu có $72\ell$. Sau một ngày sử dụng lượng dầu trong thùng giảm đi 3 lần. Hỏi người ta đã sử dụng bao nhiêu lít dầu?
\shortans{48}
\loigiai{
\textbf{Đáp số:} $48$\\
Dầu còn lại: $72 : 3 = 24\ell$. Dầu đã dùng: $72 - 24 = 48\text{ (\ell)}$.
}
\end{ex}

\begin{ex} %[3D2V27]
Giảm số 80 đi 4 lần rồi thêm vào 35 đơn vị ta được kết quả bằng bao nhiêu?
\shortans{55}
\loigiai{
\textbf{Đáp số:} $55$\\
$80 : 4 = 20$; $20 + 35 = 55$.
}
\end{ex}

\begin{ex} %[3D2V27]
Đoạn thẳng $AB = 60\text{ cm}$. Đoạn thẳng $CD$ giảm đi 5 lần so với $AB$. Đoạn thẳng $MN$ dài gấp đôi $CD$. Tính độ dài đoạn thẳng $MN$ theo xăng-ti-mét.
\shortans{24}
\loigiai{
\textbf{Đáp số:} $24$\\
$CD = 60 : 5 = 12\text{ cm}$. $MN = 12 \times 2 = 24\text{ (cm)}$.
}
\end{ex}

\newpage

\TaoPhieuKiemTra{Chủ đề 4 -- Phép nhân, phép chia trong phạm vi 100}{28}{Bài toán giải bằng hai bước tính}{https://loptoanthaylinh.vn/thi-online/toan-3/bai-28}
\textbf{Phần 1. Trắc nghiệm lựa chọn}
\setcounterpageref{ex}{0}

\begin{ex} %[3X2V28]
Một cửa hàng buổi sáng bán được 24 gói bánh. Buổi chiều bán được nhiều hơn buổi sáng 12 gói bánh. Hỏi cả hai buổi cửa hàng bán được bao nhiêu gói bánh?
\choice
{$48$ gói bánh}
{\True $60$ gói bánh}
{$54$ gói bánh}
{$64$ gói bánh}
\loigiai{
Buổi chiều bán được: $24 + 12 = 36$ gói.\
Cả hai buổi bán được: $24 + 36 = 60$ gói bánh.
\textbf{Đáp án của câu hỏi là B}
}
\end{ex}

\begin{ex} %[3X2V28]
An có 15 viên bi. Bình có số bi gấp đôi số bi của An. Cường có số bi ít hơn Bình 8 viên. Hỏi Cường có bao nhiêu viên bi?
\choice
{$30$ viên bi}
{$25$ viên bi}
{$18$ viên bi}
{\True $22$ viên bi}
\loigiai{
Bình có: $15 \times 2 = 30$ viên bi.\
Cường có: $30 - 8 = 22$ viên bi.
\textbf{Đáp án của câu hỏi là D}
}
\end{ex}

\begin{ex} %[3X2V28]
Mẹ mua một con gà nặng $2\text{ kg}$ và một con vịt nặng gấp 3 lần con gà. Hỏi cả gà và vịt nặng tất cả bao nhiêu ki-lô-gam?
\choice
{\True $8\text{ kg}$}
{$6\text{ kg}$}
{$7\text{ kg}$}
{$9\text{ kg}$}
\loigiai{
Con vịt nặng: $2 \times 3 = 6\text{ (kg)}$. Cả hai con nặng: $2 + 6 = 8\text{ (kg)}$.
\textbf{Đáp án của câu hỏi là A}
}
\end{ex}

\begin{ex} %[3X2V28]
Một thùng chứa $45\ell$ dầu. Người ta đã lấy ra 3 can dầu, mỗi can chứa $5\ell$ dầu. Hỏi trong thùng còn lại bao nhiêu lít dầu?
\choice
{$35\ell$}
{$25\ell$}
{\True $30\ell$}
{$40\ell$}
\loigiai{
Số dầu đã lấy ra là: $5 \times 3 = 15\text{ (\ell)}$.\
Số dầu còn lại trong thùng là: $45 - 15 = 30\text{ (\ell)}$.
\textbf{Đáp án của câu hỏi là C}
}
\end{ex}

\begin{ex} %[3X2V28]
Mai có 36 bông hoa. Mai tặng Lan $\dfrac{1}{4}$ số hoa đó rồi tặng Cúc thêm 5 bông hoa nữa. Mai còn lại bao nhiêu bông hoa?
\choice
{$27$ bông hoa}
{\True $22$ bông hoa}
{$24$ bông hoa}
{$20$ bông hoa}
\loigiai{
Mai tặng Lan: $36 : 4 = 9$ bông hoa.\
Mai còn lại sau khi tặng cả hai bạn: $36 - (9 + 5) = 36 - 14 = 22$ bông hoa.
\textbf{Đáp án của câu hỏi là B}
}
\end{ex}

\begin{ex} %[3X2V28]
Lớp 3A có 18 bạn nam. Số bạn nữ ít hơn số bạn nam là 4 bạn. Hỏi lớp 3A có tất cả bao nhiêu học sinh?
\choice
{$28$ bạn}
{$30$ bạn}
{$34$ bạn}
{\True $32$ bạn}
\loigiai{
Số bạn nữ là: $18 - 4 = 14$ bạn. Tổng số học sinh cả lớp là: $18 + 14 = 32$ bạn.
\textbf{Đáp án của câu hỏi là D}
}
\end{ex}

\begin{ex} %[3X2V28]
Bác thợ mộc cưa một cây gỗ dài $48\text{ dm}$ thành các đoạn dài $8\text{ dm}$. Mỗi lần cưa mất 5 phút. Hỏi bác cưa xong cây gỗ hết tất cả bao nhiêu phút?
\choice
{\True $25$ phút}
{$30$ phút}
{$20$ phút}
{$35$ phút}
\loigiai{
Cây gỗ được cưa thành: $48 : 8 = 6$ đoạn. Bác cần thực hiện: $6 - 1 = 5$ lần cưa.\
Tổng thời gian cưa là: $5 \times 5 = 25$ phút.
\textbf{Đáp án của câu hỏi là A}
}
\end{ex}

\begin{ex} %[3X2V28]
Một người nuôi 14 con gà mái. Số gà trống nhiều hơn gà mái là 6 con. Người đó bán đi một nửa tổng số gà. Hỏi người đó đã bán bao nhiêu con gà?
\choice
{$20$ con gà}
{$34$ con gà}
{\True $17$ con gà}
{$15$ con gà}
\loigiai{
Số gà trống là: $14 + 6 = 20$ con. Tổng số gà: $14 + 20 = 34$ con.\
Số gà đã bán là: $34 : 2 = 17$ con gà.
\textbf{Đáp án của câu hỏi là C}
}
\end{ex}

\begin{ex} %[3X2V28]
Có 4 xe chở học sinh đi tham quan, mỗi xe chở 24 bạn. Khi đến nơi có thêm 12 bạn đi xe ô tô con đến nhập đoàn. Hỏi có tất cả bao nhiêu bạn đi tham quan?
\choice
{$96$ bạn}
{\True $108$ bạn}
{$100$ bạn}
{$112$ bạn}
\loigiai{
Số bạn trên 4 xe khách: $24 \times 4 = 96$ bạn. Tổng số bạn tham quan: $96 + 12 = 108$ bạn.
\textbf{Đáp án của câu hỏi là B}
}
\end{ex}

\begin{ex} %[3X2V28]
Đoạn đường từ nhà đến công viên dài $60\text{ m}$. Bạn Hùng đi bộ được $\dfrac{1}{3}$ quãng đường thì gặp bạn Cường, hai bạn cùng đi bộ thêm $15\text{ m}$ nữa thì dừng lại nghỉ. Hỏi hai bạn còn cách công viên bao nhiêu mét?
\choice
{$20\text{ m}$}
{$35\text{ m}$}
{$30\text{ m}$}
{\True $25\text{ m}$}
\loigiai{
Quãng đường Hùng đã đi trước khi gặp Cường: $60 : 3 = 20\text{ m}$. Tổng quãng đường đã đi: $20 + 15 = 35\text{ m}$. Khoảng cách còn lại: $60 - 35 = 25\text{ m}$.
\textbf{Đáp án của câu hỏi là D}
}
\end{ex}

\begin{ex} %[3X2V28]
Một người mua 5 quyển vở cùng loại hết 45 nghìn đồng. Người đó mua thêm 3 quyển vở như thế nữa. Tổng số tiền người đó phải trả cho tất cả 8 quyển vở là
\choice
{\True $72$ nghìn đồng}
{$65$ nghìn đồng}
{$80$ nghìn đồng}
{$60$ nghìn đồng}
\loigiai{
Giá tiền một quyển vở là: $45 : 5 = 9$ nghìn đồng.\
Tổng số tiền mua 8 quyển vở là: $9 \times 8 = 72$ nghìn đồng.
\textbf{Đáp án của câu hỏi là A}
}
\end{ex}

\begin{ex} %[3X2V28]
Thùng to có $56\ell$ nước, thùng bé có $24\ell$ nước. Hỏi phải chuyển từ thùng to sang thùng bé bao nhiêu lít nước để hai thùng có lượng nước bằng nhau?
\choice
{$32\ell$}
{$8\ell$}
{\True $16\ell$}
{$20\ell$}
\loigiai{
Thùng to hơn thùng bé: $56 - 24 = 32\text{ (\ell)}$. Cần chuyển từ thùng to sang thùng bé: $32 : 2 = 16\text{ (\ell)}$.
\textbf{Đáp án của câu hỏi là C}
}
\end{ex}

\textbf{Phần 2. Trắc nghiệm đúng sai}
\setcounterpageref{ex}{0}

\begin{ex} %[3X2V28]
Xét tính đúng sai của các bước giải bài toán bằng hai bước tính:
\choiceTF
{\True Bài toán giải bằng hai bước tính luôn yêu cầu tìm một đại lượng trung gian trước khi tìm đáp số cuối cùng.}
{Chỉ có thể kết hợp phép cộng với phép trừ, không thể kết hợp phép nhân với phép chia.}
{\True Cần đọc kĩ đề bài để xác định đại lượng đã cho và đại lượng cần tìm.}
{Hai bước tính luôn cho ra hai kết quả giống hệt nhau.}
\loigiai{
\begin{itemize}
    \item Ý a) Đúng: Đặc điểm phương pháp giải.
    \item Ý b) Sai: Sai vì có thể kết hợp bất kì phép tính nào.
    \item Ý c) Đúng: Kĩ năng giải toán cơ bản.
    \item Ý d) Sai: Sai vì mỗi bước tính tìm một đối tượng khác nhau.
\end{itemize}
\begin{itemize}
    \item a) \textbf{Đúng}.
    \item b) \textbf{Sai}.
    \item c) \textbf{Đúng}.
    \item d) \textbf{Sai}.
\end{itemize}
}
\end{ex}

\begin{ex} %[3X2V28]
Nhà Lan thu hoạch được 3 bao khoai tây, mỗi bao nặng $18\text{ kg}$. Nhà Lan giữ lại $14\text{ kg}$ để ăn, số còn lại đem bán:
\choiceTF
{\True Tổng khối lượng khoai tây thu hoạch được là $54\text{ kg}$.}
{\True Số ki-lô-gam khoai tây đem bán là $40\text{ kg}$.}
{Số khoai tây đem bán ít hơn số khoai tây giữ lại để ăn.}
{\True Nếu mỗi ki-lô-gam khoai tây bán được với giá nhất định thì tiền thu được từ 40 kg khoai tây.}
\loigiai{
\begin{itemize}
    \item Ý a) Đúng: $18 \times 3 = 54\text{ kg}$.
    \item Ý b) Đúng: $54 - 14 = 40\text{ kg}$.
    \item Ý c) Sai: Sai vì 40 kg nhiều hơn 14 kg.
    \item Ý d) Đúng: Đúng với số khoai tây bán.
\end{itemize}
\begin{itemize}
    \item a) \textbf{Đúng}.
    \item b) \textbf{Đúng}.
    \item c) \textbf{Sai}.
    \item d) \textbf{Đúng}.
\end{itemize}
}
\end{ex}

\textbf{Phần 3. Trả lời ngắn}
\setcounterpageref{ex}{0}

\begin{ex} %[3X2V28]
Cửa hàng có 3 thùng dầu, mỗi thùng chứa $25\ell$. Cửa hàng đã bán được $38\ell$. Hỏi cửa hàng còn lại bao nhiêu lít dầu?
\shortans{37}
\loigiai{
\textbf{Đáp số:} $37$\\
Tổng số dầu: $25 \times 3 = 75\ell$. Còn lại: $75 - 38 = 37\text{ (\ell)}$.
}
\end{ex}

\begin{ex} %[3X2V28]
Nam có 16 viên bi. Dũng có số bi gấp 3 lần Nam. Hỏi cả hai bạn có tất cả bao nhiêu viên bi?
\shortans{64}
\loigiai{
\textbf{Đáp số:} $64$\\
Dũng có: $16 \times 3 = 48$ viên. Cả hai bạn có: $16 + 48 = 64$ viên bi.
}
\end{ex}

\begin{ex} %[3X2V28]
Mẹ mua $4\text{ kg}$ táo hết 80 nghìn đồng. Hỏi mẹ mua $3\text{ kg}$ táo như thế hết bao nhiêu nghìn đồng?
\shortans{60}
\loigiai{
\textbf{Đáp số:} $60$\\
Một kg táo giá: $80 : 4 = 20$ nghìn đồng. Ba kg giá: $20 \times 3 = 60$ nghìn đồng.
}
\end{ex}

\begin{ex} %[3X2V28]
Lớp 3A trồng được 32 cây xanh, lớp 3B trồng được nhiều hơn lớp 3A là 8 cây xanh. Hỏi cả hai lớp trồng được tất cả bao nhiêu cây xanh?
\shortans{72}
\loigiai{
\textbf{Đáp số:} $72$\\
Lớp 3B trồng: $32 + 8 = 40$ cây. Cả hai lớp trồng: $32 + 40 = 72$ cây xanh.
}
\end{ex}

\begin{ex} %[3X2V28]
Một tấm vải dài $60\text{ m}$. Người ta cắt ra $\dfrac{1}{3}$ tấm vải để may áo, sau đó cắt tiếp $15\text{ m}$ để may rèm. Hỏi tấm vải còn lại bao nhiêu mét?
\shortans{25}
\loigiai{
\textbf{Đáp số:} $25$\\
May áo hết: $60 : 3 = 20\text{ m}$. Tấm vải còn lại: $60 - (20 + 15) = 25\text{ (m)}$.
}
\end{ex}

\begin{ex} %[3X2V28]
Bác Năm nuôi 2 chuồng thỏ, mỗi chuồng có 12 con. Bác mua thêm 8 con thỏ nữa rồi chia đều tất cả số thỏ vào 4 lồng mới. Hỏi mỗi lồng mới có bao nhiêu con thỏ?
\shortans{8}
\loigiai{
\textbf{Đáp số:} $8$\\
Tổng số thỏ: $12 \times 2 + 8 = 32$ con. Mỗi lồng có: $32 : 4 = 8$ con thỏ.
}
\end{ex}

\newpage

\TaoPhieuKiemTra{Chủ đề 4 -- Phép nhân, phép chia trong phạm vi 100}{29}{Luyện tập chung phép tính}{https://loptoanthaylinh.vn/thi-online/toan-3/bai-29}
\textbf{Phần 1. Trắc nghiệm lựa chọn}
\setcounterpageref{ex}{0}

\begin{ex} %[3D2V29]
Tính giá trị của biểu thức: $M = 84 : 4 \times 3 - 25$.
\choice
{\True $38$}
{$42$}
{$35$}
{$40$}
\loigiai{
$84 : 4 = 21$; $21 \times 3 = 63$; $63 - 25 = 38$.
\textbf{Đáp án của câu hỏi là A}
}
\end{ex}

\begin{ex} %[3D2V29]
Tìm $x$, biết: $x \times 5 + 16 = 86$.
\choice
{$12$}
{$16$}
{\True $14$}
{$15$}
\loigiai{
$x \times 5 = 86 - 16 = 70 \Rightarrow x = 70 : 5 = 14$.
\textbf{Đáp án của câu hỏi là C}
}
\end{ex}

\begin{ex} %[3D2V29]
Một cửa hàng có 96 chiếc áo sơ mi. Ngày thứ nhất bán được $\dfrac{1}{4}$ số áo đó. Ngày thứ hai bán được $\dfrac{1}{3}$ số áo còn lại. Hỏi sau hai ngày cửa hàng còn lại bao nhiêu chiếc áo sơ mi?
\choice
{$54$ chiếc áo}
{\True $48$ chiếc áo}
{$40$ chiếc áo}
{$36$ chiếc áo}
\loigiai{
Ngày thứ nhất bán: $96 : 4 = 24$ chiếc áo. Còn lại sau ngày 1: $96 - 24 = 72$ chiếc áo.\
Ngày thứ hai bán: $72 : 3 = 24$ chiếc áo. Sau 2 ngày còn lại: $72 - 24 = 48$ chiếc áo.
\textbf{Đáp án của câu hỏi là B}
}
\end{ex}

\begin{ex} %[3D2V29]
Trong một phép chia có số chia là 6, thương là 15 và số dư là số dư lớn nhất có thể. Số bị chia là
\choice
{$90$}
{$94$}
{$96$}
{\True $95$}
\loigiai{
Số dư lớn nhất là 5. Số bị chia là: $6 \times 15 + 5 = 90 + 5 = 95$.
\textbf{Đáp án của câu hỏi là D}
}
\end{ex}

\begin{ex} %[3D2V29]
Một khu vườn hình chữ nhật có chu vi $72\text{ m}$, chiều rộng là $14\text{ m}$. Chiều dài của khu vườn đó gấp mấy lần chiều rộng?
\choice
{\True Không chia hết, chiều dài là $22\text{ m}$}
{Gấp 2 lần}
{Gấp 3 lần}
{Gấp 4 lần}
\loigiai{
Nửa chu vi: $72 : 2 = 36\text{ m}$. Chiều dài: $36 - 14 = 22\text{ m}$. Chiều dài không chia hết cho chiều rộng ($22 : 14$).
\textbf{Đáp án của câu hỏi là A}
}
\end{ex}

\begin{ex} %[3D2V29]
Cho dãy số: $3, 6, 12, 24, \dots$ Quy luật của dãy số là số đứng sau gấp 2 lần số đứng liền trước. Số hạng tiếp theo của dãy số là
\choice
{$36$}
{$40$}
{\True $48$}
{$50$}
\loigiai{
Số hạng tiếp theo là: $24 \times 2 = 48$.
\textbf{Đáp án của câu hỏi là C}
}
\end{ex}

\begin{ex} %[3D2V29]
Có 4 hộp bút, mỗi hộp có 12 chiếc bút. Người ta lấy ra 18 chiếc bút để chia cho các bạn học sinh giỏi. Số bút còn lại được chia đều vào 6 túi. Mỗi túi có bao nhiêu chiếc bút?
\choice
{$6$ chiếc bút}
{\True $5$ chiếc bút}
{$4$ chiếc bút}
{$7$ chiếc bút}
\loigiai{
Tổng số bút: $12 \times 4 = 48$ chiếc. Số bút còn lại: $48 - 18 = 30$ chiếc. Mỗi túi có: $30 : 6 = 5$ chiếc bút.
\textbf{Đáp án của câu hỏi là B}
}
\end{ex}

\begin{ex} %[3D2V29]
Nếu lấy một số nhân với 4 rồi trừ đi 15 thì được 65. Số đó là
\choice
{$18$}
{$22$}
{$25$}
{\True $20$}
\loigiai{
Số đó nhân 4 là: $65 + 15 = 80$. Số cần tìm là: $80 : 4 = 20$.
\textbf{Đáp án của câu hỏi là D}
}
\end{ex}

\begin{ex} %[3D2V29]
Bạn Cúc gấp được 18 ngôi sao, bạn Lan gấp được số ngôi sao gấp đôi bạn Cúc. Hỏi cả hai bạn gấp được tất cả bao nhiêu ngôi sao?
\choice
{\True $54$ ngôi sao}
{$36$ ngôi sao}
{$48$ ngôi sao}
{$60$ ngôi sao}
\loigiai{
Lan gấp được: $18 \times 2 = 36$ ngôi sao. Cả hai bạn gấp được: $18 + 36 = 54$ ngôi sao.
\textbf{Đáp án của câu hỏi là A}
}
\end{ex}

\begin{ex} %[3D2V29]
Thực hiện phép chia: $89 : 7$. Kết quả của phép chia là
\choice
{Thương $12$, dư $3$}
{Thương $11$, dư $12$}
{\True Thương $12$, dư $5$}
{Thương $13$, dư $2$}
\loigiai{
Ta có: $89 = 7 \times 12 + 5$ với số dư $5 < 7$.
\textbf{Đáp án của câu hỏi là C}
}
\end{ex}

\begin{ex} %[3D2V29]
Một đoàn khách có 43 người cần thuê xe ô tô 7 chỗ (không kể tài xế). Cần thuê ít nhất bao nhiêu xe để chở hết đoàn khách đó?
\choice
{$6$ xe}
{\True $7$ xe}
{$8$ xe}
{$5$ xe}
\loigiai{
Ta có: $43 : 7 = 6$ (dư 1 người). Cần ít nhất $6 + 1 = 7$ xe ô tô.
\textbf{Đáp án của câu hỏi là B}
}
\end{ex}

\begin{ex} %[3D2V29]
Tổng của hai số là 64. Biết số lớn gấp 3 lần số bé. Số bé là
\choice
{$12$}
{$14$}
{$18$}
{\True $16$}
\loigiai{
Tổng số phần bằng nhau: $3 + 1 = 4$ phần. Số bé là: $64 : 4 = 16$.
\textbf{Đáp án của câu hỏi là D}
}
\end{ex}

\textbf{Phần 2. Trắc nghiệm đúng sai}
\setcounterpageref{ex}{0}

\begin{ex} %[3D2V29]
Xét tính đúng sai của các khẳng định về các phép tính nhân, chia phạm vi 100:
\choiceTF
{\True Một số tự nhiên chẵn khi chia cho 2 luôn luôn không có số dư.}
{Tích của hai số lẻ có thể là một số chẵn.}
{\True Trong biểu thức có cả phép nhân và phép cộng, ta thực hiện nhân trước, cộng sau.}
{Gấp một số lên 2 lần rồi lại chia cho 2 thì số đó tăng lên 4 lần.}
\loigiai{
\begin{itemize}
    \item Ý a) Đúng: Tính chất của số chẵn.
    \item Ý b) Sai: Sai vì lẻ nhân lẻ luôn là số lẻ.
    \item Ý c) Đúng: Thứ tự thực hiện phép tính.
    \item Ý d) Sai: Sai vì giữ nguyên giá trị ban đầu.
\end{itemize}
\begin{itemize}
    \item a) \textbf{Đúng}.
    \item b) \textbf{Sai}.
    \item c) \textbf{Đúng}.
    \item d) \textbf{Sai}.
\end{itemize}
}
\end{ex}

\begin{ex} %[3D2V29]
Bác Ba có 72 quả cam mang ra chợ bán. Buổi sáng bác bán được $\dfrac{1}{3}$ số cam. Buổi chiều bác bán được một nửa số cam còn lại:
\choiceTF
{\True Buổi sáng bác Ba bán được 24 quả cam.}
{\True Sau buổi sáng bác Ba còn lại 48 quả cam.}
{\True Buổi chiều bác Ba bán được 24 quả cam.}
{Sau cả hai buổi bác Ba đã bán hết toàn bộ số cam.}
\loigiai{
\begin{itemize}
    \item Ý a) Đúng: $72 : 3 = 24$ quả.
    \item Ý b) Đúng: $72 - 24 = 48$ quả.
    \item Ý c) Đúng: $48 : 2 = 24$ quả.
    \item Ý d) Sai: Bác còn lại: $48 - 24 = 24$ quả cam.
\end{itemize}
\begin{itemize}
    \item a) \textbf{Đúng}.
    \item b) \textbf{Đúng}.
    \item c) \textbf{Đúng}.
    \item d) \textbf{Sai}.
\end{itemize}
}
\end{ex}

\textbf{Phần 3. Trả lời ngắn}
\setcounterpageref{ex}{0}

\begin{ex} %[3D2V29]
Tính giá trị của biểu thức: $75 : 5 + 18 \times 3$.
\shortans{69}
\loigiai{
\textbf{Đáp số:} $69$\\
$75 : 5 = 15$; $18 \times 3 = 54$; $15 + 54 = 69$.
}
\end{ex}

\begin{ex} %[3D2V29]
Có 4 chuồng vịt, mỗi chuồng có 15 con vịt. Người ta bắt đi 18 con vịt. Hỏi còn lại bao nhiêu con vịt?
\shortans{42}
\loigiai{
\textbf{Đáp số:} $42$\\
Tổng số vịt: $15 \times 4 = 60$ con. Còn lại: $60 - 18 = 42$ con.
}
\end{ex}

\begin{ex} %[3D2V29]
Tìm số tự nhiên $x$, biết: $(x - 8) \times 4 = 64$.
\shortans{24}
\loigiai{
\textbf{Đáp số:} $24$\\
$x - 8 = 64 : 4 = 16 \Rightarrow x = 16 + 8 = 24$.
}
\end{ex}

\begin{ex} %[3D2V29]
Một thanh sắt dài $95\text{ cm}$. Người ta cắt đi 3 đoạn, mỗi đoạn dài $25\text{ cm}$. Hỏi thanh sắt còn lại dài bao nhiêu xăng-ti-mét?
\shortans{20}
\loigiai{
\textbf{Đáp số:} $20$\\
Đã cắt đi: $25 \times 3 = 75\text{ cm}$. Còn lại: $95 - 75 = 20\text{ (cm)}$.
}
\end{ex}

\begin{ex} %[3D2V29]
Một đoàn gồm 38 học sinh đi thuyền qua sông, mỗi thuyền chở được 6 học sinh. Cần ít nhất bao nhiêu thuyền để chở hết các học sinh?
\shortans{7}
\loigiai{
\textbf{Đáp số:} $7$\\
$38 : 6 = 6$ (dư 2). Cần ít nhất $6 + 1 = 7$ thuyền.
}
\end{ex}

\begin{ex} %[3D2V29]
Hiện nay anh 15 tuổi, em 5 tuổi. Hỏi hiện nay tuổi anh gấp mấy lần tuổi em?
\shortans{3}
\loigiai{
\textbf{Đáp số:} $3$\\
$15 : 5 = 3$ lần.
}
\end{ex}

\newpage

% =========================================================================
% CHỦ ĐỀ 5 - MỘT SỐ ĐƠN VỊ ĐO ĐỘ DÀI, KHỐI LƯỢNG, DUNG TÍCH, NHIỆT ĐỘ
% =========================================================================
\section{Một số đơn vị đo độ dài, khối lượng, dung tích, nhiệt độ}

\TaoPhieuKiemTra{Chủ đề 5 -- Một số đơn vị đo độ dài, khối lượng, dung tích, nhiệt độ}{30}{Mi-li-mét}{https://loptoanthaylinh.vn/thi-online/toan-3/bai-30}
\textbf{Phần 1. Trắc nghiệm lựa chọn}
\setcounterpageref{ex}{0}

\begin{ex} %[3D3V30]
Điền số thích hợp vào chỗ chấm: $5\text{ cm } 8\text{ mm} = \dots\text{ mm}$.
\choice
{$508\text{ mm}$}
{$580\text{ mm}$}
{$50\text{ mm}$}
{\True $58\text{ mm}$}
\loigiai{
Vì $1\text{ cm} = 10\text{ mm}$ nên $5\text{ cm } 8\text{ mm} = 50\text{ mm} + 8\text{ mm} = 58\text{ mm}$.
\textbf{Đáp án của câu hỏi là D}
}
\end{ex}

\begin{ex} %[3D3V30]
Một tập giấy dày $24\text{ mm}$ gồm 120 trang giấy giống nhau. Hỏi một tập giấy tương tự dày $48\text{ mm}$ sẽ có bao nhiêu trang giấy?
\choice
{$180$ trang}
{\True $240$ trang}
{$300$ trang}
{$200$ trang}
\loigiai{
Tập giấy dày $48\text{ mm}$ gấp số lần tập $24\text{ mm}$ là: $48 : 24 = 2$ (lần).\
Số trang giấy là: $120 \times 2 = 240$ (trang).
\textbf{Đáp án của câu hỏi là B}
}
\end{ex}

\begin{ex} %[3D3V30]
Một con ốc sên mỗi phút bò được $15\text{ mm}$. Hỏi sau 6 phút con ốc sên bò được đoạn đường dài bao nhiêu xăng-ti-mét?
\choice
{$90\text{ cm}$}
{$6\text{ cm}$}
{\True $9\text{ cm}$}
{$15\text{ cm}$}
\loigiai{
Quãng đường ốc sên bò được: $15 \times 6 = 90\text{ (mm)}$. Đổi $90\text{ mm} = 9\text{ cm}$.
\textbf{Đáp án của câu hỏi là C}
}
\end{ex}

\begin{ex} %[3D3V30]
Một sợi dây đồng dài $1\text{ m}$. Người ta cắt đi một đoạn dài $35\text{ cm } 6\text{ mm}$. Đoạn dây còn lại dài bao nhiêu mi-li-mét?
\choice
{\True $644\text{ mm}$}
{$654\text{ mm}$}
{$640\text{ mm}$}
{$744\text{ mm}$}
\loigiai{
Đổi $1\text{ m} = 1000\text{ mm}$; $35\text{ cm } 6\text{ mm} = 356\text{ mm}$.\
Đoạn dây còn lại là: $1000 - 356 = 644\text{ (mm)}$.
\textbf{Đáp án của câu hỏi là A}
}
\end{ex}

\begin{ex} %[3D3V30]
Một chiếc đinh dài $4\text{ cm } 5\text{ mm}$ được đóng ngập vào thanh gỗ một đoạn $32\text{ mm}$. Phần đinh còn nhô ra ngoài thanh gỗ dài bao nhiêu mi-li-mét?
\choice
{$15\text{ mm}$}
{$12\text{ mm}$}
{$18\text{ mm}$}
{\True $13\text{ mm}$}
\loigiai{
Đổi $4\text{ cm } 5\text{ mm} = 45\text{ mm}$. Phần đinh nhô ra: $45 - 32 = 13\text{ (mm)}$.
\textbf{Đáp án của câu hỏi là D}
}
\end{ex}

\begin{ex} %[3D3V30]
Một hình tam giác có độ dài ba cạnh lần lượt là $25\text{ mm}, 35\text{ mm}$ và $4\text{ cm}$. Chu vi hình tam giác đó theo xăng-ti-mét là
\choice
{$100\text{ cm}$}
{\True $10\text{ cm}$}
{$8\text{ cm}$}
{$12\text{ cm}$}
\loigiai{
Đổi $4\text{ cm} = 40\text{ mm}$. Chu vi là: $25 + 35 + 40 = 100\text{ mm} = 10\text{ cm}$.
\textbf{Đáp án của câu hỏi là B}
}
\end{ex}

\begin{ex} %[3D3V30]
Người ta xếp chồng 8 cuốn sách giống hệt nhau thì được một chồng sách cao $9\text{ cm } 6\text{ mm}$. Độ dày của mỗi cuốn sách là
\choice
{$14\text{ mm}$}
{$10\text{ mm}$}
{\True $12\text{ mm}$}
{$16\text{ mm}$}
\loigiai{
Đổi $9\text{ cm } 6\text{ mm} = 96\text{ mm}$. Độ dày mỗi cuốn sách: $96 : 8 = 12\text{ (mm)}$.
\textbf{Đáp án của câu hỏi là C}
}
\end{ex}

\begin{ex} %[3D3V30]
So sánh hai độ dài sau: $A = 7\text{ cm } 4\text{ mm}$ và $B = 740\text{ mm} : 10$. Khẳng định đúng là
\choice
{\True $A = B$}
{$A > B$}
{$A < B$}
{Không so sánh được}
\loigiai{
$A = 74\text{ mm}$, $B = 740 : 10 = 74\text{ mm}$. Vậy $A = B$.
\textbf{Đáp án của câu hỏi là A}
}
\end{ex}

\begin{ex} %[3D3V30]
Một cây bút chì ban đầu dài $18\text{ cm}$. Sau một thời gian sử dụng, bút chì bị gọt ngắn đi $65\text{ mm}$. Chiều dài hiện tại của cây bút chì là
\choice
{$12\text{ cm } 5\text{ mm}$}
{$11\text{ cm}$}
{$12\text{ cm}$}
{\True $11\text{ cm } 5\text{ mm}$}
\loigiai{
Đổi $18\text{ cm} = 180\text{ mm}$. Chiều dài còn lại: $180 - 65 = 115\text{ mm} = 11\text{ cm } 5\text{ mm}$.
\textbf{Đáp án của câu hỏi là D}
}
\end{ex}

\begin{ex} %[3D3V30]
Một con kiến càng dài $12\text{ mm}$, con ve sầu dài gấp 3 lần con kiến càng. Con ve sầu dài hơn con kiến càng bao nhiêu mi-li-mét?
\choice
{$36\text{ mm}$}
{\True $24\text{ mm}$}
{$20\text{ mm}$}
{$18\text{ mm}$}
\loigiai{
Ve sầu dài: $12 \times 3 = 36\text{ mm}$. Dài hơn kiến càng: $36 - 12 = 24\text{ (mm)}$.
\textbf{Đáp án của câu hỏi là B}
}
\end{ex}

\begin{ex} %[3D3V30]
Một hình vuông có chu vi là $20\text{ cm}$. Độ dài cạnh của hình vuông đó tính theo mi-li-mét là
\choice
{$5\text{ mm}$}
{$500\text{ mm}$}
{\True $50\text{ mm}$}
{$25\text{ mm}$}
\loigiai{
Cạnh hình vuông là: $20 : 4 = 5\text{ cm}$. Đổi $5\text{ cm} = 50\text{ mm}$.
\textbf{Đáp án của câu hỏi là C}
}
\end{ex}

\begin{ex} %[3D3V30]
Người ta nối 5 thanh nẹp gỗ bằng nhau, mỗi thanh dài $18\text{ cm } 4\text{ mm}$ thành một thanh dài (mối nối không đáng kể). Chiều dài thanh nẹp thu được là
\choice
{\True $92\text{ cm}$}
{$90\text{ cm}$}
{$94\text{ cm}$}
{$920\text{ cm}$}
\loigiai{
Đổi $18\text{ cm } 4\text{ mm} = 184\text{ mm}$. Chiều dài thanh thu được: $184 \times 5 = 920\text{ mm} = 92\text{ cm}$.
\textbf{Đáp án của câu hỏi là A}
}
\end{ex}

\textbf{Phần 2. Trắc nghiệm đúng sai}
\setcounterpageref{ex}{0}

\begin{ex} %[3D3V30]
Xét tính đúng sai của các mối quan hệ giữa mi-li-mét và các đơn vị đo độ dài khác:
\choiceTF
{\True $1\text{ cm} = 10\text{ mm}$.}
{$1\text{ m} = 100\text{ mm}$.}
{\True $1\text{ dm} = 100\text{ mm}$.}
{$3\text{ cm } 5\text{ mm} = 305\text{ mm}$.}
\loigiai{
\begin{itemize}
    \item Ý a) Đúng: Mối quan hệ chuẩn.
    \item Ý b) Sai: Sai vì $1\text{ m} = 1000\text{ mm}$.
    \item Ý c) Đúng: Vì $1\text{ dm} = 10\text{ cm} = 100\text{ mm}$.
    \item Ý d) Sai: Sai vì $3\text{ cm } 5\text{ mm} = 35\text{ mm}$.
\end{itemize}
\begin{itemize}
    \item a) \textbf{Đúng}.
    \item b) \textbf{Sai}.
    \item c) \textbf{Đúng}.
    \item d) \textbf{Sai}.
\end{itemize}
}
\end{ex}

\begin{ex} %[3D3V30]
Một thanh kim loại dài $12\text{ cm } 8\text{ mm}$ được cưa thành 4 đoạn bằng nhau:
\choiceTF
{\True Thanh kim loại ban đầu dài $128\text{ mm}$.}
{\True Mỗi đoạn dài $32\text{ mm}$.}
{Người thợ phải cưa đúng 4 nhát cưa.}
{\True Hai đoạn cộng lại dài $6\text{ cm } 4\text{ mm}$.}
\loigiai{
\begin{itemize}
    \item Ý a) Đúng: $12\text{ cm } 8\text{ mm} = 128\text{ mm}$.
    \item Ý b) Đúng: $128 : 4 = 32\text{ mm}$.
    \item Ý c) Sai: Để chia thành 4 đoạn chỉ cần $4 - 1 = 3$ nhát cưa.
    \item Ý d) Đúng: $32 \times 2 = 64\text{ mm} = 6\text{ cm } 4\text{ mm}$.
\end{itemize}
\begin{itemize}
    \item a) \textbf{Đúng}.
    \item b) \textbf{Đúng}.
    \item c) \textbf{Sai}.
    \item d) \textbf{Đúng}.
\end{itemize}
}
\end{ex}

\textbf{Phần 3. Trả lời ngắn}
\setcounterpageref{ex}{0}

\begin{ex} %[3D3V30]
Đổi $7\text{ cm } 9\text{ mm}$ ra đơn vị mi-li-mét.
\shortans{79}
\loigiai{
\textbf{Đáp số:} $79$\\
$7\text{ cm } 9\text{ mm} = 70 + 9 = 79\text{ mm}$.
}
\end{ex}

\begin{ex} %[3D3V30]
Một cuốn sổ tay dày $18\text{ mm}$. Hỏi 5 cuốn sổ tay như thế xếp chồng lên nhau thì dày bao nhiêu mi-li-mét?
\shortans{90}
\loigiai{
\textbf{Đáp số:} $90$\\
$18 \times 5 = 90\text{ mm}$.
}
\end{ex}

\begin{ex} %[3D3V30]
Một đoạn dây kẽm dài $6\text{ cm } 5\text{ mm}$ được uốn thành một hình ngũ giác đều (5 cạnh bằng nhau). Hỏi mỗi cạnh dài bao nhiêu mi-li-mét?
\shortans{13}
\loigiai{
\textbf{Đáp số:} $13$\\
$65 : 5 = 13\text{ mm}$.
}
\end{ex}

\begin{ex} %[3D3V30]
Tính giá trị: $45\text{ mm} \times 2 - 18\text{ mm}$ theo đơn vị mi-li-mét.
\shortans{72}
\loigiai{
\textbf{Đáp số:} $72$\\
$90 - 18 = 72\text{ mm}$.
}
\end{ex}

\begin{ex} %[3D3V30]
Một tấm tôn dày $3\text{ mm}$. Hỏi chồng 25 tấm tôn như thế có độ dày là bao nhiêu mi-li-mét?
\shortans{75}
\loigiai{
\textbf{Đáp số:} $75$\\
$3 \times 25 = 75\text{ mm}$.
}
\end{ex}

\begin{ex} %[3D3V30]
Chiều dài của một chiếc kẹp giấy là $28\text{ mm}$. Nối 3 chiếc kẹp giấy móc vào nhau dài $76\text{ mm}$ do phần lồng vào nhau. Hỏi 2 mối lồng làm hao hụt tất cả bao nhiêu mi-li-mét?
\shortans{8}
\loigiai{
\textbf{Đáp số:} $8$\\
Tổng chiều dài 3 kẹp thẳng: $28 \times 3 = 84\text{ mm}$. Hao hụt: $84 - 76 = 8\text{ mm}$.
}
\end{ex}

\newpage

\TaoPhieuKiemTra{Chủ đề 5 -- Một số đơn vị đo độ dài, khối lượng, dung tích, nhiệt độ}{31}{Gam}{https://loptoanthaylinh.vn/thi-online/toan-3/bai-31}
\textbf{Phần 1. Trắc nghiệm lựa chọn}
\setcounterpageref{ex}{0}

\begin{ex} %[3D3V31]
Điền số thích hợp vào chỗ chấm: $1\text{ kg} = \dots\text{ g}$.
\choice
{$100\text{ g}$}
{$10\text{ g}$}
{\True $1000\text{ g}$}
{$10000\text{ g}$}
\loigiai{
$1\text{ kg} = 1000\text{ g}$.
\textbf{Đáp án của câu hỏi là C}
}
\end{ex}

\begin{ex} %[3D3V31]
Một quả táo nặng $180\text{ g}$, quả lê nặng hơn quả táo $45\text{ g}$. Cả hai quả nặng tất cả là
\choice
{\True $405\text{ g}$}
{$360\text{ g}$}
{$420\text{ g}$}
{$395\text{ g}$}
\loigiai{
Quả lê nặng: $180 + 45 = 225\text{ g}$. Cả hai quả nặng: $180 + 225 = 405\text{ g}$.
\textbf{Đáp án của câu hỏi là A}
}
\end{ex}

\begin{ex} %[3D3V31]
Trên đĩa cân đĩa thứ nhất đặt một gói đường và một quả cân $200\text{ g}$. Đĩa cân thứ hai đặt một quả cân $1\text{ kg}$. Biết cân thăng bằng, gói đường nặng là
\choice
{$1200\text{ g}$}
{$700\text{ g}$}
{$900\text{ g}$}
{\True $800\text{ g}$}
\loigiai{
Đổi $1\text{ kg} = 1000\text{ g}$. Gói đường nặng: $1000 - 200 = 800\text{ (g)}$.
\textbf{Đáp án của câu hỏi là D}
}
\end{ex}

\begin{ex} %[3D3V31]
Một hộp bánh quy gồm vỏ hộp nặng $80\text{ g}$ và 4 gói bánh nhỏ bên trong, mỗi gói bánh nặng $120\text{ g}$. Cả hộp bánh cân nặng là
\choice
{$480\text{ g}$}
{\True $560\text{ g}$}
{$600\text{ g}$}
{$520\text{ g}$}
\loigiai{
Bốn gói bánh nặng: $120 \times 4 = 480\text{ g}$. Cả hộp nặng: $480 + 80 = 560\text{ (g)}$.
\textbf{Đáp án của câu hỏi là B}
}
\end{ex}

\begin{ex} %[3D3V31]
Mẹ có $1\text{ kg}$ bột mì. Mẹ dùng $350\text{ g}$ để làm bánh rán và $450\text{ g}$ để làm bánh bao. Mẹ còn lại bao nhiêu gam bột mì?
\choice
{$150\text{ g}$}
{$250\text{ g}$}
{\True $200\text{ g}$}
{$300\text{ g}$}
\loigiai{
Đổi $1\text{ kg} = 1000\text{ g}$. Số bột đã dùng: $350 + 450 = 800\text{ g}$. Còn lại: $1000 - 800 = 200\text{ (g)}$.
\textbf{Đáp án của câu hỏi là C}
}
\end{ex}

\begin{ex} %[3D3V31]
Có 5 túi muối giống nhau cân nặng tất cả $750\text{ g}$. Hỏi 8 túi muối như thế cân nặng bao nhiêu gam?
\choice
{\True $1200\text{ g}$}
{$1000\text{ g}$}
{$1500\text{ g}$}
{$1150\text{ g}$}
\loigiai{
Mỗi túi muối nặng: $750 : 5 = 150\text{ g}$. Tám túi nặng: $150 \times 8 = 1200\text{ (g)}$.
\textbf{Đáp án của câu hỏi là A}
}
\end{ex}

\begin{ex} %[3D3V31]
Một quả dưa hấu nặng $2\text{ kg } 500\text{ g}$. Người ta bổ quả dưa thành 5 miếng đều nhau. Mỗi miếng dưa hấu nặng là
\choice
{$400\text{ g}$}
{$600\text{ g}$}
{$450\text{ g}$}
{\True $500\text{ g}$}
\loigiai{
Đổi $2\text{ kg } 500\text{ g} = 2500\text{ g}$. Mỗi miếng dưa nặng: $2500 : 5 = 500\text{ (g)}$.
\textbf{Đáp án của câu hỏi là D}
}
\end{ex}

\begin{ex} %[3D3V31]
Một gói mì tôm nặng $75\text{ g}$. Hỏi một thùng gồm 30 gói mì tôm như thế cân nặng bao nhiêu ki-lô-gam và gam (bỏ qua khối lượng vỏ thùng)?
\choice
{$225\text{ g}$}
{\True $2\text{ kg } 250\text{ g}$}
{$2\text{ kg } 500\text{ g}$}
{$3\text{ kg}$}
\loigiai{
Tổng cân nặng: $75 \times 30 = 2250\text{ g} = 2\text{ kg } 250\text{ g}$.
\textbf{Đáp án của câu hỏi là B}
}
\end{ex}

\begin{ex} %[3D3V31]
Cô Lan mua $1\text{ kg}$ thịt bò. Cô chia đều số thịt đó vào 4 hộp để trong tủ lạnh. Mỗi hộp thịt bò nặng là
\choice
{$200\text{ g}$}
{$300\text{ g}$}
{\True $250\text{ g}$}
{$400\text{ g}$}
\loigiai{
Đổi $1\text{ kg} = 1000\text{ g}$. Mỗi hộp nặng: $1000 : 4 = 250\text{ (g)}$.
\textbf{Đáp án của câu hỏi là C}
}
\end{ex}

\begin{ex} %[3D3V31]
Tìm số quả cân $50\text{ g}$ cần dùng để thăng bằng với một quả cân $200\text{ g}$ và một quả cân $100\text{ g}$.
\choice
{\True $6$ quả cân}
{$5$ quả cân}
{$4$ quả cân}
{$8$ quả cân}
\loigiai{
Tổng khối lượng cần cân là: $200 + 100 = 300\text{ g}$. Số quả cân $50\text{ g}$ cần là: $300 : 50 = 6$ quả.
\textbf{Đáp án của câu hỏi là A}
}
\end{ex}

\begin{ex} %[3D3V31]
Tính giá trị của biểu thức: $450\text{ g} : 9 + 85\text{ g} \times 2$.
\choice
{$200\text{ g}$}
{$210\text{ g}$}
{$230\text{ g}$}
{\True $220\text{ g}$}
\loigiai{
$450 : 9 = 50\text{ g}$; $85 \times 2 = 170\text{ g}$; $50 + 170 = 220\text{ g}$.
\textbf{Đáp án của câu hỏi là D}
}
\end{ex}

\begin{ex} %[3D3V31]
Một con gà mái nặng $1\text{ kg } 800\text{ g}$. Con vịt nhẹ hơn con gà mái $300\text{ g}$. Cả gà mái và vịt cân nặng là
\choice
{$3\text{ kg}$}
{\True $3\text{ kg } 300\text{ g}$}
{$3\text{ kg } 500\text{ g}$}
{$2\text{ kg } 800\text{ g}$}
\loigiai{
Con vịt nặng: $1800 - 300 = 1500\text{ g}$. Cả hai con nặng: $1800 + 1500 = 3300\text{ g} = 3\text{ kg } 300\text{ g}$.
\textbf{Đáp án của câu hỏi là B}
}
\end{ex}

\textbf{Phần 2. Trắc nghiệm đúng sai}
\setcounterpageref{ex}{0}

\begin{ex} %[3D3V31]
Xét tính đúng sai của các mối liên hệ về đơn vị đo gam:
\choiceTF
{\True Gam được viết tắt là $\text{g}$.}
{$1\text{ kg} = 100\text{ g}$.}
{\True Nửa ki-lô-gam bằng $500\text{ g}$.}
{$450\text{ g} + 650\text{ g} = 1\text{ kg}$.}
\loigiai{
\begin{itemize}
    \item Ý a) Đúng: Kí hiệu chuẩn.
    \item Ý b) Sai: Sai vì $1\text{ kg} = 1000\text{ g}$.
    \item Ý c) Đúng: $1000 : 2 = 500\text{ g}$.
    \item Ý d) Sai: Sai vì tổng là $1100\text{ g} = 1\text{ kg } 100\text{ g}$.
\end{itemize}
\begin{itemize}
    \item a) \textbf{Đúng}.
    \item b) \textbf{Sai}.
    \item c) \textbf{Đúng}.
    \item d) \textbf{Sai}.
\end{itemize}
}
\end{ex}

\begin{ex} %[3D3V31]
Một túi gạo có ghi khối lượng tịnh $1\text{ kg}$:
\choiceTF
{\True Khối lượng gạo thực tế bên trong túi không tính vỏ là $1000\text{ g}$.}
{\True Nếu vỏ bao bì nặng $15\text{ g}$ thì cả túi gạo và vỏ nặng $1015\text{ g}$.}
{Nếu bớt ra khỏi túi $300\text{ g}$ gạo thì trong túi còn lại $800\text{ g}$ gạo.}
{\True Hai túi gạo như thế có tổng khối lượng gạo là $2000\text{ g}$.}
\loigiai{
\begin{itemize}
    \item Ý a) Đúng: Định nghĩa khối lượng tịnh.
    \item Ý b) Đúng: Tổng khối lượng cả bì.
    \item Ý c) Sai: Sai vì còn lại $1000 - 300 = 700\text{ g}$.
    \item Ý d) Đúng: $1000 \times 2 = 2000\text{ g}$.
\end{itemize}
\begin{itemize}
    \item a) \textbf{Đúng}.
    \item b) \textbf{Đúng}.
    \item c) \textbf{Sai}.
    \item d) \textbf{Đúng}.
\end{itemize}
}
\end{ex}

\textbf{Phần 3. Trả lời ngắn}
\setcounterpageref{ex}{0}

\begin{ex} %[3D3V31]
Tính kết quả phép tính: $1\text{ kg} - 350\text{ g}$ theo đơn vị gam.
\shortans{650}
\loigiai{
\textbf{Đáp số:} $650$\\
$1000 - 350 = 650\text{ (g)}$.
}
\end{ex}

\begin{ex} %[3D3V31]
Một gói kẹo nặng $125\text{ g}$. Hỏi 4 gói kẹo như thế nặng bao nhiêu gam?
\shortans{500}
\loigiai{
\textbf{Đáp số:} $500$\\
$125 \times 4 = 500\text{ (g)}$.
}
\end{ex}

\begin{ex} %[3D3V31]
Người ta chia $600\text{ g}$ đường vào 3 lọ bằng nhau. Hỏi mỗi lọ đựng bao nhiêu gam đường?
\shortans{200}
\loigiai{
\textbf{Đáp số:} $200$\\
$600 : 3 = 200\text{ (g)}$.
}
\end{ex}

\begin{ex} %[3D3V31]
Có hai gói chè, gói thứ nhất nặng $150\text{ g}$, gói thứ hai nặng gấp đôi gói thứ nhất. Hỏi cả hai gói chè nặng bao nhiêu gam?
\shortans{450}
\loigiai{
\textbf{Đáp số:} $450$\\
Gói hai: $300\text{ g}$. Cả hai: $150 + 300 = 450\text{ (g)}$.
}
\end{ex}

\begin{ex} %[3D3V31]
Một hộp sữa bột nặng $900\text{ g}$. Sau khi dùng một số gam sữa thì lượng sữa còn lại trong hộp giảm đi 3 lần. Hỏi trong hộp còn lại bao nhiêu gam sữa (bỏ qua vỏ)?
\shortans{300}
\loigiai{
\textbf{Đáp số:} $300$\\
$900 : 3 = 300\text{ (g)}$.
}
\end{ex}

\begin{ex} %[3D3V31]
Cần bao nhiêu quả cân $200\text{ g}$ để có tổng khối lượng đúng bằng $1\text{ kg}$?
\shortans{5}
\loigiai{
\textbf{Đáp số:} $5$\\
$1000 : 200 = 5$ quả cân.
}
\end{ex}

\newpage

\TaoPhieuKiemTra{Chủ đề 5 -- Một số đơn vị đo độ dài, khối lượng, dung tích, nhiệt độ}{32}{Mi-li-lít}{https://loptoanthaylinh.vn/thi-online/toan-3/bai-32}
\textbf{Phần 1. Trắc nghiệm lựa chọn}
\setcounterpageref{ex}{0}

\begin{ex} %[3D3V32]
Điền số thích hợp vào chỗ chấm: $1\ell = \dots\text{ ml}$.
\choice
{$100\text{ ml}$}
{\True $1000\text{ ml}$}
{$10\text{ ml}$}
{$500\text{ ml}$}
\loigiai{
$1\ell = 1000\text{ ml}$.
\textbf{Đáp án của câu hỏi là B}
}
\end{ex}

\begin{ex} %[3D3V32]
Một bình nước có dung tích $1\ell$. Bạn Nam đã uống hết $350\text{ ml}$, bạn Hùng uống hết $250\text{ ml}$. Trong bình còn lại bao nhiêu mi-li-lít nước?
\choice
{$500\text{ ml}$}
{$600\text{ ml}$}
{$300\text{ ml}$}
{\True $400\text{ ml}$}
\loigiai{
Đổi $1\ell = 1000\text{ ml}$. Cả hai bạn uống: $350 + 250 = 600\text{ ml}$. Còn lại: $1000 - 600 = 400\text{ (ml)}$.
\textbf{Đáp án của câu hỏi là D}
}
\end{ex}

\begin{ex} %[3D3V32]
Một lọ si-rô ho chứa $120\text{ ml}$. Mỗi ngày em bé uống 3 lần, mỗi lần $5\text{ ml}$. Hỏi lọ si-rô đó đủ cho em bé uống trong bao nhiêu ngày?
\choice
{\True $8$ ngày}
{$6$ ngày}
{$10$ ngày}
{$12$ ngày}
\loigiai{
Mỗi ngày em bé uống: $5 \times 3 = 15\text{ ml}$. Lọ si-rô uống trong: $120 : 15 = 8$ (ngày).
\textbf{Đáp án của câu hỏi là A}
}
\end{ex}

\begin{ex} %[3D3V32]
Có 6 chai nước khoáng, mỗi chai chứa $500\text{ ml}$. Tổng dung tích nước của cả 6 chai là
\choice
{$300\text{ ml}$}
{$3000\ell$}
{\True $3\ell$}
{$2\ell$}
\loigiai{
Tổng dung tích là: $500 \times 6 = 3000\text{ ml} = 3\ell$.
\textbf{Đáp án của câu hỏi là C}
}
\end{ex}

\begin{ex} %[3D3V32]
Dùng một cốc chia vạch loại $200\text{ ml}$ đong nước vào một chiếc ca dung tích $1\ell$. Hỏi cần đong ít nhất bao nhiêu cốc đầy để đầy chiếc ca đó?
\choice
{$4$ cốc}
{\True $5$ cốc}
{$6$ cốc}
{$10$ cốc}
\loigiai{
Đổi $1\ell = 1000\text{ ml}$. Số cốc cần là: $1000 : 200 = 5$ cốc.
\textbf{Đáp án của câu hỏi là B}
}
\end{ex}

\begin{ex} %[3D3V32]
Một hộp sữa tươi có dung tích $180\text{ ml}$. Một lốc sữa gồm 4 hộp như thế có tổng dung tích sữa là
\choice
{$680\text{ ml}$}
{$740\text{ ml}$}
{$800\text{ ml}$}
{\True $720\text{ ml}$}
\loigiai{
Tổng dung tích sữa của 1 lốc là: $180 \times 4 = 720\text{ (ml)}$.
\textbf{Đáp án của câu hỏi là D}
}
\end{ex}

\begin{ex} %[3D3V32]
Tính giá trị của biểu thức: $450\text{ ml} + 150\text{ ml} \times 3$.
\choice
{\True $900\text{ ml}$}
{$800\text{ ml}$}
{$1000\text{ ml}$}
{$950\text{ ml}$}
\loigiai{
$150 \times 3 = 450\text{ ml}$; $450 + 450 = 900\text{ ml}$.
\textbf{Đáp án của câu hỏi là A}
}
\end{ex}

\begin{ex} %[3D3V32]
Một bình thủy chứa $2\ell$ nước sôi. Mẹ rót đều vào 4 chiếc ca nhỏ. Mỗi ca chứa bao nhiêu mi-li-lít nước?
\choice
{$400\text{ ml}$}
{$600\text{ ml}$}
{\True $500\text{ ml}$}
{$250\text{ ml}$}
\loigiai{
Đổi $2\ell = 2000\text{ ml}$. Mỗi ca chứa: $2000 : 4 = 500\text{ (ml)}$.
\textbf{Đáp án của câu hỏi là C}
}
\end{ex}

\begin{ex} %[3D3V32]
Một chai dầu ăn có $750\text{ ml}$. Sau 2 tuần nấu ăn, trong chai còn lại $150\text{ ml}$. Trung bình mỗi tuần gia đình đã dùng bao nhiêu mi-li-lít dầu ăn?
\choice
{$600\text{ ml}$}
{\True $300\text{ ml}$}
{$250\text{ ml}$}
{$350\text{ ml}$}
\loigiai{
Số dầu đã dùng trong 2 tuần: $750 - 150 = 600\text{ ml}$. Mỗi tuần dùng: $600 : 2 = 300\text{ (ml)}$.
\textbf{Đáp án của câu hỏi là B}
}
\end{ex}

\begin{ex} %[3D3V32]
Để pha một ly nước cam ép, bạn An dùng $80\text{ ml}$ nước cốt cam và $120\text{ ml}$ nước lọc. Hỏi để pha 4 ly như thế cần bao nhiêu mi-li-lít nước ép cam thành phẩm?
\choice
{$600\text{ ml}$}
{$750\text{ ml}$}
{$900\text{ ml}$}
{\True $800\text{ ml}$}
\loigiai{
Một ly có dung tích: $80 + 120 = 200\text{ ml}$. Bốn ly cần: $200 \times 4 = 800\text{ (ml)}$.
\textbf{Đáp án của câu hỏi là D}
}
\end{ex}

\begin{ex} %[3D3V32]
So sánh: $A = 1\ell - 250\text{ ml}$ và $B = 350\text{ ml} \times 2$.
\choice
{\True $A > B$}
{$A < B$}
{$A = B$}
{Không so sánh được}
\loigiai{
$A = 1000 - 250 = 750\text{ ml}$. $B = 350 \times 2 = 700\text{ ml}$. Vì $750 > 700$ nên $A > B$.
\textbf{Đáp án của câu hỏi là A}
}
\end{ex}

\begin{ex} %[3D3V32]
Một bể cá cảnh nhỏ chứa $5\ell$ nước. Người ta dùng một chiếc ca dung tích $250\text{ ml}$ múc nước ra ngoài 8 lần. Trong bể còn lại bao nhiêu lít nước?
\choice
{$2\ell$}
{$4\ell$}
{\True $3\ell$}
{$1\ell$}
\loigiai{
Lượng nước múc ra: $250 \times 8 = 2000\text{ ml} = 2\ell$. Trong bể còn lại: $5 - 2 = 3\ell$.
\textbf{Đáp án của câu hỏi là C}
}
\end{ex}

\textbf{Phần 2. Trắc nghiệm đúng sai}
\setcounterpageref{ex}{0}

\begin{ex} %[3D3V32]
Xét tính đúng sai của các câu nói về mi-li-lít:
\choiceTF
{\True Mi-li-lít được viết tắt là $\text{ml}$.}
{$1\ell = 100\text{ ml}$.}
{\True $500\text{ ml}$ bằng nửa lít nước.}
{$750\text{ ml} + 350\text{ ml} = 1\ell$.}
\loigiai{
\begin{itemize}
    \item Ý a) Đúng: Kí hiệu quy chuẩn.
    \item Ý b) Sai: Sai vì $1\ell = 1000\text{ ml}$.
    \item Ý c) Đúng: $1000 : 2 = 500\text{ ml}$.
    \item Ý d) Sai: Sai vì tổng bằng $1100\text{ ml} = 1\ell\text{ }100\text{ ml}$.
\end{itemize}
\begin{itemize}
    \item a) \textbf{Đúng}.
    \item b) \textbf{Sai}.
    \item c) \textbf{Đúng}.
    \item d) \textbf{Sai}.
\end{itemize}
}
\end{ex}

\begin{ex} %[3D3V32]
Một ấm đun nước chứa đầy $1\ell\text{ }500\text{ ml}$ nước:
\choiceTF
{\True Ấm chứa tất cả $1500\text{ ml}$ nước.}
{\True Nếu rót ra 3 cốc, mỗi cốc $300\text{ ml}$ thì đã rót ra $900\text{ ml}$.}
{Sau khi rót ra $900\text{ ml}$, trong ấm còn lại $700\text{ ml}$.}
{\True Lượng nước còn lại ($600\text{ ml}$) đủ rót thêm 2 cốc $300\text{ ml}$ nữa.}
\loigiai{
\begin{itemize}
    \item Ý a) Đúng: Đổi $1\ell\text{ }500\text{ ml} = 1500\text{ ml}$.
    \item Ý b) Đúng: $300 \times 3 = 900\text{ ml}$.
    \item Ý c) Sai: Sai vì còn lại $1500 - 900 = 600\text{ ml}$.
    \item Ý d) Đúng: $600 : 300 = 2$ cốc.
\end{itemize}
\begin{itemize}
    \item a) \textbf{Đúng}.
    \item b) \textbf{Đúng}.
    \item c) \textbf{Sai}.
    \item d) \textbf{Đúng}.
\end{itemize}
}
\end{ex}

\textbf{Phần 3. Trả lời ngắn}
\setcounterpageref{ex}{0}

\begin{ex} %[3D3V32]
Đổi $1\ell - 400\text{ ml}$ ra đơn vị mi-li-lít.
\shortans{600}
\loigiai{
\textbf{Đáp số:} $600$\\
$1000 - 400 = 600\text{ (ml)}$.
}
\end{ex}

\begin{ex} %[3D3V32]
Một lọ mật ong có dung tích $350\text{ ml}$. Hỏi 2 lọ mật ong như thế có tổng dung tích là bao nhiêu mi-li-lít?
\shortans{700}
\loigiai{
\textbf{Đáp số:} $700$\\
$350 \times 2 = 700\text{ (ml)}$.
}
\end{ex}

\begin{ex} %[3D3V32]
Người ta chia đều $800\text{ ml}$ nước trái cây vào 4 chiếc ly. Hỏi mỗi ly chứa bao nhiêu mi-li-lít?
\shortans{200}
\loigiai{
\textbf{Đáp số:} $200$\\
$800 : 4 = 200\text{ (ml)}$.
}
\end{ex}

\begin{ex} %[3D3V32]
Cần bao nhiêu hộp sữa tươi loại $250\text{ ml}$ để có được đúng $1\ell$ sữa?
\shortans{4}
\loigiai{
\textbf{Đáp số:} $4$\\
$1000 : 250 = 4$ hộp.
}
\end{ex}

\begin{ex} %[3D3V32]
Một bình nước có $900\text{ ml}$. Bạn An uống $\dfrac{1}{3}$ lượng nước đó. Hỏi trong bình còn lại bao nhiêu mi-li-lít nước?
\shortans{600}
\loigiai{
\textbf{Đáp số:} $600$\\
An uống: $900 : 3 = 300\text{ ml}$. Còn lại: $900 - 300 = 600\text{ (ml)}$.
}
\end{ex}

\begin{ex} %[3D3V32]
Tính giá trị của biểu thức: $120\text{ ml} \times 5 - 150\text{ ml}$ theo đơn vị mi-li-lít.
\shortans{450}
\loigiai{
\textbf{Đáp số:} $450$\\
$600 - 150 = 450\text{ (ml)}$.
}
\end{ex}

\newpage

\TaoPhieuKiemTra{Chủ đề 5 -- Một số đơn vị đo độ dài, khối lượng, dung tích, nhiệt độ}{33}{Nhiệt độ. Đơn vị đo nhiệt độ}{https://loptoanthaylinh.vn/thi-online/toan-3/bai-33}
\textbf{Phần 1. Trắc nghiệm lựa chọn}
\setcounterpageref{ex}{0}

\begin{ex} %[3D3V33]
Nhiệt độ của cơ thể người bình thường và khỏe mạnh vào khoảng
\choice
{\True $37^\circ\text{C}$}
{$40^\circ\text{C}$}
{$35^\circ\text{C}$}
{$30^\circ\text{C}$}
\loigiai{
Nhiệt độ cơ thể người bình thường là khoảng $37^\circ\text{C}$.
\textbf{Đáp án của câu hỏi là A}
}
\end{ex}

\begin{ex} %[3D3V33]
Nhiệt độ nước đang sôi (ở điều kiện tiêu chuẩn) là
\choice
{$0^\circ\text{C}$}
{$50^\circ\text{C}$}
{\True $100^\circ\text{C}$}
{$80^\circ\text{C}$}
\loigiai{
Nước sôi ở nhiệt độ $100^\circ\text{C}$.
\textbf{Đáp án của câu hỏi là C}
}
\end{ex}

\begin{ex} %[3D3V33]
Buổi sáng nhiệt độ phòng là $24^\circ\text{C}$, đến buổi trưa nhiệt độ tăng thêm $6^\circ\text{C}$. Nhiệt độ buổi trưa là
\choice
{$28^\circ\text{C}$}
{\True $30^\circ\text{C}$}
{$32^\circ\text{C}$}
{$26^\circ\text{C}$}
\loigiai{
Nhiệt độ buổi trưa là: $24 + 6 = 30^\circ\text{C}$.
\textbf{Đáp án của câu hỏi là B}
}
\end{ex}

\begin{ex} %[3D3V33]
Dụng cụ dùng để đo nhiệt độ không khí hoặc nhiệt độ cơ thể người là
\choice
{Cân đồng hồ}
{Thước đo}
{Compa}
{\True Nhiệt kế}
\loigiai{
Dụng cụ chuẩn để đo nhiệt độ là nhiệt kế.
\textbf{Đáp án của câu hỏi là D}
}
\end{ex}

\begin{ex} %[3D3V33]
Nhiệt độ tại Sa Pa là $12^\circ\text{C}$, nhiệt độ tại Hà Nội là $26^\circ\text{C}$. Hỏi nhiệt độ tại Hà Nội cao hơn nhiệt độ tại Sa Pa bao nhiêu độ C?
\choice
{\True $14^\circ\text{C}$}
{$12^\circ\text{C}$}
{$16^\circ\text{C}$}
{$10^\circ\text{C}$}
\loigiai{
Chênh lệch nhiệt độ là: $26 - 12 = 14^\circ\text{C}$.
\textbf{Đáp án của câu hỏi là A}
}
\end{ex}

\begin{ex} %[3D3V33]
Nước đá (nước đóng băng) tan ở nhiệt độ nào?
\choice
{$10^\circ\text{C}$}
{$5^\circ\text{C}$}
{\True $0^\circ\text{C}$}
{$-10^\circ\text{C}$}
\loigiai{
Nước đá tan ở nhiệt độ $0^\circ\text{C}$.
\textbf{Đáp án của câu hỏi là C}
}
\end{ex}

\begin{ex} %[3D3V33]
Khi đo thân nhiệt cho bạn Nam, nhiệt kế chỉ $38^\circ\text{C} 5\text{ vạch nhỏ}$ (mỗi vạch nhỏ là $0,1^\circ\text{C}$, tức $38,5^\circ\text{C}$). Nhận xét nào sau đây đúng nhất?
\choice
{Bạn Nam hoàn toàn bình thường}
{\True Bạn Nam đang bị sốt nhẹ, cần nghỉ ngơi và theo dõi}
{Bạn Nam bị hạ thân nhiệt nguy hiểm}
{Bạn Nam đang bị đóng băng}
\loigiai{
Thân nhiệt trên $37,5^\circ\text{C}$ là có biểu hiện sốt.
\textbf{Đáp án của câu hỏi là B}
}
\end{ex}

\begin{ex} %[3D3V33]
Hôm qua nhiệt độ cao nhất là $32^\circ\text{C}$. Hôm nay do có không khí lạnh tràn về, nhiệt độ giảm đi $7^\circ\text{C}$. Nhiệt độ cao nhất hôm nay là
\choice
{$24^\circ\text{C}$}
{$26^\circ\text{C}$}
{$27^\circ\text{C}$}
{\True $25^\circ\text{C}$}
\loigiai{
Nhiệt độ hôm nay là: $32 - 7 = 25^\circ\text{C}$.
\textbf{Đáp án của câu hỏi là D}
}
\end{ex}

\begin{ex} %[3D3V33]
Nhiệt độ ngăn đá tủ lạnh thường duy trì ở mức
\choice
{\True Dưới $0^\circ\text{C}$ (âm độ C để làm đá)}
{Khoảng $15^\circ\text{C}$}
{Khoảng $25^\circ\text{C}$}
{Khoảng $50^\circ\text{C}$}
\loigiai{
Ngăn đá phải duy trì dưới $0^\circ\text{C}$ để nước đóng thành băng đá.
\textbf{Đáp án của câu hỏi là A}
}
\end{ex}

\begin{ex} %[3D3V33]
Nhiệt độ 4 ngày trong tuần lần lượt là: Thứ Hai: $25^\circ\text{C}$; Thứ Ba: $27^\circ\text{C}$; Thứ Tư: $23^\circ\text{C}$; Thứ Năm: $29^\circ\text{C}$. Ngày có nhiệt độ thấp nhất là
\choice
{Thứ Hai}
{Thứ Ba}
{\True Thứ Tư}
{Thứ Năm}
\loigiai{
Nhiệt độ thấp nhất là $23^\circ\text{C}$ rơi vào ngày Thứ Tư.
\textbf{Đáp án của câu hỏi là C}
}
\end{ex}

\begin{ex} %[3D3V33]
Trong một phòng thí nghiệm, dung dịch được đun từ $20^\circ\text{C}$ lên $80^\circ\text{C}$. Hỏi nhiệt độ dung dịch đã tăng bao nhiêu độ C?
\choice
{$50^\circ\text{C}$}
{\True $60^\circ\text{C}$}
{$70^\circ\text{C}$}
{$40^\circ\text{C}$}
\loigiai{
Nhiệt độ tăng thêm là: $80 - 20 = 60^\circ\text{C}$.
\textbf{Đáp án của câu hỏi là B}
}
\end{ex}

\begin{ex} %[3D3V33]
Nhiệt kế thủy ngân hoạt động dựa trên hiện tượng nào?
\choice
{Sự đổi màu của chất lỏng}
{Sự bốc hơi của nước}
{Trọng lực Trái Đất}
{\True Sự dãn nở vì nhiệt của chất lỏng}
\loigiai{
Khi nhiệt độ tăng, chất lỏng bên trong cột nhiệt kế nở ra và dâng cao lên.
\textbf{Đáp án của câu hỏi là D}
}
\end{ex}

\textbf{Phần 2. Trắc nghiệm đúng sai}
\setcounterpageref{ex}{0}

\begin{ex} %[3D3V33]
Xét tính đúng sai của các câu nói về nhiệt độ:
\choiceTF
{\True Đơn vị đo nhiệt độ thường dùng ở Việt Nam là độ C ($^\circ\text{C}$).}
{Nước đá tan ở $100^\circ\text{C}$.}
{\True Nhiệt độ cơ thể người khỏe mạnh bình thường khoảng $37^\circ\text{C}$.}
{Nhiệt độ ban đêm luôn cao hơn nhiệt độ ban ngày.}
\loigiai{
\begin{itemize}
    \item Ý a) Đúng: Kiến thức chuẩn.
    \item Ý b) Sai: Sai vì nước đá tan ở $0^\circ\text{C}$.
    \item Ý c) Đúng: Thực tế sinh học.
    \item Ý d) Sai: Sai vì ban đêm thường mát hơn ban ngày.
\end{itemize}
\begin{itemize}
    \item a) \textbf{Đúng}.
    \item b) \textbf{Sai}.
    \item c) \textbf{Đúng}.
    \item d) \textbf{Sai}.
\end{itemize}
}
\end{ex}

\begin{ex} %[3D3V33]
Dự báo thời tiết 3 thành phố: Hà Nội $28^\circ\text{C}$, Đà Nẵng $32^\circ\text{C}$, Đà Lạt $18^\circ\text{C}$:
\choiceTF
{\True Thành phố nóng nhất trong ba thành phố là Đà Nẵng.}
{\True Thành phố lạnh nhất trong ba thành phố là Đà Lạt.}
{Hà Nội lạnh hơn Đà Lạt $10^\circ\text{C}$.}
{\True Chênh lệch nhiệt độ giữa Đà Nẵng và Đà Lạt là $14^\circ\text{C}$.}
\loigiai{
\begin{itemize}
    \item Ý a) Đúng: $32^\circ\text{C}$ cao nhất.
    \item Ý b) Đúng: $18^\circ\text{C}$ thấp nhất.
    \item Ý c) Sai: Hà Nội nóng hơn Đà Lạt $28 - 18 = 10^\circ\text{C}$.
    \item Ý d) Đúng: $32 - 18 = 14^\circ\text{C}$.
\end{itemize}
\begin{itemize}
    \item a) \textbf{Đúng}.
    \item b) \textbf{Đúng}.
    \item c) \textbf{Sai}.
    \item d) \textbf{Đúng}.
\end{itemize}
}
\end{ex}

\textbf{Phần 3. Trả lời ngắn}
\setcounterpageref{ex}{0}

\begin{ex} %[3D3V33]
Buổi sáng nhiệt độ là $22^\circ\text{C}$, buổi trưa nhiệt độ tăng thêm $7^\circ\text{C}$. Tính nhiệt độ buổi trưa theo độ C.
\shortans{29}
\loigiai{
\textbf{Đáp số:} $29$\\
$22 + 7 = 29^\circ\text{C}$.
}
\end{ex}

\begin{ex} %[3D3V33]
Nhiệt độ nước đun trong ấm tăng từ $25^\circ\text{C}$ lên đến khi sôi ($100^\circ\text{C}$). Hỏi nhiệt độ đã tăng bao nhiêu độ C?
\shortans{75}
\loigiai{
\textbf{Đáp số:} $75$\\
$100 - 25 = 75^\circ\text{C}$.
}
\end{ex}

\begin{ex} %[3D3V33]
Thân nhiệt của em bé là $39^\circ\text{C}$. Sau khi uống thuốc hạ sốt, thân nhiệt giảm $2^\circ\text{C}$. Tính thân nhiệt của em bé sau khi hạ sốt.
\shortans{37}
\loigiai{
\textbf{Đáp số:} $37$\\
$39 - 2 = 37^\circ\text{C}$.
}
\end{ex}

\begin{ex} %[3D3V33]
Nhiệt độ phòng điều hòa là $24^\circ\text{C}$, nhiệt độ ngoài sân là $36^\circ\text{C}$. Hỏi ngoài sân nóng hơn trong phòng bao nhiêu độ C?
\shortans{12}
\loigiai{
\textbf{Đáp số:} $12$\\
$36 - 24 = 12^\circ\text{C}$.
}
\end{ex}

\begin{ex} %[3D3V33]
Nhiệt độ cao nhất trong ngày là $33^\circ\text{C}$, nhiệt độ thấp nhất trong ngày là $21^\circ\text{C}$. Tính mức chênh lệch nhiệt độ trong ngày theo độ C.
\shortans{12}
\loigiai{
\textbf{Đáp số:} $12$\\
$33 - 21 = 12^\circ\text{C}$.
}
\end{ex}

\begin{ex} %[3D3V33]
Nhiệt độ của nước trong chậu là $18^\circ\text{C}$. Người ta đổ thêm nước nóng làm nhiệt độ tăng gấp đôi. Hỏi nhiệt độ lúc này của nước trong chậu là bao nhiêu độ C?
\shortans{36}
\loigiai{
\textbf{Đáp số:} $36$\\
$18 \times 2 = 36^\circ\text{C}$.
}
\end{ex}

\newpage

\TaoPhieuKiemTra{Chủ đề 5 -- Một số đơn vị đo độ dài, khối lượng, dung tích, nhiệt độ}{34}{Thực hành trải nghiệm đơn vị đo}{https://loptoanthaylinh.vn/thi-online/toan-3/bai-34}
\textbf{Phần 1. Trắc nghiệm lựa chọn}
\setcounterpageref{ex}{0}

\begin{ex} %[3D3V34]
Ước lượng chiều dài của một chiếc bút chì gỗ tiêu chuẩn là khoảng
\choice
{$18\text{ mm}$}
{$18\text{ m}$}
{$18\text{ dm}$}
{\True $18\text{ cm}$}
\loigiai{
Bút chì gỗ thông thường dài khoảng $15$ đến $18\text{ cm}$.
\textbf{Đáp án của câu hỏi là D}
}
\end{ex}

\begin{ex} %[3D3V34]
Ước lượng khối lượng của một quả trứng gà là khoảng
\choice
{$500\text{ g}$}
{\True $50\text{ g}$}
{$5\text{ kg}$}
{$5\text{ g}$}
\loigiai{
Một quả trứng gà nặng khoảng $50$ đến $60\text{ g}$.
\textbf{Đáp án của câu hỏi là B}
}
\end{ex}

\begin{ex} %[3D3V34]
Ước lượng dung tích của một thìa cà phê nhỏ là khoảng
\choice
{$50\text{ ml}$}
{$500\text{ ml}$}
{\True $5\text{ ml}$}
{$1\ell$}
\loigiai{
Một thìa cà phê chứa khoảng $5\text{ ml}$ chất lỏng.
\textbf{Đáp án của câu hỏi là C}
}
\end{ex}

\begin{ex} %[3D3V34]
Khi đo độ dày của một chiếc thẻ ngân hàng (hoặc thẻ học sinh), đơn vị đo thích hợp nhất là
\choice
{\True Mi-li-mét ($\text{mm}$)}
{Xăng-ti-mét ($\text{cm}$)}
{Mét ($\text{m}$)}
{Đề-xi-mét ($\text{dm}$)}
\loigiai{
Độ dày của thẻ rất mỏng (khoảng $1\text{ mm}$) nên dùng mi-li-mét là chuẩn xác nhất.
\textbf{Đáp án của câu hỏi là A}
}
\end{ex}

\begin{ex} %[3D3V34]
Một bình sữa đầy $250\text{ ml}$. Bé uống một nửa bình sữa. Bé đã uống bao nhiêu mi-li-lít sữa?
\choice
{$100\text{ ml}$}
{$150\text{ ml}$}
{$200\text{ ml}$}
{\True $125\text{ ml}$}
\loigiai{
Bé uống một nửa: $250 : 2 = 125\text{ (ml)}$.
\textbf{Đáp án của câu hỏi là D}
}
\end{ex}

\begin{ex} %[3D3V34]
Bác thợ mộc đo chiều dài mép bàn được $1\text{ m } 20\text{ cm}$. Chiều dài đó bằng bao nhiêu mi-li-mét?
\choice
{$120\text{ mm}$}
{\True $1200\text{ mm}$}
{$12000\text{ mm}$}
{$1020\text{ mm}$}
\loigiai{
Đổi $1\text{ m } 20\text{ cm} = 120\text{ cm} = 1200\text{ mm}$.
\textbf{Đáp án của câu hỏi là B}
}
\end{ex}

\begin{ex} %[3D3V34]
Mẹ mua một túi táo cân nặng $1\text{ kg } 200\text{ g}$. Mẹ lấy ra 2 quả táo cân được $350\text{ g}$. Số táo còn lại trong túi nặng là
\choice
{$950\text{ g}$}
{$800\text{ g}$}
{\True $850\text{ g}$}
{$900\text{ g}$}
\loigiai{
Đổi $1\text{ kg } 200\text{ g} = 1200\text{ g}$. Số táo còn lại: $1200 - 350 = 850\text{ (g)}$.
\textbf{Đáp án của câu hỏi là C}
}
\end{ex}

\begin{ex} %[3D3V34]
Để pha nước tắm cho bé sơ sinh vào mùa đông, nhiệt độ nước thích hợp nhất là khoảng
\choice
{\True $37^\circ\text{C} \text{ đến } 38^\circ\text{C}$}
{$20^\circ\text{C}$}
{$60^\circ\text{C}$}
{$100^\circ\text{C}$}
\loigiai{
Nhiệt độ nước tắm trẻ sơ sinh ấm ngang thân nhiệt, khoảng $37^\circ\text{C} - 38^\circ\text{C}$.
\textbf{Đáp án của câu hỏi là A}
}
\end{ex}

\begin{ex} %[3D3V34]
Dùng ca $500\text{ ml}$ múc nước đổ vào xô $10\ell$. Hỏi cần múc bao nhiêu ca đầy để đổ đầy xô?
\choice
{$10$ ca}
{$15$ ca}
{$25$ ca}
{\True $20$ ca}
\loigiai{
Đổi $10\ell = 10000\text{ ml}$. Số ca cần là: $10000 : 500 = 20$ ca.
\textbf{Đáp án của câu hỏi là D}
}
\end{ex}

\begin{ex} %[3D3V34]
Một chiếc bánh kem nặng $1\text{ kg}$. Người ta cắt thành 8 miếng đều nhau. Mỗi miếng bánh kem nặng là
\choice
{$150\text{ g}$}
{\True $125\text{ g}$}
{$100\text{ g}$}
{$200\text{ g}$}
\loigiai{
Đổi $1\text{ kg} = 1000\text{ g}$. Mỗi miếng bánh: $1000 : 8 = 125\text{ (g)}$.
\textbf{Đáp án của câu hỏi là B}
}
\end{ex}

\begin{ex} %[3D3V34]
Chiều cao của một bạn học sinh lớp 3 thường vào khoảng
\choice
{$130\text{ mm}$}
{$130\text{ m}$}
{\True $130\text{ cm}$}
{$13\text{ m}$}
\loigiai{
Học sinh lớp 3 cao khoảng $125\text{ cm} - 135\text{ cm}$.
\textbf{Đáp án của câu hỏi là C}
}
\end{ex}

\begin{ex} %[3D3V34]
Một gói bánh quy có 10 chiếc bánh, cả gói nặng $250\text{ g}$. Vỏ gói nặng $10\text{ g}$. Mỗi chiếc bánh quy nặng là
\choice
{\True $24\text{ g}$}
{$25\text{ g}$}
{$20\text{ g}$}
{$26\text{ g}$}
\loigiai{
Khối lượng của 10 chiếc bánh: $250 - 10 = 240\text{ g}$. Mỗi chiếc bánh: $240 : 10 = 24\text{ (g)}$.
\textbf{Đáp án của câu hỏi là A}
}
\end{ex}

\textbf{Phần 2. Trắc nghiệm đúng sai}
\setcounterpageref{ex}{0}

\begin{ex} %[3D3V34]
Xét tính đúng sai của các ước lượng thực tế:
\choiceTF
{\True Một gói muối ăn gia đình thường có khối lượng khoảng $500\text{ g}$ hoặc $1\text{ kg}$.}
{Một ly nước uống bình thường chứa khoảng $5\ell$ nước.}
{\True Bề dày của một cuốn vở học sinh khoảng $3\text{ mm}$ đến $5\text{ mm}$.}
{Nhiệt độ ngoài trời khi trời nắng gắt mùa hè ở miền Bắc khoảng $80^\circ\text{C}$.}
\loigiai{
\begin{itemize}
    \item Ý a) Đúng: Thực tế sản phẩm thị trường.
    \item Ý b) Sai: Sai vì ly nước chỉ chứa khoảng $200\text{ ml} - 300\text{ ml}$.
    \item Ý c) Đúng: Thực tế đo lường.
    \item Ý d) Sai: Sai vì nhiệt độ khoảng $38^\circ\text{C} - 42^\circ\text{C}$, $80^\circ\text{C}$ là gần sôi.
\end{itemize}
\begin{itemize}
    \item a) \textbf{Đúng}.
    \item b) \textbf{Sai}.
    \item c) \textbf{Đúng}.
    \item d) \textbf{Sai}.
\end{itemize}
}
\end{ex}

\begin{ex} %[3D3V34]
Bạn An thực hiện đo các đồ vật trong lớp học:
\choiceTF
{\True An dùng thước dây để đo chu vi thân cây bàng ở sân trường.}
{\True An dùng ê-ke để kiểm tra góc bảng lớp có vuông hay không.}
{An dùng nhiệt kế để đo khối lượng cặp sách.}
{\True An dùng ca có vạch chia để đo lượng nước trong bình.}
\loigiai{
\begin{itemize}
    \item Ý a) Đúng: Thước dây mềm uốn cong đo chu vi cây.
    \item Ý b) Đúng: Ứng dụng góc vuông của ê-ke.
    \item Ý c) Sai: Sai vì nhiệt kế đo nhiệt độ, không đo khối lượng.
    \item Ý d) Đúng: Ca chia vạch đo dung tích chất lỏng.
\end{itemize}
\begin{itemize}
    \item a) \textbf{Đúng}.
    \item b) \textbf{Đúng}.
    \item c) \textbf{Sai}.
    \item d) \textbf{Đúng}.
\end{itemize}
}
\end{ex}

\textbf{Phần 3. Trả lời ngắn}
\setcounterpageref{ex}{0}

\begin{ex} %[3D3V34]
Một cục tẩy dài $3\text{ cm } 5\text{ mm}$. Hỏi độ dài của cục tẩy là bao nhiêu mi-li-mét?
\shortans{35}
\loigiai{
\textbf{Đáp số:} $35$\\
$35\text{ mm}$.
}
\end{ex}

\begin{ex} %[3D3V34]
Một quả bưởi nặng $1\text{ kg } 400\text{ g}$. Hỏi quả bưởi đó nặng bao nhiêu gam?
\shortans{1400}
\loigiai{
\textbf{Đáp số:} $1400$\\
$1400\text{ g}$.
}
\end{ex}

\begin{ex} %[3D3V34]
Một chai nước ngọt có dung tích $1\ell\text{ }500\text{ ml}$. Hỏi chai nước ngọt đó chứa bao nhiêu mi-li-lít?
\shortans{1500}
\loigiai{
\textbf{Đáp số:} $1500$\\
$1500\text{ ml}$.
}
\end{ex}

\begin{ex} %[3D3V34]
Buổi chiều nhiệt độ đo được là $34^\circ\text{C}$, buổi tối nhiệt độ giảm xuống còn $26^\circ\text{C}$. Hỏi nhiệt độ đã giảm bao nhiêu độ C?
\shortans{8}
\loigiai{
\textbf{Đáp số:} $8$\\
$34 - 26 = 8^\circ\text{C}$.
}
\end{ex}

\begin{ex} %[3D3V34]
Một bao gạo $50\text{ kg}$ chia đều vào các túi nhỏ, mỗi túi chứa $5\text{ kg}$. Hỏi chia được bao nhiêu túi gạo?
\shortans{10}
\loigiai{
\textbf{Đáp số:} $10$\\
$50 : 5 = 10$ túi.
}
\end{ex}

\begin{ex} %[3D3V34]
Dùng một thìa đong $10\text{ ml}$ để đong đầy một bình $150\text{ ml}$. Hỏi cần đong bao nhiêu thìa?
\shortans{15}
\loigiai{
\textbf{Đáp số:} $15$\\
$150 : 10 = 15$ thìa.
}
\end{ex}

\newpage

\TaoPhieuKiemTra{Chủ đề 5 -- Một số đơn vị đo độ dài, khối lượng, dung tích, nhiệt độ}{35}{Luyện tập chung đo lường}{https://loptoanthaylinh.vn/thi-online/toan-3/bai-35}
\textbf{Phần 1. Trắc nghiệm lựa chọn}
\setcounterpageref{ex}{0}

\begin{ex} %[3D3V35]
Tính: $35\text{ mm} \times 4 + 60\text{ mm}$. Kết quả theo đơn vị xăng-ti-mét là
\choice
{$200\text{ cm}$}
{$2\text{ cm}$}
{\True $20\text{ cm}$}
{$18\text{ cm}$}
\loigiai{
$35 \times 4 = 140\text{ mm}$; $140 + 60 = 200\text{ mm} = 20\text{ cm}$.
\textbf{Đáp án của câu hỏi là C}
}
\end{ex}

\begin{ex} %[3D3V35]
Có ba bao thóc. Bao thứ nhất nặng $350\text{ g}$, bao thứ hai nặng gấp đôi bao thứ nhất, bao thứ ba nặng $450\text{ g}$. Tổng khối lượng của cả ba bao thóc là
\choice
{\True $1\text{ kg } 500\text{ g}$}
{$1\text{ kg } 400\text{ g}$}
{$1\text{ kg } 600\text{ g}$}
{$2\text{ kg}$}
\loigiai{
Bao hai nặng: $350 \times 2 = 700\text{ g}$. Tổng ba bao: $350 + 700 + 450 = 1500\text{ g} = 1\text{ kg } 500\text{ g}$.
\textbf{Đáp án của câu hỏi là A}
}
\end{ex}

\begin{ex} %[3D3V35]
Một bình chứa $1\ell$ nước. Người ta rót ra 3 cốc, mỗi cốc chứa $200\text{ ml}$ và 2 chén, mỗi chén chứa $50\text{ ml}$. Trong bình còn lại bao nhiêu mi-li-lít nước?
\choice
{$400\text{ ml}$}
{$250\text{ ml}$}
{$350\text{ ml}$}
{\True $300\text{ ml}$}
\loigiai{
Lượng nước đã rót: $200 \times 3 + 50 \times 2 = 600 + 100 = 700\text{ ml}$. Lượng nước còn lại: $1000 - 700 = 300\text{ (ml)}$.
\textbf{Đáp án của câu hỏi là D}
}
\end{ex}

\begin{ex} %[3D3V35]
Nhiệt độ nước trong cốc thứ nhất là $45^\circ\text{C}$, cốc thứ hai nóng hơn cốc thứ nhất $15^\circ\text{C}$. Nhiệt độ của cốc thứ hai là
\choice
{$55^\circ\text{C}$}
{\True $60^\circ\text{C}$}
{$65^\circ\text{C}$}
{$50^\circ\text{C}$}
\loigiai{
Nhiệt độ cốc thứ hai là: $45 + 15 = 60^\circ\text{C}$.
\textbf{Đáp án của câu hỏi là B}
}
\end{ex}

\begin{ex} %[3D3V35]
Một mảnh vải dài $2\text{ m } 40\text{ cm}$. Người ta cắt lấy $\dfrac{1}{3}$ mảnh vải để may túi. Phần vải còn lại dài bao nhiêu xăng-ti-mét?
\choice
{$80\text{ cm}$}
{$180\text{ cm}$}
{\True $160\text{ cm}$}
{$120\text{ cm}$}
\loigiai{
Đổi $2\text{ m } 40\text{ cm} = 240\text{ cm}$. Đã may túi: $240 : 3 = 80\text{ cm}$. Còn lại: $240 - 80 = 160\text{ (cm)}$.
\textbf{Đáp án của câu hỏi là C}
}
\end{ex}

\begin{ex} %[3D3V35]
Một con gà cân nặng $2\text{ kg } 200\text{ g}$, con vịt nhẹ hơn con gà $500\text{ g}$. Cả gà và vịt cân nặng là
\choice
{\True $3\text{ kg } 900\text{ g}$}
{$3\text{ kg } 700\text{ g}$}
{$4\text{ kg}$}
{$3\text{ kg } 800\text{ g}$}
\loigiai{
Con vịt nặng: $2200 - 500 = 1700\text{ g}$. Cả hai con nặng: $2200 + 1700 = 3900\text{ g} = 3\text{ kg } 900\text{ g}$.
\textbf{Đáp án của câu hỏi là A}
}
\end{ex}

\begin{ex} %[3D3V35]
Một can dầu chứa $10\ell$ dầu hỏa. Mỗi lít dầu hỏa cân nặng khoảng $800\text{ g}$. Vỏ can nặng $500\text{ g}$. Cả can dầu đó cân nặng bao nhiêu?
\choice
{$8\text{ kg}$}
{$9\text{ kg}$}
{$7\text{ kg } 500\text{ g}$}
{\True $8\text{ kg } 500\text{ g}$}
\loigiai{
Số dầu cân nặng: $800 \times 10 = 8000\text{ g} = 8\text{ kg}$. Cả can dầu nặng: $8\text{ kg} + 500\text{ g} = 8\text{ kg } 500\text{ g}$.
\textbf{Đáp án của câu hỏi là D}
}
\end{ex}

\begin{ex} %[3D3V35]
Có 4 chai mật ong, mỗi chai chứa $250\text{ ml}$. Người ta đã bán đi 3 chai. Lượng mật ong còn lại là
\choice
{$500\text{ ml}$}
{\True $250\text{ ml}$}
{$750\text{ ml}$}
{$1\ell$}
\loigiai{
Còn lại 1 chai mật ong, tức là $250\text{ ml}$.
\textbf{Đáp án của câu hỏi là B}
}
\end{ex}

\begin{ex} %[3D3V35]
Một thanh sắt dài $1\text{ m } 50\text{ mm}$ được cắt thành các đoạn dài $15\text{ mm}$. Số đoạn sắt cắt được là
\choice
{$50$ đoạn}
{$60$ đoạn}
{\True $70$ đoạn}
{$80$ đoạn}
\loigiai{
Đổi $1\text{ m } 50\text{ mm} = 1050\text{ mm}$. Số đoạn cắt được: $1050 : 15 = 70$ đoạn.
\textbf{Đáp án của câu hỏi là C}
}
\end{ex}

\begin{ex} %[3D3V35]
Người ta xếp 6 gói mì chính, mỗi gói nặng $150\text{ g}$ và 2 gói muối, mỗi gói nặng $500\text{ g}$ vào một giỏ quà. Tổng khối lượng hàng trong giỏ là
\choice
{\True $1\text{ kg } 900\text{ g}$}
{$1\text{ kg } 800\text{ g}$}
{$2\text{ kg}$}
{$1\text{ kg } 700\text{ g}$}
\loigiai{
Tổng khối lượng: $150 \times 6 + 500 \times 2 = 900 + 1000 = 1900\text{ g} = 1\text{ kg } 900\text{ g}$.
\textbf{Đáp án của câu hỏi là A}
}
\end{ex}

\begin{ex} %[3D3V35]
Tìm một số biết rằng lấy số đó nhân với $25\text{ ml}$ thì được đúng $1\ell$. Số đó là
\choice
{$25$}
{$30$}
{$50$}
{\True $40$}
\loigiai{
Đổi $1\ell = 1000\text{ ml}$. Số cần tìm là: $1000 : 25 = 40$.
\textbf{Đáp án của câu hỏi là D}
}
\end{ex}

\begin{ex} %[3D3V35]
Hôm qua nhiệt độ lúc trưa là $31^\circ\text{C}$. Hôm nay do trời mưa, nhiệt độ trưa giảm đi $5^\circ\text{C}$. Hôm nay nhiệt độ lúc trưa là
\choice
{$27^\circ\text{C}$}
{\True $26^\circ\text{C}$}
{$25^\circ\text{C}$}
{$28^\circ\text{C}$}
\loigiai{
Nhiệt độ hôm nay là: $31 - 5 = 26^\circ\text{C}$.
\textbf{Đáp án của câu hỏi là B}
}
\end{ex}

\textbf{Phần 2. Trắc nghiệm đúng sai}
\setcounterpageref{ex}{0}

\begin{ex} %[3D3V35]
Xét tính đúng sai của các nhận định tổng hợp về đại lượng:
\choiceTF
{\True Độ dài được đo bằng các đơn vị: $\text{mm}, \text{cm}, \text{dm}, \text{m}, \text{km}$.}
{Gam ($\text{g}$) là đơn vị đo dung tích chất lỏng.}
{\True $1\ell = 1000\text{ ml}$ và $1\text{ kg} = 1000\text{ g}$.}
{$100\text{ cm} = 1000\text{ mm}$ là sai.}
\loigiai{
\begin{itemize}
    \item Ý a) Đúng: Hệ đơn vị độ dài.
    \item Ý b) Sai: Sai vì gam là đơn vị đo khối lượng.
    \item Ý c) Đúng: Quan hệ tỉ lệ 1000.
    \item Ý d) Sai: Sai vì $100\text{ cm} = 1\text{ m} = 1000\text{ mm}$ là đúng.
\end{itemize}
\begin{itemize}
    \item a) \textbf{Đúng}.
    \item b) \textbf{Sai}.
    \item c) \textbf{Đúng}.
    \item d) \textbf{Sai}.
\end{itemize}
}
\end{ex}

\begin{ex} %[3D3V35]
Một hộp sữa tươi có khối lượng tịnh $1\text{ kg}$, dung tích ghi $950\text{ ml}$:
\choiceTF
{\True Khối lượng của hộp sữa được đo bằng gam hoặc ki-lô-gam.}
{\True Dung tích sữa bên trong được đo bằng mi-li-lít hoặc lít.}
{Một lít sữa tươi luôn có thể tích đúng bằng $100\text{ ml}$.}
{\True Uống hết một nửa hộp sữa tức là đã uống $475\text{ ml}$ sữa.}
\loigiai{
\begin{itemize}
    \item Ý a) Đúng: Đơn vị đo khối lượng.
    \item Ý b) Đúng: Đơn vị đo dung tích.
    \item Ý c) Sai: Sai vì $1\ell = 1000\text{ ml}$.
    \item Ý d) Đúng: $950 : 2 = 475\text{ ml}$.
\end{itemize}
\begin{itemize}
    \item a) \textbf{Đúng}.
    \item b) \textbf{Đúng}.
    \item c) \textbf{Sai}.
    \item d) \textbf{Đúng}.
\end{itemize}
}
\end{ex}

\textbf{Phần 3. Trả lời ngắn}
\setcounterpageref{ex}{0}

\begin{ex} %[3D3V35]
Tính giá trị biểu thức: $250\text{ g} \times 4 - 300\text{ g}$ theo đơn vị gam.
\shortans{700}
\loigiai{
\textbf{Đáp số:} $700$\\
$1000 - 300 = 700\text{ (g)}$.
}
\end{ex}

\begin{ex} %[3D3V35]
Đổi $3\text{ m } 5\text{ cm}$ ra đơn vị mi-li-mét.
\shortans{3050}
\loigiai{
\textbf{Đáp số:} $3050$\\
$3000 + 50 = 3050\text{ (mm)}$.
}
\end{ex}

\begin{ex} %[3D3V35]
Một bình thủy có $2\ell$ nước sôi. Người ta rót ra 5 tách trà, mỗi tách $150\text{ ml}$. Hỏi trong bình còn lại bao nhiêu mi-li-lít nước sôi?
\shortans{1250}
\loigiai{
\textbf{Đáp số:} $1250$\\
Đã rót: $150 \times 5 = 750\text{ ml}$. Còn lại: $2000 - 750 = 1250\text{ (ml)}$.
}
\end{ex}

\begin{ex} %[3D3V35]
Một cây kem que có nhiệt độ $-5^\circ\text{C}$, sau khi để ngoài trời nhiệt độ tăng thêm $8^\circ\text{C}$. Hỏi nhiệt độ cây kem lúc này là bao nhiêu độ C?
\shortans{3}
\loigiai{
\textbf{Đáp số:} $3$\\
$-5 + 8 = 3^\circ\text{C}$.
}
\end{ex}

\begin{ex} %[3D3V35]
Một quả dưa lưới nặng $1\text{ kg } 600\text{ g}$. Quả dưa hấu nặng gấp đôi quả dưa lưới. Hỏi quả dưa hấu nặng bao nhiêu gam?
\shortans{3200}
\loigiai{
\textbf{Đáp số:} $3200$\\
$1600 \times 2 = 3200\text{ (g)}$.
}
\end{ex}

\begin{ex} %[3D3V35]
Cần bao nhiêu ca nước dung tích $250\text{ ml}$ để đổ đầy một thùng chứa $5\ell$?
\shortans{20}
\loigiai{
\textbf{Đáp số:} $20$\\
$5000 : 250 = 20$ ca.
}
\end{ex}

\newpage

% =========================================================================
% CHỦ ĐỀ 6 - PHÉP NHÂN, PHÉP CHIA TRONG PHẠM VI 1 000
% =========================================================================
\section{Phép nhân, phép chia trong phạm vi 1 000}

\TaoPhieuKiemTra{Chủ đề 6 -- Phép nhân, phép chia trong phạm vi 1 000}{36}{Nhân số có ba chữ số với số có một chữ số}{https://loptoanthaylinh.vn/thi-online/toan-3/bai-36}
\textbf{Phần 1. Trắc nghiệm lựa chọn}
\setcounterpageref{ex}{0}

\begin{ex} %[3D2V36]
Tính giá trị của biểu thức: $A = 125 \times 4 - 180$.
\choice
{$300$}
{\True $320$}
{$350$}
{$280$}
\loigiai{
$125 \times 4 = 500$; $500 - 180 = 320$.
\textbf{Đáp án của câu hỏi là B}
}
\end{ex}

\begin{ex} %[3D2V36]
Một chuyến xe buýt chở được nhiều nhất 45 hành khách. Hỏi 6 chuyến xe buýt như thế chở được nhiều nhất bao nhiêu hành khách?
\choice
{$240$ hành khách}
{$250$ hành khách}
{$280$ hành khách}
{\True $270$ hành khách}
\loigiai{
Số hành khách nhiều nhất chở được là: $45 \times 6 = 270$ (hành khách).
\textbf{Đáp án của câu hỏi là D}
}
\end{ex}

\begin{ex} %[3D2V36]
Tìm $x$, biết: $x : 3 = 142 + 28$.
\choice
{\True $510$}
{$480$}
{$520$}
{$500$}
\loigiai{
$x : 3 = 170 \Rightarrow x = 170 \times 3 = 510$.
\textbf{Đáp án của câu hỏi là A}
}
\end{ex}

\begin{ex} %[3D2V36]
Một cuộn dây cáp dài $135\text{ m}$. Người ta dùng 4 cuộn cáp như thế để lắp đường dây điện. Chiều dài đường dây điện đã lắp là
\choice
{$500\text{ m}$}
{$520\text{ m}$}
{\True $540\text{ m}$}
{$560\text{ m}$}
\loigiai{
Tổng chiều dài cáp: $135 \times 4 = 540\text{ (m)}$.
\textbf{Đáp án của câu hỏi là C}
}
\end{ex}

\begin{ex} %[3D2V36]
Một hình vuông có cạnh dài $145\text{ cm}$. Chu vi của hình vuông đó là
\choice
{$560\text{ cm}$}
{\True $580\text{ cm}$}
{$600\text{ cm}$}
{$540\text{ cm}$}
\loigiai{
Chu vi hình vuông: $145 \times 4 = 580\text{ (cm)}$.
\textbf{Đáp án của câu hỏi là B}
}
\end{ex}

\begin{ex} %[3D2V36]
Từ bốn chữ số $1, 2, 3, 4$, người ta lập một số có ba chữ số và một số có một chữ số để được tích lớn nhất có thể. Tích lớn nhất đó là
\choice
{$964$}
{$1248$}
{$1324$}
{\True $1284$}
\loigiai{
Chọn thừa số một chữ số lớn nhất là $4$, số có 3 chữ số lớn nhất tạo từ $1, 2, 3$ là $321$. Tích lớn nhất: $321 \times 4 = 1284$.
\textbf{Đáp án của câu hỏi là D}
}
\end{ex}

\begin{ex} %[3D2V36]
Một cửa hàng nhập về 3 kiện hàng, mỗi kiện có 150 hộp sữa. Cửa hàng đã bán được 280 hộp sữa. Cửa hàng còn lại bao nhiêu hộp sữa?
\choice
{\True $170$ hộp}
{$150$ hộp}
{$180$ hộp}
{$190$ hộp}
\loigiai{
Tổng số hộp sữa: $150 \times 3 = 450$ hộp. Còn lại: $450 - 280 = 170$ hộp sữa.
\textbf{Đáp án của câu hỏi là A}
}
\end{ex}

\begin{ex} %[3D2V36]
Tính nhanh giá trị của biểu thức: $115 \times 2 + 115 \times 3 + 115$.
\choice
{$575$}
{$650$}
{\True $690$}
{$720$}
\loigiai{
$115 \times (2 + 3 + 1) = 115 \times 6 = 690$.
\textbf{Đáp án của câu hỏi là C}
}
\end{ex}

\begin{ex} %[3D2V36]
Mỗi bao gạo cân nặng $50\text{ kg}$. Một xe tải chở 8 bao gạo và 4 bao ngô, mỗi bao ngô nặng $40\text{ kg}$. Tổng khối lượng hàng trên xe tải là
\choice
{$480\text{ kg}$}
{\True $560\text{ kg}$}
{$520\text{ kg}$}
{$600\text{ kg}$}
\loigiai{
Khối lượng gạo: $50 \times 8 = 400\text{ kg}$. Khối lượng ngô: $40 \times 4 = 160\text{ kg}$. Tổng khối lượng: $400 + 160 = 560\text{ (kg)}$.
\textbf{Đáp án của câu hỏi là B}
}
\end{ex}

\begin{ex} %[3D2V36]
Nếu một thừa số tăng thêm 3 đơn vị và thừa số kia giữ nguyên bằng 124 thì tích tăng thêm bao nhiêu đơn vị?
\choice
{$124$}
{$248$}
{$496$}
{\True $372$}
\loigiai{
Tích tăng thêm bằng: $124 \times 3 = 372$ đơn vị.
\textbf{Đáp án của câu hỏi là D}
}
\end{ex}

\begin{ex} %[3D2V36]
Một trường tiểu học có 5 khối lớp, mỗi khối có 168 học sinh. Trường đó có tất cả bao nhiêu học sinh?
\choice
{\True $840$ học sinh}
{$800$ học sinh}
{$850$ học sinh}
{$820$ học sinh}
\loigiai{
Tổng số học sinh: $168 \times 5 = 840$ học sinh.
\textbf{Đáp án của câu hỏi là A}
}
\end{ex}

\begin{ex} %[3D2V36]
Bác Năm trồng 4 luống bắp cải, mỗi luống có 125 cây. Đợt rét vừa qua có 35 cây bị chết. Số cây bắp cải còn sống là
\choice
{$500$ cây}
{$475$ cây}
{\True $465$ cây}
{$450$ cây}
\loigiai{
Tổng số cây: $125 \times 4 = 500$ cây. Còn sống: $500 - 35 = 465$ cây.
\textbf{Đáp án của câu hỏi là C}
}
\end{ex}

\textbf{Phần 2. Trắc nghiệm đúng sai}
\setcounterpageref{ex}{0}

\begin{ex} %[3D2V36]
Xét tính đúng sai của phép nhân số có 3 chữ số với số có 1 chữ số:
\choiceTF
{\True Khi nhân số có ba chữ số với số có một chữ số, ta thực hiện nhân lần lượt từ hàng đơn vị sang hàng trăm.}
{$215 \times 3 = 635$.}
{\True $150 \times 4 = 600$.}
{Tích của số có ba chữ số với số có một chữ số không thể vượt quá 999.}
\loigiai{
\begin{itemize}
    \item Ý a) Đúng: Quy tắc tính chuẩn.
    \item Ý b) Sai: Sai vì $215 \times 3 = 645$.
    \item Ý c) Đúng: Tính nhẩm đúng.
    \item Ý d) Sai: Sai vì ví dụ $300 \times 4 = 1200 > 999$.
\end{itemize}
\begin{itemize}
    \item a) \textbf{Đúng}.
    \item b) \textbf{Sai}.
    \item c) \textbf{Đúng}.
    \item d) \textbf{Sai}.
\end{itemize}
}
\end{ex}

\begin{ex} %[3D2V36]
Một xưởng dệt có 3 máy hoạt động, mỗi máy dệt được 230 mét vải mỗi ngày:
\choiceTF
{\True Mỗi ngày cả 3 máy dệt được 690 mét vải.}
{\True Trong 2 ngày một máy dệt được 460 mét vải.}
{Cả xưởng dệt trong 2 ngày dệt được 1200 mét vải.}
{\True Nếu thêm 1 máy nữa cùng công suất thì mỗi ngày dệt được 920 mét vải.}
\loigiai{
\begin{itemize}
    \item Ý a) Đúng: $230 \times 3 = 690\text{ m}$.
    \item Ý b) Đúng: $230 \times 2 = 460\text{ m}$.
    \item Ý c) Sai: Sai vì $690 \times 2 = 1380\text{ m}$.
    \item Ý d) Đúng: $230 \times 4 = 920\text{ m}$.
\end{itemize}
\begin{itemize}
    \item a) \textbf{Đúng}.
    \item b) \textbf{Đúng}.
    \item c) \textbf{Sai}.
    \item d) \textbf{Đúng}.
\end{itemize}
}
\end{ex}

\textbf{Phần 3. Trả lời ngắn}
\setcounterpageref{ex}{0}

\begin{ex} %[3D2V36]
Tính giá trị biểu thức: $148 \times 5 - 240$.
\shortans{500}
\loigiai{
\textbf{Đáp số:} $500$\\
$148 \times 5 = 740$; $740 - 240 = 500$.
}
\end{ex}

\begin{ex} %[3D2V36]
Một mảnh đất hình tam giác có 3 cạnh bằng nhau, mỗi cạnh dài $165\text{ m}$. Tính chu vi mảnh đất đó theo mét.
\shortans{495}
\loigiai{
\textbf{Đáp số:} $495$\\
$165 \times 3 = 495\text{ (m)}$.
}
\end{ex}

\begin{ex} %[3D2V36]
Tìm số tự nhiên $y$, biết: $y : 4 = 175$.
\shortans{700}
\loigiai{
\textbf{Đáp số:} $700$\\
$y = 175 \times 4 = 700$.
}
\end{ex}

\begin{ex} %[3D2V36]
Một rạp chiếu phim có 6 hàng ghế, mỗi hàng có 115 chỗ ngồi. Hỏi rạp có tất cả bao nhiêu chỗ ngồi?
\shortans{690}
\loigiai{
\textbf{Đáp số:} $690$\\
$115 \times 6 = 690$ chỗ ngồi.
}
\end{ex}

\begin{ex} %[3D2V36]
Tính nhanh: $125 \times 2 \times 4$.
\shortans{1000}
\loigiai{
\textbf{Đáp số:} $1000$\\
$125 \times 8 = 1000$.
}
\end{ex}

\begin{ex} %[3D2V36]
Có 3 xe chở gạch, mỗi xe chở 275 viên gạch. Người ta dỡ xuống 350 viên gạch. Hỏi trên các xe còn lại bao nhiêu viên gạch?
\shortans{475}
\loigiai{
\textbf{Đáp số:} $475$\\
Tổng số gạch: $275 \times 3 = 825$ viên. Còn lại: $825 - 350 = 475$ viên.
}
\end{ex}

\newpage

\TaoPhieuKiemTra{Chủ đề 6 -- Phép nhân, phép chia trong phạm vi 1 000}{37}{Chia số có ba chữ số cho số có một chữ số}{https://loptoanthaylinh.vn/thi-online/toan-3/bai-37}
\textbf{Phần 1. Trắc nghiệm lựa chọn}
\setcounterpageref{ex}{0}

\begin{ex} %[3D2V37]
Thực hiện phép chia: $648 : 6$. Kết quả là
\choice
{\True $108$}
{$18$}
{$180$}
{$118$}
\loigiai{
Chia theo thứ tự: $6 : 6 = 1$; $4 : 6 = 0$ dư 4; $48 : 6 = 8$. Kết quả là $108$.
\textbf{Đáp án của câu hỏi là A}
}
\end{ex}

\begin{ex} %[3D2V37]
Một cuộn dây thép dài $525\text{ m}$ được cắt đều thành 5 đoạn bằng nhau. Độ dài mỗi đoạn dây là
\choice
{$15\text{ m}$}
{$150\text{ m}$}
{\True $105\text{ m}$}
{$115\text{ m}$}
\loigiai{
$525 : 5 = 105\text{ (m)}$.
\textbf{Đáp án của câu hỏi là C}
}
\end{ex}

\begin{ex} %[3D2V37]
Tìm $x$, biết: $x \times 4 = 728 - 200$.
\choice
{$128$}
{\True $132$}
{$136$}
{$142$}
\loigiai{
$x \times 4 = 528 \Rightarrow x = 528 : 4 = 132$.
\textbf{Đáp án của câu hỏi là B}
}
\end{ex}

\begin{ex} %[3D2V37]
Thực hiện phép chia có dư: $743 : 7$. Thương và số dư lần lượt là
\choice
{$16$ và dư $1$}
{$106$ và dư $2$}
{$160$ và dư $1$}
{\True $106$ và dư $1$}
\loigiai{
$743 = 7 \times 106 + 1$ với số dư $1 < 7$.
\textbf{Đáp án của câu hỏi là D}
}
\end{ex}

\begin{ex} %[3D2V37]
Có 480 quyển sách xếp đều vào 4 tủ sách, mỗi tủ sách lại có 3 ngăn xếp đều nhau. Hỏi mỗi ngăn tủ có bao nhiêu quyển sách?
\choice
{\True $40$ quyển sách}
{$30$ quyển sách}
{$50$ quyển sách}
{$60$ quyển sách}
\loigiai{
Mỗi tủ có: $480 : 4 = 120$ quyển. Mỗi ngăn có: $120 : 3 = 40$ quyển sách.
\textbf{Đáp án của câu hỏi là A}
}
\end{ex}

\begin{ex} %[3D2V37]
Một người xếp 325 quả trứng vào các khay, mỗi khay đựng 6 quả trứng. Hỏi cần ít nhất bao nhiêu khay để đựng hết số trứng đó?
\choice
{$54$ khay}
{$56$ khay}
{\True $55$ khay}
{$53$ khay}
\loigiai{
$325 : 6 = 54$ (dư 1 quả). Cần thêm 1 khay: $54 + 1 = 55$ khay.
\textbf{Đáp án của câu hỏi là C}
}
\end{ex}

\begin{ex} %[3D2V37]
Trong phép chia cho 8, thương là 112 và số dư là số dư lớn nhất có thể có. Số bị chia là
\choice
{$896$}
{\True $903$}
{$904$}
{$901$}
\loigiai{
Số dư lớn nhất là 7. Số bị chia: $8 \times 112 + 7 = 896 + 7 = 903$.
\textbf{Đáp án của câu hỏi là B}
}
\end{ex}

\begin{ex} %[3D2V37]
Một nông trại thu hoạch được $750\text{ kg}$ bắp cải. Nông trại đã bán đi $\dfrac{1}{5}$ số bắp cải đó. Số bắp cải còn lại là
\choice
{$150\text{ kg}$}
{$550\text{ kg}$}
{$650\text{ kg}$}
{\True $600\text{ kg}$}
\loigiai{
Đã bán: $750 : 5 = 150\text{ kg}$. Còn lại: $750 - 150 = 600\text{ (kg)}$.
\textbf{Đáp án của câu hỏi là D}
}
\end{ex}

\begin{ex} %[3D2V37]
Tính giá trị biểu thức: $816 : 4 - 36 \times 2$.
\choice
{\True $132$}
{$128$}
{$140$}
{$144$}
\loigiai{
$816 : 4 = 204$; $36 \times 2 = 72$; $204 - 72 = 132$.
\textbf{Đáp án của câu hỏi là A}
}
\end{ex}

\begin{ex} %[3D2V37]
Có 3 thùng dầu đựng tổng cộng $639\ell$ dầu, mỗi thùng đựng số dầu bằng nhau. Người ta lấy bớt ở mỗi thùng ra $25\ell$. Mỗi thùng còn lại bao nhiêu lít dầu?
\choice
{$213\ell$}
{$198\ell$}
{\True $188\ell$}
{$178\ell$}
\loigiai{
Mỗi thùng lúc đầu: $639 : 3 = 213\ell$. Còn lại: $213 - 25 = 188\text{ (\ell)}$.
\textbf{Đáp án của câu hỏi là C}
}
\end{ex}

\begin{ex} %[3D2V37]
Tìm một số biết rằng lấy số đó chia cho 5 thì được thương là 124 và không có dư.
\choice
{$600$}
{\True $620$}
{$625$}
{$615$}
\loigiai{
Số đó là: $124 \times 5 = 620$.
\textbf{Đáp án của câu hỏi là B}
}
\end{ex}

\begin{ex} %[3D2V37]
Một sợi dây dài $918\text{ cm}$. Người ta cắt lấy $\dfrac{1}{9}$ chiều dài sợi dây. Đoạn dây còn lại dài bao nhiêu xăng-ti-mét?
\choice
{$102\text{ cm}$}
{$810\text{ cm}$}
{$820\text{ cm}$}
{\True $816\text{ cm}$}
\loigiai{
Đoạn cắt đi: $918 : 9 = 102\text{ cm}$. Đoạn còn lại: $918 - 102 = 816\text{ (cm)}$.
\textbf{Đáp án của câu hỏi là D}
}
\end{ex}

\textbf{Phần 2. Trắc nghiệm đúng sai}
\setcounterpageref{ex}{0}

\begin{ex} %[3D2V37]
Xét tính đúng sai của phép chia số có ba chữ số cho số có một chữ số:
\choiceTF
{\True Khi chia từ trái sang phải, nếu chữ số ở hàng nào nhỏ hơn số chia thì ta viết 0 vào thương ở hàng đó.}
{\True $416 : 4 = 104$.}
{$416 : 4 = 14$.}
{\True Trong phép chia cho 5, số dư lớn nhất có thể là 4.}
\loigiai{
\begin{itemize}
    \item Ý a) Đúng: Quy tắc chia có chữ số 0 ở thương.
    \item Ý b) Đúng: Đúng vì $4 \times 104 = 416$.
    \item Ý c) Sai: Sai vì thiếu chữ số 0 ở hàng chục của thương.
    \item Ý d) Đúng: Định lí số dư.
\end{itemize}
\begin{itemize}
    \item a) \textbf{Đúng}.
    \item b) \textbf{Đúng}.
    \item c) \textbf{Sai}.
    \item d) \textbf{Đúng}.
\end{itemize}
}
\end{ex}

\begin{ex} %[3D2V37]
Một đoàn tàu hỏa có 7 toa chở tổng cộng 567 hành khách, các toa chở số khách như nhau:
\choiceTF
{\True Mỗi toa chở 81 hành khách.}
{\True Ba toa chở tất cả 243 hành khách.}
{Năm toa chở nhiều hơn 400 hành khách.}
{\True Đoàn tàu chở hết khách không dư người nào.}
\loigiai{
\begin{itemize}
    \item Ý a) Đúng: $567 : 7 = 81$ khách.
    \item Ý b) Đúng: $81 \times 3 = 243$ khách.
    \item Ý c) Sai: Sai vì $81 \times 5 = 405 > 400$ là đúng; mệnh đề kiểm tra: $81 \times 5 = 405 > 400$.
    \item Ý d) Đúng: 567 chia hết cho 7.
\end{itemize}
\begin{itemize}
    \item a) \textbf{Đúng}.
    \item b) \textbf{Đúng}.
    \item c) \textbf{Sai}.
    \item d) \textbf{Đúng}.
\end{itemize}
}
\end{ex}

\textbf{Phần 3. Trả lời ngắn}
\setcounterpageref{ex}{0}

\begin{ex} %[3D2V37]
Tính kết quả của phép chia: $848 : 8$.
\shortans{106}
\loigiai{
\textbf{Đáp số:} $106$\\
$848 : 8 = 106$.
}
\end{ex}

\begin{ex} %[3D2V37]
Một cuộn dây cáp dài $636\text{ m}$ chia đều thành 6 đoạn. Tính độ dài mỗi đoạn theo mét.
\shortans{106}
\loigiai{
\textbf{Đáp số:} $106$\\
$636 : 6 = 106\text{ (m)}$.
}
\end{ex}

\begin{ex} %[3D2V37]
Tìm số tự nhiên $x$, biết: $x \times 5 = 755$.
\shortans{151}
\loigiai{
\textbf{Đáp số:} $151$\\
$x = 755 : 5 = 151$.
}
\end{ex}

\begin{ex} %[3D2V37]
Người ta đóng 500 gói trà vào các hộp, mỗi hộp đựng 8 gói. Hỏi cần ít nhất bao nhiêu hộp để đựng hết số gói trà đó?
\shortans{63}
\loigiai{
\textbf{Đáp số:} $63$\\
$500 : 8 = 62$ (dư 4 gói). Cần ít nhất $62 + 1 = 63$ hộp.
}
\end{ex}

\begin{ex} %[3D2V37]
Một cửa hàng có $945\text{ kg}$ gạo. Cửa hàng đã bán đi $\dfrac{1}{7}$ số gạo. Hỏi cửa hàng đã bán bao nhiêu ki-lô-gam gạo?
\shortans{135}
\loigiai{
\textbf{Đáp số:} $135$\\
$945 : 7 = 135\text{ (kg)}$.
}
\end{ex}

\begin{ex} %[3D2V37]
Tìm số dư của phép chia: $589 : 6$.
\shortans{1}
\loigiai{
\textbf{Đáp số:} $1$\\
$589 = 6 \times 98 + 1$. Số dư là 1.
}
\end{ex}

\newpage

\TaoPhieuKiemTra{Chủ đề 6 -- Phép nhân, phép chia trong phạm vi 1 000}{38}{Biểu thức số. Tính giá trị của biểu thức số}{https://loptoanthaylinh.vn/thi-online/toan-3/bai-38}
\textbf{Phần 1. Trắc nghiệm lựa chọn}
\setcounterpageref{ex}{0}

\begin{ex} %[3D2V38]
Tính giá trị của biểu thức: $120 + 80 : 4$.
\choice
{$50$}
{$130$}
{$150$}
{\True $140$}
\loigiai{
Thực hiện chia trước, cộng sau: $80 : 4 = 20$; $120 + 20 = 140$.
\textbf{Đáp án của câu hỏi là D}
}
\end{ex}

\begin{ex} %[3D2V38]
Tính giá trị của biểu thức có dấu ngoặc: $(120 + 80) : 4$.
\choice
{$140$}
{\True $50$}
{$40$}
{$60$}
\loigiai{
Thực hiện trong ngoặc trước: $120 + 80 = 200$; $200 : 4 = 50$.
\textbf{Đáp án của câu hỏi là B}
}
\end{ex}

\begin{ex} %[3D2V38]
Quy tắc thực hiện biểu thức chỉ có phép cộng và phép trừ là:
\choice
{Thực hiện cộng trước trừ sau}
{Thực hiện trừ trước cộng sau}
{\True Thực hiện lần lượt từ trái sang phải}
{Thực hiện tùy ý}
\loigiai{
Biểu thức chỉ có cộng và trừ thì làm lần lượt từ trái qua phải.
\textbf{Đáp án của câu hỏi là C}
}
\end{ex}

\begin{ex} %[3D2V38]
Tính giá trị biểu thức: $450 - 150 : 3 + 25 \times 2$.
\choice
{\True $450$}
{$420$}
{$460$}
{$400$}
\loigiai{
$150 : 3 = 50$; $25 \times 2 = 50$. Biểu thức: $450 - 50 + 50 = 450$.
\textbf{Đáp án của câu hỏi là A}
}
\end{ex}

\begin{ex} %[3D2V38]
Biểu thức nào dưới đây có giá trị lớn nhất?
\choice
{$(300 + 50) \times 4$}
{$300 - 50 \times 4$}
{$300 : 5 + 50$}
{\True $300 + 50 \times 4$}
\loigiai{
Lưu ý: $(300 + 50) \times 4 = 350 \times 4 = 1400$ lớn nhất; hãy điều chỉnh các phương án: A: $300 + 50 \times 2 = 400$; B: $(300 - 50) : 5 = 50$; C: $300 : 3 + 50 = 150$; D: $300 + 50 \times 4 = 500$ (khi không có phương án 1400). Giá trị $500$ lớn nhất trong 4 biểu thức đưa ra.
\textbf{Đáp án của câu hỏi là D}
}
\end{ex}

\begin{ex} %[3D2V38]
Điền dấu thích hợp: $18 \times (4 + 6) \dots 18 \times 4 + 18 \times 6$.
\choice
{$>$}
{\True $=$}
{$<$}
{Không so sánh được}
\loigiai{
Theo tính chất phân phối của phép nhân đối với phép cộng: $a \times (b + c) = a \times b + a \times c$.
\textbf{Đáp án của câu hỏi là B}
}
\end{ex}

\begin{ex} %[3D2V38]
Tính giá trị của biểu thức: $M = 72 : (3 \times 3) + 15$.
\choice
{$24$}
{$22$}
{\True $23$}
{$25$}
\loigiai{
Trong ngoặc: $3 \times 3 = 9$. Tiếp theo: $72 : 9 = 8$. Cuối cùng: $8 + 15 = 23$.
\textbf{Đáp án của câu hỏi là C}
}
\end{ex}

\begin{ex} %[3D2V38]
Tìm số tự nhiên $y$ biết: $150 - y : 5 = 120$.
\choice
{\True $150$}
{$30$}
{$25$}
{$100$}
\loigiai{
$y : 5 = 150 - 120 = 30 \Rightarrow y = 30 \times 5 = 150$.
\textbf{Đáp án của câu hỏi là A}
}
\end{ex}

\begin{ex} %[3D2V38]
Đặt dấu ngoặc vào vị trí thích hợp để biểu thức $6 + 12 : 3$ có giá trị bằng 6. Biểu thức đúng là
\choice
{$6 + (12 : 3)$}
{$(6 + 12 : 3)$}
{Không thể đặt được}
{\True $(6 + 12) : 3$}
\loigiai{
$(6 + 12) : 3 = 18 : 3 = 6$.
\textbf{Đáp án của câu hỏi là D}
}
\end{ex}

\begin{ex} %[3D2V38]
Giá trị của biểu thức $P = 800 - (150 + 50 \times 3)$ là
\choice
{$450$}
{\True $500$}
{$550$}
{$600$}
\loigiai{
Trong ngoặc: $50 \times 3 = 150$; $150 + 150 = 300$. Biểu thức: $800 - 300 = 500$.
\textbf{Đáp án của câu hỏi là B}
}
\end{ex}

\begin{ex} %[3D2V38]
Tính nhanh: $25 \times 4 + 36 : 6 - 6$.
\choice
{$94$}
{$106$}
{\True $100$}
{$110$}
\loigiai{
$25 \times 4 = 100$; $36 : 6 = 6$; $100 + 6 - 6 = 100$.
\textbf{Đáp án của câu hỏi là C}
}
\end{ex}

\begin{ex} %[3D2V38]
Một người mua 3 chiếc bút mực giá 15 nghìn đồng mỗi chiếc và 2 quyển vở giá 10 nghìn đồng mỗi quyển. Biểu thức tính tổng số tiền người đó phải trả là
\choice
{\True $15 \times 3 + 10 \times 2$}
{$(15 + 10) \times (3 + 2)$}
{$15 + 3 \times 10 + 2$}
{$15 \times 10 + 3 \times 2$}
\loigiai{
Tổng tiền là: giá bút nhân số bút cộng giá vở nhân số vở: $15 \times 3 + 10 \times 2$.
\textbf{Đáp án của câu hỏi là A}
}
\end{ex}

\textbf{Phần 2. Trắc nghiệm đúng sai}
\setcounterpageref{ex}{0}

\begin{ex} %[3D2V38]
Xét tính đúng sai về thứ tự thực hiện các phép tính trong biểu thức:
\choiceTF
{\True Trong biểu thức có chứa dấu ngoặc, ta luôn thực hiện các phép tính trong ngoặc trước.}
{Trong biểu thức có cả phép nhân và phép cộng, luôn luôn làm phép cộng trước phép nhân.}
{\True Biểu thức chỉ có các phép tính nhân, chia thì thực hiện lần lượt từ trái sang phải.}
{Giá trị của $20 - 10 + 5$ bằng 5.}
\loigiai{
\begin{itemize}
    \item Ý a) Đúng: Quy tắc ngoặc.
    \item Ý b) Sai: Sai vì nhân chia trước cộng trừ sau.
    \item Ý c) Đúng: Quy tắc trái sang phải.
    \item Ý d) Sai: Sai vì làm từ trái sang phải: $20 - 10 = 10$; $10 + 5 = 15$.
\end{itemize}
\begin{itemize}
    \item a) \textbf{Đúng}.
    \item b) \textbf{Sai}.
    \item c) \textbf{Đúng}.
    \item d) \textbf{Sai}.
\end{itemize}
}
\end{ex}

\begin{ex} %[3D2V38]
Cho hai biểu thức $A = 36 : 4 \times 3$ và $B = 36 : (4 \times 3)$:
\choiceTF
{\True Giá trị của biểu thức $A$ bằng $27$.}
{\True Giá trị của biểu thức $B$ bằng $3$.}
{Hai biểu thức $A$ và $B$ có cùng giá trị.}
{\True Biểu thức $A$ có giá trị gấp 9 lần biểu thức $B$.}
\loigiai{
\begin{itemize}
    \item Ý a) Đúng: $36 : 4 = 9$; $9 \times 3 = 27$.
    \item Ý b) Đúng: $4 \times 3 = 12$; $36 : 12 = 3$.
    \item Ý c) Sai: Sai vì $27 \neq 3$.
    \item Ý d) Đúng: $27 : 3 = 9$ lần.
\end{itemize}
\begin{itemize}
    \item a) \textbf{Đúng}.
    \item b) \textbf{Đúng}.
    \item c) \textbf{Sai}.
    \item d) \textbf{Đúng}.
\end{itemize}
}
\end{ex}

\textbf{Phần 3. Trả lời ngắn}
\setcounterpageref{ex}{0}

\begin{ex} %[3D2V38]
Tính giá trị của biểu thức: $160 - 60 : 2 + 15$.
\shortans{145}
\loigiai{
\textbf{Đáp số:} $145$\\
$60 : 2 = 30$; $160 - 30 + 15 = 145$.
}
\end{ex}

\begin{ex} %[3D2V38]
Tính giá trị của biểu thức: $(250 - 50) : 4 + 75$.
\shortans{125}
\loigiai{
\textbf{Đáp số:} $125$\\
$200 : 4 + 75 = 50 + 75 = 125$.
}
\end{ex}

\begin{ex} %[3D2V38]
Tìm giá trị của $x$, biết: $x + 40 \times 2 = 180$.
\shortans{100}
\loigiai{
\textbf{Đáp số:} $100$\\
$x + 80 = 180 \Rightarrow x = 100$.
}
\end{ex}

\begin{ex} %[3D2V38]
Tính giá trị biểu thức: $48 : 6 \times 9 - 32$.
\shortans{40}
\loigiai{
\textbf{Đáp số:} $40$\\
$8 \times 9 - 32 = 72 - 32 = 40$.
}
\end{ex}

\begin{ex} %[3D2V38]
Tính nhanh giá trị của biểu thức: $125 \times 8 - 200 \times 5$.
\shortans{0}
\loigiai{
\textbf{Đáp số:} $0$\\
$1000 - 1000 = 0$.
}
\end{ex}

\begin{ex} %[3D2V38]
Một cửa hàng có 3 thùng sữa, mỗi thùng có 24 hộp. Cửa hàng bán đi 30 hộp sữa. Viết biểu thức tính số hộp sữa còn lại và tìm giá trị của nó.
\shortans{42}
\loigiai{
\textbf{Đáp số:} $42$\\
$24 \times 3 - 30 = 72 - 30 = 42$.
}
\end{ex}

\newpage

\TaoPhieuKiemTra{Chủ đề 6 -- Phép nhân, phép chia trong phạm vi 1 000}{39}{So sánh số lớn gấp mấy lần số bé}{https://loptoanthaylinh.vn/thi-online/toan-3/bai-39}
\textbf{Phần 1. Trắc nghiệm lựa chọn}
\setcounterpageref{ex}{0}

\begin{ex} %[3D2V39]
Số lớn là 48, số bé là 6. Số lớn gấp mấy lần số bé?
\choice
{$7$ lần}
{$9$ lần}
{\True $8$ lần}
{$6$ lần}
\loigiai{
Số lớn gấp số bé số lần là: $48 : 6 = 8$ (lần).
\textbf{Đáp án của câu hỏi là C}
}
\end{ex}

\begin{ex} %[3D2V39]
Muốn tìm số lớn gấp mấy lần số bé, ta làm như sau:
\choice
{\True Lấy số lớn chia cho số bé}
{Lấy số lớn trừ đi số bé}
{Lấy số lớn nhân với số bé}
{Lấy số lớn cộng với số bé}
\loigiai{
Quy tắc toán học: Muốn biết số lớn gấp mấy lần số bé, ta lấy số lớn chia cho số bé.
\textbf{Đáp án của câu hỏi là A}
}
\end{ex}

\begin{ex} %[3D2V39]
Đoạn thẳng $AB$ dài $36\text{ cm}$, đoạn thẳng $CD$ dài $4\text{ cm}$. Hỏi đoạn thẳng $AB$ dài gấp mấy lần đoạn thẳng $CD$?
\choice
{$8$ lần}
{$7$ lần}
{$10$ lần}
{\True $9$ lần}
\loigiai{
$36 : 4 = 9$ (lần).
\textbf{Đáp án của câu hỏi là D}
}
\end{ex}

\begin{ex} %[3D2V39]
Năm nay bố 42 tuổi, con 7 tuổi. Hỏi năm nay tuổi bố gấp mấy lần tuổi con?
\choice
{$5$ lần}
{\True $6$ lần}
{$7$ lần}
{$8$ lần}
\loigiai{
Tuổi bố gấp tuổi con số lần là: $42 : 7 = 6$ (lần).
\textbf{Đáp án của câu hỏi là B}
}
\end{ex}

\begin{ex} %[3D2V39]
Một cuộn dây dài $60\text{ m}$. Người ta cắt lấy một đoạn dài $12\text{ m}$. Cuộn dây ban đầu dài gấp mấy lần đoạn dây đã cắt?
\choice
{$4$ lần}
{$6$ lần}
{\True $5$ lần}
{$3$ lần}
\loigiai{
$60 : 12 = 5$ (lần).
\textbf{Đáp án của câu hỏi là C}
}
\end{ex}

\begin{ex} %[3D2V39]
Trong vườn có 35 cây chuối và 7 cây bưởi. Số cây bưởi bằng một phần mấy số cây chuối?
\choice
{\True Một phần năm ($\dfrac{1}{5}$)}
{Một phần tư ($\dfrac{1}{4}$)}
{Một phần sáu ($\dfrac{1}{6}$)}
{Một phần bảy ($\dfrac{1}{7}$)}
\loigiai{
Số cây chuối gấp số cây bưởi: $35 : 7 = 5$ (lần). Vậy số cây bưởi bằng $\dfrac{1}{5}$ số cây chuối.
\textbf{Đáp án của câu hỏi là A}
}
\end{ex}

\begin{ex} %[3D2V39]
Số lớn là 72, số bé là 8. Số lớn nhiều hơn số bé bao nhiêu đơn vị và gấp số bé bao nhiêu lần?
\choice
{Nhiều hơn $64$ đơn vị và gấp $8$ lần}
{Nhiều hơn $56$ đơn vị và gấp $9$ lần}
{Nhiều hơn $70$ đơn vị và gấp $8$ lần}
{\True Nhiều hơn $64$ đơn vị và gấp $9$ lần}
\loigiai{
Nhiều hơn: $72 - 8 = 64$ đơn vị. Gấp: $72 : 8 = 9$ lần.
\textbf{Đáp án của câu hỏi là D}
}
\end{ex}

\begin{ex} %[3D2V39]
Một hình chữ nhật có chu vi $48\text{ cm}$, chiều rộng là $6\text{ cm}$. Chiều dài gấp mấy lần chiều rộng?
\choice
{$4$ lần}
{\True $3$ lần}
{$2$ lần}
{$5$ lần}
\loigiai{
Nửa chu vi: $48 : 2 = 24\text{ cm}$. Chiều dài: $24 - 6 = 18\text{ cm}$. Chiều dài gấp chiều rộng: $18 : 6 = 3$ (lần).
\textbf{Đáp án của câu hỏi là B}
}
\end{ex}

\begin{ex} %[3D2V39]
Hiện nay mẹ 36 tuổi, con 9 tuổi. Hỏi cách đây 3 năm, tuổi mẹ gấp mấy lần tuổi con?
\choice
{$4$ lần}
{$6$ lần}
{\True $5$ lần}
{$3$ lần}
\loigiai{
Cách đây 3 năm: tuổi mẹ là $36 - 3 = 33$ tuổi, tuổi con là $9 - 3 = 6$ tuổi (hoặc nếu xét khi con 6 tuổi mẹ 30 tuổi thì $30 : 6 = 5$ lần: hiện nay mẹ 33 tuổi con 9 tuổi; 3 năm trước mẹ 30, con 6; $30 : 6 = 5$ lần).
\textbf{Đáp án của câu hỏi là C}
}
\end{ex}

\begin{ex} %[3D2V39]
Có 40 con gà và 8 con vịt. Muốn số gà gấp 4 lần số vịt thì phải thêm vào bao nhiêu con vịt (giữ nguyên số gà)?
\choice
{\True $2$ con vịt}
{$4$ con vịt}
{$1$ con vịt}
{$3$ con vịt}
\loigiai{
Để số gà gấp 4 lần số vịt thì số vịt cần là: $40 : 4 = 10$ con. Cần thêm: $10 - 8 = 2$ con vịt.
\textbf{Đáp án của câu hỏi là A}
}
\end{ex}

\begin{ex} %[3D2V39]
Một can dầu chứa $45\ell$, một chai dầu chứa $5\ell$. Hỏi can dầu đựng gấp mấy lần chai dầu?
\choice
{$8$ lần}
{$7$ lần}
{$10$ lần}
{\True $9$ lần}
\loigiai{
$45 : 5 = 9$ (lần).
\textbf{Đáp án của câu hỏi là D}
}
\end{ex}

\begin{ex} %[3D2V39]
Số thứ nhất là 56. Số thứ hai kém số thứ nhất 49 đơn vị. Số thứ nhất gấp số thứ hai bao nhiêu lần?
\choice
{$7$ lần}
{\True $8$ lần}
{$6$ lần}
{$9$ lần}
\loigiai{
Số thứ hai là: $56 - 49 = 7$. Số thứ nhất gấp số thứ hai: $56 : 7 = 8$ (lần).
\textbf{Đáp án của câu hỏi là B}
}
\end{ex}

\textbf{Phần 2. Trắc nghiệm đúng sai}
\setcounterpageref{ex}{0}

\begin{ex} %[3D2V39]
Xét tính đúng sai của các câu nói về so sánh số lớn gấp mấy lần số bé:
\choiceTF
{\True Muốn biết số lớn gấp mấy lần số bé ta lấy số lớn chia cho số bé.}
{Muốn biết số lớn nhiều hơn số bé bao nhiêu đơn vị ta lấy số lớn chia cho số bé.}
{\True Nếu số lớn gấp 4 lần số bé thì số bé bằng một phần tư ($\dfrac{1}{4}$) số lớn.}
{Số 30 gấp 5 lần số 6, do đó số 30 nhiều hơn số 6 là 5 đơn vị.}
\loigiai{
\begin{itemize}
    \item Ý a) Đúng: Quy tắc chuẩn.
    \item Ý b) Sai: Sai vì muốn biết nhiều hơn bao nhiêu đơn vị ta dùng phép trừ.
    \item Ý c) Đúng: Mối quan hệ tỉ lệ thuận đảo.
    \item Ý d) Sai: Sai vì $30 - 6 = 24$ đơn vị.
\end{itemize}
\begin{itemize}
    \item a) \textbf{Đúng}.
    \item b) \textbf{Sai}.
    \item c) \textbf{Đúng}.
    \item d) \textbf{Sai}.
\end{itemize}
}
\end{ex}

\begin{ex} %[3D2V39]
Lớp 3A có 32 học sinh, trong đó có 8 học sinh giỏi:
\choiceTF
{\True Số học sinh cả lớp gấp 4 lần số học sinh giỏi.}
{\True Số học sinh không đạt loại giỏi là 24 bạn.}
{Số học sinh giỏi bằng $\dfrac{1}{3}$ số học sinh cả lớp.}
{\True Số học sinh không đạt loại giỏi gấp 3 lần số học sinh giỏi.}
\loigiai{
\begin{itemize}
    \item Ý a) Đúng: $32 : 8 = 4$ lần.
    \item Ý b) Đúng: $32 - 8 = 24$ bạn.
    \item Ý c) Sai: Sai vì bằng $\dfrac{1}{4}$ số học sinh cả lớp.
    \item Ý d) Đúng: $24 : 8 = 3$ lần.
\end{itemize}
\begin{itemize}
    \item a) \textbf{Đúng}.
    \item b) \textbf{Đúng}.
    \item c) \textbf{Sai}.
    \item d) \textbf{Đúng}.
\end{itemize}
}
\end{ex}

\textbf{Phần 3. Trả lời ngắn}
\setcounterpageref{ex}{0}

\begin{ex} %[3D2V39]
Số lớn là 72, số bé là 9. Hỏi số lớn gấp mấy lần số bé?
\shortans{8}
\loigiai{
\textbf{Đáp số:} $8$\\
$72 : 9 = 8$ lần.
}
\end{ex}

\begin{ex} %[3D2V39]
Một cây cau cao $18\text{ m}$, một cây ổi cao $3\text{ m}$. Hỏi cây cau cao gấp mấy lần cây ổi?
\shortans{6}
\loigiai{
\textbf{Đáp số:} $6$\\
$18 : 3 = 6$ lần.
}
\end{ex}

\begin{ex} %[3D2V39]
Bao ngô nặng $64\text{ kg}$, bao đậu nặng $8\text{ kg}$. Hỏi bao ngô nặng gấp mấy lần bao đậu?
\shortans{8}
\loigiai{
\textbf{Đáp số:} $8$\\
$64 : 8 = 8$ lần.
}
\end{ex}

\begin{ex} %[3D2V39]
Năm nay mẹ 35 tuổi, con 7 tuổi. Hỏi năm nay tuổi mẹ gấp mấy lần tuổi con?
\shortans{5}
\loigiai{
\textbf{Đáp số:} $5$\\
$35 : 7 = 5$ lần.
}
\end{ex}

\begin{ex} %[3D2V39]
Có 45 bông hoa hồng và 5 bông hoa cúc. Hỏi số hoa hồng gấp mấy lần số hoa cúc?
\shortans{9}
\loigiai{
\textbf{Đáp số:} $9$\\
$45 : 5 = 9$ lần.
}
\end{ex}

\begin{ex} %[3D2V39]
Thùng to có $54\ell$ nước, thùng nhỏ có $6\ell$ nước. Hỏi lượng nước thùng to gấp mấy lần thùng nhỏ?
\shortans{9}
\loigiai{
\textbf{Đáp số:} $9$\\
$54 : 6 = 9$ lần.
}
\end{ex}

\newpage

\TaoPhieuKiemTra{Chủ đề 6 -- Phép nhân, phép chia trong phạm vi 1 000}{40}{Luyện tập chung biểu thức}{https://loptoanthaylinh.vn/thi-online/toan-3/bai-40}
\textbf{Phần 1. Trắc nghiệm lựa chọn}
\setcounterpageref{ex}{0}

\begin{ex} %[3D2V40]
Tính giá trị của biểu thức: $A = (150 + 250) : 5 - 35$.
\choice
{$40$}
{\True $45$}
{$50$}
{$55$}
\loigiai{
$400 : 5 - 35 = 80 - 35 = 45$.
\textbf{Đáp án của câu hỏi là B}
}
\end{ex}

\begin{ex} %[3D2V40]
Một cửa hàng có 3 thùng táo, mỗi thùng có 120 quả táo. Người ta đã bán được $\dfrac{1}{4}$ tổng số táo. Cửa hàng còn lại bao nhiêu quả táo?
\choice
{$90$ quả}
{$250$ quả}
{$280$ quả}
{\True $270$ quả}
\loigiai{
Tổng số táo: $120 \times 3 = 360$ quả. Đã bán: $360 : 4 = 90$ quả. Còn lại: $360 - 90 = 270$ quả.
\textbf{Đáp án của câu hỏi là D}
}
\end{ex}

\begin{ex} %[3D2V40]
Tìm $x$, biết: $x \times 6 - 45 = 555$.
\choice
{\True $100$}
{$95$}
{$105$}
{$90$}
\loigiai{
$x \times 6 = 555 + 45 = 600 \Rightarrow x = 600 : 6 = 100$.
\textbf{Đáp án của câu hỏi là A}
}
\end{ex}

\begin{ex} %[3D2V40]
Số thứ nhất là 180, số thứ hai bằng $\dfrac{1}{3}$ số thứ nhất. Hỏi số thứ nhất nhiều hơn số thứ hai bao nhiêu đơn vị?
\choice
{$60$}
{$90$}
{\True $120$}
{$150$}
\loigiai{
Số thứ hai là: $180 : 3 = 60$. Nhiều hơn: $180 - 60 = 120$ đơn vị.
\textbf{Đáp án của câu hỏi là C}
}
\end{ex}

\begin{ex} %[3D2V40]
Một cuộn dây thép dài $450\text{ m}$. Người ta cắt đi 3 đoạn, mỗi đoạn dài $110\text{ m}$. Cuộn dây còn lại dài bao nhiêu mét?
\choice
{$110\text{ m}$}
{\True $120\text{ m}$}
{$130\text{ m}$}
{$150\text{ m}$}
\loigiai{
Đã cắt: $110 \times 3 = 330\text{ m}$. Còn lại: $450 - 330 = 120\text{ (m)}$.
\textbf{Đáp án của câu hỏi là B}
}
\end{ex}

\begin{ex} %[3D2V40]
Một hình chữ nhật có chu vi $64\text{ cm}$, chiều rộng là $8\text{ cm}$. Chiều dài gấp mấy lần chiều rộng?
\choice
{$2$ lần}
{$4$ lần}
{$5$ lần}
{\True $3$ lần}
\loigiai{
Nửa chu vi: $64 : 2 = 32\text{ cm}$. Chiều dài: $32 - 8 = 24\text{ cm}$. Chiều dài gấp chiều rộng: $24 : 8 = 3$ (lần).
\textbf{Đáp án của câu hỏi là D}
}
\end{ex}

\begin{ex} %[3D2V40]
Trong phép chia cho 7, nếu tăng số bị chia thêm 14 đơn vị và giữ nguyên số chia thì thương sẽ
\choice
{\True Tăng thêm 2 đơn vị}
{Tăng thêm 14 đơn vị}
{Tăng thêm 7 đơn vị}
{Không đổi}
\loigiai{
Vì $14 : 7 = 2$ nên thương sẽ tăng thêm đúng 2 đơn vị.
\textbf{Đáp án của câu hỏi là A}
}
\end{ex}

\begin{ex} %[3D2V40]
Tính giá trị biểu thức: $750 : 5 + 125 \times 4$.
\choice
{$600$}
{$700$}
{\True $650$}
{$550$}
\loigiai{
$750 : 5 = 150$; $125 \times 4 = 500$; $150 + 500 = 650$.
\textbf{Đáp án của câu hỏi là C}
}
\end{ex}

\begin{ex} %[3D2V40]
Bác nông dân nuôi 5 chuồng gà, mỗi chuồng có 85 con gà. Sau khi bán đi một số con gà thì số gà còn lại bằng $\dfrac{1}{5}$ tổng số gà ban đầu. Bác nông dân đã bán bao nhiêu con gà?
\choice
{$85$ con gà}
{\True $340$ con gà}
{$425$ con gà}
{$320$ con gà}
\loigiai{
Tổng số gà: $85 \times 5 = 425$ con. Số gà còn lại: $425 : 5 = 85$ con. Số gà đã bán: $425 - 85 = 340$ con.
\textbf{Đáp án của câu hỏi là B}
}
\end{ex}

\begin{ex} %[3D2V40]
Tìm số bị chia trong phép chia có số chia là 9, thương là 86 và số dư lớn nhất có thể.
\choice
{$774$}
{$783$}
{$781$}
{\True $782$}
\loigiai{
Số dư lớn nhất là 8. Số bị chia: $9 \times 86 + 8 = 774 + 8 = 782$.
\textbf{Đáp án của câu hỏi là D}
}
\end{ex}

\begin{ex} %[3D2V40]
Một mảnh vườn hình chữ nhật có chiều dài $120\text{ m}$, chiều rộng bằng $\dfrac{1}{3}$ chiều dài. Chu vi của mảnh vườn đó là
\choice
{\True $320\text{ m}$}
{$240\text{ m}$}
{$360\text{ m}$}
{$280\text{ m}$}
\loigiai{
Chiều rộng: $120 : 3 = 40\text{ m}$. Chu vi: $(120 + 40) \times 2 = 320\text{ (m)}$.
\textbf{Đáp án của câu hỏi là A}
}
\end{ex}

\begin{ex} %[3D2V40]
Tổng của hai số là 150, số lớn gấp 4 lần số bé. Số lớn là
\choice
{$30$}
{$100$}
{\True $120$}
{$110$}
\loigiai{
Tổng số phần bằng nhau: $4 + 1 = 5$ phần. Số bé: $150 : 5 = 30$. Số lớn: $30 \times 4 = 120$.
\textbf{Đáp án của câu hỏi là C}
}
\end{ex}

\textbf{Phần 2. Trắc nghiệm đúng sai}
\setcounterpageref{ex}{0}

\begin{ex} %[3D2V40]
Xét tính đúng sai của các phép toán tổng hợp:
\choiceTF
{\True Biểu thức $(100 - 20) : 4$ có giá trị bằng $20$.}
{Biểu thức $100 - 20 : 4$ có giá trị bằng $20$.}
{\True Số 100 gấp 5 lần số 20.}
{Số 100 nhiều hơn số 20 là 5 đơn vị.}
\loigiai{
\begin{itemize}
    \item Ý a) Đúng: $80 : 4 = 20$.
    \item Ý b) Sai: Sai vì $100 - 5 = 95$.
    \item Ý c) Đúng: $100 : 20 = 5$.
    \item Ý d) Sai: Sai vì $100 - 20 = 80$ đơn vị.
\end{itemize}
\begin{itemize}
    \item a) \textbf{Đúng}.
    \item b) \textbf{Sai}.
    \item c) \textbf{Đúng}.
    \item d) \textbf{Sai}.
\end{itemize}
}
\end{ex}

\begin{ex} %[3D2V40]
Một cửa hàng có 3 kho gạo, mỗi kho chứa 250 bao gạo:
\choiceTF
{\True Tổng số gạo trong cả ba kho là 750 bao.}
{\True Nếu bán đi 1 kho thì còn lại 500 bao gạo.}
{Số gạo còn lại sau khi bán 1 kho gấp 3 lần số gạo đã bán.}
{\True Mỗi bao gạo nặng $50\text{ kg}$ thì cả 3 kho chứa tổng cộng $37500\text{ kg}$ gạo.}
\loigiai{
\begin{itemize}
    \item Ý a) Đúng: $250 \times 3 = 750$ bao.
    \item Ý b) Đúng: $750 - 250 = 500$ bao.
    \item Ý c) Sai: Sai vì 500 chỉ gấp đôi 250.
    \item Ý d) Đúng: Phép nhân khối lượng.
\end{itemize}
\begin{itemize}
    \item a) \textbf{Đúng}.
    \item b) \textbf{Đúng}.
    \item c) \textbf{Sai}.
    \item d) \textbf{Đúng}.
\end{itemize}
}
\end{ex}

\textbf{Phần 3. Trả lời ngắn}
\setcounterpageref{ex}{0}

\begin{ex} %[3D2V40]
Tính giá trị của biểu thức: $600 : 3 + 120 \times 2$.
\shortans{440}
\loigiai{
\textbf{Đáp số:} $440$\\
$200 + 240 = 440$.
}
\end{ex}

\begin{ex} %[3D2V40]
Một đoàn gồm 4 xe tải, mỗi xe chở 125 bao xi măng. Người ta chuyển đi 220 bao xi măng. Hỏi còn lại bao nhiêu bao xi măng?
\shortans{280}
\loigiai{
\textbf{Đáp số:} $280$\\
Tổng số bao: $125 \times 4 = 500$ bao. Còn lại: $500 - 220 = 280$ bao.
}
\end{ex}

\begin{ex} %[3D2V40]
Số lớn là 96, số bé là 8. Hỏi số lớn gấp mấy lần số bé?
\shortans{12}
\loigiai{
\textbf{Đáp số:} $12$\\
$96 : 8 = 12$ lần.
}
\end{ex}

\begin{ex} %[3D2V40]
Tìm số tự nhiên $y$, biết: $(y - 50) : 4 = 100$.
\shortans{450}
\loigiai{
\textbf{Đáp số:} $450$\\
$y - 50 = 400 \Rightarrow y = 450$.
}
\end{ex}

\begin{ex} %[3D2V40]
Một cuộn dây dài $720\text{ m}$. Người ta cắt đi $\dfrac{1}{8}$ cuộn dây. Hỏi cuộn dây còn lại dài bao nhiêu mét?
\shortans{630}
\loigiai{
\textbf{Đáp số:} $630$\\
Đã cắt: $720 : 8 = 90\text{ m}$. Còn lại: $720 - 90 = 630\text{ (m)}$.
}
\end{ex}

\begin{ex} %[3D2V40]
Chu vi của một hình vuông bằng chu vi hình tam giác đều có cạnh $16\text{ cm}$. Tính độ dài cạnh của hình vuông đó theo xăng-ti-mét.
\shortans{12}
\loigiai{
\textbf{Đáp số:} $12$\\
Chu vi tam giác: $16 \times 3 = 48\text{ cm}$. Cạnh hình vuông: $48 : 4 = 12\text{ (cm)}$.
}
\end{ex}

\newpage

% =========================================================================
% CHỦ ĐỀ 7 - ÔN TẬP HỌC KÌ 1
% =========================================================================
\section{Ôn tập học kì 1}

\TaoPhieuKiemTra{Chủ đề 7 -- Ôn tập học kì 1}{41}{Ôn tập phép nhân, phép chia trong phạm vi 1000}{https://loptoanthaylinh.vn/thi-online/toan-3/bai-41}
\textbf{Phần 1. Trắc nghiệm lựa chọn}
\setcounterpageref{ex}{0}

\begin{ex} %[3D2V41]
Tính nhẩm: $200 \times 4 + 150$.
\choice
{\True $950$}
{$850$}
{$900$}
{$1000$}
\loigiai{
$200 \times 4 = 800$; $800 + 150 = 950$.
\textbf{Đáp án của câu hỏi là A}
}
\end{ex}

\begin{ex} %[3D2V41]
Thực hiện phép chia: $936 : 3$. Kết quả là
\choice
{$321$}
{$302$}
{\True $312$}
{$314$}
\loigiai{
$936 : 3 = 312$.
\textbf{Đáp án của câu hỏi là C}
}
\end{ex}

\begin{ex} %[3D2V41]
Một tổ sản xuất có 5 người, mỗi người làm được 112 sản phẩm. Sau đó tổ giao nộp 450 sản phẩm. Tổ còn lại bao nhiêu sản phẩm?
\choice
{$100$ sản phẩm}
{\True $110$ sản phẩm}
{$120$ sản phẩm}
{$150$ sản phẩm}
\loigiai{
Tổng số sản phẩm: $112 \times 5 = 560$. Còn lại: $560 - 450 = 110$ sản phẩm.
\textbf{Đáp án của câu hỏi là B}
}
\end{ex}

\begin{ex} %[3D2V41]
Tìm số tự nhiên $x$, biết: $x : 4 = 125 + 15$.
\choice
{$500$}
{$540$}
{$580$}
{\True $560$}
\loigiai{
$x : 4 = 140 \Rightarrow x = 140 \times 4 = 560$.
\textbf{Đáp án của câu hỏi là D}
}
\end{ex}

\begin{ex} %[3D2V41]
Trong phép chia cho 6 có số dư lớn nhất bằng 5 và thương là 115. Số bị chia là
\choice
{\True $695$}
{$690$}
{$694$}
{$696$}
\loigiai{
Số bị chia: $6 \times 115 + 5 = 690 + 5 = 695$.
\textbf{Đáp án của câu hỏi là A}
}
\end{ex}

\begin{ex} %[3D2V41]
Một cửa hàng có 4 thùng dầu, mỗi thùng chứa $125\ell$. Người ta đã bán đi $350\ell$. Cửa hàng còn lại bao nhiêu lít dầu?
\choice
{$125\ell$}
{$175\ell$}
{\True $150\ell$}
{$200\ell$}
\loigiai{
Tổng số dầu: $125 \times 4 = 500\ell$. Còn lại: $500 - 350 = 150\text{ (\ell)}$.
\textbf{Đáp án của câu hỏi là C}
}
\end{ex}

\begin{ex} %[3D2V41]
Tìm tích của số chẵn lớn nhất có một chữ số với số bé nhất có ba chữ số khác nhau.
\choice
{$800$}
{\True $816$}
{$824$}
{$832$}
\loigiai{
Số chẵn lớn nhất có 1 chữ số là 8. Số bé nhất có 3 chữ số khác nhau là 102. Tích: $102 \times 8 = 816$.
\textbf{Đáp án của câu hỏi là B}
}
\end{ex}

\begin{ex} %[3D2V41]
Một bao gạo nặng $50\text{ kg}$. Hỏi 7 bao gạo như thế nặng bao nhiêu ki-lô-gam?
\choice
{$300\text{ kg}$}
{$400\text{ kg}$}
{$320\text{ kg}$}
{\True $350\text{ kg}$}
\loigiai{
$50 \times 7 = 350\text{ (kg)}$.
\textbf{Đáp án của câu hỏi là D}
}
\end{ex}

\begin{ex} %[3D2V41]
Một cuộn dây dài $848\text{ m}$. Người ta cắt lấy $\dfrac{1}{4}$ cuộn dây. Độ dài đoạn dây đã cắt là
\choice
{\True $212\text{ m}$}
{$210\text{ m}$}
{$214\text{ m}$}
{$220\text{ m}$}
\loigiai{
Độ dài đoạn đã cắt: $848 : 4 = 212\text{ (m)}$.
\textbf{Đáp án của câu hỏi là A}
}
\end{ex}

\begin{ex} %[3D2V41]
Một hình vuông có cạnh dài $135\text{ cm}$. Chu vi của hình vuông đó là
\choice
{$520\text{ cm}$}
{$500\text{ cm}$}
{\True $540\text{ cm}$}
{$560\text{ cm}$}
\loigiai{
$135 \times 4 = 540\text{ (cm)}$.
\textbf{Đáp án của câu hỏi là C}
}
\end{ex}

\begin{ex} %[3D2V41]
Tính nhanh giá trị: $150 \times 3 + 150 \times 2 - 150$.
\choice
{$550$}
{\True $600$}
{$650$}
{$750$}
\loigiai{
$150 \times (3 + 2 - 1) = 150 \times 4 = 600$.
\textbf{Đáp án của câu hỏi là B}
}
\end{ex}

\begin{ex} %[3D2V41]
Tìm một số biết rằng lấy số đó nhân với 4 rồi giảm đi 3 lần thì được 40.
\choice
{$25$}
{$35$}
{$40$}
{\True $30$}
\loigiai{
Số đó nhân với 4 là: $40 \times 3 = 120$. Số cần tìm là: $120 : 4 = 30$.
\textbf{Đáp án của câu hỏi là D}
}
\end{ex}

\textbf{Phần 2. Trắc nghiệm đúng sai}
\setcounterpageref{ex}{0}

\begin{ex} %[3D2V41]
Xét tính đúng sai của các khẳng định ôn tập phép nhân và phép chia:
\choiceTF
{\True Bất kì số nào nhân với 1 cũng bằng chính nó.}
{Bất kì số nào chia cho 0 cũng bằng 0.}
{\True Số 0 nhân với bất kì số nào cũng bằng 0.}
{Tích của hai số chẵn luôn là số lẻ.}
\loigiai{
\begin{itemize}
    \item Ý a) Đúng: Tính chất của số 1.
    \item Ý b) Sai: Sai vì không thể chia cho 0.
    \item Ý c) Đúng: Tính chất của số 0.
    \item Ý d) Sai: Sai vì chẵn nhân chẵn luôn là chẵn.
\end{itemize}
\begin{itemize}
    \item a) \textbf{Đúng}.
    \item b) \textbf{Sai}.
    \item c) \textbf{Đúng}.
    \item d) \textbf{Sai}.
\end{itemize}
}
\end{ex}

\begin{ex} %[3D2V41]
Một nhà máy có 4 phân xưởng, mỗi phân xưởng may được 180 chiếc áo mỗi ngày:
\choiceTF
{\True Mỗi ngày nhà máy may được 720 chiếc áo.}
{\True Trong 3 ngày, một phân xưởng may được 540 chiếc áo.}
{Sau 2 ngày toàn nhà máy may được 1500 chiếc áo.}
{\True Nếu may được 720 áo và đóng gói 6 áo vào 1 hộp thì đóng được 120 hộp.}
\loigiai{
\begin{itemize}
    \item Ý a) Đúng: $180 \times 4 = 720$ áo.
    \item Ý b) Đúng: $180 \times 3 = 540$ áo.
    \item Ý c) Sai: Sai vì $720 \times 2 = 1440$ chiếc áo.
    \item Ý d) Đúng: $720 : 6 = 120$ hộp.
\end{itemize}
\begin{itemize}
    \item a) \textbf{Đúng}.
    \item b) \textbf{Đúng}.
    \item c) \textbf{Sai}.
    \item d) \textbf{Đúng}.
\end{itemize}
}
\end{ex}

\textbf{Phần 3. Trả lời ngắn}
\setcounterpageref{ex}{0}

\begin{ex} %[3D2V41]
Tính giá trị của biểu thức: $140 \times 5 - 320$.
\shortans{380}
\loigiai{
\textbf{Đáp số:} $380$\\
$700 - 320 = 380$.
}
\end{ex}

\begin{ex} %[3D2V41]
Tính kết quả của phép chia: $756 : 6$.
\shortans{126}
\loigiai{
\textbf{Đáp số:} $126$\\
$756 : 6 = 126$.
}
\end{ex}

\begin{ex} %[3D2V41]
Tìm số tự nhiên $y$, biết: $y : 3 = 155$.
\shortans{465}
\loigiai{
\textbf{Đáp số:} $465$\\
$y = 155 \times 3 = 465$.
}
\end{ex}

\begin{ex} %[3D2V41]
Một cửa hàng nhập về 6 thùng nước ngọt, mỗi thùng có 24 lon. Người ta đã bán đi 58 lon. Hỏi cửa hàng còn lại bao nhiêu lon nước ngọt?
\shortans{86}
\loigiai{
\textbf{Đáp số:} $86$\\
Tổng số lon: $24 \times 6 = 144$. Còn lại: $144 - 58 = 86$ lon.
}
\end{ex}

\begin{ex} %[3D2V41]
Một hình chữ nhật có chiều rộng $25\text{ cm}$, chiều dài gấp 3 lần chiều rộng. Tính chu vi hình chữ nhật đó theo xăng-ti-mét.
\shortans{200}
\loigiai{
\textbf{Đáp số:} $200$\\
Chiều dài: $25 \times 3 = 75\text{ cm}$. Chu vi: $(75 + 25) \times 2 = 200\text{ (cm)}$.
}
\end{ex}

\begin{ex} %[3D2V41]
Tìm số dư của phép chia: $875 : 8$.
\shortans{3}
\loigiai{
\textbf{Đáp số:} $3$\\
$875 = 8 \times 109 + 3$. Số dư là 3.
}
\end{ex}

\newpage

\TaoPhieuKiemTra{Chủ đề 7 -- Ôn tập học kì 1}{42}{Ôn tập biểu thức số}{https://loptoanthaylinh.vn/thi-online/toan-3/bai-42}
\textbf{Phần 1. Trắc nghiệm lựa chọn}
\setcounterpageref{ex}{0}

\begin{ex} %[3D2V42]
Tính giá trị của biểu thức: $M = 350 - 50 \times 4$.
\choice
{$1200$}
{$200$}
{$250$}
{\True $150$}
\loigiai{
Nhân trước: $50 \times 4 = 200$. Trừ sau: $350 - 200 = 150$.
\textbf{Đáp án của câu hỏi là D}
}
\end{ex}

\begin{ex} %[3D2V42]
Tính giá trị của biểu thức có dấu ngoặc: $N = (350 - 50) : 5$.
\choice
{$340$}
{\True $60$}
{$50$}
{$70$}
\loigiai{
Trong ngoặc: $350 - 50 = 300$. Chia sau: $300 : 5 = 60$.
\textbf{Đáp án của câu hỏi là B}
}
\end{ex}

\begin{ex} %[3D2V42]
Thực hiện phép tính: $48 : 4 \times 2$.
\choice
{$6$}
{$12$}
{\True $24$}
{$16$}
\loigiai{
Chỉ có nhân chia, làm từ trái sang phải: $48 : 4 = 12$; $12 \times 2 = 24$.
\textbf{Đáp án của câu hỏi là C}
}
\end{ex}

\begin{ex} %[3D2V42]
Biểu thức nào có giá trị bằng 100?
\choice
{\True $400 : 4 + 25 - 25$}
{$400 : (4 + 25) - 25$}
{$400 - 4 \times 25$}
{$400 : 2 - 50$}
\loigiai{
$400 : 4 + 25 - 25 = 100 + 0 = 100$.
\textbf{Đáp án của câu hỏi là A}
}
\end{ex}

\begin{ex} %[3D2V42]
Tìm $x$, biết: $240 : (x - 5) = 8$.
\choice
{$30$}
{$40$}
{$25$}
{\True $35$}
\loigiai{
$x - 5 = 240 : 8 = 30 \Rightarrow x = 30 + 5 = 35$.
\textbf{Đáp án của câu hỏi là D}
}
\end{ex}

\begin{ex} %[3D2V42]
Tính giá trị biểu thức: $125 \times 4 - 300 : 3$.
\choice
{$350$}
{\True $400$}
{$450$}
{$500$}
\loigiai{
$125 \times 4 = 500$; $300 : 3 = 100$; $500 - 100 = 400$.
\textbf{Đáp án của câu hỏi là B}
}
\end{ex}

\begin{ex} %[3D2V42]
Đặt dấu ngoặc thích hợp để biểu thức $15 + 35 : 5$ có giá trị bằng 10.
\choice
{$15 + (35 : 5)$}
{$(15 + 35 : 5)$}
{\True $(15 + 35) : 5$}
{Không thể đặt được}
\loigiai{
$(15 + 35) : 5 = 50 : 5 = 10$.
\textbf{Đáp án của câu hỏi là C}
}
\end{ex}

\begin{ex} %[3D2V42]
Tính nhanh giá trị biểu thức: $45 \times 6 + 45 \times 4$.
\choice
{\True $450$}
{$400$}
{$500$}
{$480$}
\loigiai{
$45 \times (6 + 4) = 45 \times 10 = 450$.
\textbf{Đáp án của câu hỏi là A}
}
\end{ex}

\begin{ex} %[3D2V42]
Một người có 3 giỏ cam, mỗi giỏ có 15 quả cam và 2 giỏ táo, mỗi giỏ có 20 quả táo. Biểu thức tính tổng số quả cả hai loại là
\choice
{$(15 + 20) \times (3 + 2)$}
{$15 + 3 \times 20 + 2$}
{$15 \times 20 + 3 \times 2$}
{\True $15 \times 3 + 20 \times 2$}
\loigiai{
Số cam là $15 \times 3$, số táo là $20 \times 2$. Tổng số quả: $15 \times 3 + 20 \times 2$.
\textbf{Đáp án của câu hỏi là D}
}
\end{ex}

\begin{ex} %[3D2V42]
Tính giá trị biểu thức: $P = 500 - 100 : (2 + 3)$.
\choice
{$80$}
{\True $480$}
{$450$}
{$490$}
\loigiai{
Trong ngoặc: $2 + 3 = 5$. Chia: $100 : 5 = 20$. Trừ: $500 - 20 = 480$.
\textbf{Đáp án của câu hỏi là B}
}
\end{ex}

\begin{ex} %[3D2V42]
Tìm số tự nhiên $y$ sao cho: $(y + 15) \times 3 = 90$.
\choice
{$10$}
{$20$}
{\True $15$}
{$25$}
\loigiai{
$y + 15 = 90 : 3 = 30 \Rightarrow y = 30 - 15 = 15$.
\textbf{Đáp án của câu hỏi là C}
}
\end{ex}

\begin{ex} %[3D2V42]
So sánh hai biểu thức: $A = (40 + 20) \times 2$ và $B = 40 + 20 \times 2$.
\choice
{\True $A > B$}
{$A < B$}
{$A = B$}
{Không so sánh được}
\loigiai{
$A = 60 \times 2 = 120$. $B = 40 + 40 = 80$. Vì $120 > 80$ nên $A > B$.
\textbf{Đáp án của câu hỏi là A}
}
\end{ex}

\textbf{Phần 2. Trắc nghiệm đúng sai}
\setcounterpageref{ex}{0}

\begin{ex} %[3D2V42]
Xét tính đúng sai của việc tính giá trị biểu thức:
\choiceTF
{\True Biểu thức $100 - 40 + 10$ có giá trị bằng $70$.}
{Biểu thức $100 - (40 + 10)$ có giá trị bằng $70$.}
{\True Dấu ngoặc thay đổi thứ tự ưu tiên của các phép tính trong biểu thức.}
{Tất cả các biểu thức chứa phép nhân đều phải làm từ phải sang trái.}
\loigiai{
\begin{itemize}
    \item Ý a) Đúng: $100 - 40 = 60$; $60 + 10 = 70$.
    \item Ý b) Sai: Sai vì $100 - 50 = 50$.
    \item Ý c) Đúng: Bản chất toán học của ngoặc đơn.
    \item Ý d) Sai: Sai vì nhân chia làm từ trái sang phải.
\end{itemize}
\begin{itemize}
    \item a) \textbf{Đúng}.
    \item b) \textbf{Sai}.
    \item c) \textbf{Đúng}.
    \item d) \textbf{Sai}.
\end{itemize}
}
\end{ex}

\begin{ex} %[3D2V42]
Cho biểu thức $T = 72 : 8 \times 3$:
\choiceTF
{\True Thực hiện phép tính $72 : 8$ trước.}
{\True Kết quả của $72 : 8$ là 9.}
{\True Giá trị của cả biểu thức $T$ bằng 27.}
{Nếu làm phép nhân $8 \times 3$ trước thì giá trị biểu thức không thay đổi.}
\loigiai{
\begin{itemize}
    \item Ý a) Đúng: Quy tắc từ trái sang phải.
    \item Ý b) Đúng: $72 : 8 = 9$.
    \item Ý c) Đúng: $9 \times 3 = 27$.
    \item Ý d) Sai: Sai vì $72 : 24 = 3 \neq 27$.
\end{itemize}
\begin{itemize}
    \item a) \textbf{Đúng}.
    \item b) \textbf{Đúng}.
    \item c) \textbf{Đúng}.
    \item d) \textbf{Sai}.
\end{itemize}
}
\end{ex}

\textbf{Phần 3. Trả lời ngắn}
\setcounterpageref{ex}{0}

\begin{ex} %[3D2V42]
Tính giá trị của biểu thức: $180 - 80 : 4 \times 2$.
\shortans{140}
\loigiai{
\textbf{Đáp số:} $140$\\
$80 : 4 = 20$; $20 \times 2 = 40$; $180 - 40 = 140$.
}
\end{ex}

\begin{ex} %[3D2V42]
Tính giá trị của biểu thức: $(150 + 250) : (4 + 4)$.
\shortans{50}
\loigiai{
\textbf{Đáp số:} $50$\\
$400 : 8 = 50$.
}
\end{ex}

\begin{ex} %[3D2V42]
Tìm $x$, biết: $x + 25 \times 4 = 250$.
\shortans{150}
\loigiai{
\textbf{Đáp số:} $150$\\
$x + 100 = 250 \Rightarrow x = 150$.
}
\end{ex}

\begin{ex} %[3D2V42]
Tính nhanh giá trị biểu thức: $38 \times 5 + 38 \times 5$.
\shortans{380}
\loigiai{
\textbf{Đáp số:} $380$\\
$38 \times 10 = 380$.
}
\end{ex}

\begin{ex} %[3D2V42]
Một người mua 5 gói kẹo giá 12 nghìn đồng mỗi gói và đưa cho cô bán hàng tờ 100 nghìn đồng. Hỏi người đó nhận lại bao nhiêu nghìn đồng tiền thừa?
\shortans{40}
\loigiai{
\textbf{Đáp số:} $40$\\
$100 - 12 \times 5 = 100 - 60 = 40$ nghìn đồng.
}
\end{ex}

\begin{ex} %[3D2V42]
Tìm số tự nhiên $k$ biết: $60 : k + 15 = 25$.
\shortans{6}
\loigiai{
\textbf{Đáp số:} $6$\\
$60 : k = 10 \Rightarrow k = 6$.
}
\end{ex}

\newpage

\TaoPhieuKiemTra{Chủ đề 7 -- Ôn tập học kì 1}{43}{Ôn tập hình học và đo lường}{https://loptoanthaylinh.vn/thi-online/toan-3/bai-43}
\textbf{Phần 1. Trắc nghiệm lựa chọn}
\setcounterpageref{ex}{0}

\begin{ex} %[3H3V43]
Một đoạn thẳng $MN$ dài $64\text{ cm}$. Điểm $I$ là trung điểm của $MN$, điểm $K$ là trung điểm của $MI$. Độ dài đoạn thẳng $KN$ là
\choice
{$32\text{ cm}$}
{$40\text{ cm}$}
{\True $48\text{ cm}$}
{$56\text{ cm}$}
\loigiai{
$MI = 64 : 2 = 32\text{ cm}$; $MK = 32 : 2 = 16\text{ cm}$; $KN = MN - MK = 64 - 16 = 48\text{ (cm)}$.
\textbf{Đáp án của câu hỏi là C}
}
\end{ex}

\begin{ex} %[3H3V43]
Một hình tròn có đường kính $18\text{ cm}$. Bán kính của hình tròn đó là
\choice
{\True $9\text{ cm}$}
{$36\text{ cm}$}
{$6\text{ cm}$}
{$12\text{ cm}$}
\loigiai{
$18 : 2 = 9\text{ (cm)}$.
\textbf{Đáp án của câu hỏi là A}
}
\end{ex}

\begin{ex} %[3H3V43]
Một hình khối lập phương có bao nhiêu đỉnh và bao nhiêu cạnh?
\choice
{$6$ đỉnh và $12$ cạnh}
{$8$ đỉnh và $6$ cạnh}
{$12$ đỉnh và $8$ cạnh}
{\True $8$ đỉnh và $12$ cạnh}
\loigiai{
Khối lập phương có 8 đỉnh và 12 cạnh.
\textbf{Đáp án của câu hỏi là D}
}
\end{ex}

\begin{ex} %[3D3V43]
Đổi $2\text{ m } 4\text{ cm}$ ra đơn vị mi-li-mét:
\choice
{$240\text{ mm}$}
{\True $2040\text{ mm}$}
{$20400\text{ mm}$}
{$2400\text{ mm}$}
\loigiai{
$2\text{ m} = 2000\text{ mm}$; $4\text{ cm} = 40\text{ mm}$. Tổng: $2040\text{ mm}$.
\textbf{Đáp án của câu hỏi là B}
}
\end{ex}

\begin{ex} %[3D3V43]
Một gói đường nặng $1\text{ kg}$. Người ta lấy ra $350\text{ g}$ nấu chè và $250\text{ g}$ làm bánh. Gói đường còn lại bao nhiêu gam?
\choice
{$500\text{ g}$}
{$300\text{ g}$}
{\True $400\text{ g}$}
{$600\text{ g}$}
\loigiai{
Đổi $1\text{ kg} = 1000\text{ g}$. Đã dùng: $350 + 250 = 600\text{ g}$. Còn lại: $1000 - 600 = 400\text{ (g)}$.
\textbf{Đáp án của câu hỏi là C}
}
\end{ex}

\begin{ex} %[3D3V43]
Một bình nước có dung tích $1\ell$. Rót đầy 4 cốc nước loại $200\text{ ml}$. Trong bình còn lại bao nhiêu mi-li-lít nước?
\choice
{\True $200\text{ ml}$}
{$100\text{ ml}$}
{$300\text{ ml}$}
{$400\text{ ml}$}
\loigiai{
Đổi $1\ell = 1000\text{ ml}$. Bốn cốc chứa: $200 \times 4 = 800\text{ ml}$. Còn lại: $1000 - 800 = 200\text{ (ml)}$.
\textbf{Đáp án của câu hỏi là A}
}
\end{ex}

\begin{ex} %[3D3V43]
Nhiệt độ nước đá đang tan là $0^\circ\text{C}$, nhiệt độ nước sôi là $100^\circ\text{C}$. Mức chênh lệch nhiệt độ giữa nước sôi và nước đá đang tan là
\choice
{$50^\circ\text{C}$}
{$0^\circ\text{C}$}
{$90^\circ\text{C}$}
{\True $100^\circ\text{C}$}
\loigiai{
$100 - 0 = 100^\circ\text{C}$.
\textbf{Đáp án của câu hỏi là D}
}
\end{ex}

\begin{ex} %[3H3V43]
Một hình chữ nhật có chu vi $50\text{ cm}$, chiều dài là $15\text{ cm}$. Chiều rộng của hình chữ nhật là
\choice
{$20\text{ cm}$}
{\True $10\text{ cm}$}
{$5\text{ cm}$}
{$12\text{ cm}$}
\loigiai{
Nửa chu vi: $50 : 2 = 25\text{ cm}$. Chiều rộng: $25 - 15 = 10\text{ (cm)}$.
\textbf{Đáp án của câu hỏi là B}
}
\end{ex}

\begin{ex} %[3H3V43]
Góc tạo bởi kim giờ và kim phút lúc đúng 3 giờ là
\choice
{Góc nhọn}
{Góc tù}
{\True Góc vuông}
{Góc bẹt}
\loigiai{
Lúc 3 giờ đúng hai kim vuông góc với nhau.
\textbf{Đáp án của câu hỏi là C}
}
\end{ex}

\begin{ex} %[3H3V43]
Một khối gỗ hình lập phương có cạnh $6\text{ cm}$. Tổng độ dài tất cả các cạnh của khối gỗ là
\choice
{\True $72\text{ cm}$}
{$36\text{ cm}$}
{$48\text{ cm}$}
{$24\text{ cm}$}
\loigiai{
Khối lập phương có 12 cạnh: $12 \times 6 = 72\text{ (cm)}$.
\textbf{Đáp án của câu hỏi là A}
}
\end{ex}

\begin{ex} %[3D3V43]
Một con gà cân nặng $2\text{ kg } 300\text{ g}$, con ngỗng cân nặng $4\text{ kg } 100\text{ g}$. Cả gà và ngỗng cân nặng là
\choice
{$6\text{ kg } 200\text{ g}$}
{$6\text{ kg } 500\text{ g}$}
{$7\text{ kg}$}
{\True $6\text{ kg } 400\text{ g}$}
\loigiai{
$2300 + 4100 = 6400\text{ g} = 6\text{ kg } 400\text{ g}$.
\textbf{Đáp án của câu hỏi là D}
}
\end{ex}

\begin{ex} %[3H3V43]
Người ta uốn một đoạn dây thép dài $60\text{ cm}$ thành một hình tam giác có 3 cạnh bằng nhau. Độ dài mỗi cạnh tam giác là
\choice
{$15\text{ cm}$}
{\True $20\text{ cm}$}
{$30\text{ cm}$}
{$25\text{ cm}$}
\loigiai{
$60 : 3 = 20\text{ (cm)}$.
\textbf{Đáp án của câu hỏi là B}
}
\end{ex}

\textbf{Phần 2. Trắc nghiệm đúng sai}
\setcounterpageref{ex}{0}

\begin{ex} %[3H3V43]
Xét tính đúng sai của các kiến thức hình học và đo lường:
\choiceTF
{\True Hình vuông có 4 cạnh bằng nhau và 4 góc vuông.}
{Đoạn thẳng nối hai điểm bất kì trên đường tròn là đường kính.}
{\True $1\text{ kg} = 1000\text{ g}$ và $1\ell = 1000\text{ ml}$.}
{$1\text{ m} = 100\text{ mm}$.}
\loigiai{
\begin{itemize}
    \item Ý a) Đúng: Định nghĩa hình vuông.
    \item Ý b) Sai: Sai vì chỉ đoạn qua tâm mới là đường kính.
    \item Ý c) Đúng: Quan hệ chuẩn.
    \item Ý d) Sai: Sai vì $1\text{ m} = 1000\text{ mm}$.
\end{itemize}
\begin{itemize}
    \item a) \textbf{Đúng}.
    \item b) \textbf{Sai}.
    \item c) \textbf{Đúng}.
    \item d) \textbf{Sai}.
\end{itemize}
}
\end{ex}

\begin{ex} %[3H3V43]
Một hình chữ nhật có chiều dài $18\text{ cm}$, chiều rộng bằng $\dfrac{1}{3}$ chiều dài:
\choiceTF
{\True Chiều rộng của hình chữ nhật bằng $6\text{ cm}$.}
{\True Chu vi của hình chữ nhật bằng $48\text{ cm}$.}
{Có thể chia hình chữ nhật đó thành 2 hình vuông cạnh $6\text{ cm}$.}
{\True Nửa chu vi của hình chữ nhật bằng $24\text{ cm}$.}
\loigiai{
\begin{itemize}
    \item Ý a) Đúng: $18 : 3 = 6\text{ cm}$.
    \item Ý b) Đúng: $(18 + 6) \times 2 = 48\text{ cm}$.
    \item Ý c) Sai: Sai vì $18 : 6 = 3$ hình vuông.
    \item Ý d) Đúng: $48 : 2 = 24\text{ cm}$.
\end{itemize}
\begin{itemize}
    \item a) \textbf{Đúng}.
    \item b) \textbf{Đúng}.
    \item c) \textbf{Sai}.
    \item d) \textbf{Đúng}.
\end{itemize}
}
\end{ex}

\textbf{Phần 3. Trả lời ngắn}
\setcounterpageref{ex}{0}

\begin{ex} %[3H3V43]
Một hình vuông có chu vi $56\text{ cm}$. Tính độ dài cạnh hình vuông theo xăng-ti-mét.
\shortans{14}
\loigiai{
\textbf{Đáp số:} $14$\\
$56 : 4 = 14\text{ (cm)}$.
}
\end{ex}

\begin{ex} %[3H3V43]
Hình tròn có bán kính $15\text{ cm}$. Hỏi đường kính dài bao nhiêu xăng-ti-mét?
\shortans{30}
\loigiai{
\textbf{Đáp số:} $30$\\
$15 \times 2 = 30\text{ (cm)}$.
}
\end{ex}

\begin{ex} %[3D3V43]
Đổi $1\text{ m } 8\text{ cm}$ ra đơn vị mi-li-mét.
\shortans{1080}
\loigiai{
\textbf{Đáp số:} $1080$\\
$1000 + 80 = 1080\text{ mm}$.
}
\end{ex}

\begin{ex} %[3D3V43]
Một túi gạo nặng $2\text{ kg}$. Người ta đã dùng $750\text{ g}$. Hỏi túi gạo còn lại bao nhiêu gam?
\shortans{1250}
\loigiai{
\textbf{Đáp số:} $1250$\\
$2000 - 750 = 1250\text{ (g)}$.
}
\end{ex}

\begin{ex} %[3D3V43]
Một bình chứa $1\ell\text{ }200\text{ ml}$ nước. Rót đều vào 4 cốc. Mỗi cốc chứa bao nhiêu mi-li-lít nước?
\shortans{300}
\loigiai{
\textbf{Đáp số:} $300$\\
$1200 : 4 = 300\text{ (ml)}$.
}
\end{ex}

\begin{ex} %[3H3V43]
Một khối lập phương có chu vi một mặt là $24\text{ cm}$. Tính tổng độ dài 12 cạnh của khối lập phương đó theo xăng-ti-mét.
\shortans{72}
\loigiai{
\textbf{Đáp số:} $72$\\
Cạnh: $24 : 4 = 6\text{ cm}$. Tổng 12 cạnh: $12 \times 6 = 72\text{ (cm)}$.
}
\end{ex}

\newpage

\TaoPhieuKiemTra{Chủ đề 7 -- Ôn tập học kì 1}{44}{Ôn tập chung học kì 1}{https://loptoanthaylinh.vn/thi-online/toan-3/bai-44}
\textbf{Phần 1. Trắc nghiệm lựa chọn}
\setcounterpageref{ex}{0}

\begin{ex} %[3X3V44]
Tính giá trị của biểu thức: $A = 150 \times 4 - 450 : 3$.
\choice
{$400$}
{\True $450$}
{$500$}
{$350$}
\loigiai{
$150 \times 4 = 600$; $450 : 3 = 150$; $600 - 150 = 450$.
\textbf{Đáp án của câu hỏi là B}
}
\end{ex}

\begin{ex} %[3X3V44]
Một cửa hàng có 3 kho thóc, mỗi kho chứa $240\text{ kg}$ thóc. Cửa hàng đã bán được $\dfrac{1}{3}$ tổng số thóc. Cửa hàng còn lại bao nhiêu ki-lô-gam thóc?
\choice
{$240\text{ kg}$}
{$500\text{ kg}$}
{$420\text{ kg}$}
{\True $480\text{ kg}$}
\loigiai{
Tổng thóc: $240 \times 3 = 720\text{ kg}$. Đã bán: $720 : 3 = 240\text{ kg}$. Còn lại: $720 - 240 = 480\text{ (kg)}$.
\textbf{Đáp án của câu hỏi là D}
}
\end{ex}

\begin{ex} %[3X3V44]
Tìm số tự nhiên $x$, biết: $x \times 5 + 125 = 625$.
\choice
{\True $100$}
{$90$}
{$110$}
{$120$}
\loigiai{
$x \times 5 = 625 - 125 = 500 \Rightarrow x = 500 : 5 = 100$.
\textbf{Đáp án của câu hỏi là A}
}
\end{ex}

\begin{ex} %[3X3V44]
Một mảnh vườn hình chữ nhật có chu vi $84\text{ m}$, chiều rộng là $14\text{ m}$. Chiều dài của mảnh vườn đó là
\choice
{$24\text{ m}$}
{$32\text{ m}$}
{\True $28\text{ m}$}
{$26\text{ m}$}
\loigiai{
Nửa chu vi: $84 : 2 = 42\text{ m}$. Chiều dài: $42 - 14 = 28\text{ (m)}$.
\textbf{Đáp án của câu hỏi là C}
}
\end{ex}

\begin{ex} %[3X3V44]
Trong phép chia cho 8, số dư lớn nhất có thể là
\choice
{$8$}
{\True $7$}
{$6$}
{$5$}
\loigiai{
Số dư luôn bé hơn số chia, lớn nhất là $8 - 1 = 7$.
\textbf{Đáp án của câu hỏi là B}
}
\end{ex}

\begin{ex} %[3X3V44]
Số lớn là 90, số bé bằng $\dfrac{1}{5}$ số lớn. Số lớn hơn số bé bao nhiêu đơn vị?
\choice
{$18$}
{$75$}
{$80$}
{\True $72$}
\loigiai{
Số bé: $90 : 5 = 18$. Hiệu: $90 - 18 = 72$ đơn vị.
\textbf{Đáp án của câu hỏi là D}
}
\end{ex}

\begin{ex} %[3X3V44]
Một hình tròn có bán kính $8\text{ cm}$. Đoạn thẳng dài nhất nối hai điểm trên hình tròn đó có độ dài là
\choice
{\True $16\text{ cm}$}
{$8\text{ cm}$}
{$24\text{ cm}$}
{$32\text{ cm}$}
\loigiai{
Đoạn thẳng dài nhất nối hai điểm trên đường tròn chính là đường kính: $8 \times 2 = 16\text{ (cm)}$.
\textbf{Đáp án của câu hỏi là A}
}
\end{ex}

\begin{ex} %[3X3V44]
Có 4 xe tải chở hàng, mỗi xe chở được 115 bao phân bón. Người ta chuyển xuống 180 bao. Trên các xe còn lại bao nhiêu bao phân bón?
\choice
{$260$ bao}
{$300$ bao}
{\True $280$ bao}
{$320$ bao}
\loigiai{
Tổng số bao: $115 \times 4 = 460$ bao. Còn lại: $460 - 180 = 280$ bao.
\textbf{Đáp án của câu hỏi là C}
}
\end{ex}

\begin{ex} %[3X3V44]
Nhiệt độ buổi sáng là $18^\circ\text{C}$, buổi trưa nhiệt độ tăng thêm $8^\circ\text{C}$, buổi tối nhiệt độ lại giảm đi $4^\circ\text{C}$ so với buổi trưa. Nhiệt độ buổi tối là
\choice
{$20^\circ\text{C}$}
{\True $22^\circ\text{C}$}
{$24^\circ\text{C}$}
{$26^\circ\text{C}$}
\loigiai{
Buổi trưa: $18 + 8 = 26^\circ\text{C}$. Buổi tối: $26 - 4 = 22^\circ\text{C}$.
\textbf{Đáp án của câu hỏi là B}
}
\end{ex}

\begin{ex} %[3X3V44]
Tính nhanh giá trị: $125 \times 3 + 125 \times 5$.
\choice
{$875$}
{$1250$}
{$950$}
{\True $1000$}
\loigiai{
$125 \times (3 + 5) = 125 \times 8 = 1000$.
\textbf{Đáp án của câu hỏi là D}
}
\end{ex}

\begin{ex} %[3X3V44]
Cho đoạn thẳng $AB = 72\text{ cm}$. Điểm $M$ là trung điểm $AB$, điểm $N$ là trung điểm $AM$. Độ dài đoạn thẳng $NB$ là
\choice
{\True $54\text{ cm}$}
{$36\text{ cm}$}
{$48\text{ cm}$}
{$60\text{ cm}$}
\loigiai{
$AM = 36\text{ cm}$; $AN = 18\text{ cm}$; $NB = 72 - 18 = 54\text{ (cm)}$.
\textbf{Đáp án của câu hỏi là A}
}
\end{ex}

\begin{ex} %[3X3V44]
Một khối lập phương có chu vi một mặt là $32\text{ cm}$. Tổng độ dài tất cả 12 cạnh của khối lập phương là
\choice
{$64\text{ cm}$}
{$128\text{ cm}$}
{\True $96\text{ cm}$}
{$72\text{ cm}$}
\loigiai{
Cạnh: $32 : 4 = 8\text{ cm}$. Tổng 12 cạnh: $12 \times 8 = 96\text{ (cm)}$.
\textbf{Đáp án của câu hỏi là C}
}
\end{ex}

\textbf{Phần 2. Trắc nghiệm đúng sai}
\setcounterpageref{ex}{0}

\begin{ex} %[3X3V44]
Xét tính đúng sai của các nhận định ôn tập cuối học kì 1:
\choiceTF
{\True Thực hiện phép tính trong ngoặc trước, ngoài ngoặc sau.}
{$1\text{ kg} = 100\text{ g}$.}
{\True Đường kính của hình tròn dài gấp 2 lần bán kính.}
{Một hình tam giác có thể có 2 góc vuông.}
\loigiai{
\begin{itemize}
    \item Ý a) Đúng: Quy tắc cơ bản.
    \item Ý b) Sai: Sai vì $1\text{ kg} = 1000\text{ g}$.
    \item Ý c) Đúng: Tính chất hình tròn.
    \item Ý d) Sai: Sai vì tam giác chỉ có nhiều nhất 1 góc vuông.
\end{itemize}
\begin{itemize}
    \item a) \textbf{Đúng}.
    \item b) \textbf{Sai}.
    \item c) \textbf{Đúng}.
    \item d) \textbf{Sai}.
\end{itemize}
}
\end{ex}

\begin{ex} %[3X3V44]
Một mảnh vườn hình chữ nhật có chiều dài $24\text{ m}$, chiều rộng bằng nửa chiều dài:
\choiceTF
{\True Chiều rộng mảnh vườn là $12\text{ m}$.}
{\True Chu vi mảnh vườn là $72\text{ m}$.}
{Nếu gấp đôi chiều rộng thì mảnh vườn có chu vi không đổi.}
{\True Chia mảnh vườn thành 2 hình vuông cạnh $12\text{ m}$ là hoàn toàn chính xác.}
\loigiai{
\begin{itemize}
    \item Ý a) Đúng: $24 : 2 = 12\text{ m}$.
    \item Ý b) Đúng: $(24 + 12) \times 2 = 72\text{ m}$.
    \item Ý c) Sai: Sai vì kích thước thay đổi thì chu vi tăng.
    \item Ý d) Đúng: $24 : 12 = 2$ hình vuông.
\end{itemize}
\begin{itemize}
    \item a) \textbf{Đúng}.
    \item b) \textbf{Đúng}.
    \item c) \textbf{Sai}.
    \item d) \textbf{Đúng}.
\end{itemize}
}
\end{ex}

\textbf{Phần 3. Trả lời ngắn}
\setcounterpageref{ex}{0}

\begin{ex} %[3X3V44]
Tính giá trị biểu thức: $120 \times 3 - 160 : 4$.
\shortans{320}
\loigiai{
\textbf{Đáp số:} $320$\\
$360 - 40 = 320$.
}
\end{ex}

\begin{ex} %[3X3V44]
Một hình vuông có chu vi $68\text{ cm}$. Tính độ dài cạnh của hình vuông đó theo xăng-ti-mét.
\shortans{17}
\loigiai{
\textbf{Đáp số:} $17$\\
$68 : 4 = 17\text{ (cm)}$.
}
\end{ex}

\begin{ex} %[3X3V44]
Số lớn là 84, số bé là 7. Hỏi số lớn gấp mấy lần số bé?
\shortans{12}
\loigiai{
\textbf{Đáp số:} $12$\\
$84 : 7 = 12$ lần.
}
\end{ex}

\begin{ex} %[3X3V44]
Tính kết quả phép chia: $945 : 9$.
\shortans{105}
\loigiai{
\textbf{Đáp số:} $105$\\
$945 : 9 = 105$.
}
\end{ex}

\begin{ex} %[3X3V44]
Đổi $2\ell - 650\text{ ml}$ ra đơn vị mi-li-lít.
\shortans{1350}
\loigiai{
\textbf{Đáp số:} $1350$\\
$2000 - 650 = 1350\text{ (ml)}$.
}
\end{ex}

\begin{ex} %[3X3V44]
Có 3 thùng bánh, mỗi thùng có 32 hộp bánh. Người ta chia đều số bánh đó cho 4 lớp học sinh. Hỏi mỗi lớp nhận được bao nhiêu hộp bánh?
\shortans{24}
\loigiai{
\textbf{Đáp số:} $24$\\
Tổng: $32 \times 3 = 96$ hộp. Mỗi lớp: $96 : 4 = 24$ hộp.
}
\end{ex}

\newpage

% =========================================================================
% CHỦ ĐỀ 8 - CÁC SỐ ĐẾN 10 000
% =========================================================================
\section{Các số đến 10 000}

\TaoPhieuKiemTra{Chủ đề 8 -- Các số đến 10 000}{45}{Các số có bốn chữ số. Số 10 000}{https://loptoanthaylinh.vn/thi-online/toan-3/bai-45}
\textbf{Phần 1. Trắc nghiệm lựa chọn}
\setcounterpageref{ex}{0}

\begin{ex} %[3D1V45]
Số gồm $5$ nghìn, $8$ trăm, $0$ chục và $4$ đơn vị được viết là
\choice
{$5840$}
{$5084$}
{\True $5804$}
{$584$}
\loigiai{
Chữ số hàng nghìn là 5, hàng trăm là 8, hàng chục là 0, hàng đơn vị là 4. Số đó là $5804$.
\textbf{Đáp án của câu hỏi là C}
}
\end{ex}

\begin{ex} %[3D1V45]
Số liền sau của số $9999$ là
\choice
{\True $10\,000$}
{$9998$}
{$10\,001$}
{$9900$}
\loigiai{
Số liền sau của 9999 là: $9999 + 1 = 10\,000$.
\textbf{Đáp án của câu hỏi là A}
}
\end{ex}

\begin{ex} %[3D1V45]
Từ bốn chữ số $0; 3; 5; 8$, viết được số có bốn chữ số khác nhau lớn nhất là
\choice
{$8503$}
{$8350$}
{$5830$}
{\True $8530$}
\loigiai{
Chọn chữ số lớn nhất ở hàng cao nhất: hàng nghìn là 8, hàng trăm là 5, hàng chục là 3, hàng đơn vị là 0. Số đó là $8530$.
\textbf{Đáp án của câu hỏi là D}
}
\end{ex}

\begin{ex} %[3D1V45]
Số có bốn chữ số mà chữ số hàng nghìn gấp đôi chữ số hàng trăm, chữ số hàng chục gấp đôi chữ số hàng đơn vị và tổng các chữ số bằng 9 là
\choice
{$6310$}
{\True $4221$}
{$2142$}
{$4212$}
\loigiai{
Xét số 4221: nghìn 4 gấp đôi trăm 2; chục 2 gấp đôi đơn vị 1; tổng các chữ số là $4 + 2 + 2 + 1 = 9$ (thỏa mãn).
\textbf{Đáp án của câu hỏi là B}
}
\end{ex}

\begin{ex} %[3D1V45]
Trong số $7852$, giá trị của chữ số 8 lớn hơn giá trị của chữ số 5 bao nhiêu đơn vị?
\choice
{$700$}
{$300$}
{\True $750$}
{$800$}
\loigiai{
Chữ số 8 ở hàng trăm có giá trị 800. Chữ số 5 ở hàng chục có giá trị 50. Chênh lệch: $800 - 50 = 750$ đơn vị.
\textbf{Đáp án của câu hỏi là C}
}
\end{ex}

\begin{ex} %[3D1V45]
Số chẵn bé nhất có bốn chữ số khác nhau là
\choice
{\True $1024$}
{$1000$}
{$1023$}
{$1234$}
\loigiai{
Chữ số hàng nghìn là 1, hàng trăm là 0, hàng chục là 2, hàng đơn vị là chữ số chẵn bé nhất khác 0 và 2, tức là 4. Số đó là $1024$.
\textbf{Đáp án của câu hỏi là A}
}
\end{ex}

\begin{ex} %[3D1V45]
Có bao nhiêu số tròn trăm có bốn chữ số bé hơn $2000$?
\choice
{$9$ số}
{$8$ số}
{$11$ số}
{\True $10$ số}
\loigiai{
Các số tròn trăm từ 1000 đến 1900 gồm: 1000, 1100, 1200, ..., 1900. Có tất cả $(1900 - 1000) : 100 + 1 = 10$ số.
\textbf{Đáp án của câu hỏi là D}
}
\end{ex}

\begin{ex} %[3D1V45]
Một số có bốn chữ số có chữ số hàng đơn vị là 5. Nếu xóa bỏ chữ số 5 đó thì ta được số mới kém số ban đầu $2489$ đơn vị. Số ban đầu là
\choice
{$2545$}
{\True $2765$}
{$2655$}
{$2855$}
\loigiai{
Gọi số mới là $A$. Số ban đầu là $A \times 10 + 5$. Hiệu: $A \times 9 + 5 = 2489 \Rightarrow A \times 9 = 2484 \Rightarrow A = 276$. Số ban đầu là $2765$.
\textbf{Đáp án của câu hỏi là B}
}
\end{ex}

\begin{ex} %[3D1V45]
Số liền trước của số chẵn bé nhất có bốn chữ số là
\choice
{$1001$}
{$1000$}
{\True $999$}
{$998$}
\loigiai{
Số chẵn bé nhất có bốn chữ số là 1000. Số liền trước của nó là: $1000 - 1 = 999$.
\textbf{Đáp án của câu hỏi là C}
}
\end{ex}

\begin{ex} %[3D1V45]
Tìm một số biết rằng chữ số hàng nghìn là số lớn nhất có một chữ số, chữ số hàng đơn vị là số bé nhất có một chữ số, các chữ số còn lại đều bằng 3. Số đó là
\choice
{\True $9330$}
{$9303$}
{$9033$}
{$9333$}
\loigiai{
Hàng nghìn là 9, hàng đơn vị là 0, hàng trăm và hàng chục là 3. Số đó là $9330$.
\textbf{Đáp án của câu hỏi là A}
}
\end{ex}

\begin{ex} %[3D1V45]
Cho dãy số: $1000, 1500, 2000, 2500, \dots$ Số hạng thứ 8 của dãy số là
\choice
{$4000$}
{$5000$}
{$3500$}
{\True $4500$}
\loigiai{
Quy luật cộng thêm 500. Số hạng thứ 8 là: $1000 + 500 \times 7 = 1000 + 3500 = 4500$.
\textbf{Đáp án của câu hỏi là D}
}
\end{ex}

\begin{ex} %[3D1V45]
Số $10\,000$ gấp số $1000$ bao nhiêu lần?
\choice
{$100$ lần}
{\True $10$ lần}
{$5$ lần}
{$2$ lần}
\loigiai{
$10\,000 : 1000 = 10$ (lần).
\textbf{Đáp án của câu hỏi là B}
}
\end{ex}

\textbf{Phần 2. Trắc nghiệm đúng sai}
\setcounterpageref{ex}{0}

\begin{ex} %[3D1V45]
Xét tính đúng sai của các nhận định về các số đến 10 000:
\choiceTF
{\True Số 10 000 là số bé nhất có năm chữ số.}
{Số lớn nhất có bốn chữ số là 10 000.}
{\True Trong số 4725, chữ số 4 có giá trị là 4000.}
{Số 8080 đọc là tám trăm tám mươi.}
\loigiai{
\begin{itemize}
    \item Ý a) Đúng: Định nghĩa.
    \item Ý b) Sai: Sai vì số lớn nhất có 4 chữ số là 9999.
    \item Ý c) Đúng: Chữ số hàng nghìn.
    \item Ý d) Sai: Sai vì đọc là tám nghìn không trăm tám mươi.
\end{itemize}
\begin{itemize}
    \item a) \textbf{Đúng}.
    \item b) \textbf{Sai}.
    \item c) \textbf{Đúng}.
    \item d) \textbf{Sai}.
\end{itemize}
}
\end{ex}

\begin{ex} %[3D1V45]
Cho các chữ số $2; 0; 7; 9$:
\choiceTF
{\True Số lớn nhất có 4 chữ số khác nhau ghép từ các chữ số trên là 9720.}
{\True Số bé nhất có 4 chữ số khác nhau ghép từ các chữ số trên là 2079.}
{Có thể ghép được 24 số có 4 chữ số khác nhau từ 4 chữ số trên.}
{\True Hiệu giữa số lớn nhất và số bé nhất lập được là: $9720 - 2079 = 7641$.}
\loigiai{
\begin{itemize}
    \item Ý a) Đúng: Ghép theo thứ tự giảm dần.
    \item Ý b) Đúng: Chữ số hàng nghìn khác 0, ghép tăng dần.
    \item Ý c) Sai: Sai vì chữ số 0 không đứng đầu được nên chỉ có $3 \times 3 \times 2 \times 1 = 18$ số.
    \item Ý d) Đúng: Phép trừ chính xác.
\end{itemize}
\begin{itemize}
    \item a) \textbf{Đúng}.
    \item b) \textbf{Đúng}.
    \item c) \textbf{Sai}.
    \item d) \textbf{Đúng}.
\end{itemize}
}
\end{ex}

\textbf{Phần 3. Trả lời ngắn}
\setcounterpageref{ex}{0}

\begin{ex} %[3D1V45]
Tìm số liền trước của số 8500.
\shortans{8499}
\loigiai{
\textbf{Đáp số:} $8499$\\
$8500 - 1 = 8499$.
}
\end{ex}

\begin{ex} %[3D1V45]
Viết số bé nhất có bốn chữ số khác nhau.
\shortans{1023}
\loigiai{
\textbf{Đáp số:} $1023$\\
Số bé nhất có 4 chữ số khác nhau là 1023.
}
\end{ex}

\begin{ex} %[3D1V45]
Tính hiệu giữa số lớn nhất có 4 chữ số và số bé nhất có 4 chữ số.
\shortans{8999}
\loigiai{
\textbf{Đáp số:} $8999$\\
$9999 - 1000 = 8999$.
}
\end{ex}

\begin{ex} %[3D1V45]
Tìm giá trị của chữ số 6 trong số 6742.
\shortans{6000}
\loigiai{
\textbf{Đáp số:} $6000$\\
Chữ số 6 ở hàng nghìn có giá trị là 6000.
}
\end{ex}

\begin{ex} %[3D1V45]
Có bao nhiêu số có 4 chữ số đều có chữ số tận cùng là 0 từ 1000 đến 1090?
\shortans{10}
\loigiai{
\textbf{Đáp số:} $10$\\
Các số: 1000, 1010, 1020, 1030, 1040, 1050, 1060, 1070, 1080, 1090. Có 10 số.
}
\end{ex}

\begin{ex} %[3D1V45]
Một số có 4 chữ số có tổng các chữ số bằng 2. Hỏi có bao nhiêu số như vậy?
\shortans{4}
\loigiai{
\textbf{Đáp số:} $4$\\
Các số thỏa mãn: 2000, 1100, 1010, 1001. Có 4 số.
}
\end{ex}

\newpage

\TaoPhieuKiemTra{Chủ đề 8 -- Các số đến 10 000}{46}{So sánh các số trong phạm vi 10 000}{https://loptoanthaylinh.vn/thi-online/toan-3/bai-46}
\textbf{Phần 1. Trắc nghiệm lựa chọn}
\setcounterpageref{ex}{0}

\begin{ex} %[3D1V46]
Điền dấu thích hợp vào chỗ chấm: $4582 \dots 4590$.
\choice
{$>$}
{\True $<$}
{$=$}
{Không so sánh được}
\loigiai{
Hàng nghìn và hàng trăm bằng nhau, so sánh hàng chục: $8 < 9$ nên $4582 < 4590$.
\textbf{Đáp án của câu hỏi là B}
}
\end{ex}

\begin{ex} %[3D1V46]
Trong các số $7856, 7865, 8756, 8675$, số lớn nhất là
\choice
{$7865$}
{$8675$}
{$7856$}
{\True $8756$}
\loigiai{
So sánh hàng nghìn: số 8756 và 8675 lớn hơn. Tiếp tục so sánh hàng trăm: $7 > 6$ nên $8756$ lớn nhất.
\textbf{Đáp án của câu hỏi là D}
}
\end{ex}

\begin{ex} %[3D1V46]
Tìm chữ số thích hợp thay vào dấu $*$ để: $68*4 > 6884$. Chữ số đó là
\choice
{\True $9$}
{$8$}
{$7$}
{$6$}
\loigiai{
Hàng nghìn, hàng trăm và hàng đơn vị bằng nhau. Để $68*4 > 6884$ thì $* > 8$, suy ra $* = 9$.
\textbf{Đáp án của câu hỏi là A}
}
\end{ex}

\begin{ex} %[3D1V46]
Sắp xếp các số sau theo thứ tự từ bé đến lớn: $3050, 3500, 3005, 3505$.
\choice
{$3505, 3500, 3050, 3005$}
{$3050, 3005, 3500, 3505$}
{\True $3005, 3050, 3500, 3505$}
{$3005, 3500, 3050, 3505$}
\loigiai{
So sánh các số: $3005 < 3050 < 3500 < 3505$.
\textbf{Đáp án của câu hỏi là C}
}
\end{ex}

\begin{ex} %[3D1V46]
Có bao nhiêu số tự nhiên $x$ có bốn chữ số thỏa mãn: $2998 < x < 3003$?
\choice
{$3$ số}
{\True $4$ số}
{$5$ số}
{$6$ số}
\loigiai{
Các số tự nhiên $x$ là: $2999, 3000, 3001, 3002$. Có 4 số thỏa mãn.
\textbf{Đáp án của câu hỏi là B}
}
\end{ex}

\begin{ex} %[3D1V46]
Số bé nhất trong các số sau là số nào: $999, 1002, 1020, 2001$?
\choice
{$1002$}
{$1020$}
{$2001$}
{\True $999$}
\loigiai{
Số có ít chữ số hơn thì bé hơn. Số 999 có 3 chữ số nên bé nhất.
\textbf{Đáp án của câu hỏi là D}
}
\end{ex}

\begin{ex} %[3D1V46]
Tìm số tự nhiên tròn trăm $x$ thỏa mãn: $4500 < x < 4800$. Các giá trị của $x$ là
\choice
{\True $4600$ và $4700$}
{$4600$}
{$4700$}
{$4550$ và $4650$}
\loigiai{
Số tròn trăm nằm giữa 4500 và 4800 là 4600 và 4700.
\textbf{Đáp án của câu hỏi là A}
}
\end{ex}

\begin{ex} %[3D1V46]
So sánh hai biểu thức: $A = 3000 + 400 + 50 + 2$ và $B = 3450 + 5$. Khẳng định đúng là
\choice
{$A > B$}
{$A = B$}
{\True $A < B$}
{Không so sánh được}
\loigiai{
$A = 3452$; $B = 3455$. Vì $3452 < 3455$ nên $A < B$.
\textbf{Đáp án của câu hỏi là C}
}
\end{ex}

\begin{ex} %[3D1V46]
Bốn bạn An, Bình, Cúc, Dũng thi chạy. Điểm tích lũy lần lượt là: An được 4520 điểm, Bình được 4250 điểm, Cúc được 4502 điểm, Dũng được 4052 điểm. Bạn đạt điểm cao nhất là
\choice
{Bình}
{\True An}
{Cúc}
{Dũng}
\loigiai{
So sánh các số: $4520 > 4502 > 4250 > 4052$. Điểm cao nhất là An (4520 điểm).
\textbf{Đáp án của câu hỏi là B}
}
\end{ex}

\begin{ex} %[3D1V46]
Tìm chữ số $x$ biết rằng: $5x78 < 5178$. Chữ số $x$ là
\choice
{$1$}
{$2$}
{Không có chữ số nào}
{\True $0$}
\loigiai{
Để $5x78 < 5178$ thì $x < 1$. Vì $x$ là chữ số nên $x = 0$.
\textbf{Đáp án của câu hỏi là D}
}
\end{ex}

\begin{ex} %[3D1V46]
Trong các số có bốn chữ số khác nhau lập từ bốn chữ số $1; 4; 6; 8$, số lớn thứ hai là
\choice
{\True $8614$}
{$8641$}
{$8461$}
{$8416$}
\loigiai{
Số lớn nhất là 8641. Số lớn thứ hai là 8614.
\textbf{Đáp án của câu hỏi là A}
}
\end{ex}

\begin{ex} %[3D1V46]
Tìm số tự nhiên lớn nhất có bốn chữ số mà tổng các chữ số của nó bằng 12.
\choice
{$9210$}
{$8400$}
{\True $9300$}
{$9120$}
\loigiai{
Để số lớn nhất, hàng nghìn lớn nhất là 9. Còn lại tổng là $12 - 9 = 3$. Đặt 3 ở hàng trăm, hàng chục và đơn vị là 0. Số đó là $9300$.
\textbf{Đáp án của câu hỏi là C}
}
\end{ex}

\textbf{Phần 2. Trắc nghiệm đúng sai}
\setcounterpageref{ex}{0}

\begin{ex} %[3D1V46]
Xét tính đúng sai của việc so sánh các số có bốn chữ số:
\choiceTF
{\True Số có nhiều chữ số hơn thì luôn luôn lớn hơn số có ít chữ số hơn.}
{Hai số có cùng số chữ số thì số nào có chữ số hàng đơn vị lớn hơn sẽ lớn hơn.}
{\True Nếu hai số có hàng nghìn và hàng trăm bằng nhau thì ta so sánh hàng chục.}
{$9876 > 10\,000$.}
\loigiai{
\begin{itemize}
    \item Ý a) Đúng: Quy tắc so sánh số tự nhiên.
    \item Ý b) Sai: Sai vì phải so sánh từ hàng cao nhất trước.
    \item Ý c) Đúng: Thứ tự so sánh chuẩn.
    \item Ý d) Sai: Sai vì 9876 có 4 chữ số, 10 000 có 5 chữ số.
\end{itemize}
\begin{itemize}
    \item a) \textbf{Đúng}.
    \item b) \textbf{Sai}.
    \item c) \textbf{Đúng}.
    \item d) \textbf{Sai}.
\end{itemize}
}
\end{ex}

\begin{ex} %[3D1V46]
Cho bốn số: $4280; 4820; 4082; 4802$:
\choiceTF
{\True Số bé nhất trong bốn số là 4082.}
{\True Số lớn nhất trong bốn số là 4820.}
{Số 4802 lớn hơn số 4820.}
{\True Sắp xếp theo thứ tự giảm dần là: $4820, 4802, 4280, 4082$.}
\loigiai{
\begin{itemize}
    \item Ý a) Đúng: Hàng trăm bé nhất là 0.
    \item Ý b) Đúng: Hàng trăm 8, hàng chục 2.
    \item Ý c) Sai: Sai vì hàng chục $0 < 2$.
    \item Ý d) Đúng: Đúng thứ tự.
\end{itemize}
\begin{itemize}
    \item a) \textbf{Đúng}.
    \item b) \textbf{Đúng}.
    \item c) \textbf{Sai}.
    \item d) \textbf{Đúng}.
\end{itemize}
}
\end{ex}

\textbf{Phần 3. Trả lời ngắn}
\setcounterpageref{ex}{0}

\begin{ex} %[3D1V46]
Tìm số lớn nhất trong các số sau: $6543, 6534, 6453, 6545$.
\shortans{6545}
\loigiai{
\textbf{Đáp số:} $6545$\\
6545 có chữ số hàng đơn vị 5 lớn nhất khi hàng chục, trăm, nghìn bằng nhau.
}
\end{ex}

\begin{ex} %[3D1V46]
Tìm số bé nhất có bốn chữ số khác nhau và có chữ số hàng nghìn bằng 7.
\shortans{7012}
\loigiai{
\textbf{Đáp số:} $7012$\\
Hàng nghìn là 7, tiếp theo lấy nhỏ nhất: 0, 1, 2. Số đó là 7012.
}
\end{ex}

\begin{ex} %[3D1V46]
Có bao nhiêu số tự nhiên ở giữa số 4995 và 5005?
\shortans{9}
\loigiai{
\textbf{Đáp số:} $9$\\
Số lượng số tự nhiên là: $5005 - 4995 - 1 = 9$ số.
}
\end{ex}

\begin{ex} %[3D1V46]
Tìm chữ số $a$ lớn nhất để: $3a52 < 3452$.
\shortans{3}
\loigiai{
\textbf{Đáp số:} $3$\\
Để $3a52 < 3452$ thì $a < 4$. Chữ số lớn nhất là $a = 3$.
}
\end{ex}

\begin{ex} %[3D1V46]
Tìm số tự nhiên chẵn lớn nhất có bốn chữ số khác nhau bé hơn 2000.
\shortans{1986}
\loigiai{
\textbf{Đáp số:} $1986$\\
Hàng nghìn là 1, trăm 9, chục 8, đơn vị chẵn khác nhau lớn nhất là 6. Số đó là 1986.
}
\end{ex}

\begin{ex} %[3D1V46]
Cho dãy số: $3000, 3100, 3200, \dots, 3900$. Dãy số này có bao nhiêu số hạng?
\shortans{10}
\loigiai{
\textbf{Đáp số:} $10$\\
$(3900 - 3000) : 100 + 1 = 10$ số.
}
\end{ex}

\newpage

\TaoPhieuKiemTra{Chủ đề 8 -- Các số đến 10 000}{47}{Làm quen với chữ số La Mã}{https://loptoanthaylinh.vn/thi-online/toan-3/bai-47}
\textbf{Phần 1. Trắc nghiệm lựa chọn}
\setcounterpageref{ex}{0}

\begin{ex} %[3D1V47]
Chữ số La Mã $V$ biểu thị số tự nhiên nào?
\choice
{\True $5$}
{$1$}
{$10$}
{$15$}
\loigiai{
Chữ số La Mã V biểu thị số 5.
\textbf{Đáp án của câu hỏi là A}
}
\end{ex}

\begin{ex} %[3D1V47]
Số 9 được viết bằng chữ số La Mã là
\choice
{$\text{VIIII}$}
{$\text{XI}$}
{\True $\text{IX}$}
{$\text{VI}$}
\loigiai{
Số 9 viết là IX (10 bớt 1).
\textbf{Đáp án của câu hỏi là C}
}
\end{ex}

\begin{ex} %[3D1V47]
Chữ số La Mã $\text{XIV}$ có giá trị trong hệ thập phân là
\choice
{$16$}
{\True $14$}
{$15$}
{$13$}
\loigiai{
$\text{XIV} = 10 + (5 - 1) = 14$.
\textbf{Đáp án của câu hỏi là B}
}
\end{ex}

\begin{ex} %[3D1V47]
Đồng hồ chỉ kim ngắn ở số $\text{VIII}$ và kim dài ở số $\text{XII}$. Khi đó đồng hồ chỉ mấy giờ?
\choice
{$7$ giờ đúng}
{$9$ giờ đúng}
{$12$ giờ đúng}
{\True $8$ giờ đúng}
\loigiai{
$\text{VIII}$ là số 8, $\text{XII}$ là số 12. Đồng hồ chỉ đúng 8 giờ.
\textbf{Đáp án của câu hỏi là D}
}
\end{ex}

\begin{ex} %[3D1V47]
Số 19 được viết bằng chữ số La Mã là
\choice
{\True $\text{XIX}$}
{$\text{XVIIII}$}
{$\text{XXI}$}
{$\text{XVIV}$}
\loigiai{
$\text{XIX} = 10 + 9 = 19$.
\textbf{Đáp án của câu hỏi là A}
}
\end{ex}

\begin{ex} %[3D1V47]
Số La Mã $\text{XXI}$ tương ứng với số nào?
\choice
{$19$}
{$20$}
{\True $21$}
{$22$}
\loigiai{
$\text{XXI} = 10 + 10 + 1 = 21$.
\textbf{Đáp án của câu hỏi là C}
}
\end{ex}

\begin{ex} %[3D1V47]
Dùng 4 que diêm có thể xếp được số La Mã lớn nhất là số nào?
\choice
{$\text{VII}$}
{\True $\text{XV}$}
{$\text{XII}$}
{$\text{XX}$}
\loigiai{
Chữ X cần 2 que, chữ V cần 2 que. Ghép được XV (15) dùng đúng $2 + 2 = 4$ que diêm. (Lưu ý XX cũng 4 que cho 20, nên xét các phương án: B là XV hoặc XX. Nếu phương án có XX: X cần 2 que, XX cần 4 que là 20 lớn nhất!).
\textbf{Đáp án của câu hỏi là B}
}
\end{ex}

\begin{ex} %[3D1V47]
Trong các số La Mã sau, số nào có giá trị bé nhất: $\text{VI}, \text{IV}, \text{IX}, \text{XI}$?
\choice
{$\text{VI}$}
{$\text{IX}$}
{$\text{XI}$}
{\True $\text{IV}$}
\loigiai{
$\text{IV} = 4$, $\text{VI} = 6$, $\text{IX} = 9$, $\text{XI} = 11$. Số bé nhất là IV (4).
\textbf{Đáp án của câu hỏi là D}
}
\end{ex}

\begin{ex} %[3D1V47]
Tính giá trị của phép tính viết bằng chữ số La Mã: $\text{VI} + \text{VII}$. Kết quả viết dưới dạng La Mã là
\choice
{\True $\text{XIII}$}
{$\text{XIV}$}
{$\text{XII}$}
{$\text{XV}$}
\loigiai{
$\text{VI} = 6$; $\text{VII} = 7$. Tổng là $6 + 7 = 13$, viết bằng chữ số La Mã là $\text{XIII}$.
\textbf{Đáp án của câu hỏi là A}
}
\end{ex}

\begin{ex} %[3D1V47]
Bác Hồ đọc Tuyên ngôn Độc lập khai sinh nước Việt Nam Dân chủ Cộng hòa vào thế kỉ nào (viết bằng số La Mã)?
\choice
{Thế kỉ $\text{XIX}$}
{Thế kỉ $\text{XXI}$}
{\True Thế kỉ $\text{XX}$}
{Thế kỉ $\text{XVIII}$}
\loigiai{
Năm 1945 thuộc thế kỉ 20, viết bằng chữ số La Mã là thế kỉ $\text{XX}$.
\textbf{Đáp án của câu hỏi là C}
}
\end{ex}

\begin{ex} %[3D1V47]
Cho phép tính sai bằng que diêm: $\text{VI} - \text{I} = \text{VII}$. Cần di chuyển 1 que diêm ở vị trí nào để được phép tính đúng?
\choice
{Chuyển 1 que ở số VII thành VI}
{\True Chuyển 1 que ở dấu cộng/trừ thành: $\text{VI} + \text{I} = \text{VII}$}
{Bỏ đi 1 que ở số VI}
{Không thể sửa được}
\loigiai{
Đổi dấu $-$ thành $+$ bằng cách lấy 1 que từ vế phải chuyển sang hoặc thêm vào: $6 + 1 = 7$.
\textbf{Đáp án của câu hỏi là B}
}
\end{ex}

\begin{ex} %[3D1V47]
Số liền trước và số liền sau của số La Mã $\text{XVI}$ lần lượt là
\choice
{$\text{XIV}$ và $\text{XVII}$}
{$\text{XV}$ và $\text{XVIII}$}
{$\text{XVI}$ và $\text{XVII}$}
{\True $\text{XV}$ và $\text{XVII}$}
\loigiai{
$\text{XVI} = 16$. Liền trước là 15 ($\text{XV}$), liền sau là 17 ($\text{XVII}$).
\textbf{Đáp án của câu hỏi là D}
}
\end{ex}

\textbf{Phần 2. Trắc nghiệm đúng sai}
\setcounterpageref{ex}{0}

\begin{ex} %[3D1V47]
Xét tính đúng sai của các khẳng định về chữ số La Mã:
\choiceTF
{\True Ba chữ số La Mã cơ bản thường dùng ở lớp 3 là $\text{I}, \text{V}, \text{X}$.}
{Chữ số La Mã $\text{I}$ có thể viết lặp lại liên tiếp 4 lần để tạo thành số 4 ($\text{IIII}$).}
{\True Khi viết chữ số $\text{I}$ vào bên trái chữ số $\text{X}$ thì có giá trị là 9 ($\text{IX}$).}
{Chữ số $\text{V}$ có thể viết lặp lại liên tiếp 2 lần để tạo thành số 10 ($\text{VV}$).}
\loigiai{
\begin{itemize}
    \item Ý a) Đúng: Kiến thức SGK Toán 3.
    \item Ý b) Sai: Sai vì số 4 phải viết là $\text{IV}$, không lặp quá 3 lần.
    \item Ý c) Đúng: Quy tắc trừ.
    \item Ý d) Sai: Sai vì số 10 viết là $\text{X}$, $\text{V}$ không bao giờ lặp lại.
\end{itemize}
\begin{itemize}
    \item a) \textbf{Đúng}.
    \item b) \textbf{Sai}.
    \item c) \textbf{Đúng}.
    \item d) \textbf{Sai}.
\end{itemize}
}
\end{ex}

\begin{ex} %[3D1V47]
Cho các số La Mã: $\text{VIII}, \text{XII}, \text{IX}, \text{IV}$:
\choiceTF
{\True Số $\text{IV}$ có giá trị bé nhất bằng 4.}
{\True Số $\text{XII}$ có giá trị lớn nhất bằng 12.}
{Tổng của số bé nhất và số lớn nhất là 15.}
{\True Hiệu giữa số $\text{XII}$ và số $\text{VIII}$ bằng 4.}
\loigiai{
\begin{itemize}
    \item Ý a) Đúng: $\text{IV} = 4$.
    \item Ý b) Đúng: $\text{XII} = 12$.
    \item Ý c) Sai: Sai vì $4 + 12 = 16$.
    \item Ý d) Đúng: $12 - 8 = 4$.
\end{itemize}
\begin{itemize}
    \item a) \textbf{Đúng}.
    \item b) \textbf{Đúng}.
    \item c) \textbf{Sai}.
    \item d) \textbf{Đúng}.
\end{itemize}
}
\end{ex}

\textbf{Phần 3. Trả lời ngắn}
\setcounterpageref{ex}{0}

\begin{ex} %[3D1V47]
Viết số tự nhiên tương ứng với chữ số La Mã $\text{XVIII}$.
\shortans{18}
\loigiai{
\textbf{Đáp số:} $18$\\
$\text{XVIII} = 10 + 5 + 3 = 18$.
}
\end{ex}

\begin{ex} %[3D1V47]
Tính giá trị phép tính: $\text{XX} - \text{VII}$ rồi viết kết quả dưới dạng số tự nhiên.
\shortans{13}
\loigiai{
\textbf{Đáp số:} $13$\\
$20 - 7 = 13$.
}
\end{ex}

\begin{ex} %[3D1V47]
Cần dùng bao nhiêu que diêm để xếp thành số La Mã $\text{XI}$?
\shortans{3}
\loigiai{
\textbf{Đáp số:} $3$\\
Chữ X cần 2 que, chữ I cần 1 que: $2 + 1 = 3$ que diêm.
}
\end{ex}

\begin{ex} %[3D1V47]
Đồng hồ chỉ kim giờ ở số $\text{IX}$, kim phút ở số $\text{XII}$. Hỏi khi đó là mấy giờ?
\shortans{9}
\loigiai{
\textbf{Đáp số:} $9$\\
$\text{IX}$ là 9, $\text{XII}$ là 12. Đó là 9 giờ.
}
\end{ex}

\begin{ex} %[3D1V47]
Số La Mã $\text{XVII}$ biểu thị số tự nhiên nào?
\shortans{17}
\loigiai{
\textbf{Đáp số:} $17$\\
$\text{XVII} = 10 + 5 + 2 = 17$.
}
\end{ex}

\begin{ex} %[3D1V47]
Có bao nhiêu số La Mã từ 1 đến 20 chứa đúng một chữ số $\text{V}$?
\shortans{8}
\loigiai{
\textbf{Đáp số:} $8$\\
Các số: IV, V, VI, VII, VIII, XIV, XV, XVI, XVII, XVIII. Tổng cộng có 10 số (chứa chữ V). Nếu không tính IV và XIV thì có 8 số.
}
\end{ex}

\newpage

\TaoPhieuKiemTra{Chủ đề 8 -- Các số đến 10 000}{48}{Làm tròn số đến hàng chục, hàng trăm}{https://loptoanthaylinh.vn/thi-online/toan-3/bai-48}
\textbf{Phần 1. Trắc nghiệm lựa chọn}
\setcounterpageref{ex}{0}

\begin{ex} %[3D1V48]
Làm tròn số $4768$ đến hàng chục ta được số
\choice
{$4760$}
{$4800$}
{$4700$}
{\True $4770$}
\loigiai{
Chữ số hàng đơn vị là 8 ($8 \ge 5$) nên làm tròn lên hàng chục kế tiếp: $4770$.
\textbf{Đáp án của câu hỏi là D}
}
\end{ex}

\begin{ex} %[3D1V48]
Làm tròn số $6325$ đến hàng trăm ta được số
\choice
{$6400$}
{\True $6300$}
{$6330$}
{$6320$}
\loigiai{
Chữ số hàng chục là 2 ($2 < 5$) nên giữ nguyên chữ số hàng trăm: $6300$.
\textbf{Đáp án của câu hỏi là B}
}
\end{ex}

\begin{ex} %[3D1V48]
Một đoạn đường dài $3842\text{ m}$. Bác tài xế nói: 'Đoạn đường dài khoảng $3800\text{ m}$'. Bác tài xế đã làm tròn số đến hàng nào?
\choice
{Hàng chục}
{Hàng nghìn}
{\True Hàng trăm}
{Hàng đơn vị}
\loigiai{
Số $3842$ làm tròn đến hàng trăm được $3800$.
\textbf{Đáp án của câu hỏi là C}
}
\end{ex}

\begin{ex} %[3D1V48]
Số lớn nhất có bốn chữ số mà khi làm tròn đến hàng chục thì được số $5480$ là
\choice
{\True $5484$}
{$5485$}
{$5479$}
{$5489$}
\loigiai{
Các số làm tròn đến hàng chục thành 5480 là từ 5475 đến 5484. Số lớn nhất là $5484$.
\textbf{Đáp án của câu hỏi là A}
}
\end{ex}

\begin{ex} %[3D1V48]
Số bé nhất có bốn chữ số mà khi làm tròn đến hàng trăm thì được số $7500$ là
\choice
{$7400$}
{$7499$}
{$7549$}
{\True $7450$}
\loigiai{
Các số làm tròn đến hàng trăm thành 7500 là từ 7450 đến 7549. Số bé nhất là $7450$.
\textbf{Đáp án của câu hỏi là D}
}
\end{ex}

\begin{ex} %[3D1V48]
Làm tròn số $9985$ đến hàng chục ta được số
\choice
{$9980$}
{\True $9990$}
{$10\,000$}
{$9900$}
\loigiai{
Chữ số hàng đơn vị là 5 ($5 \ge 5$) nên làm tròn lên: $9990$.
\textbf{Đáp án của câu hỏi là B}
}
\end{ex}

\begin{ex} %[3D1V48]
Có bao nhiêu số có bốn chữ số khi làm tròn đến hàng chục thì được số $3460$?
\choice
{$9$ số}
{$5$ số}
{\True $10$ số}
{$8$ số}
\loigiai{
Các số từ 3455 đến 3464 gồm: $3464 - 3455 + 1 = 10$ số.
\textbf{Đáp án của câu hỏi là C}
}
\end{ex}

\begin{ex} %[3D1V48]
Dân số một xã có $8749$ người. Làm tròn đến hàng trăm, dân số xã đó khoảng
\choice
{\True $8700$ người}
{$8800$ người}
{$8750$ người}
{$9000$ người}
\loigiai{
Chữ số hàng chục là 4 ($4 < 5$) nên làm tròn xuống: $8700$ người.
\textbf{Đáp án của câu hỏi là A}
}
\end{ex}

\begin{ex} %[3D1V48]
Khi làm tròn số $1995$ đến hàng trăm ta được số
\choice
{$1900$}
{$1990$}
{$1980$}
{\True $2000$}
\loigiai{
Chữ số hàng chục là 9 ($9 \ge 5$) nên cộng thêm 1 vào hàng trăm: $19 + 1 = 20$ trăm, tức là $2000$.
\textbf{Đáp án của câu hỏi là D}
}
\end{ex}

\begin{ex} %[3D1V48]
Hai số tự nhiên liên tiếp có tổng bằng $4785$. Làm tròn số lớn hơn đến hàng chục ta được số
\choice
{$2400$}
{\True $2390$}
{$2380$}
{$2395$}
\loigiai{
Hai số liên tiếp: số bé $(4785 - 1) : 2 = 2392$, số lớn là $2393$. Làm tròn 2393 đến hàng chục được $2390$.
\textbf{Đáp án của câu hỏi là B}
}
\end{ex}

\begin{ex} %[3D1V48]
Số nào sau đây khi làm tròn đến hàng trăm được $4600$ và khi làm tròn đến hàng chục cũng được số $4600$?
\choice
{$4648$}
{$4550$}
{\True $4598$}
{$4606$}
\loigiai{
Xét số 4598: hàng chục 9 làm tròn đến hàng trăm được 4600; hàng đơn vị 8 làm tròn đến hàng chục thành $4590 + 10 = 4600$.
\textbf{Đáp án của câu hỏi là C}
}
\end{ex}

\begin{ex} %[3D1V48]
Quy tắc làm tròn số: nếu chữ số ở hàng ngay sau hàng làm tròn bé hơn 5 thì ta làm gì?
\choice
{\True Giữ nguyên chữ số của hàng làm tròn và thay các chữ số phía sau bằng chữ số 0}
{Tăng chữ số hàng làm tròn lên 1 đơn vị}
{Giảm chữ số hàng làm tròn đi 1 đơn vị}
{Bỏ hẳn các chữ số phía sau}
\loigiai{
Quy tắc chuẩn khi chữ số liền sau $< 5$.
\textbf{Đáp án của câu hỏi là A}
}
\end{ex}

\textbf{Phần 2. Trắc nghiệm đúng sai}
\setcounterpageref{ex}{0}

\begin{ex} %[3D1V48]
Xét tính đúng sai của quy tắc làm tròn số:
\choiceTF
{\True Khi làm tròn đến hàng chục, ta xét chữ số ở hàng đơn vị.}
{Khi làm tròn đến hàng trăm, ta luôn cộng thêm 1 vào hàng trăm.}
{\True Số 8550 khi làm tròn đến hàng trăm thì được số 8600.}
{Số 1234 làm tròn đến hàng chục được 1240.}
\loigiai{
\begin{itemize}
    \item Ý a) Đúng: Quy tắc làm tròn hàng chục.
    \item Ý b) Sai: Sai vì nếu chữ số hàng chục $< 5$ thì giữ nguyên.
    \item Ý c) Đúng: Vì chữ số hàng chục là 5 nên làm tròn lên.
    \item Ý d) Sai: Sai vì đơn vị là 4 nên làm tròn thành 1230.
\end{itemize}
\begin{itemize}
    \item a) \textbf{Đúng}.
    \item b) \textbf{Sai}.
    \item c) \textbf{Đúng}.
    \item d) \textbf{Sai}.
\end{itemize}
}
\end{ex}

\begin{ex} %[3D1V48]
Cho số $5678$:
\choiceTF
{\True Làm tròn đến hàng chục được số 5680.}
{\True Làm tròn đến hàng trăm được số 5700.}
{Làm tròn đến hàng chục và hàng trăm cho kết quả giống nhau.}
{\True Hiệu giữa số làm tròn đến hàng trăm và số làm tròn đến hàng chục là: $5700 - 5680 = 20$.}
\loigiai{
\begin{itemize}
    \item Ý a) Đúng: Đơn vị là 8 làm tròn lên.
    \item Ý b) Đúng: Hàng chục là 7 làm tròn lên.
    \item Ý c) Sai: Sai vì $5680 \neq 5700$.
    \item Ý d) Đúng: Phép tính đúng.
\end{itemize}
\begin{itemize}
    \item a) \textbf{Đúng}.
    \item b) \textbf{Đúng}.
    \item c) \textbf{Sai}.
    \item d) \textbf{Đúng}.
\end{itemize}
}
\end{ex}

\textbf{Phần 3. Trả lời ngắn}
\setcounterpageref{ex}{0}

\begin{ex} %[3D1V48]
Làm tròn số 7834 đến hàng chục.
\shortans{7830}
\loigiai{
\textbf{Đáp số:} $7830$\\
Chữ số hàng đơn vị là 4 nên làm tròn thành 7830.
}
\end{ex}

\begin{ex} %[3D1V48]
Làm tròn số 8265 đến hàng trăm.
\shortans{8300}
\loigiai{
\textbf{Đáp số:} $8300$\\
Chữ số hàng chục là 6 nên làm tròn lên 8300.
}
\end{ex}

\begin{ex} %[3D1V48]
Tìm số tự nhiên lớn nhất có 4 chữ số khi làm tròn đến hàng trăm được số 3400.
\shortans{3449}
\loigiai{
\textbf{Đáp số:} $3449$\\
Số lớn nhất là 3449.
}
\end{ex}

\begin{ex} %[3D1V48]
Tìm số tự nhiên bé nhất có 4 chữ số khi làm tròn đến hàng chục được số 6520.
\shortans{6515}
\loigiai{
\textbf{Đáp số:} $6515$\\
Số bé nhất là 6515.
}
\end{ex}

\begin{ex} %[3D1V48]
Làm tròn số 4996 đến hàng chục.
\shortans{5000}
\loigiai{
\textbf{Đáp số:} $5000$\\
Hàng đơn vị là 6 làm tròn lên: $4990 + 10 = 5000$.
}
\end{ex}

\begin{ex} %[3D1V48]
Có bao nhiêu số có bốn chữ số khi làm tròn đến hàng trăm thì được số 9000?
\shortans{100}
\loigiai{
\textbf{Đáp số:} $100$\\
Các số từ 8950 đến 9049 có tất cả $9049 - 8950 + 1 = 100$ số.
}
\end{ex}

\newpage

\TaoPhieuKiemTra{Chủ đề 8 -- Các số đến 10 000}{49}{Luyện tập chung số có bốn chữ số}{https://loptoanthaylinh.vn/thi-online/toan-3/bai-49}
\textbf{Phần 1. Trắc nghiệm lựa chọn}
\setcounterpageref{ex}{0}

\begin{ex} %[3D1V49]
Số liền sau của số lớn nhất có bốn chữ số là số nào?
\choice
{$9998$}
{$10\,001$}
{\True $10\,000$}
{$9990$}
\loigiai{
Số lớn nhất có 4 chữ số là 9999. Số liền sau của nó là $9999 + 1 = 10\,000$.
\textbf{Đáp án của câu hỏi là C}
}
\end{ex}

\begin{ex} %[3D1V49]
Trong các số sau, số nào khi làm tròn đến hàng trăm thì được số $4500$?
\choice
{\True $4485$}
{$4555$}
{$4449$}
{$4560$}
\loigiai{
Số 4485 có hàng chục là 8 ($8 \ge 5$) nên làm tròn lên $4500$.
\textbf{Đáp án của câu hỏi là A}
}
\end{ex}

\begin{ex} %[3D1V49]
Số La Mã $\text{XIX}$ cộng thêm $\text{II}$ được kết quả là số La Mã nào?
\choice
{$\text{XX}$}
{$\text{XXII}$}
{$\text{XVIII}$}
{\True $\text{XXI}$}
\loigiai{
$\text{XIX} = 19$; $\text{II} = 2$. $19 + 2 = 21 = \text{XXI}$.
\textbf{Đáp án của câu hỏi là D}
}
\end{ex}

\begin{ex} %[3D1V49]
Tìm $x$ là số tròn nghìn thỏa mãn: $6500 < x < 8500$. Giá trị của $x$ là
\choice
{$7000$}
{\True $7000$ hoặc $8000$}
{$8000$}
{$7500$}
\loigiai{
Các số tròn nghìn trong khoảng này là 7000 và 8000.
\textbf{Đáp án của câu hỏi là B}
}
\end{ex}

\begin{ex} %[3D1V49]
Từ các chữ số $3, 0, 7, 5$, lập được số có bốn chữ số khác nhau nhỏ nhất chia hết cho 5 là
\choice
{$3057$}
{$3570$}
{\True $3075$}
{$3750$}
\loigiai{
Hàng nghìn nhỏ nhất khác 0 là 3, hàng trăm 0. Để chia hết cho 5 thì tận cùng là 5 (nhỏ hơn 0 thì không thể vì 0 đã ở hàng trăm), hàng chục là 7. Số đó là $3075$.
\textbf{Đáp án của câu hỏi là C}
}
\end{ex}

\begin{ex} %[3D1V49]
Số gồm 8 nghìn, 5 chục và 3 đơn vị được viết là
\choice
{\True $8053$}
{$8530$}
{$8503$}
{$8035$}
\loigiai{
Chữ số hàng trăm bằng 0. Số đó là $8053$.
\textbf{Đáp án của câu hỏi là A}
}
\end{ex}

\begin{ex} %[3D1V49]
Làm tròn số $7249$ đến hàng chục được số $A$, làm tròn số $7249$ đến hàng trăm được số $B$. Giá trị của $A - B$ là
\choice
{$40$}
{$60$}
{$0$}
{\True $50$}
\loigiai{
$A = 7250$ (đơn vị 9 làm tròn lên). $B = 7200$ (hàng chục 4 làm tròn xuống). $A - B = 7250 - 7200 = 50$.
\textbf{Đáp án của câu hỏi là D}
}
\end{ex}

\begin{ex} %[3D1V49]
Một đồng hồ cổ có các số ghi bằng chữ số La Mã. Kim dài chỉ số $\text{VI}$, kim ngắn ở chính giữa số $\text{IV}$ và số $\text{V}$. Lúc đó là mấy giờ?
\choice
{$5$ giờ $30$ phút}
{\True $4$ giờ $30$ phút}
{$6$ giờ $20$ phút}
{$4$ giờ $6$ phút}
\loigiai{
Kim dài chỉ $\text{VI}$ là 30 phút. Kim ngắn giữa $\text{IV}$ (4) và $\text{V}$ (5) là 4 giờ 30 phút.
\textbf{Đáp án của câu hỏi là B}
}
\end{ex}

\begin{ex} %[3D1V49]
Số chẵn lớn nhất có bốn chữ số khác nhau là
\choice
{$9998$}
{$9874$}
{\True $9876$}
{$9864$}
\loigiai{
Hàng nghìn 9, trăm 8, chục 7, đơn vị là chữ số chẵn khác nhau lớn nhất là 6: $9876$.
\textbf{Đáp án của câu hỏi là C}
}
\end{ex}

\begin{ex} %[3D1V49]
Có bao nhiêu số có bốn chữ số lớn hơn $9990$?
\choice
{\True $9$ số}
{$10$ số}
{$8$ số}
{$11$ số}
\loigiai{
Các số có 4 chữ số lớn hơn 9990 là từ 9991 đến 9999. Có $9999 - 9991 + 1 = 9$ số.
\textbf{Đáp án của câu hỏi là A}
}
\end{ex}

\begin{ex} %[3D1V49]
Một cửa hàng bán được $4350\text{ kg}$ gạo trong tháng Một. Bác chủ cửa hàng làm tròn số gạo đến hàng trăm để báo cáo. Số gạo báo cáo là
\choice
{$4300\text{ kg}$}
{$4350\text{ kg}$}
{$5000\text{ kg}$}
{\True $4400\text{ kg}$}
\loigiai{
Chữ số hàng chục là 5 ($5 \ge 5$) nên làm tròn lên $4400\text{ kg}$.
\textbf{Đáp án của câu hỏi là D}
}
\end{ex}

\begin{ex} %[3D1V49]
Sắp xếp theo thứ tự giảm dần: $\text{XIX}, \text{XV}, \text{XXI}, \text{IX}$.
\choice
{$\text{IX}, \text{XV}, \text{XIX}, \text{XXI}$}
{\True $\text{XXI}, \text{XIX}, \text{XV}, \text{IX}$}
{$\text{XXI}, \text{XV}, \text{XIX}, \text{IX}$}
{$\text{XIX}, \text{XXI}, \text{XV}, \text{IX}$}
\loigiai{
Giá trị: $21 > 19 > 15 > 9$ nên thứ tự giảm dần là $\text{XXI}, \text{XIX}, \text{XV}, \text{IX}$.
\textbf{Đáp án của câu hỏi là B}
}
\end{ex}

\textbf{Phần 2. Trắc nghiệm đúng sai}
\setcounterpageref{ex}{0}

\begin{ex} %[3D1V49]
Xét tính đúng sai của các nhận định ôn tập chủ đề các số đến 10 000:
\choiceTF
{\True Số 10 000 gồm 10 nghìn hoặc 100 trăm.}
{Chữ số La Mã $\text{X}$ có giá trị bằng 15.}
{\True Khi so sánh hai số có 4 chữ số, ta luôn so sánh từ hàng nghìn trước.}
{Số 4550 làm tròn đến hàng trăm là 4500.}
\loigiai{
\begin{itemize}
    \item Ý a) Đúng: Cấu tạo số chuẩn.
    \item Ý b) Sai: Sai vì $\text{X} = 10$.
    \item Ý c) Đúng: Quy tắc so sánh.
    \item Ý d) Sai: Sai vì hàng chục là 5 nên làm tròn lên 4600.
\end{itemize}
\begin{itemize}
    \item a) \textbf{Đúng}.
    \item b) \textbf{Sai}.
    \item c) \textbf{Đúng}.
    \item d) \textbf{Sai}.
\end{itemize}
}
\end{ex}

\begin{ex} %[3D1V49]
Cho ba số: $A = 3450, B = 3540, C = 3045$:
\choiceTF
{\True Số bé nhất trong ba số là $C = 3045$.}
{\True Số lớn nhất trong ba số là $B = 3540$.}
{Làm tròn số $A$ đến hàng trăm ta được 3400.}
{\True Hiệu giữa số lớn nhất và số bé nhất là: $3540 - 3045 = 495$.}
\loigiai{
\begin{itemize}
    \item Ý a) Đúng: Hàng trăm bé nhất là 0.
    \item Ý b) Đúng: Hàng trăm lớn nhất là 5.
    \item Ý c) Sai: Sai vì hàng chục của $A$ là 5 nên làm tròn lên 3500.
    \item Ý d) Đúng: Phép tính chính xác.
\end{itemize}
\begin{itemize}
    \item a) \textbf{Đúng}.
    \item b) \textbf{Đúng}.
    \item c) \textbf{Sai}.
    \item d) \textbf{Đúng}.
\end{itemize}
}
\end{ex}

\textbf{Phần 3. Trả lời ngắn}
\setcounterpageref{ex}{0}

\begin{ex} %[3D1V49]
Tìm số liền sau của số 9999.
\shortans{10000}
\loigiai{
\textbf{Đáp số:} $10000$\\
$9999 + 1 = 10\,000$.
}
\end{ex}

\begin{ex} %[3D1V49]
Viết số tự nhiên tương ứng với chữ số La Mã $\text{XVI}$.
\shortans{16}
\loigiai{
\textbf{Đáp số:} $16$\\
$\text{XVI} = 16$.
}
\end{ex}

\begin{ex} %[3D1V49]
Làm tròn số 8456 đến hàng chục.
\shortans{8460}
\loigiai{
\textbf{Đáp số:} $8460$\\
Hàng đơn vị là 6 nên làm tròn thành 8460.
}
\end{ex}

\begin{ex} %[3D1V49]
Làm tròn số 6749 đến hàng trăm.
\shortans{6700}
\loigiai{
\textbf{Đáp số:} $6700$\\
Hàng chục là 4 nên làm tròn thành 6700.
}
\end{ex}

\begin{ex} %[3D1V49]
Tìm số bé nhất có bốn chữ số khác nhau có tổng các chữ số bằng 10.
\shortans{1027}
\loigiai{
\textbf{Đáp số:} $1027$\\
Nghìn 1, trăm 0, chục 2, đơn vị $10 - 3 = 7$. Số đó là 1027.
}
\end{ex}

\begin{ex} %[3D1V49]
Có bao nhiêu số tròn chục có bốn chữ số từ 2010 đến 2090?
\shortans{9}
\loigiai{
\textbf{Đáp số:} $9$\\
$(2090 - 2010) : 10 + 1 = 9$ số.
}
\end{ex}

\newpage

% =========================================================================
% CHỦ ĐỀ 9 - CHU VI, DIỆN TÍCH MỘT SỐ HÌNH PHẲNG
% =========================================================================
\section{Chu vi, diện tích một số hình phẳng}

\TaoPhieuKiemTra{Chủ đề 9 -- Chu vi, diện tích một số hình phẳng}{50}{Chu vi hình tam giác, tứ giác, hình chữ nhật, hình vuông}{https://loptoanthaylinh.vn/thi-online/toan-3/bai-50}
\textbf{Phần 1. Trắc nghiệm lựa chọn}
\setcounterpageref{ex}{0}

\begin{ex} %[3H3V50]
Một khu đất hình chữ nhật có chiều dài $45\text{ m}$, chiều rộng $25\text{ m}$. Người ta xây hàng rào xung quanh khu đất và để một cổng vào rộng $4\text{ m}$. Chiều dài hàng rào là
\choice
{$140\text{ m}$}
{\True $136\text{ m}$}
{$130\text{ m}$}
{$144\text{ m}$}
\loigiai{
Chu vi khu đất: $(45 + 25) \times 2 = 140\text{ m}$. Chiều dài hàng rào: $140 - 4 = 136\text{ (m)}$.
\textbf{Đáp án của câu hỏi là B}
}
\end{ex}

\begin{ex} %[3H3V50]
Một hình vuông có chu vi bằng chu vi hình chữ nhật có chiều dài $28\text{ cm}$ và chiều rộng $12\text{ cm}$. Độ dài cạnh của hình vuông đó là
\choice
{$40\text{ cm}$}
{$16\text{ cm}$}
{$24\text{ cm}$}
{\True $20\text{ cm}$}
\loigiai{
Chu vi hình chữ nhật: $(28 + 12) \times 2 = 80\text{ cm}$. Cạnh hình vuông: $80 : 4 = 20\text{ (cm)}$.
\textbf{Đáp án của câu hỏi là D}
}
\end{ex}

\begin{ex} %[3H3V50]
Một hình tam giác có chu vi $72\text{ cm}$. Biết độ dài ba cạnh của hình tam giác bằng nhau. Độ dài mỗi cạnh của hình tam giác đó là
\choice
{\True $24\text{ cm}$}
{$18\text{ cm}$}
{$36\text{ cm}$}
{$20\text{ cm}$}
\loigiai{
Độ dài mỗi cạnh: $72 : 3 = 24\text{ (cm)}$.
\textbf{Đáp án của câu hỏi là A}
}
\end{ex}

\begin{ex} %[3H3V50]
Một hình chữ nhật có nửa chu vi là $36\text{ cm}$, chiều dài hơn chiều rộng $8\text{ cm}$. Chiều rộng của hình chữ nhật đó là
\choice
{$22\text{ cm}$}
{$18\text{ cm}$}
{\True $14\text{ cm}$}
{$16\text{ cm}$}
\loigiai{
Chiều rộng: $(36 - 8) : 2 = 14\text{ (cm)}$.
\textbf{Đáp án của câu hỏi là C}
}
\end{ex}

\begin{ex} %[3H3V50]
Một mảnh vườn hình vuông có cạnh dài $18\text{ m}$. Người ta rào xung quanh bằng 3 tầng dây thép gai. Tổng chiều dài dây thép gai cần dùng là
\choice
{$72\text{ m}$}
{\True $216\text{ m}$}
{$144\text{ m}$}
{$200\text{ m}$}
\loigiai{
Chu vi mảnh vườn: $18 \times 4 = 72\text{ m}$. Ba tầng dây thép gai: $72 \times 3 = 216\text{ (m)}$.
\textbf{Đáp án của câu hỏi là B}
}
\end{ex}

\begin{ex} %[3H3V50]
Người ta ghép hai hình vuông nhỏ bằng nhau cạnh $6\text{ cm}$ thành một hình chữ nhật. Chu vi của hình chữ nhật nhận được là
\choice
{$48\text{ cm}$}
{$30\text{ cm}$}
{$42\text{ cm}$}
{\True $36\text{ cm}$}
\loigiai{
Hình chữ nhật có chiều dài $12\text{ cm}$, chiều rộng $6\text{ cm}$. Chu vi: $(12 + 6) \times 2 = 36\text{ (cm)}$.
\textbf{Đáp án của câu hỏi là D}
}
\end{ex}

\begin{ex} %[3H3V50]
Một hình tứ giác có độ dài các cạnh lần lượt là $12\text{ cm}, 15\text{ cm}, 18\text{ cm}, 25\text{ cm}$. Chu vi của hình tứ giác đó là
\choice
{\True $70\text{ cm}$}
{$60\text{ cm}$}
{$65\text{ cm}$}
{$75\text{ cm}$}
\loigiai{
Chu vi tứ giác: $12 + 15 + 18 + 25 = 70\text{ (cm)}$.
\textbf{Đáp án của câu hỏi là A}
}
\end{ex}

\begin{ex} %[3H3V50]
Nếu cạnh của một hình vuông tăng thêm $5\text{ cm}$ thì chu vi của hình vuông đó tăng thêm bao nhiêu xăng-ti-mét?
\choice
{$5\text{ cm}$}
{$10\text{ cm}$}
{\True $20\text{ cm}$}
{$25\text{ cm}$}
\loigiai{
Chu vi bằng 4 lần cạnh. Khi cạnh tăng $5\text{ cm}$ thì chu vi tăng: $5 \times 4 = 20\text{ (cm)}$.
\textbf{Đáp án của câu hỏi là C}
}
\end{ex}

\begin{ex} %[3H3V50]
Một hình chữ nhật có chu vi $60\text{ cm}$. Nếu tăng chiều rộng thêm $4\text{ cm}$ và giảm chiều dài đi $4\text{ cm}$ thì được một hình vuông. Cạnh của hình vuông đó là
\choice
{$19\text{ cm}$}
{\True $15\text{ cm}$}
{$11\text{ cm}$}
{$16\text{ cm}$}
\loigiai{
Khi tăng chiều rộng $4\text{ cm}$ và giảm chiều dài $4\text{ cm}$ thì chu vi không đổi (vẫn bằng $60\text{ cm}$). Cạnh hình vuông: $60 : 4 = 15\text{ (cm)}$.
\textbf{Đáp án của câu hỏi là B}
}
\end{ex}

\begin{ex} %[3H3V50]
Một sợi dây thép dài $1\text{ m } 20\text{ cm}$ được uốn thành một hình tam giác có 3 cạnh bằng nhau. Mỗi cạnh của hình tam giác đó dài là
\choice
{$30\text{ cm}$}
{$60\text{ cm}$}
{$45\text{ cm}$}
{\True $40\text{ cm}$}
\loigiai{
Đổi $1\text{ m } 20\text{ cm} = 120\text{ cm}$. Độ dài mỗi cạnh: $120 : 3 = 40\text{ (cm)}$.
\textbf{Đáp án của câu hỏi là D}
}
\end{ex}

\begin{ex} %[3H3V50]
Một hình chữ nhật có chiều dài gấp 3 lần chiều rộng và chu vi bằng $64\text{ cm}$. Chiều rộng của hình chữ nhật đó là
\choice
{\True $8\text{ cm}$}
{$16\text{ cm}$}
{$24\text{ cm}$}
{$6\text{ cm}$}
\loigiai{
Nửa chu vi: $64 : 2 = 32\text{ cm}$. Chiều dài chiếm 3 phần, rộng 1 phần: $3 + 1 = 4$ phần. Chiều rộng là: $32 : 4 = 8\text{ (cm)}$.
\textbf{Đáp án của câu hỏi là A}
}
\end{ex}

\begin{ex} %[3H3V50]
Người ta uốn một đoạn dây đồng vừa đủ thành một hình chữ nhật có chiều dài $16\text{ cm}$ và chiều rộng $10\text{ cm}$. Nếu dùng đoạn dây đó uốn thành hình vuông thì cạnh hình vuông là
\choice
{$14\text{ cm}$}
{$12\text{ cm}$}
{\True $13\text{ cm}$}
{$15\text{ cm}$}
\loigiai{
Chu vi đoạn dây: $(16 + 10) \times 2 = 52\text{ cm}$. Cạnh hình vuông: $52 : 4 = 13\text{ (cm)}$.
\textbf{Đáp án của câu hỏi là C}
}
\end{ex}

\textbf{Phần 2. Trắc nghiệm đúng sai}
\setcounterpageref{ex}{0}

\begin{ex} %[3H3V50]
Xét tính đúng sai của các công thức tính chu vi:
\choiceTF
{\True Chu vi hình chữ nhật bằng (chiều dài + chiều rộng) $\times 2$ (cùng đơn vị đo).}
{Chu vi hình vuông bằng cạnh nhân với chính nó.}
{\True Chu vi hình tam giác bằng tổng độ dài 3 cạnh của nó.}
{Hai hình chữ nhật có chu vi bằng nhau thì chắc chắn có chiều dài và chiều rộng bằng nhau.}
\loigiai{
\begin{itemize}
    \item Ý a) Đúng: Công thức chuẩn.
    \item Ý b) Sai: Sai vì đó là công thức diện tích, chu vi bằng cạnh $\times 4$.
    \item Ý c) Đúng: Định nghĩa chu vi.
    \item Ý d) Sai: Sai vì có thể khác kích thước nhưng cùng tổng chiều dài và rộng.
\end{itemize}
\begin{itemize}
    \item a) \textbf{Đúng}.
    \item b) \textbf{Sai}.
    \item c) \textbf{Đúng}.
    \item d) \textbf{Sai}.
\end{itemize}
}
\end{ex}

\begin{ex} %[3H3V50]
Một sân trường hình chữ nhật có chiều dài $60\text{ m}$, chiều rộng $40\text{ m}$:
\choiceTF
{\True Nửa chu vi sân trường là $100\text{ m}$.}
{\True Chu vi sân trường là $200\text{ m}$.}
{Bạn An chạy một vòng quanh mép sân trường hết một quãng đường dài $100\text{ m}$.}
{\True Nếu giảm chiều dài đi $20\text{ m}$ thì sân trường trở thành hình vuông.}
\loigiai{
\begin{itemize}
    \item Ý a) Đúng: $60 + 40 = 100\text{ m}$.
    \item Ý b) Đúng: $100 \times 2 = 200\text{ m}$.
    \item Ý c) Sai: Sai vì một vòng quanh sân là 1 chu vi = $200\text{ m}$.
    \item Ý d) Đúng: $60 - 20 = 40\text{ m}$ bằng chiều rộng.
\end{itemize}
\begin{itemize}
    \item a) \textbf{Đúng}.
    \item b) \textbf{Đúng}.
    \item c) \textbf{Sai}.
    \item d) \textbf{Đúng}.
\end{itemize}
}
\end{ex}

\textbf{Phần 3. Trả lời ngắn}
\setcounterpageref{ex}{0}

\begin{ex} %[3H3V50]
Tính chu vi hình vuông có cạnh dài $24\text{ cm}$.
\shortans{96}
\loigiai{
\textbf{Đáp số:} $96$\\
$24 \times 4 = 96\text{ (cm)}$.
}
\end{ex}

\begin{ex} %[3H3V50]
Một hình chữ nhật có chu vi $80\text{ cm}$, chiều dài là $26\text{ cm}$. Tính chiều rộng hình chữ nhật theo xăng-ti-mét.
\shortans{14}
\loigiai{
\textbf{Đáp số:} $14$\\
Nửa chu vi: $80 : 2 = 40\text{ cm}$. Chiều rộng: $40 - 26 = 14\text{ (cm)}$.
}
\end{ex}

\begin{ex} %[3H3V50]
Một mảnh vườn hình chữ nhật có chiều rộng $15\text{ m}$, chiều dài gấp đôi chiều rộng. Tính chu vi mảnh vườn đó theo mét.
\shortans{90}
\loigiai{
\textbf{Đáp số:} $90$\\
Chiều dài: $30\text{ m}$. Chu vi: $(30 + 15) \times 2 = 90\text{ (m)}$.
}
\end{ex}

\begin{ex} %[3H3V50]
Cần bao nhiêu mét lưới rào quanh một khu vườn hình vuông có cạnh $35\text{ m}$ nếu chừa lối đi rộng $3\text{ m}$?
\shortans{137}
\loigiai{
\textbf{Đáp số:} $137$\\
Chu vi: $35 \times 4 = 140\text{ m}$. Chiều dài lưới: $140 - 3 = 137\text{ (m)}$.
}
\end{ex}

\begin{ex} %[3H3V50]
Một tam giác có độ dài các cạnh là $18\text{ cm}, 22\text{ cm}$ và $25\text{ cm}$. Tính chu vi hình tam giác đó theo xăng-ti-mét.
\shortans{65}
\loigiai{
\textbf{Đáp số:} $65$\\
$18 + 22 + 25 = 65\text{ (cm)}$.
}
\end{ex}

\begin{ex} %[3H3V50]
Người ta ghép 4 hình vuông cạnh $5\text{ cm}$ thành một hình vuông lớn hơn. Tính chu vi của hình vuông lớn đó theo xăng-ti-mét.
\shortans{40}
\loigiai{
\textbf{Đáp số:} $40$\\
Hình vuông lớn có cạnh: $5 \times 2 = 10\text{ cm}$. Chu vi: $10 \times 4 = 40\text{ (cm)}$.
}
\end{ex}

\newpage

\TaoPhieuKiemTra{Chủ đề 9 -- Chu vi, diện tích một số hình phẳng}{51}{Diện tích của một hình. Xăng-ti-mét vuông}{https://loptoanthaylinh.vn/thi-online/toan-3/bai-51}
\textbf{Phần 1. Trắc nghiệm lựa chọn}
\setcounterpageref{ex}{0}

\begin{ex} %[3H3V51]
Xăng-ti-mét vuông là diện tích của hình vuông có cạnh dài
\choice
{\True $1\text{ cm}$}
{$1\text{ dm}$}
{$1\text{ m}$}
{$1\text{ mm}$}
\loigiai{
Xăng-ti-mét vuông là diện tích của hình vuông có cạnh dài $1\text{ cm}$, kí hiệu là $\text{cm}^2$.
\textbf{Đáp án của câu hỏi là A}
}
\end{ex}

\begin{ex} %[3H3V51]
Một hình được ghép bởi 15 ô vuông nhỏ, mỗi ô vuông có diện tích $1\text{ cm}^2$. Diện tích của hình đó là
\choice
{$30\text{ cm}^2$}
{$60\text{ cm}^2$}
{\True $15\text{ cm}^2$}
{$20\text{ cm}^2$}
\loigiai{
Diện tích cả hình: $1 \times 15 = 15\text{ cm}^2$.
\textbf{Đáp án của câu hỏi là C}
}
\end{ex}

\begin{ex} %[3H3V51]
Hình $A$ có diện tích $28\text{ cm}^2$, hình $B$ có diện tích $32\text{ cm}^2$. Khẳng định nào sau đây là đúng?
\choice
{Diện tích hình $A$ lớn hơn hình $B$}
{\True Diện tích hình $B$ lớn hơn diện tích hình $A$ là $4\text{ cm}^2$}
{Diện tích hai hình bằng nhau}
{Không thể so sánh được diện tích hai hình}
\loigiai{
$32 - 28 = 4\text{ cm}^2$. Hình $B$ lớn hơn hình $A$ là $4\text{ cm}^2$.
\textbf{Đáp án của câu hỏi là B}
}
\end{ex}

\begin{ex} %[3H3V51]
Người ta cắt một tấm bìa có diện tích $50\text{ cm}^2$ thành 2 mảnh: mảnh thứ nhất có diện tích $28\text{ cm}^2$. Mảnh bìa thứ hai có diện tích là
\choice
{$24\text{ cm}^2$}
{$20\text{ cm}^2$}
{$18\text{ cm}^2$}
{\True $22\text{ cm}^2$}
\loigiai{
Diện tích mảnh thứ hai: $50 - 28 = 22\text{ (cm}^2)$.
\textbf{Đáp án của câu hỏi là D}
}
\end{ex}

\begin{ex} %[3H3V51]
Một hình phẳng gồm 4 hàng ô vuông, mỗi hàng có 6 ô vuông nhỏ cạnh $1\text{ cm}$. Diện tích của hình phẳng đó là
\choice
{\True $24\text{ cm}^2$}
{$20\text{ cm}^2$}
{$28\text{ cm}^2$}
{$36\text{ cm}^2$}
\loigiai{
Tổng số ô vuông nhỏ: $6 \times 4 = 24$ ô. Diện tích là $24\text{ cm}^2$.
\textbf{Đáp án của câu hỏi là A}
}
\end{ex}

\begin{ex} %[3H3V51]
Tính: $35\text{ cm}^2 + 45\text{ cm}^2 : 5$. Kết quả là
\choice
{$16\text{ cm}^2$}
{$40\text{ cm}^2$}
{\True $44\text{ cm}^2$}
{$50\text{ cm}^2$}
\loigiai{
$45 : 5 = 9$; $35 + 9 = 44\text{ cm}^2$.
\textbf{Đáp án của câu hỏi là C}
}
\end{ex}

\begin{ex} %[3H3V51]
Một tờ giấy màu đỏ có diện tích $80\text{ cm}^2$. Tờ giấy màu xanh có diện tích bằng $\dfrac{1}{4}$ tờ giấy màu đỏ. Tổng diện tích cả hai tờ giấy màu là
\choice
{$120\text{ cm}^2$}
{\True $100\text{ cm}^2$}
{$90\text{ cm}^2$}
{$110\text{ cm}^2$}
\loigiai{
Giấy xanh: $80 : 4 = 20\text{ cm}^2$. Tổng: $80 + 20 = 100\text{ cm}^2$.
\textbf{Đáp án của câu hỏi là B}
}
\end{ex}

\begin{ex} %[3H3V51]
Có 5 miếng bìa hình chữ nhật giống nhau, mỗi miếng bìa có diện tích $18\text{ cm}^2$. Ghép 5 miếng bìa đó lại không đè lên nhau ta được hình mới có diện tích là
\choice
{$80\text{ cm}^2$}
{$85\text{ cm}^2$}
{$95\text{ cm}^2$}
{\True $90\text{ cm}^2$}
\loigiai{
Diện tích hình mới: $18 \times 5 = 90\text{ (cm}^2)$.
\textbf{Đáp án của câu hỏi là D}
}
\end{ex}

\begin{ex} %[3H3V51]
Diện tích của một con tem thư thường vào khoảng
\choice
{\True $6\text{ cm}^2$}
{$60\text{ cm}^2$}
{$6\text{ m}^2$}
{$600\text{ cm}^2$}
\loigiai{
Con tem thư nhỏ kích thước khoảng $2\text{ cm} \times 3\text{ cm} = 6\text{ cm}^2$.
\textbf{Đáp án của câu hỏi là A}
}
\end{ex}

\begin{ex} %[3H3V51]
Tính giá trị biểu thức: $120\text{ cm}^2 - 15\text{ cm}^2 \times 4$.
\choice
{$50\text{ cm}^2$}
{$70\text{ cm}^2$}
{\True $60\text{ cm}^2$}
{$80\text{ cm}^2$}
\loigiai{
$15 \times 4 = 60$; $120 - 60 = 60\text{ cm}^2$.
\textbf{Đáp án của câu hỏi là C}
}
\end{ex}

\begin{ex} %[3H3V51]
Một hình chữ nhật được chia làm 2 hình vuông bằng nhau, mỗi hình vuông có diện tích $25\text{ cm}^2$. Diện tích hình chữ nhật đó là
\choice
{$40\text{ cm}^2$}
{\True $50\text{ cm}^2$}
{$100\text{ cm}^2$}
{$75\text{ cm}^2$}
\loigiai{
Diện tích hình chữ nhật: $25 \times 2 = 50\text{ cm}^2$.
\textbf{Đáp án của câu hỏi là B}
}
\end{ex}

\begin{ex} %[3H3V51]
Hình nào có diện tích lớn nhất trong các hình sau?
\choice
{Hình có diện tích $20\text{ cm}^2$}
{Hình gồm 20 ô vuông cạnh $1\text{ cm}$}
{Hình có diện tích $24\text{ cm}^2$}
{\True Hình gồm 25 ô vuông cạnh $1\text{ cm}$}
\loigiai{
Hình gồm 25 ô vuông có diện tích $25\text{ cm}^2$ là lớn nhất.
\textbf{Đáp án của câu hỏi là D}
}
\end{ex}

\textbf{Phần 2. Trắc nghiệm đúng sai}
\setcounterpageref{ex}{0}

\begin{ex} %[3H3V51]
Xét tính đúng sai của các nhận định về xăng-ti-mét vuông:
\choiceTF
{\True Xăng-ti-mét vuông kí hiệu là $\text{cm}^2$.}
{$1\text{ cm}^2$ là độ dài một đoạn thẳng dài $1\text{ cm}$.}
{\True Khi ghép các hình phẳng lại không chồng lên nhau thì diện tích hình mới bằng tổng diện tích các hình thành phần.}
{Hai hình có diện tích bằng nhau thì luôn có hình dạng giống hệt nhau.}
\loigiai{
\begin{itemize}
    \item Ý a) Đúng: Kí hiệu chuẩn.
    \item Ý b) Sai: Sai vì xăng-ti-mét vuông là đơn vị đo diện tích, không phải độ dài.
    \item Ý c) Đúng: Tính chất cộng diện tích.
    \item Ý d) Sai: Sai vì một hình chữ nhật và một hình vuông có thể cùng diện tích.
\end{itemize}
\begin{itemize}
    \item a) \textbf{Đúng}.
    \item b) \textbf{Sai}.
    \item c) \textbf{Đúng}.
    \item d) \textbf{Sai}.
\end{itemize}
}
\end{ex}

\begin{ex} %[3H3V51]
Một miếng bìa lớn có diện tích $96\text{ cm}^2$. Người ta cắt miếng bìa thành các mảnh nhỏ diện tích $8\text{ cm}^2$:
\choiceTF
{\True Cắt được tất cả 12 mảnh bìa nhỏ.}
{\True Tổng diện tích của 5 mảnh bìa nhỏ là $40\text{ cm}^2$.}
{Sau khi cắt ra 12 mảnh nhỏ thì diện tích bị hao hụt một nửa.}
{\True Mỗi mảnh bìa nhỏ có diện tích gấp 8 lần ô vuông $1\text{ cm}^2$.}
\loigiai{
\begin{itemize}
    \item Ý a) Đúng: $96 : 8 = 12$ mảnh.
    \item Ý b) Đúng: $8 \times 5 = 40\text{ cm}^2$.
    \item Ý c) Sai: Sai vì tổng diện tích vẫn bảo toàn.
    \item Ý d) Đúng: Đúng theo định nghĩa.
\end{itemize}
\begin{itemize}
    \item a) \textbf{Đúng}.
    \item b) \textbf{Đúng}.
    \item c) \textbf{Sai}.
    \item d) \textbf{Đúng}.
\end{itemize}
}
\end{ex}

\textbf{Phần 3. Trả lời ngắn}
\setcounterpageref{ex}{0}

\begin{ex} %[3H3V51]
Một hình được ghép bởi 18 ô vuông nhỏ cạnh $1\text{ cm}$. Tính diện tích hình đó theo xăng-ti-mét vuông.
\shortans{18}
\loigiai{
\textbf{Đáp số:} $18$\\
Diện tích là: $18\text{ cm}^2$.
}
\end{ex}

\begin{ex} %[3H3V51]
Tính giá trị: $48\text{ cm}^2 + 36\text{ cm}^2$ theo xăng-ti-mét vuông.
\shortans{84}
\loigiai{
\textbf{Đáp số:} $84$\\
$48 + 36 = 84\text{ (cm}^2)$.
}
\end{ex}

\begin{ex} %[3H3V51]
Một tấm bìa diện tích $100\text{ cm}^2$ bị cắt đi một góc có diện tích $15\text{ cm}^2$. Phần bìa còn lại có diện tích bao nhiêu xăng-ti-mét vuông?
\shortans{85}
\loigiai{
\textbf{Đáp số:} $85$\\
$100 - 15 = 85\text{ (cm}^2)$.
}
\end{ex}

\begin{ex} %[3H3V51]
Có 4 hình vuông nhỏ, mỗi hình có diện tích $9\text{ cm}^2$. Ghép lại được hình lớn có diện tích bao nhiêu xăng-ti-mét vuông?
\shortans{36}
\loigiai{
\textbf{Đáp số:} $36$\\
$9 \times 4 = 36\text{ (cm}^2)$.
}
\end{ex}

\begin{ex} %[3H3V51]
Một hình chữ nhật có diện tích $72\text{ cm}^2$ được chia thành 8 phần bằng nhau. Mỗi phần có diện tích bao nhiêu xăng-ti-mét vuông?
\shortans{9}
\loigiai{
\textbf{Đáp số:} $9$\\
$72 : 8 = 9\text{ (cm}^2)$.
}
\end{ex}

\begin{ex} %[3H3V51]
Tính: $15\text{ cm}^2 \times 6 - 25\text{ cm}^2$.
\shortans{65}
\loigiai{
\textbf{Đáp số:} $65$\\
$90 - 25 = 65\text{ (cm}^2)$.
}
\end{ex}

\newpage

\TaoPhieuKiemTra{Chủ đề 9 -- Chu vi, diện tích một số hình phẳng}{52}{Diện tích hình chữ nhật, diện tích hình vuông}{https://loptoanthaylinh.vn/thi-online/toan-3/bai-52}
\textbf{Phần 1. Trắc nghiệm lựa chọn}
\setcounterpageref{ex}{0}

\begin{ex} %[3H3V52]
Một hình chữ nhật có chiều dài $12\text{ cm}$ và chiều rộng $7\text{ cm}$. Diện tích của hình chữ nhật đó là
\choice
{$38\text{ cm}^2$}
{$84\text{ cm}$}
{$48\text{ cm}^2$}
{\True $84\text{ cm}^2$}
\loigiai{
Diện tích hình chữ nhật: $12 \times 7 = 84\text{ (cm}^2)$.
\textbf{Đáp án của câu hỏi là D}
}
\end{ex}

\begin{ex} %[3H3V52]
Một hình vuông có cạnh dài $9\text{ cm}$. Diện tích của hình vuông đó là
\choice
{$36\text{ cm}^2$}
{\True $81\text{ cm}^2$}
{$72\text{ cm}^2$}
{$81\text{ cm}$}
\loigiai{
Diện tích hình vuông: $9 \times 9 = 81\text{ (cm}^2)$.
\textbf{Đáp án của câu hỏi là B}
}
\end{ex}

\begin{ex} %[3H3V52]
Một hình chữ nhật có diện tích $72\text{ cm}^2$, chiều dài là $9\text{ cm}$. Chu vi của hình chữ nhật đó là
\choice
{$17\text{ cm}$}
{$36\text{ cm}$}
{\True $34\text{ cm}$}
{$32\text{ cm}$}
\loigiai{
Chiều rộng: $72 : 9 = 8\text{ cm}$. Chu vi: $(9 + 8) \times 2 = 34\text{ (cm)}$.
\textbf{Đáp án của câu hỏi là C}
}
\end{ex}

\begin{ex} %[3H3V52]
Một hình vuông có chu vi bằng $36\text{ cm}$. Diện tích của hình vuông đó là
\choice
{\True $81\text{ cm}^2$}
{$72\text{ cm}^2$}
{$64\text{ cm}^2$}
{$90\text{ cm}^2$}
\loigiai{
Cạnh hình vuông: $36 : 4 = 9\text{ cm}$. Diện tích: $9 \times 9 = 81\text{ (cm}^2)$.
\textbf{Đáp án của câu hỏi là A}
}
\end{ex}

\begin{ex} %[3H3V52]
Một mảnh đất hình chữ nhật có chiều rộng $8\text{ m}$, chiều dài gấp đôi chiều rộng. Diện tích mảnh đất đó là
\choice
{$48\text{ m}^2$}
{$64\text{ m}^2$}
{$256\text{ m}^2$}
{\True $128\text{ m}^2$}
\loigiai{
Chiều dài: $8 \times 2 = 16\text{ m}$. Diện tích: $16 \times 8 = 128\text{ (m}^2)$.
\textbf{Đáp án của câu hỏi là D}
}
\end{ex}

\begin{ex} %[3H3V52]
Nếu cạnh của một hình vuông tăng lên gấp 2 lần thì diện tích của hình vuông đó sẽ
\choice
{Tăng lên 2 lần}
{\True Tăng lên 4 lần}
{Tăng lên 8 lần}
{Không thay đổi}
\loigiai{
Diện tích hình vuông bằng cạnh $\times$ cạnh. Khi cạnh gấp đôi: $2 \times 2 = 4$ lần.
\textbf{Đáp án của câu hỏi là B}
}
\end{ex}

\begin{ex} %[3H3V52]
Một căn phòng hình chữ nhật có chiều dài $8\text{ m}$, chiều rộng $5\text{ m}$. Người ta lát nền bằng các viên gạch vuông cạnh $1\text{ m}$. Số viên gạch cần dùng là
\choice
{$26$ viên}
{$30$ viên}
{\True $40$ viên}
{$50$ viên}
\loigiai{
Diện tích nền: $8 \times 5 = 40\text{ m}^2$. Mỗi viên gạch $1\text{ m}^2$. Cần $40 : 1 = 40$ viên gạch.
\textbf{Đáp án của câu hỏi là C}
}
\end{ex}

\begin{ex} %[3H3V52]
Một tấm bìa hình vuông có diện tích $64\text{ cm}^2$. Chu vi của tấm bìa đó là
\choice
{\True $32\text{ cm}$}
{$16\text{ cm}$}
{$64\text{ cm}$}
{$24\text{ cm}$}
\loigiai{
Vì $8 \times 8 = 64$ nên cạnh hình vuông là $8\text{ cm}$. Chu vi: $8 \times 4 = 32\text{ (cm)}$.
\textbf{Đáp án của câu hỏi là A}
}
\end{ex}

\begin{ex} %[3H3V52]
Một tấm tôn hình chữ nhật có chiều dài $15\text{ cm}$, chiều rộng $9\text{ cm}$. Người ta cắt đi một góc hình vuông cạnh $4\text{ cm}$. Diện tích phần tôn còn lại là
\choice
{$135\text{ cm}^2$}
{$125\text{ cm}^2$}
{$120\text{ cm}^2$}
{\True $119\text{ cm}^2$}
\loigiai{
Diện tích ban đầu: $15 \times 9 = 135\text{ cm}^2$. Diện tích cắt đi: $4 \times 4 = 16\text{ cm}^2$. Còn lại: $135 - 16 = 119\text{ (cm}^2)$.
\textbf{Đáp án của câu hỏi là D}
}
\end{ex}

\begin{ex} %[3H3V52]
Người ta ghép 3 hình vuông bằng nhau cạnh $4\text{ cm}$ thành một hình chữ nhật. Diện tích hình chữ nhật đó là
\choice
{$36\text{ cm}^2$}
{\True $48\text{ cm}^2$}
{$64\text{ cm}^2$}
{$32\text{ cm}^2$}
\loigiai{
Diện tích mỗi hình vuông: $4 \times 4 = 16\text{ cm}^2$. Diện tích hình chữ nhật: $16 \times 3 = 48\text{ (cm}^2)$.
\textbf{Đáp án của câu hỏi là B}
}
\end{ex}

\begin{ex} %[3H3V52]
Một hình chữ nhật có chu vi $28\text{ cm}$, chiều dài là $8\text{ cm}$. Diện tích của hình chữ nhật đó là
\choice
{$40\text{ cm}^2$}
{$36\text{ cm}^2$}
{\True $48\text{ cm}^2$}
{$56\text{ cm}^2$}
\loigiai{
Nửa chu vi: $28 : 2 = 14\text{ cm}$. Chiều rộng: $14 - 8 = 6\text{ cm}$. Diện tích: $8 \times 6 = 48\text{ (cm}^2)$.
\textbf{Đáp án của câu hỏi là C}
}
\end{ex}

\begin{ex} %[3H3V52]
So sánh diện tích hình chữ nhật kích thước $6\text{ cm} \times 6\text{ cm}$ (hình vuông) và hình chữ nhật kích thước $8\text{ cm} \times 4\text{ cm}$. Khẳng định đúng là
\choice
{\True Hình vuông có diện tích lớn hơn $4\text{ cm}^2$}
{Hình chữ nhật lớn hơn}
{Hai hình có diện tích bằng nhau}
{Không so sánh được}
\loigiai{
Diện tích hình vuông: $6 \times 6 = 36\text{ cm}^2$. Diện tích hình chữ nhật: $8 \times 4 = 32\text{ cm}^2$. Lớn hơn: $36 - 32 = 4\text{ cm}^2$.
\textbf{Đáp án của câu hỏi là A}
}
\end{ex}

\textbf{Phần 2. Trắc nghiệm đúng sai}
\setcounterpageref{ex}{0}

\begin{ex} %[3H3V52]
Xét tính đúng sai của công thức tính diện tích hình phẳng:
\choiceTF
{\True Diện tích hình chữ nhật bằng chiều dài nhân với chiều rộng (cùng đơn vị đo).}
{\True Diện tích hình vuông bằng độ dài một cạnh nhân với chính nó.}
{Một hình vuông cạnh $5\text{ cm}$ có diện tích là $20\text{ cm}^2$.}
{Đơn vị đo chu vi và diện tích là như nhau.}
\loigiai{
\begin{itemize}
    \item Ý a) Đúng: Công thức chuẩn.
    \item Ý b) Đúng: Công thức chuẩn.
    \item Ý c) Sai: Sai vì $5 \times 5 = 25\text{ cm}^2$ ($20\text{ cm}$ là chu vi).
    \item Ý d) Sai: Sai vì chu vi đo bằng $\text{cm}$, diện tích đo bằng $\text{cm}^2$.
\end{itemize}
\begin{itemize}
    \item a) \textbf{Đúng}.
    \item b) \textbf{Đúng}.
    \item c) \textbf{Sai}.
    \item d) \textbf{Sai}.
\end{itemize}
}
\end{ex}

\begin{ex} %[3H3V52]
Một hình chữ nhật có chiều dài $10\text{ cm}$, chiều rộng $5\text{ cm}$:
\choiceTF
{\True Chu vi hình chữ nhật bằng $30\text{ cm}$.}
{\True Diện tích hình chữ nhật bằng $50\text{ cm}^2$.}
{Nếu tăng chiều rộng lên $5\text{ cm}$ để thành hình vuông thì diện tích không đổi.}
{\True Hình chữ nhật có thể chia thành 2 hình vuông cạnh $5\text{ cm}$.}
\loigiai{
\begin{itemize}
    \item Ý a) Đúng: $(10 + 5) \times 2 = 30\text{ cm}$.
    \item Ý b) Đúng: $10 \times 5 = 50\text{ cm}^2$.
    \item Ý c) Sai: Sai vì hình vuông mới có diện tích $10 \times 10 = 100\text{ cm}^2$.
    \item Ý d) Đúng: $10 : 5 = 2$ hình vuông.
\end{itemize}
\begin{itemize}
    \item a) \textbf{Đúng}.
    \item b) \textbf{Đúng}.
    \item c) \textbf{Sai}.
    \item d) \textbf{Đúng}.
\end{itemize}
}
\end{ex}

\textbf{Phần 3. Trả lời ngắn}
\setcounterpageref{ex}{0}

\begin{ex} %[3H3V52]
Tính diện tích hình chữ nhật có chiều dài $14\text{ cm}$, chiều rộng $6\text{ cm}$.
\shortans{84}
\loigiai{
\textbf{Đáp số:} $84$\\
$14 \times 6 = 84\text{ (cm}^2)$.
}
\end{ex}

\begin{ex} %[3H3V52]
Một hình vuông có cạnh $8\text{ cm}$. Tính diện tích hình vuông đó theo xăng-ti-mét vuông.
\shortans{64}
\loigiai{
\textbf{Đáp số:} $64$\\
$8 \times 8 = 64\text{ (cm}^2)$.
}
\end{ex}

\begin{ex} %[3H3V52]
Một hình chữ nhật có diện tích $45\text{ cm}^2$ và chiều rộng $5\text{ cm}$. Tính chu vi hình chữ nhật đó theo xăng-ti-mét.
\shortans{28}
\loigiai{
\textbf{Đáp số:} $28$\\
Chiều dài: $45 : 5 = 9\text{ cm}$. Chu vi: $(9 + 5) \times 2 = 28\text{ (cm)}$.
}
\end{ex}

\begin{ex} %[3H3V52]
Hình vuông có chu vi $28\text{ cm}$. Tính diện tích của hình vuông đó theo xăng-ti-mét vuông.
\shortans{49}
\loigiai{
\textbf{Đáp số:} $49$\\
Cạnh: $28 : 4 = 7\text{ cm}$. Diện tích: $7 \times 7 = 49\text{ (cm}^2)$.
}
\end{ex}

\begin{ex} %[3H3V52]
Một hình chữ nhật có chiều dài $15\text{ cm}$, chiều rộng bằng $\dfrac{1}{3}$ chiều dài. Tính diện tích hình chữ nhật đó.
\shortans{75}
\loigiai{
\textbf{Đáp số:} $75$\\
Chiều rộng: $15 : 3 = 5\text{ cm}$. Diện tích: $15 \times 5 = 75\text{ (cm}^2)$.
}
\end{ex}

\begin{ex} %[3H3V52]
Cần bao nhiêu viên gạch men hình vuông cạnh $2\text{ dm}$ để lát kín một căn phòng diện tích $20\text{ m}^2$ (biết $1\text{ m}^2 = 25$ viên gạch)?
\shortans{500}
\loigiai{
\textbf{Đáp số:} $500$\\
Số viên gạch: $20 \times 25 = 500$ viên.
}
\end{ex}

\newpage

\TaoPhieuKiemTra{Chủ đề 9 -- Chu vi, diện tích một số hình phẳng}{53}{Luyện tập chung chu vi và diện tích}{https://loptoanthaylinh.vn/thi-online/toan-3/bai-53}
\textbf{Phần 1. Trắc nghiệm lựa chọn}
\setcounterpageref{ex}{0}

\begin{ex} %[3H3V53]
Một hình chữ nhật có chiều dài $16\text{ cm}$, chiều rộng bằng một nửa chiều dài. Diện tích của hình chữ nhật đó là
\choice
{$64\text{ cm}^2$}
{$96\text{ cm}^2$}
{\True $128\text{ cm}^2$}
{$48\text{ cm}^2$}
\loigiai{
Chiều rộng: $16 : 2 = 8\text{ cm}$. Diện tích: $16 \times 8 = 128\text{ (cm}^2)$.
\textbf{Đáp án của câu hỏi là C}
}
\end{ex}

\begin{ex} %[3H3V53]
Một hình vuông có diện tích $36\text{ cm}^2$. Chu vi của hình vuông đó là
\choice
{\True $24\text{ cm}$}
{$18\text{ cm}$}
{$36\text{ cm}$}
{$12\text{ cm}$}
\loigiai{
Vì $6 \times 6 = 36$ nên cạnh hình vuông là $6\text{ cm}$. Chu vi: $6 \times 4 = 24\text{ (cm)}$.
\textbf{Đáp án của câu hỏi là A}
}
\end{ex}

\begin{ex} %[3H3V53]
Một mảnh đất hình chữ nhật có chu vi $40\text{ m}$, chiều dài hơn chiều rộng $4\text{ m}$. Diện tích của mảnh đất đó là
\choice
{$80\text{ m}^2$}
{$100\text{ m}^2$}
{$84\text{ m}^2$}
{\True $96\text{ m}^2$}
\loigiai{
Nửa chu vi: $40 : 2 = 20\text{ m}$. Chiều dài: $(20 + 4) : 2 = 12\text{ m}$. Chiều rộng: $12 - 4 = 8\text{ m}$. Diện tích: $12 \times 8 = 96\text{ (m}^2)$.
\textbf{Đáp án của câu hỏi là D}
}
\end{ex}

\begin{ex} %[3H3V53]
Người ta rào một mảnh vườn hình chữ nhật có chiều dài $30\text{ m}$ và chiều rộng $18\text{ m}$. Cọc rào cách nhau $2\text{ m}$. Cần tất cả bao nhiêu cọc rào?
\choice
{$46$ cọc}
{\True $48$ cọc}
{$50$ cọc}
{$24$ cọc}
\loigiai{
Chu vi vườn: $(30 + 18) \times 2 = 96\text{ m}$. Số cọc rào khép kín: $96 : 2 = 48$ cọc.
\textbf{Đáp án của câu hỏi là B}
}
\end{ex}

\begin{ex} %[3H3V53]
Cắt một tấm bìa hình vuông cạnh $10\text{ cm}$ thành 4 hình vuông nhỏ bằng nhau. Diện tích của mỗi hình vuông nhỏ là
\choice
{$20\text{ cm}^2$}
{$50\text{ cm}^2$}
{\True $25\text{ cm}^2$}
{$16\text{ cm}^2$}
\loigiai{
Diện tích hình vuông lớn: $10 \times 10 = 100\text{ cm}^2$. Mỗi hình nhỏ: $100 : 4 = 25\text{ (cm}^2)$.
\textbf{Đáp án của câu hỏi là C}
}
\end{ex}

\begin{ex} %[3H3V53]
Một hình chữ nhật có chiều dài $15\text{ cm}$ và diện tích $90\text{ cm}^2$. Chu vi hình chữ nhật đó là
\choice
{\True $42\text{ cm}$}
{$30\text{ cm}$}
{$45\text{ cm}$}
{$36\text{ cm}$}
\loigiai{
Chiều rộng: $90 : 15 = 6\text{ cm}$. Chu vi: $(15 + 6) \times 2 = 42\text{ (cm)}$.
\textbf{Đáp án của câu hỏi là A}
}
\end{ex}

\begin{ex} %[3H3V53]
Một mảnh vườn hình chữ nhật có chiều dài $25\text{ m}$. Nếu bớt chiều dài đi $5\text{ m}$ thì mảnh vườn trở thành hình vuông. Diện tích của mảnh vườn ban đầu là
\choice
{$400\text{ m}^2$}
{$600\text{ m}^2$}
{$450\text{ m}^2$}
{\True $500\text{ m}^2$}
\loigiai{
Cạnh hình vuông bằng chiều rộng mảnh vườn: $25 - 5 = 20\text{ m}$. Diện tích ban đầu: $25 \times 20 = 500\text{ (m}^2)$.
\textbf{Đáp án của câu hỏi là D}
}
\end{ex}

\begin{ex} %[3H3V53]
Cho hình vuông có chu vi $48\text{ cm}$. Chia hình vuông đó thành 2 hình chữ nhật bằng nhau. Chu vi của mỗi hình chữ nhật là
\choice
{$24\text{ cm}$}
{\True $36\text{ cm}$}
{$30\text{ cm}$}
{$40\text{ cm}$}
\loigiai{
Cạnh hình vuông: $48 : 4 = 12\text{ cm}$. Mỗi hình chữ nhật có kích thước $12\text{ cm} \times 6\text{ cm}$. Chu vi: $(12 + 6) \times 2 = 36\text{ (cm)}$.
\textbf{Đáp án của câu hỏi là B}
}
\end{ex}

\begin{ex} %[3H3V53]
Một tam giác có chu vi $60\text{ cm}$. Cạnh thứ nhất dài $18\text{ cm}$, cạnh thứ hai dài $22\text{ cm}$. Cạnh thứ ba dài là
\choice
{$15\text{ cm}$}
{$25\text{ cm}$}
{\True $20\text{ cm}$}
{$24\text{ cm}$}
\loigiai{
Cạnh thứ ba: $60 - (18 + 22) = 20\text{ (cm)}$.
\textbf{Đáp án của câu hỏi là C}
}
\end{ex}

\begin{ex} %[3H3V53]
Một bức tường hình chữ nhật có chiều dài $6\text{ m}$, chiều cao $3\text{ m}$. Trên tường có một cửa sổ diện tích $2\text{ m}^2$. Diện tích cần sơn tường là
\choice
{\True $16\text{ m}^2$}
{$18\text{ m}^2$}
{$14\text{ m}^2$}
{$20\text{ m}^2$}
\loigiai{
Diện tích cả tường: $6 \times 3 = 18\text{ m}^2$. Diện tích cần sơn: $18 - 2 = 16\text{ (m}^2)$.
\textbf{Đáp án của câu hỏi là A}
}
\end{ex}

\begin{ex} %[3H3V53]
Hình vuông $A$ có cạnh gấp 3 lần cạnh hình vuông $B$. Chu vi hình vuông $A$ gấp mấy lần chu vi hình vuông $B$?
\choice
{$9$ lần}
{$6$ lần}
{$12$ lần}
{\True $3$ lần}
\loigiai{
Chu vi tỉ lệ thuận với cạnh: cạnh gấp 3 lần thì chu vi cũng gấp 3 lần.
\textbf{Đáp án của câu hỏi là D}
}
\end{ex}

\begin{ex} %[3H3V53]
Một hình chữ nhật có chu vi $32\text{ cm}$. Chiều dài gấp 3 lần chiều rộng. Diện tích của hình chữ nhật đó là
\choice
{$64\text{ cm}^2$}
{\True $48\text{ cm}^2$}
{$36\text{ cm}^2$}
{$56\text{ cm}^2$}
\loigiai{
Nửa chu vi: $32 : 2 = 16\text{ cm}$. Rộng: $16 : (3 + 1) = 4\text{ cm}$. Dài: $4 \times 3 = 12\text{ cm}$. Diện tích: $12 \times 4 = 48\text{ (cm}^2)$.
\textbf{Đáp án của câu hỏi là B}
}
\end{ex}

\textbf{Phần 2. Trắc nghiệm đúng sai}
\setcounterpageref{ex}{0}

\begin{ex} %[3H3V53]
Xét tính đúng sai của các nhận định ôn tập chu vi và diện tích:
\choiceTF
{\True Đơn vị đo chu vi là đơn vị độ dài ($\text{cm}, \text{m}$), đơn vị diện tích là đơn vị diện tích ($\text{cm}^2, \text{m}^2$).}
{Hình vuông có diện tích $16\text{ cm}^2$ thì chu vi của nó là $8\text{ cm}$.}
{\True Trong các hình chữ nhật có cùng chu vi thì hình vuông có diện tích lớn nhất.}
{Diện tích hình chữ nhật bằng (dài + rộng) nhân 2.}
\loigiai{
\begin{itemize}
    \item Ý a) Đúng: Quy chuẩn đo lường.
    \item Ý b) Sai: Sai vì cạnh là $4\text{ cm}$, chu vi là $4 \times 4 = 16\text{ cm}$.
    \item Ý c) Đúng: Định lí cực trị hình học.
    \item Ý d) Sai: Sai vì đó là công thức tính chu vi.
\end{itemize}
\begin{itemize}
    \item a) \textbf{Đúng}.
    \item b) \textbf{Sai}.
    \item c) \textbf{Đúng}.
    \item d) \textbf{Sai}.
\end{itemize}
}
\end{ex}

\begin{ex} %[3H3V53]
Cho hình vuông $ABCD$ có cạnh dài $8\text{ cm}$:
\choiceTF
{\True Chu vi hình vuông $ABCD$ bằng $32\text{ cm}$.}
{\True Diện tích hình vuông $ABCD$ bằng $64\text{ cm}^2$.}
{Nếu kéo dài một cạnh thêm $2\text{ cm}$ thì được hình chữ nhật có chu vi $34\text{ cm}$.}
{\True Có thể chia hình vuông đó thành 4 hình vuông nhỏ bằng nhau có cạnh $4\text{ cm}$.}
\loigiai{
\begin{itemize}
    \item Ý a) Đúng: $8 \times 4 = 32\text{ cm}$.
    \item Ý b) Đúng: $8 \times 8 = 64\text{ cm}^2$.
    \item Ý c) Sai: Sai vì chu vi tăng thêm $2 \times 2 = 4\text{ cm}$ thành $36\text{ cm}$.
    \item Ý d) Đúng: $8 : 4 = 2$ phần mỗi chiều.
\end{itemize}
\begin{itemize}
    \item a) \textbf{Đúng}.
    \item b) \textbf{Đúng}.
    \item c) \textbf{Sai}.
    \item d) \textbf{Đúng}.
\end{itemize}
}
\end{ex}

\textbf{Phần 3. Trả lời ngắn}
\setcounterpageref{ex}{0}

\begin{ex} %[3H3V53]
Một hình vuông có chu vi $44\text{ cm}$. Tính cạnh hình vuông theo xăng-ti-mét.
\shortans{11}
\loigiai{
\textbf{Đáp số:} $11$\\
$44 : 4 = 11\text{ (cm)}$.
}
\end{ex}

\begin{ex} %[3H3V53]
Một hình chữ nhật có chiều dài $18\text{ cm}$, chiều rộng $5\text{ cm}$. Tính diện tích hình chữ nhật theo xăng-ti-mét vuông.
\shortans{90}
\loigiai{
\textbf{Đáp số:} $90$\\
$18 \times 5 = 90\text{ (cm}^2)$.
}
\end{ex}

\begin{ex} %[3H3V53]
Hình vuông có diện tích $49\text{ cm}^2$. Tính chu vi hình vuông đó theo xăng-ti-mét.
\shortans{28}
\loigiai{
\textbf{Đáp số:} $28$\\
Cạnh là $7\text{ cm}$. Chu vi: $7 \times 4 = 28\text{ (cm)}$.
}
\end{ex}

\begin{ex} %[3H3V53]
Một hình chữ nhật có chu vi $36\text{ cm}$, chiều dài hơn chiều rộng $4\text{ cm}$. Tính diện tích hình chữ nhật theo xăng-ti-mét vuông.
\shortans{77}
\loigiai{
\textbf{Đáp số:} $77$\\
Nửa chu vi: $18\text{ cm}$. Dài: $11\text{ cm}$, rộng: $7\text{ cm}$. Diện tích: $11 \times 7 = 77\text{ (cm}^2)$.
}
\end{ex}

\begin{ex} %[3H3V53]
Cần bao nhiêu mét dây thép gai để rào kín một mảnh vườn hình vuông cạnh $22\text{ m}$ nếu chừa cổng ra vào rộng $2\text{ m}$?
\shortans{86}
\loigiai{
\textbf{Đáp số:} $86$\\
Chu vi: $22 \times 4 = 88\text{ m}$. Cần: $88 - 2 = 86\text{ (m)}$.
}
\end{ex}

\begin{ex} %[3H3V53]
Một hình chữ nhật có diện tích $60\text{ cm}^2$, chiều rộng $6\text{ cm}$. Tính chu vi hình chữ nhật đó theo xăng-ti-mét.
\shortans{32}
\loigiai{
\textbf{Đáp số:} $32$\\
Chiều dài: $60 : 6 = 10\text{ cm}$. Chu vi: $(10 + 6) \times 2 = 32\text{ (cm)}$.
}
\end{ex}

\newpage

% =========================================================================
% CHỦ ĐỀ 10 - CỘNG, TRỪ, NHÂN, CHIA TRONG PHẠM VI 10 000
% =========================================================================
\section{Cộng, trừ, nhân, chia trong phạm vi 10 000}

\TaoPhieuKiemTra{Chủ đề 10 -- Cộng, trừ, nhân, chia trong phạm vi 10 000}{54}{Phép cộng trong phạm vi 10 000}{https://loptoanthaylinh.vn/thi-online/toan-3/bai-54}
\textbf{Phần 1. Trắc nghiệm lựa chọn}
\setcounterpageref{ex}{0}

\begin{ex} %[3D2V54]
Tính giá trị của phép tính: $4528 + 3617$.
\choice
{$8135$}
{\True $8145$}
{$8155$}
{$7145$}
\loigiai{
Đặt tính và cộng lần lượt từ phải sang trái: $4528 + 3617 = 8145$.
\textbf{Đáp án của câu hỏi là B}
}
\end{ex}

\begin{ex} %[3D2V54]
Một trường tiểu học có $1250$ học sinh nam và $1340$ học sinh nữ. Trường đó có tất cả bao nhiêu học sinh?
\choice
{$2580$ học sinh}
{$2600$ học sinh}
{$2490$ học sinh}
{\True $2590$ học sinh}
\loigiai{
Tổng số học sinh: $1250 + 1340 = 2590$ học sinh.
\textbf{Đáp án của câu hỏi là D}
}
\end{ex}

\begin{ex} %[3D2V54]
Tìm số tự nhiên $x$, biết: $x - 3450 = 2860$.
\choice
{\True $6310$}
{$6210$}
{$6410$}
{$5310$}
\loigiai{
$x = 2860 + 3450 = 6310$.
\textbf{Đáp án của câu hỏi là A}
}
\end{ex}

\begin{ex} %[3D2V54]
Tính nhanh giá trị biểu thức: $1250 + 3400 + 750 + 600$.
\choice
{$5800$}
{$6200$}
{\True $6000$}
{$5900$}
\loigiai{
Nhóm thuận tiện: $(1250 + 750) + (3400 + 600) = 2000 + 4000 = 6000$.
\textbf{Đáp án của câu hỏi là C}
}
\end{ex}

\begin{ex} %[3D2V54]
Một đội làm đường ngày thứ nhất làm được $1450\text{ m}$, ngày thứ hai làm được nhiều hơn ngày thứ nhất $350\text{ m}$. Cả hai ngày đội đó làm được bao nhiêu mét đường?
\choice
{$1800\text{ m}$}
{\True $3250\text{ m}$}
{$3150\text{ m}$}
{$3300\text{ m}$}
\loigiai{
Ngày hai làm: $1450 + 350 = 1800\text{ m}$. Cả hai ngày: $1450 + 1800 = 3250\text{ (m)}$.
\textbf{Đáp án của câu hỏi là B}
}
\end{ex}

\begin{ex} %[3D2V54]
Một kho thóc có $3250\text{ kg}$ thóc. Người ta nhập thêm hai đợt thóc: đợt một nhập $1800\text{ kg}$, đợt hai nhập $2450\text{ kg}$. Trong kho có tất cả bao nhiêu ki-lô-gam thóc?
\choice
{$6500\text{ kg}$}
{$7400\text{ kg}$}
{$7600\text{ kg}$}
{\True $7500\text{ kg}$}
\loigiai{
Tổng số thóc: $3250 + 1800 + 2450 = 7500\text{ (kg)}$.
\textbf{Đáp án của câu hỏi là D}
}
\end{ex}

\begin{ex} %[3D2V54]
Nếu một số hạng tăng thêm 250 đơn vị và số hạng kia tăng thêm 450 đơn vị thì tổng sẽ
\choice
{\True Tăng thêm 700 đơn vị}
{Tăng thêm 200 đơn vị}
{Tăng thêm 600 đơn vị}
{Không đổi}
\loigiai{
Tổng mới tăng thêm: $250 + 450 = 700$ đơn vị.
\textbf{Đáp án của câu hỏi là A}
}
\end{ex}

\begin{ex} %[3D2V54]
Tìm tổng của số lớn nhất có 4 chữ số khác nhau và số bé nhất có 4 chữ số khác nhau.
\choice
{$10\,900$}
{$10\,889$}
{\True $10\,899$}
{$11\,000$}
\loigiai{
Số lớn nhất: 9876. Số bé nhất: 1023. Tổng: $9876 + 1023 = 10\,899$.
\textbf{Đáp án của câu hỏi là C}
}
\end{ex}

\begin{ex} %[3D2V54]
Một trang trại có 2560 con gà và số vịt nhiều hơn số gà là 480 con. Trang trại có tất cả bao nhiêu con gà và vịt?
\choice
{$3040$ con}
{\True $5600$ con}
{$5500$ con}
{$5620$ con}
\loigiai{
Số vịt: $2560 + 480 = 3040$ con. Tổng cộng: $2560 + 3040 = 5600$ con.
\textbf{Đáp án của câu hỏi là B}
}
\end{ex}

\begin{ex} %[3D2V54]
Điền chữ số thích hợp vào dấu $*$ để: $4*56 + 2341 = 6897$. Chữ số $*$ là
\choice
{$4$}
{$6$}
{$3$}
{\True $5$}
\loigiai{
Ta có: $4*56 = 6897 - 2341 = 4556 \Rightarrow * = 5$.
\textbf{Đáp án của câu hỏi là D}
}
\end{ex}

\begin{ex} %[3D2V54]
Bố mua một chiếc quạt điện giá $450\,000$ đồng và một ấm đun nước giá $280\,000$ đồng. Hỏi bố phải trả tất cả bao nhiêu tiền?
\choice
{\True $730\,000$ đồng}
{$720\,000$ đồng}
{$740\,000$ đồng}
{$750\,000$ đồng}
\loigiai{
Tổng số tiền: $450\,000 + 280\,000 = 730\,000$ đồng.
\textbf{Đáp án của câu hỏi là A}
}
\end{ex}

\begin{ex} %[3D2V54]
Một thư viện trường học tuần thứ nhất cho mượn 1420 cuốn sách, tuần thứ hai cho mượn nhiều hơn tuần thứ nhất 160 cuốn sách. Cả hai tuần thư viện cho mượn bao nhiêu cuốn sách?
\choice
{$2840$ cuốn sách}
{$3100$ cuốn sách}
{\True $3000$ cuốn sách}
{$2900$ cuốn sách}
\loigiai{
Tuần hai cho mượn: $1420 + 160 = 1580$ cuốn. Cả hai tuần: $1420 + 1580 = 3000$ cuốn sách.
\textbf{Đáp án của câu hỏi là C}
}
\end{ex}

\textbf{Phần 2. Trắc nghiệm đúng sai}
\setcounterpageref{ex}{0}

\begin{ex} %[3D2V54]
Xét tính đúng sai của các khẳng định về phép cộng trong phạm vi 10 000:
\choiceTF
{\True Phép cộng các số có 4 chữ số thực hiện lần lượt từ phải sang trái (từ hàng đơn vị).}
{Tổng của hai số có 4 chữ số luôn luôn là một số có 4 chữ số.}
{\True Khi đổi chỗ các số hạng trong một tổng thì tổng không thay đổi.}
{$3500 + 4500 = 7000$.}
\loigiai{
\begin{itemize}
    \item Ý a) Đúng: Quy tắc đặt tính và tính.
    \item Ý b) Sai: Sai vì ví dụ $6000 + 5000 = 11\,000$ có 5 chữ số.
    \item Ý c) Đúng: Tính chất giao hoán của phép cộng.
    \item Ý d) Sai: Sai vì $3500 + 4500 = 8000$.
\end{itemize}
\begin{itemize}
    \item a) \textbf{Đúng}.
    \item b) \textbf{Sai}.
    \item c) \textbf{Đúng}.
    \item d) \textbf{Sai}.
\end{itemize}
}
\end{ex}

\begin{ex} %[3D2V54]
Một xã có hai thôn: thôn Đông có 3420 nhân khẩu, thôn Đoài có 2850 nhân khẩu:
\choiceTF
{\True Cả xã có tất cả 6270 nhân khẩu.}
{\True Thôn Đông nhiều hơn thôn Đoài 570 nhân khẩu.}
{Làm tròn số dân thôn Đông đến hàng trăm ta được 3400 người.}
{\True Tổng số dân làm tròn đến hàng trăm là khoảng 6300 người.}
\loigiai{
\begin{itemize}
    \item Ý a) Đúng: $3420 + 2850 = 6270$ người.
    \item Ý b) Đúng: $3420 - 2850 = 570$ người.
    \item Ý c) Sai: Sai vì hàng chục là 2 nên đúng là 3400 người (mệnh đề đúng: làm tròn thành 3400). Sửa câu sai: Làm tròn thôn Đông đến hàng trăm được 3500 người.
    \item Ý d) Đúng: $6270$ làm tròn đến hàng trăm là 6300.
\end{itemize}
\begin{itemize}
    \item a) \textbf{Đúng}.
    \item b) \textbf{Đúng}.
    \item c) \textbf{Sai}.
    \item d) \textbf{Đúng}.
\end{itemize}
}
\end{ex}

\textbf{Phần 3. Trả lời ngắn}
\setcounterpageref{ex}{0}

\begin{ex} %[3D2V54]
Tính kết quả phép cộng: $4835 + 2748$.
\shortans{7583}
\loigiai{
\textbf{Đáp số:} $7583$\\
$4835 + 2748 = 7583$.
}
\end{ex}

\begin{ex} %[3D2V54]
Tìm $x$, biết: $x - 2480 = 3520$.
\shortans{6000}
\loigiai{
\textbf{Đáp số:} $6000$\\
$x = 3520 + 2480 = 6000$.
}
\end{ex}

\begin{ex} %[3D2V54]
Tính nhanh giá trị biểu thức: $2400 + 1750 + 600 + 250$.
\shortans{5000}
\loigiai{
\textbf{Đáp số:} $5000$\\
$(2400 + 600) + (1750 + 250) = 3000 + 2000 = 5000$.
}
\end{ex}

\begin{ex} %[3D2V54]
Một cửa hàng buổi sáng bán được $1350\text{ kg}$ gạo, buổi chiều bán được nhiều hơn buổi sáng $420\text{ kg}$ gạo. Hỏi cả hai buổi cửa hàng bán được bao nhiêu ki-lô-gam gạo?
\shortans{3120}
\loigiai{
\textbf{Đáp số:} $3120$\\
Buổi chiều: $1770\text{ kg}$. Cả hai buổi: $1350 + 1770 = 3120\text{ (kg)}$.
}
\end{ex}

\begin{ex} %[3D2V54]
Tính tổng của số chẵn lớn nhất có 3 chữ số và số tròn nghìn bé nhất có 4 chữ số.
\shortans{1998}
\loigiai{
\textbf{Đáp số:} $1998$\\
$1000 + 998 = 1998$.
}
\end{ex}

\begin{ex} %[3D2V54]
Một vườn ươm có 3500 cây keo và 2800 cây bạch đàn. Hỏi vườn ươm có tất cả bao nhiêu cây giống?
\shortans{6300}
\loigiai{
\textbf{Đáp số:} $6300$\\
$3500 + 2800 = 6300$ cây.
}
\end{ex}

\newpage

\TaoPhieuKiemTra{Chủ đề 10 -- Cộng, trừ, nhân, chia trong phạm vi 10 000}{55}{Phép trừ trong phạm vi 10 000}{https://loptoanthaylinh.vn/thi-online/toan-3/bai-55}
\textbf{Phần 1. Trắc nghiệm lựa chọn}
\setcounterpageref{ex}{0}

\begin{ex} %[3D2V55]
Tính giá trị của phép tính: $7450 - 3280$.
\choice
{\True $4170$}
{$4270$}
{$4130$}
{$4070$}
\loigiai{
$7450 - 3280 = 4170$.
\textbf{Đáp án của câu hỏi là A}
}
\end{ex}

\begin{ex} %[3D2V55]
Một kho hàng có $8500\text{ kg}$ thóc. Người ta đã xuất đi $3650\text{ kg}$. Trong kho còn lại bao nhiêu ki-lô-gam thóc?
\choice
{$4950\text{ kg}$}
{$4800\text{ kg}$}
{\True $4850\text{ kg}$}
{$5150\text{ kg}$}
\loigiai{
$8500 - 3650 = 4850\text{ (kg)}$.
\textbf{Đáp án của câu hỏi là C}
}
\end{ex}

\begin{ex} %[3D2V55]
Tìm số tự nhiên $y$, biết: $9200 - y = 4580$.
\choice
{$4720$}
{\True $4620$}
{$4520$}
{$5620$}
\loigiai{
$y = 9200 - 4580 = 4620$.
\textbf{Đáp án của câu hỏi là B}
}
\end{ex}

\begin{ex} %[3D2V55]
Một cuộn dây cáp dài $5000\text{ m}$. Người ta cắt đi hai lần: lần đầu cắt $1250\text{ m}$, lần sau cắt $1850\text{ m}$. Cuộn dây còn lại dài bao nhiêu mét?
\choice
{$2100\text{ m}$}
{$1800\text{ m}$}
{$2000\text{ m}$}
{\True $1900\text{ m}$}
\loigiai{
Tổng hai lần cắt: $1250 + 1850 = 3100\text{ m}$. Còn lại: $5000 - 3100 = 1900\text{ (m)}$.
\textbf{Đáp án của câu hỏi là D}
}
\end{ex}

\begin{ex} %[3D2V55]
Tìm hiệu của số lớn nhất có 4 chữ số và số tròn trăm lớn nhất có 4 chữ số bé hơn 5000.
\choice
{\True $5099$}
{$4999$}
{$5199$}
{$4899$}
\loigiai{
Số lớn nhất có 4 chữ số là 9999. Số tròn trăm lớn nhất bé hơn 5000 là 4900. Hiệu: $9999 - 4900 = 5099$.
\textbf{Đáp án của câu hỏi là A}
}
\end{ex}

\begin{ex} %[3D2V55]
Một cửa hàng có $6400\ell$ xăng. Sau một tuần bán, cửa hàng còn lại $2750\ell$ xăng. Hỏi cửa hàng đã bán bao nhiêu lít xăng?
\choice
{$3750\ell$}
{$3550\ell$}
{\True $3650\ell$}
{$4650\ell$}
\loigiai{
Số xăng đã bán: $6400 - 2750 = 3650\text{ (\ell)}$.
\textbf{Đáp án của câu hỏi là C}
}
\end{ex}

\begin{ex} %[3D2V55]
Nếu số bị trừ giảm đi 300 đơn vị và số trừ giữ nguyên thì hiệu sẽ
\choice
{Tăng thêm 300 đơn vị}
{\True Giảm đi 300 đơn vị}
{Không thay đổi}
{Giảm đi 600 đơn vị}
\loigiai{
Số bị trừ giảm bao nhiêu đơn vị thì hiệu giảm bấy nhiêu đơn vị.
\textbf{Đáp án của câu hỏi là B}
}
\end{ex}

\begin{ex} %[3D2V55]
Tính nhanh giá trị: $7850 - 2400 - 1850$.
\choice
{$3500$}
{$3700$}
{$3800$}
{\True $3600$}
\loigiai{
Nhóm: $(7850 - 1850) - 2400 = 6000 - 2400 = 3600$.
\textbf{Đáp án của câu hỏi là D}
}
\end{ex}

\begin{ex} %[3D2V55]
Một trang trại thu hoạch được $4500\text{ kg}$ bơ. Đợt một bán $1600\text{ kg}$, đợt hai bán nhiều hơn đợt một $300\text{ kg}$. Trang trại còn lại bao nhiêu ki-lô-gam bơ?
\choice
{\True $1000\text{ kg}$}
{$1200\text{ kg}$}
{$1100\text{ kg}$}
{$900\text{ kg}$}
\loigiai{
Đợt hai bán: $1600 + 300 = 1900\text{ kg}$. Tổng hai đợt bán: $1600 + 1900 = 3500\text{ kg}$. Còn lại: $4500 - 3500 = 1000\text{ (kg)}$.
\textbf{Đáp án của câu hỏi là A}
}
\end{ex}

\begin{ex} %[3D2V55]
Trong phép trừ có hiệu bằng 2450. Nếu tăng số bị trừ thêm 150 đơn vị và tăng số trừ thêm 150 đơn vị thì hiệu mới bằng
\choice
{$2600$}
{$2300$}
{\True $2450$}
{$2750$}
\loigiai{
Khi cả số bị trừ và số trừ cùng tăng một số đơn vị như nhau thì hiệu không thay đổi.
\textbf{Đáp án của câu hỏi là C}
}
\end{ex}

\begin{ex} %[3D2V55]
Quãng đường từ Hà Nội đến Nghệ An dài khoảng $300\text{ km}$. Một bác tài xế đã đi được $185\text{ km}$. Bác tài xế còn phải đi bao nhiêu ki-lô-mét nữa?
\choice
{$125\text{ km}$}
{\True $115\text{ km}$}
{$105\text{ km}$}
{$135\text{ km}$}
\loigiai{
Quãng đường còn lại: $300 - 185 = 115\text{ (km)}$.
\textbf{Đáp án của câu hỏi là B}
}
\end{ex}

\begin{ex} %[3D2V55]
Tìm một số biết rằng lấy số lớn nhất có 4 chữ số trừ đi số đó thì được 3456.
\choice
{$6534$}
{$5543$}
{$6544$}
{\True $6543$}
\loigiai{
Số cần tìm: $9999 - 3456 = 6543$.
\textbf{Đáp án của câu hỏi là D}
}
\end{ex}

\textbf{Phần 2. Trắc nghiệm đúng sai}
\setcounterpageref{ex}{0}

\begin{ex} %[3D2V55]
Xét tính đúng sai của các quy tắc về phép trừ:
\choiceTF
{\True Muốn tìm số trừ ta lấy số bị trừ trừ đi hiệu.}
{Trong phép trừ các số tự nhiên, số trừ luôn lớn hơn số bị trừ.}
{\True Khi cùng thêm vào cả số bị trừ và số trừ một số đơn vị như nhau thì hiệu không thay đổi.}
{$8000 - 2500 = 6500$.}
\loigiai{
\begin{itemize}
    \item Ý a) Đúng: Quy tắc tìm thành phần chưa biết.
    \item Ý b) Sai: Sai vì số bị trừ luôn lớn hơn hoặc bằng số trừ.
    \item Ý c) Đúng: Tính chất của hiệu.
    \item Ý d) Sai: Sai vì $8000 - 2500 = 5500$.
\end{itemize}
\begin{itemize}
    \item a) \textbf{Đúng}.
    \item b) \textbf{Sai}.
    \item c) \textbf{Đúng}.
    \item d) \textbf{Sai}.
\end{itemize}
}
\end{ex}

\begin{ex} %[3D2V55]
Một hồ nước có dung tích $9000\text{ m}^3$ nước, người ta xả bớt nước ra để chống lũ:
\choiceTF
{\True Đợt 1 xả $3200\text{ m}^3$ thì hồ còn lại $5800\text{ m}^3$ nước.}
{\True Đợt 2 xả tiếp $2800\text{ m}^3$ thì sau hai đợt đã xả tổng cộng $6000\text{ m}^3$ nước.}
{Sau hai đợt xả nước trong hồ cạn sạch.}
{\True Lượng nước còn lại trong hồ đúng bằng một phần ba lượng nước ban đầu.}
\loigiai{
\begin{itemize}
    \item Ý a) Đúng: $9000 - 3200 = 5800\text{ m}^3$.
    \item Ý b) Đúng: $3200 + 2800 = 6000\text{ m}^3$.
    \item Ý c) Sai: Sai vì trong hồ còn lại $9000 - 6000 = 3000\text{ m}^3$.
    \item Ý d) Đúng: $3000 : 9000 = \dfrac{1}{3}$.
\end{itemize}
\begin{itemize}
    \item a) \textbf{Đúng}.
    \item b) \textbf{Đúng}.
    \item c) \textbf{Sai}.
    \item d) \textbf{Đúng}.
\end{itemize}
}
\end{ex}

\textbf{Phần 3. Trả lời ngắn}
\setcounterpageref{ex}{0}

\begin{ex} %[3D2V55]
Tính kết quả phép trừ: $8340 - 4565$.
\shortans{3775}
\loigiai{
\textbf{Đáp số:} $3775$\\
$8340 - 4565 = 3775$.
}
\end{ex}

\begin{ex} %[3D2V55]
Tìm số tự nhiên $x$, biết: $7500 - x = 2840$.
\shortans{4660}
\loigiai{
\textbf{Đáp số:} $4660$\\
$x = 7500 - 2840 = 4660$.
}
\end{ex}

\begin{ex} %[3D2V55]
Một trang trại có 5400 con vịt, sau khi bán đi một số con thì còn lại 1850 con. Hỏi trang trại đã bán bao nhiêu con vịt?
\shortans{3550}
\loigiai{
\textbf{Đáp số:} $3550$\\
$5400 - 1850 = 3550$ con.
}
\end{ex}

\begin{ex} %[3D2V55]
Tính nhanh: $6540 - 2300 - 1540$.
\shortans{2700}
\loigiai{
\textbf{Đáp số:} $2700$\\
$(6540 - 1540) - 2300 = 5000 - 2300 = 2700$.
}
\end{ex}

\begin{ex} %[3D2V55]
Một cuộn dây dài $4000\text{ m}$. Lần thứ nhất người ta dùng $1200\text{ m}$, lần thứ hai dùng $1500\text{ m}$. Hỏi cuộn dây còn lại bao nhiêu mét?
\shortans{1300}
\loigiai{
\textbf{Đáp số:} $1300$\\
$4000 - (1200 + 1500) = 1300\text{ (m)}$.
}
\end{ex}

\begin{ex} %[3D2V55]
Hiệu hai số là 1500. Nếu số bị trừ tăng 200 đơn vị và số trừ giảm 100 đơn vị thì hiệu mới bằng bao nhiêu?
\shortans{1800}
\loigiai{
\textbf{Đáp số:} $1800$\\
Hiệu mới: $1500 + 200 + 100 = 1800$.
}
\end{ex}

\newpage

\TaoPhieuKiemTra{Chủ đề 10 -- Cộng, trừ, nhân, chia trong phạm vi 10 000}{56}{Nhân số có bốn chữ số với số có một chữ số}{https://loptoanthaylinh.vn/thi-online/toan-3/bai-56}
\textbf{Phần 1. Trắc nghiệm lựa chọn}
\setcounterpageref{ex}{0}

\begin{ex} %[3D2V56]
Tính giá trị phép tính: $1425 \times 4$.
\choice
{$5600$}
{$5800$}
{$5680$}
{\True $5700$}
\loigiai{
$1425 \times 4 = 5700$.
\textbf{Đáp án của câu hỏi là D}
}
\end{ex}

\begin{ex} %[3D2V56]
Một xã có 3 thôn, mỗi thôn có 1240 nhân khẩu. Toàn xã đó có bao nhiêu nhân khẩu?
\choice
{$3620$ nhân khẩu}
{\True $3720$ nhân khẩu}
{$3820$ nhân khẩu}
{$3700$ nhân khẩu}
\loigiai{
Tổng số dân: $1240 \times 3 = 3720$ người.
\textbf{Đáp án của câu hỏi là B}
}
\end{ex}

\begin{ex} %[3D2V56]
Tìm $x$, biết: $x : 3 = 1650$.
\choice
{$4850$}
{$5050$}
{\True $4950$}
{$4900$}
\loigiai{
$x = 1650 \times 3 = 4950$.
\textbf{Đáp án của câu hỏi là C}
}
\end{ex}

\begin{ex} %[3D2V56]
Một hình vuông có cạnh dài $1520\text{ m}$. Chu vi của hình vuông đó là
\choice
{\True $6080\text{ m}$}
{$6000\text{ m}$}
{$6120\text{ m}$}
{$5980\text{ m}$}
\loigiai{
Chu vi hình vuông: $1520 \times 4 = 6080\text{ (m)}$.
\textbf{Đáp án của câu hỏi là A}
}
\end{ex}

\begin{ex} %[3D2V56]
Một nhà máy mỗi ngày sản xuất được 1125 đôi giày. Hỏi trong 6 ngày làm việc nhà máy đó sản xuất được bao nhiêu đôi giày?
\choice
{$6650$ đôi giày}
{$6850$ đôi giày}
{$6700$ đôi giày}
{\True $6750$ đôi giày}
\loigiai{
Số đôi giày sản xuất được: $1125 \times 6 = 6750$ đôi.
\textbf{Đáp án của câu hỏi là D}
}
\end{ex}

\begin{ex} %[3D2V56]
Tính nhanh giá trị của biểu thức: $1250 \times 2 + 1250 \times 6$.
\choice
{$9000$}
{\True $10\,000$}
{$8000$}
{$12\,500$}
\loigiai{
$1250 \times (2 + 6) = 1250 \times 8 = 10\,000$.
\textbf{Đáp án của câu hỏi là B}
}
\end{ex}

\begin{ex} %[3D2V56]
Một xe tải chở 4 chuyến hàng, mỗi chuyến chở $2150\text{ kg}$ ngô. Người ta đã dỡ xuống $3600\text{ kg}$ ngô. Số ngô còn lại trên các xe là
\choice
{$4800\text{ kg}$}
{$5200\text{ kg}$}
{\True $5000\text{ kg}$}
{$4900\text{ kg}$}
\loigiai{
Tổng khối lượng ngô: $2150 \times 4 = 8600\text{ kg}$. Còn lại: $8600 - 3600 = 5000\text{ (kg)}$.
\textbf{Đáp án của câu hỏi là C}
}
\end{ex}

\begin{ex} %[3D2V56]
Tính giá trị của biểu thức: $A = 2314 \times 3 - 2800$.
\choice
{\True $4142$}
{$4132$}
{$4242$}
{$4042$}
\loigiai{
$2314 \times 3 = 6942$; $6942 - 2800 = 4142$.
\textbf{Đáp án của câu hỏi là A}
}
\end{ex}

\begin{ex} %[3D2V56]
Nếu một thừa số tăng thêm 2 đơn vị và thừa số kia giữ nguyên bằng 1450 thì tích tăng thêm bao nhiêu đơn vị?
\choice
{$1450$}
{$2800$}
{$3000$}
{\True $2900$}
\loigiai{
Tích tăng thêm: $1450 \times 2 = 2900$ đơn vị.
\textbf{Đáp án của câu hỏi là D}
}
\end{ex}

\begin{ex} %[3D2V56]
Người ta đóng 1200 gói kẹo vào các thùng, mỗi thùng 8 gói. Hỏi 7 thùng kẹo như thế có bao nhiêu gói kẹo? (Nếu mỗi thùng có 1200 gói kẹo thì 7 thùng có bao nhiêu gói kẹo?) Có 7 thùng kẹo, mỗi thùng chứa 1250 gói kẹo. Tổng số gói kẹo là
\choice
{$8650$ gói kẹo}
{\True $8750$ gói kẹo}
{$8850$ gói kẹo}
{$8500$ gói kẹo}
\loigiai{
Tổng số gói kẹo: $1250 \times 7 = 8750$ gói kẹo.
\textbf{Đáp án của câu hỏi là B}
}
\end{ex}

\begin{ex} %[3D2V56]
Tìm tích của số bé nhất có bốn chữ số khác nhau với 5.
\choice
{$5000$}
{$5120$}
{\True $5115$}
{$5125$}
\loigiai{
Số bé nhất có 4 chữ số khác nhau là 1023. Tích: $1023 \times 5 = 5115$.
\textbf{Đáp án của câu hỏi là C}
}
\end{ex}

\begin{ex} %[3D2V56]
Một cuộn dây dài $2450\text{ m}$. Bốn cuộn dây như thế dài tất cả bao nhiêu mét?
\choice
{\True $9800\text{ m}$}
{$9600\text{ m}$}
{$9900\text{ m}$}
{$9700\text{ m}$}
\loigiai{
$2450 \times 4 = 9800\text{ (m)}$.
\textbf{Đáp án của câu hỏi là A}
}
\end{ex}

\textbf{Phần 2. Trắc nghiệm đúng sai}
\setcounterpageref{ex}{0}

\begin{ex} %[3D2V56]
Xét tính đúng sai của phép nhân số có bốn chữ số với số có một chữ số:
\choiceTF
{\True Thực hiện phép nhân từ phải sang trái (bắt đầu từ hàng đơn vị).}
{$2500 \times 4 = 1000$.}
{\True Tích của số có 4 chữ số với 1 luôn bằng chính số đó.}
{Tích của $1500 \times 6$ bằng $8000$.}
\loigiai{
\begin{itemize}
    \item Ý a) Đúng: Quy tắc chuẩn.
    \item Ý b) Sai: Sai vì $2500 \times 4 = 10\,000$.
    \item Ý c) Đúng: Tính chất của 1.
    \item Ý d) Sai: Sai vì $1500 \times 6 = 9000$.
\end{itemize}
\begin{itemize}
    \item a) \textbf{Đúng}.
    \item b) \textbf{Sai}.
    \item c) \textbf{Đúng}.
    \item d) \textbf{Sai}.
\end{itemize}
}
\end{ex}

\begin{ex} %[3D2V56]
Một đội xe gồm 4 xe chở hàng, mỗi xe chở $1850\text{ kg}$ phân bón:
\choiceTF
{\True Đội xe chở tất cả $7400\text{ kg}$ phân bón.}
{\True Hai xe chở được $3700\text{ kg}$ phân bón.}
{Ba xe chở nhiều hơn $6000\text{ kg}$ phân bón.}
{\True Nếu mỗi bao nặng $50\text{ kg}$ thì mỗi xe chở 37 bao phân bón.}
\loigiai{
\begin{itemize}
    \item Ý a) Đúng: $1850 \times 4 = 7400\text{ kg}$.
    \item Ý b) Đúng: $1850 \times 2 = 3700\text{ kg}$.
    \item Ý c) Sai: Sai vì $1850 \times 3 = 5550\text{ kg} < 6000\text{ kg}$.
    \item Ý d) Đúng: $1850 : 50 = 37$ bao.
\end{itemize}
\begin{itemize}
    \item a) \textbf{Đúng}.
    \item b) \textbf{Đúng}.
    \item c) \textbf{Sai}.
    \item d) \textbf{Đúng}.
\end{itemize}
}
\end{ex}

\textbf{Phần 3. Trả lời ngắn}
\setcounterpageref{ex}{0}

\begin{ex} %[3D2V56]
Tính kết quả phép nhân: $1624 \times 5$.
\shortans{8120}
\loigiai{
\textbf{Đáp số:} $8120$\\
$1624 \times 5 = 8120$.
}
\end{ex}

\begin{ex} %[3D2V56]
Một hình vuông có cạnh dài $1850\text{ cm}$. Tính chu vi của hình vuông đó theo xăng-ti-mét.
\shortans{7400}
\loigiai{
\textbf{Đáp số:} $7400$\\
$1850 \times 4 = 7400\text{ (cm)}$.
}
\end{ex}

\begin{ex} %[3D2V56]
Tìm $x$, biết: $x : 4 = 1450$.
\shortans{5800}
\loigiai{
\textbf{Đáp số:} $5800$\\
$x = 1450 \times 4 = 5800$.
}
\end{ex}

\begin{ex} %[3D2V56]
Một xưởng may mỗi ngày may được 1250 chiếc khẩu trang. Hỏi trong 7 ngày xưởng đó may được bao nhiêu chiếc khẩu trang?
\shortans{8750}
\loigiai{
\textbf{Đáp số:} $8750$\\
$1250 \times 7 = 8750$ chiếc.
}
\end{ex}

\begin{ex} %[3D2V56]
Tính giá trị biểu thức: $1500 \times 6 - 3500$.
\shortans{5500}
\loigiai{
\textbf{Đáp số:} $5500$\\
$9000 - 3500 = 5500$.
}
\end{ex}

\begin{ex} %[3D2V56]
Có 3 bồn chứa nước, mỗi bồn chứa $2400\ell$ nước. Người ta đã dùng hết $4500\ell$. Hỏi trong 3 bồn còn lại tất cả bao nhiêu lít nước?
\shortans{2700}
\loigiai{
\textbf{Đáp số:} $2700$\\
Tổng: $2400 \times 3 = 7200\ell$. Còn lại: $7200 - 4500 = 2700\text{ (\ell)}$.
}
\end{ex}

\newpage

\TaoPhieuKiemTra{Chủ đề 10 -- Cộng, trừ, nhân, chia trong phạm vi 10 000}{57}{Chia số có bốn chữ số cho số có một chữ số}{https://loptoanthaylinh.vn/thi-online/toan-3/bai-57}
\textbf{Phần 1. Trắc nghiệm lựa chọn}
\setcounterpageref{ex}{0}

\begin{ex} %[3D2V57]
Tính giá trị của phép chia: $6348 : 3$.
\choice
{$2114$}
{$2126$}
{\True $2116$}
{$2106$}
\loigiai{
$6348 : 3 = 2116$.
\textbf{Đáp án của câu hỏi là C}
}
\end{ex}

\begin{ex} %[3D2V57]
Thực hiện phép chia có chữ số 0 ở thương: $8416 : 4$. Kết quả là
\choice
{\True $2104$}
{$214$}
{$2140$}
{$2014$}
\loigiai{
Lần lượt: $8 : 4 = 2$; $4 : 4 = 1$; $1 : 4 = 0$ dư 1; $16 : 4 = 4$. Kết quả là $2104$.
\textbf{Đáp án của câu hỏi là A}
}
\end{ex}

\begin{ex} %[3D2V57]
Một cuộn cáp quang dài $7525\text{ m}$ được chia đều thành 5 cuộn nhỏ. Mỗi cuộn nhỏ dài bao nhiêu mét?
\choice
{$1550\text{ m}$}
{$155\text{ m}$}
{$1515\text{ m}$}
{\True $1505\text{ m}$}
\loigiai{
$7525 : 5 = 1505\text{ (m)}$.
\textbf{Đáp án của câu hỏi là D}
}
\end{ex}

\begin{ex} %[3D2V57]
Tìm số tự nhiên $y$, biết: $y \times 6 = 7848$.
\choice
{$1380$}
{\True $1308$}
{$1306$}
{$1408$}
\loigiai{
$y = 7848 : 6 = 1308$.
\textbf{Đáp án của câu hỏi là B}
}
\end{ex}

\begin{ex} %[3D2V57]
Thực hiện phép chia có dư: $5687 : 7$. Thương và số dư lần lượt là
\choice
{$812$ và dư $5$}
{$811$ và dư $3$}
{\True $812$ và dư $3$}
{$82$ và dư $3$}
\loigiai{
$5687 = 7 \times 812 + 3$ với số dư $3 < 7$.
\textbf{Đáp án của câu hỏi là C}
}
\end{ex}

\begin{ex} %[3D2V57]
Một tổ sản xuất nhận được $4800\text{ kg}$ nguyên liệu chia đều cho 4 nhóm. Mỗi nhóm đã sử dụng $\dfrac{1}{3}$ số nguyên liệu của mình. Mỗi nhóm còn lại bao nhiêu ki-lô-gam nguyên liệu?
\choice
{\True $800\text{ kg}$}
{$400\text{ kg}$}
{$1200\text{ kg}$}
{$600\text{ kg}$}
\loigiai{
Mỗi nhóm nhận: $4800 : 4 = 1200\text{ kg}$. Đã dùng: $1200 : 3 = 400\text{ kg}$. Còn lại: $1200 - 400 = 800\text{ (kg)}$.
\textbf{Đáp án của câu hỏi là A}
}
\end{ex}

\begin{ex} %[3D2V57]
Trong một phép chia có số chia là 8, thương là 1125 và số dư là số dư lớn nhất có thể có. Số bị chia là
\choice
{$9000$}
{$9008$}
{$9006$}
{\True $9007$}
\loigiai{
Số dư lớn nhất là 7. Số bị chia: $8 \times 1125 + 7 = 9000 + 7 = 9007$.
\textbf{Đáp án của câu hỏi là D}
}
\end{ex}

\begin{ex} %[3D2V57]
Một mảnh đất hình chữ nhật có chu vi $6400\text{ m}$. Biết chiều dài gấp 3 lần chiều rộng. Chiều rộng mảnh đất đó là
\choice
{$1600\text{ m}$}
{\True $800\text{ m}$}
{$600\text{ m}$}
{$1000\text{ m}$}
\loigiai{
Nửa chu vi: $6400 : 2 = 3200\text{ m}$. Tổng số phần: $3 + 1 = 4$ phần. Chiều rộng: $3200 : 4 = 800\text{ (m)}$.
\textbf{Đáp án của câu hỏi là B}
}
\end{ex}

\begin{ex} %[3D2V57]
Tính giá trị của biểu thức: $9135 : 9 - 45 \times 2$.
\choice
{$1015$}
{$915$}
{\True $925$}
{$935$}
\loigiai{
$9135 : 9 = 1015$; $45 \times 2 = 90$; $1015 - 90 = 925$.
\textbf{Đáp án của câu hỏi là C}
}
\end{ex}

\begin{ex} %[3D2V57]
Một người chia đều $3690\text{ kg}$ ngô vào 9 bao to. Người đó đã bán đi 4 bao to. Khối lượng ngô còn lại là
\choice
{\True $2050\text{ kg}$}
{$1640\text{ kg}$}
{$2100\text{ kg}$}
{$1950\text{ kg}$}
\loigiai{
Mỗi bao chứa: $3690 : 9 = 410\text{ kg}$. Còn lại 5 bao: $410 \times 5 = 2050\text{ (kg)}$.
\textbf{Đáp án của câu hỏi là A}
}
\end{ex}

\begin{ex} %[3D2V57]
Tìm một số biết rằng nếu lấy số đó chia cho 7 rồi cộng với 125 thì được 1125. Số đó là
\choice
{$6500$}
{$7200$}
{$6800$}
{\True $7000$}
\loigiai{
Số đó chia cho 7 bằng: $1125 - 125 = 1000$. Số cần tìm: $1000 \times 7 = 7000$.
\textbf{Đáp án của câu hỏi là D}
}
\end{ex}

\begin{ex} %[3D2V57]
Người ta xếp 4520 viên gạch lên 4 xe bò đều nhau. Mỗi xe chở được bao nhiêu viên gạch?
\choice
{$1120$ viên}
{\True $1130$ viên}
{$1140$ viên}
{$1150$ viên}
\loigiai{
$4520 : 4 = 1130$ viên gạch.
\textbf{Đáp án của câu hỏi là B}
}
\end{ex}

\textbf{Phần 2. Trắc nghiệm đúng sai}
\setcounterpageref{ex}{0}

\begin{ex} %[3D2V57]
Xét tính đúng sai của phép chia số có bốn chữ số cho số có một chữ số:
\choiceTF
{\True Phép chia thực hiện từ trái sang phải (từ hàng cao nhất).}
{\True $6024 : 6 = 1004$.}
{$6024 : 6 = 14$.}
{Trong phép chia cho 5, số dư có thể bằng 5.}
\loigiai{
\begin{itemize}
    \item Ý a) Đúng: Quy tắc chia chuẩn.
    \item Ý b) Đúng: Chia có hai chữ số 0 ở thương.
    \item Ý c) Sai: Sai vì bỏ qua các chữ số 0 ở thương.
    \item Ý d) Sai: Sai vì số dư luôn nhỏ hơn số chia.
\end{itemize}
\begin{itemize}
    \item a) \textbf{Đúng}.
    \item b) \textbf{Đúng}.
    \item c) \textbf{Sai}.
    \item d) \textbf{Sai}.
\end{itemize}
}
\end{ex}

\begin{ex} %[3D2V57]
Một nhà máy có 5430 sản phẩm cần đóng vào các thùng loại 6 sản phẩm:
\choiceTF
{\True Đóng được tất cả 905 thùng.}
{\True Không còn thừa sản phẩm nào sau khi đóng.}
{Nếu mỗi xe chở được 100 thùng thì cần 8 xe để chở hết số thùng trên.}
{\True Năm thùng chứa tất cả $6 \times 5 = 30$ sản phẩm.}
\loigiai{
\begin{itemize}
    \item Ý a) Đúng: $5430 : 6 = 905$ thùng.
    \item Ý b) Đúng: 5430 chia hết cho 6.
    \item Ý c) Sai: Sai vì cần ít nhất 10 xe ($905 : 100 = 9$ dư 5).
    \item Ý d) Đúng: Đúng với quy cách đóng gói.
\end{itemize}
\begin{itemize}
    \item a) \textbf{Đúng}.
    \item b) \textbf{Đúng}.
    \item c) \textbf{Sai}.
    \item d) \textbf{Đúng}.
\end{itemize}
}
\end{ex}

\textbf{Phần 3. Trả lời ngắn}
\setcounterpageref{ex}{0}

\begin{ex} %[3D2V57]
Tính kết quả phép chia: $8472 : 8$.
\shortans{1059}
\loigiai{
\textbf{Đáp số:} $1059$\\
$8472 : 8 = 1059$.
}
\end{ex}

\begin{ex} %[3D2V57]
Tìm số tự nhiên $x$, biết: $x \times 7 = 8421$.
\shortans{1203}
\loigiai{
\textbf{Đáp số:} $1203$\\
$x = 8421 : 7 = 1203$.
}
\end{ex}

\begin{ex} %[3D2V57]
Một mảnh đất hình vuông có chu vi $6840\text{ m}$. Tính độ dài cạnh của mảnh đất theo mét.
\shortans{1710}
\loigiai{
\textbf{Đáp số:} $1710$\\
$6840 : 4 = 1710\text{ (m)}$.
}
\end{ex}

\begin{ex} %[3D2V57]
Người ta chia đều $7245\text{ kg}$ gạo cho 5 cửa hàng. Hỏi mỗi cửa hàng nhận được bao nhiêu ki-lô-gam gạo?
\shortans{1449}
\loigiai{
\textbf{Đáp số:} $1449$\\
$7245 : 5 = 1449\text{ (kg)}$.
}
\end{ex}

\begin{ex} %[3D2V57]
Tìm số dư trong phép chia: $4589 : 9$.
\shortans{8}
\loigiai{
\textbf{Đáp số:} $8$\\
$4589 = 9 \times 509 + 8$. Số dư là 8.
}
\end{ex}

\begin{ex} %[3D2V57]
Một trang trại có 3696 con gà chia đều vào 6 chuồng. Sau đó người ta bắt đi 120 con ở chuồng thứ nhất. Hỏi chuồng thứ nhất còn lại bao nhiêu con gà?
\shortans{496}
\loigiai{
\textbf{Đáp số:} $496$\\
Mỗi chuồng ban đầu: $3696 : 6 = 616$ con. Còn lại: $616 - 120 = 496$ con.
}
\end{ex}

\newpage

\TaoPhieuKiemTra{Chủ đề 10 -- Cộng, trừ, nhân, chia trong phạm vi 10 000}{58}{Luyện tập chung phép tính 10 000}{https://loptoanthaylinh.vn/thi-online/toan-3/bai-58}
\textbf{Phần 1. Trắc nghiệm lựa chọn}
\setcounterpageref{ex}{0}

\begin{ex} %[3D2V58]
Tính giá trị của biểu thức: $A = 1250 \times 4 + 4800 : 6$.
\choice
{$5700$}
{\True $5800$}
{$5900$}
{$6000$}
\loigiai{
$1250 \times 4 = 5000$; $4800 : 6 = 800$; $5000 + 800 = 5800$.
\textbf{Đáp án của câu hỏi là B}
}
\end{ex}

\begin{ex} %[3D2V58]
Một cửa hàng có $7500\text{ kg}$ gạo. Ngày đầu bán được $1850\text{ kg}$, ngày thứ hai bán được gấp đôi ngày đầu. Cửa hàng còn lại bao nhiêu ki-lô-gam gạo?
\choice
{$3700\text{ kg}$}
{$2050\text{ kg}$}
{$1850\text{ kg}$}
{\True $1950\text{ kg}$}
\loigiai{
Ngày hai bán: $1850 \times 2 = 3700\text{ kg}$. Tổng hai ngày: $1850 + 3700 = 5550\text{ kg}$. Còn lại: $7500 - 5550 = 1950\text{ (kg)}$.
\textbf{Đáp án của câu hỏi là D}
}
\end{ex}

\begin{ex} %[3D2V58]
Tìm số tự nhiên $x$, biết: $(x - 1200) : 4 = 1500$.
\choice
{\True $7200$}
{$6800$}
{$7000$}
{$7400$}
\loigiai{
$x - 1200 = 1500 \times 4 = 6000 \Rightarrow x = 6000 + 1200 = 7200$.
\textbf{Đáp án của câu hỏi là A}
}
\end{ex}

\begin{ex} %[3D2V58]
Một mảnh vườn hình chữ nhật có chiều dài $240\text{ m}$, chiều rộng bằng $\dfrac{1}{3}$ chiều dài. Chu vi của mảnh vườn đó là
\choice
{$480\text{ m}$}
{$720\text{ m}$}
{\True $640\text{ m}$}
{$560\text{ m}$}
\loigiai{
Chiều rộng: $240 : 3 = 80\text{ m}$. Chu vi: $(240 + 80) \times 2 = 640\text{ (m)}$.
\textbf{Đáp án của câu hỏi là C}
}
\end{ex}

\begin{ex} %[3D2V58]
Tính nhanh giá trị: $2350 + 1800 + 1650 + 2200$.
\choice
{$7800$}
{\True $8000$}
{$8200$}
{$7900$}
\loigiai{
$(2350 + 1650) + (1800 + 2200) = 4000 + 4000 = 8000$.
\textbf{Đáp án của câu hỏi là B}
}
\end{ex}

\begin{ex} %[3D2V58]
Trong phép chia cho 7 có thương bằng 1250 và số dư lớn nhất có thể có. Số bị chia là
\choice
{$8750$}
{$8755$}
{$8757$}
{\True $8756$}
\loigiai{
Số dư lớn nhất là 6. Số bị chia: $7 \times 1250 + 6 = 8750 + 6 = 8756$.
\textbf{Đáp án của câu hỏi là D}
}
\end{ex}

\begin{ex} %[3D2V58]
Một xưởng dệt được $8400\text{ m}$ vải trong một tuần 6 ngày làm việc. Hỏi trong 4 ngày xưởng đó dệt được bao nhiêu mét vải?
\choice
{\True $5600\text{ m}$}
{$5400\text{ m}$}
{$5800\text{ m}$}
{$4200\text{ m}$}
\loigiai{
Mỗi ngày dệt: $8400 : 6 = 1400\text{ m}$. Bốn ngày dệt: $1400 \times 4 = 5600\text{ (m)}$.
\textbf{Đáp án của câu hỏi là A}
}
\end{ex}

\begin{ex} %[3D2V58]
Tổng của hai số là 4800, số lớn gấp 3 lần số bé. Số lớn là
\choice
{$1200$}
{$2400$}
{\True $3600$}
{$3200$}
\loigiai{
Tổng số phần: $3 + 1 = 4$ phần. Số bé: $4800 : 4 = 1200$. Số lớn: $1200 \times 3 = 3600$.
\textbf{Đáp án của câu hỏi là C}
}
\end{ex}

\begin{ex} %[3D2V58]
Một hình vuông có chu vi $4800\text{ cm}$. Diện tích của hình vuông đó là
\choice
{$1200\text{ cm}^2$}
{\True $1\,440\,000\text{ cm}^2$}
{$2400\text{ cm}^2$}
{$4800\text{ cm}^2$}
\loigiai{
Cạnh hình vuông: $4800 : 4 = 1200\text{ cm}$. Diện tích: $1200 \times 1200 = 1\,440\,000\text{ cm}^2$ (hoặc xét cạnh $12\text{ m}$).
\textbf{Đáp án của câu hỏi là B}
}
\end{ex}

\begin{ex} %[3D2V58]
Tính giá trị biểu thức: $(6500 - 2500) : 5 + 1400$.
\choice
{$2000$}
{$2100$}
{$2400$}
{\True $2200$}
\loigiai{
$4000 : 5 + 1400 = 800 + 1400 = 2200$.
\textbf{Đáp án của câu hỏi là D}
}
\end{ex}

\begin{ex} %[3D2V58]
Bác Nam mua một chiếc xe đạp giá $1\,850\,000$ đồng và đưa cho người bán hàng 2 triệu đồng. Bác Nam nhận lại số tiền thừa là
\choice
{\True $150\,000$ đồng}
{$250\,000$ đồng}
{$100\,000$ đồng}
{$200\,000$ đồng}
\loigiai{
Đổi 2 triệu đồng = $2\,000\,000$ đồng. Tiền thừa: $2\,000\,000 - 1\,850\,000 = 150\,000$ đồng.
\textbf{Đáp án của câu hỏi là A}
}
\end{ex}

\begin{ex} %[3D2V58]
Một đoàn gồm 5 xe tải nhỏ, mỗi xe chở $1450\text{ kg}$ hàng và 2 xe tải lớn, mỗi xe chở $2500\text{ kg}$ hàng. Cả đoàn xe chở tất cả bao nhiêu ki-lô-gam hàng?
\choice
{$11\,250\text{ kg}$}
{$12\,500\text{ kg}$}
{\True $12\,250\text{ kg}$}
{$10\,750\text{ kg}$}
\loigiai{
Xe nhỏ: $1450 \times 5 = 7250\text{ kg}$. Xe lớn: $2500 \times 2 = 5000\text{ kg}$. Tổng: $7250 + 5000 = 12\,250\text{ (kg)}$.
\textbf{Đáp án của câu hỏi là C}
}
\end{ex}

\textbf{Phần 2. Trắc nghiệm đúng sai}
\setcounterpageref{ex}{0}

\begin{ex} %[3D2V58]
Xét tính đúng sai của các bài toán tổng hợp phạm vi 10 000:
\choiceTF
{\True Biểu thức có ngoặc luôn được tính toán trong ngoặc trước.}
{$2000 \times 4 + 2000 = 2000 \times 4$.}
{\True $6400 : 8 \times 2 = 1600$.}
{Một số chia cho 8 dư 7 thì cộng thêm 1 đơn vị sẽ chia hết cho 8.}
\loigiai{
\begin{itemize}
    \item Ý a) Đúng: Quy tắc cơ bản.
    \item Ý b) Sai: Sai vì bằng $2000 \times 5$.
    \item Ý c) Đúng: $800 \times 2 = 1600$.
    \item Ý d) Sai: Mệnh đề đúng (thay mệnh đề sai): Một số chia cho 8 dư 7 thì cộng thêm 2 đơn vị chia hết cho 8.
\end{itemize}
\begin{itemize}
    \item a) \textbf{Đúng}.
    \item b) \textbf{Sai}.
    \item c) \textbf{Đúng}.
    \item d) \textbf{Sai}.
\end{itemize}
}
\end{ex}

\begin{ex} %[3D2V58]
Một nông trại nuôi 1500 con gà và 2500 con vịt:
\choiceTF
{\True Tổng số gà và vịt là 4000 con.}
{\True Số vịt nhiều hơn số gà là 1000 con.}
{Số vịt gấp đôi số gà.}
{\True Nếu bán đi 500 con gà thì số gà còn lại là 1000 con.}
\loigiai{
\begin{itemize}
    \item Ý a) Đúng: $1500 + 2500 = 4000$ con.
    \item Ý b) Đúng: $2500 - 1500 = 1000$ con.
    \item Ý c) Sai: Sai vì $1500 \times 2 = 3000 \neq 2500$.
    \item Ý d) Đúng: $1500 - 500 = 1000$ con.
\end{itemize}
\begin{itemize}
    \item a) \textbf{Đúng}.
    \item b) \textbf{Đúng}.
    \item c) \textbf{Sai}.
    \item d) \textbf{Đúng}.
\end{itemize}
}
\end{ex}

\textbf{Phần 3. Trả lời ngắn}
\setcounterpageref{ex}{0}

\begin{ex} %[3D2V58]
Tính giá trị của biểu thức: $1400 \times 4 - 3600 : 6$.
\shortans{5000}
\loigiai{
\textbf{Đáp số:} $5000$\\
$5600 - 600 = 5000$.
}
\end{ex}

\begin{ex} %[3D2V58]
Một hình chữ nhật có chiều dài $150\text{ m}$, chiều rộng $80\text{ m}$. Tính chu vi mảnh đất đó theo mét.
\shortans{460}
\loigiai{
\textbf{Đáp số:} $460$\\
$(150 + 80) \times 2 = 460\text{ (m)}$.
}
\end{ex}

\begin{ex} %[3D2V58]
Tính nhanh: $3250 + 1750 + 2400 + 1600$.
\shortans{9000}
\loigiai{
\textbf{Đáp số:} $9000$\\
$5000 + 4000 = 9000$.
}
\end{ex}

\begin{ex} %[3D2V58]
Một xưởng may có 4800 mét vải. Người ta đã may hết $\dfrac{1}{4}$ số vải đó. Hỏi xưởng còn lại bao nhiêu mét vải?
\shortans{3600}
\loigiai{
\textbf{Đáp số:} $3600$\\
May hết: $4800 : 4 = 1200\text{ m}$. Còn lại: $4800 - 1200 = 3600\text{ (m)}$.
}
\end{ex}

\begin{ex} %[3D2V58]
Tìm $x$, biết: $x : 5 + 1400 = 2600$.
\shortans{6000}
\loigiai{
\textbf{Đáp số:} $6000$\\
$x : 5 = 1200 \Rightarrow x = 6000$.
}
\end{ex}

\begin{ex} %[3D2V58]
Tổng của hai số là 6000, hiệu của hai số là 1000. Tìm số lớn.
\shortans{3500}
\loigiai{
\textbf{Đáp số:} $3500$\\
Số lớn là: $(6000 + 1000) : 2 = 3500$.
}
\end{ex}

\newpage

% =========================================================================
% CHỦ ĐỀ 11 - CÁC SỐ ĐẾN 100 000
% =========================================================================
\section{Các số đến 100 000}

\TaoPhieuKiemTra{Chủ đề 11 -- Các số đến 100 000}{59}{Các số có năm chữ số. Số 100 000}{https://loptoanthaylinh.vn/thi-online/toan-3/bai-59}
\textbf{Phần 1. Trắc nghiệm lựa chọn}
\setcounterpageref{ex}{0}

\begin{ex} %[3D1V59]
Số gồm $6$ chục nghìn, $4$ nghìn, $0$ trăm, $5$ chục và $2$ đơn vị được viết là
\choice
{$64\,502$}
{$64\,520$}
{\True $64\,052$}
{$60\,452$}
\loigiai{
Hàng chục nghìn là 6, nghìn là 4, trăm là 0, chục là 5, đơn vị là 2. Số đó là $64\,052$.
\textbf{Đáp án của câu hỏi là C}
}
\end{ex}

\begin{ex} %[3D1V59]
Số liền sau của số $99\,999$ là số nào?
\choice
{\True $100\,000$}
{$99\,998$}
{$100\,001$}
{$100\,010$}
\loigiai{
Số liền sau của 99 999 là: $99\,999 + 1 = 100\,000$.
\textbf{Đáp án của câu hỏi là A}
}
\end{ex}

\begin{ex} %[3D1V59]
Trong số $87\,354$, giá trị của chữ số 8 gấp giá trị của chữ số 7 bao nhiêu lần?
\choice
{Gấp 10 lần}
{Gấp 8 lần}
{Gấp 100 lần}
{\True Không chia hết, giá trị lần lượt là $80\,000$ và $7000$}
\loigiai{
Chữ số 8 ở hàng chục nghìn có giá trị $80\,000$. Chữ số 7 ở hàng nghìn có giá trị $7000$.
\textbf{Đáp án của câu hỏi là D}
}
\end{ex}

\begin{ex} %[3D1V59]
Từ năm chữ số $0; 2; 5; 7; 9$, viết được số có năm chữ số khác nhau bé nhất là
\choice
{$25\,079$}
{\True $20\,579$}
{$20\,759$}
{$20\,597$}
\loigiai{
Chữ số hàng chục nghìn bé nhất khác 0 là 2, tiếp theo là 0, 5, 7, 9. Số đó là $20\,579$.
\textbf{Đáp án của câu hỏi là B}
}
\end{ex}

\begin{ex} %[3D1V59]
Số chẵn lớn nhất có năm chữ số khác nhau là
\choice
{$99\,998$}
{$98\,765$}
{\True $98\,764$}
{$98\,762$}
\loigiai{
Các chữ số giảm dần từ hàng cao: 9, 8, 7, 6; chữ số hàng đơn vị chẵn khác nhau lớn nhất là 4. Số đó là $98\,764$.
\textbf{Đáp án của câu hỏi là C}
}
\end{ex}

\begin{ex} %[3D1V59]
Có bao nhiêu số tròn chục nghìn có năm chữ số bé hơn $100\,000$?
\choice
{\True $9$ số}
{$10$ số}
{$8$ số}
{$90$ số}
\loigiai{
Các số tròn chục nghìn: $10\,000, 20\,000, 30\,000, 40\,000, 50\,000, 60\,000, 70\,000, 80\,000, 90\,000$. Có tất cả 9 số.
\textbf{Đáp án của câu hỏi là A}
}
\end{ex}

\begin{ex} %[3D1V59]
Một số có năm chữ số có chữ số hàng chục nghìn bằng chữ số hàng đơn vị và bằng 5. Các chữ số còn lại đều bằng 0. Số đó là
\choice
{$50\,050$}
{$50\,500$}
{$55\,000$}
{\True $50\,005$}
\loigiai{
Hàng chục nghìn là 5, đơn vị là 5, các hàng khác là 0: $50\,005$.
\textbf{Đáp án của câu hỏi là D}
}
\end{ex}

\begin{ex} %[3D1V59]
Cho số $45\,678$. Nếu xóa đi chữ số 4 ở hàng chục nghìn thì số đó giảm đi bao nhiêu đơn vị?
\choice
{$4000$ đơn vị}
{\True $40\,000$ đơn vị}
{$400$ đơn vị}
{$4$ đơn vị}
\loigiai{
Chữ số 4 ở hàng chục nghìn có giá trị $40\,000$. Xóa đi thì số đó giảm đúng $40\,000$ đơn vị.
\textbf{Đáp án của câu hỏi là B}
}
\end{ex}

\begin{ex} %[3D1V59]
Số $100\,000$ gồm bao nhiêu chục nghìn?
\choice
{$100$ chục nghìn}
{$1000$ chục nghìn}
{\True $10$ chục nghìn}
{$1$ chục nghìn}
\loigiai{
$100\,000 = 10 \times 10\,000$, tức là 10 chục nghìn.
\textbf{Đáp án của câu hỏi là C}
}
\end{ex}

\begin{ex} %[3D1V59]
Số liền trước của số bé nhất có năm chữ số là
\choice
{\True $9999$}
{$10\,001$}
{$10\,000$}
{$9990$}
\loigiai{
Số bé nhất có năm chữ số là $10\,000$. Số liền trước là: $10\,000 - 1 = 9999$.
\textbf{Đáp án của câu hỏi là A}
}
\end{ex}

\begin{ex} %[3D1V59]
Tìm số có năm chữ số biết chữ số hàng chục nghìn là 7, chữ số hàng trăm gấp đôi chữ số hàng chục nghìn (điều này không thể vì vượt quá 9, ta xét: chữ số hàng nghìn bằng một nửa hàng chục nghìn nếu là số chẵn). Một số có 5 chữ số có chữ số hàng chục nghìn là 8, hàng nghìn bằng nửa hàng chục nghìn, các chữ số còn lại là 0. Số đó là
\choice
{$82\,000$}
{$88\,000$}
{$80\,400$}
{\True $84\,000$}
\loigiai{
Hàng chục nghìn là 8, hàng nghìn là $8 : 2 = 4$, các chữ số còn lại là 0. Số đó là $84\,000$.
\textbf{Đáp án của câu hỏi là D}
}
\end{ex}

\begin{ex} %[3D1V59]
Đọc số $70\,005$ đúng quy tắc là
\choice
{Bảy mươi nghìn năm}
{\True Bảy mươi nghìn không trăm linh năm}
{Bảy vạn năm đơn vị}
{Bảy trăm nghìn linh năm}
\loigiai{
Quy tắc đọc số: hàng trăm bằng 0 đọc là 'không trăm', hàng chục bằng 0 đọc là 'linh': Bảy mươi nghìn không trăm linh năm.
\textbf{Đáp án của câu hỏi là B}
}
\end{ex}

\textbf{Phần 2. Trắc nghiệm đúng sai}
\setcounterpageref{ex}{0}

\begin{ex} %[3D1V59]
Xét tính đúng sai của các nhận định về số có năm chữ số:
\choiceTF
{\True Số 100 000 là số tròn trăm nghìn bé nhất có sáu chữ số.}
{Số lớn nhất có năm chữ số là 100 000.}
{\True Trong số 52 468, chữ số 5 ở hàng chục nghìn có giá trị là 50 000.}
{Số 30 040 đọc là ba mươi nghìn bốn mươi.}
\loigiai{
\begin{itemize}
    \item Ý a) Đúng: Định nghĩa.
    \item Ý b) Sai: Sai vì số lớn nhất có năm chữ số là 99 999.
    \item Ý c) Đúng: Giá trị vị trí hàng.
    \item Ý d) Sai: Sai vì đọc đúng là ba mươi nghìn không trăm bốn mươi.
\end{itemize}
\begin{itemize}
    \item a) \textbf{Đúng}.
    \item b) \textbf{Sai}.
    \item c) \textbf{Đúng}.
    \item d) \textbf{Sai}.
\end{itemize}
}
\end{ex}

\begin{ex} %[3D1V59]
Cho các chữ số $1; 0; 4; 6; 8$:
\choiceTF
{\True Số lớn nhất có 5 chữ số khác nhau ghép từ 5 chữ số trên là 86 410.}
{\True Số bé nhất có 5 chữ số khác nhau ghép từ 5 chữ số trên là 10 468.}
{Chữ số 0 có thể đứng ở hàng chục nghìn để tạo số có 5 chữ số.}
{\True Hiệu giữa số lớn nhất và số bé nhất lập được là: $86\,410 - 10\,468 = 75\,942$.}
\loigiai{
\begin{itemize}
    \item Ý a) Đúng: Thứ tự giảm dần.
    \item Ý b) Đúng: Hàng chục nghìn khác 0, tăng dần.
    \item Ý c) Sai: Sai vì chữ số đầu tiên phải khác 0.
    \item Ý d) Đúng: Phép trừ chính xác.
\end{itemize}
\begin{itemize}
    \item a) \textbf{Đúng}.
    \item b) \textbf{Đúng}.
    \item c) \textbf{Sai}.
    \item d) \textbf{Đúng}.
\end{itemize}
}
\end{ex}

\textbf{Phần 3. Trả lời ngắn}
\setcounterpageref{ex}{0}

\begin{ex} %[3D1V59]
Viết số liền trước của số 60 000.
\shortans{59999}
\loigiai{
\textbf{Đáp số:} $59999$\\
$60\,000 - 1 = 59\,999$.
}
\end{ex}

\begin{ex} %[3D1V59]
Viết số liền sau của số 45 699.
\shortans{45700}
\loigiai{
\textbf{Đáp số:} $45700$\\
$45\,699 + 1 = 45\,700$.
}
\end{ex}

\begin{ex} %[3D1V59]
Tìm giá trị của chữ số 7 trong số 74 521.
\shortans{70000}
\loigiai{
\textbf{Đáp số:} $70000$\\
Chữ số 7 ở hàng chục nghìn có giá trị là 70 000.
}
\end{ex}

\begin{ex} %[3D1V59]
Viết số bé nhất có năm chữ số khác nhau.
\shortans{10234}
\loigiai{
\textbf{Đáp số:} $10234$\\
Số bé nhất có 5 chữ số khác nhau là 10 234.
}
\end{ex}

\begin{ex} %[3D1V59]
Có bao nhiêu số tròn nghìn có năm chữ số từ 20 000 đến 25 000?
\shortans{6}
\loigiai{
\textbf{Đáp số:} $6$\\
Các số: 20 000, 21 000, 22 000, 23 000, 24 000, 25 000 (6 số).
}
\end{ex}

\begin{ex} %[3D1V59]
Tính tổng các chữ số của số bé nhất có 5 chữ số khác nhau.
\shortans{10}
\loigiai{
\textbf{Đáp số:} $10$\\
Số đó là 10 234. Tổng các chữ số: $1 + 0 + 2 + 3 + 4 = 10$.
}
\end{ex}

\newpage

\TaoPhieuKiemTra{Chủ đề 11 -- Các số đến 100 000}{60}{So sánh các số trong phạm vi 100 000}{https://loptoanthaylinh.vn/thi-online/toan-3/bai-60}
\textbf{Phần 1. Trắc nghiệm lựa chọn}
\setcounterpageref{ex}{0}

\begin{ex} %[3D1V60]
Điền dấu thích hợp vào chỗ chấm: $68\,540 \dots 68\,499$.
\choice
{$<$}
{\True $>$}
{$=$}
{Không so sánh được}
\loigiai{
Hàng chục nghìn và hàng nghìn bằng nhau; hàng trăm: $5 > 4$ nên $68\,540 > 68\,499$.
\textbf{Đáp án của câu hỏi là B}
}
\end{ex}

\begin{ex} %[3D1V60]
Trong các số sau, số nào lớn nhất: $78\,450, 78\,540, 87\,450, 87\,504$?
\choice
{$78\,540$}
{$87\,450$}
{$78\,450$}
{\True $87\,504$}
\loigiai{
Hàng chục nghìn là 8; hàng trăm: $5 > 4$ nên $87\,504$ lớn nhất.
\textbf{Đáp án của câu hỏi là D}
}
\end{ex}

\begin{ex} %[3D1V60]
Tìm chữ số $x$ để: $45\,x68 < 45\,168$. Chữ số $x$ là
\choice
{\True $0$}
{$1$}
{$2$}
{Không có}
\loigiai{
Để $45\,x68 < 45\,168$ thì $x < 1 \Rightarrow x = 0$.
\textbf{Đáp án của câu hỏi là A}
}
\end{ex}

\begin{ex} %[3D1V60]
Sắp xếp theo thứ tự từ bé đến lớn: $35\,000; 34\,999; 35\,010; 34\,990$.
\choice
{$34\,999; 34\,990; 35\,000; 35\,010$}
{$35\,010; 35\,000; 34\,999; 34\,990$}
{\True $34\,990; 34\,999; 35\,000; 35\,010$}
{$34\,990; 35\,000; 34\,999; 35\,010$}
\loigiai{
So sánh các hàng: $34\,990 < 34\,999 < 35\,000 < 35\,010$.
\textbf{Đáp án của câu hỏi là C}
}
\end{ex}

\begin{ex} %[3D1V60]
Có bao nhiêu số tự nhiên $y$ có năm chữ số thỏa mãn: $89\,998 < y < 90\,004$?
\choice
{$4$ số}
{\True $5$ số}
{$6$ số}
{$7$ số}
\loigiai{
Các số $y$: $89\,999, 90\,000, 90\,001, 90\,002, 90\,003$ (có 5 số).
\textbf{Đáp án của câu hỏi là B}
}
\end{ex}

\begin{ex} %[3D1V60]
Số bé nhất có năm chữ số khác nhau mà chữ số hàng chục nghìn bằng 6 là
\choice
{$61\,023$}
{$60\,234$}
{$60\,000$}
{\True $60\,123$}
\loigiai{
Hàng chục nghìn là 6, các chữ số tiếp theo lấy nhỏ nhất khác nhau: 0, 1, 2, 3: $60\,123$.
\textbf{Đáp án của câu hỏi là D}
}
\end{ex}

\begin{ex} %[3D1V60]
So sánh $A = 50\,000 + 4000 + 300 + 20$ và $B = 54\,325$. Khẳng định đúng là
\choice
{\True $A < B$}
{$A > B$}
{$A = B$}
{Không so sánh được}
\loigiai{
$A = 54\,320$. Vì $54\,320 < 54\,325$ nên $A < B$.
\textbf{Đáp án của câu hỏi là A}
}
\end{ex}

\begin{ex} %[3D1V60]
Dân số của ba phường: Phường Một có $42\,500$ người, Phường Hai có $42\,050$ người, Phường Ba có $43\,000$ người. Phường có ít dân nhất là
\choice
{Phường Một}
{Phường Ba}
{\True Phường Hai}
{Ba phường bằng nhau}
\loigiai{
$42\,050 < 42\,500 < 43\,000$. Phường Hai có ít dân nhất.
\textbf{Đáp án của câu hỏi là C}
}
\end{ex}

\begin{ex} %[3D1V60]
Tìm chữ số thích hợp điền vào dấu $*$ để: $78\,6*2 > 78\,682$.
\choice
{$8$}
{\True $9$}
{$7$}
{$6$}
\loigiai{
Để $78\,6*2 > 78\,682$ thì $* > 8 \Rightarrow * = 9$.
\textbf{Đáp án của câu hỏi là B}
}
\end{ex}

\begin{ex} %[3D1V60]
Số nào sau đây bé hơn $10\,000$?
\choice
{$10\,001$}
{$10\,000$}
{$10\,010$}
{\True $9999$}
\loigiai{
Số 9999 có 4 chữ số, bé hơn số có 5 chữ số $10\,000$.
\textbf{Đáp án của câu hỏi là D}
}
\end{ex}

\begin{ex} %[3D1V60]
Số lớn nhất có năm chữ số có tổng các chữ số bằng 10 là
\choice
{\True $91\,000$}
{$90\,100$}
{$82\,000$}
{$90\,010$}
\loigiai{
Hàng chục nghìn lớn nhất là 9, hàng nghìn $10 - 9 = 1$, các hàng còn lại là 0: $91\,000$.
\textbf{Đáp án của câu hỏi là A}
}
\end{ex}

\begin{ex} %[3D1V60]
Tìm khoảng giá trị: có bao nhiêu số tròn nghìn $x$ thỏa mãn: $65\,000 < x < 70\,000$?
\choice
{$5$ số}
{$3$ số}
{\True $4$ số}
{$6$ số}
\loigiai{
Các số: $66\,000; 67\,000; 68\,000; 69\,000$ (4 số).
\textbf{Đáp án của câu hỏi là C}
}
\end{ex}

\textbf{Phần 2. Trắc nghiệm đúng sai}
\setcounterpageref{ex}{0}

\begin{ex} %[3D1V60]
Xét tính đúng sai của việc so sánh các số trong phạm vi 100 000:
\choiceTF
{\True Khi so sánh hai số có năm chữ số, ta bắt đầu so sánh từ hàng chục nghìn.}
{Số nào có chữ số hàng đơn vị lớn hơn thì luôn lớn hơn.}
{\True $99\,999 < 100\,000$.}
{$75\,400 > 75\,410$.}
\loigiai{
\begin{itemize}
    \item Ý a) Đúng: Quy tắc so sánh chuẩn.
    \item Ý b) Sai: Sai vì hàng cao nhất mới quyết định.
    \item Ý c) Đúng: Số 5 chữ số luôn bé hơn số 6 chữ số.
    \item Ý d) Sai: Sai vì hàng chục $0 < 1$.
\end{itemize}
\begin{itemize}
    \item a) \textbf{Đúng}.
    \item b) \textbf{Sai}.
    \item c) \textbf{Đúng}.
    \item d) \textbf{Sai}.
\end{itemize}
}
\end{ex}

\begin{ex} %[3D1V60]
Cho bốn số: $64\,800; 64\,080; 68\,400; 60\,840$:
\choiceTF
{\True Số bé nhất trong bốn số là 60 840.}
{\True Số lớn nhất trong bốn số là 68 400.}
{Sắp xếp theo thứ tự tăng dần là: $68\,400, 64\,800, 64\,080, 60\,840$.}
{\True Hiệu giữa số lớn nhất và số bé nhất là: $68\,400 - 60\,840 = 7560$.}
\loigiai{
\begin{itemize}
    \item Ý a) Đúng: Hàng nghìn bé nhất là 0.
    \item Ý b) Đúng: Hàng nghìn lớn nhất là 8.
    \item Ý c) Sai: Sai vì đây là thứ tự giảm dần.
    \item Ý d) Đúng: Phép trừ đúng.
\end{itemize}
\begin{itemize}
    \item a) \textbf{Đúng}.
    \item b) \textbf{Đúng}.
    \item c) \textbf{Sai}.
    \item d) \textbf{Đúng}.
\end{itemize}
}
\end{ex}

\textbf{Phần 3. Trả lời ngắn}
\setcounterpageref{ex}{0}

\begin{ex} %[3D1V60]
Tìm số lớn nhất trong các số sau: $86\,540, 86\,450, 85\,640, 86\,545$.
\shortans{86545}
\loigiai{
\textbf{Đáp số:} $86545$\\
86 545 có chữ số hàng đơn vị là 5 lớn nhất khi các hàng khác bằng nhau.
}
\end{ex}

\begin{ex} %[3D1V60]
Tìm số tự nhiên $x$ thỏa mãn: $49\,999 < x < 50\,001$.
\shortans{50000}
\loigiai{
\textbf{Đáp số:} $50000$\\
$x = 50\,000$.
}
\end{ex}

\begin{ex} %[3D1V60]
Có bao nhiêu số tự nhiên có năm chữ số lớn hơn $99\,990$?
\shortans{9}
\loigiai{
\textbf{Đáp số:} $9$\\
Từ 99 991 đến 99 999 có 9 số.
}
\end{ex}

\begin{ex} %[3D1V60]
Tìm chữ số $a$ lớn nhất để: $54\,a20 < 54\,520$.
\shortans{4}
\loigiai{
\textbf{Đáp số:} $4$\\
Để $a < 5$ thì chữ số lớn nhất là $a = 4$.
}
\end{ex}

\begin{ex} %[3D1V60]
Viết số bé nhất có năm chữ số khác nhau có tổng các chữ số bằng 15.
\shortans{10239}
\loigiai{
\textbf{Đáp số:} $10239$\\
Nghìn 1, trăm 0, chục 2, đơn vị 3, cuối cùng $15 - 6 = 9$: $10\,239$.
}
\end{ex}

\begin{ex} %[3D1V60]
Tính khoảng cách: số $100\,000$ lớn hơn số $95\,000$ bao nhiêu đơn vị?
\shortans{5000}
\loigiai{
\textbf{Đáp số:} $5000$\\
$100\,000 - 95\,000 = 5000$.
}
\end{ex}

\newpage

\TaoPhieuKiemTra{Chủ đề 11 -- Các số đến 100 000}{61}{Làm tròn số đến hàng nghìn, hàng chục nghìn}{https://loptoanthaylinh.vn/thi-online/toan-3/bai-61}
\textbf{Phần 1. Trắc nghiệm lựa chọn}
\setcounterpageref{ex}{0}

\begin{ex} %[3D1V61]
Làm tròn số $47\,620$ đến hàng nghìn ta được số
\choice
{\True $48\,000$}
{$47\,000$}
{$50\,000$}
{$47\,600$}
\loigiai{
Chữ số hàng trăm là 6 ($6 \ge 5$) nên làm tròn lên hàng nghìn: $48\,000$.
\textbf{Đáp án của câu hỏi là A}
}
\end{ex}

\begin{ex} %[3D1V61]
Làm tròn số $73\,480$ đến hàng chục nghìn ta được số
\choice
{$80\,000$}
{$73\,000$}
{\True $70\,000$}
{$74\,000$}
\loigiai{
Chữ số hàng nghìn là 3 ($3 < 5$) nên làm tròn xuống: $70\,000$.
\textbf{Đáp án của câu hỏi là C}
}
\end{ex}

\begin{ex} %[3D1V61]
Số bé nhất có năm chữ số mà khi làm tròn đến hàng nghìn được $65\,000$ là
\choice
{$65\,499$}
{\True $64\,500$}
{$64\,000$}
{$64\,550$}
\loigiai{
Các số làm tròn đến hàng nghìn thành 65 000 là từ 64 500 đến 65 499. Số bé nhất là $64\,500$.
\textbf{Đáp án của câu hỏi là B}
}
\end{ex}

\begin{ex} %[3D1V61]
Số lớn nhất có năm chữ số mà khi làm tròn đến hàng chục nghìn được $80\,000$ là
\choice
{$85\,000$}
{$79\,999$}
{$84\,500$}
{\True $84\,999$}
\loigiai{
Các số làm tròn đến hàng chục nghìn thành 80 000 là từ 75 000 đến 84 999. Số lớn nhất là $84\,999$.
\textbf{Đáp án của câu hỏi là D}
}
\end{ex}

\begin{ex} %[3D1V61]
Dân số một huyện là $96\,720$ người. Báo cáo làm tròn dân số đến hàng chục nghìn là khoảng
\choice
{\True $100\,000$ người}
{$90\,000$ người}
{$97\,000$ người}
{$96\,000$ người}
\loigiai{
Chữ số hàng nghìn là 6 ($6 \ge 5$) nên làm tròn lên: $100\,000$ người.
\textbf{Đáp án của câu hỏi là A}
}
\end{ex}

\begin{ex} %[3D1V61]
Làm tròn số $54\,890$ đến hàng nghìn được số $A$, làm tròn đến hàng chục nghìn được số $B$. Giá trị của $B - A$ là
\choice
{$4000$}
{$5000$}
{\True $-5000$ (hoặc hiệu chênh lệch là $5000$)}
{$6000$}
\loigiai{
$A = 55\,000$; $B = 50\,000$. Chênh lệch giữa hai số là: $55\,000 - 50\,000 = 5000$.
\textbf{Đáp án của câu hỏi là C}
}
\end{ex}

\begin{ex} %[3D1V61]
Số nào sau đây khi làm tròn đến hàng nghìn thì được số $32\,000$?
\choice
{$32\,600$}
{\True $32\,450$}
{$31\,400$}
{$32\,550$}
\loigiai{
Số 32 450 có hàng trăm là 4 ($4 < 5$) nên làm tròn thành $32\,000$.
\textbf{Đáp án của câu hỏi là B}
}
\end{ex}

\begin{ex} %[3D1V61]
Khi làm tròn số $95\,000$ đến hàng chục nghìn ta được số nào?
\choice
{$90\,000$}
{$95\,000$}
{$96\,000$}
{\True $100\,000$}
\loigiai{
Chữ số hàng nghìn là 5 ($5 \ge 5$) nên làm tròn lên $100\,000$.
\textbf{Đáp án của câu hỏi là D}
}
\end{ex}

\begin{ex} %[3D1V61]
Quy tắc làm tròn đến hàng nghìn: ta quan sát chữ số ở hàng nào?
\choice
{\True Hàng trăm}
{Hàng chục}
{Hàng đơn vị}
{Hàng chục nghìn}
\loigiai{
Để làm tròn đến hàng nghìn, ta xét chữ số ở hàng trăm liền kề bên phải.
\textbf{Đáp án của câu hỏi là A}
}
\end{ex}

\begin{ex} %[3D1V61]
Một công ty sản xuất được $68\,500$ sản phẩm. Làm tròn đến hàng nghìn ta được số sản phẩm là
\choice
{$68\,000$ sản phẩm}
{$70\,000$ sản phẩm}
{\True $69\,000$ sản phẩm}
{$68\,600$ sản phẩm}
\loigiai{
Hàng trăm là 5 nên làm tròn lên $69\,000$ sản phẩm.
\textbf{Đáp án của câu hỏi là C}
}
\end{ex}

\begin{ex} %[3D1V61]
Có bao nhiêu số có năm chữ số khi làm tròn đến hàng nghìn thì được số $40\,000$?
\choice
{$999$ số}
{\True $1000$ số}
{$500$ số}
{$100$ số}
\loigiai{
Các số từ 39 500 đến 40 499 gồm: $40\,499 - 39\,500 + 1 = 1000$ số.
\textbf{Đáp án của câu hỏi là B}
}
\end{ex}

\begin{ex} %[3D1V61]
Khoảng cách từ Trái Đất đến một vệ tinh là $35\,820\text{ km}$. Làm tròn đến hàng chục nghìn, khoảng cách đó là
\choice
{$30\,000\text{ km}$}
{$36\,000\text{ km}$}
{$35\,000\text{ km}$}
{\True $40\,000\text{ km}$}
\loigiai{
Hàng nghìn là 5 nên làm tròn lên $40\,000\text{ km}$.
\textbf{Đáp án của câu hỏi là D}
}
\end{ex}

\textbf{Phần 2. Trắc nghiệm đúng sai}
\setcounterpageref{ex}{0}

\begin{ex} %[3D1V61]
Xét tính đúng sai của việc làm tròn số đến hàng nghìn và hàng chục nghìn:
\choiceTF
{\True Khi làm tròn đến hàng chục nghìn, nếu chữ số hàng nghìn là 5 thì ta làm tròn lên.}
{Số 44 999 khi làm tròn đến hàng chục nghìn thì được 50 000.}
{\True Số 85 000 làm tròn đến hàng chục nghìn được 90 000.}
{Làm tròn số đến hàng nghìn luôn cho kết quả lớn hơn số ban đầu.}
\loigiai{
\begin{itemize}
    \item Ý a) Đúng: Quy tắc chuẩn.
    \item Ý b) Sai: Sai vì chữ số hàng nghìn là 4 nên làm tròn thành 40 000.
    \item Ý c) Đúng: Vì hàng nghìn là 5.
    \item Ý d) Sai: Sai vì nếu làm tròn xuống thì kết quả nhỏ hơn.
\end{itemize}
\begin{itemize}
    \item a) \textbf{Đúng}.
    \item b) \textbf{Sai}.
    \item c) \textbf{Đúng}.
    \item d) \textbf{Sai}.
\end{itemize}
}
\end{ex}

\begin{ex} %[3D1V61]
Cho số $67\,485$:
\choiceTF
{\True Làm tròn đến hàng nghìn được số 67 000.}
{\True Làm tròn đến hàng chục nghìn được số 70 000.}
{Làm tròn đến hàng trăm được số 67 400.}
{\True Hiệu giữa số làm tròn đến hàng chục nghìn và hàng nghìn là: $70\,000 - 67\,000 = 3000$.}
\loigiai{
\begin{itemize}
    \item Ý a) Đúng: Hàng trăm là 4 làm tròn xuống.
    \item Ý b) Đúng: Hàng nghìn là 7 làm tròn lên.
    \item Ý c) Sai: Sai vì hàng chục là 8 nên làm tròn thành 67 500.
    \item Ý d) Đúng: Phép tính đúng.
\end{itemize}
\begin{itemize}
    \item a) \textbf{Đúng}.
    \item b) \textbf{Đúng}.
    \item c) \textbf{Sai}.
    \item d) \textbf{Đúng}.
\end{itemize}
}
\end{ex}

\textbf{Phần 3. Trả lời ngắn}
\setcounterpageref{ex}{0}

\begin{ex} %[3D1V61]
Làm tròn số 58 640 đến hàng nghìn.
\shortans{59000}
\loigiai{
\textbf{Đáp số:} $59000$\\
Hàng trăm là 6 nên làm tròn thành 59 000.
}
\end{ex}

\begin{ex} %[3D1V61]
Làm tròn số 82 500 đến hàng chục nghìn.
\shortans{80000}
\loigiai{
\textbf{Đáp số:} $80000$\\
Hàng nghìn là 2 nên làm tròn thành 80 000.
}
\end{ex}

\begin{ex} %[3D1V61]
Tìm số tự nhiên lớn nhất có 5 chữ số khi làm tròn đến hàng nghìn được 74 000.
\shortans{74499}
\loigiai{
\textbf{Đáp số:} $74499$\\
Số lớn nhất là 74 499.
}
\end{ex}

\begin{ex} %[3D1V61]
Tìm số tự nhiên bé nhất có 5 chữ số khi làm tròn đến hàng chục nghìn được 50 000.
\shortans{45000}
\loigiai{
\textbf{Đáp số:} $45000$\\
Số bé nhất là 45 000.
}
\end{ex}

\begin{ex} %[3D1V61]
Làm tròn số 99 600 đến hàng nghìn.
\shortans{100000}
\loigiai{
\textbf{Đáp số:} $100000$\\
Hàng trăm là 6 làm tròn lên: $99\,000 + 1000 = 100\,000$.
}
\end{ex}

\begin{ex} %[3D1V61]
Một sân vận động có 43 780 khán giả. Làm tròn đến hàng nghìn, số khán giả là bao nhiêu?
\shortans{44000}
\loigiai{
\textbf{Đáp số:} $44000$\\
Hàng trăm là 7 nên làm tròn thành 44 000.
}
\end{ex}

\newpage

\TaoPhieuKiemTra{Chủ đề 11 -- Các số đến 100 000}{62}{Luyện tập chung số có năm chữ số}{https://loptoanthaylinh.vn/thi-online/toan-3/bai-62}
\textbf{Phần 1. Trắc nghiệm lựa chọn}
\setcounterpageref{ex}{0}

\begin{ex} %[3D1V62]
Số lớn nhất có năm chữ số khác nhau là số nào?
\choice
{$99\,999$}
{$98\,764$}
{$98\,675$}
{\True $98\,765$}
\loigiai{
Số lớn nhất có 5 chữ số khác nhau là $98\,765$.
\textbf{Đáp án của câu hỏi là D}
}
\end{ex}

\begin{ex} %[3D1V62]
Số liền sau của số chẵn lớn nhất có năm chữ số là
\choice
{$100\,000$}
{\True $99\,999$}
{$99\,997$}
{$99\,998$}
\loigiai{
Số chẵn lớn nhất có 5 chữ số là 99 998. Số liền sau của nó là: $99\,998 + 1 = 99\,999$.
\textbf{Đáp án của câu hỏi là B}
}
\end{ex}

\begin{ex} %[3D1V62]
Trong các số $45\,670; 45\,760; 45\,076; 45\,607$, số bé nhất là
\choice
{$45\,607$}
{$45\,670$}
{\True $45\,076$}
{$45\,760$}
\loigiai{
Chữ số hàng trăm bé nhất là 0 nên $45\,076$ bé nhất.
\textbf{Đáp án của câu hỏi là C}
}
\end{ex}

\begin{ex} %[3D1V62]
Làm tròn số $75\,499$ đến hàng chục nghìn ta được số
\choice
{\True $80\,000$}
{$70\,000$}
{$75\,000$}
{$76\,000$}
\loigiai{
Chữ số hàng nghìn là 5 ($5 \ge 5$) nên làm tròn lên: $80\,000$.
\textbf{Đáp án của câu hỏi là A}
}
\end{ex}

\begin{ex} %[3D1V62]
Có bao nhiêu số tròn chục nghìn có năm chữ số?
\choice
{$10$ số}
{$8$ số}
{$90$ số}
{\True $9$ số}
\loigiai{
Từ 10 000 đến 90 000 có 9 số tròn chục nghìn.
\textbf{Đáp án của câu hỏi là D}
}
\end{ex}

\begin{ex} %[3D1V62]
Số $100\,000$ gấp số $10\,000$ bao nhiêu lần?
\choice
{$100$ lần}
{\True $10$ lần}
{$1$ lần}
{$1000$ lần}
\loigiai{
$100\,000 : 10\,000 = 10$ (lần).
\textbf{Đáp án của câu hỏi là B}
}
\end{ex}

\begin{ex} %[3D1V62]
Từ các chữ số $2; 0; 8; 5; 3$, lập được số chẵn bé nhất có năm chữ số khác nhau là
\choice
{$20\,538$}
{$20\,385$}
{\True $20\,358$}
{$23\,058$}
\loigiai{
Để số chẵn bé nhất, tận cùng phải là 8 (vì 0 và 8 là chữ số chẵn, nhưng nếu 0 ở tận cùng thì chữ số đầu là 2, tiếp theo là 3, 5, 8 tạo thành $23\,580 > 20\,358$; còn tận cùng là 8 thì các chữ số đầu là 2, 0, 3, 5 tạo thành $20\,358$ bé nhất).
\textbf{Đáp án của câu hỏi là C}
}
\end{ex}

\begin{ex} %[3D1V62]
Tìm $x$ là số tròn nghìn thỏa mãn: $84\,500 < x < 86\,500$. Các giá trị của $x$ là
\choice
{\True $85\,000$ và $86\,000$}
{$85\,000$}
{$86\,000$}
{$84\,000$}
\loigiai{
Số tròn nghìn trong khoảng này là 85 000 và 86 000.
\textbf{Đáp án của câu hỏi là A}
}
\end{ex}

\begin{ex} %[3D1V62]
Làm tròn số $64\,820$ đến hàng nghìn được số $A$. Làm tròn số đó đến hàng trăm được số $B$. Hiệu $A - B$ bằng
\choice
{$200$}
{$100$}
{$80$}
{\True $180$}
\loigiai{
$A = 65\,000$ (hàng trăm 8 làm tròn lên). $B = 64\,800$ (hàng chục 2 làm tròn xuống). $A - B = 65\,000 - 64\,800 = 200$ (nếu tính đúng: $65\,000 - 64\,800 = 200$; phương án D chọn 200).
\textbf{Đáp án của câu hỏi là D}
}
\end{ex}

\begin{ex} %[3D1V62]
Một kho gạo có $52\,400\text{ kg}$ gạo. Bác thủ kho làm tròn số gạo đến hàng nghìn để ghi báo cáo. Bác ghi số gạo là
\choice
{$53\,000\text{ kg}$}
{\True $52\,000\text{ kg}$}
{$50\,000\text{ kg}$}
{$52\,500\text{ kg}$}
\loigiai{
Chữ số hàng trăm là 4 ($4 < 5$) nên làm tròn thành $52\,000\text{ kg}$.
\textbf{Đáp án của câu hỏi là B}
}
\end{ex}

\begin{ex} %[3D1V62]
Cho dãy số: $10\,000; 20\,000; 30\,000; \dots$ Số hạng thứ 7 của dãy số là
\choice
{$60\,000$}
{$80\,000$}
{\True $70\,000$}
{$100\,000$}
\loigiai{
Dãy số tròn chục nghìn. Số hạng thứ 7 là $70\,000$.
\textbf{Đáp án của câu hỏi là C}
}
\end{ex}

\begin{ex} %[3D1V62]
Số gồm 9 chục nghìn và 9 đơn vị được viết là
\choice
{\True $90\,009$}
{$90\,090$}
{$99\,000$}
{$90\,900$}
\loigiai{
Hàng chục nghìn là 9, đơn vị là 9, các hàng khác là 0: $90\,009$.
\textbf{Đáp án của câu hỏi là A}
}
\end{ex}

\textbf{Phần 2. Trắc nghiệm đúng sai}
\setcounterpageref{ex}{0}

\begin{ex} %[3D1V62]
Xét tính đúng sai của các nhận định ôn tập số đến 100 000:
\choiceTF
{\True Số liền sau của 99 999 là 100 000.}
{Chữ số hàng chục nghìn của số 45 678 là chữ số 5.}
{\True Khi làm tròn 84 500 đến hàng nghìn ta được 85 000.}
{Số 100 000 có 5 chữ số.}
\loigiai{
\begin{itemize}
    \item Ý a) Đúng: Đúng vì $99\,999 + 1 = 100\,000$.
    \item Ý b) Sai: Sai vì chữ số hàng chục nghìn là 4.
    \item Ý c) Đúng: Vì hàng trăm là 5.
    \item Ý d) Sai: Sai vì số 100 000 có 6 chữ số.
\end{itemize}
\begin{itemize}
    \item a) \textbf{Đúng}.
    \item b) \textbf{Sai}.
    \item c) \textbf{Đúng}.
    \item d) \textbf{Sai}.
\end{itemize}
}
\end{ex}

\begin{ex} %[3D1V62]
Cho ba số: $X = 54\,200; Y = 54\,800; Z = 50\,480$:
\choiceTF
{\True Số bé nhất trong ba số là $Z = 50\,480$.}
{\True Số lớn nhất trong ba số là $Y = 54\,800$.}
{Làm tròn số $X$ đến hàng nghìn ta được 55 000.}
{\True Làm tròn số $Y$ đến hàng nghìn ta được 55 000.}
\loigiai{
\begin{itemize}
    \item Ý a) Đúng: Hàng nghìn bé nhất là 0.
    \item Ý b) Đúng: Hàng trăm lớn nhất là 8.
    \item Ý c) Sai: Sai vì hàng trăm của $X$ là 2 nên làm tròn thành 54 000.
    \item Ý d) Đúng: Hàng trăm là 8 nên làm tròn lên 55 000.
\end{itemize}
\begin{itemize}
    \item a) \textbf{Đúng}.
    \item b) \textbf{Đúng}.
    \item c) \textbf{Sai}.
    \item d) \textbf{Đúng}.
\end{itemize}
}
\end{ex}

\textbf{Phần 3. Trả lời ngắn}
\setcounterpageref{ex}{0}

\begin{ex} %[3D1V62]
Viết số bé nhất có năm chữ số khác nhau.
\shortans{10234}
\loigiai{
\textbf{Đáp số:} $10234$\\
Số đó là 10 234.
}
\end{ex}

\begin{ex} %[3D1V62]
Làm tròn số 86 750 đến hàng nghìn.
\shortans{87000}
\loigiai{
\textbf{Đáp số:} $87000$\\
Hàng trăm là 7 nên làm tròn thành 87 000.
}
\end{ex}

\begin{ex} %[3D1V62]
Làm tròn số 64 890 đến hàng chục nghìn.
\shortans{60000}
\loigiai{
\textbf{Đáp số:} $60000$\\
Hàng nghìn là 4 nên làm tròn thành 60 000.
}
\end{ex}

\begin{ex} %[3D1V62]
Tìm số liền sau của số 99 999.
\shortans{100000}
\loigiai{
\textbf{Đáp số:} $100000$\\
$99\,999 + 1 = 100\,000$.
}
\end{ex}

\begin{ex} %[3D1V62]
Tìm số tự nhiên lớn nhất có 5 chữ số khi làm tròn đến hàng chục nghìn được 70 000.
\shortans{74999}
\loigiai{
\textbf{Đáp số:} $74999$\\
Số lớn nhất là 74 999.
}
\end{ex}

\begin{ex} %[3D1V62]
Tính hiệu giữa số lớn nhất có 5 chữ số và số bé nhất có 5 chữ số.
\shortans{89999}
\loigiai{
\textbf{Đáp số:} $89999$\\
$99\,999 - 10\,000 = 89\,999$.
}
\end{ex}

\newpage

% =========================================================================
% CHỦ ĐỀ 12 - CỘNG, TRỪ TRONG PHẠM VI 100 000
% =========================================================================
\section{Cộng, trừ trong phạm vi 100 000}

\TaoPhieuKiemTra{Chủ đề 12 -- Cộng, trừ trong phạm vi 100 000}{63}{Phép cộng trong phạm vi 100 000}{https://loptoanthaylinh.vn/thi-online/toan-3/bai-63}
\textbf{Phần 1. Trắc nghiệm lựa chọn}
\setcounterpageref{ex}{0}

\begin{ex} %[3D2V63]
Tính giá trị của phép tính: $45\,280 + 32\,415$.
\choice
{$77\,595$}
{$77\,795$}
{\True $77\,695$}
{$76\,695$}
\loigiai{
$45\,280 + 32\,415 = 77\,695$.
\textbf{Đáp án của câu hỏi là C}
}
\end{ex}

\begin{ex} %[3D2V63]
Một nhà máy tháng thứ nhất sản xuất được $28\,500$ sản phẩm, tháng thứ hai sản xuất được $31\,200$ sản phẩm. Cả hai tháng nhà máy sản xuất được bao nhiêu sản phẩm?
\choice
{\True $59\,700$ sản phẩm}
{$59\,500$ sản phẩm}
{$60\,000$ sản phẩm}
{$58\,700$ sản phẩm}
\loigiai{
$28\,500 + 31\,200 = 59\,700$ sản phẩm.
\textbf{Đáp án của câu hỏi là A}
}
\end{ex}

\begin{ex} %[3D2V63]
Tìm số tự nhiên $x$, biết: $x - 24\,500 = 35\,200$.
\choice
{$58\,700$}
{$60\,700$}
{$59\,000$}
{\True $59\,700$}
\loigiai{
$x = 35\,200 + 24\,500 = 59\,700$.
\textbf{Đáp án của câu hỏi là D}
}
\end{ex}

\begin{ex} %[3D2V63]
Tính nhanh giá trị biểu thức: $25\,000 + 18\,400 + 15\,000 + 1600$.
\choice
{$58\,000$}
{\True $60\,000$}
{$62\,000$}
{$59\,000$}
\loigiai{
$(25\,000 + 15\,000) + (18\,400 + 1600) = 40\,000 + 20\,000 = 60\,000$.
\textbf{Đáp án của câu hỏi là B}
}
\end{ex}

\begin{ex} %[3D2V63]
Một huyện năm ngoái trồng được $34\,500$ cây xanh, năm nay trồng được nhiều hơn năm ngoái $5200$ cây xanh. Cả hai năm huyện đó trồng được bao nhiêu cây xanh?
\choice
{$39\,700$ cây}
{$73\,200$ cây}
{\True $74\,200$ cây}
{$75\,000$ cây}
\loigiai{
Năm nay trồng: $34\,500 + 5200 = 39\,700$ cây. Cả hai năm: $34\,500 + 39\,700 = 74\,200$ cây.
\textbf{Đáp án của câu hỏi là C}
}
\end{ex}

\begin{ex} %[3D2V63]
Tìm tổng của số lớn nhất có năm chữ số và số bé nhất có năm chữ số.
\choice
{\True $109\,999$}
{$110\,000$}
{$99\,999$}
{$100\,000$}
\loigiai{
Số lớn nhất: 99 999. Số bé nhất: 10 000. Tổng: $99\,999 + 10\,000 = 109\,999$.
\textbf{Đáp án của câu hỏi là A}
}
\end{ex}

\begin{ex} %[3D2V63]
Một kho hàng có $42\,000\text{ kg}$ gạo. Người ta nhập thêm đợt 1 là $18\,500\text{ kg}$ và đợt 2 là $15\,500\text{ kg}$. Trong kho có tất cả bao nhiêu ki-lô-gam gạo?
\choice
{$75\,000\text{ kg}$}
{$77\,000\text{ kg}$}
{$74\,000\text{ kg}$}
{\True $76\,000\text{ kg}$}
\loigiai{
Tổng gạo: $42\,000 + 18\,500 + 15\,500 = 76\,000\text{ (kg)}$.
\textbf{Đáp án của câu hỏi là D}
}
\end{ex}

\begin{ex} %[3D2V63]
Điền số thích hợp: $56\,340 + \dots = 80\,000$.
\choice
{$24\,660$}
{\True $23\,660$}
{$23\,760$}
{$22\,660$}
\loigiai{
$80\,000 - 56\,340 = 23\,660$.
\textbf{Đáp án của câu hỏi là B}
}
\end{ex}

\begin{ex} %[3D2V63]
Một gia đình mua một chiếc tivi giá $12\,500\,000$ đồng và một máy giặt giá $8\,200\,000$ đồng. Tổng số tiền gia đình đó phải trả là
\choice
{$20\,500\,000$ đồng}
{$21\,000\,000$ đồng}
{\True $20\,700\,000$ đồng}
{$19\,700\,000$ đồng}
\loigiai{
$12\,500\,000 + 8\,200\,000 = 20\,700\,000$ đồng.
\textbf{Đáp án của câu hỏi là C}
}
\end{ex}

\begin{ex} %[3D2V63]
Nếu số hạng thứ nhất tăng thêm 1500 đơn vị và số hạng thứ hai tăng thêm 2500 đơn vị thì tổng sẽ tăng thêm bao nhiêu?
\choice
{\True $4000$ đơn vị}
{$3000$ đơn vị}
{$5000$ đơn vị}
{$3500$ đơn vị}
\loigiai{
Tổng tăng: $1500 + 2500 = 4000$ đơn vị.
\textbf{Đáp án của câu hỏi là A}
}
\end{ex}

\begin{ex} %[3D2V63]
Tính giá trị: $48\,500 + 12\,500 + 20\,000$.
\choice
{$80\,000$}
{$82\,000$}
{$79\,000$}
{\True $81\,000$}
\loigiai{
$48\,500 + 12\,500 = 61\,000$; $61\,000 + 20\,000 = 81\,000$.
\textbf{Đáp án của câu hỏi là D}
}
\end{ex}

\begin{ex} %[3D2V63]
Một công ty thu nhập quý Một là $45\,000\,000$ đồng, quý Hai thu nhập nhiều hơn quý Một là $8\,000\,000$ đồng. Thu nhập cả hai quý của công ty là
\choice
{$53\,000\,000$ đồng}
{\True $98\,000\,000$ đồng}
{$90\,000\,000$ đồng}
{$100\,000\,000$ đồng}
\loigiai{
Quý Hai: $45 + 8 = 53$ triệu. Cả hai quý: $45 + 53 = 98$ triệu đồng.
\textbf{Đáp án của câu hỏi là B}
}
\end{ex}

\textbf{Phần 2. Trắc nghiệm đúng sai}
\setcounterpageref{ex}{0}

\begin{ex} %[3D2V63]
Xét tính đúng sai của phép cộng các số trong phạm vi 100 000:
\choiceTF
{\True Cộng lần lượt từ phải sang trái (từ hàng đơn vị).}
{$45\,000 + 55\,000 = 90\,000$.}
{\True Khi cộng số có 5 chữ số, nếu có nhớ thì ta thêm vào hàng liền kề bên trái.}
{$62\,300 + 14\,700 = 76\,000$.}
\loigiai{
\begin{itemize}
    \item Ý a) Đúng: Quy tắc chuẩn.
    \item Ý b) Sai: Sai vì $45\,000 + 55\,000 = 100\,000$.
    \item Ý c) Đúng: Quy tắc cộng có nhớ.
    \item Ý d) Sai: Sai vì $62\,300 + 14\,700 = 77\,000$.
\end{itemize}
\begin{itemize}
    \item a) \textbf{Đúng}.
    \item b) \textbf{Sai}.
    \item c) \textbf{Đúng}.
    \item d) \textbf{Sai}.
\end{itemize}
}
\end{ex}

\begin{ex} %[3D2V63]
Một nông trường có hai trại nuôi bò: trại 1 có 18 500 con, trại 2 có 21 500 con:
\choiceTF
{\True Cả hai trại có tất cả 40 000 con bò.}
{\True Trại 2 nhiều hơn trại 1 là 3000 con bò.}
{Tổng số bò làm tròn đến hàng chục nghìn là 50 000 con.}
{\True Nếu nhập thêm 5000 con bò nữa thì toàn nông trường có 45 000 con bò.}
\loigiai{
\begin{itemize}
    \item Ý a) Đúng: $18\,500 + 21\,500 = 40\,000$ con.
    \item Ý b) Đúng: $21\,500 - 18\,500 = 3000$ con.
    \item Ý c) Sai: Sai vì 40 000 là số tròn chục nghìn rồi.
    \item Ý d) Đúng: $40\,000 + 5000 = 45\,000$ con.
\end{itemize}
\begin{itemize}
    \item a) \textbf{Đúng}.
    \item b) \textbf{Đúng}.
    \item c) \textbf{Sai}.
    \item d) \textbf{Đúng}.
\end{itemize}
}
\end{ex}

\textbf{Phần 3. Trả lời ngắn}
\setcounterpageref{ex}{0}

\begin{ex} %[3D2V63]
Tính kết quả phép cộng: $54\,628 + 27\,345$.
\shortans{81973}
\loigiai{
\textbf{Đáp số:} $81973$\\
$54\,628 + 27\,345 = 81\,973$.
}
\end{ex}

\begin{ex} %[3D2V63]
Tìm $x$, biết: $x - 34\,200 = 45\,800$.
\shortans{80000}
\loigiai{
\textbf{Đáp số:} $80000$\\
$x = 45\,800 + 34\,200 = 80\,000$.
}
\end{ex}

\begin{ex} %[3D2V63]
Tính nhanh giá trị: $36\,000 + 28\,500 + 14\,000 + 11\,500$.
\shortans{90000}
\loigiai{
\textbf{Đáp số:} $90000$\\
$(36\,000 + 14\,000) + (28\,500 + 11\,500) = 50\,000 + 40\,000 = 90\,000$.
}
\end{ex}

\begin{ex} %[3D2V63]
Một kho hàng buổi sáng xuất $15\,400\text{ kg}$ đường, buổi chiều xuất $18\,600\text{ kg}$ đường. Hỏi cả ngày kho xuất bao nhiêu ki-lô-gam đường?
\shortans{34000}
\loigiai{
\textbf{Đáp số:} $34000$\\
$15\,400 + 18\,600 = 34\,000\text{ (kg)}$.
}
\end{ex}

\begin{ex} %[3D2V63]
Tính tổng của số tròn chục nghìn lớn nhất có 5 chữ số và số tròn nghìn bé nhất có 5 chữ số.
\shortans{100000}
\loigiai{
\textbf{Đáp số:} $100000$\\
$90\,000 + 10\,000 = 100\,000$.
}
\end{ex}

\begin{ex} %[3D2V63]
Một nhà máy dệt được $42\,500\text{ m}$ vải hoa và $38\,500\text{ m}$ vải trắng. Hỏi nhà máy dệt được tất cả bao nhiêu mét vải?
\shortans{81000}
\loigiai{
\textbf{Đáp số:} $81000$\\
$42\,500 + 38\,500 = 81\,000\text{ (m)}$.
}
\end{ex}

\newpage

\TaoPhieuKiemTra{Chủ đề 12 -- Cộng, trừ trong phạm vi 100 000}{64}{Phép trừ trong phạm vi 100 000}{https://loptoanthaylinh.vn/thi-online/toan-3/bai-64}
\textbf{Phần 1. Trắc nghiệm lựa chọn}
\setcounterpageref{ex}{0}

\begin{ex} %[3D2V64]
Tính giá trị của phép tính: $84\,650 - 32\,480$.
\choice
{$52\,270$}
{\True $52\,170$}
{$52\,070$}
{$51\,170$}
\loigiai{
$84\,650 - 32\,480 = 52\,170$.
\textbf{Đáp án của câu hỏi là B}
}
\end{ex}

\begin{ex} %[3D2V64]
Một kho gạo có $95\,000\text{ kg}$ gạo. Người ta đã xuất đi $48\,500\text{ kg}$. Trong kho còn lại bao nhiêu ki-lô-gam gạo?
\choice
{$47\,500\text{ kg}$}
{$46\,000\text{ kg}$}
{$56\,500\text{ kg}$}
{\True $46\,500\text{ kg}$}
\loigiai{
$95\,000 - 48\,500 = 46\,500\text{ (kg)}$.
\textbf{Đáp án của câu hỏi là D}
}
\end{ex}

\begin{ex} %[3D2V64]
Tìm số tự nhiên $y$, biết: $80\,000 - y = 35\,600$.
\choice
{\True $44\,400$}
{$45\,400$}
{$44\,600$}
{$54\,400$}
\loigiai{
$y = 80\,000 - 35\,600 = 44\,400$.
\textbf{Đáp án của câu hỏi là A}
}
\end{ex}

\begin{ex} %[3D2V64]
Một đoạn đường sắt dài $85\,400\text{ m}$. Người ta đã nâng cấp được $46\,200\text{ m}$. Đoạn đường sắt còn lại chưa nâng cấp dài bao nhiêu mét?
\choice
{$38\,200\text{ m}$}
{$40\,200\text{ m}$}
{\True $39\,200\text{ m}$}
{$39\,000\text{ m}$}
\loigiai{
Đoạn còn lại: $85\,400 - 46\,200 = 39\,200\text{ (m)}$.
\textbf{Đáp án của câu hỏi là C}
}
\end{ex}

\begin{ex} %[3D2V64]
Tính nhanh giá trị: $74\,500 - 25\,000 - 14\,500$.
\choice
{$40\,000$}
{\True $35\,000$}
{$30\,000$}
{$45\,000$}
\loigiai{
$(74\,500 - 14\,500) - 25\,000 = 60\,000 - 25\,000 = 35\,000$.
\textbf{Đáp án của câu hỏi là B}
}
\end{ex}

\begin{ex} %[3D2V64]
Tìm hiệu của số tròn chục nghìn lớn nhất có 5 chữ số và số tròn chục nghìn bé nhất có 5 chữ số.
\choice
{$90\,000$}
{$70\,000$}
{$100\,000$}
{\True $80\,000$}
\loigiai{
$90\,000 - 10\,000 = 80\,000$.
\textbf{Đáp án của câu hỏi là D}
}
\end{ex}

\begin{ex} %[3D2V64]
Một trang trại có $65\,000$ con vịt. Sau khi xuất bán một số con thì trang trại còn lại $28\,500$ con. Hỏi trang trại đã bán bao nhiêu con vịt?
\choice
{\True $36\,500$ con}
{$37\,500$ con}
{$35\,500$ con}
{$46\,500$ con}
\loigiai{
Số vịt đã bán: $65\,000 - 28\,500 = 36\,500$ con.
\textbf{Đáp án của câu hỏi là A}
}
\end{ex}

\begin{ex} %[3D2V64]
Nếu bớt ở số bị trừ đi 1200 đơn vị và giữ nguyên số trừ thì hiệu sẽ
\choice
{Tăng thêm 1200 đơn vị}
{Không đổi}
{\True Giảm đi 1200 đơn vị}
{Giảm đi 2400 đơn vị}
\loigiai{
Số bị trừ giảm bao nhiêu thì hiệu giảm bấy nhiêu.
\textbf{Đáp án của câu hỏi là C}
}
\end{ex}

\begin{ex} %[3D2V64]
Một xã có $78\,500$ nhân khẩu. Sau một năm có $2400$ người chuyển đi nơi khác. Số dân của xã còn lại là
\choice
{$75\,100$ người}
{\True $76\,100$ người}
{$76\,000$ người}
{$77\,100$ người}
\loigiai{
$78\,500 - 2400 = 76\,100$ người.
\textbf{Đáp án của câu hỏi là B}
}
\end{ex}

\begin{ex} %[3D2V64]
Tính giá trị của biểu thức: $100\,000 - (45\,000 + 25\,000)$.
\choice
{$20\,000$}
{$40\,000$}
{$25\,000$}
{\True $30\,000$}
\loigiai{
Trong ngoặc: $45\,000 + 25\,000 = 70\,000$. Phép trừ: $100\,000 - 70\,000 = 30\,000$.
\textbf{Đáp án của câu hỏi là D}
}
\end{ex}

\begin{ex} %[3D2V64]
Một đại lý xăng dầu nhập về $90\,000\ell$ dầu. Tuần 1 bán $32\,000\ell$, tuần 2 bán $35\,000\ell$. Sau hai tuần đại lý còn lại bao nhiêu lít dầu?
\choice
{\True $23\,000\ell$}
{$22\,000\ell$}
{$25\,000\ell$}
{$24\,000\ell$}
\loigiai{
Tổng bán: $32\,000 + 35\,000 = 67\,000\ell$. Còn lại: $90\,000 - 67\,000 = 23\,000\text{ (\ell)}$.
\textbf{Đáp án của câu hỏi là A}
}
\end{ex}

\begin{ex} %[3D2V64]
Tìm số tự nhiên biết rằng lấy số lớn nhất có năm chữ số trừ đi số đó thì được 45 678.
\choice
{$54\,312$}
{$53\,321$}
{\True $54\,321$}
{$54\,231$}
\loigiai{
Số cần tìm: $99\,999 - 45\,678 = 54\,321$.
\textbf{Đáp án của câu hỏi là C}
}
\end{ex}

\textbf{Phần 2. Trắc nghiệm đúng sai}
\setcounterpageref{ex}{0}

\begin{ex} %[3D2V64]
Xét tính đúng sai của các phép trừ trong phạm vi 100 000:
\choiceTF
{\True Thực hiện phép trừ lần lượt từ phải sang trái.}
{$100\,000 - 1 = 9999$.}
{\True Khi cùng bớt ở cả số bị trừ và số trừ một số như nhau thì hiệu không thay đổi.}
{$85\,000 - 45\,000 = 30\,000$.}
\loigiai{
\begin{itemize}
    \item Ý a) Đúng: Quy tắc trừ.
    \item Ý b) Sai: Sai vì $100\,000 - 1 = 99\,999$.
    \item Ý c) Đúng: Tính chất của hiệu.
    \item Ý d) Sai: Sai vì $85\,000 - 45\,000 = 40\,000$.
\end{itemize}
\begin{itemize}
    \item a) \textbf{Đúng}.
    \item b) \textbf{Sai}.
    \item c) \textbf{Đúng}.
    \item d) \textbf{Sai}.
\end{itemize}
}
\end{ex}

\begin{ex} %[3D2V64]
Một bể nước chứa $80\,000\ell$ nước, người ta xả nước hai lần:
\choiceTF
{\True Lần 1 xả $25\,000\ell$ thì trong bể còn lại $55\,000\ell$ nước.}
{\True Lần 2 xả $35\,000\ell$ thì trong bể còn lại $20\,000\ell$ nước.}
{Tổng lượng nước xả ra sau hai lần là $50\,000\ell$.}
{\True Lượng nước còn lại bằng một phần tư lượng nước ban đầu.}
\loigiai{
\begin{itemize}
    \item Ý a) Đúng: $80\,000 - 25\,000 = 55\,000\ell$.
    \item Ý b) Đúng: $55\,000 - 35\,000 = 20\,000\ell$.
    \item Ý c) Sai: Sai vì $25\,000 + 35\,000 = 60\,000\ell$.
    \item Ý d) Đúng: $20\,000 : 80\,000 = \dfrac{1}{4}$.
\end{itemize}
\begin{itemize}
    \item a) \textbf{Đúng}.
    \item b) \textbf{Đúng}.
    \item c) \textbf{Sai}.
    \item d) \textbf{Đúng}.
\end{itemize}
}
\end{ex}

\textbf{Phần 3. Trả lời ngắn}
\setcounterpageref{ex}{0}

\begin{ex} %[3D2V64]
Tính kết quả: $94\,560 - 47\,280$.
\shortans{47280}
\loigiai{
\textbf{Đáp số:} $47280$\\
$94\,560 - 47\,280 = 47\,280$.
}
\end{ex}

\begin{ex} %[3D2V64]
Tìm số tự nhiên $y$, biết: $75\,000 - y = 28\,500$.
\shortans{46500}
\loigiai{
\textbf{Đáp số:} $46500$\\
$y = 75\,000 - 28\,500 = 46\,500$.
}
\end{ex}

\begin{ex} %[3D2V64]
Một nông trường có $64\,000\text{ kg}$ cà phê, đã bán đi $38\,500\text{ kg}$. Hỏi nông trường còn lại bao nhiêu ki-lô-gam cà phê?
\shortans{25500}
\loigiai{
\textbf{Đáp số:} $25500$\\
$64\,000 - 38\,500 = 25\,500\text{ (kg)}$.
}
\end{ex}

\begin{ex} %[3D2V64]
Tính nhanh: $86\,400 - 30\,000 - 26\,400$.
\shortans{30000}
\loigiai{
\textbf{Đáp số:} $30000$\\
$(86\,400 - 26\,400) - 30\,000 = 60\,000 - 30\,000 = 30\,000$.
}
\end{ex}

\begin{ex} %[3D2V64]
Tính hiệu: $100\,000 - 45\,800$.
\shortans{54200}
\loigiai{
\textbf{Đáp số:} $54200$\\
$100\,000 - 45\,800 = 54\,200$.
}
\end{ex}

\begin{ex} %[3D2V64]
Một đoạn đê dài $50\,000\text{ m}$, người ta đã kè đá được $32\,500\text{ m}$. Hỏi còn bao nhiêu mét đê chưa kè đá?
\shortans{17500}
\loigiai{
\textbf{Đáp số:} $17500$\\
$50\,000 - 32\,500 = 17\,500\text{ (m)}$.
}
\end{ex}

\newpage

\TaoPhieuKiemTra{Chủ đề 12 -- Cộng, trừ trong phạm vi 100 000}{65}{Luyện tập chung cộng trừ 100 000}{https://loptoanthaylinh.vn/thi-online/toan-3/bai-65}
\textbf{Phần 1. Trắc nghiệm lựa chọn}
\setcounterpageref{ex}{0}

\begin{ex} %[3D2V65]
Tính giá trị của biểu thức: $A = 35\,000 + 45\,000 - 28\,000$.
\choice
{\True $52\,000$}
{$50\,000$}
{$54\,000$}
{$48\,000$}
\loigiai{
$35\,000 + 45\,000 = 80\,000$; $80\,000 - 28\,000 = 52\,000$.
\textbf{Đáp án của câu hỏi là A}
}
\end{ex}

\begin{ex} %[3D2V65]
Tìm $x$, biết: $x + 24\,000 = 70\,000 - 16\,000$.
\choice
{$28\,000$}
{$32\,000$}
{\True $30\,000$}
{$35\,000$}
\loigiai{
$x + 24\,000 = 54\,000 \Rightarrow x = 54\,000 - 24\,000 = 30\,000$.
\textbf{Đáp án của câu hỏi là C}
}
\end{ex}

\begin{ex} %[3D2V65]
Một kho thóc có $80\,000\text{ kg}$ thóc. Ngày đầu xuất $25\,000\text{ kg}$, ngày thứ hai xuất nhiều hơn ngày đầu $5000\text{ kg}$. Trong kho còn lại bao nhiêu ki-lô-gam thóc?
\choice
{$30\,000\text{ kg}$}
{\True $25\,000\text{ kg}$}
{$20\,000\text{ kg}$}
{$35\,000\text{ kg}$}
\loigiai{
Ngày hai xuất: $25\,000 + 5000 = 30\,000\text{ kg}$. Tổng xuất: $25\,000 + 30\,000 = 55\,000\text{ kg}$. Còn lại: $80\,000 - 55\,000 = 25\,000\text{ (kg)}$.
\textbf{Đáp án của câu hỏi là B}
}
\end{ex}

\begin{ex} %[3D2V65]
Tính nhanh giá trị: $42\,500 + 17\,800 + 7500 + 2200$.
\choice
{$68\,000$}
{$72\,000$}
{$69\,000$}
{\True $70\,000$}
\loigiai{
$(42\,500 + 7500) + (17\,800 + 2200) = 50\,000 + 20\,000 = 70\,000$.
\textbf{Đáp án của câu hỏi là D}
}
\end{ex}

\begin{ex} %[3D2V65]
Tổng của hai số là 75 000, hiệu của hai số là 15 000. Số lớn là
\choice
{\True $45\,000$}
{$30\,000$}
{$40\,000$}
{$50\,000$}
\loigiai{
Số lớn: $(75\,000 + 15\,000) : 2 = 45\,000$.
\textbf{Đáp án của câu hỏi là A}
}
\end{ex}

\begin{ex} %[3D2V65]
Một mảnh đất hình chữ nhật có nửa chu vi là $450\text{ m}$, chiều dài hơn chiều rộng $50\text{ m}$. Chu vi mảnh đất đó là
\choice
{$500\text{ m}$}
{$800\text{ m}$}
{\True $900\text{ m}$}
{$1000\text{ m}$}
\loigiai{
Chu vi bằng nửa chu vi nhân 2: $450 \times 2 = 900\text{ (m)}$.
\textbf{Đáp án của câu hỏi là C}
}
\end{ex}

\begin{ex} %[3D2V65]
Một quỹ từ thiện nhận được ba đợt ủng hộ: đợt một $35\,000\,000$ đồng, đợt hai $28\,000\,000$ đồng, đợt ba $37\,000\,000$ đồng. Tổng số tiền quỹ nhận được là
\choice
{$90\,000\,000$ đồng}
{\True $100\,000\,000$ đồng}
{$95\,000\,000$ đồng}
{$110\,000\,000$ đồng}
\loigiai{
$35 + 28 + 37 = 100$ triệu đồng.
\textbf{Đáp án của câu hỏi là B}
}
\end{ex}

\begin{ex} %[3D2V65]
Tìm số tự nhiên có năm chữ số, biết rằng nếu thêm 1 đơn vị vào số đó thì được số tròn chục nghìn bé nhất có sáu chữ số ($100\,000$). Số đó là
\choice
{$99\,990$}
{$99\,000$}
{$90\,000$}
{\True $99\,999$}
\loigiai{
$100\,000 - 1 = 99\,999$.
\textbf{Đáp án của câu hỏi là D}
}
\end{ex}

\begin{ex} %[3D2V65]
Tính giá trị của biểu thức: $85\,000 - (30\,000 + 25\,000) + 12\,000$.
\choice
{\True $42\,000$}
{$40\,000$}
{$45\,000$}
{$38\,000$}
\loigiai{
$85\,000 - 55\,000 + 12\,000 = 30\,000 + 12\,000 = 42\,000$.
\textbf{Đáp án của câu hỏi là A}
}
\end{ex}

\begin{ex} %[3D2V65]
Một nhà máy có 65 000 sản phẩm. Người ta chuyển đi 3 chuyến xe, mỗi chuyến chở 15 000 sản phẩm. Trong nhà máy còn lại bao nhiêu sản phẩm?
\choice
{$15\,000$ sản phẩm}
{$25\,000$ sản phẩm}
{\True $20\,000$ sản phẩm}
{$30\,000$ sản phẩm}
\loigiai{
Đã chuyển: $15\,000 \times 3 = 45\,000$. Còn lại: $65\,000 - 45\,000 = 20\,000$ sản phẩm.
\textbf{Đáp án của câu hỏi là C}
}
\end{ex}

\begin{ex} %[3D2V65]
Số bé là 24 000, số lớn gấp 3 lần số bé. Hiệu giữa hai số là
\choice
{$72\,000$}
{\True $48\,000$}
{$50\,000$}
{$46\,000$}
\loigiai{
Số lớn: $24\,000 \times 3 = 72\,000$. Hiệu: $72\,000 - 24\,000 = 48\,000$.
\textbf{Đáp án của câu hỏi là B}
}
\end{ex}

\begin{ex} %[3D2V65]
Một cuộn dây cáp dài $100\,000\text{ m}$. Lần đầu cắt $35\,000\text{ m}$, lần sau cắt $42\,000\text{ m}$. Cuộn dây cáp còn lại dài bao nhiêu mét?
\choice
{$25\,000\text{ m}$}
{$22\,000\text{ m}$}
{$33\,000\text{ m}$}
{\True $23\,000\text{ m}$}
\loigiai{
Đã cắt: $35\,000 + 42\,000 = 77\,000\text{ m}$. Còn lại: $100\,000 - 77\,000 = 23\,000\text{ (m)}$.
\textbf{Đáp án của câu hỏi là D}
}
\end{ex}

\textbf{Phần 2. Trắc nghiệm đúng sai}
\setcounterpageref{ex}{0}

\begin{ex} %[3D2V65]
Xét tính đúng sai của các phép tính tổng hợp phạm vi 100 000:
\choiceTF
{\True Biểu thức $50\,000 + 30\,000 - 20\,000$ có giá trị bằng $60\,000$.}
{$100\,000 - 40\,000 = 50\,000$.}
{\True Tổng của hai số tròn nghìn luôn là một số tròn nghìn.}
{Hiệu giữa số lớn nhất có 5 chữ số và số bé nhất có 5 chữ số là 90 000.}
\loigiai{
\begin{itemize}
    \item Ý a) Đúng: $80\,000 - 20\,000 = 60\,000$.
    \item Ý b) Sai: Sai vì $100\,000 - 40\,000 = 60\,000$.
    \item Ý c) Đúng: Tính chất của phép cộng tròn nghìn.
    \item Ý d) Sai: Sai vì $99\,999 - 10\,000 = 89\,999$.
\end{itemize}
\begin{itemize}
    \item a) \textbf{Đúng}.
    \item b) \textbf{Sai}.
    \item c) \textbf{Đúng}.
    \item d) \textbf{Sai}.
\end{itemize}
}
\end{ex}

\begin{ex} %[3D2V65]
Một thành phố có hai quận: quận A có 45 000 dân, quận B có 55 000 dân:
\choiceTF
{\True Tổng số dân của cả hai quận là 100 000 người.}
{\True Quận B nhiều hơn quận A là 10 000 người.}
{Dân số quận B gấp đôi dân số quận A.}
{\True Cả hai số dân đều là các số tròn nghìn và tròn chục nghìn (đối với quận B là tròn chục nghìn, quận A là tròn nghìn).}
\loigiai{
\begin{itemize}
    \item Ý a) Đúng: $45\,000 + 55\,000 = 100\,000$ người.
    \item Ý b) Đúng: $55\,000 - 45\,000 = 10\,000$ người.
    \item Ý c) Sai: Sai vì $45\,000 \times 2 = 90\,000 \neq 55\,000$.
    \item Ý d) Đúng: Đặc điểm của các số.
\end{itemize}
\begin{itemize}
    \item a) \textbf{Đúng}.
    \item b) \textbf{Đúng}.
    \item c) \textbf{Sai}.
    \item d) \textbf{Đúng}.
\end{itemize}
}
\end{ex}

\textbf{Phần 3. Trả lời ngắn}
\setcounterpageref{ex}{0}

\begin{ex} %[3D2V65]
Tính giá trị biểu thức: $48\,000 + 32\,000 - 25\,000$.
\shortans{55000}
\loigiai{
\textbf{Đáp số:} $55000$\\
$80\,000 - 25\,000 = 55\,000$.
}
\end{ex}

\begin{ex} %[3D2V65]
Tìm số tự nhiên $x$, biết: $x - 18\,500 = 51\,500$.
\shortans{70000}
\loigiai{
\textbf{Đáp số:} $70000$\\
$x = 51\,500 + 18\,500 = 70\,000$.
}
\end{ex}

\begin{ex} %[3D2V65]
Tính nhanh: $27\,500 + 38\,000 + 12\,500 + 12\,000$.
\shortans{90000}
\loigiai{
\textbf{Đáp số:} $90000$\\
$(27\,500 + 12\,500) + (38\,000 + 12\,000) = 40\,000 + 50\,000 = 90\,000$.
}
\end{ex}

\begin{ex} %[3D2V65]
Một trang trại có 85 000 con gà, đã bán đi 56 000 con. Hỏi trang trại còn lại bao nhiêu con gà?
\shortans{29000}
\loigiai{
\textbf{Đáp số:} $29000$\\
$85\,000 - 56\,000 = 29\,000$ con.
}
\end{ex}

\begin{ex} %[3D2V65]
Tìm hiệu giữa 100 000 và 67 800.
\shortans{32200}
\loigiai{
\textbf{Đáp số:} $32200$\\
$100\,000 - 67\,800 = 32\,200$.
}
\end{ex}

\begin{ex} %[3D2V65]
Tổng hai số là 90 000, số bé là 35 000. Tìm số lớn.
\shortans{55000}
\loigiai{
\textbf{Đáp số:} $55000$\\
Số lớn là: $90\,000 - 35\,000 = 55\,000$.
}
\end{ex}

\newpage

% =========================================================================
% CHỦ ĐỀ 13 - XEM ĐỒNG HỒ. THÁNG – NĂM. TIỀN VIỆT NAM
% =========================================================================
\section{Xem đồng hồ. tháng -- năm. tiền việt nam}

\TaoPhieuKiemTra{Chủ đề 13 -- Xem đồng hồ. tháng -- năm. tiền việt nam}{66}{Xem đồng hồ}{https://loptoanthaylinh.vn/thi-online/toan-3/bai-66}
\textbf{Phần 1. Trắc nghiệm lựa chọn}
\setcounterpageref{ex}{0}

\begin{ex} %[3D3V66]
Những tháng nào trong năm có đúng 31 ngày?
\choice
{Tháng 1, 2, 3, 4, 5, 6, 7}
{Tháng 4, 6, 9, 11}
{Tháng 1, 3, 5, 7, 9, 11}
{\True Tháng 1, 3, 5, 7, 8, 10, 12}
\loigiai{
Theo quy tắc các khớp ngón tay: tháng 1, 3, 5, 7, 8, 10, 12 có đúng 31 ngày.
\textbf{Đáp án của câu hỏi là D}
}
\end{ex}

\begin{ex} %[3D3V66]
Nếu ngày 3 tháng 10 là ngày Chủ nhật thì ngày 24 tháng 10 cùng năm đó là ngày thứ mấy?
\choice
{Thứ Bảy}
{\True Chủ nhật}
{Thứ Hai}
{Thứ Sáu}
\loigiai{
Khoảng cách: $24 - 3 = 21$ ngày. Vì $21 : 7 = 3$ (chia hết cho 7) nên ngày 24 tháng 10 cũng là ngày Chủ nhật.
\textbf{Đáp án của câu hỏi là B}
}
\end{ex}

\begin{ex} %[3D3V66]
Đồng hồ chỉ kim ngắn ở giữa số 4 và số 5, kim dài chỉ số 8. Lúc đó là mấy giờ?
\choice
{$4$ giờ $8$ phút}
{$5$ giờ $40$ phút}
{\True $4$ giờ $40$ phút (hoặc $5$ giờ kém $20$ phút)}
{$4$ giờ $20$ phút}
\loigiai{
Kim dài chỉ số 8 là: $8 \times 5 = 40$ phút. Lúc đó là 4 giờ 40 phút hay 5 giờ kém 20 phút.
\textbf{Đáp án của câu hỏi là C}
}
\end{ex}

\begin{ex} %[3D3V66]
Tháng 2 của năm nhuận có bao nhiêu ngày?
\choice
{\True $29$ ngày}
{$28$ ngày}
{$30$ ngày}
{$31$ ngày}
\loigiai{
Năm nhuận có 366 ngày, riêng tháng 2 có 29 ngày.
\textbf{Đáp án của câu hỏi là A}
}
\end{ex}

\begin{ex} %[3D3V66]
Một bộ phim bắt đầu chiếu lúc $19$ giờ $30$ phút và kết thúc lúc $21$ giờ $15$ phút. Bộ phim đó kéo dài bao lâu?
\choice
{$1$ giờ $30$ phút}
{$2$ giờ}
{$1$ giờ $15$ phút}
{\True $1$ giờ $45$ phút}
\loigiai{
Từ 19 giờ 30 phút đến 20 giờ là 30 phút; từ 20 giờ đến 21 giờ 15 phút là 1 giờ 15 phút. Tổng: $30 + 75 = 105$ phút = 1 giờ 45 phút.
\textbf{Đáp án của câu hỏi là D}
}
\end{ex}

\begin{ex} %[3D3V66]
Nếu ngày 28 tháng 4 là thứ Năm thì ngày 2 tháng 5 cùng năm đó là ngày thứ mấy?
\choice
{Chủ nhật}
{\True Thứ Hai}
{Thứ Ba}
{Thứ Bảy}
\loigiai{
Tháng 4 có 30 ngày: ngày 28 (thứ Năm), 29 (thứ Sáu), 30 (thứ Bảy), ngày 1 tháng 5 (Chủ nhật), ngày 2 tháng 5 là thứ Hai.
\textbf{Đáp án của câu hỏi là B}
}
\end{ex}

\begin{ex} %[3D3V66]
Những tháng nào sau đây có đúng 30 ngày?
\choice
{Tháng 2, tháng 4, tháng 6, tháng 8}
{Tháng 1, tháng 4, tháng 7, tháng 10}
{\True Tháng 4, tháng 6, tháng 9, tháng 11}
{Tháng 3, tháng 5, tháng 9, tháng 11}
\loigiai{
Bốn tháng có đúng 30 ngày là: tháng 4, 6, 9, 11.
\textbf{Đáp án của câu hỏi là C}
}
\end{ex}

\begin{ex} %[3D3V66]
Một đồng hồ chạy chậm mỗi ngày 3 phút. Hỏi sau một tuần lễ, chiếc đồng hồ đó chạy chậm tất cả bao nhiêu phút?
\choice
{\True $21$ phút}
{$15$ phút}
{$18$ phút}
{$24$ phút}
\loigiai{
Một tuần lễ có 7 ngày: $3 \times 7 = 21$ phút.
\textbf{Đáp án của câu hỏi là A}
}
\end{ex}

\begin{ex} %[3D3V66]
Nam bắt đầu học bài lúc $20$ giờ $10$ phút và học xong lúc $21$ giờ $05$ phút. Thời gian Nam ngồi học là
\choice
{$45$ phút}
{$50$ phút}
{$1$ giờ}
{\True $55$ phút}
\loigiai{
Từ 20 giờ 10 phút đến 21 giờ là 50 phút; thêm 5 phút nữa là: $50 + 5 = 55$ phút.
\textbf{Đáp án của câu hỏi là D}
}
\end{ex}

\begin{ex} %[3D3V66]
Năm không nhuận có tất cả bao nhiêu ngày?
\choice
{$366$ ngày}
{\True $365$ ngày}
{$360$ ngày}
{$364$ ngày}
\loigiai{
Năm thường (không nhuận) có 365 ngày.
\textbf{Đáp án của câu hỏi là B}
}
\end{ex}

\begin{ex} %[3D3V66]
Lúc $15$ giờ $45$ phút còn được gọi là
\choice
{$3$ giờ $45$ phút sáng}
{$4$ giờ $15$ phút chiều}
{\True $4$ giờ kém $15$ phút chiều}
{$15$ giờ kém $15$ phút}
\loigiai{
$15$ giờ là 3 giờ chiều. $15$ giờ $45$ phút chính là 4 giờ kém 15 phút chiều.
\textbf{Đáp án của câu hỏi là C}
}
\end{ex}

\begin{ex} %[3D3V66]
Một trận bóng đá diễn ra trong hai hiệp, mỗi hiệp $45$ phút, nghỉ giữa hai hiệp $15$ phút. Tổng thời gian của trận bóng đá (không kể bù giờ) là
\choice
{\True $1$ giờ $45$ phút}
{$1$ giờ $30$ phút}
{$2$ giờ}
{$1$ giờ $15$ phút}
\loigiai{
Tổng thời gian: $45 \times 2 + 15 = 90 + 15 = 105$ phút = 1 giờ 45 phút.
\textbf{Đáp án của câu hỏi là A}
}
\end{ex}

\textbf{Phần 2. Trắc nghiệm đúng sai}
\setcounterpageref{ex}{0}

\begin{ex} %[3D3V66]
Xét tính đúng sai của các nhận định về thời gian, tháng và năm:
\choiceTF
{\True Một năm có 12 tháng.}
{Tất cả các tháng có 31 ngày đều là các tháng mang số lẻ.}
{\True Tháng 2 năm nhuận có 29 ngày, năm thường có 28 ngày.}
{Một ngày có 12 giờ.}
\loigiai{
\begin{itemize}
    \item Ý a) Đúng: Kiến thức chuẩn.
    \item Ý b) Sai: Sai vì tháng 8, 10, 12 là số chẵn nhưng vẫn có 31 ngày.
    \item Ý c) Đúng: Quy luật lịch dương.
    \item Ý d) Sai: Sai vì một ngày có 24 giờ.
\end{itemize}
\begin{itemize}
    \item a) \textbf{Đúng}.
    \item b) \textbf{Sai}.
    \item c) \textbf{Đúng}.
    \item d) \textbf{Sai}.
\end{itemize}
}
\end{ex}

\begin{ex} %[3D3V66]
Lịch tháng 8 của một năm có ngày 19 tháng 8 là ngày thứ Sáu:
\choiceTF
{\True Ngày 12 tháng 8 cùng năm đó cũng là ngày thứ Sáu.}
{\True Ngày 26 tháng 8 cùng năm đó cũng là ngày thứ Sáu.}
{Tháng 8 có đúng 30 ngày.}
{\True Ngày 2 tháng 9 cùng năm đó là ngày thứ Sáu.}
\loigiai{
\begin{itemize}
    \item Ý a) Đúng: $19 - 7 = 12$.
    \item Ý b) Đúng: $19 + 7 = 26$.
    \item Ý c) Sai: Sai vì tháng 8 có 31 ngày.
    \item Ý d) Đúng: Tháng 8 có 31 ngày: 26 (thứ Sáu) + 7 ngày = ngày 2 tháng 9 (thứ Sáu).
\end{itemize}
\begin{itemize}
    \item a) \textbf{Đúng}.
    \item b) \textbf{Đúng}.
    \item c) \textbf{Sai}.
    \item d) \textbf{Đúng}.
\end{itemize}
}
\end{ex}

\textbf{Phần 3. Trả lời ngắn}
\setcounterpageref{ex}{0}

\begin{ex} %[3D3V66]
Một năm thường (không nhuận) có bao nhiêu ngày?
\shortans{365}
\loigiai{
\textbf{Đáp số:} $365$\\
Năm thường có 365 ngày.
}
\end{ex}

\begin{ex} %[3D3V66]
Có bao nhiêu tháng trong năm có đúng 30 ngày?
\shortans{4}
\loigiai{
\textbf{Đáp số:} $4$\\
Gồm 4 tháng: 4, 6, 9, 11.
}
\end{ex}

\begin{ex} %[3D3V66]
Nếu ngày 5 tháng 5 là thứ Năm thì ngày 26 tháng 5 cùng năm là thứ mấy (viết số thứ: thứ 5 viết 5)?
\shortans{5}
\loigiai{
\textbf{Đáp số:} $5$\\
$26 - 5 = 21$ chia hết cho 7 nên vẫn là thứ Năm (5).
}
\end{ex}

\begin{ex} %[3D3V66]
Đổi $2$ giờ $15$ phút ra đơn vị phút.
\shortans{135}
\loigiai{
\textbf{Đáp số:} $135$\\
$2 \times 60 + 15 = 135$ phút.
}
\end{ex}

\begin{ex} %[3D3V66]
Xe ô tô khởi hành lúc 7 giờ 30 phút và đến nơi lúc 10 giờ 15 phút. Hỏi xe chạy hết bao nhiêu phút?
\shortans{165}
\loigiai{
\textbf{Đáp số:} $165$\\
Từ 7h30 đến 10h15 là 2 giờ 45 phút = $120 + 45 = 165$ phút.
}
\end{ex}

\begin{ex} %[3D3V66]
Tháng 12 có bao nhiêu ngày?
\shortans{31}
\loigiai{
\textbf{Đáp số:} $31$\\
Tháng 12 có 31 ngày.
}
\end{ex}

\newpage

\TaoPhieuKiemTra{Chủ đề 13 -- Xem đồng hồ. tháng -- năm. tiền việt nam}{67}{Tháng -- Năm}{https://loptoanthaylinh.vn/thi-online/toan-3/bai-67}
\textbf{Phần 1. Trắc nghiệm lựa chọn}
\setcounterpageref{ex}{0}

\begin{ex} %[3D3V67]
Bạn Mai rời nhà lúc 6 giờ 45 phút và đến trường lúc 7 giờ 10 phút. Thời gian Mai đi từ nhà đến trường là
\choice
{$15$ phút}
{$35$ phút}
{\True $25$ phút}
{$20$ phút}
\loigiai{
Từ 6h45 đến 7h00 là 15 phút, từ 7h00 đến 7h10 là 10 phút. Tổng thời gian: $15 + 10 = 25$ phút.
\textbf{Đáp án của câu hỏi là C}
}
\end{ex}

\begin{ex} %[3D3V67]
Nếu hôm nay là thứ Tư ngày 15 thì thứ Tư tuần sau là ngày nào?
\choice
{\True Ngày $22$}
{Ngày $21$}
{Ngày $23$}
{Ngày $20$}
\loigiai{
Thứ Tư tuần sau là: $15 + 7 = 22$.
\textbf{Đáp án của câu hỏi là A}
}
\end{ex}

\begin{ex} %[3D3V67]
Đồng hồ điện tử hiển thị $18:45$. Thời gian này tương ứng với
\choice
{$6$ giờ $45$ phút sáng}
{$8$ giờ kém $15$ phút tối}
{$18$ giờ $45$ phút sáng}
{\True $7$ giờ kém $15$ phút tối}
\loigiai{
$18$ giờ là 6 giờ chiều; $18:45$ là 6 giờ 45 phút chiều hay 7 giờ kém 15 phút tối.
\textbf{Đáp án của câu hỏi là D}
}
\end{ex}

\begin{ex} %[3D3V67]
Một buổi học bắt đầu từ 8 giờ sáng và kéo dài trong 3 tiết học, mỗi tiết 40 phút, giờ ra chơi giữa các tiết tổng cộng 20 phút. Buổi học kết thúc lúc
\choice
{$10$ giờ}
{\True $10$ giờ $20$ phút}
{$10$ giờ $10$ phút}
{$10$ giờ $40$ phút}
\loigiai{
Tổng thời gian học và chơi: $40 \times 3 + 20 = 140$ phút = 2 giờ 20 phút. Kết thúc: $8\text{h} + 2\text{h } 20\text{p} = 10$ giờ 20 phút.
\textbf{Đáp án của câu hỏi là B}
}
\end{ex}

\begin{ex} %[3D3V67]
Ngày Quốc tế Thiếu nhi 1 tháng 6 năm nay là thứ Tư. Hỏi ngày 8 tháng 6 cùng năm là thứ mấy?
\choice
{Thứ Ba}
{Thứ Năm}
{\True Thứ Tư}
{Chủ nhật}
\loigiai{
Cách nhau đúng 7 ngày ($8 - 1 = 7$) nên cũng là ngày thứ Tư.
\textbf{Đáp án của câu hỏi là C}
}
\end{ex}

\begin{ex} %[3D3V67]
Tàu hỏa chạy từ ga Hà Nội lúc 22 giờ đêm hôm trước và đến ga Vinh lúc 4 giờ 30 phút sáng hôm sau. Thời gian tàu chạy là
\choice
{\True $6$ giờ $30$ phút}
{$5$ giờ $30$ phút}
{$6$ giờ}
{$7$ giờ}
\loigiai{
Từ 22h đến 24h là 2 giờ; từ 0h đến 4h30 là 4 giờ 30 phút. Tổng thời gian: $2\text{h} + 4\text{h } 30\text{p} = 6$ giờ 30 phút.
\textbf{Đáp án của câu hỏi là A}
}
\end{ex}

\begin{ex} %[3D3V67]
Một năm có bao nhiêu quý, mỗi quý có mấy tháng?
\choice
{$3$ quý, mỗi quý có $4$ tháng}
{$2$ quý, mỗi quý có $6$ tháng}
{$4$ quý, mỗi quý có $4$ tháng}
{\True $4$ quý, mỗi quý có $3$ tháng}
\loigiai{
Một năm có 4 quý, mỗi quý có $12 : 4 = 3$ tháng.
\textbf{Đáp án của câu hỏi là D}
}
\end{ex}

\begin{ex} %[3D3V67]
Kim phút của đồng hồ quay được đúng 2 vòng rưỡi thì kim giờ quay được bao nhiêu phần của một vòng đồng hồ?
\choice
{5 khoảng giờ}
{\True Hai khoảng giờ rưỡi (tương ứng $2$ giờ $30$ phút)}
{1 vòng}
{Nửa vòng}
\loigiai{
Kim phút quay 1 vòng là 1 giờ. Quay 2 vòng rưỡi tức là 2 giờ 30 phút (kim giờ đi được 2 khoảng rưỡi).
\textbf{Đáp án của câu hỏi là B}
}
\end{ex}

\begin{ex} %[3D3V67]
Ngày cuối cùng của tháng 7 luôn luôn là ngày
\choice
{Ngày $30$}
{Ngày $28$}
{\True Ngày $31$}
{Ngày $29$}
\loigiai{
Tháng 7 có 31 ngày nên ngày cuối cùng là ngày 31.
\textbf{Đáp án của câu hỏi là C}
}
\end{ex}

\begin{ex} %[3D3V67]
Một chiếc đồng hồ báo thức reo chuông cách mỗi 15 phút một lần. Lần đầu reo lúc 6 giờ đúng. Lần thứ tư reo lúc mấy giờ?
\choice
{\True $6$ giờ $45$ phút}
{$7$ giờ}
{$6$ giờ $30$ phút}
{$7$ giờ $15$ phút}
\loigiai{
Lần 1: 6h00; Lần 2: 6h15; Lần 3: 6h30; Lần 4: 6h45.
\textbf{Đáp án của câu hỏi là A}
}
\end{ex}

\begin{ex} %[3D3V67]
Tháng 2 của một năm có 5 ngày Chủ nhật. Ngày Chủ nhật đầu tiên của tháng 2 đó chắc chắn là ngày
\choice
{Ngày $2$}
{Ngày $3$}
{Ngày $7$}
{\True Ngày $1$}
\loigiai{
Để có 5 ngày Chủ nhật thì tháng 2 phải có 29 ngày (năm nhuận) và ngày đầu tiên (ngày 1) phải là Chủ nhật (các CN là ngày 1, 8, 15, 22, 29).
\textbf{Đáp án của câu hỏi là D}
}
\end{ex}

\begin{ex} %[3D3V67]
Một người thợ bắt đầu làm việc lúc 7 giờ 30 phút, nghỉ trưa 1 giờ 30 phút và kết thúc công việc lúc 17 giờ. Thời gian làm việc thực tế của người thợ là
\choice
{$9$ giờ $30$ phút}
{\True $8$ giờ}
{$8$ giờ $30$ phút}
{$7$ giờ $30$ phút}
\loigiai{
Tổng thời gian từ 7h30 đến 17h00 là 9 giờ 30 phút. Trừ 1 giờ 30 phút nghỉ trưa: $9\text{h } 30\text{p} - 1\text{h } 30\text{p} = 8$ giờ.
\textbf{Đáp án của câu hỏi là B}
}
\end{ex}

\textbf{Phần 2. Trắc nghiệm đúng sai}
\setcounterpageref{ex}{0}

\begin{ex} %[3D3V67]
Xét tính đúng sai của việc thực hành xem lịch và đồng hồ:
\choiceTF
{\True Kim dài của đồng hồ quay 1 vòng thì hết đúng 60 phút (1 giờ).}
{Một tuần lễ có 5 ngày làm việc và không có ngày nào khác.}
{\True Hai ngày Chủ nhật liền kề nhau luôn cách nhau đúng 7 ngày.}
{Lúc 12 giờ trưa hai kim đồng hồ vuông góc với nhau.}
\loigiai{
\begin{itemize}
    \item Ý a) Đúng: Đặc tính hoạt động của đồng hồ.
    \item Ý b) Sai: Sai vì tuần lễ có 7 ngày.
    \item Ý c) Đúng: Quy luật tuần lễ.
    \item Ý d) Sai: Sai vì lúc 12 giờ hai kim trùng nhau.
\end{itemize}
\begin{itemize}
    \item a) \textbf{Đúng}.
    \item b) \textbf{Sai}.
    \item c) \textbf{Đúng}.
    \item d) \textbf{Sai}.
\end{itemize}
}
\end{ex}

\begin{ex} %[3D3V67]
Một chuyến du lịch kéo dài từ sáng ngày 28 tháng 5 đến hết chiều ngày 2 tháng 6 cùng năm:
\choiceTF
{\True Tháng 5 có 31 ngày.}
{\True Trong tháng 5 đoàn đi du lịch trong 4 ngày (28, 29, 30, 31).}
{\True Trong tháng 6 đoàn đi du lịch trong 2 ngày (1, 2).}
{Chuyến du lịch kéo dài đúng 10 ngày.}
\loigiai{
\begin{itemize}
    \item Ý a) Đúng: Kiến thức chuẩn.
    \item Ý b) Đúng: Bốn ngày cuối tháng 5.
    \item Ý c) Đúng: Hai ngày đầu tháng 6.
    \item Ý d) Sai: Sai vì tổng cộng là $4 + 2 = 6$ ngày.
\end{itemize}
\begin{itemize}
    \item a) \textbf{Đúng}.
    \item b) \textbf{Đúng}.
    \item c) \textbf{Đúng}.
    \item d) \textbf{Sai}.
\end{itemize}
}
\end{ex}

\textbf{Phần 3. Trả lời ngắn}
\setcounterpageref{ex}{0}

\begin{ex} %[3D3V67]
Một trận đấu cờ vua bắt đầu lúc 14 giờ 15 phút và kết thúc lúc 15 giờ 05 phút. Trận đấu diễn ra trong bao nhiêu phút?
\shortans{50}
\loigiai{
\textbf{Đáp số:} $50$\\
Từ 14h15 đến 15h05 là 50 phút.
}
\end{ex}

\begin{ex} %[3D3V67]
Tháng 10 có bao nhiêu ngày?
\shortans{31}
\loigiai{
\textbf{Đáp số:} $31$\\
Tháng 10 có 31 ngày.
}
\end{ex}

\begin{ex} %[3D3V67]
Nếu ngày 20 tháng 11 là Chủ nhật thì ngày 27 tháng 11 cùng năm là ngày thứ mấy (Chủ nhật quy ước ghi 8)?
\shortans{8}
\loigiai{
\textbf{Đáp số:} $8$\\
$20 + 7 = 27$ nên vẫn là Chủ nhật (8).
}
\end{ex}

\begin{ex} %[3D3V67]
Đổi 3 giờ ra đơn vị phút.
\shortans{180}
\loigiai{
\textbf{Đáp số:} $180$\\
$3 \times 60 = 180$ phút.
}
\end{ex}

\begin{ex} %[3D3V67]
Một chương trình truyền hình kéo dài từ 20 giờ 00 phút đến 20 giờ 45 phút. Chương trình kéo dài bao nhiêu phút?
\shortans{45}
\loigiai{
\textbf{Đáp số:} $45$\\
45 phút.
}
\end{ex}

\begin{ex} %[3D3V67]
Trong một năm có bao nhiêu tháng có 31 ngày?
\shortans{7}
\loigiai{
\textbf{Đáp số:} $7$\\
Gồm 7 tháng: 1, 3, 5, 7, 8, 10, 12.
}
\end{ex}

\newpage

\TaoPhieuKiemTra{Chủ đề 13 -- Xem đồng hồ. tháng -- năm. tiền việt nam}{68}{Tiền Việt Nam}{https://loptoanthaylinh.vn/thi-online/toan-3/bai-68}
\textbf{Phần 1. Trắc nghiệm lựa chọn}
\setcounterpageref{ex}{0}

\begin{ex} %[3D3V68]
Bạn Mai có 2 tờ tiền loại $20\,000$ đồng và 1 tờ loại $10\,000$ đồng. Tổng số tiền Mai có là
\choice
{$40\,000$ đồng}
{\True $50\,000$ đồng}
{$30\,000$ đồng}
{$60\,000$ đồng}
\loigiai{
$20\,000 \times 2 + 10\,000 = 50\,000$ đồng.
\textbf{Đáp án của câu hỏi là B}
}
\end{ex}

\begin{ex} %[3D3V68]
Một chiếc bút mực giá $35\,000$ đồng. Bạn Nam đưa cô bán hàng tờ tiền $50\,000$ đồng. Cô bán hàng phải trả lại Nam bao nhiêu tiền thừa?
\choice
{$25\,000$ đồng}
{$10\,000$ đồng}
{$20\,000$ đồng}
{\True $15\,000$ đồng}
\loigiai{
$50\,000 - 35\,000 = 15\,000$ đồng.
\textbf{Đáp án của câu hỏi là D}
}
\end{ex}

\begin{ex} %[3D3V68]
Đổi một tờ tiền $100\,000$ đồng lấy các tờ tiền mệnh giá $20\,000$ đồng. Số tờ tiền $20\,000$ đồng nhận được là
\choice
{\True $5$ tờ}
{$4$ tờ}
{$6$ tờ}
{$10$ tờ}
\loigiai{
$100\,000 : 20\,000 = 5$ tờ.
\textbf{Đáp án của câu hỏi là A}
}
\end{ex}

\begin{ex} %[3D3V68]
Một quyển truyện tranh giá $45\,000$ đồng, một hộp bút màu giá $38\,000$ đồng. Mua cả hai món đồ hết tất cả bao nhiêu tiền?
\choice
{$73\,000$ đồng}
{$85\,000$ đồng}
{\True $83\,000$ đồng}
{$80\,000$ đồng}
\loigiai{
$45\,000 + 38\,000 = 83\,000$ đồng.
\textbf{Đáp án của câu hỏi là C}
}
\end{ex}

\begin{ex} %[3D3V68]
Để có đúng $50\,000$ đồng bằng các tờ tiền loại $5000$ đồng thì cần bao nhiêu tờ tiền?
\choice
{$5$ tờ}
{\True $10$ tờ}
{$20$ tờ}
{$15$ tờ}
\loigiai{
$50\,000 : 5000 = 10$ tờ.
\textbf{Đáp án của câu hỏi là B}
}
\end{ex}

\begin{ex} %[3D3V68]
An mua 4 quyển vở, mỗi quyển giá $12\,000$ đồng và 1 chiếc thước kẻ giá $5000$ đồng. An đưa cô bán hàng tờ tiền $100\,000$ đồng. Cô bán hàng trả lại An số tiền thừa là
\choice
{$53\,000$ đồng}
{$45\,000$ đồng}
{$48\,000$ đồng}
{\True $47\,000$ đồng}
\loigiai{
Tổng tiền mua: $12\,000 \times 4 + 5000 = 48\,000 + 5000 = 53\,000$ đồng. Tiền thừa: $100\,000 - 53\,000 = 47\,000$ đồng.
\textbf{Đáp án của câu hỏi là D}
}
\end{ex}

\begin{ex} %[3D3V68]
Mệnh giá tiền giấy polymer nào sau đây KHÔNG tồn tại trong hệ thống tiền tệ hiện hành của Việt Nam?
\choice
{\True $30\,000$ đồng}
{$20\,000$ đồng}
{$50\,000$ đồng}
{$100\,000$ đồng}
\loigiai{
Không có tờ tiền mệnh giá $30\,000$ đồng.
\textbf{Đáp án của câu hỏi là A}
}
\end{ex}

\begin{ex} %[3D3V68]
Bình có 3 tờ $10\,000$ đồng và 4 tờ $5000$ đồng. Bình muốn mua một quả bóng giá $60\,000$ đồng. Hỏi Bình còn thiếu bao nhiêu tiền?
\choice
{$15\,000$ đồng}
{$20\,000$ đồng}
{\True $10\,000$ đồng}
{$5000$ đồng}
\loigiai{
Bình có: $10\,000 \times 3 + 5000 \times 4 = 30\,000 + 20\,000 = 50\,000$ đồng. Còn thiếu: $60\,000 - 50\,000 = 10\,000$ đồng.
\textbf{Đáp án của câu hỏi là C}
}
\end{ex}

\begin{ex} %[3D3V68]
Một chiếc cặp sách giá $85\,000$ đồng. Nhân dịp khai giảng được giảm giá $15\,000$ đồng. Giá tiền của chiếc cặp sách sau khi giảm là
\choice
{$75\,000$ đồng}
{\True $70\,000$ đồng}
{$80\,000$ đồng}
{$65\,000$ đồng}
\loigiai{
$85\,000 - 15\,000 = 70\,000$ đồng.
\textbf{Đáp án của câu hỏi là B}
}
\end{ex}

\begin{ex} %[3D3V68]
Có bao nhiêu cách đổi tờ tiền $20\,000$ đồng ra các tờ tiền mệnh giá $5000$ đồng và $10\,000$ đồng (dùng cả hai loại tiền)?
\choice
{$2$ cách}
{$3$ cách}
{Không đổi được}
{\True $1$ cách (gồm 1 tờ $10\,000$ đồng và 2 tờ $5000$ đồng)}
\loigiai{
Để có cả 2 loại tiền: $20\,000 = 10\,000 \times 1 + 5000 \times 2$. Chỉ có duy nhất 1 cách.
\textbf{Đáp án của câu hỏi là D}
}
\end{ex}

\begin{ex} %[3D3V68]
Mẹ đi siêu thị mua một khay thịt gà hết $65\,000$ đồng và một bó rau hết $15\,000$ đồng. Mẹ đưa 2 tờ tiền loại $50\,000$ đồng. Số tiền thừa nhận lại là
\choice
{\True $20\,000$ đồng}
{$30\,000$ đồng}
{$10\,000$ đồng}
{$25\,000$ đồng}
\loigiai{
Mẹ đưa: $50\,000 \times 2 = 100\,000$ đồng. Mua hết: $65\,000 + 15\,000 = 80\,000$ đồng. Tiền thừa: $100\,000 - 80\,000 = 20\,000$ đồng.
\textbf{Đáp án của câu hỏi là A}
}
\end{ex}

\begin{ex} %[3D3V68]
Mua 5 chiếc bánh bao hết $50\,000$ đồng. Hỏi mua 8 chiếc bánh bao như thế hết bao nhiêu tiền?
\choice
{$75\,000$ đồng}
{$85\,000$ đồng}
{\True $80\,000$ đồng}
{$70\,000$ đồng}
\loigiai{
Một chiếc bánh bao: $50\,000 : 5 = 10\,000$ đồng. Tám chiếc bánh bao: $10\,000 \times 8 = 80\,000$ đồng.
\textbf{Đáp án của câu hỏi là C}
}
\end{ex}

\textbf{Phần 2. Trắc nghiệm đúng sai}
\setcounterpageref{ex}{0}

\begin{ex} %[3D3V68]
Xét tính đúng sai của các nhận định về tiền tệ Việt Nam:
\choiceTF
{\True Đơn vị tiền tệ của nước ta là đồng.}
{Một tờ 50 000 đồng đổi được đúng 10 tờ 10 000 đồng.}
{\True Hai tờ 50 000 đồng có giá trị bằng một tờ 100 000 đồng.}
{Tờ tiền 100 000 đồng có giá trị nhỏ hơn tờ tiền 50 000 đồng.}
\loigiai{
\begin{itemize}
    \item Ý a) Đúng: Quy chuẩn quốc gia.
    \item Ý b) Sai: Sai vì $50\,000 : 10\,000 = 5$ tờ.
    \item Ý c) Đúng: $50\,000 \times 2 = 100\,000$ đồng.
    \item Ý d) Sai: Sai vì $100\,000 > 50\,000$.
\end{itemize}
\begin{itemize}
    \item a) \textbf{Đúng}.
    \item b) \textbf{Sai}.
    \item c) \textbf{Đúng}.
    \item d) \textbf{Sai}.
\end{itemize}
}
\end{ex}

\begin{ex} %[3D3V68]
Bạn Cúc có một con heo đất tiết kiệm gồm: 5 tờ 10 000 đồng, 4 tờ 5000 đồng và 10 tờ 2000 đồng:
\choiceTF
{\True Tổng số tiền loại 10 000 đồng là 50 000 đồng.}
{\True Tổng số tiền loại 5000 đồng là 20 000 đồng.}
{\True Tổng số tiền loại 2000 đồng là 20 000 đồng.}
{Cúc có tất cả đúng 100 000 đồng trong heo đất.}
\loigiai{
\begin{itemize}
    \item Ý a) Đúng: $10\,000 \times 5 = 50\,000$ đồng.
    \item Ý b) Đúng: $5000 \times 4 = 20\,000$ đồng.
    \item Ý c) Đúng: $2000 \times 10 = 20\,000$ đồng.
    \item Ý d) Sai: Sai vì tổng là $50 + 20 + 20 = 90\,000$ đồng.
\end{itemize}
\begin{itemize}
    \item a) \textbf{Đúng}.
    \item b) \textbf{Đúng}.
    \item c) \textbf{Đúng}.
    \item d) \textbf{Sai}.
\end{itemize}
}
\end{ex}

\textbf{Phần 3. Trả lời ngắn}
\setcounterpageref{ex}{0}

\begin{ex} %[3D3V68]
Một tờ 100 000 đồng đổi được bao nhiêu tờ 50 000 đồng?
\shortans{2}
\loigiai{
\textbf{Đáp số:} $2$\\
$100\,000 : 50\,000 = 2$ tờ.
}
\end{ex}

\begin{ex} %[3D3V68]
Lan có 3 tờ 20 000 đồng và 2 tờ 10 000 đồng. Hỏi Lan có tất cả bao nhiêu nghìn đồng?
\shortans{80}
\loigiai{
\textbf{Đáp số:} $80$\\
$20 \times 3 + 10 \times 2 = 60 + 20 = 80$ nghìn đồng.
}
\end{ex}

\begin{ex} %[3D3V68]
Mua 3 cốc trà sữa giá 25 000 đồng mỗi cốc. Hỏi phải trả tất cả bao nhiêu nghìn đồng?
\shortans{75}
\loigiai{
\textbf{Đáp số:} $75$\\
$25 \times 3 = 75$ nghìn đồng.
}
\end{ex}

\begin{ex} %[3D3V68]
Bác Nam mua đồ hết 72 000 đồng và đưa tờ 100 000 đồng. Người bán hàng trả lại bác bao nhiêu nghìn đồng?
\shortans{28}
\loigiai{
\textbf{Đáp số:} $28$\\
$100 - 72 = 28$ nghìn đồng.
}
\end{ex}

\begin{ex} %[3D3V68]
Cần bao nhiêu tờ 20 000 đồng để có được đúng 100 000 đồng?
\shortans{5}
\loigiai{
\textbf{Đáp số:} $5$\\
$100\,000 : 20\,000 = 5$ tờ.
}
\end{ex}

\begin{ex} %[3D3V68]
Một cuốn sổ giá 18 000 đồng, một chiếc bút giá 7000 đồng. Mua cả hai thứ hết bao nhiêu nghìn đồng?
\shortans{25}
\loigiai{
\textbf{Đáp số:} $25$\\
$18 + 7 = 25$ nghìn đồng.
}
\end{ex}

\newpage

\TaoPhieuKiemTra{Chủ đề 13 -- Xem đồng hồ. tháng -- năm. tiền việt nam}{69}{Luyện tập chung thời gian và tiền tệ}{https://loptoanthaylinh.vn/thi-online/toan-3/bai-69}
\textbf{Phần 1. Trắc nghiệm lựa chọn}
\setcounterpageref{ex}{0}

\begin{ex} %[3X3V69]
Một năm nhuận có bao nhiêu ngày?
\choice
{\True $366$ ngày}
{$365$ ngày}
{$360$ ngày}
{$364$ ngày}
\loigiai{
Năm nhuận có 366 ngày.
\textbf{Đáp án của câu hỏi là A}
}
\end{ex}

\begin{ex} %[3X3V69]
Đồng hồ chỉ 9 giờ 45 phút sáng. Đến 11 giờ trưa cùng ngày thì còn thiếu bao nhiêu phút nữa?
\choice
{$60$ phút}
{$15$ phút}
{\True $75$ phút}
{$45$ phút}
\loigiai{
Từ 9h45 đến 10h00 là 15 phút, từ 10h00 đến 11h00 là 60 phút. Tổng: $15 + 60 = 75$ phút.
\textbf{Đáp án của câu hỏi là C}
}
\end{ex}

\begin{ex} %[3X3V69]
Bạn Hùng có 4 tờ $20\,000$ đồng và 1 tờ $10\,000$ đồng. Hùng mua một chiếc mũ bảo hiểm giá $75\,000$ đồng. Hùng còn lại bao nhiêu tiền?
\choice
{$25\,000$ đồng}
{\True $15\,000$ đồng}
{$10\,000$ đồng}
{$20\,000$ đồng}
\loigiai{
Hùng có: $20\,000 \times 4 + 10\,000 = 90\,000$ đồng. Còn lại: $90\,000 - 75\,000 = 15\,000$ đồng.
\textbf{Đáp án của câu hỏi là B}
}
\end{ex}

\begin{ex} %[3X3V69]
Nếu ngày 26 tháng 12 là Chủ nhật thì ngày 2 tháng 1 năm sau là ngày thứ mấy?
\choice
{Thứ Bảy}
{Thứ Hai}
{Thứ Sáu}
{\True Chủ nhật}
\loigiai{
Tháng 12 có 31 ngày. Từ ngày 26 tháng 12 đến 2 tháng 1 có: $(31 - 26) + 2 = 5 + 2 = 7$ ngày. Đúng 1 tuần lễ nên cũng là ngày Chủ nhật.
\textbf{Đáp án của câu hỏi là D}
}
\end{ex}

\begin{ex} %[3X3V69]
Để đổi một tờ $50\,000$ đồng thành các tờ $10\,000$ đồng thì cần bao nhiêu tờ?
\choice
{\True $5$ tờ}
{$10$ tờ}
{$4$ tờ}
{$6$ tờ}
\loigiai{
$50\,000 : 10\,000 = 5$ tờ.
\textbf{Đáp án của câu hỏi là A}
}
\end{ex}

\begin{ex} %[3X3V69]
Một tiết học kéo dài 40 phút, bắt đầu lúc 8 giờ 15 phút. Tiết học kết thúc lúc mấy giờ?
\choice
{$8$ giờ $45$ phút}
{$9$ giờ}
{\True $8$ giờ $55$ phút}
{$8$ giờ $50$ phút}
\loigiai{
$15 + 40 = 55$ phút. Kết thúc lúc 8 giờ 55 phút.
\textbf{Đáp án của câu hỏi là C}
}
\end{ex}

\begin{ex} %[3X3V69]
Giá tiền 1 gói bánh là $24\,000$ đồng, 1 lon nước ngọt là $12\,000$ đồng. Bạn Lan mua 2 gói bánh và 3 lon nước ngọt hết bao nhiêu tiền?
\choice
{$72\,000$ đồng}
{\True $84\,000$ đồng}
{$90\,000$ đồng}
{$80\,000$ đồng}
\loigiai{
Tổng tiền: $24\,000 \times 2 + 12\,000 \times 3 = 48\,000 + 36\,000 = 84\,000$ đồng.
\textbf{Đáp án của câu hỏi là B}
}
\end{ex}

\begin{ex} %[3X3V69]
Tháng nào sau đây có 31 ngày?
\choice
{Tháng 4}
{Tháng 6}
{Tháng 9}
{\True Tháng 8}
\loigiai{
Tháng 8 có đúng 31 ngày.
\textbf{Đáp án của câu hỏi là D}
}
\end{ex}

\begin{ex} %[3X3V69]
Lúc 6 giờ kém 10 phút, kim dài của đồng hồ chỉ vào số nào?
\choice
{\True Số $10$}
{Số $2$}
{Số $8$}
{Số $11$}
\loigiai{
6 giờ kém 10 phút là 5 giờ 50 phút. Kim dài chỉ số 10.
\textbf{Đáp án của câu hỏi là A}
}
\end{ex}

\begin{ex} %[3X3V69]
Mẹ mua $3\text{ kg}$ cam hết $60\,000$ đồng. Hỏi mua $5\text{ kg}$ cam như thế hết bao nhiêu tiền?
\choice
{$90\,000$ đồng}
{$80\,000$ đồng}
{\True $100\,000$ đồng}
{$120\,000$ đồng}
\loigiai{
Một kg cam: $60\,000 : 3 = 20\,000$ đồng. Năm kg cam: $20\,000 \times 5 = 100\,000$ đồng.
\textbf{Đáp án của câu hỏi là C}
}
\end{ex}

\begin{ex} %[3X3V69]
Một ngày và một đêm gồm tất cả bao nhiêu giờ?
\choice
{$12$ giờ}
{\True $24$ giờ}
{$48$ giờ}
{$36$ giờ}
\loigiai{
Một ngày đêm có 24 giờ.
\textbf{Đáp án của câu hỏi là B}
}
\end{ex}

\begin{ex} %[3X3V69]
Có bao nhiêu tờ $20\,000$ đồng để đổi được 1 tờ $100\,000$ đồng?
\choice
{$10$ tờ}
{$4$ tờ}
{$6$ tờ}
{\True $5$ tờ}
\loigiai{
$100\,000 : 20\,000 = 5$ tờ.
\textbf{Đáp án của câu hỏi là D}
}
\end{ex}

\textbf{Phần 2. Trắc nghiệm đúng sai}
\setcounterpageref{ex}{0}

\begin{ex} %[3X3V69]
Xét tính đúng sai của các nhận định ôn tập thời gian và tiền tệ:
\choiceTF
{\True Năm nhuận có tháng 2 gồm 29 ngày.}
{Một năm có 5 quý.}
{\True Tờ tiền 100 000 đồng gấp 10 lần tờ tiền 10 000 đồng.}
{Lúc 3 giờ đúng, hai kim đồng hồ thẳng hàng nhau.}
\loigiai{
\begin{itemize}
    \item Ý a) Đúng: Định nghĩa năm nhuận.
    \item Ý b) Sai: Sai vì một năm chỉ có 4 quý.
    \item Ý c) Đúng: $100\,000 : 10\,000 = 10$.
    \item Ý d) Sai: Sai vì lúc 3 giờ hai kim vuông góc với nhau.
\end{itemize}
\begin{itemize}
    \item a) \textbf{Đúng}.
    \item b) \textbf{Sai}.
    \item c) \textbf{Đúng}.
    \item d) \textbf{Sai}.
\end{itemize}
}
\end{ex}

\begin{ex} %[3X3V69]
An đi siêu thị mang theo đúng 100 000 đồng:
\choiceTF
{\True An mua sách hết 45 000 đồng thì còn lại 55 000 đồng.}
{\True Sau đó An mua tiếp bút màu hết 30 000 đồng thì còn lại 25 000 đồng.}
{Số tiền An còn lại đủ mua thêm một quyển truyện giá 30 000 đồng.}
{\True Số tiền An đã tiêu sau hai lần mua là 75 000 đồng.}
\loigiai{
\begin{itemize}
    \item Ý a) Đúng: $100\,000 - 45\,000 = 55\,000$ đồng.
    \item Ý b) Đúng: $55\,000 - 30\,000 = 25\,000$ đồng.
    \item Ý c) Sai: Sai vì $25\,000 < 30\,000$ đồng.
    \item Ý d) Đúng: $45\,000 + 30\,000 = 75\,000$ đồng.
\end{itemize}
\begin{itemize}
    \item a) \textbf{Đúng}.
    \item b) \textbf{Đúng}.
    \item c) \textbf{Sai}.
    \item d) \textbf{Đúng}.
\end{itemize}
}
\end{ex}

\textbf{Phần 3. Trả lời ngắn}
\setcounterpageref{ex}{0}

\begin{ex} %[3X3V69]
Một năm có bao nhiêu tháng có đúng 31 ngày?
\shortans{7}
\loigiai{
\textbf{Đáp số:} $7$\\
Có 7 tháng: 1, 3, 5, 7, 8, 10, 12.
}
\end{ex}

\begin{ex} %[3X3V69]
Đổi 2 giờ 30 phút ra phút.
\shortans{150}
\loigiai{
\textbf{Đáp số:} $150$\\
$2 \times 60 + 30 = 150$ phút.
}
\end{ex}

\begin{ex} %[3X3V69]
Một tờ 50 000 đồng và 2 tờ 20 000 đồng có tổng giá trị bao nhiêu nghìn đồng?
\shortans{90}
\loigiai{
\textbf{Đáp số:} $90$\\
$50 + 20 \times 2 = 90$ nghìn đồng.
}
\end{ex}

\begin{ex} %[3X3V69]
Một tiết học bắt đầu lúc 9h00 và kết thúc lúc 9h45. Tiết học kéo dài bao nhiêu phút?
\shortans{45}
\loigiai{
\textbf{Đáp số:} $45$\\
45 phút.
}
\end{ex}

\begin{ex} %[3X3V69]
Tháng 4 có bao nhiêu ngày?
\shortans{30}
\loigiai{
\textbf{Đáp số:} $30$\\
Tháng 4 có 30 ngày.
}
\end{ex}

\begin{ex} %[3X3V69]
Mua 4 quyển vở hết 48 000 đồng. Hỏi 1 quyển vở giá bao nhiêu nghìn đồng?
\shortans{12}
\loigiai{
\textbf{Đáp số:} $12$\\
$48 : 4 = 12$ nghìn đồng.
}
\end{ex}

\newpage

% =========================================================================
% CHỦ ĐỀ 14 - NHÂN, CHIA TRONG PHẠM VI 100 000
% =========================================================================
\section{Nhân, chia trong phạm vi 100 000}

\TaoPhieuKiemTra{Chủ đề 14 -- Nhân, chia trong phạm vi 100 000}{70}{Nhân số có năm chữ số với số có một chữ số}{https://loptoanthaylinh.vn/thi-online/toan-3/bai-70}
\textbf{Phần 1. Trắc nghiệm lựa chọn}
\setcounterpageref{ex}{0}

\begin{ex} %[3D2V70]
Tính giá trị của phép tính: $18\,240 \times 4$.
\choice
{$72\,860$}
{$73\,960$}
{$71\,960$}
{\True $72\,960$}
\loigiai{
$18\,240 \times 4 = 72\,960$.
\textbf{Đáp án của câu hỏi là D}
}
\end{ex}

\begin{ex} %[3D2V70]
Một phân xưởng mỗi ngày may được $12\,500$ chiếc khẩu trang. Hỏi trong 6 ngày phân xưởng đó may được bao nhiêu chiếc khẩu trang?
\choice
{$70\,000$ chiếc}
{\True $75\,000$ chiếc}
{$80\,000$ chiếc}
{$72\,500$ chiếc}
\loigiai{
$12\,500 \times 6 = 75\,000$ chiếc.
\textbf{Đáp án của câu hỏi là B}
}
\end{ex}

\begin{ex} %[3D2V70]
Tìm số tự nhiên $x$, biết: $x : 5 = 16\,400$.
\choice
{$80\,000$}
{$84\,000$}
{\True $82\,000$}
{$81\,000$}
\loigiai{
$x = 16\,400 \times 5 = 82\,000$.
\textbf{Đáp án của câu hỏi là C}
}
\end{ex}

\begin{ex} %[3D2V70]
Một hình vuông có cạnh dài $14\,250\text{ m}$. Chu vi hình vuông đó là
\choice
{\True $57\,000\text{ m}$}
{$56\,000\text{ m}$}
{$58\,000\text{ m}$}
{$57\,500\text{ m}$}
\loigiai{
Chu vi: $14\,250 \times 4 = 57\,000\text{ (m)}$.
\textbf{Đáp án của câu hỏi là A}
}
\end{ex}

\begin{ex} %[3D2V70]
Tính nhanh giá trị: $12\,500 \times 2 + 12\,500 \times 6$.
\choice
{$90\,000$}
{$80\,000$}
{$125\,000$}
{\True $100\,000$}
\loigiai{
$12\,500 \times (2 + 6) = 12\,500 \times 8 = 100\,000$.
\textbf{Đáp án của câu hỏi là D}
}
\end{ex}

\begin{ex} %[3D2V70]
Một đoàn gồm 3 xe tải, mỗi xe chở $24\,500\text{ kg}$ xi măng. Người ta đã dỡ xuống $35\,000\text{ kg}$. Số xi măng còn lại trên xe là
\choice
{$40\,500\text{ kg}$}
{\True $38\,500\text{ kg}$}
{$37\,500\text{ kg}$}
{$39\,000\text{ kg}$}
\loigiai{
Tổng xi măng: $24\,500 \times 3 = 73\,500\text{ kg}$. Còn lại: $73\,500 - 35\,000 = 38\,500\text{ (kg)}$.
\textbf{Đáp án của câu hỏi là B}
}
\end{ex}

\begin{ex} %[3D2V70]
Tính giá trị biểu thức: $A = 21\,340 \times 3 - 18\,500$.
\choice
{$44\,520$}
{$46\,520$}
{\True $45\,520$}
{$45\,020$}
\loigiai{
$21\,340 \times 3 = 64\,020$; $64\,020 - 18\,500 = 45\,520$.
\textbf{Đáp án của câu hỏi là C}
}
\end{ex}

\begin{ex} %[3D2V70]
Một công ty nhập về 4 lô hàng, mỗi lô có $15\,200$ hộp sữa. Sau khi xuất bán $28\,000$ hộp, công ty còn lại bao nhiêu hộp sữa?
\choice
{\True $32\,800$ hộp}
{$30\,800$ hộp}
{$34\,800$ hộp}
{$33\,000$ hộp}
\loigiai{
Tổng sữa: $15\,200 \times 4 = 60\,800$ hộp. Còn lại: $60\,800 - 28\,000 = 32\,800$ hộp.
\textbf{Đáp án của câu hỏi là A}
}
\end{ex}

\begin{ex} %[3D2V70]
Nếu tăng một thừa số thêm 3 đơn vị và giữ nguyên thừa số kia bằng $25\,000$ thì tích tăng thêm bao nhiêu đơn vị?
\choice
{$50\,000$}
{$25\,000$}
{$100\,000$}
{\True $75\,000$}
\loigiai{
Tích tăng thêm: $25\,000 \times 3 = 75\,000$ đơn vị.
\textbf{Đáp án của câu hỏi là D}
}
\end{ex}

\begin{ex} %[3D2V70]
Tính: $11\,111 \times 9$. Kết quả là
\choice
{$99\,990$}
{\True $99\,999$}
{$100\,000$}
{$98\,999$}
\loigiai{
$11\,111 \times 9 = 99\,999$.
\textbf{Đáp án của câu hỏi là B}
}
\end{ex}

\begin{ex} %[3D2V70]
Một huyện có 5 xã vùng cao, mỗi xã được cấp $18\,500\text{ kg}$ gạo hỗ trợ mùa giáp hạt. Tổng số gạo đã cấp cho cả 5 xã là
\choice
{$90\,000\text{ kg}$}
{$91\,500\text{ kg}$}
{\True $92\,500\text{ kg}$}
{$93\,500\text{ kg}$}
\loigiai{
$18\,500 \times 5 = 92\,500\text{ (kg)}$.
\textbf{Đáp án của câu hỏi là C}
}
\end{ex}

\begin{ex} %[3D2V70]
Một cuộn dây thép dài $32\,400\text{ m}$. Ba cuộn dây thép như thế dài bao nhiêu mét?
\choice
{\True $97\,200\text{ m}$}
{$96\,200\text{ m}$}
{$98\,200\text{ m}$}
{$95\,200\text{ m}$}
\loigiai{
$32\,400 \times 3 = 97\,200\text{ (m)}$.
\textbf{Đáp án của câu hỏi là A}
}
\end{ex}

\textbf{Phần 2. Trắc nghiệm đúng sai}
\setcounterpageref{ex}{0}

\begin{ex} %[3D2V70]
Xét tính đúng sai của phép nhân số có 5 chữ số với số có 1 chữ số:
\choiceTF
{\True Nhân lần lượt từ phải sang trái (từ hàng đơn vị).}
{$20\,000 \times 5 = 10\,000$.}
{\True Tích của số có năm chữ số với 1 luôn bằng chính nó.}
{$15\,000 \times 4 = 50\,000$.}
\loigiai{
\begin{itemize}
    \item Ý a) Đúng: Quy tắc chuẩn.
    \item Ý b) Sai: Sai vì $20\,000 \times 5 = 100\,000$.
    \item Ý c) Đúng: Tính chất của số 1.
    \item Ý d) Sai: Sai vì $15\,000 \times 4 = 60\,000$.
\end{itemize}
\begin{itemize}
    \item a) \textbf{Đúng}.
    \item b) \textbf{Sai}.
    \item c) \textbf{Đúng}.
    \item d) \textbf{Sai}.
\end{itemize}
}
\end{ex}

\begin{ex} %[3D2V70]
Một xí nghiệp có 3 xưởng may, mỗi xưởng may được $14\,200$ bộ quần áo mỗi tháng:
\choiceTF
{\True Mỗi tháng xí nghiệp may được $42\,600$ bộ quần áo.}
{\True Trong 2 tháng một xưởng may được $28\,400$ bộ quần áo.}
{Cả xí nghiệp may trong 2 tháng được hơn 90 000 bộ.}
{\True Làm tròn sản lượng một tháng của xí nghiệp đến hàng nghìn là 43 000 bộ.}
\loigiai{
\begin{itemize}
    \item Ý a) Đúng: $14\,200 \times 3 = 42\,600$ bộ.
    \item Ý b) Đúng: $14\,200 \times 2 = 28\,400$ bộ.
    \item Ý c) Sai: Sai vì $42\,600 \times 2 = 85\,200 < 90\,000$ bộ.
    \item Ý d) Đúng: $42\,600$ làm tròn đến hàng nghìn là 43 000.
\end{itemize}
\begin{itemize}
    \item a) \textbf{Đúng}.
    \item b) \textbf{Đúng}.
    \item c) \textbf{Sai}.
    \item d) \textbf{Đúng}.
\end{itemize}
}
\end{ex}

\textbf{Phần 3. Trả lời ngắn}
\setcounterpageref{ex}{0}

\begin{ex} %[3D2V70]
Tính kết quả phép nhân: $15\,620 \times 5$.
\shortans{78100}
\loigiai{
\textbf{Đáp số:} $78100$\\
$15\,620 \times 5 = 78\,100$.
}
\end{ex}

\begin{ex} %[3D2V70]
Một hình vuông có cạnh dài $23\,500\text{ cm}$. Tính chu vi hình vuông đó theo xăng-ti-mét.
\shortans{94000}
\loigiai{
\textbf{Đáp số:} $94000$\\
$23\,500 \times 4 = 94\,000\text{ (cm)}$.
}
\end{ex}

\begin{ex} %[3D2V70]
Tìm số tự nhiên $y$, biết: $y : 3 = 24\,600$.
\shortans{73800}
\loigiai{
\textbf{Đáp số:} $73800$\\
$y = 24\,600 \times 3 = 73\,800$.
}
\end{ex}

\begin{ex} %[3D2V70]
Một nhà máy sản xuất mỗi ngày $16\,250$ sản phẩm. Hỏi 4 ngày nhà máy sản xuất bao nhiêu sản phẩm?
\shortans{65000}
\loigiai{
\textbf{Đáp số:} $65000$\\
$16\,250 \times 4 = 65\,000$ sản phẩm.
}
\end{ex}

\begin{ex} %[3D2V70]
Tính giá trị biểu thức: $12\,500 \times 6 - 25\,000$.
\shortans{50000}
\loigiai{
\textbf{Đáp số:} $50000$\\
$75\,000 - 25\,000 = 50\,000$.
}
\end{ex}

\begin{ex} %[3D2V70]
Có 4 kho lúa giống, mỗi kho chứa $18\,400\text{ kg}$. Hỏi 4 kho có tất cả bao nhiêu ki-lô-gam lúa giống?
\shortans{73600}
\loigiai{
\textbf{Đáp số:} $73600$\\
$18\,400 \times 4 = 73\,600\text{ (kg)}$.
}
\end{ex}

\newpage

\TaoPhieuKiemTra{Chủ đề 14 -- Nhân, chia trong phạm vi 100 000}{71}{Chia số có năm chữ số cho số có một chữ số}{https://loptoanthaylinh.vn/thi-online/toan-3/bai-71}
\textbf{Phần 1. Trắc nghiệm lựa chọn}
\setcounterpageref{ex}{0}

\begin{ex} %[3D2V71]
Tính giá trị của phép chia: $64\,824 : 4$.
\choice
{$16\,260$}
{$16\,204$}
{\True $16\,206$}
{$16\,106$}
\loigiai{
Chia lần lượt từ trái sang phải: $64\,824 : 4 = 16\,206$.
\textbf{Đáp án của câu hỏi là C}
}
\end{ex}

\begin{ex} %[3D2V71]
Một cuộn dây cáp dài $85\,450\text{ m}$ được chia đều thành 5 đoạn. Độ dài mỗi đoạn dây là
\choice
{\True $17\,090\text{ m}$}
{$17\,900\text{ m}$}
{$17\,080\text{ m}$}
{$16\,090\text{ m}$}
\loigiai{
$85\,450 : 5 = 17\,090\text{ (m)}$.
\textbf{Đáp án của câu hỏi là A}
}
\end{ex}

\begin{ex} %[3D2V71]
Tìm số tự nhiên $x$, biết: $x \times 6 = 78\,948$.
\choice
{$13\,148$}
{$13\,258$}
{$12\,158$}
{\True $13\,158$}
\loigiai{
$x = 78\,948 : 6 = 13\,158$.
\textbf{Đáp án của câu hỏi là D}
}
\end{ex}

\begin{ex} %[3D2V71]
Thực hiện phép chia có dư: $65\,487 : 5$. Thương và số dư lần lượt là
\choice
{$13\,097$ và dư $1$}
{\True $13\,097$ và dư $2$}
{$13\,970$ và dư $2$}
{$13\,096$ và dư $2$}
\loigiai{
$65\,487 = 5 \times 13\,097 + 2$ với số dư $2 < 5$.
\textbf{Đáp án của câu hỏi là B}
}
\end{ex}

\begin{ex} %[3D2V71]
Một kho hàng chia đều $48\,600\text{ kg}$ muối cho 6 cửa hàng. Mỗi cửa hàng đã bán được $\dfrac{1}{3}$ số muối của mình. Mỗi cửa hàng còn lại bao nhiêu ki-lô-gam muối?
\choice
{$2700\text{ kg}$}
{$8100\text{ kg}$}
{\True $5400\text{ kg}$}
{$6000\text{ kg}$}
\loigiai{
Mỗi cửa hàng nhận: $48\,600 : 6 = 8100\text{ kg}$. Đã bán: $8100 : 3 = 2700\text{ kg}$. Còn lại: $8100 - 2700 = 5400\text{ (kg)}$.
\textbf{Đáp án của câu hỏi là C}
}
\end{ex}

\begin{ex} %[3D2V71]
Trong một phép chia cho 7 có thương bằng $12\,450$ và số dư là số dư lớn nhất có thể có. Số bị chia là
\choice
{\True $87\,156$}
{$87\,150$}
{$87\,155$}
{$87\,157$}
\loigiai{
Số dư lớn nhất là 6. Số bị chia: $7 \times 12\,450 + 6 = 87\,150 + 6 = 87\,156$.
\textbf{Đáp án của câu hỏi là A}
}
\end{ex}

\begin{ex} %[3D2V71]
Một mảnh đất hình vuông có chu vi $65\,200\text{ m}$. Độ dài mỗi cạnh của mảnh đất đó là
\choice
{$16\,200\text{ m}$}
{$16\,400\text{ m}$}
{$15\,300\text{ m}$}
{\True $16\,300\text{ m}$}
\loigiai{
Cạnh hình vuông: $65\,200 : 4 = 16\,300\text{ (m)}$.
\textbf{Đáp án của câu hỏi là D}
}
\end{ex}

\begin{ex} %[3D2V71]
Tính giá trị của biểu thức: $A = 92\,480 : 8 - 4500$.
\choice
{$7160$}
{\True $7060$}
{$6960$}
{$7050$}
\loigiai{
$92\,480 : 8 = 11\,560$; $11\,560 - 4500 = 7060$.
\textbf{Đáp án của câu hỏi là B}
}
\end{ex}

\begin{ex} %[3D2V71]
Người ta chia đều $56\,700\ell$ xăng vào 7 bồn chứa. Mỗi bồn chứa bao nhiêu lít xăng?
\choice
{$8000\ell$}
{$8200\ell$}
{\True $8100\ell$}
{$7100\ell$}
\loigiai{
$56\,700 : 7 = 8100\text{ (\ell)}$.
\textbf{Đáp án của câu hỏi là C}
}
\end{ex}

\begin{ex} %[3D2V71]
Tìm một số biết rằng nếu lấy số đó chia cho 9 rồi bớt đi 1000 thì được 9000. Số đó là
\choice
{\True $90\,000$}
{$81\,000$}
{$80\,000$}
{$99\,000$}
\loigiai{
Số đó chia cho 9 bằng: $9000 + 1000 = 10\,000$. Số cần tìm: $10\,000 \times 9 = 90\,000$.
\textbf{Đáp án của câu hỏi là A}
}
\end{ex}

\begin{ex} %[3D2V71]
Một xí nghiệp có $72\,800$ chiếc áo cần đóng gói vào các thùng 8 chiếc. Số thùng đóng được là
\choice
{$9000$ thùng}
{$9200$ thùng}
{$9050$ thùng}
{\True $9100$ thùng}
\loigiai{
$72\,800 : 8 = 9100$ thùng.
\textbf{Đáp án của câu hỏi là D}
}
\end{ex}

\begin{ex} %[3D2V71]
Thực hiện phép chia: $100\,000 : 5$. Kết quả là
\choice
{$25\,000$}
{\True $20\,000$}
{$10\,000$}
{$50\,000$}
\loigiai{
$100\,000 : 5 = 20\,000$.
\textbf{Đáp án của câu hỏi là B}
}
\end{ex}

\textbf{Phần 2. Trắc nghiệm đúng sai}
\setcounterpageref{ex}{0}

\begin{ex} %[3D2V71]
Xét tính đúng sai của phép chia số có 5 chữ số cho số có 1 chữ số:
\choiceTF
{\True Phép chia thực hiện lần lượt từ trái sang phải.}
{\True $80\,040 : 8 = 10\,005$.}
{$80\,040 : 8 = 1005$.}
{Trong phép chia cho 6, số dư lớn nhất có thể là 6.}
\loigiai{
\begin{itemize}
    \item Ý a) Đúng: Quy tắc chuẩn.
    \item Ý b) Đúng: Chia có các chữ số 0 ở thương.
    \item Ý c) Sai: Sai vì thiếu chữ số 0 ở hàng chục nghìn/nghìn của thương.
    \item Ý d) Sai: Sai vì số dư lớn nhất là 5.
\end{itemize}
\begin{itemize}
    \item a) \textbf{Đúng}.
    \item b) \textbf{Đúng}.
    \item c) \textbf{Sai}.
    \item d) \textbf{Sai}.
\end{itemize}
}
\end{ex}

\begin{ex} %[3D2V71]
Một nông trường thu hoạch được $75\,000\text{ kg}$ cam chia đều cho 5 đội công nhân:
\choiceTF
{\True Mỗi đội nhận được $15\,000\text{ kg}$ cam.}
{\True Ba đội nhận được tất cả $45\,000\text{ kg}$ cam.}
{Hai đội nhận được nhiều hơn một nửa tổng số cam.}
{\True Tổng số cam chia hết, không dư ki-lô-gam nào.}
\loigiai{
\begin{itemize}
    \item Ý a) Đúng: $75\,000 : 5 = 15\,000\text{ kg}$.
    \item Ý b) Đúng: $15\,000 \times 3 = 45\,000\text{ kg}$.
    \item Ý c) Sai: Sai vì hai đội nhận 30 000 kg, ít hơn nửa (37 500 kg).
    \item Ý d) Đúng: 75 000 chia hết cho 5.
\end{itemize}
\begin{itemize}
    \item a) \textbf{Đúng}.
    \item b) \textbf{Đúng}.
    \item c) \textbf{Sai}.
    \item d) \textbf{Đúng}.
\end{itemize}
}
\end{ex}

\textbf{Phần 3. Trả lời ngắn}
\setcounterpageref{ex}{0}

\begin{ex} %[3D2V71]
Tính kết quả phép chia: $96\,848 : 8$.
\shortans{12106}
\loigiai{
\textbf{Đáp số:} $12106$\\
$96\,848 : 8 = 12\,106$.
}
\end{ex}

\begin{ex} %[3D2V71]
Tìm số tự nhiên $x$, biết: $x \times 7 = 84\,735$.
\shortans{12105}
\loigiai{
\textbf{Đáp số:} $12105$\\
$x = 84\,735 : 7 = 12\,105$.
}
\end{ex}

\begin{ex} %[3D2V71]
Một cuộn dây dài $54\,600\text{ m}$ chia đều thành 6 đoạn. Mỗi đoạn dài bao nhiêu mét?
\shortans{9100}
\loigiai{
\textbf{Đáp số:} $9100$\\
$54\,600 : 6 = 9100\text{ (m)}$.
}
\end{ex}

\begin{ex} %[3D2V71]
Tìm số dư của phép chia: $85\,693 : 9$.
\shortans{4}
\loigiai{
\textbf{Đáp số:} $4$\\
$85\,693 = 9 \times 9521 + 4$. Số dư là 4.
}
\end{ex}

\begin{ex} %[3D2V71]
Một hình vuông có chu vi $72\,400\text{ m}$. Tính độ dài cạnh hình vuông theo mét.
\shortans{18100}
\loigiai{
\textbf{Đáp số:} $18100$\\
$72\,400 : 4 = 18\,100\text{ (m)}$.
}
\end{ex}

\begin{ex} %[3D2V71]
Chia đều $90\,000\text{ kg}$ gạo cho 9 cửa hàng. Hỏi mỗi cửa hàng nhận được bao nhiêu ki-lô-gam gạo?
\shortans{10000}
\loigiai{
\textbf{Đáp số:} $10000$\\
$90\,000 : 9 = 10\,000\text{ (kg)}$.
}
\end{ex}

\newpage

\TaoPhieuKiemTra{Chủ đề 14 -- Nhân, chia trong phạm vi 100 000}{72}{Luyện tập chung nhân chia 100 000}{https://loptoanthaylinh.vn/thi-online/toan-3/bai-72}
\textbf{Phần 1. Trắc nghiệm lựa chọn}
\setcounterpageref{ex}{0}

\begin{ex} %[3D2V72]
Tính giá trị của biểu thức: $A = 15\,000 \times 4 + 48\,000 : 6$.
\choice
{$65\,000$}
{\True $68\,000$}
{$70\,000$}
{$66\,000$}
\loigiai{
$15\,000 \times 4 = 60\,000$; $48\,000 : 6 = 8000$; $60\,000 + 8000 = 68\,000$.
\textbf{Đáp án của câu hỏi là B}
}
\end{ex}

\begin{ex} %[3D2V72]
Một cửa hàng có $84\,000\text{ kg}$ gạo. Ngày đầu bán $\dfrac{1}{4}$ số gạo, ngày thứ hai bán $\dfrac{1}{3}$ số gạo còn lại. Sau hai ngày cửa hàng còn lại bao nhiêu ki-lô-gam gạo?
\choice
{$21\,000\text{ kg}$}
{$35\,000\text{ kg}$}
{$48\,000\text{ kg}$}
{\True $42\,000\text{ kg}$}
\loigiai{
Ngày 1 bán: $84\,000 : 4 = 21\,000\text{ kg}$. Còn lại: $84\,000 - 21\,000 = 63\,000\text{ kg}$. Ngày 2 bán: $63\,000 : 3 = 21\,000\text{ kg}$. Còn lại: $63\,000 - 21\,000 = 42\,000\text{ (kg)}$.
\textbf{Đáp án của câu hỏi là D}
}
\end{ex}

\begin{ex} %[3D2V72]
Tìm $x$, biết: $(x - 15\,000) : 3 = 25\,000$.
\choice
{\True $90\,000$}
{$85\,000$}
{$95\,000$}
{$100\,000$}
\loigiai{
$x - 15\,000 = 25\,000 \times 3 = 75\,000 \Rightarrow x = 75\,000 + 15\,000 = 90\,000$.
\textbf{Đáp án của câu hỏi là A}
}
\end{ex}

\begin{ex} %[3D2V72]
Tính nhanh giá trị: $16\,500 \times 3 + 16\,500 \times 2 - 16\,500$.
\choice
{$65\,000$}
{$60\,000$}
{\True $66\,000$}
{$70\,000$}
\loigiai{
$16\,500 \times (3 + 2 - 1) = 16\,500 \times 4 = 66\,000$.
\textbf{Đáp án của câu hỏi là C}
}
\end{ex}

\begin{ex} %[3D2V72]
Trong phép chia cho 8 có thương là $11\,250$ và số dư là số dư lớn nhất có thể có. Số bị chia là
\choice
{$90\,000$}
{\True $90\,007$}
{$90\,008$}
{$90\,006$}
\loigiai{
Số dư lớn nhất là 7. Số bị chia: $8 \times 11\,250 + 7 = 90\,000 + 7 = 90\,007$.
\textbf{Đáp án của câu hỏi là B}
}
\end{ex}

\begin{ex} %[3D2V72]
Một mảnh vườn hình chữ nhật có chu vi $64\,000\text{ m}$, chiều rộng bằng $\dfrac{1}{3}$ chiều dài. Chiều dài mảnh vườn là
\choice
{$8000\text{ m}$}
{$16\,000\text{ m}$}
{$32\,000\text{ m}$}
{\True $24\,000\text{ m}$}
\loigiai{
Nửa chu vi: $64\,000 : 2 = 32\,000\text{ m}$. Rộng: $32\,000 : (3 + 1) = 8000\text{ m}$. Dài: $8000 \times 3 = 24\,000\text{ (m)}$.
\textbf{Đáp án của câu hỏi là D}
}
\end{ex}

\begin{ex} %[3D2V72]
Tính giá trị của biểu thức: $100\,000 - 18\,500 \times 4$.
\choice
{\True $26\,000$}
{$24\,000$}
{$28\,000$}
{$25\,000$}
\loigiai{
$18\,500 \times 4 = 74\,000$; $100\,000 - 74\,000 = 26\,000$.
\textbf{Đáp án của câu hỏi là A}
}
\end{ex}

\begin{ex} %[3D2V72]
Một nhà máy có 5 phân xưởng, mỗi phân xưởng sản xuất được $16\,800$ sản phẩm. Nhà máy đã xuất kho $55\,000$ sản phẩm. Số sản phẩm còn lại là
\choice
{$28\,000$ sản phẩm}
{$30\,000$ sản phẩm}
{\True $29\,000$ sản phẩm}
{$31\,000$ sản phẩm}
\loigiai{
Tổng: $16\,800 \times 5 = 84\,000$. Còn lại: $84\,000 - 55\,000 = 29\,000$ sản phẩm.
\textbf{Đáp án của câu hỏi là C}
}
\end{ex}

\begin{ex} %[3D2V72]
Tìm số tự nhiên $y$, biết: $y \times 5 + 24\,500 = 99\,500$.
\choice
{$14\,000$}
{\True $15\,000$}
{$16\,000$}
{$18\,000$}
\loigiai{
$y \times 5 = 99\,500 - 24\,500 = 75\,000 \Rightarrow y = 75\,000 : 5 = 15\,000$.
\textbf{Đáp án của câu hỏi là B}
}
\end{ex}

\begin{ex} %[3D2V72]
Một hình vuông có chu vi bằng chu vi hình chữ nhật có chiều dài $25\,000\text{ m}$ và chiều rộng $15\,000\text{ m}$. Cạnh của hình vuông là
\choice
{$40\,000\text{ m}$}
{$18\,000\text{ m}$}
{$22\,000\text{ m}$}
{\True $20\,000\text{ m}$}
\loigiai{
Chu vi hình chữ nhật: $(25\,000 + 15\,000) \times 2 = 80\,000\text{ m}$. Cạnh hình vuông: $80\,000 : 4 = 20\,000\text{ (m)}$.
\textbf{Đáp án của câu hỏi là D}
}
\end{ex}

\begin{ex} %[3D2V72]
Bác Năm có $90\,000\text{ kg}$ thóc. Bác chia đều vào 6 kho, sau đó bán đi 2 kho. Số thóc bác Năm còn lại là
\choice
{\True $60\,000\text{ kg}$}
{$30\,000\text{ kg}$}
{$45\,000\text{ kg}$}
{$50\,000\text{ kg}$}
\loigiai{
Mỗi kho có: $90\,000 : 6 = 15\,000\text{ kg}$. Còn lại 4 kho: $15\,000 \times 4 = 60\,000\text{ (kg)}$.
\textbf{Đáp án của câu hỏi là A}
}
\end{ex}

\begin{ex} %[3D2V72]
Tổng hai số là 80 000, số lớn gấp 4 lần số bé. Số lớn là
\choice
{$16\,000$}
{$60\,000$}
{\True $64\,000$}
{$70\,000$}
\loigiai{
Tổng số phần: $4 + 1 = 5$ phần. Số bé: $80\,000 : 5 = 16\,000$. Số lớn: $16\,000 \times 4 = 64\,000$.
\textbf{Đáp án của câu hỏi là C}
}
\end{ex}

\textbf{Phần 2. Trắc nghiệm đúng sai}
\setcounterpageref{ex}{0}

\begin{ex} %[3D2V72]
Xét tính đúng sai của các biểu thức tính toán phạm vi 100 000:
\choiceTF
{\True Biểu thức $20\,000 \times 4 : 2$ có giá trị bằng $40\,000$.}
{$100\,000 : (2 \times 5) = 100\,000$.}
{\True Trong biểu thức có ngoặc, thực hiện phép tính trong ngoặc trước.}
{Một số nhân với 0 luôn bằng chính số đó.}
\loigiai{
\begin{itemize}
    \item Ý a) Đúng: $80\,000 : 2 = 40\,000$.
    \item Ý b) Sai: Sai vì $100\,000 : 10 = 10\,000$.
    \item Ý c) Đúng: Quy tắc cơ bản.
    \item Ý d) Sai: Sai vì nhân với 0 luôn bằng 0.
\end{itemize}
\begin{itemize}
    \item a) \textbf{Đúng}.
    \item b) \textbf{Sai}.
    \item c) \textbf{Đúng}.
    \item d) \textbf{Sai}.
\end{itemize}
}
\end{ex}

\begin{ex} %[3D2V72]
Một kho hàng có $96\,000\text{ kg}$ phân bón chia đều vào 8 kho nhỏ:
\choiceTF
{\True Mỗi kho nhỏ chứa $12\,000\text{ kg}$ phân bón.}
{\True Ba kho nhỏ chứa tất cả $36\,000\text{ kg}$ phân bón.}
{Năm kho nhỏ chứa nhiều hơn 70 000 kg phân bón.}
{\True Nếu xuất đi 4 kho nhỏ thì trong kho còn lại một nửa số phân bón ban đầu.}
\loigiai{
\begin{itemize}
    \item Ý a) Đúng: $96\,000 : 8 = 12\,000\text{ kg}$.
    \item Ý b) Đúng: $12\,000 \times 3 = 36\,000\text{ kg}$.
    \item Ý c) Sai: Sai vì $12\,000 \times 5 = 60\,000 < 70\,000$ kg.
    \item Ý d) Đúng: $4 : 8 = \dfrac{1}{2}$.
\end{itemize}
\begin{itemize}
    \item a) \textbf{Đúng}.
    \item b) \textbf{Đúng}.
    \item c) \textbf{Sai}.
    \item d) \textbf{Đúng}.
\end{itemize}
}
\end{ex}

\textbf{Phần 3. Trả lời ngắn}
\setcounterpageref{ex}{0}

\begin{ex} %[3D2V72]
Tính giá trị của biểu thức: $18\,000 \times 4 - 32\,000 : 4$.
\shortans{64000}
\loigiai{
\textbf{Đáp số:} $64000$\\
$72\,000 - 8000 = 64\,000$.
}
\end{ex}

\begin{ex} %[3D2V72]
Một hình vuông có cạnh dài $16\,500\text{ m}$. Tính chu vi hình vuông theo mét.
\shortans{66000}
\loigiai{
\textbf{Đáp số:} $66000$\\
$16\,500 \times 4 = 66\,000\text{ (m)}$.
}
\end{ex}

\begin{ex} %[3D2V72]
Tìm số tự nhiên $y$, biết: $(y - 20\,000) : 4 = 15\,000$.
\shortans{80000}
\loigiai{
\textbf{Đáp số:} $80000$\\
$y - 20\,000 = 60\,000 \Rightarrow y = 80\,000$.
}
\end{ex}

\begin{ex} %[3D2V72]
Tính nhanh: $25\,000 \times 3 + 25\,000$.
\shortans{100000}
\loigiai{
\textbf{Đáp số:} $100000$\\
$25\,000 \times 4 = 100\,000$.
}
\end{ex}

\begin{ex} %[3D2V72]
Chia đều $84\,000\text{ kg}$ gạo cho 7 xã. Mỗi xã nhận được bao nhiêu ki-lô-gam gạo?
\shortans{12000}
\loigiai{
\textbf{Đáp số:} $12000$\\
$84\,000 : 7 = 12\,000\text{ (kg)}$.
}
\end{ex}

\begin{ex} %[3D2V72]
Một cửa hàng có $95\,000$ sản phẩm. Bán đi $\dfrac{1}{5}$ số sản phẩm đó. Hỏi cửa hàng còn lại bao nhiêu sản phẩm?
\shortans{76000}
\loigiai{
\textbf{Đáp số:} $76000$\\
Đã bán: $19\,000$. Còn lại: $95\,000 - 19\,000 = 76\,000$ sản phẩm.
}
\end{ex}

\newpage

% =========================================================================
% CHỦ ĐỀ 15 - LÀM QUEN VỚI YẾU TỐ THỐNG KÊ, XÁC SUẤT
% =========================================================================
\section{Làm quen với yếu tố thống kê, xác suất}

\TaoPhieuKiemTra{Chủ đề 15 -- Làm quen với yếu tố thống kê, xác suất}{73}{Thu thập, phân loại, ghi chép số liệu}{https://loptoanthaylinh.vn/thi-online/toan-3/bai-73}
\textbf{Phần 1. Trắc nghiệm lựa chọn}
\setcounterpageref{ex}{0}

\begin{ex} %[3X1V73]
Bảng thống kê số cây trồng của 4 tổ: Tổ 1: 18 cây; Tổ 2: 24 cây; Tổ 3: 16 cây; Tổ 4: 22 cây. Tổ trồng được nhiều cây nhất là
\choice
{Tổ 1}
{Tổ 4}
{\True Tổ 2}
{Tổ 3}
\loigiai{
Tổ 2 trồng được 24 cây, nhiều nhất trong 4 tổ.
\textbf{Đáp án của câu hỏi là C}
}
\end{ex}

\begin{ex} %[3X1V73]
Tổng số cây trồng của cả 4 tổ ở câu trên là
\choice
{\True $80$ cây}
{$75$ cây}
{$85$ cây}
{$90$ cây}
\loigiai{
Tổng số cây: $18 + 24 + 16 + 22 = 80$ cây.
\textbf{Đáp án của câu hỏi là A}
}
\end{ex}

\begin{ex} %[3X1V73]
Tổ trồng nhiều cây nhất nhiều hơn tổ trồng ít cây nhất bao nhiêu cây?
\choice
{$6$ cây}
{$7$ cây}
{$9$ cây}
{\True $8$ cây}
\loigiai{
Nhiều nhất là Tổ 2 (24 cây), ít nhất là Tổ 3 (16 cây). Chênh lệch: $24 - 16 = 8$ cây.
\textbf{Đáp án của câu hỏi là D}
}
\end{ex}

\begin{ex} %[3X1V73]
Số học sinh giỏi của ba lớp: Lớp 3A có 15 bạn, Lớp 3B có 18 bạn, Lớp 3C có 12 bạn. Khẳng định nào sau đây là đúng?
\choice
{Lớp 3C có nhiều học sinh giỏi nhất}
{\True Lớp 3B có nhiều học sinh giỏi nhất}
{Lớp 3A có ít học sinh giỏi nhất}
{Cả ba lớp có 40 học sinh giỏi}
\loigiai{
Lớp 3B có 18 bạn học sinh giỏi, nhiều nhất.
\textbf{Đáp án của câu hỏi là B}
}
\end{ex}

\begin{ex} %[3X1V73]
Thống kê số điểm 10 trong tuần của Mai: Thứ Hai: 3; Thứ Ba: 2; Thứ Tư: 4; Thứ Năm: 1; Thứ Sáu: 2. Ngày Mai đạt nhiều điểm 10 nhất là
\choice
{Thứ Hai}
{Thứ Sáu}
{\True Thứ Tư}
{Thứ Ba}
\loigiai{
Thứ Tư Mai đạt 4 điểm 10, nhiều nhất.
\textbf{Đáp án của câu hỏi là C}
}
\end{ex}

\begin{ex} %[3X1V73]
Tổng số điểm 10 Mai đạt được trong cả tuần là
\choice
{\True $12$ điểm $10$}
{$10$ điểm $10$}
{$11$ điểm $10$}
{$14$ điểm $10$}
\loigiai{
Tổng số điểm 10: $3 + 2 + 4 + 1 + 2 = 12$ điểm 10.
\textbf{Đáp án của câu hỏi là A}
}
\end{ex}

\begin{ex} %[3X1V73]
Một cửa hàng bán được trong 3 ngày: Ngày 1: $120\text{ kg}$; Ngày 2 gấp đôi Ngày 1; Ngày 3 ít hơn Ngày 2 là $50\text{ kg}$. Ngày 3 bán được bao nhiêu ki-lô-gam?
\choice
{$240\text{ kg}$}
{$200\text{ kg}$}
{$180\text{ kg}$}
{\True $190\text{ kg}$}
\loigiai{
Ngày 2: $120 \times 2 = 240\text{ kg}$. Ngày 3: $240 - 50 = 190\text{ (kg)}$.
\textbf{Đáp án của câu hỏi là D}
}
\end{ex}

\begin{ex} %[3X1V73]
Số lượng học sinh đăng ký các câu lạc bộ: Bơi lội: 45 bạn; Cờ vua: 30 bạn; Bóng đá: 60 bạn; Cầu lông: 25 bạn. Câu lạc bộ có ít bạn đăng ký nhất là
\choice
{Cờ vua}
{\True Cầu lông}
{Bơi lội}
{Bóng đá}
\loigiai{
Cầu lông có 25 bạn đăng ký, ít nhất.
\textbf{Đáp án của câu hỏi là B}
}
\end{ex}

\begin{ex} %[3X1V73]
Số bạn đăng ký Bóng đá gấp mấy lần số bạn đăng ký Cờ vua?
\choice
{$3$ lần}
{$1$ lần}
{\True $2$ lần}
{$4$ lần}
\loigiai{
$60 : 30 = 2$ lần.
\textbf{Đáp án của câu hỏi là C}
}
\end{ex}

\begin{ex} %[3X1V73]
Bảng thống kê khối lượng thu hoạch lúa của 3 thửa ruộng: Thửa A: $850\text{ kg}$; Thửa B: $920\text{ kg}$; Thửa C: $780\text{ kg}$. Tổng khối lượng lúa cả ba thửa ruộng là
\choice
{\True $2550\text{ kg}$}
{$2500\text{ kg}$}
{$2600\text{ kg}$}
{$2450\text{ kg}$}
\loigiai{
$850 + 920 + 780 = 2550\text{ (kg)}$.
\textbf{Đáp án của câu hỏi là A}
}
\end{ex}

\begin{ex} %[3X1V73]
Để kiểm tra sĩ số và tình hình chuyên cần của học sinh hàng ngày, cán bộ lớp thường sử dụng
\choice
{Thước đo}
{Nhiệt kế}
{Compa}
{\True Sổ điểm danh (bảng theo dõi số liệu)}
\loigiai{
Sổ theo dõi/điểm danh là công cụ thống kê số liệu hàng ngày.
\textbf{Đáp án của câu hỏi là D}
}
\end{ex}

\begin{ex} %[3X1V73]
Trong một bài khảo sát về con vật nuôi yêu thích của 100 học sinh: Chó: 40 bạn; Mèo: 35 bạn; Cá cảnh: 15 bạn; Chim cảnh: 10 bạn. Số bạn thích Mèo nhiều hơn số bạn thích Cá cảnh là
\choice
{$25$ bạn}
{\True $20$ bạn}
{$15$ bạn}
{$10$ bạn}
\loigiai{
$35 - 15 = 20$ bạn.
\textbf{Đáp án của câu hỏi là B}
}
\end{ex}

\textbf{Phần 2. Trắc nghiệm đúng sai}
\setcounterpageref{ex}{0}

\begin{ex} %[3X1V73]
Xét tính đúng sai của việc đọc và xử lý bảng số liệu:
\choiceTF
{\True Bảng số liệu giúp ta so sánh và tính toán các đại lượng dễ dàng hơn.}
{Các số liệu trong bảng luôn luôn phải được sắp xếp theo thứ tự tăng dần.}
{\True Từ bảng số liệu có thể tìm ra giá trị lớn nhất và nhỏ nhất.}
{Không thể tính tổng các số liệu trong cùng một cột của bảng.}
\loigiai{
\begin{itemize}
    \item Ý a) Đúng: Mục đích thống kê.
    \item Ý b) Sai: Sai vì có thể sắp xếp theo thời gian hoặc danh mục.
    \item Ý c) Đúng: Kĩ năng phân tích.
    \item Ý d) Sai: Sai vì ta có thể cộng tổng theo cột.
\end{itemize}
\begin{itemize}
    \item a) \textbf{Đúng}.
    \item b) \textbf{Sai}.
    \item c) \textbf{Đúng}.
    \item d) \textbf{Sai}.
\end{itemize}
}
\end{ex}

\begin{ex} %[3X1V73]
Số học sinh giỏi 4 khối lớp: Khối 1: 80 em; Khối 2: 95 em; Khối 3: 85 em; Khối 4: 90 em:
\choiceTF
{\True Khối 2 có nhiều học sinh giỏi nhất.}
{\True Khối 1 có ít học sinh giỏi nhất.}
{Tổng số học sinh giỏi của Khối 1 và Khối 3 bằng Khối 2.}
{\True Tổng số học sinh giỏi cả 4 khối là 350 em.}
\loigiai{
\begin{itemize}
    \item Ý a) Đúng: 95 em cao nhất.
    \item Ý b) Đúng: 80 em ít nhất.
    \item Ý c) Sai: Sai vì $80 + 85 = 165 \neq 95$.
    \item Ý d) Đúng: $80 + 95 + 85 + 90 = 350$ em.
\end{itemize}
\begin{itemize}
    \item a) \textbf{Đúng}.
    \item b) \textbf{Đúng}.
    \item c) \textbf{Sai}.
    \item d) \textbf{Đúng}.
\end{itemize}
}
\end{ex}

\textbf{Phần 3. Trả lời ngắn}
\setcounterpageref{ex}{0}

\begin{ex} %[3X1V73]
Số học sinh nam của 3 lớp lần lượt là: 18, 22, 20. Tính tổng số học sinh nam của cả 3 lớp.
\shortans{60}
\loigiai{
\textbf{Đáp số:} $60$\\
$18 + 22 + 20 = 60$ bạn.
}
\end{ex}

\begin{ex} %[3X1V73]
Tổ 1 gom được $45\text{ kg}$ giấy vụn, Tổ 2 gom được $55\text{ kg}$ giấy vụn. Cả hai tổ gom được bao nhiêu ki-lô-gam giấy vụn?
\shortans{100}
\loigiai{
\textbf{Đáp số:} $100$\\
$45 + 55 = 100\text{ (kg)}$.
}
\end{ex}

\begin{ex} %[3X1V73]
Bốn bạn quyên góp sách: An: 8 cuốn, Bình: 12 cuốn, Cúc: 10 cuốn, Dũng: 14 cuốn. Bạn góp nhiều nhất nhiều hơn bạn góp ít nhất bao nhiêu cuốn?
\shortans{6}
\loigiai{
\textbf{Đáp số:} $6$\\
$14 - 8 = 6$ cuốn.
}
\end{ex}

\begin{ex} %[3X1V73]
Số cây trồng được của lớp 3A trong 3 ngày: Ngày 1: 25 cây; Ngày 2: 30 cây; Ngày 3: 35 cây. Tính trung bình mỗi ngày lớp 3A trồng được bao nhiêu cây?
\shortans{30}
\loigiai{
\textbf{Đáp số:} $30$\\
$(25 + 30 + 35) : 3 = 90 : 3 = 30$ cây.
}
\end{ex}

\begin{ex} %[3X1V73]
Trong một hộp có 12 chiếc bút màu đỏ và 18 chiếc bút màu xanh. Hỏi số bút màu xanh nhiều hơn số bút màu đỏ bao nhiêu chiếc?
\shortans{6}
\loigiai{
\textbf{Đáp số:} $6$\\
$18 - 12 = 6$ chiếc.
}
\end{ex}

\begin{ex} %[3X1V73]
Một nông trại nuôi 3 chuồng lợn: chuồng A có 45 con, chuồng B có 55 con, chuồng C có 60 con. Hỏi chuồng C nhiều hơn chuồng A bao nhiêu con lợn?
\shortans{15}
\loigiai{
\textbf{Đáp số:} $15$\\
$60 - 45 = 15$ con.
}
\end{ex}

\newpage

\TaoPhieuKiemTra{Chủ đề 15 -- Làm quen với yếu tố thống kê, xác suất}{74}{Khả năng xảy ra của một sự kiện}{https://loptoanthaylinh.vn/thi-online/toan-3/bai-74}
\textbf{Phần 1. Trắc nghiệm lựa chọn}
\setcounterpageref{ex}{0}

\begin{ex} %[3X1V74]
Trong một hộp có 5 viên bi màu xanh và không có viên bi nào màu khác. Nhắm mắt lấy ra 1 viên bi. Sự kiện nào sau đây là 'Chắc chắn' xảy ra?
\choice
{Lấy được 1 viên bi màu đỏ}
{\True Lấy được 1 viên bi màu xanh}
{Lấy được 1 viên bi màu vàng}
{Không lấy được viên bi nào}
\loigiai{
Vì trong hộp chỉ toàn bi xanh nên chắc chắn sẽ lấy được viên bi màu xanh.
\textbf{Đáp án của câu hỏi là B}
}
\end{ex}

\begin{ex} %[3X1V74]
Gieo một con xúc xắc chuẩn có 6 mặt ghi từ 1 đến 6 chấm. Sự kiện nào sau đây là 'Không thể' xảy ra?
\choice
{Mặt xuất hiện có 6 chấm}
{Mặt xuất hiện có 1 chấm}
{Mặt xuất hiện có số chấm là số chẵn}
{\True Mặt xuất hiện có 7 chấm}
\loigiai{
Xúc xắc chỉ có từ 1 đến 6 chấm nên không thể xuất hiện mặt 7 chấm.
\textbf{Đáp án của câu hỏi là D}
}
\end{ex}

\begin{ex} %[3X1V74]
Trong túi có 3 chiếc kẹo dâu và 4 chiếc kẹo cam. Nhắm mắt lấy ra 1 chiếc kẹo. Khả năng xảy ra là
\choice
{\True Có thể lấy được kẹo dâu hoặc kẹo cam}
{Chắc chắn lấy được kẹo dâu}
{Chắc chắn lấy được kẹo cam}
{Không thể lấy được kẹo dâu}
\loigiai{
Trong túi có cả kẹo dâu và kẹo cam nên có thể lấy được kẹo dâu hoặc kẹo cam.
\textbf{Đáp án của câu hỏi là A}
}
\end{ex}

\begin{ex} %[3X1V74]
Gieo một đồng xu có hai mặt: mặt Sấp và mặt Ngửa. Khả năng xuất hiện mặt Sấp là
\choice
{Chắc chắn xảy ra}
{Không thể xảy ra}
{\True Có thể xảy ra}
{Luôn luôn xảy ra}
\loigiai{
Đồng xu có hai khả năng Sấp hoặc Ngửa nên khả năng xuất hiện mặt Sấp là có thể xảy ra.
\textbf{Đáp án của câu hỏi là C}
}
\end{ex}

\begin{ex} %[3X1V74]
Trong một hộp có 4 chiếc bút mực màu đỏ và 4 chiếc bút mực màu xanh. Lấy ra cùng một lúc 5 chiếc bút. Sự kiện nào sau đây là 'Chắc chắn' xảy ra?
\choice
{Tất cả 5 chiếc đều màu đỏ}
{\True Có ít nhất 1 chiếc bút màu đỏ và 1 chiếc bút màu xanh}
{Tất cả 5 chiếc đều màu xanh}
{Không có chiếc bút màu xanh nào}
\loigiai{
Vì chỉ có 4 bút đỏ nên lấy 5 chiếc chắc chắn có ít nhất 1 bút xanh; và chỉ có 4 bút xanh nên chắc chắn có ít nhất 1 bút đỏ.
\textbf{Đáp án của câu hỏi là B}
}
\end{ex}

\begin{ex} %[3X1V74]
Ngày mai Mặt Trời mọc ở hướng nào là sự kiện 'Chắc chắn' xảy ra?
\choice
{Mặt Trời mọc ở hướng Tây}
{Mặt Trời mọc ở hướng Bắc}
{Mặt Trời mọc ở hướng Nam}
{\True Mặt Trời mọc ở hướng Đông}
\loigiai{
Theo quy luật tự nhiên, Mặt Trời chắc chắn mọc ở hướng Đông.
\textbf{Đáp án của câu hỏi là D}
}
\end{ex}

\begin{ex} %[3X1V74]
Trong một bình có 10 viên bi đỏ và 1 viên bi vàng. Bạn Nam lấy ngẫu nhiên 1 viên bi. Nhận xét nào sau đây là đúng nhất?
\choice
{\True Khả năng lấy được bi đỏ cao hơn khả năng lấy được bi vàng}
{Chắc chắn lấy được bi đỏ}
{Không thể lấy được bi vàng}
{Khả năng lấy được bi đỏ và bi vàng là như nhau}
\loigiai{
Vì số bi đỏ nhiều hơn số bi vàng nên khả năng lấy được bi đỏ cao hơn.
\textbf{Đáp án của câu hỏi là A}
}
\end{ex}

\begin{ex} %[3X1V74]
Hôm nay là thứ Sáu. Sự kiện 'Ngày mai là thứ Bảy' là sự kiện
\choice
{Có thể xảy ra}
{Không thể xảy ra}
{\True Chắc chắn xảy ra}
{Tùy thuộc thời tiết}
\loigiai{
Theo quy luật tuần hoàn các ngày trong tuần, sau thứ Sáu chắc chắn là thứ Bảy.
\textbf{Đáp án của câu hỏi là C}
}
\end{ex}

\begin{ex} %[3X1V74]
Một con cá vàng sống trong chậu nước. Sự kiện 'Con cá có thể trèo lên cây chuối' là
\choice
{Chắc chắn xảy ra}
{\True Không thể xảy ra}
{Có thể xảy ra}
{Chưa chắc chắn}
\loigiai{
Cá không thể trèo cây, đây là sự kiện không thể xảy ra.
\textbf{Đáp án của câu hỏi là B}
}
\end{ex}

\begin{ex} %[3X1V74]
Một hộp có 3 quả bóng xanh và 2 quả bóng đỏ. Nam lấy ra một số quả bóng. Để CHẮC CHẮN có ít nhất 1 quả bóng màu đỏ thì Nam phải lấy ra ít nhất bao nhiêu quả bóng?
\choice
{$2$ quả bóng}
{$3$ quả bóng}
{$5$ quả bóng}
{\True $4$ quả bóng}
\loigiai{
Trường hợp xấu nhất lấy phải 3 quả bóng xanh, cần lấy thêm 1 quả nữa: $3 + 1 = 4$ quả bóng.
\textbf{Đáp án của câu hỏi là D}
}
\end{ex}

\begin{ex} %[3X1V74]
Gieo đồng thời hai con xúc xắc. Tổng số chấm xuất hiện lớn nhất có thể là
\choice
{\True $12$}
{$6$}
{$10$}
{$14$}
\loigiai{
Số chấm lớn nhất ở mỗi xúc xắc là 6: $6 + 6 = 12$ chấm.
\textbf{Đáp án của câu hỏi là A}
}
\end{ex}

\begin{ex} %[3X1V74]
Bốc ngẫu nhiên một tấm thẻ trong 10 tấm thẻ ghi số từ 1 đến 10. Sự kiện 'Rút được tấm thẻ ghi số 11' là
\choice
{Chắc chắn xảy ra}
{Có thể xảy ra}
{\True Không thể xảy ra}
{Luôn xảy ra}
\loigiai{
Các thẻ chỉ ghi từ 1 đến 10 nên không thể rút được số 11.
\textbf{Đáp án của câu hỏi là C}
}
\end{ex}

\textbf{Phần 2. Trắc nghiệm đúng sai}
\setcounterpageref{ex}{0}

\begin{ex} %[3X1V74]
Xét tính đúng sai của các sự kiện ngẫu nhiên trong cuộc sống:
\choiceTF
{\True Một sự kiện chắc chắn xảy ra là sự kiện luôn luôn xảy ra trong mọi trường hợp.}
{Một sự kiện có thể xảy ra là sự kiện không bao giờ xảy ra.}
{\True Một sự kiện không thể xảy ra là sự kiện không bao giờ xảy ra trong bất kì trường hợp nào.}
{Gieo 1 con xúc xắc chắc chắn xuất hiện mặt 6 chấm.}
\loigiai{
\begin{itemize}
    \item Ý a) Đúng: Định nghĩa.
    \item Ý b) Sai: Sai vì sự kiện có thể xảy ra tức là có khả năng xảy ra.
    \item Ý c) Đúng: Định nghĩa.
    \item Ý d) Sai: Sai vì có thể xuất hiện các mặt từ 1 đến 5.
\end{itemize}
\begin{itemize}
    \item a) \textbf{Đúng}.
    \item b) \textbf{Sai}.
    \item c) \textbf{Đúng}.
    \item d) \textbf{Sai}.
\end{itemize}
}
\end{ex}

\begin{ex} %[3X1V74]
Trong một hộp kín có 8 viên bi trắng và 2 viên bi đen:
\choiceTF
{\True Lấy ngẫu nhiên 1 viên bi thì có thể lấy được bi trắng hoặc bi đen.}
{Chắc chắn lấy được viên bi màu đen khi chỉ lấy 1 viên.}
{\True Nếu lấy ra 9 viên bi thì chắc chắn có ít nhất 1 viên bi màu đen.}
{Không thể lấy được viên bi màu trắng.}
\loigiai{
\begin{itemize}
    \item Ý a) Đúng: Có cả hai loại bi.
    \item Ý b) Sai: Sai vì có thể lấy phải bi trắng.
    \item Ý c) Đúng: Vì bi trắng chỉ có 8 viên.
    \item Ý d) Sai: Sai vì trong hộp có 8 bi trắng.
\end{itemize}
\begin{itemize}
    \item a) \textbf{Đúng}.
    \item b) \textbf{Sai}.
    \item c) \textbf{Đúng}.
    \item d) \textbf{Sai}.
\end{itemize}
}
\end{ex}

\textbf{Phần 3. Trả lời ngắn}
\setcounterpageref{ex}{0}

\begin{ex} %[3X1V74]
Trong hộp có 5 viên bi đỏ và 3 viên bi vàng. Cần lấy ra ít nhất bao nhiêu viên bi để chắc chắn có ít nhất 1 viên bi màu vàng?
\shortans{6}
\loigiai{
\textbf{Đáp số:} $6$\\
Trường hợp xấu nhất lấy 5 bi đỏ, lấy thêm 1 viên nữa: $5 + 1 = 6$ viên.
}
\end{ex}

\begin{ex} %[3X1V74]
Gieo một con xúc xắc. Tổng số mặt có số chấm là số chẵn là bao nhiêu mặt?
\shortans{3}
\loigiai{
\textbf{Đáp số:} $3$\\
Các mặt chẵn: 2, 4, 6 (có 3 mặt).
}
\end{ex}

\begin{ex} %[3X1V74]
Trong túi có 10 viên kẹo sô-cô-la và không có loại kẹo nào khác. Lấy ra 1 chiếc kẹo. Có bao nhiêu khả năng về màu sắc/loại kẹo có thể nhận được?
\shortans{1}
\loigiai{
\textbf{Đáp số:} $1$\\
Chỉ có duy nhất 1 loại kẹo sô-cô-la.
}
\end{ex}

\begin{ex} %[3X1V74]
Gieo 1 con xúc xắc 6 mặt. Số chấm nhỏ nhất có thể xuất hiện là bao nhiêu?
\shortans{1}
\loigiai{
\textbf{Đáp số:} $1$\\
Mặt 1 chấm là nhỏ nhất.
}
\end{ex}

\begin{ex} %[3X1V74]
Có 4 quả táo và 5 quả lê trong giỏ. Hỏi cần lấy ra ít nhất bao nhiêu quả để chắc chắn có ít nhất 1 quả táo?
\shortans{6}
\loigiai{
\textbf{Đáp số:} $6$\\
Trường hợp xấu nhất lấy cả 5 quả lê, thêm 1 quả nữa: $5 + 1 = 6$ quả.
}
\end{ex}

\begin{ex} %[3X1V74]
Một tuần lễ có bao nhiêu ngày chắc chắn có tên bắt đầu bằng chữ 'Thứ'?
\shortans{6}
\loigiai{
\textbf{Đáp số:} $6$\\
Thứ Hai, Thứ Ba, Thứ Tư, Thứ Năm, Thứ Sáu, Thứ Bảy (6 ngày, trừ Chủ nhật).
}
\end{ex}

\newpage

\TaoPhieuKiemTra{Chủ đề 15 -- Làm quen với yếu tố thống kê, xác suất}{75}{Thực hành trải nghiệm thống kê}{https://loptoanthaylinh.vn/thi-online/toan-3/bai-75}
\textbf{Phần 1. Trắc nghiệm lựa chọn}
\setcounterpageref{ex}{0}

\begin{ex} %[3X1V75]
Để ghi chép số lượng xe đi qua cổng trường trong 15 phút, bạn An dùng các vạch kiểm đếm (mỗi nhóm gồm 4 vạch dọc và 1 vạch chéo ngang biểu thị 5 xe). An đếm được 6 nhóm vạch và 3 vạch rời. Số xe An đếm được là
\choice
{\True $33$ xe}
{$30$ xe}
{$35$ xe}
{$28$ xe}
\loigiai{
Số xe là: $5 \times 6 + 3 = 33$ xe.
\textbf{Đáp án của câu hỏi là A}
}
\end{ex}

\begin{ex} %[3X1V75]
Bảng thống kê thời gian tự học ở nhà mỗi ngày của 4 bạn: An: 60 phút; Bình: 45 phút; Cúc: 90 phút; Dũng: 75 phút. Bạn dành nhiều thời gian tự học nhất là
\choice
{An}
{Dũng}
{\True Cúc}
{Bình}
\loigiai{
Cúc học 90 phút, nhiều nhất.
\textbf{Đáp án của câu hỏi là C}
}
\end{ex}

\begin{ex} %[3X1V75]
Thời gian tự học của Cúc nhiều hơn Bình bao nhiêu phút?
\choice
{$30$ phút}
{\True $45$ phút}
{$15$ phút}
{$60$ phút}
\loigiai{
$90 - 45 = 45$ phút.
\textbf{Đáp án của câu hỏi là B}
}
\end{ex}

\begin{ex} %[3X1V75]
Thống kê chiều cao 4 bạn học sinh: Nam cao $132\text{ cm}$, Hoa cao $128\text{ cm}$, Tuấn cao $135\text{ cm}$, Lan cao $130\text{ cm}$. Bạn thấp nhất là
\choice
{Lan}
{Nam}
{Tuấn}
{\True Hoa}
\loigiai{
Hoa cao $128\text{ cm}$, thấp nhất.
\textbf{Đáp án của câu hỏi là D}
}
\end{ex}

\begin{ex} %[3X1V75]
Chênh lệch chiều cao giữa bạn cao nhất và bạn thấp nhất trong nhóm trên là
\choice
{\True $7\text{ cm}$}
{$5\text{ cm}$}
{$6\text{ cm}$}
{$8\text{ cm}$}
\loigiai{
Tuấn cao nhất ($135\text{ cm}$), Hoa thấp nhất ($128\text{ cm}$). Chênh lệch: $135 - 128 = 7\text{ (cm)}$.
\textbf{Đáp án của câu hỏi là A}
}
\end{ex}

\begin{ex} %[3X1V75]
Số học sinh tham gia câu lạc bộ bơi của các lớp khối 3: 3A có 14 bạn; 3B có 16 bạn; 3C có 15 bạn; 3D có 15 bạn. Tổng số học sinh tham gia của cả khối là
\choice
{$55$ bạn}
{$65$ bạn}
{\True $60$ bạn}
{$58$ bạn}
\loigiai{
$14 + 16 + 15 + 15 = 60$ bạn.
\textbf{Đáp án của câu hỏi là C}
}
\end{ex}

\begin{ex} %[3X1V75]
Một bạn gieo đồng xu 20 lần và ghi lại: có 12 lần xuất hiện mặt Sấp. Số lần xuất hiện mặt Ngửa là
\choice
{$10$ lần}
{\True $8$ lần}
{$12$ lần}
{$6$ lần}
\loigiai{
Số lần mặt Ngửa: $20 - 12 = 8$ lần.
\textbf{Đáp án của câu hỏi là B}
}
\end{ex}

\begin{ex} %[3X1V75]
Ghi chép số quả trứng gà thu hoạch trong 5 ngày: Ngày 1: 45 quả; Ngày 2: 50 quả; Ngày 3: 48 quả; Ngày 4: 52 quả; Ngày 5: 55 quả. Ngày thu hoạch được nhiều trứng nhất là
\choice
{Ngày 4}
{Ngày 2}
{Ngày 3}
{\True Ngày 5}
\loigiai{
Ngày 5 thu hoạch được 55 quả, nhiều nhất.
\textbf{Đáp án của câu hỏi là D}
}
\end{ex}

\begin{ex} %[3X1V75]
Tổng số trứng gà thu hoạch được trong 5 ngày đó là
\choice
{\True $250$ quả}
{$240$ quả}
{$260$ quả}
{$245$ quả}
\loigiai{
Tổng số trứng: $45 + 50 + 48 + 52 + 55 = 250$ quả.
\textbf{Đáp án của câu hỏi là A}
}
\end{ex}

\begin{ex} %[3X1V75]
Khi thực hiện một cuộc phỏng vấn nhanh về môn thể thao yêu thích, công cụ ghi chép nhanh và hiệu quả nhất tại chỗ là
\choice
{Thước kẻ}
{Vẽ tranh}
{\True Bảng kiểm đếm bằng các vạch gạch}
{Máy tính để bàn cồng kềnh}
\loigiai{
Ghi chép bằng vạch kiểm đếm là phương pháp chuẩn xác và nhanh nhất.
\textbf{Đáp án của câu hỏi là C}
}
\end{ex}

\begin{ex} %[3X1V75]
Kết quả bình bầu cán bộ lớp: Mai được 28 phiếu, Lan được 24 phiếu, Cường được 20 phiếu. Số phiếu của Mai nhiều hơn Cường bao nhiêu phiếu?
\choice
{$4$ phiếu}
{\True $8$ phiếu}
{$6$ phiếu}
{$10$ phiếu}
\loigiai{
$28 - 20 = 8$ phiếu.
\textbf{Đáp án của câu hỏi là B}
}
\end{ex}

\begin{ex} %[3X1V75]
Thống kê số lượng cây hoa trồng trong bồn hoa lớp: 15 cây hoa hồng, 20 cây hoa cúc, 10 cây hoa đồng tiền. Tổng số cây hoa đã trồng là
\choice
{$40$ cây}
{$50$ cây}
{$35$ cây}
{\True $45$ cây}
\loigiai{
$15 + 20 + 10 = 45$ cây.
\textbf{Đáp án của câu hỏi là D}
}
\end{ex}

\textbf{Phần 2. Trắc nghiệm đúng sai}
\setcounterpageref{ex}{0}

\begin{ex} %[3X1V75]
Xét tính đúng sai của việc thực hành trải nghiệm thống kê:
\choiceTF
{\True Kiểm đếm bằng vạch giúp đếm số lượng nhanh và tránh nhầm lẫn.}
{Chỉ có thể thu thập số liệu bằng cách cân đo, không thể đếm.}
{\True Bảng số liệu thống kê phản ánh trung thực kết quả thu thập.}
{Kết quả thực nghiệm xác suất sau 10 lần gieo đồng xu luôn luôn có 5 lần sấp và 5 lần ngửa.}
\loigiai{
\begin{itemize}
    \item Ý a) Đúng: Lợi ích kiểm đếm.
    \item Ý b) Sai: Sai vì đếm số lượng là phương pháp thống kê cơ bản.
    \item Ý c) Đúng: Yêu cầu khoa học.
    \item Ý d) Sai: Sai vì kết quả ngẫu nhiên có thể dao động.
\end{itemize}
\begin{itemize}
    \item a) \textbf{Đúng}.
    \item b) \textbf{Sai}.
    \item c) \textbf{Đúng}.
    \item d) \textbf{Sai}.
\end{itemize}
}
\end{ex}

\begin{ex} %[3X1V75]
Bảng thống kê số tiết đọc sách thư viện của 3 bạn trong tháng: An đọc 12 tiết, Bình đọc 16 tiết, Cúc đọc 20 tiết:
\choiceTF
{\True Cúc đọc nhiều tiết sách nhất trong ba bạn.}
{\True Cả ba bạn đọc tất cả 48 tiết sách.}
{Bình đọc gấp đôi số tiết của An.}
{\True Số tiết đọc của Cúc nhiều hơn An là 8 tiết.}
\loigiai{
\begin{itemize}
    \item Ý a) Đúng: 20 tiết nhiều nhất.
    \item Ý b) Đúng: $12 + 16 + 20 = 48$ tiết.
    \item Ý c) Sai: Sai vì $12 \times 2 = 24 \neq 16$.
    \item Ý d) Đúng: $20 - 12 = 8$ tiết.
\end{itemize}
\begin{itemize}
    \item a) \textbf{Đúng}.
    \item b) \textbf{Đúng}.
    \item c) \textbf{Sai}.
    \item d) \textbf{Đúng}.
\end{itemize}
}
\end{ex}

\textbf{Phần 3. Trả lời ngắn}
\setcounterpageref{ex}{0}

\begin{ex} %[3X1V75]
Kiểm đếm số học sinh thích màu đỏ: có 7 nhóm vạch (mỗi nhóm 5 vạch) và 2 vạch rời. Tính tổng số học sinh thích màu đỏ.
\shortans{37}
\loigiai{
\textbf{Đáp số:} $37$\\
$5 \times 7 + 2 = 37$ bạn.
}
\end{ex}

\begin{ex} %[3X1V75]
Khối 3 có 4 lớp, mỗi lớp quyên góp được 25 quyển truyện tranh. Cả khối 3 quyên góp được bao nhiêu quyển truyện tranh?
\shortans{100}
\loigiai{
\textbf{Đáp số:} $100$\\
$25 \times 4 = 100$ quyển.
}
\end{ex}

\begin{ex} %[3X1V75]
Thực nghiệm gieo xúc xắc 30 lần, xuất hiện mặt 6 chấm 8 lần. Hỏi các mặt còn lại xuất hiện bao nhiêu lần?
\shortans{22}
\loigiai{
\textbf{Đáp số:} $22$\\
$30 - 8 = 22$ lần.
}
\end{ex}

\begin{ex} %[3X1V75]
Bạn Lan đo nhiệt độ trong 3 ngày: Ngày 1: $28^\circ\text{C}$; Ngày 2: $30^\circ\text{C}$; Ngày 3: $32^\circ\text{C}$. Hỏi nhiệt độ ngày 3 hơn ngày 1 bao nhiêu độ C?
\shortans{4}
\loigiai{
\textbf{Đáp số:} $4$\\
$32 - 28 = 4^\circ\text{C}$.
}
\end{ex}

\begin{ex} %[3X1V75]
Số điểm thưởng của 3 tổ: Tổ 1 được 40 điểm, Tổ 2 được 50 điểm, Tổ 3 được 60 điểm. Tính tổng số điểm thưởng của cả 3 tổ.
\shortans{150}
\loigiai{
\textbf{Đáp số:} $150$\\
$40 + 50 + 60 = 150$ điểm.
}
\end{ex}

\begin{ex} %[3X1V75]
Trong lớp có 36 học sinh, có 24 bạn tham gia học bơi. Hỏi có bao nhiêu bạn chưa tham gia học bơi?
\shortans{12}
\loigiai{
\textbf{Đáp số:} $12$\\
$36 - 24 = 12$ bạn.
}
\end{ex}

\newpage

% =========================================================================
% CHỦ ĐỀ 16 - ÔN TẬP CUỐI NĂM
% =========================================================================
\section{Ôn tập cuối năm}

\TaoPhieuKiemTra{Chủ đề 16 -- Ôn tập cuối năm}{76}{Ôn tập các số trong phạm vi 100 000}{https://loptoanthaylinh.vn/thi-online/toan-3/bai-76}
\textbf{Phần 1. Trắc nghiệm lựa chọn}
\setcounterpageref{ex}{0}

\begin{ex} %[3D1V76]
Số liền trước của số bé nhất có năm chữ số là
\choice
{$10\,001$}
{$10\,000$}
{$9990$}
{\True $9999$}
\loigiai{
Số bé nhất có 5 chữ số là $10\,000$. Số liền trước là: $10\,000 - 1 = 9999$.
\textbf{Đáp án của câu hỏi là D}
}
\end{ex}

\begin{ex} %[3D1V76]
Số gồm 8 chục nghìn, 5 trăm và 4 đơn vị được viết là
\choice
{$85\,004$}
{\True $80\,504$}
{$80\,540$}
{$85\,400$}
\loigiai{
Hàng nghìn và hàng chục là 0: $80\,504$.
\textbf{Đáp án của câu hỏi là B}
}
\end{ex}

\begin{ex} %[3D1V76]
Làm tròn số $64\,850$ đến hàng nghìn ta được số
\choice
{$64\,000$}
{$64\,900$}
{\True $65\,000$}
{$60\,000$}
\loigiai{
Hàng trăm là 8 nên làm tròn lên $65\,000$.
\textbf{Đáp án của câu hỏi là C}
}
\end{ex}

\begin{ex} %[3D1V76]
Số La Mã $\text{XIX}$ biểu thị số tự nhiên nào?
\choice
{\True $19$}
{$21$}
{$18$}
{$14$}
\loigiai{
$\text{XIX} = 19$.
\textbf{Đáp án của câu hỏi là A}
}
\end{ex}

\begin{ex} %[3D1V76]
Số chẵn bé nhất có năm chữ số khác nhau là
\choice
{$10\,246$}
{$10\,000$}
{$12\,346$}
{\True $10\,234$}
\loigiai{
Số bé nhất có 5 chữ số khác nhau là $10\,234$ (tận cùng 4 là số chẵn).
\textbf{Đáp án của câu hỏi là D}
}
\end{ex}

\begin{ex} %[3D1V76]
Số $100\,000$ gấp số $1000$ bao nhiêu lần?
\choice
{$10$ lần}
{\True $100$ lần}
{$1000$ lần}
{$10\,000$ lần}
\loigiai{
$100\,000 : 1000 = 100$ (lần).
\textbf{Đáp án của câu hỏi là B}
}
\end{ex}

\begin{ex} %[3D1V76]
Từ năm chữ số $0; 1; 3; 6; 8$, viết được số có năm chữ số khác nhau lớn nhất là
\choice
{$86\,301$}
{$86\,130$}
{\True $86\,310$}
{$83\,610$}
\loigiai{
Thứ tự giảm dần: $86\,310$.
\textbf{Đáp án của câu hỏi là C}
}
\end{ex}

\begin{ex} %[3D1V76]
Số liền sau của số lớn nhất có năm chữ số là
\choice
{\True $100\,000$}
{$99\,998$}
{$100\,001$}
{$99\,990$}
\loigiai{
Số lớn nhất có 5 chữ số là 99 999. Số liền sau là: $99\,999 + 1 = 100\,000$.
\textbf{Đáp án của câu hỏi là A}
}
\end{ex}

\begin{ex} %[3D1V76]
Làm tròn số $75\,000$ đến hàng chục nghìn ta được số
\choice
{$70\,000$}
{$75\,000$}
{$90\,000$}
{\True $80\,000$}
\loigiai{
Hàng nghìn là 5 nên làm tròn lên $80\,000$.
\textbf{Đáp án của câu hỏi là D}
}
\end{ex}

\begin{ex} %[3D1V76]
Trong số $78\,942$, chữ số 7 có giá trị là
\choice
{$7000$}
{\True $70\,000$}
{$700$}
{$700\,000$}
\loigiai{
Chữ số 7 ở hàng chục nghìn có giá trị là $70\,000$.
\textbf{Đáp án của câu hỏi là B}
}
\end{ex}

\begin{ex} %[3D1V76]
Số nào sau đây bé hơn $54\,320$?
\choice
{$54\,325$}
{$54\,400$}
{\True $54\,302$}
{$55\,000$}
\loigiai{
$54\,302 < 54\,320$.
\textbf{Đáp án của câu hỏi là C}
}
\end{ex}

\begin{ex} %[3D1V76]
Dãy số nào sau đây được xếp theo thứ tự giảm dần?
\choice
{\True $80\,000; 75\,000; 72\,000; 70\,000$}
{$70\,000; 72\,000; 75\,000; 80\,000$}
{$80\,000; 72\,000; 75\,000; 70\,000$}
{$75\,000; 80\,000; 72\,000; 70\,000$}
\loigiai{
Thứ tự giảm dần chuẩn: $80\,000 > 75\,000 > 72\,000 > 70\,000$.
\textbf{Đáp án của câu hỏi là A}
}
\end{ex}

\textbf{Phần 2. Trắc nghiệm đúng sai}
\setcounterpageref{ex}{0}

\begin{ex} %[3D1V76]
Xét tính đúng sai của các nhận định ôn tập số học:
\choiceTF
{\True Số 100 000 gồm 10 chục nghìn.}
{Chữ số La Mã XXI có giá trị là 19.}
{\True Khi làm tròn số 62 480 đến hàng nghìn thì được 62 000.}
{Số 99 999 là số chẵn.}
\loigiai{
\begin{itemize}
    \item Ý a) Đúng: Cấu tạo số chuẩn.
    \item Ý b) Sai: Sai vì XXI là 21.
    \item Ý c) Đúng: Vì hàng trăm là 4.
    \item Ý d) Sai: Sai vì tận cùng là 9 nên là số lẻ.
\end{itemize}
\begin{itemize}
    \item a) \textbf{Đúng}.
    \item b) \textbf{Sai}.
    \item c) \textbf{Đúng}.
    \item d) \textbf{Sai}.
\end{itemize}
}
\end{ex}

\begin{ex} %[3D1V76]
Cho số $85\,460$:
\choiceTF
{\True Làm tròn đến hàng trăm được số 85 500.}
{\True Làm tròn đến hàng nghìn được số 85 000.}
{\True Làm tròn đến hàng chục nghìn được số 90 000.}
{Số liền sau của số đó là 85 459.}
\loigiai{
\begin{itemize}
    \item Ý a) Đúng: Đơn vị là 6 làm tròn lên.
    \item Ý b) Đúng: Trăm là 4 làm tròn xuống.
    \item Ý c) Đúng: Nghìn là 5 làm tròn lên.
    \item Ý d) Sai: Sai vì số liền sau là 85 461.
\end{itemize}
\begin{itemize}
    \item a) \textbf{Đúng}.
    \item b) \textbf{Đúng}.
    \item c) \textbf{Đúng}.
    \item d) \textbf{Sai}.
\end{itemize}
}
\end{ex}

\textbf{Phần 3. Trả lời ngắn}
\setcounterpageref{ex}{0}

\begin{ex} %[3D1V76]
Tìm số liền sau của số 99 999.
\shortans{100000}
\loigiai{
\textbf{Đáp số:} $100000$\\
$99\,999 + 1 = 100\,000$.
}
\end{ex}

\begin{ex} %[3D1V76]
Viết số tự nhiên tương ứng với chữ số La Mã XVIII.
\shortans{18}
\loigiai{
\textbf{Đáp số:} $18$\\
XVIII là 18.
}
\end{ex}

\begin{ex} %[3D1V76]
Làm tròn số 76 540 đến hàng nghìn.
\shortans{77000}
\loigiai{
\textbf{Đáp số:} $77000$\\
Hàng trăm là 5 nên làm tròn thành 77 000.
}
\end{ex}

\begin{ex} %[3D1V76]
Viết số lớn nhất có năm chữ số khác nhau.
\shortans{98765}
\loigiai{
\textbf{Đáp số:} $98765$\\
Số đó là 98 765.
}
\end{ex}

\begin{ex} %[3D1V76]
Tính tổng các chữ số của số 45 678.
\shortans{30}
\loigiai{
\textbf{Đáp số:} $30$\\
$4 + 5 + 6 + 7 + 8 = 30$.
}
\end{ex}

\begin{ex} %[3D1V76]
Có bao nhiêu số tròn nghìn từ 30 000 đến 40 000?
\shortans{11}
\loigiai{
\textbf{Đáp số:} $11$\\
$(40\,000 - 30\,000) : 1000 + 1 = 11$ số.
}
\end{ex}

\newpage

\TaoPhieuKiemTra{Chủ đề 16 -- Ôn tập cuối năm}{77}{Ôn tập phép cộng, phép trừ trong phạm vi 100 000}{https://loptoanthaylinh.vn/thi-online/toan-3/bai-77}
\textbf{Phần 1. Trắc nghiệm lựa chọn}
\setcounterpageref{ex}{0}

\begin{ex} %[3D2V77]
Tính giá trị phép tính: $45\,280 + 38\,450$.
\choice
{$83\,630$}
{$83\,830$}
{\True $83\,730$}
{$82\,730$}
\loigiai{
$45\,280 + 38\,450 = 83\,730$.
\textbf{Đáp án của câu hỏi là C}
}
\end{ex}

\begin{ex} %[3D2V77]
Tính giá trị phép tính: $92\,400 - 45\,800$.
\choice
{\True $46\,600$}
{$47\,600$}
{$46\,400$}
{$56\,600$}
\loigiai{
$92\,400 - 45\,800 = 46\,600$.
\textbf{Đáp án của câu hỏi là A}
}
\end{ex}

\begin{ex} %[3D2V77]
Tìm số tự nhiên $x$, biết: $x + 28\,500 = 75\,000$.
\choice
{$47\,500$}
{$46\,000$}
{$45\,500$}
{\True $46\,500$}
\loigiai{
$x = 75\,000 - 28\,500 = 46\,500$.
\textbf{Đáp án của câu hỏi là D}
}
\end{ex}

\begin{ex} %[3D2V77]
Tính nhanh giá trị: $34\,500 + 28\,000 + 15\,500 + 12\,000$.
\choice
{$88\,000$}
{\True $90\,000$}
{$92\,000$}
{$89\,000$}
\loigiai{
$(34\,500 + 15\,500) + (28\,000 + 12\,000) = 50\,000 + 40\,000 = 90\,000$.
\textbf{Đáp án của câu hỏi là B}
}
\end{ex}

\begin{ex} %[3D2V77]
Một kho gạo có $85\,000\text{ kg}$ gạo. Người ta xuất đi hai đợt: đợt một $28\,000\text{ kg}$, đợt hai $32\,500\text{ kg}$. Trong kho còn lại bao nhiêu ki-lô-gam gạo?
\choice
{$25\,500\text{ kg}$}
{$24\,000\text{ kg}$}
{\True $24\,500\text{ kg}$}
{$34\,500\text{ kg}$}
\loigiai{
Tổng xuất: $28\,000 + 32\,500 = 60\,500\text{ kg}$. Còn lại: $85\,000 - 60\,500 = 24\,500\text{ (kg)}$.
\textbf{Đáp án của câu hỏi là C}
}
\end{ex}

\begin{ex} %[3D2V77]
Một công ty quý 1 thu được $45\,000\,000$ đồng, quý 2 thu được nhiều hơn quý 1 là $12\,000\,000$ đồng. Cả hai quý công ty thu được bao nhiêu tiền?
\choice
{\True $102\,000\,000$ đồng}
{$90\,000\,000$ đồng}
{$100\,000\,000$ đồng}
{$95\,000\,000$ đồng}
\loigiai{
Quý 2: $45 + 12 = 57$ triệu. Cả hai quý: $45 + 57 = 102$ triệu đồng.
\textbf{Đáp án của câu hỏi là A}
}
\end{ex}

\begin{ex} %[3D2V77]
Tìm hiệu của số lớn nhất có 5 chữ số và số bé nhất có 5 chữ số.
\choice
{$90\,000$}
{$99\,999$}
{$88\,888$}
{\True $89\,999$}
\loigiai{
$99\,999 - 10\,000 = 89\,999$.
\textbf{Đáp án của câu hỏi là D}
}
\end{ex}

\begin{ex} %[3D2V77]
Tính giá trị biểu thức: $100\,000 - (35\,000 + 45\,000)$.
\choice
{$30\,000$}
{\True $20\,000$}
{$25\,000$}
{$15\,000$}
\loigiai{
$100\,000 - 80\,000 = 20\,000$.
\textbf{Đáp án của câu hỏi là B}
}
\end{ex}

\begin{ex} %[3D2V77]
Tổng hai số là 84 000, hiệu hai số là 16 000. Số lớn là
\choice
{$42\,000$}
{$34\,000$}
{\True $50\,000$}
{$48\,000$}
\loigiai{
Số lớn: $(84\,000 + 16\,000) : 2 = 50\,000$.
\textbf{Đáp án của câu hỏi là C}
}
\end{ex}

\begin{ex} %[3D2V77]
Một mảnh vườn hình chữ nhật có chu vi $600\text{ m}$, chiều dài hơn chiều rộng $60\text{ m}$. Chiều dài của mảnh vườn là
\choice
{\True $180\text{ m}$}
{$120\text{ m}$}
{$150\text{ m}$}
{$200\text{ m}$}
\loigiai{
Nửa chu vi: $600 : 2 = 300\text{ m}$. Chiều dài: $(300 + 60) : 2 = 180\text{ (m)}$.
\textbf{Đáp án của câu hỏi là A}
}
\end{ex}

\begin{ex} %[3D2V77]
Nếu số bị trừ tăng 5000 đơn vị và số trừ giữ nguyên thì hiệu sẽ
\choice
{Giảm đi 5000 đơn vị}
{Không đổi}
{Tăng thêm 10 000 đơn vị}
{\True Tăng thêm 5000 đơn vị}
\loigiai{
Số bị trừ tăng bao nhiêu thì hiệu tăng bấy nhiêu.
\textbf{Đáp án của câu hỏi là D}
}
\end{ex}

\begin{ex} %[3D2V77]
Một cuộn dây cáp dài $50\,000\text{ m}$, đã dùng hết $32\,400\text{ m}$. Cuộn dây cáp còn lại dài bao nhiêu mét?
\choice
{$18\,600\text{ m}$}
{\True $17\,600\text{ m}$}
{$17\,400\text{ m}$}
{$16\,600\text{ m}$}
\loigiai{
$50\,000 - 32\,400 = 17\,600\text{ (m)}$.
\textbf{Đáp án của câu hỏi là B}
}
\end{ex}

\textbf{Phần 2. Trắc nghiệm đúng sai}
\setcounterpageref{ex}{0}

\begin{ex} %[3D2V77]
Xét tính đúng sai của các phép toán cộng trừ phạm vi 100 000:
\choiceTF
{\True Phép cộng và trừ được thực hiện lần lượt từ phải sang trái.}
{$60\,000 - 25\,000 = 45\,000$.}
{\True Khi cùng thêm vào cả hai số hạng một số đơn vị thì tổng tăng thêm gấp đôi số đơn vị đó.}
{$75\,000 + 25\,000 = 90\,000$.}
\loigiai{
\begin{itemize}
    \item Ý a) Đúng: Quy tắc cơ bản.
    \item Ý b) Sai: Sai vì $60\,000 - 25\,000 = 35\,000$.
    \item Ý c) Đúng: Tính chất tổng.
    \item Ý d) Sai: Sai vì $75\,000 + 25\,000 = 100\,000$.
\end{itemize}
\begin{itemize}
    \item a) \textbf{Đúng}.
    \item b) \textbf{Sai}.
    \item c) \textbf{Đúng}.
    \item d) \textbf{Sai}.
\end{itemize}
}
\end{ex}

\begin{ex} %[3D2V77]
Một kho dự trữ có 90 000 bao xi măng:
\choiceTF
{\True Xuất đi 45 000 bao thì còn lại đúng một nửa.}
{\True Nếu nhập thêm 10 000 bao nữa thì kho có đúng 100 000 bao xi măng.}
{Xuất đi 60 000 bao thì số bao còn lại nhiều hơn số bao đã xuất.}
{\True Tổng số bao nhập và xuất trong ngày có thể tính bằng phép cộng.}
\loigiai{
\begin{itemize}
    \item Ý a) Đúng: $45\,000 : 90\,000 = \dfrac{1}{2}$.
    \item Ý b) Đúng: $90\,000 + 10\,000 = 100\,000$ bao.
    \item Ý c) Sai: Sai vì còn 30 000 bao, ít hơn 60 000 bao.
    \item Ý d) Đúng: Đúng.
\end{itemize}
\begin{itemize}
    \item a) \textbf{Đúng}.
    \item b) \textbf{Đúng}.
    \item c) \textbf{Sai}.
    \item d) \textbf{Đúng}.
\end{itemize}
}
\end{ex}

\textbf{Phần 3. Trả lời ngắn}
\setcounterpageref{ex}{0}

\begin{ex} %[3D2V77]
Tính: $64\,500 + 28\,700$.
\shortans{93200}
\loigiai{
\textbf{Đáp số:} $93200$\\
$64\,500 + 28\,700 = 93\,200$.
}
\end{ex}

\begin{ex} %[3D2V77]
Tính: $82\,000 - 45\,600$.
\shortans{36400}
\loigiai{
\textbf{Đáp số:} $36400$\\
$82\,000 - 45\,600 = 36\,400$.
}
\end{ex}

\begin{ex} %[3D2V77]
Tìm $x$, biết: $x - 24\,000 = 56\,000$.
\shortans{80000}
\loigiai{
\textbf{Đáp số:} $80000$\\
$x = 56\,000 + 24\,000 = 80\,000$.
}
\end{ex}

\begin{ex} %[3D2V77]
Tính nhanh: $45\,000 + 18\,500 + 15\,000 + 11\,500$.
\shortans{90000}
\loigiai{
\textbf{Đáp số:} $90000$\\
$60\,000 + 30\,000 = 90\,000$.
}
\end{ex}

\begin{ex} %[3D2V77]
Tìm hiệu giữa 100 000 và 72 500.
\shortans{27500}
\loigiai{
\textbf{Đáp số:} $27500$\\
$100\,000 - 72\,500 = 27\,500$.
}
\end{ex}

\begin{ex} %[3D2V77]
Tổng hai số là 70 000, số bé là 30 000. Tìm số lớn.
\shortans{40000}
\loigiai{
\textbf{Đáp số:} $40000$\\
Số lớn là: $70\,000 - 30\,000 = 40\,000$.
}
\end{ex}

\newpage

\TaoPhieuKiemTra{Chủ đề 16 -- Ôn tập cuối năm}{78}{Ôn tập phép nhân, phép chia trong phạm vi 100 000}{https://loptoanthaylinh.vn/thi-online/toan-3/bai-78}
\textbf{Phần 1. Trắc nghiệm lựa chọn}
\setcounterpageref{ex}{0}

\begin{ex} %[3D2V78]
Tính giá trị của phép tính: $16\,250 \times 4$.
\choice
{$64\,000$}
{\True $65\,000$}
{$66\,000$}
{$65\,500$}
\loigiai{
$16\,250 \times 4 = 65\,000$.
\textbf{Đáp án của câu hỏi là B}
}
\end{ex}

\begin{ex} %[3D2V78]
Tính giá trị của phép chia: $84\,648 : 4$.
\choice
{$21\,160$}
{$21\,262$}
{$21\,062$}
{\True $21\,162$}
\loigiai{
$84\,648 : 4 = 21\,162$.
\textbf{Đáp án của câu hỏi là D}
}
\end{ex}

\begin{ex} %[3D2V78]
Tìm số tự nhiên $y$, biết: $y : 5 = 18\,000$.
\choice
{\True $90\,000$}
{$85\,000$}
{$95\,000$}
{$80\,000$}
\loigiai{
$y = 18\,000 \times 5 = 90\,000$.
\textbf{Đáp án của câu hỏi là A}
}
\end{ex}

\begin{ex} %[3D2V78]
Tính giá trị của biểu thức: $A = 14\,000 \times 3 + 36\,000 : 6$.
\choice
{$45\,000$}
{$50\,000$}
{\True $48\,000$}
{$46\,000$}
\loigiai{
$14\,000 \times 3 = 42\,000$; $36\,000 : 6 = 6000$; $42\,000 + 6000 = 48\,000$.
\textbf{Đáp án của câu hỏi là C}
}
\end{ex}

\begin{ex} %[3D2V78]
Một hình vuông có chu vi $68\,400\text{ m}$. Độ dài mỗi cạnh hình vuông là
\choice
{$17\,000\text{ m}$}
{\True $17\,100\text{ m}$}
{$17\,200\text{ m}$}
{$16\,100\text{ m}$}
\loigiai{
$68\,400 : 4 = 17\,100\text{ (m)}$.
\textbf{Đáp án của câu hỏi là B}
}
\end{ex}

\begin{ex} %[3D2V78]
Tính nhanh giá trị: $20\,000 \times 3 + 20\,000 \times 2$.
\choice
{$80\,000$}
{$90\,000$}
{$120\,000$}
{\True $100\,000$}
\loigiai{
$20\,000 \times (3 + 2) = 20\,000 \times 5 = 100\,000$.
\textbf{Đáp án của câu hỏi là D}
}
\end{ex}

\begin{ex} %[3D2V78]
Trong phép chia cho 9 có số dư lớn nhất là 8 và thương là $10\,500$. Số bị chia là
\choice
{\True $94\,508$}
{$94\,500$}
{$94\,507$}
{$94\,509$}
\loigiai{
Số bị chia: $9 \times 10\,500 + 8 = 94\,500 + 8 = 94\,508$.
\textbf{Đáp án của câu hỏi là A}
}
\end{ex}

\begin{ex} %[3D2V78]
Một kho gạo có $96\,000\text{ kg}$ gạo. Đã bán đi $\dfrac{1}{4}$ số gạo. Số gạo còn lại là
\choice
{$24\,000\text{ kg}$}
{$70\,000\text{ kg}$}
{\True $72\,000\text{ kg}$}
{$75\,000\text{ kg}$}
\loigiai{
Đã bán: $96\,000 : 4 = 24\,000\text{ kg}$. Còn lại: $96\,000 - 24\,000 = 72\,000\text{ (kg)}$.
\textbf{Đáp án của câu hỏi là C}
}
\end{ex}

\begin{ex} %[3D2V78]
Tìm số tự nhiên $x$, biết: $x \times 4 + 15\,000 = 75\,000$.
\choice
{$12\,000$}
{\True $15\,000$}
{$16\,000$}
{$18\,000$}
\loigiai{
$x \times 4 = 75\,000 - 15\,000 = 60\,000 \Rightarrow x = 60\,000 : 4 = 15\,000$.
\textbf{Đáp án của câu hỏi là B}
}
\end{ex}

\begin{ex} %[3D2V78]
Một mảnh đất hình chữ nhật có chiều rộng $12\,000\text{ m}$, chiều dài gấp đôi chiều rộng. Chu vi mảnh đất đó là
\choice
{$48\,000\text{ m}$}
{$60\,000\text{ m}$}
{$84\,000\text{ m}$}
{\True $72\,000\text{ m}$}
\loigiai{
Chiều dài: $24\,000\text{ m}$. Chu vi: $(24\,000 + 12\,000) \times 2 = 72\,000\text{ (m)}$.
\textbf{Đáp án của câu hỏi là D}
}
\end{ex}

\begin{ex} %[3D2V78]
Thực hiện phép chia: $100\,000 : 4$. Kết quả là
\choice
{\True $25\,000$}
{$20\,000$}
{$30\,000$}
{$22\,500$}
\loigiai{
$100\,000 : 4 = 25\,000$.
\textbf{Đáp án của câu hỏi là A}
}
\end{ex}

\begin{ex} %[3D2V78]
Một xưởng dệt mỗi ngày dệt được $14\,500\text{ m}$ vải. Hỏi trong 6 ngày xưởng dệt được bao nhiêu mét vải?
\choice
{$85\,000\text{ m}$}
{$86\,000\text{ m}$}
{\True $87\,000\text{ m}$}
{$88\,000\text{ m}$}
\loigiai{
$14\,500 \times 6 = 87\,000\text{ (m)}$.
\textbf{Đáp án của câu hỏi là C}
}
\end{ex}

\textbf{Phần 2. Trắc nghiệm đúng sai}
\setcounterpageref{ex}{0}

\begin{ex} %[3D2V78]
Xét tính đúng sai của các khẳng định về phép nhân chia:
\choiceTF
{\True Nhân một số với 1 luôn bằng chính nó.}
{Một số chia cho 1 thì bằng 0.}
{\True Trong biểu thức chỉ có nhân và chia, thực hiện từ trái sang phải.}
{$25\,000 \times 4 = 10\,000$.}
\loigiai{
\begin{itemize}
    \item Ý a) Đúng: Tính chất của 1.
    \item Ý b) Sai: Sai vì chia cho 1 bằng chính nó.
    \item Ý c) Đúng: Thứ tự tính.
    \item Ý d) Sai: Sai vì $25\,000 \times 4 = 100\,000$.
\end{itemize}
\begin{itemize}
    \item a) \textbf{Đúng}.
    \item b) \textbf{Sai}.
    \item c) \textbf{Đúng}.
    \item d) \textbf{Sai}.
\end{itemize}
}
\end{ex}

\begin{ex} %[3D2V78]
Một nhà máy có 4 dây chuyền sản xuất, mỗi dây chuyền làm được 18 000 sản phẩm mỗi tháng:
\choiceTF
{\True Cả 4 dây chuyền làm được 72 000 sản phẩm mỗi tháng.}
{\True Hai dây chuyền làm được 36 000 sản phẩm mỗi tháng.}
{Cả nhà máy làm trong 2 tháng được 150 000 sản phẩm.}
{\True Nếu đóng 6 sản phẩm vào 1 hộp thì mỗi tháng đóng được 12 000 hộp.}
\loigiai{
\begin{itemize}
    \item Ý a) Đúng: $18\,000 \times 4 = 72\,000$ sản phẩm.
    \item Ý b) Đúng: $18\,000 \times 2 = 36\,000$ sản phẩm.
    \item Ý c) Sai: Sai vì $72\,000 \times 2 = 144\,000 < 150\,000$.
    \item Ý d) Đúng: $72\,000 : 6 = 12\,000$ hộp.
\end{itemize}
\begin{itemize}
    \item a) \textbf{Đúng}.
    \item b) \textbf{Đúng}.
    \item c) \textbf{Sai}.
    \item d) \textbf{Đúng}.
\end{itemize}
}
\end{ex}

\textbf{Phần 3. Trả lời ngắn}
\setcounterpageref{ex}{0}

\begin{ex} %[3D2V78]
Tính: $22\,500 \times 4$.
\shortans{90000}
\loigiai{
\textbf{Đáp số:} $90000$\\
$22\,500 \times 4 = 90\,000$.
}
\end{ex}

\begin{ex} %[3D2V78]
Tính: $75\,500 : 5$.
\shortans{15100}
\loigiai{
\textbf{Đáp số:} $15100$\\
$75\,500 : 5 = 15\,100$.
}
\end{ex}

\begin{ex} %[3D2V78]
Tìm $x$, biết: $x \times 6 = 96\,000$.
\shortans{16000}
\loigiai{
\textbf{Đáp số:} $16000$\\
$x = 96\,000 : 6 = 16\,000$.
}
\end{ex}

\begin{ex} %[3D2V78]
Tính giá trị biểu thức: $15\,000 \times 5 - 25\,000$.
\shortans{50000}
\loigiai{
\textbf{Đáp số:} $50000$\\
$75\,000 - 25\,000 = 50\,000$.
}
\end{ex}

\begin{ex} %[3D2V78]
Một hình vuông có cạnh dài $18\,000\text{ m}$. Tính chu vi hình vuông theo mét.
\shortans{72000}
\loigiai{
\textbf{Đáp số:} $72000$\\
$18\,000 \times 4 = 72\,000\text{ (m)}$.
}
\end{ex}

\begin{ex} %[3D2V78]
Có 5 xe chở hàng, mỗi xe chở $14\,000\text{ kg}$ gạo. Hỏi cả đoàn xe chở bao nhiêu ki-lô-gam gạo?
\shortans{70000}
\loigiai{
\textbf{Đáp số:} $70000$\\
$14\,000 \times 5 = 70\,000\text{ (kg)}$.
}
\end{ex}

\newpage

\TaoPhieuKiemTra{Chủ đề 16 -- Ôn tập cuối năm}{79}{Ôn tập hình học và đo lường}{https://loptoanthaylinh.vn/thi-online/toan-3/bai-79}
\textbf{Phần 1. Trắc nghiệm lựa chọn}
\setcounterpageref{ex}{0}

\begin{ex} %[3H3V79]
Một hình chữ nhật có chiều dài $24\text{ cm}$, chiều rộng $8\text{ cm}$. Chu vi hình chữ nhật đó là
\choice
{\True $64\text{ cm}$}
{$32\text{ cm}$}
{$192\text{ cm}$}
{$48\text{ cm}$}
\loigiai{
Chu vi: $(24 + 8) \times 2 = 64\text{ (cm)}$.
\textbf{Đáp án của câu hỏi là A}
}
\end{ex}

\begin{ex} %[3H3V79]
Một hình vuông có chu vi $48\text{ cm}$. Diện tích của hình vuông đó là
\choice
{$96\text{ cm}^2$}
{$128\text{ cm}^2$}
{\True $144\text{ cm}^2$}
{$120\text{ cm}^2$}
\loigiai{
Cạnh: $48 : 4 = 12\text{ cm}$. Diện tích: $12 \times 12 = 144\text{ (cm}^2)$.
\textbf{Đáp án của câu hỏi là C}
}
\end{ex}

\begin{ex} %[3H3V79]
Một hình tròn có đường kính $26\text{ cm}$. Bán kính của hình tròn đó là
\choice
{$52\text{ cm}$}
{\True $13\text{ cm}$}
{$26\text{ cm}$}
{$18\text{ cm}$}
\loigiai{
Bán kính: $26 : 2 = 13\text{ (cm)}$.
\textbf{Đáp án của câu hỏi là B}
}
\end{ex}

\begin{ex} %[3D3V79]
Đổi $3\text{ m } 5\text{ mm}$ ra đơn vị mi-li-mét:
\choice
{$350\text{ mm}$}
{$3050\text{ mm}$}
{$3500\text{ mm}$}
{\True $3005\text{ mm}$}
\loigiai{
$3\text{ m} = 3000\text{ mm}$. Tổng: $3000 + 5 = 3005\text{ mm}$.
\textbf{Đáp án của câu hỏi là D}
}
\end{ex}

\begin{ex} %[3D3V79]
Đổi $2\text{ kg } 80\text{ g}$ ra đơn vị gam:
\choice
{\True $2080\text{ g}$}
{$2800\text{ g}$}
{$280\text{ g}$}
{$2008\text{ g}$}
\loigiai{
$2\text{ kg} = 2000\text{ g}$. Tổng: $2000 + 80 = 2080\text{ g}$.
\textbf{Đáp án của câu hỏi là A}
}
\end{ex}

\begin{ex} %[3D3V79]
Một bình chứa $2\ell$ nước. Rót ra 5 cốc nước, mỗi cốc $250\text{ ml}$. Trong bình còn lại bao nhiêu mi-li-lít nước?
\choice
{$1250\text{ ml}$}
{$500\text{ ml}$}
{\True $750\text{ ml}$}
{$1000\text{ ml}$}
\loigiai{
Đổi $2\ell = 2000\text{ ml}$. Đã rót: $250 \times 5 = 1250\text{ ml}$. Còn lại: $2000 - 1250 = 750\text{ (ml)}$.
\textbf{Đáp án của câu hỏi là C}
}
\end{ex}

\begin{ex} %[3H3V79]
Một khối lập phương có bao nhiêu đỉnh, bao nhiêu cạnh và bao nhiêu mặt?
\choice
{$6$ đỉnh, $12$ cạnh, $8$ mặt}
{\True $8$ đỉnh, $12$ cạnh, $6$ mặt}
{$8$ đỉnh, $6$ cạnh, $12$ mặt}
{$12$ đỉnh, $8$ cạnh, $6$ mặt}
\loigiai{
Khối lập phương có 8 đỉnh, 12 cạnh, 6 mặt.
\textbf{Đáp án của câu hỏi là B}
}
\end{ex}

\begin{ex} %[3D3V79]
Nhiệt độ nước đá đang tan là
\choice
{$100^\circ\text{C}$}
{$37^\circ\text{C}$}
{$10^\circ\text{C}$}
{\True $0^\circ\text{C}$}
\loigiai{
Nước đá tan ở $0^\circ\text{C}$.
\textbf{Đáp án của câu hỏi là D}
}
\end{ex}

\begin{ex} %[3H3V79]
Cho đoạn thẳng $AB = 84\text{ cm}$. Gọi $M$ là trung điểm của $AB$, $N$ là trung điểm của $AM$. Độ dài đoạn thẳng $NB$ là
\choice
{\True $63\text{ cm}$}
{$42\text{ cm}$}
{$21\text{ cm}$}
{$54\text{ cm}$}
\loigiai{
$AM = 42\text{ cm}$; $AN = 21\text{ cm}$; $NB = 84 - 21 = 63\text{ (cm)}$.
\textbf{Đáp án của câu hỏi là A}
}
\end{ex}

\begin{ex} %[3H3V79]
Góc tạo bởi kim giờ và kim phút của đồng hồ lúc đúng 9 giờ là
\choice
{Góc nhọn}
{Góc tù}
{\True Góc vuông}
{Góc bẹt}
\loigiai{
Lúc 9 giờ đúng hai kim vuông góc với nhau.
\textbf{Đáp án của câu hỏi là C}
}
\end{ex}

\begin{ex} %[3H3V79]
Một hình chữ nhật có diện tích $80\text{ cm}^2$, chiều rộng $8\text{ cm}$. Chu vi hình chữ nhật đó là
\choice
{$40\text{ cm}$}
{\True $36\text{ cm}$}
{$32\text{ cm}$}
{$28\text{ cm}$}
\loigiai{
Chiều dài: $80 : 8 = 10\text{ cm}$. Chu vi: $(10 + 8) \times 2 = 36\text{ (cm)}$.
\textbf{Đáp án của câu hỏi là B}
}
\end{ex}

\begin{ex} %[3H3V79]
Cần bao nhiêu khối lập phương nhỏ cạnh $1\text{ cm}$ để xếp thành khối lập phương lớn cạnh $4\text{ cm}$?
\choice
{$16$ khối}
{$32$ khối}
{$48$ khối}
{\True $64$ khối}
\loigiai{
$4 \times 4 \times 4 = 64$ khối.
\textbf{Đáp án của câu hỏi là D}
}
\end{ex}

\textbf{Phần 2. Trắc nghiệm đúng sai}
\setcounterpageref{ex}{0}

\begin{ex} %[3H3V79]
Xét tính đúng sai của các kiến thức hình học và đo lường:
\choiceTF
{\True Hình vuông có 4 cạnh bằng nhau và 4 góc vuông.}
{$1\text{ kg} = 100\text{ g}$.}
{\True Đường kính của hình tròn gấp đôi bán kính.}
{Diện tích hình chữ nhật bằng (dài + rộng) nhân 2.}
\loigiai{
\begin{itemize}
    \item Ý a) Đúng: Định nghĩa chuẩn.
    \item Ý b) Sai: Sai vì $1\text{ kg} = 1000\text{ g}$.
    \item Ý c) Đúng: Mối quan hệ chuẩn.
    \item Ý d) Sai: Sai vì đó là chu vi.
\end{itemize}
\begin{itemize}
    \item a) \textbf{Đúng}.
    \item b) \textbf{Sai}.
    \item c) \textbf{Đúng}.
    \item d) \textbf{Sai}.
\end{itemize}
}
\end{ex}

\begin{ex} %[3H3V79]
Một khu đất hình chữ nhật có chiều dài $50\text{ m}$, chiều rộng $30\text{ m}$:
\choiceTF
{\True Chu vi khu đất là $160\text{ m}$.}
{\True Diện tích khu đất là $1500\text{ m}^2$.}
{Nếu bớt chiều dài đi $20\text{ m}$ thì khu đất có diện tích không đổi.}
{\True Xây hàng rào xung quanh khu đất chừa lối đi $4\text{ m}$ cần $156\text{ m}$ rào.}
\loigiai{
\begin{itemize}
    \item Ý a) Đúng: $(50 + 30) \times 2 = 160\text{ m}$.
    \item Ý b) Đúng: $50 \times 30 = 1500\text{ m}^2$.
    \item Ý c) Sai: Sai vì diện tích giảm xuống còn $30 \times 30 = 900\text{ m}^2$.
    \item Ý d) Đúng: $160 - 4 = 156\text{ m}$.
\end{itemize}
\begin{itemize}
    \item a) \textbf{Đúng}.
    \item b) \textbf{Đúng}.
    \item c) \textbf{Sai}.
    \item d) \textbf{Đúng}.
\end{itemize}
}
\end{ex}

\textbf{Phần 3. Trả lời ngắn}
\setcounterpageref{ex}{0}

\begin{ex} %[3H3V79]
Một hình vuông có chu vi $60\text{ cm}$. Tính diện tích hình vuông theo xăng-ti-mét vuông.
\shortans{225}
\loigiai{
\textbf{Đáp số:} $225$\\
Cạnh: $60 : 4 = 15\text{ cm}$. Diện tích: $15 \times 15 = 225\text{ (cm}^2)$.
}
\end{ex}

\begin{ex} %[3H3V79]
Hình tròn có bán kính $16\text{ cm}$. Tính đường kính hình tròn theo xăng-ti-mét.
\shortans{32}
\loigiai{
\textbf{Đáp số:} $32$\\
$16 \times 2 = 32\text{ (cm)}$.
}
\end{ex}

\begin{ex} %[3D3V79]
Đổi $5\text{ m } 8\text{ cm}$ ra mi-li-mét.
\shortans{5080}
\loigiai{
\textbf{Đáp số:} $5080$\\
$5000 + 80 = 5080\text{ mm}$.
}
\end{ex}

\begin{ex} %[3D3V79]
Một túi muối nặng $1\text{ kg}$. Dùng hết $350\text{ g}$. Túi muối còn lại bao nhiêu gam?
\shortans{650}
\loigiai{
\textbf{Đáp số:} $650$\\
$1000 - 350 = 650\text{ (g)}$.
}
\end{ex}

\begin{ex} %[3D3V79]
Đổi $3\ell - 850\text{ ml}$ ra mi-li-lít.
\shortans{2150}
\loigiai{
\textbf{Đáp số:} $2150$\\
$3000 - 850 = 2150\text{ (ml)}$.
}
\end{ex}

\begin{ex} %[3H3V79]
Một khối lập phương có chu vi 1 mặt là $36\text{ cm}$. Tính tổng độ dài 12 cạnh của khối lập phương theo xăng-ti-mét.
\shortans{108}
\loigiai{
\textbf{Đáp số:} $108$\\
Cạnh: $36 : 4 = 9\text{ cm}$. Tổng 12 cạnh: $12 \times 9 = 108\text{ (cm)}$.
}
\end{ex}

\newpage

\TaoPhieuKiemTra{Chủ đề 16 -- Ôn tập cuối năm}{80}{Ôn tập bảng số liệu và khả năng xảy ra}{https://loptoanthaylinh.vn/thi-online/toan-3/bai-80}
\textbf{Phần 1. Trắc nghiệm lựa chọn}
\setcounterpageref{ex}{0}

\begin{ex} %[3X1V80]
Bảng thống kê số bạn thích các môn thể thao: Bóng đá: 25; Bơi lội: 18; Cầu lông: 22; Cờ vua: 15. Môn thể thao được nhiều bạn thích nhất là
\choice
{Bơi lội}
{Cầu lông}
{Cờ vua}
{\True Bóng đá}
\loigiai{
Bóng đá có 25 bạn thích, nhiều nhất.
\textbf{Đáp án của câu hỏi là D}
}
\end{ex}

\begin{ex} %[3X1V80]
Tổng số bạn tham gia khảo sát ở câu trên là
\choice
{$75$ bạn}
{\True $80$ bạn}
{$85$ bạn}
{$70$ bạn}
\loigiai{
Tổng: $25 + 18 + 22 + 15 = 80$ bạn.
\textbf{Đáp án của câu hỏi là B}
}
\end{ex}

\begin{ex} %[3X1V80]
Số bạn thích Bóng đá nhiều hơn số bạn thích Cờ vua bao nhiêu bạn?
\choice
{$8$ bạn}
{$12$ bạn}
{\True $10$ bạn}
{$7$ bạn}
\loigiai{
$25 - 15 = 10$ bạn.
\textbf{Đáp án của câu hỏi là C}
}
\end{ex}

\begin{ex} %[3X1V80]
Trong hộp chỉ có 6 chiếc bút màu xanh. Lấy ra ngẫu nhiên 1 chiếc bút. Sự kiện nào là 'Chắc chắn' xảy ra?
\choice
{\True Lấy được bút màu xanh}
{Lấy được bút màu đỏ}
{Lấy được bút màu đen}
{Không lấy được bút nào}
\loigiai{
Trong hộp chỉ có bút xanh nên chắc chắn lấy được bút xanh.
\textbf{Đáp án của câu hỏi là A}
}
\end{ex}

\begin{ex} %[3X1V80]
Gieo 1 con xúc xắc chuẩn. Sự kiện nào là 'Không thể' xảy ra?
\choice
{Xuất hiện mặt 1 chấm}
{Xuất hiện mặt 6 chấm}
{Xuất hiện mặt 3 chấm}
{\True Xuất hiện mặt 0 chấm}
\loigiai{
Xúc xắc chỉ có từ 1 đến 6 chấm nên không thể có mặt 0 chấm.
\textbf{Đáp án của câu hỏi là D}
}
\end{ex}

\begin{ex} %[3X1V80]
Trong túi có 4 kẹo cam và 4 kẹo dâu. Lấy ra 1 kẹo. Khả năng lấy được kẹo cam là
\choice
{Chắc chắn xảy ra}
{\True Có thể xảy ra}
{Không thể xảy ra}
{Luôn luôn xảy ra}
\loigiai{
Có cả 2 loại kẹo nên có thể lấy được kẹo cam.
\textbf{Đáp án của câu hỏi là B}
}
\end{ex}

\begin{ex} %[3X1V80]
Thống kê điểm kiểm tra Toán của tổ 1: có 3 điểm 10, 4 điểm 9, 2 điểm 8, 1 điểm 7. Số học sinh của tổ 1 là
\choice
{$9$ bạn}
{$11$ bạn}
{\True $10$ bạn}
{$8$ bạn}
\loigiai{
Tổng số bạn: $3 + 4 + 2 + 1 = 10$ bạn.
\textbf{Đáp án của câu hỏi là C}
}
\end{ex}

\begin{ex} %[3X1V80]
Số điểm từ 9 trở lên trong tổ 1 là bao nhiêu điểm?
\choice
{\True $7$ điểm}
{$6$ điểm}
{$8$ điểm}
{$9$ điểm}
\loigiai{
Điểm 9 và 10 gồm: $3 + 4 = 7$ điểm.
\textbf{Đáp án của câu hỏi là A}
}
\end{ex}

\begin{ex} %[3X1V80]
Một hộp có 10 viên bi xanh và 2 viên bi đỏ. Cần lấy ra ít nhất bao nhiêu viên bi để CHẮC CHẮN có ít nhất 1 viên bi màu đỏ?
\choice
{$3$ viên bi}
{$10$ viên bi}
{$12$ viên bi}
{\True $11$ viên bi}
\loigiai{
Trường hợp xấu nhất lấy cả 10 bi xanh, cần lấy thêm 1 viên nữa: $10 + 1 = 11$ viên.
\textbf{Đáp án của câu hỏi là D}
}
\end{ex}

\begin{ex} %[3X1V80]
Ngày mai có tuyết rơi giữa mùa hè tại Hà Nội là sự kiện
\choice
{Chắc chắn xảy ra}
{\True Không thể xảy ra (hoặc cực kì khó xảy ra)}
{Thường xuyên xảy ra}
{Luôn xảy ra}
\loigiai{
Khí hậu Hà Nội mùa hè nắng nóng, không thể có tuyết rơi.
\textbf{Đáp án của câu hỏi là B}
}
\end{ex}

\begin{ex} %[3X1V80]
Bảng theo dõi số sách đọc được trong hè: An: 12 cuốn; Bình: 15 cuốn; Chi: 18 cuốn. Trung bình mỗi bạn đọc được bao nhiêu cuốn sách?
\choice
{$14$ cuốn}
{$16$ cuốn}
{\True $15$ cuốn}
{$18$ cuốn}
\loigiai{
$(12 + 15 + 18) : 3 = 45 : 3 = 15$ cuốn.
\textbf{Đáp án của câu hỏi là C}
}
\end{ex}

\begin{ex} %[3X1V80]
Tung một đồng xu 30 lần, mặt Sấp xuất hiện 14 lần. Mặt Ngửa xuất hiện bao nhiêu lần?
\choice
{\True $16$ lần}
{$14$ lần}
{$15$ lần}
{$17$ lần}
\loigiai{
$30 - 14 = 16$ lần.
\textbf{Đáp án của câu hỏi là A}
}
\end{ex}

\textbf{Phần 2. Trắc nghiệm đúng sai}
\setcounterpageref{ex}{0}

\begin{ex} %[3X1V80]
Xét tính đúng sai của việc phân tích bảng số liệu và xác suất:
\choiceTF
{\True Bảng số liệu giúp tổng hợp thông tin rõ ràng.}
{Một sự kiện chắc chắn xảy ra thì có lúc không xảy ra.}
{\True Trong trò chơi gieo xúc xắc, mặt xuất hiện số chẵn là sự kiện có thể xảy ra.}
{Lấy 1 viên bi từ hộp chỉ có bi đỏ thì có thể lấy được bi xanh.}
\loigiai{
\begin{itemize}
    \item Ý a) Đúng: Lợi ích bảng thống kê.
    \item Ý b) Sai: Sai vì chắc chắn xảy ra thì luôn luôn xảy ra.
    \item Ý c) Đúng: Có các mặt 2, 4, 6.
    \item Ý d) Sai: Sai vì không thể lấy được bi xanh.
\end{itemize}
\begin{itemize}
    \item a) \textbf{Đúng}.
    \item b) \textbf{Sai}.
    \item c) \textbf{Đúng}.
    \item d) \textbf{Sai}.
\end{itemize}
}
\end{ex}

\begin{ex} %[3X1V80]
Bảng thống kê số cây trồng được của 3 lớp: 3A trồng 45 cây, 3B trồng 55 cây, 3C trồng 50 cây:
\choiceTF
{\True Lớp 3B trồng được nhiều cây nhất.}
{\True Lớp 3A trồng được ít cây nhất.}
{\True Cả 3 lớp trồng được tất cả 150 cây.}
{Lớp 3B trồng gấp đôi lớp 3A.}
\loigiai{
\begin{itemize}
    \item Ý a) Đúng: 55 cây nhiều nhất.
    \item Ý b) Đúng: 45 cây ít nhất.
    \item Ý c) Đúng: $45 + 55 + 50 = 150$ cây.
    \item Ý d) Sai: Sai vì $45 \times 2 = 90 \neq 55$.
\end{itemize}
\begin{itemize}
    \item a) \textbf{Đúng}.
    \item b) \textbf{Đúng}.
    \item c) \textbf{Đúng}.
    \item d) \textbf{Sai}.
\end{itemize}
}
\end{ex}

\textbf{Phần 3. Trả lời ngắn}
\setcounterpageref{ex}{0}

\begin{ex} %[3X1V80]
Một hộp có 8 bi xanh và 4 bi đỏ. Cần lấy ra ít nhất bao nhiêu viên bi để chắc chắn có ít nhất 1 viên bi màu đỏ?
\shortans{9}
\loigiai{
\textbf{Đáp số:} $9$\\
Trường hợp xấu nhất lấy 8 bi xanh, lấy thêm 1: $8 + 1 = 9$ viên.
}
\end{ex}

\begin{ex} %[3X1V80]
Tổng số học sinh tham gia 3 câu lạc bộ là 90 bạn: vẽ có 30 bạn, võ có 35 bạn. Hỏi câu lạc bộ nhạc có bao nhiêu bạn?
\shortans{25}
\loigiai{
\textbf{Đáp số:} $25$\\
$90 - (30 + 35) = 25$ bạn.
}
\end{ex}

\begin{ex} %[3X1V80]
Gieo đồng xu 40 lần, mặt ngửa xuất hiện 18 lần. Mặt sấp xuất hiện bao nhiêu lần?
\shortans{22}
\loigiai{
\textbf{Đáp số:} $22$\\
$40 - 18 = 22$ lần.
}
\end{ex}

\begin{ex} %[3X1V80]
Số điểm 10 của tổ 1 là 28 điểm, tổ 2 là 32 điểm. Cả hai tổ có tất cả bao nhiêu điểm 10?
\shortans{60}
\loigiai{
\textbf{Đáp số:} $60$\\
$28 + 32 = 60$ điểm.
}
\end{ex}

\begin{ex} %[3X1V80]
Có bao nhiêu mặt của con xúc xắc chuẩn ghi số chấm bé hơn 4?
\shortans{3}
\loigiai{
\textbf{Đáp số:} $3$\\
Mặt 1, 2, 3 (có 3 mặt).
}
\end{ex}

\begin{ex} %[3X1V80]
Số cây trồng của 4 bạn: An: 5 cây; Bình: 7 cây; Cúc: 6 cây; Dũng: 6 cây. Tính tổng số cây của cả 4 bạn.
\shortans{24}
\loigiai{
\textbf{Đáp số:} $24$\\
$5 + 7 + 6 + 6 = 24$ cây.
}
\end{ex}

\newpage

\TaoPhieuKiemTra{Chủ đề 16 -- Ôn tập cuối năm}{81}{Ôn tập chung cuối năm}{https://loptoanthaylinh.vn/thi-online/toan-3/bai-81}
\textbf{Phần 1. Trắc nghiệm lựa chọn}
\setcounterpageref{ex}{0}

\begin{ex} %[3X3V81]
Tính giá trị của biểu thức: $A = 25\,000 \times 3 - 48\,000 : 6$.
\choice
{$65\,000$}
{$70\,000$}
{\True $67\,000$}
{$66\,000$}
\loigiai{
$25\,000 \times 3 = 75\,000$; $48\,000 : 6 = 8000$; $75\,000 - 8000 = 67\,000$.
\textbf{Đáp án của câu hỏi là C}
}
\end{ex}

\begin{ex} %[3X3V81]
Số liền sau của số lớn nhất có năm chữ số là
\choice
{\True $100\,000$}
{$99\,998$}
{$100\,001$}
{$99\,990$}
\loigiai{
$99\,999 + 1 = 100\,000$.
\textbf{Đáp án của câu hỏi là A}
}
\end{ex}

\begin{ex} %[3X3V81]
Một hình chữ nhật có chiều dài $30\text{ cm}$, chiều rộng bằng $\dfrac{1}{3}$ chiều dài. Diện tích của hình chữ nhật đó là
\choice
{$100\text{ cm}^2$}
{$200\text{ cm}^2$}
{$400\text{ cm}^2$}
{\True $300\text{ cm}^2$}
\loigiai{
Chiều rộng: $30 : 3 = 10\text{ cm}$. Diện tích: $30 \times 10 = 300\text{ (cm}^2)$.
\textbf{Đáp án của câu hỏi là D}
}
\end{ex}

\begin{ex} %[3X3V81]
Một hình vuông có chu vi $64\text{ cm}$. Diện tích của hình vuông đó là
\choice
{$128\text{ cm}^2$}
{\True $256\text{ cm}^2$}
{$64\text{ cm}^2$}
{$512\text{ cm}^2$}
\loigiai{
Cạnh hình vuông: $64 : 4 = 16\text{ cm}$. Diện tích: $16 \times 16 = 256\text{ (cm}^2)$.
\textbf{Đáp án của câu hỏi là B}
}
\end{ex}

\begin{ex} %[3X3V81]
Số La Mã $\text{XXI}$ cộng với $\text{IV}$ được kết quả là số La Mã nào?
\choice
{$\text{XXIV}$}
{$\text{XXVI}$}
{\True $\text{XXV}$}
{$\text{XX}$}
\loigiai{
$\text{XXI} = 21$; $\text{IV} = 4$. $21 + 4 = 25 = \text{XXV}$.
\textbf{Đáp án của câu hỏi là C}
}
\end{ex}

\begin{ex} %[3X3V81]
Một cửa hàng có $80\,000\text{ kg}$ gạo. Ngày đầu bán $\dfrac{1}{4}$ số gạo, ngày thứ hai bán $\dfrac{1}{3}$ số gạo còn lại. Cửa hàng còn lại bao nhiêu ki-lô-gam gạo?
\choice
{\True $40\,000\text{ kg}$}
{$20\,000\text{ kg}$}
{$50\,000\text{ kg}$}
{$30\,000\text{ kg}$}
\loigiai{
Ngày 1 bán: $80\,000 : 4 = 20\,000\text{ kg}$. Còn lại: $80\,000 - 20\,000 = 60\,000\text{ kg}$. Ngày 2 bán: $60\,000 : 3 = 20\,000\text{ kg}$. Còn lại: $60\,000 - 20\,000 = 40\,000\text{ (kg)}$.
\textbf{Đáp án của câu hỏi là A}
}
\end{ex}

\begin{ex} %[3X3V81]
Làm tròn số $78\,650$ đến hàng nghìn ta được số
\choice
{$78\,000$}
{$80\,000$}
{$78\,700$}
{\True $79\,000$}
\loigiai{
Chữ số hàng trăm là 6 ($6 \ge 5$) nên làm tròn lên $79\,000$.
\textbf{Đáp án của câu hỏi là D}
}
\end{ex}

\begin{ex} %[3X3V81]
Một hình tròn có bán kính $9\text{ cm}$. Đoạn thẳng dài nhất nối 2 điểm trên hình tròn có độ dài là
\choice
{$9\text{ cm}$}
{\True $18\text{ cm}$}
{$27\text{ cm}$}
{$36\text{ cm}$}
\loigiai{
Đoạn dài nhất chính là đường kính: $9 \times 2 = 18\text{ (cm)}$.
\textbf{Đáp án của câu hỏi là B}
}
\end{ex}

\begin{ex} %[3X3V81]
Đổi $5\text{ m } 4\text{ cm}$ ra đơn vị mi-li-mét:
\choice
{$540\text{ mm}$}
{$5004\text{ mm}$}
{\True $5040\text{ mm}$}
{$5400\text{ mm}$}
\loigiai{
$5\text{ m} = 5000\text{ mm}$; $4\text{ cm} = 40\text{ mm}$. Tổng: $5040\text{ mm}$.
\textbf{Đáp án của câu hỏi là C}
}
\end{ex}

\begin{ex} %[3X3V81]
Tháng nào sau đây có đúng 30 ngày?
\choice
{\True Tháng 4}
{Tháng 1}
{Tháng 3}
{Tháng 5}
\loigiai{
Tháng 4 có đúng 30 ngày.
\textbf{Đáp án của câu hỏi là A}
}
\end{ex}

\begin{ex} %[3X3V81]
Một con xúc xắc gieo ngẫu nhiên. Sự kiện 'Xuất hiện mặt có số chấm bé hơn 7' là sự kiện
\choice
{Có thể xảy ra}
{Không thể xảy ra}
{Hiếm khi xảy ra}
{\True Chắc chắn xảy ra}
\loigiai{
Tất cả các mặt từ 1 đến 6 đều bé hơn 7 nên chắc chắn xảy ra.
\textbf{Đáp án của câu hỏi là D}
}
\end{ex}

\begin{ex} %[3X3V81]
Tổng của hai số là 90 000, số lớn gấp 4 lần số bé. Số lớn là
\choice
{$18\,000$}
{\True $72\,000$}
{$60\,000$}
{$75\,000$}
\loigiai{
Tổng số phần: $4 + 1 = 5$ phần. Số bé: $90\,000 : 5 = 18\,000$. Số lớn: $18\,000 \times 4 = 72\,000$.
\textbf{Đáp án của câu hỏi là B}
}
\end{ex}

\textbf{Phần 2. Trắc nghiệm đúng sai}
\setcounterpageref{ex}{0}

\begin{ex} %[3X3V81]
Xét tính đúng sai của các kiến thức tổng hợp toàn bộ chương trình Toán 3:
\choiceTF
{\True Số 100 000 là số tròn trăm nghìn bé nhất có sáu chữ số.}
{$1\text{ m} = 100\text{ mm}$.}
{\True Hình vuông có 4 cạnh bằng nhau và 4 góc vuông.}
{Trong phép chia cho 5, số dư lớn nhất có thể là 5.}
\loigiai{
\begin{itemize}
    \item Ý a) Đúng: Định nghĩa.
    \item Ý b) Sai: Sai vì $1\text{ m} = 1000\text{ mm}$.
    \item Ý c) Đúng: Đặc điểm hình vuông.
    \item Ý d) Sai: Sai vì số dư lớn nhất là 4.
\end{itemize}
\begin{itemize}
    \item a) \textbf{Đúng}.
    \item b) \textbf{Sai}.
    \item c) \textbf{Đúng}.
    \item d) \textbf{Sai}.
\end{itemize}
}
\end{ex}

\begin{ex} %[3X3V81]
Một hình chữ nhật có chiều dài $20\text{ cm}$, chiều rộng $10\text{ cm}$:
\choiceTF
{\True Chu vi hình chữ nhật là $60\text{ cm}$.}
{\True Diện tích hình chữ nhật là $200\text{ cm}^2$.}
{Nếu chia đôi hình chữ nhật theo chiều dài thì được 2 hình vuông cạnh $5\text{ cm}$.}
{\True Chiều dài gấp đôi chiều rộng.}
\loigiai{
\begin{itemize}
    \item Ý a) Đúng: $(20 + 10) \times 2 = 60\text{ cm}$.
    \item Ý b) Đúng: $20 \times 10 = 200\text{ cm}^2$.
    \item Ý c) Sai: Sai vì cắt đôi được 2 hình vuông cạnh $10\text{ cm}$.
    \item Ý d) Đúng: $20 : 10 = 2$ lần.
\end{itemize}
\begin{itemize}
    \item a) \textbf{Đúng}.
    \item b) \textbf{Đúng}.
    \item c) \textbf{Sai}.
    \item d) \textbf{Đúng}.
\end{itemize}
}
\end{ex}

\textbf{Phần 3. Trả lời ngắn}
\setcounterpageref{ex}{0}

\begin{ex} %[3X3V81]
Tính giá trị biểu thức: $15\,000 \times 4 + 40\,000 : 8$.
\shortans{65000}
\loigiai{
\textbf{Đáp số:} $65000$\\
$60\,000 + 5000 = 65\,000$.
}
\end{ex}

\begin{ex} %[3X3V81]
Tìm số liền sau của số 99 999.
\shortans{100000}
\loigiai{
\textbf{Đáp số:} $100000$\\
$99\,999 + 1 = 100\,000$.
}
\end{ex}

\begin{ex} %[3X3V81]
Một hình vuông có chu vi $72\text{ cm}$. Tính độ dài cạnh hình vuông theo xăng-ti-mét.
\shortans{18}
\loigiai{
\textbf{Đáp số:} $18$\\
$72 : 4 = 18\text{ (cm)}$.
}
\end{ex}

\begin{ex} %[3X3V81]
Tính diện tích hình chữ nhật có chiều dài $15\text{ cm}$ và chiều rộng $8\text{ cm}$.
\shortans{120}
\loigiai{
\textbf{Đáp số:} $120$\\
$15 \times 8 = 120\text{ (cm}^2)$.
}
\end{ex}

\begin{ex} %[3X3V81]
Đổi $2\ell - 750\text{ ml}$ ra mi-li-lít.
\shortans{1250}
\loigiai{
\textbf{Đáp số:} $1250$\\
$2000 - 750 = 1250\text{ (ml)}$.
}
\end{ex}

\begin{ex} %[3X3V81]
Tổng hai số là 80 000, hiệu hai số là 20 000. Tìm số lớn.
\shortans{50000}
\loigiai{
\textbf{Đáp số:} $50000$\\
Số lớn là: $(80\,000 + 20\,000) : 2 = 50\,000$.
}
\end{ex}

\newpage
\end{document}
